\documentclass[10pt]{book}
\usepackage[paperwidth=6in, paperheight=9in,
bindingoffset=5mm,%having an inner margin offset prevents text falling into the gutter
inner=20mm,
outer=20mm,
top=20mm,
bottom=20mm,
twoside,
heightrounded]{geometry}
\frenchspacing %Prevents the silly double space after each sentence
\usepackage{pifont} %Needed for \ding characters (for the crosses mainly)
\usepackage{ebgaramond} %Standard open source font

%To ensure language compatibility.
\usepackage[utf8]{inputenc}
\usepackage[greek,latin,british]{babel}

%Section Formatting
\setcounter{secnumdepth}{0} %To remove number from sections
\usepackage{sectsty} %Allows modification of sections
\sectionfont{\Large \normalfont \centering}
\subsectionfont{\normalsize \scshape \centering}
\subsubsectionfont{\normalfont \normalsize \centering \itshape}

%Alternative Section Formatting to Get Better Control
\usepackage{needspace} %\needspace causes some issue, but I think it can work if it is not initialised in a paracol but by forcing the sectioning to move (thereby moving the content).

\newcommand{\seccy}[1]{%All sections begin at the top of a page (except for the psalms), so this command is available but unused.
\needspace{8\baselineskip}
\section{#1}}

%Headings for the Propers
\newcommand{\feastday}[1]{\fancyhead[RO,LE]{\textit{#1}}}

%For drawing boxes around unapproved text:
\usepackage[most]{tcolorbox}

%Now that the Missal is being printed and, therefore, is a set text, we can establish concrete page numbers for BCP parts.

\newcommand{\SPMaryAdvent}{\pageref{SPMaryInAdvent}} %BCP 545}
\newcommand{\SPMaryChristmas}{\pageref{SPMaryPostChristmas}} %BCP 545}
\newcommand{\SPMaryEaster}{\pageref{SPMaryInEaster}} %BCP 545}
\newcommand{\SPSaints}{\pageref{SPSaints}} %BCP 546}
\newcommand{\SPAgainst}{\pageref{SPAgainst}} %BCP 546}
\newcommand{\SPChiefBishop}{\pageref{SPChiefBishop}} %BCP 547}
\newcommand{\SPHolyGhost}{\pageref{SPHolyGhost}} %BCP 547}
\newcommand{\SPLiving}{\pageref{SPLivingDeparted}} %BCP 548}
\newcommand{\Lammas}{BCP 481}


\newcommand{\psalmsec}[1]{
\needspace{9\baselineskip} %I'm still tweaking this. The psalms are unique, so this is just used for them.
\section*{#1}}

\newcommand{\subbypsalm}[1]{\setcounter{lettcount}{0}%
\needspace{4\baselineskip}%
\subsection{#1}}

\newcommand{\subby}[1]{\setcounter{lettcount}{0}\setcounter{margcount}{1}%
\needspace{8\baselineskip}%
\subsection{#1}
\global\setcounter{latinrubric}{1}}

\newcommand{\subbysub}[1]{\setcounter{lettcount}{0}%
\needspace{6\baselineskip}% I have no idea how sectioning affects the calculation of baselineskip.
\subsubsection{#1}}

%These are commands for when we want the counters reset and some provision made for spacing, but no as much needed by lettrine or other circumstances:

\newcommand{\subsec}[1]{\setcounter{lettcount}{0}\setcounter{margcount}{1}%
\needspace{2\baselineskip}%
\subsection{#1}%
\global\setcounter{latinrubric}{1}}

\newcommand{\subsubsec}[1]{\setcounter{lettcount}{0}%
\needspace{2\baselineskip}%
\subsubsection{#1}}

\newcommand{\bcpfeast}[3]{\subby{[#1]}%
%Feast Day:
\fancyhead[RO,LE]{\textit{#2}}%
%Feast Date:
\fancyhead[RE,LO]{#3}}

%For smaller subsections in the Hymnal Copyright page:
\newcommand{\hcopyrightsection}[1]{\par\needspace{1\baselineskip}
\centerline{\textbf{\small#1}}
\par}

\newcommand{\supplement}[3]{\par
%\vspace{1\baselineskip}
%\needspace{4\baselineskip}
%\centerline{[#1. #2 #3]}
\needspace{2\baselineskip}
\begin{center}
	[#1. #2 #3]
\end{center}
\feastday{#2}
%\fancyhead[RO,LE]{\textit{#2}}
\fancyhead[RE,LO]{#1}}

%URL
\usepackage[hidelinks]{hyperref}
%Remove red boxes around links:
%\hypersetup{%
%    pdfborder = {0 0 0}
%}

%Page Headers
\usepackage{fancyhdr}
\pagestyle{fancy}
\fancyhf{}
\renewcommand{\headrulewidth}{0pt} %Prevent line at top of page
\fancyfoot[C]{\thepage}
%Prevents chapter from resetting page layout
\fancypagestyle{plain}{
	\fancyhf{}
	\fancyfoot[C]{}
	\renewcommand{\headrulewidth}{0pt}
}

%Try to avoid warnings:
\setlength{\headheight}{17.13pt}


%To make sure versicle and response start with indent.
\usepackage{indentfirst}

%For hymn formatting
\usepackage{hanging}


\usepackage{paracol} %To provide balanced English-Latin dual paragraphs
\usepackage{multicol} %For non bilingual dual columns.

%Prevent lettrine misalignment:
\raggedbottom

\newcommand{\elcol}[2]{\sloppy
	\begin{paracol}{2}
		\begin{leftcolumn}
			#1
		\end{leftcolumn}
		\begin{rightcolumn}
			\selectlanguage{latin}%
			#2
		\end{rightcolumn}
		\switchcolumn*
	\end{paracol}
\selectlanguage{british}
\fussy
\setcounter{lettcount}{0}}
%Same as above but is intended for text/prayers which are neither collects nor responsory, which should not have an indent.
\newcommand{\elcoldent}[2]{\sloppy
	\begin{paracol}{2}
		\begin{leftcolumn}
		\noindent%
			#1	
		\end{leftcolumn}
		\begin{rightcolumn}
                \selectlanguage{latin}%
			\noindent
			#2
		\end{rightcolumn}
		\switchcolumn*
	\end{paracol}
\selectlanguage{british}
\fussy
\setcounter{lettcount}{0}}
%Gradual \& Alleluia
\newcommand{\gradall}[2]{\gradual{#1}\par\alleluia{#2}}%\vspace{1ex}\par\noindent
%\emph{Gradual.} #1\\
%\par\noindent
%\emph{Alleluia.} #2}

%Gradual \& Tract
\newcommand{\gradtr}[2]{\gradual{#1}\par\tract{#2}}%\vspace{1ex}\par\noindent
%\emph{Gradual.} #1\\
%\par\noindent
%\emph{Tract.} #2}

%Counter to create space from lettrine
\newcounter{prelett}
\usepackage{calc}
%Checking lett or rubric right behind
\newcounter{lettcount}
%Commands for each minor proper:
\newcommand{\alleluia}[1]{\setcounter{lettcount}{0}\vspace{1ex}\par\noindent
\emph{Alleluia.} #1}
\newcommand{\gradual}[1]{\setcounter{lettcount}{0}\vspace{1ex}\par\noindent
\emph{Gradual.} #1}
\newcommand{\tract}[1]{\setcounter{lettcount}{0}\vspace{1ex}\par\noindent
\emph{Tract.} #1}
\newcommand{\communion}[1]{\par
\ifnum\thelettcount=0
    \vspace{1\baselineskip}
    \par\noindent
    \emph{Communion.} #1\par
    \else
        \ifnum\thelettcount=1
            \ifnum\prevgraf<3
                \setcounter{prelett}{2-\prevgraf}
                \vspace{\theprelett\baselineskip}
                \medskip
                \par\noindent
                \emph{Communion.} #1\par
            \else
                \medskip
                \par\noindent
                \emph{Communion.} #1\par
            \fi
        \else
            \ifnum\thelettcount<3
                \setcounter{prelett}{2-\thelettcount}
                \vspace{\theprelett\baselineskip}
                \medskip
                \par\noindent
                \emph{Communion.} #1\par
            \else
                \medskip
                \par\noindent
                \emph{Communion.} #1\par
            \fi
        \fi
    \fi}
\newcommand{\offertory}[1]{\setcounter{lettcount}{0}\vspace{1ex}\par\noindent
\emph{Offertory.} #1}
%In order to control where the placement of the reading type and citation are placed (trying to save space by modeling after the 1928 BCP):
\newcommand{\readingcitation}[2]{\setcounter{lettcount}{0}%
\needspace{4\baselineskip}%
\subsubsection{#1. \emph{#2}}}

\newcommand{\secret}[0]{\setcounter{lettcount}{0}%
\needspace{4\baselineskip}%
\subsubsection{Secret}}

\newcommand{\introit}[0]{\setcounter{lettcount}{0}%
\needspace{4\baselineskip}%
\subsubsection{Introit}}

\newcommand{\collect}[0]{\setcounter{lettcount}{0}%
\needspace{4\baselineskip}%
\subsubsection{Collect}}

\newcommand{\postcommunion}[0]{\setcounter{lettcount}{0}%
\needspace{4\baselineskip}%
\subsubsection{Postcommunion}}

%Controls the width of the line between columns
\setlength\columnseprule{0.4pt}

%Tikhon is on the left, Gregory is on the right. `Read' indicates it is a reading, which means it needs a heading for the source.
\newcommand{\tikgregread}[4]{
\sloppy
\needspace{5\baselineskip}
\begin{paracol}{2}%
%\switchcolumn*
	\begin{leftcolumn}
		\centerline{\textit{#2}}
	\end{leftcolumn}
	\begin{rightcolumn}
		\centerline{\textit{#4}}
	\end{rightcolumn}
	\switchcolumn*
\end{paracol}
\begin{paracol}{2}%
%\switchcolumn*
	\begin{leftcolumn}
		#1
	\end{leftcolumn}
	\begin{rightcolumn}
		#3
	\end{rightcolumn}
	\switchcolumn*
\end{paracol}
\fussy
\setcounter{lettcount}{0}}

%To move Gospel and Epistle into heading without manually changing every instance late in development:

\newcommand{\tikgregreadgospel}[4]{
\sloppy
\needspace{5\baselineskip}
\begin{paracol}{2}%
%\switchcolumn*
	\begin{leftcolumn}
		\subsubsection{\textit{Gospel.} \textnormal{#2}}
	\end{leftcolumn}
	\begin{rightcolumn}
		\subsubsection{\textit{Gospel.} \textnormal{#4}}
	\end{rightcolumn}
	\switchcolumn*
\end{paracol}
\begin{paracol}{2}%
%\switchcolumn*
	\begin{leftcolumn}
		#1
	\end{leftcolumn}
	\begin{rightcolumn}
		#3
	\end{rightcolumn}
	\switchcolumn*
\end{paracol}
\fussy
\setcounter{lettcount}{0}}

\newcommand{\tikgregreadepistle}[4]{
\sloppy
\needspace{5\baselineskip}
\begin{paracol}{2}%
%\switchcolumn*
	\begin{leftcolumn}
		\subsubsection{\textit{Epistle.} \textnormal{#2}}
	\end{leftcolumn}
	\begin{rightcolumn}
		\subsubsection{\textit{Epistle.} \textnormal{#4}}
	\end{rightcolumn}
	\switchcolumn*
\end{paracol}
\begin{paracol}{2}%
%\switchcolumn*
	\begin{leftcolumn}
		#1
	\end{leftcolumn}
	\begin{rightcolumn}
		#3
	\end{rightcolumn}
	\switchcolumn*
\end{paracol}
\fussy
\setcounter{lettcount}{0}}


\newcommand{\tikgreg}[2]{\sloppy
	\begin{paracol}{2}
		\begin{leftcolumn}\noindent
			#1	
		\end{leftcolumn}
		\begin{rightcolumn}\noindent
			#2
		\end{rightcolumn}
		\switchcolumn*
	\end{paracol}
\selectlanguage{british}
\fussy
\setcounter{lettcount}{0}}

%Expedites adding dual language Oremus/Let Us Pray
\newcommand{\oremus}[0]{\par
\needspace{5\baselineskip}
	\begin{paracol}{2}
		\begin{leftcolumn}
			\selectlanguage{british}%
			\noindent%
			\centerline{\textit{Let us pray.}}
		\end{leftcolumn}
		\begin{rightcolumn}
			\selectlanguage{latin}
			\centerline{\textit{Orémus.}}
		\end{rightcolumn}
		\switchcolumn*
	\end{paracol}
\selectlanguage{british}
\setcounter{lettcount}{0}}

%Does the same as above, but for single language collects:
\newcommand{\letuspray}[0]{\par\noindent
\needspace{5\baselineskip}
\centerline{\textit{Let us pray.}}
\par
\setcounter{lettcount}{0}}

\newcommand{\oremuslatin}[0]{\par\noindent
\needspace{5\baselineskip}
\centerline{\textit{Orémus.}}
\par
\setcounter{lettcount}{0}}

\usepackage[super]{nth} %To add superscript to ordinal numeral.
%Image Configuration
\usepackage{graphicx}
\graphicspath{{resources/}}
\usepackage{wrapfig}
\usepackage{epsfig}
\usepackage{caption} %Enables text under image
%Colour
\usepackage[svgnames,dvipsnames]{xcolor}
%Hymn Pages
\usepackage{pdfpages}
%Lectionary
\usepackage{array}
\usepackage{longtable} %Allows lectionary cells to go over multiple pages
\newcolumntype{C}[1]{>{\centering\arraybackslash}p{#1}}
\newcolumntype{L}[1]{>{\RaggedRight\arraybackslash}p{#1}}
\newcommand{\evensong}[1]{Evensong} %To allow switch in naming of Vespers and Evensong in lectionary

%For Rubric Colour:
\definecolor{RubricRed}{RGB}{156, 4, 10}

%I have found it useful to have different kinds of headings in order to instruct the reader. For mere rubrical notes, which are meant to be in paragraph form, rubric is used. For headings, meant to be short, inhead (instructional heading) is used.
%Rubric Environments
\newenvironment{inhead}{\setcounter{lettcount}{0}\par\centering%
%\itshape%
\color{RubricRed}%
\small\smallskip\ignorespaces}
{\par\smallskip\ignorespacesafterend}
%The first ifnum checks the boolean (there is a lett, but no secondline or thirdline) then checks the lines preceeding:
\newenvironment{rubric}{\par\ifnum\thelettcount=1\ifnum\prevgraf=1\vspace{1\baselineskip}\smallskip\else\ifnum\prevgraf=2\medskip\smallskip\else\smallskip\fi\fi\fi\ifnum\thelettcount=2\medskip\smallskip\else\smallskip\fi\setcounter{lettcount}{0}\par\noindent%
\color{RubricRed}%
\small\ignorespaces%
%\itshape%
\P}
{\par\smallskip\ignorespacesafterend}

%Rubric after section:
\newenvironment{secrubric}{\setcounter{lettcount}{0}\par\noindent\small\ignorespaces%\itshape%%
\color{RubricRed}%
\P}
{\par\smallskip\ignorespacesafterend}

%Inline Rubric
\newcommand{\inrub}[1]{{\small{%\color{RubricRed}%
\textit{
(#1)
}
}}}

%Determine Marginnote parameters
%\usepackage{sidenotes}
\setlength{\marginparwidth}{15mm} % Set the width of the margin note
\setlength{\marginparsep}{2mm} % Set the spacing between the margin note and the text
%Fixes margin issues:
\usepackage{marginnote}
%Margin Rubrics/Notes/Description/Commentar
\newcounter{margcount}
\setcounter{margcount}{1}%
\newcommand{\margrub}[1]{\textsuperscript{\alph{margcount}}%
\marginpar{\begin{minipage}{0.7\marginparwidth}
\sloppy%
\itshape%
\scriptsize%
\textsuperscript{\alph{margcount}}%
#1%
\fussy
\end{minipage}}
\stepcounter{margcount}
\ifnum\value{margcount}>26\relax%
    \setcounter{margcount}{1}%
  \fi}
  
%Makes sure odd pages are on the right when moving between page number types.
\usepackage[strict]{changepage}
\usepackage{scrextend}

\newcommand{\margcolrub}[1]{\textsuperscript{\alph{margcount}}%
\Ifthispageodd{\reversemarginpar
\marginpar{%
\raggedright%
\sloppy%
\itshape%
\scriptsize%
    \textsuperscript{\alph{margcount}}%
    #1%
    \fussy}
\reversemarginpar}{\normalmarginpar\marginpar{%
\raggedleft%
\sloppy%
\itshape%
\scriptsize%
\textsuperscript{\alph{margcount}}%
    #1%
    \fussy}}
\stepcounter{margcount}%
\ifnum\value{margcount}>26\relax%
    \setcounter{margcount}{1}
\fi}


%Counter to add rubric indicatorlowing command  in Latin.
\newcounter{latinrubric}
\setcounter{latinrubric}{1}
%To apply counters to all paracol columns:
\globalcounter*

\newcommand{\marglatincolrub}[1]{\textsuperscript{\alph{latinrubric}}%
\stepcounter{latinrubric}%
\ifnum\value{latinrubric}>26\relax%
    \setcounter{latinrubric}{1}
\fi}
%Propers
\newcommand{\antiphon}[2]{{\centering \textit{#1 Ant.} #2}}

%For BCP Antiphons
\newcommand{\properantiphon}[2]{{\centering \textit{#1 Ant.} #2}}

%Invitatory \& Fore-Office
\newcommand{\inv}[1]{\textit{#1}}
\newcommand{\vr}[1]{{\footnotesize{ (#1)}}}

%Since I've been going back and forth on capitalisation of letters in small caps, I made this command in order to be able to change them all at once:
\newcommand{\divineName}[1]{\textsc{#1}}

\usepackage{lettrine}
\newcommand{\lett}[2]{\setcounter{lettcount}{1}\lettrine[image=true,realheight=true,lines=3]{artdeco/#1.png}{{\lowercase{#2}}}} %Specifies a specific folder for normal lettrine images
\newcommand{\lettnopic}[2]{\setcounter{lettcount}{1}\lettrine[realheight=true,lines=3]{#1}{#2}} %When there is not an image, just uses normal letter.

%These commands ensure that when there is a lettrine, following paragraphs will not conflict with the drop capital.
%We need this counter for thirdline.
\setlength{\parindent}{1em}
\newcounter{preprevgraf}

\newcommand{\secondline}[1]{\setcounter{preprevgraf}{\prevgraf}\setcounter{lettcount}{\thepreprevgraf}
\ifnum\prevgraf>2
	#1
	\par
    \addtocounter{lettcount}{\prevgraf}
	\else
		\ifnum\prevgraf=2
			\parindent=3.25em
			{#1}
			\par
            \parindent=1em
            \addtocounter{lettcount}{\prevgraf}
			\else
				\ifnum\prevgraf=1
				\noindent
				\leftskip=3.25em
				{#1}
				\par
                \addtocounter{lettcount}{\prevgraf}
				\leftskip=0em
					\else
						#1
						\par
                        \setcounter{lettcount}{\prevgraf}
				\fi
		\fi
\fi
}
%The thirdline command is to ensure proper formatting with lettrine, if it uses three lines.
\newcommand{\thirdline}[1]{
\ifnum\thelettcount=2
    \parindent=3.25em
    {#1}
    \par
    \parindent=1em
\else
    #1
    \par
\fi
\addtocounter{lettcount}{\prevgraf}
}

%Psalm Formatting
\newcommand{\psanum}[1]{{\small{#1}} } %Allows flexibility in how to handle the verse number at the beginning of each psalm verse.

%Changes to the Table of Contents
\addto\captionsbritish{\renewcommand{\contentsname}{Table of Contents}}
\setcounter{tocdepth}{1} %Chapters and Sections
\usepackage{tocloft}
\setlength{\cftbeforetoctitleskip}{-8ex}

%For the signature in the Preface
	\usepackage{ragged2e}

%For the Articles
\usepackage{enumerate}

%For the Chant Notation (Byzantine)
\usepackage[autocompile]{gregoriotex}
%Disable initial letter
\gresetheadercapture{initial-style}{gresetinitiallines}{}


%Invitatory Hymn Heading
\newcommand{\invitatoryhymn}{\needspace{4\baselineskip}
\begin{inhead}
	Invitatory Hymn
\end{inhead}}

%Office Hymn Heading
\newcommand{\officehymn}{\needspace{4\baselineskip}
\begin{inhead}
	Office Hymn
\end{inhead}}



\begin{document}
\begin{titlepage}
		\begin{center}
			{\Large{The}}
			\par
			\vspace{.4cm}
			{\Huge{\scshape Prayer Book}}
			\par
			\vspace{.4cm}
			{\Huge{\scshape Hymnal}}
			\par
			\vspace{.4cm}
			{\large{\textit{A Supplement of Hours, Hymns, \& Propers for the Book of Common Prayer}}}
			\par
			\vfill
   	        \includegraphics[scale=.17]{logo.eps}
		\end{center}
	\end{titlepage}
\cleardoublepage
\pagenumbering{gobble}
%~\clearpage
\tableofcontents
\addtocontents{toc}{\vspace*{-3ex}}%
\clearpage
%This is to ensure two odd-numbered pages are not right next to each other:
%\clearpage
%\checkoddpage
%\ifoddpage \else
%\fi
%Roman numeral page numbers for the front material.
\frontmatter
\setcounter{page}{1}
\phantomsection
\addcontentsline{toc}{chapter}{Copyright}
\fancyhead[C]{{\LARGE Copyright}}

\noindent
\copyright 2025 Apologia Anglicana, LLC\par\noindent
Edited by Augustine Watson (Augustine@ApologiaAnglicana.org), with the irreplaceable assistance of Ryan May.
The Prayer Book Hymnal, published by Apologia Anglicana, LLC, is licenced under the \href{https://creativecommons.org/licenses/by-sa/4.0/}{Creative Commons Attribution-ShareAlike 4.0 International} (\url{https://creativecommons.org/licenses/by-sa/4.0/}) licence.

\begin{center}
	\textsc{Please note that this is not currently an official text of the Antiochian Western Rite Vicariate.}
\end{center}

While great care was taken to use only works in the public domain in these United States of America, there were a few cases where this was not possible. We are thankful to have such great giants upon whose shoulders we stand. These works are used in this Book according to fair use.\par
Davis, C. Lance, ed. \textit{The Anglican Office Book}. Chester: Whithorn Press, 2023.\par\noindent
\textit{This was indispensable, for the adapted Office portions of the Feasts of Our Lady of Walsingham and St. Joseph of Arimathea.}\par
%Services for the Martyrdom of S. Charles, King and Martyr, transcribed by Mr Thomas J. W. Mason (\url{https://anglicanhistory.org/charles/texts.html})\par\noindent
%\textit{This was indispensable, for the adapted Mass portions of the Feast of Tsar St. Nicholas.}\par


And a special thanks must be given to GregoBase (\url{https://gregobase.selapa.net/}), without which the notation for the Invitatory and Office Hymns simply would not be feasible. I am especially thankful to the transcribers Albert Bloomfield, Benedictine monk, BruceL, Mateusz Ciesiółka, Dominique CROCHU, pcouderc, Pierre François, Harpa Dei, HdP, Andrew Hinkley, Nguyễn Thiện Tuấn Hoàng, Joerg Hudelmaier, Jakub Jelínek, MarkM, marteo, Br Michael, Br Michael OSB, rcarey, and Christopher Tatum.
\clearpage
\footnotesize
\section*{Devotional Hymn Copyright}
\begin{multicols}{2}

\hcopyrightsection{A Mighty Fortress is Our God}
\par\noindent
Words: Martin Luther, 1529. Translation Frederic Henry Hedge, 1853.
\par\noindent
Music: `Ein Feste Burg (Isorhythmic)' Martin Luther, 1529.  Setting: ``Common Service Book'' (ULCA), 1917, alt.
\par\noindent
Copyright: public domain. This score is a part of the Open Hymnal Project, 2013 Revision.
\par\noindent
Music source: `Common Service Book with Hymnal', ULCA 1918 Hymn 195, alt. to match syncopation of Hedge translation
\par\noindent
Lyric source: Presbyterian Hymnal, Revised, 1911 Hymn 122.

\hcopyrightsection{Alleluia, Sing to Jesus!}
\par\noindent
Words: William Chatterton Dix, 1867. 
\par\noindent
Music: `HyFrydol' Rowland H. Prichard, 1830.  Setting: ``The English Hymnal'', 1906.
\par\noindent
Arrangement: English hymnal, 1906 assumes RV Williams.
\par\noindent
Copyright: public domain. This score is a part of the Open Hymnal Project, 2006 Revision.
\par\noindent
Music source: `The English Hymnal' 1906, Hymn 563. 

\hcopyrightsection{All Creatures of Our God and King}
\par\noindent
Words: Francis of Assisi circa 1225. Translated by William H. Draper, 1919. 
\par\noindent
Music: `Lasst Uns Erfreuen' from Geistliche Kirchengesäng, Köln, 1623. Setting: Ralph Vaughan Williams, 1906.
\par\noindent
copyright: public domain. This score is a part of the Open Hymnal Project, 2005 Revision.
\par\noindent
Music source: `Lutheran Worship' Hymnal, 1982 Hymn 436.

\hcopyrightsection{Amazing Grace}
\par\noindent
Words: John Newton, 1779.  last verse author unknown, before 1829.
\par\noindent
Music: `New Britain' James P. Carrell and David L. Clayton, 1831.  Setting: Edwin Othello Excell, 1900.
\par\noindent
Copyright: public domain. This score is a part of the Open Hymnal Project, 2006 Revision.
\par\noindent
Music source: ``Joy to the World'', 1915, Hymn 209 Ed. E. O. Excell  explicit about arrangement.  Christian Classics Ethereal Library (\url{ccel.org}) Hymnary  Think setting from ``Make His praise glorious'', 1900.  Ed. Edwin Othello Excell.

\hcopyrightsection{And did those feet in ancient time}
\par\noindent
Words: First (Original) Verse, Will­iam Blake, c. 1804. Second Verse, Kenneth G.H. Bryant, 2017.
\par\noindent
Music: Je­ru­sa­lem (Par­ry) Charles H. H. Par­ry, 1916.
\par\noindent
Arrangement: Augustine Watson, 2025.
\par\noindent
Copyright: \href{https://creativecommons.org/publicdomain/zero/1.0/}{CC0 Public Domain Dedication} (\href{https://imslp.org/wiki/Jerusalem_(Parry,_Charles_Hubert_Hastings)}{Bryant dedicated his words to the public domain}).

\hcopyrightsection{Be Thou My Vision}
\par\noindent
Words: Attr. Dallan Forgaill, 8th Century.  Translated by Mary Byrne, 1905 and Eleanor Hull, 1912. 
\par\noindent
Music: `Slane' Traditional Irish.  Setting: Mark Hamilton Dewey, 2007.
\par\noindent
Copyright: public domain. This score is a part of the Open Hymnal Project, 2008 Revision. 
\par\noindent
All portions of the setting that were not already public domain were released to the public domain by the arranger on 27 July 2007. He already had released the parts and the versification (except for a few changes in the third verse, which he released to the public domain in 2007) to the public domain in 2006.

\hcopyrightsection{Come Thou Fount of Every Blessing}
\par\noindent
Words: Robert Robinson, 1758.  Music: `Nettleton' Asahel Nettleton, 1812. 
\par\noindent
Setting: ``The Evangelical Hymnal'', 1921.
\par\noindent
Copyright: public domain. This score is a part of the Open Hymnal Project, 2010 Revision.
\par\noindent
Music source: The Evangelical Hymnal, 1921, Hymn 256.  Almost the same as The Presbyterian Hymnal, 1874 Hymn 94.

\hcopyrightsection{Crown Him with Many Crowns}
\par\noindent
Words: Verses 1, 4, 5, 6, \& 9: Matthew Bridges, The Passion of Jesus, 1852. Verses 2 \& 3: Godfrey Thring, Hymns and Sacred Lyrics, 1874. 
\par\noindent
Music: `Diademata' George J. Elvey, 1868.  Setting: ``Appendix to Hymns Ancient and Modern'', 1869.
\par\noindent
Copyright: public domain. This score is a part of the Open Hymnal Project, 2008 Revision.
\par\noindent
Music source: ``Appendix to Hymns Ancient and Modern'', 1869 Hymn 318. Ed. William H. Monk.  Unknown if the arrangement is Elvey's or Monk's.

\hcopyrightsection{Dear Christians, One and All Rejoice}
\par\noindent
Words: Martin Luther, 1523.  Translated by Richard Massie, 1854, alt. 
\par\noindent
Music: `Nun Freut Euch' attr. Martin Luther from Etlich Christlich Lider, Wittenberg, 1524. 
\par\noindent
Setting: Johann Hermann Schein, 1627.
\par\noindent
Copyright: public domain. This score is a part of the Open Hymnal Project, 2009 Revision.
\par\noindent
Music source: ``The Hymns of Martin Luther'' by Bacon, 1883.  Setting originally from ``Cantional oder Gesangbuch Augburgischer Konfession'' by Johann Hermann Schein.

\hcopyrightsection{Faith of Our Fathers}
\par\noindent
Words: Frederick W. Faber, 1849. Refrain by James G. Walton, 1874. 
\par\noindent
Music: `St. Catherine' Henri F. Hemy (1818-1888).  Setting: James G. Walton, 1874.
\par\noindent
Copyright: public domain. This score is a part of the Open Hymnal Project, 2006 Revision.
\par\noindent
Music source: The Evangelical Hymnal, 1921  Hymn 408.

\hcopyrightsection{I Need Thee Every Hour}
\par\noindent
Words: Annie Sherwood Hawks, 1872.
\par\noindent
Music: `I Need Thee Every Hour' Robert Lowry, 1872.  Setting: ``Pentecostal Hymns, No. 2'', 1898.
\par\noindent
Copyright: public domain. This score is a part of the Open Hymnal Project, 2012 Revision.
\par\noindent
Music source: ``Pentecostal Hymns, No. 2'' 1898 page 155.   ABC file contributed to the Open Hymnal by Samuel Cantrell, 18 Jan 2012.

\hcopyrightsection{I Sing the Mighty Power of God}
\par\noindent
Words: Isaac Watts, 1709.
\par\noindent
Music: `Ellacombe' from Gesangbuch der Herzogl. Hofkapelle, Wurttemberg, 1784. Setting: ``Amore Dei'', 1897.
\par\noindent
Copyright: public domain. This score is a part of the Open Hymnal Project, 2014 Revision.
\par\noindent
Music and Lyrics source: Hymnal ``Amore Dei'' (published 1897)  Hymn 115.  ABC file contributed to the Open Hymnal by Tobin Strong, 09 Jan 2014.

\hcopyrightsection{I Vow to Thee My Country}
\par\noindent
Words: Cecil Spring-Rice, 1859-1918.
\par\noindent
Music: `Thaxted' Gustav Holst, 1918.
\par\noindent
Arrangement: Augustine Watson, 2025, based on ``The Winchester Hymn Supplement: with Tunes'', 1928, in consultation with the arrangement by Paul Hayward, 2011.
\par\noindent
Copyright: \href{https://creativecommons.org/publicdomain/zero/1.0/}{CC0 Public Domain Dedication}.
\par\noindent
Music source: ``The Winchester Hymn Supplement: with Tunes'', 1928.

\hcopyrightsection{If God Had Not Been on Our Side}
\par\noindent
Words: Martin Luther, 1524.  Translation composite. 
\par\noindent
Music: `Wär Gott nicht mit uns diese Zeit (1537)' from Walter's Hymnal, 1537. 
\par\noindent
Setting: ``Evangelical Lutheran Hymn-Book'', 1931.
\par\noindent
Copyright: public domain. This score is a part of the Open Hymnal Project, 2009 Revision. Translation is public domain per Project Wittenberg: http://www.iclnet.org/pub/resources/text/wittenberg/hymns/ourside.txt

\hcopyrightsection{Immaculate Mary}
\par\noindent
Words: ``Catholic Church Hymnal'', 1905. Verses added by Augustine Watson, 2025 (put into the public domain).
\par\noindent
Music: Traditional tune.
\par\noindent
Setting: ``Catholic Church Hymnal'', 1905.
\par\noindent
Copyright: \href{https://creativecommons.org/publicdomain/zero/1.0/}{CC0 Public Domain Dedication}.

\hcopyrightsection{Immortal, Invisible, God Only Wise}
\par\noindent
Words: Walter Chalmers Smith, 1876. 
\par\noindent
Music: `St. Denio' or `Joanna' or `Palestrina' traditional Welsh found in ``Caniadau y Cyssegr'' by John Roberts, 1839.
\par\noindent
Setting: ``Caniadau y Cyssegr a'r Teulu'', 1878, alt.
\par\noindent
Copyright: public domain. This score is a part of the Open Hymnal Project, 2005 Revision.
\par\noindent
Music source: Episcopal Hymnal, 1940, Hymn 301.  Also in English Hymnal, 1906 Hymn 407 (page missing in my copy). Setting from ``Caniadau y Cyssegr a'r Teulu'', 1878 hymn 442, alt.

\hcopyrightsection{It Is Well with My Soul}
\par\noindent
Words: Horatio G. Spafford, 1873.
\par\noindent
Music and Setting: `It Is Well' or `Ville Du Havre' Philip P. Bliss, 1876.
\par\noindent
Copyright: public domain. This score is a part of the Open Hymnal Project, 2010 Revision.
\par\noindent
Music source: The Evangelical Hymnal, 1921 Hymn 208.

\hcopyrightsection{Joyful, Joyful, We Adore Thee}
\par\noindent
Words: Henry J. van Dyke, 1907. 
\par\noindent
Music: `Ode to Joy' Ludwig van Beethoven; Adapted by Edward Hodges, 1824. 
\par\noindent
Setting: ``The Methodist Hymnal'', 1905.
\par\noindent
Copyright: public domain. This score is a part of the Open Hymnal Project, 2005 Revision.
\par\noindent
Music source: The Methodist Hymnal, 1905 Hymn 160.

\hcopyrightsection{Let All Mortal Flesh Keep Silence}
\par\noindent
Words: from Liturgy of St. James, 4th Century.  Translated by Gerard Moultrie, 1864. 
\par\noindent
Music: `Picardy' traditional French.  Setting: ``The English Hymnal'', 1906, alt.
\par\noindent
Copyright: public domain. This score is a part of the Open Hymnal Project, 2008 Revision.
\par\noindent
Music source: The Episcopal Hymnal, 191x Hymn 339 and English Hymnal 1906 Hymn 318.

\hcopyrightsection{Lift High the Cross}
\par\noindent
Words: George W. Kitchin (1827-1912). Modified by Michael R. Newbolt, 1916. 
\par\noindent
Music: `Crucifier' Sydney H. Nicholson, 1916.  Setting: ``Hymns Ancient and Modern'', 1922.
\par\noindent
Copyright: public domain. This score is a part of the Open Hymnal Project, 2006 Revision.
\par\noindent
Music source: `Lutheran Worship' Hymnal, 1982 Hymn 311.  From ``Hymns Ancient and Modern'', Standard Edition Hymn 745, 1922.

\hcopyrightsection{Look Down, O Lord, from Heaven Behold}
\par\noindent
Words: Martin Luther, 1524. Translated by Frances Elizabeth Cox, 1864. 
\par\noindent
Music: `Ach Gott vom Himmel' from Erfurt Enchiridion, 1524.  Setting: Karl August Haupt, 1869. This tune is also used in Klug's 1543 hymnal, and that hymnal is sometimes incorrectly referenced as the source
\par\noindent
Copyright: public domain. This score is a part of the Open Hymnal Project, 2009 Revision.
\par\noindent
Music source: ``The Hymns of Martin Luther'' by Bacon, 1883.

\hcopyrightsection{Lord, Keep Us Steadfast in Thy Word}
\par\noindent
Words: Martin Luther, 1541.  Translated by Catherine Winkworth, 1863. 
\par\noindent
Music: `Erhalt Uns, Herr, bei deinem Wort' from Klug's ``Geistliche Lieder auffs new gebessert'', 1543. 
\par\noindent
Setting: Hans Leo Hassler, unknown date.
\par\noindent
Copyright: public domain. This score is a part of the Open Hymnal Project, 2008 Revision.
\par\noindent
Music source: Evangelical Lutheran Hymn Book (LCMS), Edition of 1931 Hymn 274.

\hcopyrightsection{Now Thank We All Our God}
\par\noindent
Words: Martin Rinkart, c.1636. Translated Catherine Winkworth, 1856.  Music: `Nun Danket' Johann Crüger, 1647. 
\par\noindent
Setting: ``Evangelical Lutheran Hymn-Book'', 1931.
\par\noindent
Copyright: public domain. This score is a part of the Open Hymnal Project, 2006 Revision.
\par\noindent
Music source: `Evangelical Lutheran Hymn-Book', Edition of 1931 Hymn 64.

\hcopyrightsection{Our God, Our Help in Ages Past}
\par\noindent
Words: Isaac Watts, 1719.  Music: `St. Anne' William Croft, 1708.
\par\noindent
Setting: composite found in ``The Lutheran Hymnary'', 1913.
\par\noindent
Copyright: public domain. This score is a part of the Open Hymnal Project, 2005 Revision.
\par\noindent
Music source: `The Lutheran Hymnary', 1913, Hymn 261. Arrangement is composite, first half from ``Hymns Ancient and Modern'', 1869, Hymn 231. ed  William Henry Monk. Second half is from ``The Hymnal Companion to the Book of Common Prayer'', 1890, Hymn 279. ed. Charles Vincent, D.J. Wood, John Stainer

\hcopyrightsection{Out of the Deep I Cry to Thee}
\par\noindent
Words: Martin Luther, 1524. Translated by Arthur Tozer Russell (1806-1874). 
\par\noindent
Music: `Aus Tiefer Not (Luther)' or `Af Dybsens Nød' Martin Luther from Erfurt Enchiridion, 1524. 
\par\noindent
Setting: Johann Sebastian Bach, 1725.
\par\noindent
Copyright: public domain. This score is a part of the Open Hymnal Project, 2009 Revision.
\par\noindent
Music source: ``The Hymns of Martin Luther'' by Bacon, 1883.

\hcopyrightsection{Praise God from Whom All Blessing Flow}
\par\noindent
Words: Thomas Ken, 1674.  Music: `Old 100th' Genevan Psalter, attr. Louis Bourgeois, c. 1551.  
\par\noindent
Setting: Sternhold and Hopkins' Psalter, 1561.
\par\noindent
Copyright: public domain. This score is a part of the Open Hymnal Project, 2006 Revision.
\par\noindent
Music source:ccel from Sternhold and Hopkins' Psalter 1561.

\hcopyrightsection{Praise My Soul the King of Heaven}
\par\noindent
Words: Henry F. Lyte, 1834. 
\par\noindent
Music: `Praise My Soul' or `Lauda Anima' or `St. Paul' John Goss, 1869.  Setting: ``The Choral Hymnal'', 1888.
\par\noindent
Copyright: public domain. This score is a part of the Open Hymnal Project, 2006 Revision.
\par\noindent
Music Source: ``The Plymouth Hymnal'', 1894 Hymn 219.  ``The Choral Hymnal'', 1888 page 54.

\hcopyrightsection{Praise to the Lord, the Almighty}
\par\noindent
Words: Joachim Neander, 1680. Translated by Catherine Winkworth, 1863. 
\par\noindent
Music: `Lobe den Herren' from Ander Theil des Erneuerten Gesangbuch, 1665.
\par\noindent
Setting: William Sterndale Bennett, 1863, alt. 
\par\noindent
Copyright: public domain. This score is a part of the Open Hymnal Project, 2006 Revision.
\par\noindent
Music source: ``The Chorale Book for England'', 1863 Hymn 9.  Ed. Winkworth (words) and Bennett (music)., slightly alt. per ``Evangelical Lutheran Hymn-Book'', 1931.

\hcopyrightsection{That Men a Godly Life Might Live}
\par\noindent
Words: Martin Luther, 1524.  Translated by Richard Massie, 1854, alt. 
\par\noindent
Music: `Dies sind die heil`gen zehu Gebot' or `In Gottes Namen fahren wir' circa 1200s found in Erfurt Enchiridion, 1524.  Setting: Michael Praetorius, 1609.
\par\noindent
Copyright: public domain. This score is a part of the Open Hymnal Project, 2009 Revision.

\hcopyrightsection{The Church's One Foundation}
\par\noindent
Words: Samuel John Stone, 1866.  Music: `Aurelia' Samuel Sebastian Wesley, 1864. Verse 4 modified by Augustine Watson, 2025.
\par\noindent
Setting: ``Order of worship for the Reformed Church in the United States'', 1866.
\par\noindent
Copyright: \href{https://creativecommons.org/publicdomain/zero/1.0/}{CC0 Public Domain Dedication}.
\par\noindent
Music source: `Lutheran Worship' Hymnal, 1982 Hymn 289. Almost the same as ``Order of worship for the Reformed Church in the United States'', 1866 Hymn 441

\hcopyrightsection{The King of Love}
\par\noindent
Words: Henry W. Baker, 1868.
\par\noindent
Music: Irish Melody, harm. from ``The English Hymnal'', 1906.
\par\noindent
Copyright: public domain.
\par\noindent
Music source: Hymnary.org

\hcopyrightsection{Though in the Midst of Life We Be}
\par\noindent
Words: v.1 Medieval sequence, vs 2,3 Martin Luther, 1524.  Translated by Richard Massie, 1854, alt. 
\par\noindent
Music: `Mitten wir im Leben Sind' Medieval sequence altered by Martin Luther. Found in Walter's Geistliche Gesangbüchlein, 1524.  Setting: Erythraeus, 1608.
\par\noindent
Copyright: public domain. This score is a part of the Open Hymnal Project, 2009 Revision.

\hcopyrightsection{To Avert from Men God's Wrath}
\par\noindent
Words: Latin c. 1400, sometimes attr. John Hus.  Translated Christian Ignatius Latrobe, 1789. 
\par\noindent
Music: `Gethsemane' or `Petra' Richard Redhead, 1853. 
\par\noindent
Setting: ``The Church Hymnal, Revised and Enlarged'' (Episcopal), 1896.
\par\noindent
Copyright: public domain. This score is a part of the Open Hymnal Project, 2011 Revision.
\par\noindent
Music source: ``The Church Hymnal, Revised and Enlarged'' (Episcopal), 1896  Ed. Charles Hutchins Hymn 93. 
\par\noindent
Lyrics source:  ``Hymnal and Liturgies of the Moravian Church'', 1920.

\hcopyrightsection{Ye Watchers and Ye Holy Ones}
\par\noindent
Words: John Athelstan Laurie Riley, 1906.
\par\noindent
Music: Geistliche Kirchengesäng, 1623; harm. Ralph Vaughan Williams, 1906.
\par\noindent
Copyright: public domain.
\par\noindent
Music source: Hymnary.org

%Advent:
\hcopyrightsection{Come, Thou Long Expected Jesus}
\par\noindent
Words: Charles Wesley, 1745. 
\par\noindent
Music and Setting: `Jefferson'  from ``Southern Harmony'', 1835, alt.
\par\noindent
Copyright: public domain.  This score is a part of the Open Hymnal Project, 2008 Revision.
\par\noindent
Music source: Southern Harmony, 1835, Hymn 42.  Setting heavily altered for congregational use.


\hcopyrightsection{O Come O Come Emmanuel}
\par\noindent
Words: various, combined by unknown author approx 12th Century, Translated by John Mason Neale, 1851. Second and seventh verses translated by H. S. Coffin (modified by Augustine Watson), 1916. Last verses by Augustine Watson, 2025.
\par\noindent
Music: `Veni Emmanuel' 15th Century French processional.
\par\noindent
Setting: ``Common Service Book'' (ULCA), 1917. 
\par\noindent
Arrangement: Augustine Watson, 2025.
\par\noindent
Copyright: \href{https://creativecommons.org/publicdomain/zero/1.0/}{CC0 Public Domain Dedication}.
\par\noindent
Music source: ULCA Hymnal, 1917 Hymn 1. PECUSA Hymnal (1920), Hymn 44.

\hcopyrightsection{On Jordan's Bank the Baptist's Cry}
\par\noindent
Words: Charles Coffin, 1736. st. 1-3 translated by John Chandler, 1837; st 4-5 translator unknown.  
\par\noindent
Music: `Puer Nobis Nascitur' Michael Praetorius, 1609. 
\par\noindent
Setting: George Ratcliffe Woodward for ``The English Hymnal'', 1906.
\par\noindent
Copyright: public domain. This score is a part of the Open Hymnal Project, 2006 Revision.
\par\noindent
Music source: `The English Hymnal', 1906 Hymn 14.

\hcopyrightsection{Saviour of the Nations Come}
\par\noindent
Words: Ambrose of Milan, c. 397.  Translated to German by Martin Luther, 1524. Translated from German to English by William M. Reynolds, 1851. 
\par\noindent
Music: `Nun Komm, Der Heiden Heiland' from Walter's Geistliche Gesangbüchlein, 1524.
\par\noindent
Setting: ``Mehrstimmiges ChoralBuch'', 1906.
\par\noindent
Copyright: public domain. This score is a part of the Open Hymnal Project, 2007 Revision.
\par\noindent
Music source: ``Mehrstimmiges ChoralBuch'', 1906 Hymn \#135, page 109 Ed. Karl Brauer. slightly altered. The Evangelical Lutheran Hymn Book, 1931, Hymn 141.

%Christmas:
\hcopyrightsection{A Great and Mighty Wonder}
\par\noindent
Words: Germanus of Constantinople (634-734).  Translated by John Mason Neale, 1862. 
\par\noindent
Music: `Es Ist Ein Ros Entsprungen (Rhythmic)' German from Köln, 1599.  Setting: Michael Praetorius, 1609.
\par\noindent
Copyright: public domain. This score is a part of the Open Hymnal Project, 2010 Revision.
\par\noindent
Music source: The English Hymnal, 1906.  Hymn 19

\hcopyrightsection{All Praise to Jesus' Hallowed Name}
\par\noindent
Words: verse 1, ancient German.  verses 2-7, Martin Luther, 1524.  Translated by Richard Massie, 1854, alt. 
\par\noindent
Music: `Gelobet Seist Du' ancient German found in Walter's Geistliche Gesangbüchlein, 1524. 
\par\noindent
Setting: Karl August Haupt, 1869.
\par\noindent
Copyright: public domain. This score is a part of the Open Hymnal Project, 2009 Revision.
\par\noindent
Music source: ``The Hymns of Martin Luther by Leonard Woolsey Bacon 1883, p. 20.

\hcopyrightsection{Angels We Have Heard on High}
\par\noindent
Words: French Carol; Translated by James Chadwick, 1862. 
\par\noindent
Music: `Gloria' French carol melody.  Setting: Edward (or Edwin) S. Barnes, before 1916.
\par\noindent
Copyright: public domain. This score is a part of the Open Hymnal Project, 2005 Revision.
\par\noindent
Music source: `Lutheran Worship' Hymnal, 1982 Hymn 55. ``Carols Old And Carols New'', 1916 Carol 181.

\hcopyrightsection{Away in a Manger}
\par\noindent
Words: stanzas 1,2 anonymous published Philadelphia, 1885.  stanza 3 John T. MacFarland (1851-1913). 
\par\noindent
Music: `Mueller' James R. Murray, 1887.  Setting: ``Hymnal for American Youth'', 1919.
\par\noindent
Copyright: public domain. This score is a part of the Open Hymnal Project, 2006 Revision.
\par\noindent
Music source: ``Hymnal for American Youth'', Hymn 84 1919 ed. H. Augustine Smith. Music reputed to be first published in ``Dainty Songs for Little Lads and Lasses'' by James R. Murray, 1887, hymn 8.

\hcopyrightsection{God Rest Ye Merry Gentlemen}
\par\noindent
Words: Traditional English. 
\par\noindent
Music: `God Rest Ye Merry Gentlemen' Traditional English.  Setting: ``Carols Old And Carols New'', 1918.
\par\noindent
Copyright: public domain. This score is a part of the Open Hymnal Project, 2006 Revision.
\par\noindent
Music source: `Carols Old And Carols New', 1918 carol 722.

\hcopyrightsection{Good King Wenceslas}
\par\noindent
Words: John M. Neale, 1853.
\par\noindent
Music: `Tempus Adest Floridum' 13th Century spring carol; first published in the Swedish Piae Cantones, 1582.
\par\noindent
Setting: Carols Old And Carols New, 1916.
\par\noindent
Copyright: public domain. This score is a part of the Open Hymnal Project, 2006 Revision.

\hcopyrightsection{Hark! The Herald Angels Sing}
\par\noindent
Words: Charles Wesley, 1739, alt. 
\par\noindent
Music: `Mendelssohn' from `Festgesang' Felix Mendelssohn, 1840.  Setting: William H. Cummings, 1857.
\par\noindent
Copyright: public domain. This score is a part of the Open Hymnal Project, 2005 Revision.
\par\noindent
Music source: `Lutheran Worship' Hymnal, 1982 Hymn 49.

\hcopyrightsection{In the Bleak Mid-Winter}
\par\noindent
Words: Christina Georgina Rossetti, 1872, alt. 
\par\noindent
Music and Setting: `Cranham' Gustav Theodore Holst, 1906, alt. 
\par\noindent
Copyright: public domain. This score is a part of the Open Hymnal Project, 2011 Revision.
\par\noindent
Music source: ``The English Hymnal, 1906. 

\hcopyrightsection{It Came upon a Midnight Clear}
\par\noindent
Words: Edmund H. Sears, 1849.  Music: `Carol' Richard S. Willis, 1861.  
\par\noindent
Setting: ``Order of worship for the Reformed Church in the United States'', 1866.
\par\noindent
Copyright: public domain. This score is a part of the Open Hymnal Project, 2005 Revision.
\par\noindent
Music source: ``Order of worship for the Reformed Church in the United States'', 1866 Hymn 63. Supposed to be found in ``Church Corals and Choir Studies'' 1850 by Willis, but I don't see it in that book. Some sources say 1850, some 1859, some 1861.  Willis was editor for ``The Musical World'' magazine.  Probably first appeared there.

\hcopyrightsection{Joy to the World}
\par\noindent
Words: Isaac Watts, 1719.
\par\noindent
Music: `Antioch' pieced together from ``Messiah'' George F. Handel, 1741.  Setting: Lowell Mason, 1836.
\par\noindent
Copyright: public domain. This score is a part of the Open Hymnal Project, 2005 Revision.
\par\noindent
Music source: `Lutheran Worship' Hymnal, 1982 Hymn 53.

\hcopyrightsection{Lo, How A Rose E'er Blooming}
\par\noindent
Words: verses 1-2, 15th Century German.  Translated by Theodore Baker, 1894. Verses 3,4 Fridrich Layriz (1808-1859).  Translated by Harriet Reynolds Krauth, 1875. Verse 5, 15th Century German.  Translated by John C. Mattes, 1914.
\par\noindent
Music: `Es Ist Ein Ros Entsprungen (Rhythmic)' German from Köln, 1599.  Setting: Michael Praetorius, 1609.
\par\noindent
Copyright: public domain. This score is a part of the Open Hymnal Project, 2010 Revision.
\par\noindent
Music source: The English Hymnal, 1906.  Hymn 19

\hcopyrightsection{O Come, All Ye Faithful}
\par\noindent
Words:  John F. Wade, circa 1743. v.1-3, 6 Translated by Frederick Oakeley, 1841; v. 4, 5 Translated by William T. Brooke (1848-1917). 
\par\noindent
Music: `Adeste Fideles' or `Portuguese Hymn' John F. Wade, 1743.  Setting: ``A Hymnal'' (Episcopal), 1916.
\par\noindent
Copyright: public domain. This score is a part of the Open Hymnal Project, 2007 Revision.
\par\noindent
Music source: ``A Hymnal'' (Episcopal), 1916 Hymn 72, alt.

\hcopyrightsection{O Word of God the Father}
\par\noindent
Words: Bradford Littlejohn.
\par\noindent
Music: `Thaxted', Holst, 1921.
\par\noindent
Arrangement: Augustine Watson, 2025.
\par\noindent
Copyright: 2024 Bradford Littlejohn. Permission received for use in this hymnal.  All Rights Reserved.

\hcopyrightsection{Silent Night}
\par\noindent
Words: Josef Mohr, 1818. stanzas 1,3 Translated by John Freeman Young, 1863. Stanzas 2,4 translator anonymous.
\par\noindent
Music: `Stille Nacht' Franz Xaver Gruber, 1818.  Setting: ``Concordia Kinderchöre'', 1908.
\par\noindent
Copyright: public domain. This score is a part of the Open Hymnal Project, 2005 Revision.
\par\noindent
Music source: `Lutheran Worship' Hymnal, 1982 Hymn 68.
\par\noindent
Music source: ``Concordia Kinderchöre'', 1908 Hymn 42 page 54.

\hcopyrightsection{The First Noel}
\par\noindent
Words: Traditional English carol, possibly dating from as early as the 13th Century.
\par\noindent
Music: `The First Noel' Traditional English carol, possibly dating from as early as the 13th Century.
\par\noindent
Setting: ``The Methodist Sunday School Hymnal'', 1911.
\par\noindent
Copyright: public domain. This score is a part of the Open Hymnal Project, 2006 Revision.
\par\noindent
Music source: The Methodist Sunday School Hymnal, 1911 Hymn 66.

\hcopyrightsection{To Shepherds as They Watched by Night}
\par\noindent
Words: Martin Luther, 1543. translated by Richard Massie, 1854. 
\par\noindent
Music: `Vom Himmel Hoch' traditional German from Schumann's Geistliche Lieder, Leipzig, 1539. 
\par\noindent
Setting: ``Common Service Book'' (ULCA), 1917.
\par\noindent
Copyright: public domain. This score is a part of the Open Hymnal Project, 2009 Revision.
\par\noindent
Music source: `Common Service Book with Hymnal', ULCA 1918 Hymn 19.

\hcopyrightsection{What Child is This?}
\par\noindent
Words: William Chatterton Dix, 1865. 
\par\noindent
Music: `Greensleeves' 16th Century English Traditional. 
\par\noindent
Setting: traditional from ``The Sunday School Hymnal and Service Book'', 1871.
\par\noindent
Copyright: public domain. This score is a part of the Open Hymnal Project, 2005 Revision.
\par\noindent
Music source: `Lutheran Worship' Hymnal, 1982 Hymn 61. very tiny changes from ``The Sunday school hymnal and service book'', 1871 edited by Charles Lewis Hutchins Carol 9 (after hymns).

%Epiphany:
\hcopyrightsection{As with Gladness Men of Old}
\par\noindent
Words: William Chatterton Dix, 1860. 
\par\noindent
Music: `Dix' or `Treuer Heiland, Wir Sind Heir' Conrad Kocher, 1838. Abridged by William Henry Monk, 1861.
\par\noindent
Setting: Conrad Kocher, 1838, alt. by William Henry Monk, 1861, alt. for ``The English Hymnal'', 1906.
\par\noindent
Copyright: public domain. This score is a part of the Open Hymnal Project, 2008 Revision.
\par\noindent
Music source: The English Hymnal, 1906 Hymn 39. Music from ``Stimmen aus dem Reiche Gottes'', 1838 by Kocher, Hymn 201 page 250. Adapted by William Henry Monk from the original (removed two measures and changed parts of the arrangement) for ``Hymns Ancient and Modern, 1861. Also almost just like the Bristol Tune Book of 1863, Hymn 172.

\hcopyrightsection{Brightest and Best of the Sons of Morning}
\par\noindent
Words: James P. Harding, 1892.
\par\noindent
Music: Morning Star.
\par\noindent
Copyright: public domain.
\par\noindent
Music source: The English Hymnal, 1940 Hymn 46.

\hcopyrightsection{Lord, Who at Cana's Wedding Feast}
\par\noindent
Words: Adelaide Thrupp, 1853. Music: `St. Ursula' Frederick Westlake, 1863. 
\par\noindent
Setting: ``The Church Hymnal, Revised and Enlarged'' (Episcopal), 1896.
\par\noindent
Copyright: public domain. This score is a part of the Open Hymnal Project, 2011 Revision.
\par\noindent
Music source: ``The Church Hymnal, Revised and Enlarged'' (Episcopal), 1896 Ed. Charles Hutchins Hymn 237. Words compared against ``The Church Hymnal, Revised and Enlarged'', 1896.

\hcopyrightsection{We Three Kings}
\par\noindent
Words: John Henry Hopkins, Jr., 1857.
\par\noindent
Music: John Henry Hopkins, Jr., 1857, alt.
\par\noindent
Setting: ``Cyber Hymnal''.
\par\noindent
Arrangement: Augustine Watson, 2025.
\par\noindent
Copyright: \href{https://creativecommons.org/publicdomain/zero/1.0/}{CC0 Public Domain Dedication}.
\par\noindent
Music source: ``Cyber Hymnal''.

%Septuagesimatide
\hcopyrightsection{Maker of Earth, To Thee Alone}
\par\noindent
Words: Laeta mundi, conditor. C. Coffin. Translated by J. M. Neale.
\par\noindent
Music: Dunfermline. ``Scottish Psalter'', 1615.
\par\noindent
Copyright: public domain.
\par\noindent
Music source: ``The English Hymnal: with tunes'', 1906, Hymn 64.

\hcopyrightsection{O Love, How Deep, How Broad}
\par\noindent
Words: O Amor quam ecstaticus. \nth{15} century. Translated by B. Webb.
\par\noindent
Music: Eisenach. J. H. Schein. J. S. Bach.
\par\noindent
Copyright: public domain.
\par\noindent
Music source: ``The English Hymnal: with tunes'', 1906, Hymn 459.

\hcopyrightsection{Praise to the Holiest in the Height}
\par\noindent
Words: J. H. Newman.
\par\noindent
Music: Richmond. Adapted from T. Haweis by S. Webbe (the younger).
\par\noindent
Copyright: public domain.
\par\noindent
Music source: ``The English Hymnal: with tunes'', 1906, Hymn 471.

\hcopyrightsection{There is a Book Who Runs May Read}
\par\noindent
Words: J. Keble.
\par\noindent
Music: St. Flavian. Adapted from Psalm 132 in `Day's Psalter', 1563.
\par\noindent
Copyright: public domain.
\par\noindent
Music source: ``The English Hymnal: with tunes'', 1906, Hymn 497.

%Lent:
\hcopyrightsection{All Glory, Laud, and Honour}
\par\noindent
Words: Theodulf of Orleans, circa 820. Translated by John Mason Neale, 1851.  
\par\noindent
Music: `Valet Will Ich Dir Geben' or `St. Theodulph' Melchior Teschner, 1615. 
\par\noindent
Setting: Presbyterian Hymnal, 1911.
\par\noindent
Copyright: public domain. This score is a part of the Open Hymnal Project, 2006 Revision.
\par\noindent
Music source: Presbyterian Hymnal, Revised, 1911 Hymn 216.

\hcopyrightsection{Lord Who Throughout These Forty Days}
\par\noindent
Words: Claudia F. Hernaman, 1873.  Music: `St. Flavian' Day's Psalter, 1563. 
\par\noindent
Setting: ``The Church Hymnal, Revised and Enlarged'' (Episcopal), 1905.
\par\noindent
Copyright: public domain. This score is a part of the Open Hymnal Project, 2008 Revision.
\par\noindent
Music source: Episcopal Hymnal, 1905 Hymn 78.

\hcopyrightsection{O Jesu Christ From Thee Began}
\par\noindent
Words: \nth{9} century text. 
\par\noindent
Music: Plaistow.
\par\noindent
Copyright: Public domain.
\par\noindent
Music source: ``The English Hymnal: with tunes'', 1906, Hymn 69.


\hcopyrightsection{O Sacred Head, Now Wounded}
\par\noindent
Words: Bernard of Clairvaux, 1153. Translated by James W. Alexander, 1830. 
\par\noindent
Music: `Passion Chorale' or `Herzlich Tut Mich Verlangen' Hans Leo Hassler, 1601. Adapted by J.S. Bach, 1729.  
\par\noindent
Setting: Johann Sebastian Bach, 1729.
\par\noindent
Copyright: public domain. This score is a part of the Open Hymnal Project, 2007 Revision.
\par\noindent
Music source: The Episcopal Hymnal, 1916, Hymn 158.


\hcopyrightsection{When I Survey the Wondrous Cross}
\par\noindent
Words: Isaac Watts, 1707.
\par\noindent
Music: `Duke Street' John Hatton, 1793.  Setting: ``Finest of the Wheat No. 3'' Hymnal, 1904.
\par\noindent
Copyright: public domain. This score is a part of the Open Hymnal Project, 2014 Revision.
\par\noindent
Music source: ``Finest of the Wheat No. 3'' Hymnal, 1904 Hymn 257.  ABC file contributed to the Open Hymnal by Tobin Strong, 20 Dec 2013. Composite from ``The Anglican hymn Book'', 1871 hymn 127 and ``Hymns Ancient and Modern'', 1869 Hymn 101.

%Easter:
\hcopyrightsection{At the Lamb's High Feast}
\par\noindent
Words: Latin, circa 6th Century. Translated by Robert Campbell, 1849. 
\par\noindent
Music: `Sonne der Gerechtigkeit' Czech, Kirchengeseng, 1566.  Setting: Brian J. Dumont, 31 Dec 2009.
\par\noindent
Copyright: Words and Music, public domain. Setting: Copyright 2009 Brian J. Dumont. This setting may be freely reproduced or published for Christian worship, provided it is not altered, and this notice is on each copy. All other rights reserved.   This score is a part of the Open Hymnal Project, 2010 Revision.
\par\noindent
Music source: tune found in many places, arrangement by BJD

\hcopyrightsection{Christ the Lord Is Risen Today}
\par\noindent
Words: Stanzas 1-7, Charles Wesley, 1739. Stanzas 8-10, 14th Century; translated in Lyra Davidica. 
\par\noindent
Music: `Llanfair' Robert Williams, 1817.  Setting: John Roberts, 1837.
\par\noindent
Copyright: public domain. This score is a part of the Open Hymnal Project, 2008 Revision.
\par\noindent
Music source: Lutheran Worship, 1982, Hymn 137 (later modified to more closely match John Robert's original setting).

\hcopyrightsection{Hail Thee, Festival Day}
\par\noindent
Words: Venantius Honorius Fotunatus; translation from `The English Hymnal', 1906, alt.
\par\noindent
Music: Ralph Vaughan Williams, 1906.
\par\noindent
Copyright: public domain.
\par\noindent
Music source: `Songs of Praise', 1925, Hymn 445. PECUSA Hymnal, 1940, Hymn 86.

\hcopyrightsection{The Strife is O'er, the Battle Done}
\par\noindent
Words: from Symphonia Sirenum Selectarum, Kˆln, 1695; translated by Francis Pott, 1861.
\par\noindent
Music: `Victory' or  `Palestrina' Giovanni P. da Palestrina, 1591.
\par\noindent
Setting: William Henry Monk, 1861.
\par\noindent
Copyright: public domain.  This score is a part of the Open Hymnal Project, 2011 Revision.
\par\noindent
Music source: ``Hymns Ancient and Modern'', 1861 Hymn 114.  Words checked against the same source.


%Ascensiontide
\hcopyrightsection{See, The Lord Ascends in Triumph}
\par\noindent
Words: Christopher Wordsworth, 1862, alt.
\par\noindent
Music: 'Rex Gloriae' Henry Thomas Smart, 1868. 
\par\noindent
Setting: "Appendix to Hymns Ancient and Modern", 1869.
\par\noindent
Copyright: public domain.  This score is a part of the Open Hymnal Project, 2011 Revision.
\par\noindent
Music source: ``Appendix to Hymns Ancient and Modern'', 1869 Hymn 293.  Text is the same as ``Appendix to Hymns Ancient and Modern'' with the exception of the title.

\hcopyrightsection{Sing We Triumphant Hymns of Praise}
\par\noindent
Words: St. Bede the Venerable. Translated by B. Webb.
\par\noindent
Music: Tugwood. Nicholas Gatty.
\par\noindent
Copyright: public domain.
\par\noindent
Music source: ``The English Hymnal: with tunes'', 1906, Hymn 146.

%Whitsun
\hcopyrightsection{Blest Joys for Mighty Wonders Wrought}
\par\noindent
Words: Beata nobis gaudia. St. Hilary of Poitiers, \nth{4} century. Translated by Neale.
\par\noindent
Music: T. B. Southgate.
\par\noindent
Arrangement: Augustine Watson, 2025.
\par\noindent
Copyright: \href{https://creativecommons.org/publicdomain/zero/1.0/}{CC0 Public Domain Dedication}
\par\noindent
Music source: ``The Armagh Hymnal'', 1915, Hymn 46.

\hcopyrightsection{Let the Holy Spirit's Grace}
\par\noindent
Words: Nobis Sancti Spiritus. Benedict XII, \nth{14} century. Translated by Neale.
\par\noindent
Music: Regina Clementiae. Harleian MS. 978. Mode I.
\par\noindent
Arrangement: Augustine Watson, 2025.
\par\noindent
Copyright: \href{https://creativecommons.org/publicdomain/zero/1.0/}{CC0 Public Domain Dedication}
\par\noindent
Music source: ``The Armagh Hymnal'', 1915, Hymn 47.


%Trinity Sunday:
\hcopyrightsection{God the Father Be Our Stay}
\par\noindent
Words: 15th Century Litany, adapted by Martin Luther, 1524.  Translated by Richard Massie, 1854, alt. 
\par\noindent
Music: `Gott Der Vater, Wohn Uns Bei' from Walter's Geistliche Gesangbüchlein, 1524.
\par\noindent
Setting: composite from Landgraf Moritz, 1612 and ``Evangelical Lutheran Hymn-Book'', 1931.
\par\noindent
Copyright: public domain. This score is a part of the Open Hymnal Project, 2009 Revision.

\hcopyrightsection{Holy God, We Praise Thy Name}
\par\noindent
Words: attr. Ignaz Franz, 1774. Translated by Clarence A. Walworth, 1858. 
\par\noindent
Music: `Te Deum' or `Hursley' or `Grosser Gott, wir Loben Dich' from Katholisches Gesangbuch, Maria Theresa, 1774. Setting: ``Hymns Ancient and Modern'', 1869, alt.
\par\noindent
Copyright: public domain. This score is a part of the Open Hymnal Project, 2005 Revision.
\par\noindent
Music source: ``Hymns Ancient and Modern'', 1869 Hymn 11

\hcopyrightsection{Holy, Holy, Holy}
\par\noindent
Words: Reginald Heber, 1826. 
\par\noindent
Music: `Nicaea' John Bacchus Dykes, 1861.  Setting: ``Hymns Ancient and Modern'', 1869.
\par\noindent
Copyright: public domain. This score is a part of the Open Hymnal Project, 2005 Revision.
\par\noindent
Music source: Hymns Ancient and Modern, 1869 hymn 135.  Tune from Hymns Ancient and Modern, 1861.

\hcopyrightsection{I Bind unto Myself Today}
\par\noindent
Words: attributed to St. Patrick of Ireland (circa 387-466). Paraphrased by Cecil F. Alexander, 1889. 
\par\noindent
Music: `St. Patricks Breastplate' Charles V. Stanford, 1902.  Setting: ``The English Hymnal'', 1906.
\par\noindent
Copyright: public domain. This score is a part of the Open Hymnal Project, 2006 Revision.
\par\noindent
Music source: `The English Hymnal', 1906 Hymn 212.

\hcopyrightsection{Isaiah, Mighty Seer, in Days of Old}
\par\noindent
Words: Martin Luther, 1526 as the Sanctus of the German Mass, after Is 6:1-4.  Translation composite. 
\par\noindent
Music: `Jesaia Dem Propheten das Geschah' Martin Luther, 1526 in the German Mass.  Setting: Erythraeus, 1608.
\par\noindent
Copyright: public domain. This score is a part of the Open Hymnal Project, 2009 Revision. Translation is public domain per Project Wittenberg: http://www.iclnet.org/pub/resources/text/wittenberg/hymns/isaiah.txt

\hcopyrightsection{We All Believe in One True God}
\par\noindent
Words: Medieval text expanded by Martin Luther, 1524.  Translation composite. 
\par\noindent
Music: `Wir Glauben all an Einen Gott, Schoepfer' or `Apostolic Creed' Medieval tune altered by Martin Luther. Found in Walter's Geistliche Gesangbüchlein, 1524. 
\par\noindent
Setting: ``Eisenach Kirchenconserenz'', G.v. Tucher et. al. , Stuttgart, 1854.
\par\noindent
Copyright: public domain. This score is a part of the Open Hymnal Project, 2009 Revision. Translation is public domain per Project Wittenberg: http://www.iclnet.org/pub/resources/text/wittenberg/hymns/believe.txt

%Transfiguration:
\hcopyrightsection{O Wondrous Type, O Vision Fair}
\par\noindent
Words: From the ``Sarum Breviary'', 1495.  Translated by John Mason Neale, 1851.
\par\noindent
Music: `Deo Gracias' or `Agincourt' traditional English, circa 1415.  Setting: Charles Winfred Douglas, 1918, alt.
\par\noindent
Copyright: public domain. This score is a part of the Open Hymnal Project, 2011 Revision.
\par\noindent
Music source: ``A Hymnal'' (Episcopal), 1918 Hymn 439, alt by BJD 21 Mar 2011.  Translation found in Hymns Ancient \& Modern, 1861.

\hcopyrightsection{The Mouth of Fools Doth God Confess}
\par\noindent
Words: Martin Luther, 1524. Translated by Richard Massie, 1854, alt. 
\par\noindent
Music: `Es spricht der Unweisen Mund' from Walter's Geistliche Gesangbüchlein, 1524.  Setting: Michael Praetorius, 1610.
\par\noindent
Copyright: public domain. This score is a part of the Open Hymnal Project, 2009 Revision.

\hcopyrightsection{Tis Good, Lord, to Be Here}
\par\noindent
Words: Joseph Armitage Robinson, 1888.
\par\noindent
Music: `Potsdam' adapted from Johann Sebastian Bach, 1750, by John Goss, 1854. Setting: John Goss, 1854.
\par\noindent
Copyright: public domain. This score is a part of the Open Hymnal Project, 2011 Revision.
\par\noindent
Music source: ``The Church Psalter and Hymn Book'', 1863 hymns 212-214 (rev. from 1854 edition).  Mus. Dir and arrangements by John Goss. Words match those in ``Hymns Ancient And Modern'', 1904, Hymn 251

\end{multicols}

\normalsize
\phantomsection
\addcontentsline{toc}{chapter}{Preface}
\fancyhead[C]{{\LARGE Preface}}
\noindent
We are blessed to publish this necessary hymnal and companion to the 2024 Proposed Book of Common Prayer. The Prayer Book Hymnal supplements the Book of Common Prayer by providing liturgical elements, which may prove expedient to divine worship, along with notation for the hymns. For the convenience of the priest, the frequently said prayers of the Mass (when his BCP may not be easily accessed) are provided. For both the Office and the Mass, the most common settings and common chants are set as Office and Communion Services. Regarding the propers, the Office and Mass propers for the minor Feast Days (not found in the BCP), as well as the regular ferial hymns, are provided. Additionally, the Office propers which are omitted in the Red Letter Days are also provided. These have all been typeset to Gregorian notation, for the ease of use by cantors and the Congregation.

For all these precious liturgical elements, this book may rightly be called a hymnal. However, the chief treasure is the set of `Devotional Hymns'. For these, the best of various hymnary traditions have been chosen. Uniquely beneficial as opening and closing hymns, we hope that these provide spiritual benefits for those who use them, especially since this is the first Orthodox hymnal to include many of these musical pieces.

We hope this Prayer Book Hymnal provides benefit and utility to those living within, or interested in, the English Orthodox tradition.
%\phantomsection
\addcontentsline{toc}{chapter}{Dedication}
\fancyhead[C]{{\LARGE Dedication}}
\noindent
This \textit{Prayer Book Hymnal} is dedicated to Patriarch \textsc{John X}, \nth{171} Successor of St. Peter, Head of the Church, Patriarch of Antioch, and Defender of the Faith.\\
\par\noindent
And we are beyond thankful for our Christ-loving, God-fearing, and faithful-caring chief Bishop and Metropolitan \textsc{Saba}. For his support of the Western Rite continues our endeavours, lightens our hearts, and shows his divine charity for his flock.\\
\par\noindent
And for our faithful, compassionate, and righteous Bishop \textsc{John}, we must send up everlasting praise and thanksgiving unto the Most Holy Trinity. For through him has the daily care, discipline, and concern for the Vicariate been carried. Through him have so many of us flocked for refuge in the Holy Orthodox Catholic Church. And through him have we been intimately cared and nourished.\\
\par\noindent
Finally, we must thank our Lord and God and Saviour Jesus Christ for our beloved Priests, Deacons, and all Clergy in the Vicariate for their tireless care, concern, and work for our salvation. Chiefly among them must surely be our blessed Assistant to the Vicar General, Fr. John Fenton. He who has taken great and serious care and rigour for staying the course of the Vicariate when it is battered and afflicted from within and without.\\
\par\noindent
Let us for ever bless the Father and the Son and the Holy Ghost, together with all the Saints, chiefly of blessed Mary ever-Virgin, blessed John Baptist, and the holy Apostles Peter and Paul, for our fathers in the faith. May God, through the merits and prayers of the Saints, bless and preserve them now and for ever. Amen.
\vfill
\begin{center}
\textsc{And I say also unto thee, That thou art Peter, and upon this rock I will build my church; and the gates of hell shall not prevail against it. And I will give unto thee the keys of the kingdom of heaven: and whatsoever thou shalt bind on earth shall be bound in heaven: and whatsoever thou shalt loose on earth shall be loosed in heaven.}
\end{center}
%\fancyhead[RE,LO]{}\fancyhead[RO,LE]{}
\fancyhead[C]{}
\mainmatter
\thispagestyle{empty}
\phantomsection
\addcontentsline{toc}{chapter}{Daily Office Lectionary}
  \begin{center}
   \includegraphics[width=98mm, height=140mm]{Coronation.eps}
   \par
   \vspace{2ex}
   	\textsc{\Huge{Daily Office Lectionary}}
   \end{center}
%\fancyhead[C]{{\LARGE Lectionary}}\label{lectionary}
\section{Feasts of the Year}
\fancyhead[RO,LE]{\textit{Proper of Saints}}
\begin{longtable}{l c l l}
\multicolumn{1}{l}{Feast}&\multicolumn{1}{l}{Proper Psalms}&\multicolumn{1}{l}{Lesson 1}& \multicolumn{1}{l}{Lesson 2}\\
\hline
\endhead

%

%

%

%

%

%

%

\end{longtable}
%%\fancyhead[RE,LO]{}\fancyhead[RO,LE]{}
\fancyhead[C]{\LARGE Frequent Prayers}
\phantomsection
\addcontentsline{toc}{chapter}{Frequently Said Prayers during Mass}

\section*{Prayers at the Foot of the Altar}

\elcol{℣. In the Name of the Father, {\ding{64}} and of the Son, and of the Holy Ghost.}{℣. In nómine Patris, {\ding{64}} et Fílii, et Spíritus Sancti.}
\elcol{℟. Amen.}{℟. Amen.}

\elcol{℣. I will go unto the altar of God.}{℣. Introíbo ad altáre Dei.}
\elcol{℟. Even unto the God of my joy and gladness.}{℟. Ad Deum, qui l{\ae}tíficat juventútem meam.}
\par\noindent
	\centerline{\rule{0.5\textwidth}{0.4pt}}\\
\par\noindent
\elcol{℣. Give sentence with me, O God, and defend my cause against the ungodly people: O deliver me from the deceitful and wicked man.}{℣. Júdica me, Deus, et discérne causam meam de gente non sancta: ab hómine iníquo et dolóso érue me.}
\elcol{℟. For thou art the God of my strength; why hast thou put me from thee: and why go I so heavily, while the enemy oppresseth me?}{℟. Quia tu es, Deus, fortitúdo mea: quare me repulísti, et quare tristis incédo, dum afflígit me inimícus?}
\elcol{℣. O send out thy light and thy truth, that they may lead me: and bring me unto thy holy hill, and to thy dwelling.}{℣. Emítte lucem tuam et veritátem tuam: ipsa me deduxérunt, et adduxérunt in montem sanctum tuum et in tabernácula tua.}
\elcol{℟. And that I may go unto the altar of God, even the God of my joy and gladness: and upon the harp will I give thanks unto thee, O God, my God.}{℟. Et introíbo ad altáre Dei: ad Deum, qui l{\ae}tíficat juventútem meam.}
\elcol{℣. Why art thou so heavy, O my soul: and why art thou so disquieted within me?}{℣. Confitébor tibi in cíthara, Deus, Deus meus: quare tristis es, ánima mea, et quare contúrbas me?}
\elcol{℟. O put thy trust in God: for I will yet give him thanks, which is the help of my countenance, and my God.}{℟. Spera in Deo, quóniam adhuc confitébor illi: salutáre vultus mei, et Deus meus.}
\elcol{℣. Glory be to the Father, and to the Son: and to the Holy Ghost.}{℣. Glória Patri, et Fílio, et Spirítui Sancto.}
\elcol{℟. As it was in the beginning, is now, and ever shall be: world without end. Amen.}{℟. Sicut erat in princípio, et nunc, et semper: et in sǽcula s{\ae}culórum. Amen.}
\elcol{℣. I will go unto the altar of God:}{℣. Introíbo ad altáre Dei.}
\elcol{℟. Even unto the God of my joy and gladness.}{℟. Ad Deum, qui l{\ae}tíficat juventútem meam.}
\par\noindent
	\centerline{\rule{0.5\textwidth}{0.4pt}}\\
\par\noindent
\clearpage
\elcol{℣. Our help is in the Name {\ding{64}} of the Lord:}{℣. Adjutórium nostrum {\ding{64}} in nómine Dómini.}
\elcol{℟. Who hath made heaven and earth.}{℟. Qui fecit c{\ae}lum et terram.}

\elcol{\lett{I}{confess} to almighty God, to blessed Mary ever Virgin, to blessed Michael the Archangel, to blessed John Baptist, to the holy Apostles Peter and Paul, to all the Saints, and to you, brethren: that I have sinned exceedingly in thought, word, and deed: \inrub{Strike breast thrice.} by my fault, my own fault, my own most grievous fault. Therefore I beg blessed Mary ever Virgin, blessed Michael the Archangel, blessed John Baptist, the holy Apostles Peter and Paul, all the Saints, and you, brethren, to pray for me to the Lord our God.}{\lett{C}{onf\smash{í}teor} Deo omnipoténti, beát{\ae} Marí{\ae} semper Vírgini, beáto Michaéli Archángelo, beáto Joánni Baptíst{\ae}, sanctis Apóstolis Petro et Paulo, ómnibus Sanctis, et vobis, fratres: quia peccávi nimis cogitatióne, verbo et ópere: \inrub{Strike breast thrice.} mea culpa, mea culpa, mea máxima culpa. Ideo precor beátam Maríam semper Vírginem, beátum Michaélem Archángelum, beátum Joánnem Baptístam, sanctos Apóstolos Petrum et Paulum, omnes Sanctos, et vos, fratres, oráre pro me ad Dóminum, Deum nostrum.}
\elcol{℟. Almighty God have mercy upon thee, forgive thee thy sins, and bring thee to everlasting life.}{℟. Misereátur tui omnípotens Deus, et, dimíssis peccátis tuis, perdúcat te ad vitam {\ae}térnam.}
\elcol{℟. Amen.}{℟. Amen.}
\elcol{℟. I confess to almighty God, to blessed Mary ever Virgin, to blessed Michael the Archangel, to blessed John Baptist, to the holy Apostles Peter and Paul, to all the Saints, and to thee, father: that I have sinned exceedingly in thought, word, and deed: \inrub{Strike breast thrice.} by my fault, my own fault, my own most grievous fault. Therefore I beg blessed Mary ever Virgin, blessed Michael the Archangel, blessed John Baptist, the holy Apostles Peter and Paul, all the Saints, and thee, father, to pray for me to the Lord our God.}{℟. Confíteor Deo omnipoténti, beát{\ae} Marí{\ae} semper Vírgini, beáto Michaéli Archángelo, beáto Joánni Baptíst{\ae}, sanctis Apóstolis Petro et Paulo, ómnibus Sanctis, et tibi, pater: quia peccávi nimis cogitatióne, verbo et ópere: \inrub{Strike breast thrice.} mea culpa, mea culpa, mea máxima culpa. Ideo precor beátam Maríam semper Vírginem, beátum Michaélem Archángelum, beátum Joánnem Baptístam, sanctos Apóstolos Petrum et Paulum, omnes Sanctos, et te, pater, oráre pro me ad Dóminum, Deum nostrum.}
\elcol{℣. Almighty God have mercy upon you, forgive you your sins, and bring you to everlasting life.}{℣. Misereátur vestri omnípotens Deus, et, dimíssis peccátis vestris, perdúcat vos ad vitam {\ae}térnam.}
\elcol{℟. Amen.}{℟. Amen.}
\elcol{℣. The almighty and merciful Lord grant unto us pardon, {\ding{64}} absolution, and remission of our sins.\par
℟. Amen.}{℣. Indulgéntiam, {\ding{64}} absolutiónem et remissiónem peccatórum nostrórum tríbuat nobis omnípotens et miséricors Dóminus.\par
℟. Amen.}
\clearpage
\elcol{℣. Wilt thou not turn again and quicken us, O God?}{℣. Deus, tu convérsus vivificábis nos?}
\elcol{℟. That thy people may rejoice in thee.}{℟. Et plebs tua l{\ae}tábitur in te.}
\elcol{℣. O Lord, show thy mercy upon us.}{℣. Osténde nobis, Dómine, misericórdiam tuam.}
\elcol{℟. And grant us thy salvation.}{℟. Et salutáre tuum da nobis.}
\elcol{℣. O Lord, hear my prayer.}{℣. Dómine, exáudi oratiónem meam.}
\elcol{℟. And let my cry come unto thee.}{℟. Et clamor meus ad te véniat.}
\elcol{℣. The Lord be with you.}{℣. Dóminus vobíscum.}
\elcol{℟. And with thy spirit.}{℟. Et cum spíritu tuo.}
\elcol{℣. Let us pray.}{℣. Orémus.}

\section*{Collect for Purity}
\lett{A}{lmighty} God, unto whom all hearts are open, all desires known, and from whom no secrets are hid; Cleanse the thoughts of our hearts by the inspiration of thy Holy Spirit, that we may perfectly love thee, and worthily magnify thy holy Name; through Christ our Lord. \textit{Amen.}
\clearpage
\section*{Ten Commandments}
\begin{center}
	\textsc{God spake these words, and said:}
\end{center}
%\begin{multicols}{2}
\par\noindent
    I am the \textsc{Lord} thy God; Thou shalt have none other gods but me.\par
    \textit{Lord, have mercy upon us, and incline our hearts to keep this law.}
    \par\noindent
    Thou shalt not make to thyself any graven image, nor the likeness of any thing that is in heaven above, or in the earth beneath, or in the water under the earth; thou shalt not bow down to them, nor worship them:
    \par\noindent
    \leftskip=2em
	{\small{for I the Lord thy God am a jealous God, and visit the sins of the fathers upon the children, unto the third and fourth generation of them that hate me; and show mercy unto thousands in them that love me and keep my commandments.}}
	\par
	\leftskip=0cm
    \textit{Lord, have mercy upon us, and incline our hearts to keep this law.}
    \par\noindent
    Thou shalt not take the Name of the Lord thy God in vain;
    \par\noindent
    \leftskip=2em
	{\small{for the Lord will not hold him guiltless, that taketh his Name in vain.}}
	\par
	\leftskip=0em
    \textit{Lord, have mercy upon us, and incline our hearts to keep this law.}
    \par\noindent
    Remember that thou keep holy the Sabbath-day.
    \par\noindent
    \leftskip=2em
	{\small{Six days shalt thou labour, and do all that thou hast to do; but the seventh day is the Sabbath of the Lord thy God. In it thou shalt do no manner of work; thou, and thy son, and thy daughter, thy man-servant, and thy maid-servant, thy cattle, and the stranger that is within thy gates. For in six days the Lord made heaven and earth, the sea, and all that in them is, and rested the seventh day: wherefore the Lord blessed the seventh day, and hallowed it.}}
	\par
	\leftskip=0em
    \textit{Lord, have mercy upon us, and incline our hearts to keep this law.}
    \par\noindent
    Honour thy father and thy mother;
    \par\noindent
    \leftskip=2em
	{\small{that thy days may be long in the land which the Lord thy God giveth thee.}}
	\par
	\leftskip=0em
    \textit{Lord, have mercy upon us, and incline our hearts to keep this law.}
    \par\noindent
    Thou shalt do no murder.\par
    \textit{Lord, have mercy upon us, and incline our hearts to keep this law.}
    \par\noindent
    Thou shalt not commit adultery.\par
    \textit{Lord, have mercy upon us, and incline our hearts to keep this law.}
\par\noindent
    Thou shalt not steal.\par
    \textit{Lord, have mercy upon us, and incline our hearts to keep this law.}
\par\noindent
    Thou shalt not bear false witness against thy neighbour.\par
    \textit{Lord, have mercy upon us, and incline our hearts to keep this law.}
    \par\noindent
    Thou shalt not covet.
    \par\noindent
    \leftskip=2em
	{\small{thy neighbour's house, thou shalt not covet thy neighbour's wife, nor his servant, nor his maid, nor his ox, nor his ass, nor any thing that is his.}}
	\par
	\leftskip=0em
	\textit{Lord, have mercy upon us, and write all these thy laws in our hearts, we beseech thee.}
\clearpage
\section*{Summary of the Law}
\begin{center}
	{\textsc{Hear what our Lord Jesus Christ saith.}}
\end{center}

\lett{T}{hou} shalt love the Lord thy God with all thy heart, and with all thy soul, and with all thy mind. This is the first and great commandment. And the second is like unto it; Thou shalt love thy neighbour as thyself. On these two commandments hang all the Law and the Prophets.

\section*{Invitation to Confession}
\lett{Y}{e} who do truly and earnestly repent you of your sins, and are in love and charity with your neighbours, and intend to lead a new life, following the commandments of God, and walking from henceforth in his holy ways; Draw near with faith, and take this holy Sacrament to your comfort; and make your humble confession to Almighty God, devoutly kneeling.

\section*{Confession}
\lett{A}{lmighty} God, Father of our Lord Jesus Christ, Maker of all things, Judge of all men; We acknowledge and bewail our manifold sins and wickedness, Which we, from time to time, most grievously have committed, By thought, word, and deed, Against thy Divine Majesty, Provoking most justly thy wrath and indignation against us. We do earnestly repent, And are heartily sorry for these our misdoings; The remembrance of them is grievous unto us; The burden of them is intolerable. Have mercy upon us, Have mercy upon us, most merciful Father; For thy Son our Lord Jesus Christ's sake, Forgive us all that is past; And grant that we may ever hereafter Serve and please thee in newness of life, To the honour and glory of thy Name; Through Jesus Christ our Lord. Amen.

\section*{Absolution}
\lett{A}{lmighty} God, our heavenly Father, who of his great mercy hath promised forgiveness of sins to all those who with hearty repentance and true faith turn unto him; Have mercy upon you; pardon {\ding{64}} and deliver you from all your sins; confirm and strengthen you in all goodness; and bring you to everlasting life; through Jesus Christ our Lord.

℟. Amen.
\clearpage
\section*{Comfortable Words}
\begin{center}
	\textsc{Hear what comfortable words our Saviour Christ saith unto all who truly turn to him.}
\end{center}
\par\noindent
Come unto me, all ye that travail and are heavy laden, and I will refresh you. \vr{Matt. 11:28}
\par\noindent
    So God loved the world, that he gave his only-begotten Son, to the end that all that believe in him should not perish, but have everlasting life. \vr{John 3:16}
    \par\noindent
    \begin{center}
		\textsc{Hear also what Saint Paul saith.}
	\end{center}
\par\noindent
    This is a true saying, and worthy of all men to be received, That Christ Jesus came into the world to save sinners. \vr{1 Tim. 1:15}
\par\noindent
    \begin{center}
		\textsc{Hear also what Saint John saith.}
	\end{center}
    \par\noindent
    If any man sin, we have an Advocate with the Father, Jesus Christ the righteous; and he is the Propitiation for our sins. \vr{1 John 2:1-2}

\section*{Distribution of the Body}
\lett{T}{he} Body of our Lord Jesus Christ, which was given for thee, preserve thy body and soul unto everlasting life. Take and eat this in remembrance that Christ died for thee, and feed on him in thy heart by faith, with thanksgiving.

\section*{Distribution of the Blood}
\lett{T}{he} Blood of our Lord Jesus Christ, which was shed for thee, preserve thy body and soul unto everlasting life. Drink this in remembrance that Christ's Blood was shed for thee, and be thankful.

\section*{Distribution of the Body \& Blood}
\lett{T}{he} Body of our Lord Jesus Christ which was given for thee, and his Blood which was shed for thee. Take them in remembrance that Christ died for thee, and feed on him in thy heart by faith with thanksgiving.
%\fancyhead[RE,LO]{}\fancyhead[RO,LE]{}
\fancyhead[C]{}\thispagestyle{empty}
\phantomsection
\addcontentsline{toc}{chapter}{Daily Office}
  \begin{center}
   \includegraphics[width=98mm, height=140mm]{Resurrection.eps}
   \par
   \vspace{2ex}
   	\textsc{\Huge{Daily Office}}
   \end{center}
%\fancyhead[RO,LE]{\textit{Alternative Psalter}}
\fancyhead[RE,LO]{}
\fancyhead[C]{\LARGE Daily Office}
\section{Alternative Psalter Arrangement}
\begin{secrubric}
	For those who also pray the Minor Hours, it may be desired to have a more robust cursus of the Psalms during Mattins and Evensong, which do not overlap with the Psalms of the Minor Hours. The following arrangement is provided for such a situation.
\end{secrubric}
\begin{rubric}
		\textsc{Note,} The proper Psalms for Feast Days and obligatory Sundays take priority.
\end{rubric}
\begin{rubric}
	\textsc{Note,} The First Sunday of Advent \& Palm Sunday are always in `Week 1'.
\end{rubric}

\begin{paracol}{2}%\sloppy
	\begin{leftcolumn}
		\subby{Week I}
		\begin{inhead}
			Sunday
		\end{inhead}\par\noindent
		\textit{Mattins.} Psalms 63, 66, 67, 93, 96, 97.\par\noindent
		\textit{Evensong.} Psalms 84, 85, 104, 119:33-80.\par
		\begin{inhead}
			Monday
		\end{inhead}\par\noindent
		\textit{Mattins.} Psalms 1, 2, 3, 5, 6, 7, 8.\par\noindent
		\textit{Evensong}. Psalms 69, 70, 71, 72, 119:81-112.\par
		\begin{inhead}
			Tuesday
		\end{inhead}\par\noindent
		\textit{Mattins.} Psalms 9, 10, 11, 12, 13, 14.\par\noindent
		\textit{Evensong.} Psalms 73, 74, 75, 76, 77.\par
		\begin{inhead}
			Wednesday
		\end{inhead}\par\noindent
		\textit{Mattins.} Psalms 15, 16, 17, 18.\par\noindent
		\textit{Evensong.} Psalms 78, 79, 80, 81.\par
		\begin{inhead}
			Thursday
		\end{inhead}\par\noindent
		\textit{Mattins.} Psalms 19, 20, 21, 22, 23.\par\noindent
		\textit{Evensong.} Psalms 82, 83, 86, 87, 88.\par
		\begin{inhead}
			Friday
		\end{inhead}\par\noindent
		\textit{Mattins.} Psalms 24, 25, 26, 27, 28, 29.\par\noindent
		\textit{Evensong.} Psalms 89, 90, 92, 94.\par
		\begin{inhead}
			Saturday
		\end{inhead}\par\noindent
		\textit{Mattins.} Psalms 30, 31, 32, 33, 34.\par\noindent
		\textit{Evensong.} Psalms 101, 102, 103, 105.\par
		
	\end{leftcolumn}
	\begin{rightcolumn}
		\subby{Week II}
		\begin{inhead}
			Sunday
		\end{inhead}\par\noindent
		\textit{Mattins.} Psalms 98, 99, 100, 148, 149, 150.\par\noindent
		\textit{Evensong.} Psalms 110, 111, 112, 113, 114, 115.\par
		\begin{inhead}
			Monday
		\end{inhead}\par\noindent
		\textit{Mattins.} Psalms 35, 36, 37, 38.\par\noindent
		\textit{Evensong}. Psalms 106, 107, 108, 109.\par
		\begin{inhead}
			Tuesday
		\end{inhead}\par\noindent
		\textit{Mattins.} Psalms 39, 40, 41, 42, 43.\par\noindent
		\textit{Evensong.} Psalms 116, 117, 118, 119:113-144.\par
		\begin{inhead}
			Wednesday
		\end{inhead}\par\noindent
		\textit{Mattins.} Psalms 44, 45, 46, 47, 48, 49.\par\noindent
		\textit{Evensong.} Psalms 119:145-176, 129, 130, 131.\par
		\begin{inhead}
			Thursday
		\end{inhead}\par\noindent
		\textit{Mattins.} Psalms 50, 51, 52, 53, 55.\par\noindent
		\textit{Evensong.} Psalms 132, 133, 135, 136, 137, 138.\par
		\begin{inhead}
			Friday
		\end{inhead}\par\noindent
		\textit{Mattins.} Psalms 56, 57, 58, 59, 60, 61.\par\noindent
		\textit{Evensong.} Psalms 139, 140, 141, 142, 143.\par
		\begin{inhead}
			Saturday
		\end{inhead}\par\noindent
		\textit{Mattins.} Psalms 62, 64, 65, 68.\par\noindent
		\textit{Evensong.} Psalms 144, 145, 146, 147.\par
	\end{rightcolumn}
	\switchcolumn*
	%\fussy
\end{paracol}

%\fancyhead[RO,LE]{\textit{Terce}}
\fancyhead[RE,LO]{}
\section{Hour of Terce}
\elcol{℣. O God, {\ding{64}} make speed to save us.}{℣. Deus, {\ding{64}} in adjutórium nostrum inténde.}
\elcol{℟. O Lord, make haste to help us.}{℟. Dómine, ad adjuvándum nos festína.}
\elcol{℣. Glory be to the Father, and to the Son, and to the Holy Ghost.}{℣. Glória Patri, et Fílio, * et Spirítui Sancto:}
\elcol{℟. As it was in the beginning, is now, and ever shall be, world without end. Amen.}{℟. Sicut erat in princípio, et nunc, et semper, * et in sǽcula s{\ae}culórum. Amen.}
\elcol{℣. Praise ye the Lord.}{℣. Laudáte Dominum.}
\elcol{℟. The Lord's Name be praised.}{℟. Sit Nomen Dómini Benedíctum.}
\subby{Nunc Sancte}
\elcol{\begin{hangparas}{1.25em}{1}
\lett{C}{ome} Holy Ghost, with God the Son,\par
\secondline{And God the Father, ever one:}\par\noindent
%\thirdline{Shed forth thy grace within our breast,}
Shed forth thy grace within our breast,\par\noindent
And dwell with us, a ready Guest.\\
\par\noindent
By every power, by heart and tongue,\par\noindent
By act and deed, thy praise be sung;\par\noindent
Inflame with perfect love each sense,\par\noindent
That others' souls may kindle thence.\\
\par\noindent
O Father, that we ask be done,\par\noindent
Through Jesus Christ, thine only Son;\par\noindent
Who, with the Holy Ghost and thee,\par\noindent
Shall live and reign eternally. Amen.\\
\end{hangparas}}
{\lett{N}{unc,} Sancte, nobis, Spíritus,\par
\secondline{Unum Patri cum Fílio,}
\thirdline{Dignáre promptus íngeri}
Nostro refúsus péctori.\\

Os, lingua, mens, sensus, vigor\par
Confessiónem pérsonent.\par
Flamméscat igne cáritas,\par
Accéndat ardor próximos.\\

Præsta, Pater piíssime,\par
Patríque compar Únice,\par
Cum Spíritu Paráclito\par
Regnans per omne sǽculum. Amen.\\}


\subbypsalm{Psalm 120. \textit{Ad Dominum}}

\lett{W}{hen} I was in trouble I called upon the Lord : and he heard me.\par
\secondline{\psanum{2}Deliver my soul, O Lord, from lying lips : and from a deceitful tongue.}
\thirdline{\psanum{3}What reward shall be given or done unto thee, thou false tongue : even mighty and sharp arrows, with hot burning coals.}
\psanum{4}Woe is me, that I am constrained to dwell with Mesech : and to have my habitation among the tents of Kedar.\par
\psanum{5}My soul hath long dwelt among them : that are enemies unto peace.\par
\psanum{6}I labour for peace, but when I speak unto them thereof : they make them ready to battle.\par
℣. Glory be to the Father, and to the Son, and to the Holy Ghost.\par
℟. As it was in the beginning, is now, and ever shall be, world without end. Amen.

\subbypsalm{Psalm 121. \textit{Levavi oculus}}

\lett{I}{will} lift up mine eyes unto the hills : from whence cometh my help.\par
\secondline{\psanum{2}My help cometh even from the Lord : who hath made heaven and earth.}
\thirdline{\psanum{3}He will not suffer thy foot to be moved : and he that keepeth thee will not sleep.}
\psanum{4}Behold, he that keepeth Israel : shall neither slumber nor sleep.\par
\psanum{5}The Lord himself is thy keeper : the Lord is thy defence upon thy right hand;\par
\psanum{6}So that the sun shall not burn thee by day : neither the moon by night.\par
\psanum{7}The Lord shall preserve thee from all evil : yea, it is even he that shall keep thy soul.\par
\psanum{8}The Lord shall preserve thy going out, and thy coming in : from this time forth for evermore.\par
℣. Glory be to the Father, and to the Son, and to the Holy Ghost.\par
℟. As it was in the beginning, is now, and ever shall be, world without end. Amen.

\subbypsalm{Psalm 122. \textit{Laetatus sum}}

\lett{I}{was} glad when they said unto me : We will go into the house of the Lord.\par
\secondline{\psanum{2}Our feet shall stand in thy gates : O Jerusalem.}
\thirdline{\psanum{3}Jerusalem is built as a city : that is at unity in itself.}
\psanum{4}For thither the tribes go up, even the tribes of the Lord : to testify unto Israel, to give thanks unto the Name of the Lord.\par
\psanum{5}For there is the seat of judgement : even the seat of the house of David.\par
\psanum{6}O pray for the peace of Jerusalem : they shall prosper that love thee.\par
\psanum{7}Peace be within thy walls : and plenteousness within thy palaces.\par
\psanum{8}For my brethren and companions' sakes : I will wish thee prosperity.\par
\psanum{9}Yea, because of the house of the Lord our God : I will seek to do thee good.\par
℣. Glory be to the Father, and to the Son, and to the Holy Ghost.\par
℟. As it was in the beginning, is now, and ever shall be, world without end. Amen.

\subby{Chapter}
\begin{inhead}
	Feria Days Throughout the Year
\end{inhead}
\elcoldent{Heal me, O Lord, and I shall be healed; save me, and I shall be saved: for thou art my praise.}{Sana me, Dómine, et sanábor: salvum me fac, et salvus ero: quóniam laus mea tu es.}
\begin{inhead}
	Sundays \& Feast Days
\end{inhead}
\elcoldent{O the depth of the riches both of the wisdom and knowledge of God! How unsearchable are his judgements, and his ways past finding out!}{O altitúdo divitiárum sapiénti{\ae}, et sciénti{\ae} Dei: quam incomprehensibília sunt judícia ejus, et investigábiles vi{\ae} ejus!}
%MANUAL ADJUSTMENT:
\clearpage
\begin{inhead}
	Eastertide
\end{inhead}
\elcoldent{For whatsoever is born of God overcometh the world: and this is the victory that overcometh the world, even our faith.}{Omne, quod natum est ex Deo, vincit mundum: et h{\ae}c est victória, qu{\ae} vincit mundum, fides nostra.}
\elcol{℟. Thanks be to God.}{℟. Deo grátias.}
\subby{Short Respond}
\begin{inhead}
	Ferias Throughout the Year
\end{inhead}
\elcol{℣. Heal my soul, for I have sinned against thee.}{℣. Sana ánimam meam, Quia peccávi tibi.}

\elcol{℟. Heal my soul, for I have sinned against thee.}{℟. Sana ánimam meam, Quia peccávi tibi.}

\elcol{℣. I said, Lord, be merciful unto me.}{℣. Ego dixi: Dómine, miserére mei.}
    
\elcol{℟. For I have sinned against thee.}{℟. Quia peccávi tibi.}

\elcol{℣. Glory be to the Father, and to the Son, and to the Holy Ghost.}{℣. Glória Patri, et Fílio, et Spirítui Sancto.}
    
\elcol{℟. Heal my soul, for I have sinned against thee.\\}{℟. Sana ánimam meam, Quia peccávi tibi.\\}

\elcol{℣. Thou hast been my succour.}{℣. Adjútor meus, esto.}

\elcol{℟. Leave me not, neither forsake me, O God of my salvation.}{℟. Ne derelínquas me, Neque despícias me, Deus, salutáris meus.}

\begin{inhead}
	Sundays, Feasts, \& in Eastertide
\end{inhead}

\elcol{℣. Incline my heart, O God, unto thy testimonies (alleluia, alleluia).}{℣.  Inclína cor meum in testimónia tua (allelúja, allelúja).}

\elcol{℟. Incline my heart, O God, unto thy testimonies (alleluia, alleluia).}{℟. Inclína cor meum in testimónia tua (allelúja, allelúja).}

\elcol{℣. Turn away mine eyes, lest they behold vanity.}{℣. Avérte óculos meos ne vídeant vanitátem.}
    
\elcol{℟. Unto thy testimonies (alleluia, alleluia).}{℟. In testimónia tua (allelúja, allelúja).}

\elcol{℣. Glory be to the Father, and to the Son, and to the Holy Ghost.}{℣. Glória Patri, et Fílio, et Spirítui Sancto.}
    
\elcol{℟. Incline my heart, O God, unto thy testimonies (alleluia, alleluia).\\}{℟. Inclína cor meum in testimónia tua (allelúja, allelúja).\\}

\elcol{℣. I said, Lord, be merciful unto me (alleluia).}{℣. Ego dixi: Dómine, miserére mei (allelúja).}

\elcol{℟. For I have sinned against thee (alleluia).}{℟. Quia peccávi tibi (allelúja).}

\subby{Suffrages}
\begin{rubric}
	The Suffrages, said kneeling, are said on all Ferias outside of Eastertide. They are omitted on Sundays, Feasts, and within Principle, Primary, and Secondary Octaves.\par
	\textsc{Note,} If omitted, the Office continues with the Collect.
\end{rubric}

\elcol{℣. Lord, have mercy upon us.}{℣. Kýrie, eléison.}
\elcol{℟. Christ, have mercy upon us.}{℟. Christe, eléison.}
\elcol{℣. Lord, have mercy upon us.}{℣. Kýrie, eléison.}
%\subsec{Lord's Prayer}
\needspace{4\baselineskip}
\elcol{\lett{O}{ur} Father, who art in heaven, Hallowed be thy Name. Thy kingdom come. Thy will be done on earth, As it is in heaven. Give us this day our daily bread. And forgive us our trespasses, As we forgive those who trespass against us. And lead us not into temptation; But deliver us from evil. Amen.}
{\lett{P}{ater} noster, qui es in c{\ae}lis, sanctificétur nomen tuum: advéniat regnum tuum: fiat volúntas tua, sicut in c{\ae}lo et in terra. Panem nostrum quotidiánum da nobis hódie: et dimítte nobis débita nostra, sicut et nos dimíttimus debitóribus nostris: et ne nos indúcas in tentatiónem: sed líbera nos a malo. Amen.}

\elcol{℣. Send out thy light and thy truth, that they may lead me.}{℣. Emítte lucem tuam et veritátem tuam: ipsa me deduxérunt.}

\elcol{℟. And bring me unto thy holy hill, and to thy dwelling.}{℟. Et adduxérunt in montem sanctum tuum et in tabernácula tua.}

\elcol{℣. O let us live.}{℣. Vivificábis nos.}

\elcol{℟. And we shall call upon thy Name.}{℟. Et nomen tuum invocábimus.}

\elcol{℣. Turn us again, O Lord God of hosts.}{℣. Dómine, Deus virtútum, convérte nos.}

\elcol{℟. Show the light of thy countenance, and we shall be whole.}{℟. Et osténde fáciem tuam, et salvi érimus.}

\subby{Collects}
\elcol{℣. The Lord be with you.\par
\textit{or,} O Lord, hear our prayer.}{℣. Dóminus vobíscum.\par
\textit{vel,} Dómine, exáudi oratiónem nostram.}
\elcol{℟. And with thy spirit.\par
\textit{or,} And let our cry come unto thee.}{℟. Et cum spíritu tuo.\par
\textit{vel,} Et clamor noster ad te véniat.}
\elcol{℣. Let us pray.}{℣. Orémus.}
\begin{rubric}
	The Collect(s) of the Day is here said, followed by this memorial Collect.
\end{rubric}
\needspace{4\baselineskip}
\lett{H}{elp} us this day, O God, to serve thee devoutly and the world busily. May we do our work wisely, give succour secretly, go to our meat appetitely, sit thereat discreetly, arise temperately, please our friends duly, go to our bed merrily, and sleep surely; for the joy of our Lord Jesus Christ. \textit{Amen.}\\
%MANUAL ADJUSTMENT:
\clearpage
\elcol{℣. The Lord be with you.\par
\textit{or,} O Lord, hear our prayer.}{℣. Dóminus vobíscum.\par
\textit{vel,} Dómine, exáudi oratiónem nostram.}
\elcol{℟. And with thy spirit.\par
\textit{or,} And let our cry come unto thee.}{℟. Et cum spíritu tuo.\par
\textit{vel,} Et clamor noster ad te véniat.}
\elcol{℣. Let us bless the Lord.}{℣. Benedicámus Dómino.}
\elcol{℟. Thanks be to God.}{℟. Deo grátias}
\elcol{℣. May the souls {\ding{64}} of the faithful departed, through the mercy of God, rest in peace.}{℣. Fidélium ánim{\ae} {\ding{64}} per misericórdiam Dei requiéscant in pace.}
\needspace{4\baselineskip}
\lett{U}{nto} him that loved us, and washed us from our sins in his own blood, and hath made us kings and priests unto God and his Father, to him be glory and dominion for ever and ever. \textit{Amen.}


%\elcol{\lett{U}{nto} him that loved us, and washed us from our sins in his own blood, and hath made us kings and priests unto God and his Father, to him be glory and dominion for ever and ever. Amen.}{\lett{Q}{ui} dilexit nos, et lavit nos a peccatis nostris in sanguine suo, 6 et fecit nos regnum, et sacerdotes Deo et Patri suo: ipsi gloria et imperium in s{\ae}cula s{\ae}culorum. Amen.}
\begin{center}
	{\textsc{Here endeth the Order of Terce.}}
\end{center}
%\fancyhead[RO,LE]{\textit{Sext}}
\fancyhead[RE,LO]{}
\section{Hour of Sext}
\elcol{℣. O God, {\ding{64}} make speed to save us.}{℣. Deus, {\ding{64}} in adjutórium nostrum inténde.}
\elcol{℟. O Lord, make haste to help us.}{℟. Dómine, ad adjuvándum nos festína.}
\elcol{℣. Glory be to the Father, and to the Son, and to the Holy Ghost.}{℣. Glória Patri, et Fílio, * et Spirítui Sancto:}
\elcol{℟. As it was in the beginning, is now, and ever shall be, world without end. Amen.}{℟. Sicut erat in princípio, et nunc, et semper, * et in sǽcula s{\ae}culórum. Amen.}
\elcol{℣. Praise ye the Lord.}{℣. Laudáte Dominum.}
\elcol{℟. The Lord's Name be praised.}{℟. Sit Nomen Dómini Benedíctum.}
\subby{Rector Potens}
\elcol{\begin{hangparas}{1.25em}{1}
\lett{O}{God} of truth, O Lord of might,\par
\secondline{Who ord'rest time and change aright,}
\par\noindent
And send'st the early morning ray,\par\noindent
And light'st the glow of perfect day.\\
\par\noindent
Extinguish thou each sinful fire,\par\noindent
And banish every ill desire:\par\noindent
And while thou keep'st the body whole,\par\noindent
Shed forth thy peace upon the soul.\\
\par\noindent
O Father, that we ask be done,\par\noindent
Through Jesus Christ, thine only Son;\par\noindent
Who, with the Holy Ghost and thee,\par\noindent
Doth live and reign eternally. Amen.\\
\end{hangparas}}
{\lett{R}{ector} potens, verax Deus,\par
\secondline{Qui témperas rerum vices,}
\thirdline{Splendóre mane illúminas,}
Et ígnibus merídiem:\\
\par
Exstíngue flammas lítium,\par
Aufer calórem nóxium,\par
Confer salútem córporum,\par
Verámque pacem córdium.\\
\par
Pr{\ae}sta, Pater piíssime,\par
Patríque compar Únice,\par
Cum Spíritu Paráclito\par
Regnans per omne sǽculum. Amen.\\}


\subbypsalm{Psalm 123. \textit{Ad te levavi oculos meos}}
\lett{U}{nto} thee lift I up mine eyes : O thou that dwellest in the heavens.\par
\secondline{\psanum{2}Behold, even as the eyes of servants look unto the hand of their masters, and as the eyes of a maiden unto the hand of her mistress : even}
\noindent
so our eyes wait upon the Lord our God, until he have mercy upon us.\par
\psanum{3}Have mercy upon us, O Lord, have mercy upon us : for we are utterly despised.\par
\psanum{4}Our soul is filled with the scornful reproof of the wealthy : and with the despitefulness of the proud.\par
℣. Glory be to the Father, and to the Son, and to the Holy Ghost.\par
℟. As it was in the beginning, is now, and ever shall be, world without end. Amen.

\subbypsalm{Psalm 124. \textit{Nisi quia Dominus}}
\lett{I}{f} the Lord himself had not been on our side, now may Israel say : if the Lord himself had not been on our side, when men rose up against us;\par
\secondline{\psanum{2}They had swallowed us up quick : when they were so wrathfully displeased at us.}
\thirdline{\psanum{3}Yea, the waters had drowned us : and the stream had gone over our soul.}
\psanum{4}The deep waters of the proud : had gone even over our soul.\par
\psanum{5}But praised be the Lord : who hath not given us over for a prey unto their teeth.\par
\psanum{6}Our soul is escaped even as a bird out of the snare of the fowler : the snare is broken, and we are delivered.\par
\psanum{7}Our help standeth in the Name of the Lord : who hath made heaven and earth.\par
℣. Glory be to the Father, and to the Son, and to the Holy Ghost.\par
℟. As it was in the beginning, is now, and ever shall be, world without end. Amen.

\subbypsalm{Psalm 125. \textit{Qui confidunt}}
\lett{T}{hey} that put their trust in the Lord shall be even as the mount Sion : which may not be removed, but standeth fast for ever.\par
\secondline{\psanum{2}The hills stand about Jerusalem : even so standeth the Lord round about his people, from this time forth for evermore.}
\thirdline{\psanum{3}For the rod of the ungodly cometh not into the lot of the righteous : lest the righteous put their hand unto wickedness.}
\psanum{4}Do well, O Lord : unto those that are good and true of heart.\par
\psanum{5}As for such as turn back unto their own wickedness : the Lord shall lead them forth with the evil-doers; but peace shall be upon Israel.\par
℣. Glory be to the Father, and to the Son, and to the Holy Ghost.\par
℟. As it was in the beginning, is now, and ever shall be, world without end. Amen.

\subby{Chapter}
\begin{inhead}
	Ferias Throughout the Year
\end{inhead}
\elcoldent{Prove all things, hold fast to that which is good: abstain from all appearance of evil.}{Omnia autem probáte: quod bonum est tenéte. Ab omni spécie mala abstinéte vos.}
\begin{inhead}
	Sundays \& Feast Days
\end{inhead}
\elcoldent{There are three that bear record in heaven: The Father, the Word, and the Holy Ghost: and these three are one.}{Quóniam tres sunt, qui testimónium dant in c{\ae}lo: Pater, Verbum, et Spíritus Sanctus: et hi tres unum sunt.}
\begin{inhead}
	Eastertide
\end{inhead}
\elcoldent{Who is he that overcometh the world, but he that believeth that Jesus is the Son of God? This is he that came by water and blood, even Jesus Christ.}{Quis est, qui vincit mundum, nisi qui credit quóniam Jesus est Fílius Dei? Hic est, qui venit per aquam et sánguinem, Jesus Christus.}
\elcol{℟. Thanks be to God.}{℟. Deo grátias.}
\subby{Short Respond}
\begin{inhead}
	Ferias Throughout the Year
\end{inhead}
\elcol{℣. I will bless the Lord at all times.}{℣. Benedícam Dóminum in omni témpore.}

\elcol{℟. I will bless the Lord at all times.}{℟. Benedícam Dóminum in omni témpore.}

\elcol{℣. His praise shall ever be in my mouth.}{℣. Semper laus ejus in ore meo.}
    
\elcol{℟. At all times.}{℟. In omni témpore.}

\elcol{℣. Glory be to the Father, and to the Son, and to the Holy Ghost.}{℣. Glória Patri, et Fílio, et Spirítui Sancto.}
    
\elcol{℟. I will bless the Lord at all times.\\}{℟. Benedícam Dóminum in omni témpore.\\}

\elcol{℣. Behold, God is my helper.}{℣. Ecce enim, Deus ádjuvat me.}

\elcol{℟. The Lord is with them that uphold my soul.}{℟. Et Dóminus suscéptor est ánim{\ae} me{\ae}.}

\begin{inhead}
	Sundays, Feasts, \& in Eastertide
\end{inhead}

\elcol{℣. O Lord, thy Word endureth for ever in heaven (alleluia, alleluia).}{℣. In {\ae}térnum, Dómine, * verbum tuum pérmanet in c{\ae}lo (allelúja, allelúja).}

\elcol{℟. O Lord, thy Word endureth for ever in heaven (alleluia, alleluia).}{℟. In {\ae}térnum, Dómine, * verbum tuum pérmanet in c{\ae}lo (allelúja, allelúja).}

\elcol{℣. Thy truth also remaineth from one generation to another.}{℣. In generatiónem et generatiónem véritas tua.}

\elcol{℟. For ever in heaven (alleluia, alleluia).}{℟. Pérmanet in c{\ae}lo (allelúja, allelúja).}

\elcol{℣. Glory be to the Father, and to the Son, and to the Holy Ghost.}{℣. Glória Patri, et Fílio, et Spirítui Sancto.}
    
\elcol{℟. O Lord, thy Word endureth for ever in heaven (alleluia, alleluia).}{℟. In {\ae}térnum, Dómine, * verbum tuum pérmanet in c{\ae}lo (allelúja, allelúja).}

\elcol{℣. The Lord is my shepherd, therefore can I lack nothing (alleluia).}{℣. Dóminus regit me, et nihil mihi déerit.}

\elcol{℟. He shall feed me in a green pasture (alleluia).}{℟. In loco páscu{\ae} ibi me collocávit (allelúja).}

\subby{Suffrages}
\begin{rubric}
	The Suffrages, said kneeling, are said on all Ferias outside of Eastertide. They are omitted on Sundays, Feasts, and within Principle, Primary, and Secondary Octaves.\par
	\textsc{Note,} If omitted, the Office continues with the Collect.
\end{rubric}

\elcol{℣. Lord, have mercy upon us.}{℣. Kýrie, eléison.}
\elcol{℟. Christ, have mercy upon us.}{℟. Christe, eléison.}
\elcol{℣. Lord, have mercy upon us.}{℣. Kýrie, eléison.}
%\subsec{Lord's Prayer}
\needspace{4\baselineskip}
\elcol{\lett{O}{ur} Father, who art in heaven, Hallowed be thy Name. Thy kingdom come. Thy will be done on earth, As it is in heaven. Give us this day our daily bread. And forgive us our trespasses, As we forgive those who trespass against us. And lead us not into temptation; But deliver us from evil. Amen.}
{\lett{P}{ater} noster, qui es in c{\ae}lis, sanctificétur nomen tuum: advéniat regnum tuum: fiat volúntas tua, sicut in c{\ae}lo et in terra. Panem nostrum quotidiánum da nobis hódie: et dimítte nobis débita nostra, sicut et nos dimíttimus debitóribus nostris: et ne nos indúcas in tentatiónem: sed líbera nos a malo. Amen.}

\elcol{℣. Hide not thy face from thy servant, for I am in trouble.}{℣. Et ne avértas fáciem tuam a púero tuo: quóniam tríbulor.}

\elcol{℟. O haste thee and hear me.}{℟. Velóciter exáudi me.}

\elcol{℣. Draw nigh unto my soul and save it.}{℣. Inténde ánim{\ae} me{\ae}, et líbera eam.}

\elcol{℟. O deliver me, because of mine enemies.}{℟. Propter inimícos meos éripe me.}

\elcol{℣. Arise, O Christ, and help us.}{℣. Exsúrge, Christe, ádjuva nos.}

\elcol{℟. And deliver us for thy Name's sake.}{℟. Et líbera nos propter nomen tuum.}

\subby{Collects}
\elcol{℣. The Lord be with you.\par
\textit{or,} O Lord, hear our prayer.}{℣. Dóminus vobíscum.\par
\textit{vel,} Dómine, exáudi oratiónem nostram.}
\elcol{℟. And with thy spirit.\par
\textit{or,} And let our cry come unto thee.}{℟. Et cum spíritu tuo.\par
\textit{vel,} Et clamor noster ad te véniat.}
\elcol{℣. Let us pray.}{℣. Orémus.}
\begin{rubric}
	The Collect(s) of the Day is here said, followed by this memorial Collect.
\end{rubric}
\needspace{4\baselineskip}
\lett{L}{ord} Jesus, who didst stretch out thine arms of love on the hard wood of the Cross, that all men might come within the reach of thy saving embrace; clothe us in thy Spirit, that we, stretching forth our hands in loving labour for others, may bring those who know thee not to the knowledge and love of thee, who with the Father and the Holy Ghost livest and reignest ever one God, world without end. \textit{Amen.}\\

\elcol{℣. The Lord be with you.\par
\textit{or,} O Lord, hear our prayer.}{℣. Dóminus vobíscum.\par
\textit{vel,} Dómine, exáudi oratiónem nostram.}
\elcol{℟. And with thy spirit.\par
\textit{or,} And let our cry come unto thee.}{℟. Et cum spíritu tuo.\par
\textit{vel,} Et clamor noster ad te véniat.}
\elcol{℣. Let us bless the Lord.}{℣. Benedicámus Dómino.}
\elcol{℟. Thanks be to God.}{℟. Deo grátias}
\elcol{℣. May the souls {\ding{64}} of the faithful departed, through the mercy of God, rest in peace.}{℣. Fidélium ánim{\ae} {\ding{64}} per misericórdiam Dei requiéscant in pace.}
\needspace{4\baselineskip}
\lett{N}{ow} unto the King eternal, immortal, invisible, the only wise God, be honour and glory, for ever and ever. \textit{Amen.}


\begin{center}
	{\textsc{Here endeth the Order of Sext.}}
\end{center}
%\fancyhead[RO,LE]{\textit{None}}
\fancyhead[RE,LO]{}
\section{Hour of None}
\elcol{℣. O God, {\ding{64}} make speed to save us.}{℣. Deus, {\ding{64}} in adjutórium nostrum inténde.}
\elcol{℟. O Lord, make haste to help us.}{℟. Dómine, ad adjuvándum nos festína.}
\elcol{℣. Glory be to the Father, and to the Son, and to the Holy Ghost.}{℣. Glória Patri, et Fílio, * et Spirítui Sancto:}
\elcol{℟. As it was in the beginning, is now, and ever shall be, world without end. Amen.}{℟. Sicut erat in princípio, et nunc, et semper, * et in sǽcula s{\ae}culórum. Amen.}
\elcol{℣. Praise ye the Lord.}{℣. Laudáte Dominum.}
\elcol{℟. The Lord's Name be praised.}{℟. Sit Nomen Dómini Benedíctum.}
\subby{Rerum Deus}
\elcol{\begin{hangparas}{1.25em}{1}
\lett{O}{God,} creation’s secret force,\par
\secondline{Thyself unmoved, all motion's source,}
%\thirdline{Who from the morn till evening ray,}
\par\noindent
Who from the morn till evening ray,\par\noindent
Through all its changes guid'st the day:\\
\par\noindent
Grant us, when this short life is past,\par\noindent
The glorious evening that shall last;\par\noindent
That, by a holy death attained,\par\noindent
Eternal glory may be gained.\\
\par\noindent
O Father, that we ask be done,\par\noindent
Through Jesus Christ, thine only Son;\par\noindent
Who, with the Holy Ghost and thee,\par\noindent
Doth live and reign eternally. Amen.\\
\end{hangparas}}
{\lett{R}{erum,} Deus, tenax vigor,\par
\secondline{Immótus in te pérmanens,}
\thirdline{Lucis diúrnæ témpora}
Succéssibus detérminans:\\
\par
Largíre lumen véspere,\par
Quo vita nusquam décidat,\par
Sed prǽmium mortis sacr{\ae}\par
Perénnis instet glória.\\
\par
Præsta, Pater piíssime,\par
Patríque compar Únice,\par
Cum Spíritu Paráclito\par
Regnans per omne. Amen.\\}

\subbypsalm{Psalm 126. \textit{In convertendo}}

\lett{W}{hen} the Lord turned again the captivity of Sion : then were we like unto them that dream.\par
\secondline{\psanum{2}Then was our mouth filled with laughter : and our tongue with joy.}
\thirdline{\psanum{3}Then said they among the heathen : The Lord hath done great things for them.}
\psanum{4}Yea, the Lord hath done great things for us already : whereof we rejoice.\par
\psanum{5}Turn our captivity, O Lord : as the rivers in the south.\par
\psanum{6}They that sow in tears : shall reap in joy.\par
\psanum{7}He that now goeth on his way weeping, and beareth forth good seed : shall doubtless come again with joy, and bring his sheaves with him.\par
℣. Glory be to the Father, and to the Son, and to the Holy Ghost.\par
℟. As it was in the beginning, is now, and ever shall be, world without end. Amen.

\subbypsalm{Psalm 127. \textit{Nisi Dominus}}

\lett{E}{xcept} the Lord build the house : their labour is but lost that build it.\par
\secondline{\psanum{2}Except the Lord keep the city : the watchman waketh but in vain.}
\thirdline{\psanum{3}It is but lost labour that ye haste to rise up early, and so late take rest, and eat the bread of carefulness : for so he giveth his beloved sleep.}
\psanum{4}Lo, children and the fruit of the womb : are an heritage and gift that cometh of the Lord.\par
\psanum{5}Like as the arrows in the hand of the giant : even so are the young children.\par
\psanum{6}Happy is the man that hath his quiver full of them : they shall not be ashamed when they speak with their enemies in the gate.\par
℣. Glory be to the Father, and to the Son, and to the Holy Ghost.\par
℟. As it was in the beginning, is now, and ever shall be, world without end. Amen.

%MANUAL ADJUSTMENT:
\vspace{-1ex}

\subbypsalm{Psalm 128. \textit{Beati omnes}}

\lett{B}{lessed} are all they that fear the Lord : and walk in his ways.\par
\secondline{\psanum{2}For thou shalt eat the labours of thine hands : O well is thee, and happy shalt thou be.}
\thirdline{\psanum{3}Thy wife shall be as the fruitful vine : upon the walls of thine house.}
\psanum{4}Thy children like the olive-branches : round about thy table.\par
\psanum{5}Lo, thus shall the man be blessed : that feareth the Lord.\par
\psanum{6}The Lord from out of Sion shall so bless thee : that thou shalt see Jerusalem in prosperity all thy life long.\par
\psanum{7}Yea, that thou shalt see thy children's children : and peace upon Israel.\par
℣. Glory be to the Father, and to the Son, and to the Holy Ghost.\par
℟. As it was in the beginning, is now, and ever shall be, world without end. Amen.

%MANUAL ADJUSTMENT:
\vspace{-1ex}

\subby{Chapter}
\begin{inhead}
	Ferias Throughout the Year
\end{inhead}
\elcoldent{Bear ye one another's burdens, and so fulfil the law of Christ.}{Alter alterius ónera portáte, et sic adimplébitis legem Christi.}
\begin{inhead}
	Sundays \& Feast Days
\end{inhead}
\elcoldent{One Lord, one Faith, one Baptism, one God and Father of all, who is above all, and through all, and in you all.}{Unus Dóminus, una fides, unum baptisma. Unus Deus et Pater ómnium, qui est super omnes, et per ómnia, et in ómnibus nobis.}
\begin{inhead}
	Eastertide
\end{inhead}
\elcoldent{There are three that bear witness in earth: the Spirit, and the Water, and the Blood: and these three agree in one.}{Et tres sunt, qui testimónium dant in terra: spíritus, et aqua, et sanguis: et hi tres unum sunt.}
\elcol{℟. Thanks be to God.}{℟. Deo grátias.}
\subby{Short Respond}
\begin{inhead}
	Ferias Throughout the Year
\end{inhead}
\elcol{℣. Deliver me, O Lord, and be merciful unto me.}{℣. Rédime me, et miserére mei.}

\elcol{℟. Deliver me, O Lord, and be merciful unto me.}{℟. Rédime me, et miserére mei.}

\elcol{℣. My foot standeth right; I will praise the Lord in the congregations.}{℣. Pes meus stetit in dirécto: in ecclésiis benedícam te, Dómine.}

\elcol{℟. And be merciful unto me.}{℟. Et miserére mei.}

\elcol{℣. Glory be to the Father, and to the Son, and to the Holy Ghost.}{℣. Glória Patri, et Fílio, et Spirítui Sancto.}
    
\elcol{℟. Deliver me, O Lord, and be merciful unto me.}{℟. Rédime me, et miserére mei.}

\elcol{℣. I will praise thy Name, O Lord.}{℣. Et confitébor nómini tuo, Dómine.}

\elcol{℟. For thou hast delivered me out of all my trouble.}{℟. Quóniam ex omni tribulatióne eripuísti me.}

\begin{inhead}
	Sundays, Feasts, \& in Eastertide
\end{inhead}

\elcol{℣. I call with my whole heart; hear me, O Lord (alleluia, alleluia).}{℣. Clamávi in toto corde meo, exáudi me, Dómine (allelúja, allelúja).}

\elcol{℟. I call with my whole heart; hear me, O Lord (alleluia, alleluia).}{℟. Clamávi in toto corde meo, exáudi me, Dómine (allelúja, allelúja).}

\elcol{℣. I will keep thy statutes.}{℣. Justificatiónes tuas requíram.}

\elcol{℟. Hear me, O Lord (alleluia, alleluia).}{℟. Exáudi me, Dómine (allelúja, allelúja).}

\elcol{℣. Glory be to the Father, and to the Son, and to the Holy Ghost.}{℣. Glória Patri, et Fílio, et Spirítui Sancto.}

\elcol{℟. I will call with my whole heart; hear me, O Lord (alleluia, alleluia).}{℟. Clamávi in toto corde meo, exáudi me, Dómine (allelúja, allelúja).}

\elcol{℣. Cleanse me, O Lord, from my secret faults (alleluia).}{℣. Ab occúltis meis munda me (allelúja).}

\elcol{℟. Keep thy servant also from presumptuous sins (alleluia).}{℟. Et ab aliénis parce servo tuo (allelúja).}

\subby{Suffrages}
\begin{rubric}
	The Suffrages, said kneeling, are said on all Ferias outside of Eastertide. They are omitted on Sundays, Feasts, and within Principle, Primary, and Secondary Octaves.\par
	\textsc{Note,} If omitted, the Office continues with the Collect.
\end{rubric}

\elcol{℣. Lord, have mercy upon us.}{℣. Kýrie, eléison.}
\elcol{℟. Christ, have mercy upon us.}{℟. Christe, eléison.}
\elcol{℣. Lord, have mercy upon us.}{℣. Kýrie, eléison.}
%\subsec{Lord's Prayer}
\needspace{4\baselineskip}
\elcol{\lett{O}{ur} Father, who art in heaven, Hallowed be thy Name. Thy kingdom come. Thy will be done on earth, As it is in heaven. Give us this day our daily bread. And forgive us our trespasses, As we forgive those who trespass against us. And lead us not into temptation; But deliver us from evil. Amen.}
{\lett{P}{ater} noster, qui es in c{\ae}lis, sanctificétur nomen tuum: advéniat regnum tuum: fiat volúntas tua, sicut in c{\ae}lo et in terra. Panem nostrum quotidiánum da nobis hódie: et dimítte nobis débita nostra, sicut et nos dimíttimus debitóribus nostris: et ne nos indúcas in tentatiónem: sed líbera nos a malo. Amen.}

\elcol{℣. Cast me not away in the time of age.}{℣. Ne proícias me in témpore senectútis.}

\elcol{℟. Forsake me not when my strength faileth me.}{℟. Cum defécerit virtus mea, ne derelínquas me.}

\elcol{℣. Hide not thy face from me.}{℣. Non avértas fáciem tuam a me.}

\elcol{℟. Lest I be like unto them that go down into the pit.}{℟. Et símilis ero descendéntibus in lacum.}

\elcol{℣. Quicken me, O Lord, for thy Name's sake.}{℣. Propter nomen tuum, Dómine, vivificábis me.}

\elcol{℟. And for thy righteousness' sake bring my soul out of trouble.}{℟. In {\ae}quitáte tua edúces de tribulatióne ánimam meam}

\subby{Collects}
\elcol{℣. The Lord be with you.\par
\textit{or,} O Lord, hear our prayer.}{℣. Dóminus vobíscum.\par
\textit{vel,} Dómine, exáudi oratiónem nostram.}
\elcol{℟. And with thy spirit.\par
\textit{or,} And let our cry come unto thee.}{℟. Et cum spíritu tuo.\par
\textit{vel,} Et clamor noster ad te véniat.}
\elcol{℣. Let us pray.}{℣. Orémus.}
\begin{rubric}
	The Collect(s) of the Day is here said, followed by this memorial Collect.
\end{rubric}
\needspace{4\baselineskip}
\lett{L}{ord} Jesus Christ, who for our sakes didst tread the paths of death; make known to us the way of life; that as thou wast reckoned with the transgressors in thy death, and with the rich in thy burial, so we, who are dead in trespasses and sins, may be raised up by thee to the land of true riches. Who livest and reignest with the Father and the Holy Ghost, ever one God, world without end. \textit{Amen.}\\

\elcol{℣. The Lord be with you.\par
\textit{or,} O Lord, hear our prayer.}{℣. Dóminus vobíscum.\par
\textit{vel,} Dómine, exáudi oratiónem nostram.}
\elcol{℟. And with thy spirit.\par
\textit{or,} And let our cry come unto thee.}{℟. Et cum spíritu tuo.\par
\textit{vel,} Et clamor noster ad te véniat.}
\elcol{℣. Let us bless the Lord.}{℣. Benedicámus Dómino.}
\elcol{℟. Thanks be to God.}{℟. Deo grátias}
\elcol{℣. May the souls {\ding{64}} of the faithful departed, through the mercy of God, rest in peace.}{℣. Fidélium ánim{\ae} {\ding{64}} per misericórdiam Dei requiéscant in pace.}
\needspace{4\baselineskip}
\lett{N}{ow} unto him that is able to keep us from falling and to present us faultless before the presence of his glory with exceeding joy; to the only wise God our Saviour be glory and majesty, dominion, and power, both now and ever. \textit{Amen.}


\begin{center}
	{\textsc{Here endeth the Order of None.}}
\end{center}
%Office Hymns

\fancyhead[RE,LO]{}\fancyhead[RO,LE]{}
\fancyhead[C]{}\thispagestyle{empty}
\mainmatter
\phantomsection
\addcontentsline{toc}{chapter}{Daily Hymns}
  \begin{center}
   \includegraphics[width=98mm, height=140mm]{VirginMonks.eps}
   \par
   \vspace{2ex}
   	\textsc{\Huge{Daily Hymns}}
   \end{center}
\fancyhead[C]{\LARGE Daily Hymns}
\fancyhead[RO,LE]{}
\fancyhead[RE,LO]{}
\label{DailyHymns}

\section{Invitatory Hymns}
\subby{Sunday Invitatory Hymn in Summer}\label{SundayInvitatorySummer}
\fancyhead[RO,LE]{\textit{Summer Sunday}}
\fancyhead[RE,LO]{Invitatory}
\gregorioscore{resources/gabc/Ferial/SundayInvitatorySummer.gabc}


\subby{Sunday Invitatory Hymn in Winter}\label{SundayInvitatoryWinter}
\fancyhead[RO,LE]{\textit{Winter Sunday}}
\fancyhead[RE,LO]{Invitatory}
\gregorioscore{resources/gabc/Ferial/SundayInvitatoryWinter.gabc}


\subby{Weekday Invitatory Hymn}\label{WeekdayInvitatory}
\fancyhead[RO,LE]{\textit{Weekday}}
\fancyhead[RE,LO]{Invitatory}
\begin{rubric}
	The following Invitatory Hymn is said on weekdays when Mattins is anticipated the evening prior.
\end{rubric}
\gregorioscore{resources/gabc/Ferial/WeekdayInvitatory.gabc}


\subby{Monday Invitatory Hymn}\label{MondayInvitatory}
\fancyhead[RO,LE]{\textit{Monday}}
\fancyhead[RE,LO]{Invitatory}
\gregorioscore{resources/gabc/Ferial/MondayInvitatory.gabc}


\subby{Tuesday Invitatory Hymn}\label{TuesdayInvitatory}
\fancyhead[RO,LE]{\textit{Tuesday}}
\fancyhead[RE,LO]{Invitatory}
\gregorioscore{resources/gabc/Ferial/TuesdayInvitatory.gabc}


\subby{Wednesday Invitatory Hymn}\label{WednesdayInvitatory}
\fancyhead[RO,LE]{\textit{Wednesday}}
\fancyhead[RE,LO]{Invitatory}
\gregorioscore{resources/gabc/Ferial/WednesdayInvitatory.gabc}


\subby{Thursday Invitatory Hymn}\label{ThursdayInvitatory}
\fancyhead[RO,LE]{\textit{Thursday}}
\fancyhead[RE,LO]{Invitatory}
\gregorioscore{resources/gabc/Ferial/ThursdayInvitatory.gabc}


\subby{Friday Invitatory Hymn}\label{FridayInvitatory}
\fancyhead[RO,LE]{\textit{Friday}}
\fancyhead[RE,LO]{Invitatory}
\gregorioscore{resources/gabc/Ferial/FridayInvitatory.gabc}


\subby{Saturday Invitatory Hymn}\label{SaturdayInvitatory}
\fancyhead[RO,LE]{\textit{Saturday}}
\fancyhead[RE,LO]{Invitatory}
\gregorioscore{resources/gabc/Ferial/SaturdayInvitatory.gabc}

\clearpage
\section{Summer Hymns}
\begin{secrubric}
	The Summer Hymns are said from the Second Sunday after Trinity until the Sunday within 28 September - 4 October, exclusive.
\end{secrubric}

\subby{Sunday Mattins Office Hymn}\label{SundayMattinsSummer}
\fancyhead[RO,LE]{\textit{Summer Sunday}}
\fancyhead[RE,LO]{Mattins}
\gregorioscore{resources/gabc/Ferial/SundayMattinsSummer.gabc}

℣. The Lord is King, and hath put on glorious apparel.

℟. The Lord hath put on his apparel, and girded himself with strength.

\subby{Sunday Evensong Office Hymn}\label{SundayEvensongSummer}
\fancyhead[RO,LE]{\textit{Summer Sunday}}
\fancyhead[RE,LO]{Evensong}
\gregorioscore{resources/gabc/Ferial/SundayEvensongSummer.gabc}

℣. Lord, let my prayer be set forth

℟. In thy sight as the incense.

\subby{Monday Mattins Office Hymn}\label{MondayMattinsSummer}
\fancyhead[RO,LE]{\textit{Summer Monday}}
\fancyhead[RE,LO]{Mattins}
\gregorioscore{resources/gabc/Ferial/MondayMattinsSummer.gabc}

℣. O satisfy us with thy mercy, and that soon.

℟. So shall we rejoice and be glad all the days of our life.

\antiphon{Ben.}{Blessed {\dag} be the Lord God of Israel.}

\subby{Monday Evensong Office Hymn}\label{MondayEvensongSummer}
\fancyhead[RO,LE]{\textit{Summer Monday}}
\fancyhead[RE,LO]{Evensong}
\gregorioscore{resources/gabc/Ferial/MondayEvensongSummer.gabc}

℣. Lord, let my prayer be set forth.

℟. In thy sight as the incense.

\antiphon{Mag.}{My soul {\dag} doth magnify the Lord, for he hath regarded my lowliness.}

\subby{Tuesday Mattins Office Hymn}\label{TuesdayMattinsSummer}
\fancyhead[RO,LE]{\textit{Summer Tuesday}}
\fancyhead[RE,LO]{Mattins}
\gregorioscore{resources/gabc/Ferial/TuesdayMattinsSummer.gabc}

℣. O satisfy us with thy mercy, and that soon.

℟. So shall we rejoice and be glad all the days of our life.

\antiphon{Ben.}{The Lord hath raised up for us {\dag} an horn of salvation in the house of his servant David.}

\subby{Tuesday Evensong Office Hymn}\label{TuesdayEvensongSummer}
\fancyhead[RO,LE]{\textit{Summer Tuesday}}
\fancyhead[RE,LO]{Evensong}
\gregorioscore{resources/gabc/Ferial/TuesdayEvensongSummer.gabc}

℣. Lord, let my prayer be set forth.

℟. In thy sight as the incense.

\antiphon{Mag.}{Let my spirit rejoice {\dag} in God my Saviour.}

\subby{Wednesday Mattins Office Hymn}\label{WednesdayMattinsSummer}
\fancyhead[RO,LE]{\textit{Summer Wednesday}}
\fancyhead[RE,LO]{Mattins}
\gregorioscore{resources/gabc/Ferial/WednesdayMattinsSummer.gabc}

℣. O satisfy us with thy mercy, and that soon.

℟. So shall we rejoice and be glad all the days of our life.

\antiphon{Ben.}{From the hand of all {\dag} that hate us, the Lord hath delivered us.}

\subby{Wednesday Evensong Office Hymn}\label{WednesdayEvensongSummer}
\fancyhead[RO,LE]{\textit{Summer Wednesday}}
\fancyhead[RE,LO]{Evensong}
\gregorioscore{resources/gabc/Ferial/WednesdayEvensongSummer.gabc}

℣. Lord, let my prayer be set forth.

℟. In thy sight as the incense.

\antiphon{Mag.}{The Lord {\dag} hath regarded my lowliness: and he that is mighty hath done in me great things.}

\subby{Thursday Mattins Office Hymn}\label{ThursdayMattinsSummer}
\fancyhead[RO,LE]{\textit{Summer Thursday}}
\fancyhead[RE,LO]{Mattins}
\gregorioscore{resources/gabc/Ferial/ThursdayMattinsSummer.gabc}

℣. O satisfy us with thy mercy, and that soon.

℟. So shall we rejoice and be glad all the days of our life.

\antiphon{Ben.}{Let us serve {\dag} the Lord in holiness, and he will deliver us from the hand of our enemies.}

%MANUAL ADJUSTMENT:
\clearpage
\subby{Thursday Evensong Office Hymn}\label{ThursdayEvensongSummer}
\fancyhead[RO,LE]{\textit{Summer Thursday}}
\fancyhead[RE,LO]{Evensong}
\gregorioscore{resources/gabc/Ferial/ThursdayEvensongSummer.gabc}

℣. Lord, let my prayer be set forth.

℟. In thy sight as the incense.

\antiphon{Mag.}{Show strength, O God, {\dag} with thine arm: scatter the proud, and exalt the humble.}

\subby{Friday Mattins Office Hymn}\label{FridayMattinsSummer}
\fancyhead[RO,LE]{\textit{Summer Friday}}
\fancyhead[RE,LO]{Mattins}
\gregorioscore{resources/gabc/Ferial/FridayMattinsSummer.gabc}

℣. O satisfy us with thy mercy, and that soon.

℟. So shall we rejoice and be glad all the days of our life.

\antiphon{Ben.}{Through the tender mercy {\dag} of our God: whereby the day-spring from on high hath visited us.}

%MANUAL ADJUSTMENT:
\clearpage
\subby{Friday Evensong Office Hymn}\label{FridayEvensongSummer}
\fancyhead[RO,LE]{\textit{Summer Friday}}
\fancyhead[RE,LO]{Evensong}
\gregorioscore{resources/gabc/Ferial/FridayEvensongSummer.gabc}

℣. Lord, let my prayer be set forth.

℟. In thy sight as the incense.

\antiphon{Mag.}{He hath put down the mighty {\dag} that persecute the holy; and hath exalted the humble that confess his Christ.}

\subby{Saturday Mattins Office Hymn}\label{SaturdayMattinsSummer}
\fancyhead[RO,LE]{\textit{Summer Saturday}}
\fancyhead[RE,LO]{Mattins}
\gregorioscore{resources/gabc/Ferial/SaturdayMattinsSummer.gabc}

℣. O satisfy us with thy mercy, and that soon.

℟. So shall we rejoice and be glad all the days of our life.

\antiphon{Ben.}{Give light, O Lord, {\dag} to them that sit in darkness: and guide our feet into the way of peace, thou God of Israel.}

\subby{Saturday Evensong Office Hymn}\label{SaturdayEvensongSummer}
\fancyhead[RO,LE]{\textit{Summer Saturday}}
\fancyhead[RE,LO]{Evensong}
\gregorioscore{resources/gabc/Ferial/SaturdayEvensongSummer.gabc}

℣. Let our evening prayer come up before thee, O Lord.

℟. And let thy mercy come down on us.


\clearpage
\section{Winter Hymns}
\begin{secrubric}
	The Winter Hymns are said from the Octave of Epiphany through Quinquagesima Sunday, inclusive, and from the Sunday within 28 September - 4 October until Advent, exclusive.
\end{secrubric}

\subby{Sunday Mattins Office Hymn}\label{SundayMattinsWinter}
\fancyhead[RO,LE]{\textit{Winter Sunday}}
\fancyhead[RE,LO]{Mattins}
\gregorioscore{resources/gabc/Ferial/SundayMattinsWinter.gabc}

℣. The Lord is King, and hath put on glorious apparel.

℟. The Lord hath put on his apparel, and girded himself with strength.

\subby{Sunday Evensong Office Hymn}\label{SundayEvensongWinter}
\fancyhead[RO,LE]{\textit{Winter Sunday}}
\fancyhead[RE,LO]{Evensong}
\gregorioscore{resources/gabc/Ferial/SundayEvensongWinter.gabc}

℣. Lord, let my prayer be set forth

℟. In thy sight as the incense.

\subby{Monday Mattins Office Hymn}\label{MondayMattinsWinter}
\fancyhead[RO,LE]{\textit{Winter Monday}}
\fancyhead[RE,LO]{Mattins}
\gregorioscore{resources/gabc/Ferial/MondayMattinsWinter.gabc}

℣. O satisfy us with thy mercy, and that soon.

℟. So shall we rejoice and be glad all the days of our life.

\antiphon{Ben.}{Blessed {\dag} be the Lord God of Israel.}

\subby{Monday Evensong Office Hymn}\label{MondayEvensongWinter}
\fancyhead[RO,LE]{\textit{Winter Monday}}
\fancyhead[RE,LO]{Evensong}
\gregorioscore{resources/gabc/Ferial/MondayEvensongWinter.gabc}

℣. Lord, let my prayer be set forth.

℟. In thy sight as the incense.

\antiphon{Mag.}{My soul {\dag} doth magnify the Lord, for he hath regarded my lowliness.}

\subby{Tuesday Mattins Office Hymn}\label{TuesdayMattinsWinter}
\fancyhead[RO,LE]{\textit{Winter Tuesday}}
\fancyhead[RE,LO]{Mattins}
\gregorioscore{resources/gabc/Ferial/TuesdayMattinsWinter.gabc}

℣. O satisfy us with thy mercy, and that soon.

℟. So shall we rejoice and be glad all the days of our life.

\antiphon{Ben.}{The Lord hath raised up for us {\dag} an horn of salvation in the house of his servant David.}

%MANUAL ADJUSTMENT:
\clearpage
\subby{Tuesday Evensong Office Hymn}\label{TuesdayEvensongWinter}
\fancyhead[RO,LE]{\textit{Winter Tuesday}}
\fancyhead[RE,LO]{Evensong}
\gregorioscore{resources/gabc/Ferial/TuesdayEvensongWinter.gabc}

℣. Lord, let my prayer be set forth.

℟. In thy sight as the incense.

\antiphon{Mag.}{Let my spirit rejoice {\dag} in God my Saviour.}

\subby{Wednesday Mattins Office Hymn}\label{WednesdayMattinsWinter}
\fancyhead[RO,LE]{\textit{Winter Wednesday}}
\fancyhead[RE,LO]{Mattins}
\gregorioscore{resources/gabc/Ferial/WednesdayMattinsWinter.gabc}

℣. O satisfy us with thy mercy, and that soon.

℟. So shall we rejoice and be glad all the days of our life.

\antiphon{Ben.}{From the hand of all {\dag} that hate us, the Lord hath delivered us.}

\subby{Wednesday Evensong Office Hymn}\label{WednesdayEvensongWinter}
\fancyhead[RO,LE]{\textit{Winter Wednesday}}
\fancyhead[RE,LO]{Evensong}
\gregorioscore{resources/gabc/Ferial/WednesdayEvensongWinter.gabc}

℣. Lord, let my prayer be set forth.

℟. In thy sight as the incense.

\antiphon{Mag.}{The Lord {\dag} hath regarded my lowliness: and he that is mighty hath done in me great things.}

\subby{Thursday Mattins Office Hymn}\label{ThursdayMattinsWinter}
\fancyhead[RO,LE]{\textit{Winter Thursday}}
\fancyhead[RE,LO]{Mattins}
\gregorioscore{resources/gabc/Ferial/ThursdayMattinsWinter.gabc}

℣. O satisfy us with thy mercy, and that soon.

℟. So shall we rejoice and be glad all the days of our life.

\antiphon{Ben.}{Let us serve {\dag} the Lord in holiness, and he will deliver us from the hand of our enemies.}

\subby{Thursday Evensong Office Hymn}\label{ThursdayEvensongWinter}
\fancyhead[RO,LE]{\textit{Winter Thursday}}
\fancyhead[RE,LO]{Evensong}
\gregorioscore{resources/gabc/Ferial/ThursdayEvensongWinter.gabc}

℣. Lord, let my prayer be set forth.

℟. In thy sight as the incense.

\antiphon{Mag.}{Show strength, O God, {\dag} with thine arm: scatter the proud, and exalt the humble.}

\subby{Friday Mattins Office Hymn}\label{FridayMattinsWinter}
\fancyhead[RO,LE]{\textit{Winter Friday}}
\fancyhead[RE,LO]{Mattins}
\gregorioscore{resources/gabc/Ferial/FridayMattinsWinter.gabc}

℣. O satisfy us with thy mercy, and that soon.

℟. So shall we rejoice and be glad all the days of our life.

\antiphon{Ben.}{Through the tender mercy {\dag} of our God: whereby the day-spring from on high hath visited us.}

%MANUAL ADJUSTMENT:
\clearpage
\subby{Friday Evensong Office Hymn}\label{FridayEvensongWinter}
\fancyhead[RO,LE]{\textit{Winter Friday}}
\fancyhead[RE,LO]{Evensong}
\gregorioscore{resources/gabc/Ferial/FridayEvensongWinter.gabc}

℣. Lord, let my prayer be set forth.

℟. In thy sight as the incense.

\antiphon{Mag.}{He hath put down the mighty {\dag} that persecute the holy; and hath exalted the humble that confess his Christ.}

\subby{Saturday Mattins Office Hymn}\label{SaturdayMattinsWinter}
\fancyhead[RO,LE]{\textit{Winter Saturday}}
\fancyhead[RE,LO]{Mattins}
\gregorioscore{resources/gabc/Ferial/SaturdayMattinsWinter.gabc}

℣. O satisfy us with thy mercy, and that soon.

℟. So shall we rejoice and be glad all the days of our life.

\antiphon{Ben.}{Give light, O Lord, {\dag} to them that sit in darkness: and guide our feet into the way of peace, thou God of Israel.}

\subby{Saturday Evensong Office Hymn}\label{SaturdayEvensongWinter}
\fancyhead[RO,LE]{\textit{Winter Saturday}}
\fancyhead[RE,LO]{Evensong}
\gregorioscore{resources/gabc/Ferial/SaturdayEvensongWinter.gabc}

℣. Let our evening prayer come up before thee, O Lord.

℟. And let thy mercy come down on us.


%Liturgical Hymns

\fancyhead[RE,LO]{}\fancyhead[RO,LE]{}
\fancyhead[C]{}\thispagestyle{empty}
\phantomsection
\addcontentsline{toc}{chapter}{Liturgical Hymns}
  \begin{center}
   \includegraphics[width=98mm, height=140mm]{Resurrection.eps}
   \par
   \vspace{2ex}
   	\textsc{\Huge{Liturgical Hymns}}
   \end{center}
\fancyhead[C]{\LARGE Liturgical Hymns}

\section{Daily Office Services}
	
	\subsection{Office Service One -\\Our Worship, Praise, and Thanksgiving\\unto the Glory of God Alone}\fancyhead[RO,LE]{\textit{Office I}}
	
		\begin{inhead}
			Mattins
		\end{inhead}
		
		\subsubsection{O Come, Let Us Sing}\fancyhead[RE,LO]{O Come, Let Us Sing}
			\gregorioscore{resources/gabc/Extra/Venite.gabc}
			
			\psanum{2}Let us come before his presence with • thanksgiving : and shew ourselves • glad in him with psalms.
			
			\psanum{3}For the Lord is a • great God : and a great • King above all gods.
			
			\psanum{4}In his hand are all the corners • of the earth : and the strength of the hills is • his also.
			
			\psanum{5}The sea is his, and • he made it : and his hands prepared • the dry land.
		
			\psanum{6}O come, let us worship and • fall down : and kneel before the • Lord our Maker.
			
			\psanum{7}For he is the • Lord our God : and we are the people of his pasture, and the • sheep of his hand.
			
			\psanum{8}To-day if ye will hear his voice, harden • not your hearts : as in the provocation, and as in the day of temptation • in the wilderness.
			
			\psanum{9}When your fathers • tempted me : proved me, • and saw my works.
			
			\psanum{10}Forty years long was I grieved with this generation, • and said : It is a people that do err in their hearts, for they • have not known my ways;
			
			\psanum{11}Unto whom I sware • in my wrath : that they should not enter • into my rest.
			
			℣. Glory be to the Father, and • to the Son * and • to the Holy Ghost.
			
			℟. As it was in the beginning, is now, and ever • shall be * world without • end. Amen.
			
	
		\clearpage\subsubsection{Christ Our Passover}\fancyhead[RE,LO]{Christ Our Passover}
			\gregorioscore{resources/gabc/Extra/PaschaNostrum.gabc}
	
			Not with old leaven, neither with the leaven of malice and • wickedness * but with the unleavened bread of sin • cerity and truth.
	
			Christ being raised from the dead dieth • no more * death hath no more do • minion over him.
    
    		For in that he died, he died unto • sin once * but in that he liveth, he • liveth unto God.
    
    		Likewise reckon ye also yourselves to be dead indeed • unto sin * but alive unto God through • Jesus Christ our Lord.
    
    		Christ is risen • from the dead * and become the first • fruits of them that slept.

    		For since by • man came death * by man came also the resur • rection of the dead.
    
    		For as in Adam • all die : even so in Christ shall • all be made alive.
    
    		℣. Glory be to the Father, and • to the Son * and • to the Holy Ghost;
    
    		℟. As it was in the beginning, is now, and ever • shall be * world • without end. Amen.
	
		\clearpage\subsubsection{We Praise Thee, O God}\fancyhead[RE,LO]{We Praise Thee, O God}
			\gregorioscore{resources/gabc/Extra/TeDeum1.gabc}
			
			To thee all Angels cry aloud; the Heavens, and all the • Pow'rs therein;  * To thee Cherubim and Seraphim contin • ually do cry,\par
    		Holy, Holy, Holy, Lord God of • Sabaoth; * Heaven and earth are full of the Majesty • of thy glory.\par
    		The glorious company of the Apostles • praise thee. * The goodly fellowship of the • Prophets praise thee.\par
    		The noble army of Martyrs • praise thee. * The holy Church throughout all the world • doth acknowledge thee;\par
    		The Father of an infinite Majesty; Thine honourable, true and only • Son; * Also the Holy • Ghost the Comforter.\par
    		\gregorioscore{resources/gabc/Extra/TeDeum2.gabc}
    		When thou tookest upon thee • to deliver man, * thou didst not ab • hor the Virgin's womb.\par
    		When thou hadst overcome the • sharpness of death, * thou didst open the Kingdom of Heaven to • all believers.\par
    		Thou sittest at the • right hand of God, * in the glory • of the Father.\par
    		We believe that thou shalt • come to be our Judge. * We therefore pray thee, help thy servants, whom thou hast redeemed • with thy precious blood.\par
    		Make them to be • numbered with thy Saints, * in glory • everlasting.\par
    		\clearpage
    		\gregorioscore{resources/gabc/Extra/TeDeum3.gabc}
			Day by day we magnify • thee; * And we worship thy Name ever, • world without end.\par
			Vouchsafe, O Lord, to keep us this day without • sin. * O Lord, have mercy upon us, have mer • cy upon us.\par
			O Lord, let thy mercy lighten upon us, as our trust is in • thee. * O Lord, in thee have I trusted; let me never • be confounded.
			

		\clearpage\subsubsection{Bless Ye the Lord}\fancyhead[RE,LO]{Bless Ye the Lord}
			\gregorioscore{resources/gabc/Extra/Benedicite.gabc}
			
			O ye Angels of the Lord, bless • ye the Lord: * praise him, and magni • fy him for ever.\par
			
			\lett{O}{ye} Heavens, bless • ye the Lord: * praise him, and magni • fy him for ever.\par
			\secondline{O ye Waters that be above the firmament, bless • ye the Lord: * praise him, and magni • fy him for ever.}
			\thirdline{O all ye Powers of the Lord, bless • ye the Lord: * praise him, and magni • fy him for ever.}
			O ye Sun and Moon, bless • ye the Lord: * praise him, and magni • fy him for ever.\par
    		O ye Stars of heaven, bless • ye the Lord: * praise him, and magni • fy him for ever.\par
    		O ye Showers and Dew, bless • ye the Lord: * praise him, and magni • fy him for ever.\par
    		O ye winds of God, bless • ye the Lord: * praise him, and magni • fy him for ever.\par
    		O ye Fire and Heat, bless • ye the Lord: * praise him, and magni • fy him for ever.\par
    		O ye Winter and Summer, bless • ye the Lord: * praise him, and magni • fy him for ever.\par
    		O ye Dews and Frosts, bless • ye the Lord: * praise him, and magni • fy him for ever.\par
    		O ye Frost and Cold, bless • ye the Lord: * praise him, and magni • fy him for ever.\par
    		O ye Ice and Snow, bless ye the Lord * praise him, and magni • fy him for ever.\par
    		O ye Nights and Days, bless • ye the Lord: * praise him, and magni • fy him for ever.\par
    		O ye Light and Darkness, bless • ye the Lord: * praise him, and magni • fy him for ever.\par
    		O ye Lightnings and Clouds, bless ye the Lord * praise him, and magni • fy him for ever.\par
			\lett{O}{let} the Earth • bless the Lord: * yea, let it praise him, and magni • fy him for ever.\par
			\secondline{O ye Mountains and Hills, bless • ye the Lord: * praise him, and magni • fy him for ever.}
			\clearpage
			\gregorioscore{resources/gabc/Extra/Benedicite_secondpage.gabc}
			%\thirdline{O all ye Green Things upon the earth, bless • ye the Lord: * praise him, and magni • fy him for ever.}
			O ye Wells, bless • ye the Lord: * praise him, and magni • fy him for ever.\par
    		O ye Seas and Floods, bless • ye the Lord: * praise him, and magni • fy him for ever.\par
    		O ye Whales, and all that move in the waters, bless • ye the Lord: * praise him, and magni • fy him for ever.\par
    		O all ye Fowls of the air, bless • ye the Lord: * praise him, and magni • fy him for ever.\par
    		O all ye Beasts and Cattle, bless • ye the Lord: * praise him, and magni • fy him for ever.\par
    		O ye Children of Men, bless • ye the Lord: * praise him, and magni • fy him for ever.\par
    		\lett{O}{let} Israel • bless the Lord: * praise him, and magni • fy him for ever.\par
    		\secondline{O ye Priests of the Lord, bless • ye the Lord: * praise him, and magni • fy him for ever.}
    		\thirdline{O ye Servants of the Lord, bless • ye the Lord: * praise him, and magni • fy him for ever.}
    		O ye Spirits and Souls of the Righteous, bless • ye the Lord: * praise him, and magni • fy him for ever.\par
    		O ye holy and humble Men of heart, bless • ye the Lord: * praise him, and magni • fy him for ever.\par
    		\lett{L}{et} us bless the Father, and the Son, and the • Holy Ghost: * praise him, and magnif • y him for ever.
			
			
		\clearpage\subsubsection{Blessed Be the Lord}\fancyhead[RE,LO]{Blessed Be the Lord}
			\gregorioscore{resources/gabc/Extra/Benedictus.gabc}
			
		\clearpage\subsubsection{Athanasian Creed}\fancyhead[RE,LO]{Athanasian Creed}
			\gregorioscore{resources/gabc/Extra/Athanasian.gabc}
	
			\psanum{2}Which Faith except a man keep whole • and undefiled : without doubt he will perish • eternally.
	
			\psanum{3}Now the • Catholick Faith is this : that we worship one God in Trinity, and the Trinity • in Unity;
	
    		\psanum{4}Neither • confusing the Persons : nor divid • ing the substance.\par
    		\psanum{5}For there is one Person of the Father, • another of the Son : another • of the Holy Ghost;\par
    		\psanum{6}But the Godhead of the Father, and of the Son, and of the • Holy Ghost is all one : the glory equal, the majesty • co-eternal.\par
    		\psanum{7}Such as the Father is, • such is the Son : and such • is the Holy Ghost.\par
    		\psanum{8}The Father uncreated, the Son • uncreated : the Holy Ghost • uncreated;\par
    		\psanum{9}The Father infinite, • the Son infinite : the Holy • Ghost infinite;\par
    		\psanum{10}The Father eternal, the • Son eternal : the Holy • Ghost eternal.\par
    		\psanum{11}And yet there are not • three eternals : • but one eternal\par
    		\psanum{12}As also there are not three uncreated, nor • three infinites : but one infinite, and one • uncreated.\par
    		\psanum{13}So likewise the Father is almighty, • the Son almighty : the Holy • Ghost almighty;\par
    		\psanum{14}And yet there are not • three almighties : • but one almighty.\par
    		\psanum{15}So the Father • is God, the Son God : • the Holy Ghost God;\par
    		\psanum{16}And yet there • are not three Gods : • but one God.\par
    		\psanum{17}So the Father is • Lord, the Son Lord : • the Holy Ghost Lord;\par
    		\psanum{18}And yet there • are not three Lords : • but one Lord.\par
    		\psanum{19}For like as we are compelled by the • Christian verity : to confess each Person by himself to • be both God and Lord;\par
    		\psanum{20}So are we forbidden by the • Catholick religion : to speak of • three Gods or three Lords.\par
    		\psanum{21}The Father • is made of none : nor created, • nor begotten.\par
    		\psanum{22}The Son is • of the Father : not made, nor created, • but begotten.\par
    		\psanum{23}The Holy Ghost is • of the Father : not made, nor created, nor begotten, • but proceeding.\par
    		\psanum{24}There is therefore one Father, not three Fathers; • one Son, not three Sons : one Holy Ghost, • not three Holy Ghosts.\par
    		\psanum{25}And in this Trinity there is no • before or after : • no greater or less;\par
    		\gregorioscore{resources/gabc/Extra/Athanasian2.gabc}\par
    		%\psanum{26}But all three Persons are co • eternal together : • and co-equal.\par
    		\psanum{27}So that in all ways, • as is aforesaid : both the Trinity is to be worshipped in Unity, and the Unity • in Trinity.\par
    		\psanum{28}He therefore that • would be saved : let him thus think • of the Trinity.

    		\lett{F}{urthermore} it is necessary to • eternal salvation : that he also believe faithfully the Incarnation of • our Lord Jesus Christ.\par
    		\secondline{\psanum{30}Now the right faith is that we • believe and confess : that our Lord Jesus Christ, the Son of God, • is both God and man.}
    		\thirdline{\psanum{31}He is God, of the substance of the Father, begotten • before the worlds : and he is man, of the substance of his Mother, • born in the world;}
    		\psanum{32}Perfect • God perfect man : of reasoning soul and human • flesh subsisting;\par
    		\psanum{33}Equal to the Father as • touching his Godhead : less than the Father as • touching his manhood.\par
    		\psanum{34}Who although • he be God and man : yet he is not two, • but is one Christ;\par
    		\psanum{35}One however, not by conversion of • Godhead into flesh : but by taking • manhood into God;\par
    		\psanum{36}One • altogether : not by confusion of substance, but by unit • y of person.\par
    		\psanum{37}For as reasoning soul and • flesh is one man : so God and • man is one Christ;\par
    		\psanum{38}Who suffered for • our salvation : descended into hell, rose • again from the dead;\par
    		\psanum{39}Ascended into heaven, sat down at the • right hand of the Father : from whence he shall come to judge • the quick and the dead.\par
    		\psanum{40}At whose coming all men must rise again • with their bodies : and shall give account • for their own deeds.\par
    		\psanum{41}And they that have done good will go into • life eternal : they that have done evil into • eternal fire.\par

    		\lett{T}{his} is the • Catholick Faith : which except a man do faithfully and steadfastly believe, he • cannot be saved.\par
    		\secondline{℣. Glory be to the Father, • and to the Son : and • to the Holy Ghost;}
    		\thirdline{℟. As it was in the beginning, is now, and • ever shall be : world • without end. Amen.}
    		
    	\clearpage
    	\begin{inhead}
    		Evensong
    	\end{inhead}
    	
    	\subsubsection{Glory Be to God on High}\fancyhead[RE,LO]{Glory Be to God}
    		%Old Scottish Chant
			\gregorioscore{resources/gabc/MassParts/Scottish/ScottishGloria.gabc}
    
		\clearpage\subsubsection{My Soul Doth Magnify the Lord}\fancyhead[RE,LO]{My Soul Doth Magnify}
			\gregorioscore{resources/gabc/Extra/Magnificat.gabc}
			
		\subsubsection{Lord, Now Lettest Thou Thy Servant}\fancyhead[RE,LO]{Lord, Now Lettest}
	
			\gregorioscore{resources/gabc/Extra/NuncDimittis.gabc}

			For mine eyes • have seen * • thy salvation,

			Which thou hast pre • pared * before the face • of all people;

			To be a light to lighten the • Gentiles, * and to be the glory of thy • people Israel.
	
			℣. Glory be to the Father, and to • the Son, * and • to the Holy Ghost.
	
			℟. As it was in the beginning, {\dag} is now, and ever • shall be, * world • without end. Amen.
	
	\clearpage
	
	\selectlanguage{latin}
	
	\subsection{Office Service Two - Officium Omnium Sanctorum Monachorum et Eremitarum}\fancyhead[RO,LE]{\textit{Office II}}
	
		\begin{inhead}
			Matutinum
		\end{inhead}
		
		\subsubsection{Venite}\fancyhead[RE,LO]{Venite}
			\gregorioscore{resources/gabc/Extra/VeniteLatin.gabc}
			
		\clearpage\subsubsection{Pascha Nostrum}
			\gregorioscore{resources/gabc/Extra/PaschaNostrumLatin.gabc}
			
			Non in fermento veteri, neque in fermento maliti{\ae} et ne • quiti{\ae} * sed in Azymis sinceritatis et • veritatis.
	
			Christus, resurgens ex mortuis, jam • non moritur * mors illi ultra non do • minabitur.
			
			Quod enim mortuus est, peccato mortuus • est semel * quod autem vivit, • vivit Deo.
			
			Ita et vos existimate, vos mortuos quidem esse • peccato, * viventes autem Deo, in Christo Jesu, Do • mino nostro.
    
			Nunc Christus resurrexit a • mortuis * primiti{\ae} dorm • ientium.
			
			Quoniam quidem per ho • minem mors * et per hominem resurrectio • mortuorum. 
			
			Et, sicut in Adam omnes mor • iuntur * ita et in Christo omnes vivi • ficabuntur. 
    		
    		℣. Glória Patri, et • Fílio, * et Spirí • tui Sancto:
    
    		℟. ℟. Sicut erat in princípio, et nunc, • et semper, * et in sǽcula s{\ae}cu • lórum. Amen.
			
		\clearpage\subsubsection{Te Deum}\fancyhead[RE,LO]{Te Deum}
			\gregorioscore{resources/gabc/Extra/TeDeumLatin.gabc}
			
		\clearpage\subsubsection{Benedicite}\fancyhead[RE,LO]{Benedicite}
			\gregorioscore{resources/gabc/Extra/BenediciteLatin.gabc}
			
			Benedícite, Ángeli • Dómini, Dómino: * laudáte et superexaltáte e • um in sǽcula.\par
			Benedícite, • c{\ae}li, Dómino:  * laudáte et superexaltáte e • um in sǽcula.\par
			Benedícite, aqu{\ae} omnes, qu{\ae} super • c{\ae}los sunt, Dómino: * laudáte et superexaltáte e • um in sǽcula.\par
			Benedícite, omnes virtútes • Dómini, Dómino: * laudáte et superexaltáte e • um in sǽcula.\par
			Benedícite, sol et • luna, Dómino: * laudáte et superexaltáte e • um in sǽcula.\par
			Benedícite, stell{\ae} • c{\ae}li, Dómino: * laudáte et superexaltáte e • um in sǽcula.\par
			Benedícite, omnis imber • et ros, Dómino: * laudáte et superexaltáte e • um in sǽcula.\par
			Benedícite, omnes spíritus • Dei, Dómino: * laudáte et superexaltáte e • um in sǽcula.\par
			Benedícite, ignis et • {\ae}stus, Dómino: * laudáte et superexaltáte e • um in sǽcula.\par
			Benedícite, frigus et • {\ae}stus, Dómino: * laudáte et superexaltáte e • um in sǽcula.\par
			Benedícite, rores et pru • ína, Dómino: * laudáte et superexaltáte e • um in sǽcula.\par
			Benedícite, gelu et • frigus, Dómino: * laudáte et superexaltáte e • um in sǽcula.\par
			Benedícite, glácies et • nives, Dómino: * laudáte et superexaltáte e • um in sǽcula.\par
			Benedícite, noctes et • dies, Dómino: * laudáte et superexaltáte e • um in sǽcula.\par
			Benedícite, lux et • ténebr{\ae}, Dómino: * laudáte et superexaltáte e • um in sǽcula.\par
			Benedícite, fúlgura et • nubes, Dómino: * laudáte et superexaltáte e • um in sǽcula.\par

			Benedícat • terra Dóminum: * laudet et superexáltet e • um in sǽcula.\par
			Benedícite, montes et • colles, Dómino: * laudáte et superexaltáte e • um in sǽcula.\par
			Benedícite, univérsa germinántia in • terra, Dómino: * laudáte et superexaltáte e • um in sǽcula.\par
			Benedícite, • fontes, Dómino: * laudáte et superexaltáte e • um in sǽcula.\par
			Benedícite, mária et • flúmina, Dómino: * laudáte et superexaltáte e • um in sǽcula.\par
			Benedícite, cete, et ómnia, qu{\ae} movéntur in • aquis, Dómino: * laudáte et superexaltáte e • um in sǽcula.\par
			\gregorioscore{resources/gabc/Extra/BenediciteLatin2.gabc}
			%Benedícite, omnes vólucres • c{\ae}li, Dómino: * laudáte et superexaltáte e • um in sǽcula.\par
			Benedícite, omnes bésti{\ae} et • pécora, Dómino: * laudáte et superexaltáte e • um in sǽcula.\par
			Benedícite, fílii • hóminum, Dómino: * laudáte et superexaltáte e • um in sǽcula.\par
			Benedícat • Israël Dóminum: * laudet et superexáltet e • um in sǽcula.\par
			Benedícite, sacerdótes • Dómini, Dómino: * laudáte et superexaltáte e • um in sǽcula.\par
			Benedícite, servi • Dómini, Dómino: * laudáte et superexaltáte e • um in sǽcula.\par
			Benedícite, spíritus, et ánim{\ae} just • órum, Dómino: * laudáte et superexaltáte e • um in sǽcula.\par
			Benedícite, sancti, et húmiles • corde, Dómino: * laudáte et superexaltáte e • um in sǽcula.\par
			Benedicámus Patrem et Fílium cum • Sancto Spíritu: * laudémus et superexaltémus e • um in sǽcula.
			
		\clearpage\subsubsection{Benedictus}\fancyhead[RE,LO]{Benedictus}
			\gregorioscore{resources/gabc/Extra/BenedictusLatin.gabc}
			
		%\clearpage\subsubsection{Symbolum Apostolorum}\fancyhead[RE,LO]{Symbolum Apostolorum}
		%	\gregorioscore{resources/gabc/Extra/CredoApostolorum.gabc}
		
		\clearpage\subsubsection{Symbolum Athanasianum}\fancyhead[RE,LO]{Athanasianum}
			\gregorioscore{resources/gabc/Extra/AthanasianLatin.gabc}
	
			Quam nisi quisque íntegram inviolatámque serváverit, * absque dúbio in {\ae}térnum períbit.
	
			Fides autem cathólica h{\ae}c est: * ut unum Deum in Trinitáte, et Trinitátem in unitáte venerémur.
	
			Neque confundéntes persónas, * neque substántiam separántes.\par
			Alia est enim persóna Patris, ália Fílii, * ália Spíritus Sancti:\par
			Sed Patris, et Fílii, et Spíritus Sancti una est divínitas, * {\ae}quális glória, co{\ae}térna majéstas.\par
			Qualis Pater, talis Fílius, * talis Spíritus Sanctus.\par
			Increátus Pater, increátus Fílius, * increátus Spíritus Sanctus.\par
			Imménsus Pater, imménsus Fílius, * imménsus Spíritus Sanctus.\par
			{\AE}térnus Pater, {\ae}térnus Fílius, * {\ae}térnus Spíritus Sanctus.\par
			Et tamen non tres {\ae}térni, * sed unus {\ae}térnus.\par
			Sicut non tres increáti, nec tres imménsi, * sed unus increátus, et unus imménsus.\par
			Simíliter omnípotens Pater, omnípotens Fílius, * omnípotens Spíritus Sanctus.\par
			Et tamen non tres omnipoténtes, * sed unus omnípotens.\par
			Ita Deus Pater, Deus Fílius, * Deus Spíritus Sanctus.\par
			Ut tamen non tres Dii, * sed unus est Deus.\par
			Ita Dóminus Pater, Dóminus Fílius, * Dóminus Spíritus Sanctus.\par
			Et tamen non tres Dómini, * sed unus est Dóminus.\par
			Quia, sicut singillátim unamquámque persónam Deum ac Dóminum confitéri christiána veritáte compéllimur: * ita tres Deos aut Dóminos dícere cathólica religióne prohibémur.\par
			Pater a nullo est factus: * nec creátus, nec génitus.\par
			Fílius a Patre est: * non factus, nec creátus, sed génitus.\par
			Spíritus Sanctus a Patre: * non factus, nec creátus, nec génitus, sed procédens.\par
			Unus ergo Pater, non tres Patres: unus Fílius, non tres Fílii: * unus Spíritus Sanctus, non tres Spíritus Sancti.\par
			Et in hac Trinitáte nihil prius aut postérius, nihil majus aut minus: * sed tot{\ae} tres persón{\ae} co{\ae}térn{\ae} sibi sunt et co{\ae}quáles.\par
			Ita ut per ómnia, sicut jam supra dictum est, * et únitas in Trinitáte, et Trínitas in unitáte veneránda sit.\par
			Qui vult ergo salvus esse, * ita de Trinitáte séntiat.\par
			\gregorioscore{resources/gabc/Extra/AthanasianLatin2.gabc}\par
			%\lett{S}{ed} necessárium est ad {\ae}térnam salútem, * ut Incarnatiónem quoque Dómini nostri Jesu Christi fidéliter credat.\par
			Est ergo fides recta ut credámus et confiteámur, * quia Dóminus noster Jesus Christus, Dei Fílius, Deus et homo est.
			Deus est ex substántia Patris ante sǽcula génitus: * et homo est ex substántia matris in sǽculo natus.\par
			Perféctus Deus, perféctus homo: * ex ánima rationáli et humána carne subsístens.\par
			{\AE}quális Patri secúndum divinitátem: * minor Patre secúndum humanitátem.\par
			Qui licet Deus sit et homo, * non duo tamen, sed unus est Christus.\par
			Unus autem non conversióne divinitátis in carnem, * sed assumptióne humanitátis in Deum.\par
			Unus omníno, non confusióne substánti{\ae}, * sed unitáte persón{\ae}.\par
			Nam sicut ánima rationális et caro unus est homo: * ita Deus et homo unus est Christus.\par
			Qui passus est pro salúte nostra: descéndit ad ínferos: * tértia die resurréxit a mórtuis.\par
			Ascéndit ad c{\ae}los, sedet ad déxteram Dei Patris omnipoténtis: * inde ventúrus est judicáre vivos et mórtuos.\par
			Ad cujus advéntum omnes hómines resúrgere habent cum corpóribus suis; * et redditúri sunt de factis própriis ratiónem.\par
			Et qui bona egérunt, ibunt in vitam {\ae}térnam: * qui vero mala, in ignem {\ae}térnum.\par
			\lett{H}{{\ae}c} est fides cathólica, * quam nisi quisque fidéliter firmitérque credíderit, salvus esse non póterit.\par
			\secondline{℣. Glória Patri, et Fílio, * et Spirítui Sancto:}
			\thirdline{℟. Sicut erat in princípio, et nunc, et semper, * et in sǽcula s{\ae}culórum. Amen.}
		
		\clearpage
		
		\begin{inhead}
			Vesper{\ae}
		\end{inhead}
		
		\subsubsection{Magnificat}\fancyhead[RE,LO]{Magnificat}
			\gregorioscore{resources/gabc/Extra/MagnificatLatin.gabc}
			
		\subsubsection{Nunc Dimittis}\fancyhead[RE,LO]{Nunc Dimittis}
			\gregorioscore{resources/gabc/Extra/NuncDimittisLatin_simple.gabc}
			Quia vidérunt • óculi mei * salu • táre tuum,\par
			Quod • parásti * ante fáciem ómnium • populórum,\par
			Lumen ad revela • tiónem géntium, * et glóriam plebis • tu{\ae} Isra{\"e}l.\par
			℣. Glória Patri, • et Fílio, * et Spir • ítui Sancto:\par
			℟. Sicut erat in princípio, et • nunc, et semper, * et in sǽcula • s{\ae}culórum. Amen.
			
\clearpage

\selectlanguage{british}

\section{Marian Anthems}
	%Alma Redemptoris Mater
	\subsection{Alma Redemptoris Mater}\fancyhead[RO,LE]{\textit{Gracious Mother}}
	\begin{inhead}
	From I Evensong of Advent I until II Evensong of the Purification, inclusive.
	\end{inhead}
		\gregorioscore{resources/gabc/Extra/AlmaSimple.gabc}

	%Ave, Regina C{\ae}lorum
	\clearpage\subsection{Ave, Regina C{\ae}lorum}\fancyhead[RO,LE]{\textit{Queen of the Heavens}}
	\begin{inhead}
	From Compline of the Purification until Compline of Holy Wednesday, inclusive.
	\end{inhead}
		\gregorioscore{resources/gabc/Extra/AveSimple.gabc}

	%Regina C{\ae}li
	\clearpage\subsection{Regina C{\ae}li}\fancyhead[RO,LE]{\textit{O Queen of Heaven}}
	\begin{inhead}
	From Compline of Holy Saturday until I Evensong of Trinity Sunday, exclusive.
	\end{inhead}
		\gregorioscore{resources/gabc/Extra/ReginaSimple.gabc}

	%Salve Regina
	\clearpage\subsection{Mary, We Hail Thee, Simple Tone}\fancyhead[RO,LE]{\textit{Mary, We Hail Thee}}\label{SalveReginaEnglish}
	\begin{inhead}
	From I Evensong of Trinity Sunday until I Evensong of Advent I, exclusive.
	\end{inhead}
		\gregorioscore{resources/gabc/Extra/SalveReginaSimple.gabc}

	\clearpage\subsection{Mary, We Hail Thee, Solemn Tone}\fancyhead[RO,LE]{\textit{Mary, We Hail Thee}}
	\begin{inhead}
	From I Evensong of Trinity Sunday until I Evensong of Advent I, exclusive.
	\end{inhead}
		\gregorioscore{resources/gabc/Extra/SalveReginaSolemn.gabc}
		
\clearpage

\section{Communion Services}
	
	%Merbecke:
	\subsection{Communion Service One - Mass of the Grace of God}\fancyhead[RO,LE]{\textit{Communion I}}

		\subsubsection{Ten Commandments}\fancyhead[RE,LO]{10 Commandments}
			\begin{rubric}
				In response to the first nine Commandments, the following is sung.
			\end{rubric}

			\gregorioscore{resources/gabc/MassParts/Merbecke/Decalogue1.gabc}

			\begin{rubric}
				In response to the final Commandment, the following is sung.
			\end{rubric}

			\gregorioscore{resources/gabc/MassParts/Merbecke/Decalogue2.gabc}
			
		\subsubsection{Lord, Have Mercy Upon Us}\fancyhead[RE,LO]{Lord, Have Mercy}
			\gregorioscore{resources/gabc/MassParts/Merbecke/MerbeckeKyrie.gabc}
		
		\clearpage\subsubsection{Glory be to God on High}\fancyhead[RE,LO]{Glory Be to God}
			\gregorioscore{resources/gabc/MassParts/Merbecke/MerbeckeGloria.gabc}
		
		\clearpage\subsubsection{Nicene Creed}\fancyhead[RE,LO]{Nicene Creed}
			\gregorioscore{resources/gabc/MassParts/Merbecke/MerbeckeCreed.gabc}
		
		\clearpage\subsubsection{Holy, Holy, Holy}\fancyhead[RE,LO]{Holy, Holy, Holy}
			\gregorioscore{resources/gabc/MassParts/Merbecke/MerbeckeSanctus.gabc}
		
		\subsubsection{O Lamb of God}\fancyhead[RE,LO]{O Lamb of God}
			\gregorioscore{resources/gabc/MassParts/Merbecke/MerbeckeAgnus.gabc}
		
		
	%\subsection{Communion Service Two - Mass of St. Margaret}\fancyhead[RO,LE]{\textit{Communion II}}
	%	\subsubsection{Ten Commandments}\fancyhead[RE,LO]{10 Commandments}
	%		\begin{rubric}
	%			In response to the first nine Commandments, the following is sung.
	%		\end{rubric}
	%		\gregorioscore{resources/gabc/MassParts/Scottish/ScottishDecalogue1.gabc}
	%
	%		\begin{rubric}
	%			In response to the final Commandment, the following is sung.
	%		\end{rubric}
	%		\gregorioscore{resources/gabc/MassParts/Scottish/ScottishDecalogue2.gabc}
	%		
	%	\clearpage\subsubsection{Kyrie Eleison}\fancyhead[RE,LO]{Kyrie}
	%		\gregorioscore{resources/gabc/MassParts/Scottish/ScottishKyrie.gabc}
	%	
	%	\clearpage\subsubsection{Gloria in excelsis}\fancyhead[RE,LO]{Gloria}
	%		%Old Scottish Chant
	%		\gregorioscore{resources/gabc/MassParts/Scottish/ScottishGloria.gabc}
	%	
	%	\clearpage\subsubsection{Nicene Creed}\fancyhead[RE,LO]{Credo}
	%		
	%	
	%	\clearpage\subsubsection{Sanctus}\fancyhead[RE,LO]{Sanctus}
	%		
	%	
	%	\clearpage\subsubsection{Agnus Dei}\fancyhead[RE,LO]{Agnus}
			
	%Missa Marialis
	\clearpage\subsection{Communion Service Two - Mass of Our Lady}\fancyhead[RO,LE]{\textit{Communion II}}
		\subsubsection{Ten Commandments}\fancyhead[RE,LO]{10 Commandments}
			\begin{rubric}
				In response to the first nine Commandments, the following is sung.
			\end{rubric}
			\gregorioscore{resources/gabc/MassParts/Scottish/ScottishDecalogue1.gabc}
	
			\begin{rubric}
				In response to the final Commandment, the following is sung.
			\end{rubric}
			\gregorioscore{resources/gabc/MassParts/Scottish/ScottishDecalogue2.gabc}
			
		\clearpage\subsubsection{Lord, Have Mercy Upon Us}\fancyhead[RE,LO]{Lord, Have Mercy}
			\gregorioscore{resources/gabc/MassParts/Marialis/MarialisKyrie.gabc}
		
		\clearpage\subsubsection{Glory Be to God on High}\fancyhead[RE,LO]{Glory Be to God}
			\gregorioscore{resources/gabc/MassParts/Marialis/MarialisGloria.gabc}
		
		\clearpage\subsubsection{Nicene Creed}\fancyhead[RE,LO]{Nicene Creed}
			\gregorioscore{resources/gabc/MassParts/Marialis/MarialisCreed.gabc}
		
		\clearpage\subsubsection{Holy, Holy, Holy}\fancyhead[RE,LO]{Holy, Holy, Holy}
			\gregorioscore{resources/gabc/MassParts/Marialis/MarialisSanctus.gabc}
		
		\clearpage\subsubsection{O Lamb of God}\fancyhead[RE,LO]{O Lamb of God}
			\gregorioscore{resources/gabc/MassParts/Marialis/MarialisAgnus.gabc}
			
			
		%Missa de Angelis in English
	\clearpage\subsection{Communion Service Three - Mass of the Angels}\fancyhead[RO,LE]{\textit{Communion III}}
		%\subsubsection{Ten Commandments}\fancyhead[RE,LO]{10 Commandments}
		%	\begin{rubric}
		%		In response to the first nine Commandments, the following is sung.
		%	\end{rubric}
		%
			%
		%
		%	\begin{rubric}
		%		In response to the final Commandment, the following is sung.
		%	\end{rubric}
		%
			%
		%
		\subsubsection{Lord, Have Mercy Upon Us}\fancyhead[RE,LO]{Lord, Have Mercy}
			\gregorioscore{resources/gabc/MassParts/Angels/KyrieVIII.gabc}
		
		\clearpage\subsubsection{Glory Be to God on High}\fancyhead[RE,LO]{Glory Be to God}
			\gregorioscore{resources/gabc/MassParts/Angels/GloriaVIII.gabc}
		
		\subsubsection{Holy, Holy, Holy}\fancyhead[RE,LO]{Holy, Holy, Holy}
			\gregorioscore{resources/gabc/MassParts/Angels/SanctusVIII.gabc}
		
		\clearpage\subsubsection{O Lamb of God}\fancyhead[RE,LO]{O Lamb of God}
			\gregorioscore{resources/gabc/MassParts/Angels/AgnusVIII.gabc}
			

		%Missa Orbis Factor in English
	\clearpage\subsection{Communion Service Four - Mass of the Maker of the World}\fancyhead[RO,LE]{\textit{Communion IV}}
		%\subsubsection{Ten Commandments}\fancyhead[RE,LO]{10 Commandments}
		%	\begin{rubric}
		%		In response to the first nine Commandments, the following is sung.
		%	\end{rubric}
		%
			%
		%
		%	\begin{rubric}
		%		In response to the final Commandment, the following is sung.
		%	\end{rubric}
		%
			%
		%
		\subsubsection{Lord, Have Mercy Upon Us}\fancyhead[RE,LO]{Lord, Have Mercy}
			\gregorioscore{resources/gabc/MassParts/World/KyrieXI.gabc}
			
		\clearpage\subsubsection{Glory Be to God on High}\fancyhead[RE,LO]{Glory Be to God}
			\gregorioscore{resources/gabc/MassParts/World/GloriaXI.gabc}
		
		\clearpage\subsubsection{Holy, Holy, Holy}\fancyhead[RE,LO]{Holy, Holy, Holy}
			\gregorioscore{resources/gabc/MassParts/World/SanctusXI.gabc}
		
		\subsubsection{O Lamb of God}\fancyhead[RE,LO]{O Lamb of God}
			\gregorioscore{resources/gabc/MassParts/World/AgnusXI.gabc}


		%Missa Deus Genitor Alme in English
	\clearpage\clearpage\subsection{Communion Service Five - Mass of God Our Sweet Father}\fancyhead[RO,LE]{\textit{Communion V}}
		%\subsubsection{Ten Commandments}\fancyhead[RE,LO]{10 Commandments}
		%	\begin{rubric}
		%		In response to the first nine Commandments, the following is sung.
		%	\end{rubric}
		%
			%
		%
		%	\begin{rubric}
		%		In response to the final Commandment, the following is sung.
		%	\end{rubric}
		%
			%
		%
		\subsubsection{Lord, Have Mercy Upon Us}\fancyhead[RE,LO]{Lord, Have Mercy}
			\gregorioscore{resources/gabc/MassParts/Sweet/KyrieXVIII.gabc}
		
		\subsubsection{Holy, Holy, Holy}\fancyhead[RE,LO]{Holy, Holy, Holy}
			\gregorioscore{resources/gabc/MassParts/Sweet/SanctusXVIII.gabc}
		
		\clearpage\subsubsection{O Lamb of God}\fancyhead[RE,LO]{O Lamb of God}
			\gregorioscore{resources/gabc/MassParts/Sweet/AgnusXVIII.gabc}
			
	
\clearpage\section{Creeds}\fancyhead[RO,LE]{\textit{}}\fancyhead[RE,LO]{}

	%Credo III in English
	\subsection{Creed III}\fancyhead[RO,LE]{\textit{Creed III}}
		\gregorioscore{resources/gabc/MassParts/Creeds/CredoIII.gabc}
		
	%Credo III
	\clearpage\subsection{Credo III}\fancyhead[RO,LE]{\textit{Credo III}}
		\gregorioscore{resources/gabc/MassParts/Credos/CredoIIILatin.gabc}


\clearpage\section{Sequences}
	%Easter
	\subsection{Christians, To the Paschal Victim}\fancyhead[RO,LE]{\textit{To the Paschal Victim}}\fancyhead[RE,LO]{Easter Sunday}
		\gregorioscore{resources/gabc/Sequences/Victimae.gabc}
		
	\clearpage\subsection{Victimae paschali laudes}\fancyhead[RO,LE]{\textit{Victimae}}
		\gregorioscore{resources/gabc/Sequences/VictimaeLatin.gabc}
	
	%Whitsun
	\clearpage\subsection{Come, Thou Holy Paraclete}\fancyhead[RO,LE]{\textit{Come, Paraclete}}\fancyhead[RE,LO]{Whitsunday}
		\gregorioscore{resources/gabc/Sequences/VeniSancte.gabc}
	
	\clearpage\subsection{Veni, Sancte Spiritus}\fancyhead[RO,LE]{\textit{Veni, Sancte}}
		\gregorioscore{resources/gabc/Sequences/VeniSancteLatin.gabc}

	%All Souls
	\clearpage\subsection{Day of Wrath and Doom Impending}\fancyhead[RO,LE]{\textit{Day of Wrath}}\fancyhead[RE,LO]{Requiem}
		\gregorioscore{resources/gabc/Sequences/DayOfWrath.gabc}
		
	\clearpage\subsection{Dies irae}\fancyhead[RO,LE]{\textit{Dies Irae}}
		\gregorioscore{resources/gabc/Sequences/DiesIrae.gabc}

\clearpage
\section{Paraliturgical Hymns}\fancyhead[RE,LO]{}

	\subsection{Therefore, We Before Him Bending}\fancyhead[RO,LE]{\textit{Before Him Bending}}
		\gregorioscore{resources/gabc/Extra/Tantum.gabc}
	
	\clearpage\subsection{Tantum Ergo}\fancyhead[RO,LE]{\textit{Tantum Ergo}}
		\gregorioscore{resources/gabc/Extra/TantumLatin.gabc}
	
	\clearpage\subsection{O Saving Victim Opening Wide}\fancyhead[RO,LE]{\textit{O Saving Victim}}
		\gregorioscore{resources/gabc/Extra/OSavingVictim.gabc}
	
	\subsection{O Salutaris Hostia}\fancyhead[RO,LE]{\textit{O Salutaris}}
		\gregorioscore{resources/gabc/Extra/OSalutarisLatin.gabc}
	
	\clearpage\subsection{Prostrate, I Adore Thee}\fancyhead[RO,LE]{\textit{Prostrate, I Adore}}
		\gregorioscore{resources/gabc/Extra/Adoro.gabc}
	
	\clearpage\subsection{Adoro Te Devote}\fancyhead[RO,LE]{\textit{Adoro}}
		\gregorioscore{resources/gabc/Extra/AdoroLatin.gabc}


	\clearpage\subsection{Advent Prose}\fancyhead[RO,LE]{\textit{Advent Prose}}\label{AdventProse}

	\gregorioscore{resources/gabc/Extra/Rorate.gabc}

%\include{pages/Holy_Mass/Mass.tex}
%\include{pages/Holy_Mass/Gregory_2012.tex}

\fancyhead[RE,LO]{}\fancyhead[RO,LE]{}
\fancyhead[C]{}\thispagestyle{empty}
\phantomsection
\addcontentsline{toc}{chapter}{Devotional Hymns}
  \begin{center}
   \includegraphics[width=98mm, height=140mm]{Apocalypse.eps}
   \par
   \vspace{2ex}
   	\textsc{\Huge{Devotional Hymns}}
   \end{center}
\includepdf[pages=-]{hymns/embeddedHymnal.pdf}

%\section*{Appendix: Prayer of Manasseh}\label{PrayerOfManasseh}
%\lett{O}{Lord} Almighty, that art in heaven, thou God of our fathers, of Abraham, and Isaac, and Jacob, and of their righteous seed; who hast made heaven and earth, with all the ornament thereof; who hast bound the sea by the word of thy commandment; who hast shut up the deep, and sealed it by thy terrible and glorious name; whom all things fear, yea, tremble before thy power; for the majesty of thy glory cannot be borne, and the anger of thy threatening toward sinners is importable: thy merciful promise is unmeasurable and unsearchable; for thou art the Lord Most High, of great compassion, longsuffering and abundant in mercy, and repentest of bringing evils upon men. Thou, O Lord, according to thy great goodness hast promised repentance and forgiveness to them that have sinned against thee: and of thine infinite mercies hast appointed repentance unto sinners, that they may be saved. Thou therefore, O Lord, that art the God of the just, hast not appointed repentance to the just, to Abraham, and Isaac, and Jacob, which have not sinned against thee; but thou hast appointed repentance unto me that am a sinner: for I have sinned above the number of the sands of the sea. My transgressions are multiplied, O Lord: my transgressions are multiplied, and I am not worthy to behold and see the height of heaven for the multitude of mine iniquities. I am bowed down with many iron bands, that I cannot lift up mine head by reason of my sins, neither have I any respite: for I have provoked thy wrath, and done that which is evil before thee: I did not thy will, neither kept I thy commandments: I have set up abominations, and have multiplied detestable things. Now therefore I bow the knee of mine heart, beseeching thee of grace. I have sinned, O Lord, I have sinned, and I acknowledge mine iniquities: but, I humbly beseech thee, forgive me, O Lord, forgive me, and destroy me not with mine iniquities. Be not angry with me for ever, by reserving evil for me; neither condemn me into the lower parts of the earth. For thou, O Lord, art the God of them that repent; and in me thou wilt shew all thy goodness: for thou wilt save me, that am unworthy, according to thy great mercy. And I will praise thee for ever all the days of my life: for all the host of heaven doth sing thy praise, and thine is the glory for ever and ever. Amen. 

%\phantomsection
\addcontentsline{toc}{section}{Gregorian Propers}
\fancyhead[C]{\LARGE Gregorian Propers}
\fancyhead[RO,LE]{\textit{Of Saints}}
\fancyhead[RE,LO]{}

\subby{St. Andrew}
\fancyhead[RE,LO]{St. Andrew}
\begin{inhead}
    {Second Class Double\\
30 November}
\end{inhead}

%Considered for BCP:
%\collect
%\lett{W}{e} humbly entreat thy majesty, O Lord: that as blessed Andrew the Apostle was to thy Church a preacher and governor; so he may be a perpetual intercessor for us in thy sight. Through.


\readingcitation{Epistle}{Romans 10:9}
\lett{B}{rethren:} With the heart man believeth unto righteousness; and with the mouth confession is made unto salvation. For the scripture saith, Whosoever believeth on him shall not be ashamed. For there is no difference between the Jew and the Greek: for the same Lord over all is rich unto all that call upon him. For whosoever shall call upon the name of the Lord shall be saved. How then shall they call on him in whom they have not believed? and how shall they believe in him of whom they have not heard? and how shall they hear without a preacher? And how shall they preach, except they be sent? as it is written, How beautiful are the feet of them that preach the gospel of peace, and bring glad tidings of good things! But they have not all obeyed the gospel. For Esaias saith, Lord, who hath believed our report? So then faith cometh by hearing, and hearing by the word of God. But I say, Have they not heard? Yes verily, their sound went into all the earth, and their words unto the ends of the world.

\subby{St. Thomas}
\fancyhead[RE,LO]{St. Thomas}
\begin{inhead}
    {Second Class Double\\
21 December}
\end{inhead}
%Considered for BCP:
%\collect
%\lett{G}{rant} us, we beseech thee, O Lord, to glory in the solemnity of thy blessed Apostle Thomas: that we may ever be succoured by his protection; and follow his faith with worthy devotion. Through.
%\begin{rubric}
%    Commemoration of the Feria, except in Conventual Masses if an Ember Day occur.
%\end{rubric}

%\readingcitation{Epistle}{Ephesians 2:19}
%\lett{B}{rethren:} Now ye are no more strangers and foreigners, but fellowcitizens with the saints, and of the household of God; And are built upon the foundation of the apostles and prophets, Jesus Christ himself being the chief corner stone; In whom all the building fitly framed together groweth unto an holy temple in the Lord: In whom ye also are builded together for an habitation of God through the Spirit.

\readingcitation{Gospel}{John 20:24}
\lett{A}{t that time:} Thomas, one of the twelve, called Didymus, was not with them when Jesus came. The other disciples therefore said unto him, We have seen the Lord. But he said unto them, Except I shall see in his hands the print of the nails, and put my finger into the print of the nails, and thrust my hand into his side, I will not believe. And after eight days again his disciples were within, and Thomas with them: then came Jesus, the doors being shut, and stood in the midst, and said, Peace be unto you. Then saith he to Thomas, Reach hither thy finger, and behold my hands; and reach hither thy hand, and thrust it into my side: and be not faithless, but believing. And Thomas answered and said unto him, My Lord and my God. Jesus saith unto him, Thomas, because thou hast seen me, thou hast believed: blessed are they that have not seen, and yet have believed.

\subby{Conversion of St. Paul}
\fancyhead[RE,LO]{Conversion of St. Paul}
\begin{inhead}
    {Greater Double\\
25 January}
\end{inhead}
\collect
\lett{O}{God,} who by the preaching of blessed Paul the Apostle hast taught the whole world: grant us, we beseech thee; that we, who this day celebrate his Conversion, may through his example draw nearer unto thee. Through.
\needspace{4\baselineskip}
\lett{O}{God,} who didst bestow upon thy blessed Apostle Peter the keys of the kingdom of heaven, and didst appoint unto him the high priesthood of binding and loosing: vouchsafe; that by the help of his intercession we may be delivered from the bonds of our iniquities. Who livest.

\readingcitation{Gospel}{Matthew 19:27}
\lett{A}{t that time:} Peter said unto Jesus: Behold, we have forsaken all, and followed thee; what shall we have therefore? And Jesus said unto them, Verily I say unto you, That ye which have followed me, in the regeneration when the Son of man shall sit in the throne of his glory, ye also shall sit upon twelve thrones, judging the twelve tribes of Israel. And every one that hath forsaken houses, or brethren, or sisters, or father, or mother, or wife, or children, or lands, for my name's sake, shall receive an hundredfold, and shall inherit everlasting life.

\subby{Purification of the Blessed Virgin Mary}
\fancyhead[RE,LO]{Purification of the B.V.M.}
\begin{inhead}
    {Second Class Double\\
2 February}
\end{inhead}

\readingcitation{Epistle}{Malachi 3:1}
\lett{B}{ehold,} I will send my messenger, and he shall prepare the way before me: and the Lord, whom ye seek, shall suddenly come to his temple, even the messenger of the covenant, whom ye delight in: behold, he shall come, saith the \textsc{Lord} of hosts. But who may abide the day of his coming? and who shall stand when he appeareth? for he is like a refiner's fire, and like fullers' soap: And he shall sit as a refiner and purifier of silver: and he shall purify the sons of Levi, and purge them as gold and silver, that they may offer unto the \textsc{Lord} an offering in righteousness. Then shall the offering of Judah and Jerusalem be pleasant unto the \textsc{Lord}, as in the days of old, and as in former years.

\readingcitation{Gospel}{Luke 2:22}
\lett{A}{t that time:} When the days of her purification according to the law of Moses were accomplished, they brought him to Jerusalem, to present him to the Lord; (As it is written in the law of the Lord, Every male that openeth the womb shall be called holy to the Lord;) And to offer a sacrifice according to that which is said in the law of the Lord, A pair of turtledoves, or two young pigeons. And, behold, there was a man in Jerusalem, whose name was Simeon; and the same man was just and devout, waiting for the consolation of Israel: and the Holy Ghost was upon him. And it was revealed unto him by the Holy Ghost, that he should not see death, before he had seen the Lord's Christ. And he came by the Spirit into the temple: and when the parents brought in the child Jesus, to do for him after the custom of the law, Then took he him up in his arms, and blessed God, and said, Lord, now lettest thou thy servant depart in peace, according to thy word: For mine eyes have seen thy salvation, Which thou hast prepared before the face of all people; A light to lighten the Gentiles, and the glory of thy people Israel.

%\subby{St. Matthias}
%\fancyhead[RE,LO]{St. Matthias}
%\begin{inhead}
%    {Second Class Double\\
%24 February (25 February in a Leap Year)}
%\end{inhead}
%Considered for BCP:
%\collect
%\lett{O}{God,} who didst join blessed Matthias to the college of thine Apostles: grant, we beseech thee; that through his mediation we may ever perceive thy tender mercy towards us. Through.
%\begin{rubric}
%    In Lent, Commemoration is made of the Feria in non-conventual Masses.
%\end{rubric}


\subby{St. Benedict}
\fancyhead[RE,LO]{St. Benedict}
\begin{inhead}
    {Second Class Double\\
21 March}
\end{inhead}

\collect
\lett{W}{e} beseech thee, O Lord, that the prayers of the blessed Abbot, Benedict may commend us unto thee: that by his advocacy thou wouldest vouchsafe unto us those things to which by our own merits we cannot attain. Through.
\begin{rubric}
    In Lent, Commemoration is made of the Feria.
\end{rubric}
\readingcitation{Epistle}{Ecclesiasticus 45:1}
\lett{B}{eloved} of God and men, whose memorial is blessed. He made him like to the glorious saints, and magnified him, so that his enemies stood in fear of him. By his words he caused the wonders to cease, and he made him glorious in the sight of kings, and gave him a commandment for his people, and shewed him part of his glory. He sanctified him in his faithfuless and meekness, and chose him out of all men. He made him to hear his voice, and brought him into the dark cloud, and gave him commandments before his face, even the law of life and knowledge.

\gradtr{Thou hast prevented him, O Lord, with the blessings of goodness: thou hast set a crown of pure gold upon his head. ℣. He asked life of thee, and thou gavest him a long life, even for ever and ever.}{Blessed is the man that feareth the Lord: he hath great delight in his commandments. ℣. His seed shall be mighty upon earth: the generation of the faithful shall be blessed. ℣. Riches and plenteousness shall be in his house: and his righteousness endureth for ever.}

\readingcitation{Gospel}{Matthew 19:27}
\lett{A}{t that time:} Peter said unto Jesus: Behold, we have forsaken all, and followed thee; what shall we have therefore? And Jesus said unto them, Verily I say unto you, That ye which have followed me, in the regeneration when the Son of man shall sit in the throne of his glory, ye also shall sit upon twelve thrones, judging the twelve tribes of Israel. And every one that hath forsaken houses, or brethren, or sisters, or father, or mother, or wife, or children, or lands, for my name's sake, shall receive an hundredfold, and shall inherit everlasting life.

\secret
\lett{W}{e} beseech thee, O Lord, that thy holy Abbot Benedict may intercede for us: that this sacrifice which we offer and present upon thy holy altar may be profitable unto us for our salvation. Through.
\begin{rubric}
    In Lent, Commemoration is made of the Feria.
\end{rubric}
\communion{A faithful and wise steward, whom his lord hath made ruler over his household, to give them their portion of meat in due season.}
\postcommunion
\lett{L}{et} thy sacrament, O Lord, which we have now received and the prayers of the blessed Abbot, Benedict, effectually defend us: that we may both imitate the example of his conversation, and receive the succour of his intercession. Through.
\begin{rubric}
    In Lent, Commemoration is made of the Feria.
\end{rubric}
\begin{rubric}
    In Lent, the Last Gospel is of the Feria.
\end{rubric}

\subby{Annunciation of the Blessed Virgin Mary}
\fancyhead[RE,LO]{Annunciation of the B.V.M.}
\begin{inhead}
    {First Class Double\\
25 March}
\end{inhead}

%Considered for BCP:
%\collect
%\lett{O}{God,} who wast pleased that thy Word should take flesh of the womb of the Blessed Virgin Mary at the message of an Angel: grant to thy humble servants; that we who believe her to be truly the Mother of God may be aided by her intercession in thy sight. Through the same.
%\begin{rubric}
%    In Lent, Commemoration is made of the Feria in non-conventual Masses.
%\end{rubric}

\readingcitation{Gospel}{Luke 1:26}
\lett{A}{t that time:} The angel Gabriel was sent from God unto a city of Galilee, named Nazareth, To a virgin espoused to a man whose name was Joseph, of the house of David; and the virgin's name was Mary. And the angel came in unto her, and said, Hail, thou that art highly favoured, the Lord is with thee: blessed art thou among women. And when she saw him, she was troubled at his saying, and cast in her mind what manner of salutation this should be. And the angel said unto her, Fear not, Mary: for thou hast found favour with God. And, behold, thou shalt conceive in thy womb, and bring forth a son, and shalt call his name \divineName{Jesus}. He shall be great, and shall be called the Son of the Highest: and the Lord God shall give unto him the throne of his father David: And he shall reign over the house of Jacob for ever; and of his kingdom there shall be no end. Then said Mary unto the angel, How shall this be, seeing I know not a man? And the angel answered and said unto her, The Holy Ghost shall come upon thee, and the power of the Highest shall overshadow thee: therefore also that holy thing which shall be born of thee shall be called the Son of God. And, behold, thy cousin Elisabeth, she hath also conceived a son in her old age: and this is the sixth month with her, who was called barren. For with God nothing shall be impossible. And Mary said, Behold the handmaid of the Lord; be it unto me according to thy word. 

\postcommunion
\lett{W}{e} beseech thee, O Lord, pour thy grace into our hearts; that, as we have known the incarnation of thy Son Jesus Christ by the message of an Angel, so by his cross and passion we may be brought unto the glory of his resurrection. Through the same.

\begin{rubric}
%CHECK:
    If any Feast fall during Holy Week or the Easter Octave, it is transferred to the next available day.
\end{rubric}

\subby{St. George}
\fancyhead[RE,LO]{St. George}
\begin{inhead}
    {First Class Double with a Common Octave\\
23 April}
\end{inhead}
%Sarum-Roman: 0.2

\collect
%Roman collect with indirect invocation of a saint.
\lett{O}{God,} who makest us glad by the merits and intercession of blessed George, thy Martyr: mercifully grant; that we, who through him ask for thy benefits may by the gift of thy grace obtain the same. Through.

\readingcitation{Epistle}{2 Timothy 2:8}
\lett{D}{early beloved:} Remember that Jesus Christ of the seed of David was raised from the dead according to my gospel: Wherein I suffer trouble, as an evil doer, even unto bonds; but the word of God is not bound. Therefore I endure all things for the elect's sakes, that they may also obtain the salvation which is in Christ Jesus with eternal glory. But thou hast fully known my doctrine, manner of life, purpose, faith, longsuffering, charity, patience, Persecutions, afflictions, which came unto me at Antioch, at Iconium, at Lystra; what persecutions I endured: but out of them all the Lord delivered me. Yea, and all that will live godly in Christ Jesus shall suffer persecution.

\gradual{Blessed is the man that feareth the Lord: he hath great delight in his commandments. ℣. His seed shall be mighty upon earth: the generation of the faithful shall be blessed.}

\tract{Thou hast given him his heart's desire, and hast not denied him the request of his lips. ℣. For thou shalt prevent him with the blessing of goodness. ℣. Thou shalt set a crown of pure gold upon his head.}

\begin{rubric}
    If it be Easter, the Alleluia Verse is said instead of the Gradual \& Tract.
\end{rubric}

\alleluia{Alleluia, alleluia. ℣. O Lord, the very heavens shall praise thy wondrous works: and thy truth in the congregation of the Saints. Alleluia. ℣. Thou hast set, O Lord, a crown of pure gold upon his head. Alleluia.}

\readingcitation{Gospel}{John 15:1}
\lett{A}{t that time:} Jesus said to his disciples: I am the true vine, and my Father is the husbandman. Every branch in me that beareth not fruit he taketh away: and every branch that beareth fruit, he purgeth it, that it may bring forth more fruit. Now ye are clean through the word which I have spoken unto you. Abide in me, and I in you. As the branch cannot bear fruit of itself, except it abide in the vine; no more can ye, except ye abide in me. I am the vine, ye are the branches: He that abideth in me, and I in him, the same bringeth forth much fruit: for without me ye can do nothing. If a man abide not in me, he is cast forth as a branch, and is withered; and men gather them, and cast them into the fire, and they are burned. If ye abide in me, and my words abide in you, ye shall ask what ye will, and it shall be done unto you.
\offertory{O Lord, the very heavens shall praise thy wondrous works: and thy truth in the congregation of the saints, (alleluia, alleluia.)}
\secret
\lett{S}{anctify,} O Lord, the gifts which we offer: and at the intercession of blessed George, thy Martyr, cleanse us thereby from the defilements of our iniquities. Through.
\begin{rubric}
    In Lent, Commemoration is made of the Feria in non-conventual Masses.
\end{rubric}
\communion{The righteous shall rejoice in the Lord, and put his trust in him : and all they that are true of heart shall be glad, (alleluia, alleluia.)}
\postcommunion
\lett{W}{e} humbly beseech thee, almighty God: that as thou dost refresh us with thy sacraments, so at the intercession of blessed George, thy Martyr, thou wouldest grant unto us to do thee worthy and acceptable service. Through.
\begin{rubric}
    In Lent, Commemoration is made of the Feria in non-conventual Masses.
\end{rubric}
\begin{rubric}
    In Lent, the Last Gospel is of the Feria in non-conventual Masses.
\end{rubric}

\subby{St. Mark}
\fancyhead[RE,LO]{St. Mark}
\begin{inhead}
    {Second Class Double\\
25 April}
\end{inhead}
%S-R: N/A
%Latest Date of Gregorian Easter
\collect
%Considered for BCP:
%\lett{O}{God,} who hast exalted blessed Mark thine Evangelist by the grace of the preaching of the Gospel: grant, we beseech thee; that we may ever both profit by his learning, and be defended by his prayer. Through.
%\begin{rubric}
%    In Lent, Commemoration is made of the Feria in non-conventual Masses.
%\end{rubric}

\readingcitation{Epistle}{Ezekiel 1:10}
\lett{A}{s} for the likeness of their faces, they four had the face of a man, and the face of a lion, on the right side: and they four had the face of an ox on the left side; they four also had the face of an eagle. Thus were their faces: and their wings were stretched upward; two wings of every one were joined one to another, and two covered their bodies. And they went every one straight forward: whither the spirit was to go, they went; and they turned not when they went. As for the likeness of the living creatures, their appearance was like burning coals of fire, and like the appearance of lamps: it went up and down among the living creatures; and the fire was bright, and out of the fire went forth lightning. And the living creatures ran and returned as the appearance of a flash of lightning.

\readingcitation{Gospel}{Luke 10:1}
\lett{A}{t that time:} The Lord appointed other seventy also, and sent them two and two before his face into every city and place, whither he himself would come. Therefore said he unto them, The harvest truly is great, but the labourers are few: pray ye therefore the Lord of the harvest, that he would send forth labourers into his harvest. Go your ways: behold, I send you forth as lambs among wolves. Carry neither purse, nor scrip, nor shoes: and salute no man by the way. And into whatsoever house ye enter, first say, Peace be to this house. And if the son of peace be there, your peace shall rest upon it: if not, it shall turn to you again. And in the same house remain, eating and drinking such things as they give: for the labourer is worthy of his hire. Go not from house to house. And into whatsoever city ye enter, and they receive you, eat such things as are set before you: And heal the sick that are therein, and say unto them, The kingdom of God is come nigh unto you.

\subby{Sts. Philip \& James}
\fancyhead[RE,LO]{Sts. Philip \& James}
\begin{inhead}
    {Second Class Double\\
1 May}
\end{inhead}

\collect
\lett{O}{God,} who makest us to rejoice in the yearly solemnity of thine Apostles Philip and James: grant, we beseech thee; that, as we rejoice in their merits, so we may be taught by their example. Through.

\readingcitation{Epistle}{Wisdom 5:1}
\lett{T}{he} righteous man shall stand in great boldness before the face of such as have afflicted him, and made no account of his labours. When they see it, they shall be troubled with terrible fear, and shall be amazed at the strangeness of his salvation, so far beyond all that they looked for. And they repenting and groaning for anguish of spirit shall say within themselves, This was he, whom we had sometimes in derision, and a proverb of reproach: We fools accounted his life madness, and his end to be without honour: How is he numbered among the children of God, and his lot is among the saints!

%DOUBLE CHECK GOSPEL


\subby{St. Barnabas}
\fancyhead[RE,LO]{St. Barnabas}
\begin{inhead}
    {Greater Double\\
11 June}
\end{inhead}
%S-R: N/A
%Sarum considered for BCP.
\collect
\lett{O}{God,} who dost gladden us by the merits and intercession of blessed Barnabas thine Apostle: mercifully grant; that we, who through him ask thy benefits, may obtain them by the gift of thy grace. Through.
\readingcitation{Epistle}{Acts 11:21}%DOUBLE CHECK:
\lett{I}{n those days:} A great number believed, and turned unto the Lord. Tidings of these things came unto the ears of the church which was in Jerusalem: and they sent forth Barnabas, that he should go as far as Antioch. Who, when he came, and had seen the grace of God, was glad, and exhorted them all, that with purpose of heart they would cleave unto the Lord. For he was a good man, and full of the Holy Ghost and of faith: and much people was added unto the Lord. Then departed Barnabas to Tarsus, for to seek Saul. And when he had found him, he brought him unto Antioch. And it came to pass that a whole year they assembled themselves with the church, and taught much people. And the disciples were called Christians first in Antioch. %
Now there were in the church that was at Antioch certain prophets and teachers; as Barnabas, and Simeon that was called Niger, and Lucius of Cyrene, and Manaen, which had been brought up with Herod the tetrarch, and Saul. As they ministered to the Lord, and fasted, the Holy Ghost said, Separate me Barnabas and Saul for the work whereunto I have called them. And when they had fasted and prayed, and laid their hands on them, they sent them away.

\readingcitation{Gospel}{Matthew 10:16}
\lett{A}{t that time:} Jesus said unto his disciples: Behold, I send you forth as sheep in the midst of wolves: be ye therefore wise as serpents, and harmless as doves. But beware of men: for they will deliver you up to the councils, and they will scourge you in their synagogues; And ye shall be brought before governors and kings for my sake, for a testimony against them and the Gentiles. But when they deliver you up, take no thought how or what ye shall speak: for it shall be given you in that same hour what ye shall speak. For it is not ye that speak, but the Spirit of your Father which speaketh in you. And the brother shall deliver up the brother to death, and the father the child: and the children shall rise up against their parents, and cause them to be put to death. And ye shall be hated of all men for my name's sake: but he that endureth to the end shall be saved.

\subby{St. John Baptist}\label{StJohnBaptist}
\fancyhead[RE,LO]{St. John Baptist}
\begin{inhead}
    {First Class Double with a Common Octave\\
24 June}
%(25 June is the last possible day of Eastertide)
\end{inhead}

\collect
\lett{O}{God,} who hast made this day honourable unto us by the nativity of blessed John: grant unto thy people the grace of spiritual joys: and direct the minds of all the faithful into the way of everlasting salvation. Through.

\readingcitation{Epistle}{Isaiah 49:1}
\lett{L}{isten,} O isles, unto me; and hearken, ye people, from far; The \divineName{Lord} hath called me from the womb; from the bowels of my mother hath he made mention of my name. And he hath made my mouth like a sharp sword; in the shadow of his hand hath he hid me, and made me a polished shaft; in his quiver hath he hid me; And said unto me, Thou art my servant, O Israel, in whom I will be glorified. Then I said, I have laboured in vain, I have spent my strength for nought, and in vain: yet surely my judgment is with the \divineName{Lord}, and my work with my God. And now, saith the \divineName{Lord} that formed me from the womb to be his servant, to bring Jacob again to him, Though Israel be not gathered, yet shall I be glorious in the eyes of the \divineName{Lord}, and my God shall be my strength. And he said, It is a light thing that thou shouldest be my servant to raise up the tribes of Jacob, and to restore the preserved of Israel: I will also give thee for a light to the Gentiles, that thou mayest be my salvation unto the end of the earth. Thus saith the \divineName{Lord}, the Redeemer of Israel, and his Holy One, to him whom man despiseth, to him whom the nation abhorreth, to a servant of rulers, Kings shall see and arise, princes also shall worship, because of the \divineName{Lord} that is faithful, and the Holy One of Israel, and he shall choose thee.

\readingcitation{Gospel}{Luke 1:57}
\lett{E}{lisabeth's} full time came that she should be delivered; and she brought forth a son. And her neighbours and her cousins heard how the Lord had shewed great mercy upon her; and they rejoiced with her. And it came to pass, that on the eighth day they came to circumcise the child; and they called him Zacharias, after the name of his father. And his mother answered and said, Not so; but he shall be called John. And they said unto her, There is none of thy kindred that is called by this name. And they made signs to his father, how he would have him called. And he asked for a writing table, and wrote, saying, His name is John. And they marvelled all. And his mouth was opened immediately, and his tongue loosed, and he spake, and praised God. And fear came on all that dwelt round about them: and all these sayings were noised abroad throughout all the hill country of Judaea. And all they that heard them laid them up in their hearts, saying, What manner of child shall this be! And the hand of the Lord was with him. And his father Zacharias was filled with the Holy Ghost, and prophesied, saying, Blessed be the Lord God of Israel; for he hath visited and redeemed his people.

\subby{Sts. Peter \& Paul}
\fancyhead[RE,LO]{Sts. Peter \& Paul}
\begin{inhead}
    {First Class Double with a Common Octave\\
29 June}
\end{inhead}

\collect
\lett{O}{God,} who hast hallowed this day by the martyrdom of thine Apostles Peter and Paul: grant unto thy Church in all things to follow the commandment of those; through whom she received the beginning of religion. Through.

\subby{St. Mary Magdalene}
\fancyhead[RE,LO]{St. Mary Magdalene}
\begin{inhead}
    {Greater Double\\
22 July}
\end{inhead}

\collect
\lett{W}{e} beseech thee, O Lord, that as thou didst hearken to the prayers of blessed Mary Magdalene, and on the fourth day didst raise her brother Lazarus from the dead to life: so we may be succoured by her intercession. Who livest.

\readingcitation{Epistle}{Song of Songs 3:2}
\lett{I}{will} rise, and go about the city in the streets, and in the broad ways I will seek him whom my soul loveth: I sought him, but I found him not. The watchmen that go about the city found me: to whom I said, Saw ye him whom my soul loveth? It was but a little that I passed from them, but I found him whom my soul loveth: I held him, and would not let him go, until I had brought him into my mother's house, and into the chamber of her that conceived me. I charge you, O ye daughters of Jerusalem, by the roes, and by the hinds of the field, that ye stir not up, nor awake my love, till he please. Set me as a seal upon thine heart, as a seal upon thine arm: for love is strong as death; jealousy is cruel as the grave: the coals thereof are coals of fire, which hath a most vehement flame. Many waters cannot quench love, neither can the floods drown it: if a man would give all the substance of his house for love, it would utterly be contemned.

\subby{St. James}
\fancyhead[RE,LO]{St. James}
\begin{inhead}
    {Second Class Double\\
25 July}
\end{inhead}

\collect
\lett{B}{e} thou, O Lord, the sanctifier and defender of thy people: that being aided by the protection of thine Apostle James, they may both please thee by their conversation, and serve thee with a quiet mind. Through.

\readingcitation{Epistle}{1 Corinthians 4:9}
\lett{B}{rethren:} I think that God hath set forth us the apostles last, as it were appointed to death: for we are made a spectacle unto the world, and to angels, and to men. We are fools for Christ's sake, but ye are wise in Christ; we are weak, but ye are strong; ye are honourable, but we are despised. Even unto this present hour we both hunger, and thirst, and are naked, and are buffeted, and have no certain dwellingplace; And labour, working with our own hands: being reviled, we bless; being persecuted, we suffer it: Being defamed, we intreat: we are made as the filth of the world, and are the offscouring of all things unto this day. I write not these things to shame you, but as my beloved sons I warn you. For though ye have ten thousand instructors in Christ, yet have ye not many fathers: for in Christ Jesus I have begotten you through the gospel.

\readingcitation{Gospel}{Matthew 20:20}
\lett{A}{t that time:} There came unto Jesus the mother of Zebedee's children with her sons, worshipping him, and desiring a certain thing of him. And he said unto her, What wilt thou? She saith unto him, Grant that these my two sons may sit, the one on thy right hand, and the other on the left, in thy kingdom. But Jesus answered and said, Ye know not what ye ask. Are ye able to drink of the cup that I shall drink of, and to be baptized with the baptism that I am baptized with? They say unto him, We are able. And he saith unto them, Ye shall drink indeed of my cup, and be baptized with the baptism that I am baptized with: but to sit on my right hand, and on my left, is not mine to give, but it shall be given to them for whom it is prepared of my Father.

\subby{Transfiguration of Our Lord Jesus Christ}
\fancyhead[RE,LO]{Transfiguration}
\begin{inhead}
    {Second Class Double\\
6 August}
\end{inhead}

\introit
\lett{T}{he} lightnings shone upon the ground: the earth was moved, and shook withal. \textit{Ps.} O how amiable are thy dwellings, thou Lord of hosts! My soul hath a desire and longing to enter into the courts of the Lord.
\collect
\lett{O}{God,} who in the glorious Transfiguration of thine only-begotten Son didst confirm the mysteries of the faith by the testimony of the fathers, and in the voice proceeding from the bright cloud didst wondrously foreshow the perfect adoption of thy sons; mercifully vouchsafe; that we, being made fellow-heirs of the same the King of glory, may likewise attain to be partakers of that glory. Through the same.
\readingcitation{Epistle}{2 Peter 1:16}
%The 1928 BCP reading begins three verses earlier than the Roman.
\lett{D}{early beloved:} We have not followed cunningly devised fables, when we made known unto you the power and coming of our Lord Jesus Christ, but were eyewitnesses of his majesty. For he received from God the Father honour and glory, when there came such a voice to him from the excellent glory, This is my beloved Son, in whom I am well pleased. And this voice which came from heaven we heard, when we were with him in the holy mount. We have also a more sure word of prophecy; whereunto ye do well that ye take heed, as unto a light that shineth in a dark place, until the day dawn, and the day star arise in your hearts.

\gradall{Thou art fairer than the children of men: full of grace are thy lips. ℣. My heart is inditing of a good matter: I speak of the things which I have made unto the King.}{Alleluia, alleluia. ℣. He is the brightness of the everlasting light, the unspotted mirror, and the image of his goodness. Alleluia.}

\readingcitation{Gospel}{Matthew 17:1}
\lett{A}{t that time:} Jesus taketh Peter, James, and John his brother, and bringeth them up into an high mountain apart, And was transfigured before them: and his face did shine as the sun, and his raiment was white as the light. And, behold, there appeared unto them Moses and Elias talking with him. Then answered Peter, and said unto Jesus, Lord, it is good for us to be here: if thou wilt, let us make here three tabernacles; one for thee, and one for Moses, and one for Elias. While he yet spake, behold, a bright cloud overshadowed them: and behold a voice out of the cloud, which said, This is my beloved Son, in whom I am well pleased; hear ye him. And when the disciples heard it, they fell on their face, and were sore afraid. And Jesus came and touched them, and said, Arise, and be not afraid. And when they had lifted up their eyes, they saw no man, save Jesus only. And as they came down from the mountain, Jesus charged them, saying, Tell the vision to no man, until the Son of man be risen again from the dead.

\offertory{Riches and plenteousness shall be in his house: and his righteousness endureth for ever, alleluia.}

\secret
\lett{S}{anctify,} we beseech thee, O Lord, by the glorious Transfiguration of thine only-begotten Son, the gifts which we offer unto thee: and by the splendour of his brightness, cleanse us from the defilements of our iniquities. Through the same.

\communion{Tell the vision which ye have seen to no man, until the Son of man be risen again from the dead.}

\postcommunion
\lett{G}{rant,} we beseech thee, almighty God: that we, who celebrate in solemn worship the most holy mysteries of the Transfiguration of thy Son, may attain thereto in pureness of mind and understanding. Through the same.

\subby{Assumption of the Blessed Virgin Mary}
\fancyhead[RE,LO]{Assumption of the B.V.M.}
\begin{inhead}
    {First Class Double with Common Octave\\
15 August}
\end{inhead}

\introit
\lett{A}{great} wonder appeared in heaven: a woman clothed with the sun, and the moon under her feet, and upon her head a crown of twelve stars. \textit{Ps.} O sing unto the Lord a new song: for he hath done marvellous things.
\collect
\lett{A}{lmighty} and everlasting God, who didst assume the immaculate Virgin Mary, mother of thy Son, in body and soul to the glory of heaven: grant, we beseech thee; that, being ever intent on things above, we may be worthy to be partakers of her glory. Through the same.
%Old Roman: {\lett{F}{orgive,} O Lord, we beseech thee, the offences of thy servants: that we, who of ourselves can do nothing which is acceptable unto thee; may be saved by the intercession of the Mother of thy Son our Lord. Who liveth.}
\readingcitation{Epistle}{Judith 13:17}
\lett{B}{lessed} be thou, O our God, which hast this day brought to nought the enemies of thy people. O daughter, blessed art thou of the most high God above all the women upon the earth; and blessed be the Lord God, which hath created the heavens and the earth, which hath directed thee to the cutting off of the head of the chief of our enemies. For this thy confidence shall not depart from the heart of men, which remember the power of God for ever. And God turn these things to thee for a perpetual praise, to visit thee in good things because thou hast not spared thy life for the affliction of our nation, but hast revenged our ruin, walking a straight way before our God.

\gradall{Hearken, O daughter, and consider, and incline thine ear, so shall the king have pleasure in thy beauty. ℣. All glorious the king's daughter entereth in, her clothing is of wrought gold.}{Alleluia, alleluia. ℣. Mary is taken up into heaven: the host of Angels rejoice. Alleluia.}

\readingcitation{Gospel}{Luke 1:41}

\lett{A}{t that time:} And it came to pass, that, when Elisabeth heard the salutation of Mary, the babe leaped in her womb; and Elisabeth was filled with the Holy Ghost: and she spake out with a loud voice, and said, Blessed art thou among women, and blessed is the fruit of thy womb. And whence is this to me, that the mother of my Lord should come to me? For, lo, as soon as the voice of thy salutation sounded in mine ears, the babe leaped in my womb for joy. And blessed is she that believed: for there shall be a performance of those things which were told her from the Lord. And Mary said, My soul doth magnify the Lord, and my spirit hath rejoiced in God my Saviour. For he hath regarded the low estate of his handmaiden: for, behold, from henceforth all generations shall call me blessed. For he that is mighty hath done to me great things; and holy is his name. And his mercy is on them that fear him from generation to generation.

\offertory{I will put enmity between thee and the Woman, and between thy seed and her Seed.}

\secret
\lett{L}{et} this oblation of our bounden duty ascend unto thee, O Lord, and, at the intercession of the most blessed Virgin Mary, assumed into heaven, may our hearts, enkindled with the fire of charity, continually yearn after thee. Through.

\communion{All generations shall call me blessed, for he that is mighty hath done to me great things.}
\postcommunion
\lett{G}{rant,} we beseech thee, O Lord: that we, who have received the sacraments of salvation; may, through the merits and intercession of the blessed Virgin Mary, assumed into heaven, be brought unto the glory of the resurrection. Through.

\begin{rubric}
    Within the Octave, Mass is said as on the Feast. And the \nth{2} Collect is of the Holy Ghost (p. \pageref{SPHolyGhost}) \& \nth{3} Against the Persecutors of the Church (p. \pageref{SPAgainst}) or for the Chief Bishop (p. \pageref{SPChiefBishop}).
\end{rubric}

\subby{St. Bartholomew}
\fancyhead[RE,LO]{St. Bartholomew}
\begin{inhead}
    {Second Class Double\\
24 August}
\end{inhead}

\collect
\lett{A}{lmighty} and everlasting God, who hast given us on this day in the festival of blessed Bartholomew thine Apostle venerable and holy joy: grant unto thy Church, we beseech thee; both to love that which he believed, and to preach that which he taught. Through.

\readingcitation{Epistle}{1 Corinthians 12:27}
\lett{B}{rethren:} Ye are the body of Christ, and members in particular. And God hath set some in the church, first apostles, secondarily prophets, thirdly teachers, after that miracles, then gifts of healings, helps, governments, diversities of tongues. Are all apostles? are all prophets? are all teachers? are all workers of miracles? Have all the gifts of healing? do all speak with tongues? do all interpret? But covet earnestly the best gifts.


\readingcitation{Gospel}{Luke 6:12}
\lett{A}{t that time:} Jesus went out into a mountain to pray, and continued all night in prayer to God. And when it was day, he called unto him his disciples: and of them he chose twelve, whom also he named apostles; Simon, (whom he also named Peter,) and Andrew his brother, James and John, Philip and Bartholomew, Matthew and Thomas, James the son of Alphaeus, and Simon called Zelotes, And Judas the brother of James, and Judas Iscariot, which also was the traitor. And he came down with them, and stood in the plain, and the company of his disciples, and a great multitude of people out of all Judaea and Jerusalem, and from the sea coast of Tyre and Sidon, which came to hear him, and to be healed of their diseases; And they that were vexed with unclean spirits: and they were healed. And the whole multitude sought to touch him: for there went virtue out of him, and healed them all.

\subby{St. Matthew}
\fancyhead[RE,LO]{St. Matthew}
\begin{inhead}
    {Second Class Double\\
21 September}
\end{inhead}

%Considered for BCP:
%\collect
%\lett{M}{ay} we be assisted, O Lord, by the prayers of blessed Matthew, the Apostle and Evangelist: that those things, which of ourselves we cannot obtain, may be vouchsafed unto us by his intercession. Through.

\readingcitation{Epistle}{Ezekiel 1:10}
\lett{T}{he} likeness of the faces of the four animals: they four had the face of a man, and the face of a lion, on the right side: and they four had the face of an ox on the left side; they four also had the face of an eagle. Thus were their faces: and their wings were stretched upward; two wings of every one were joined one to another, and two covered their bodies. And they went every one straight forward: whither the spirit was to go, they went; and they turned not when they went. As for the likeness of the living creatures, their appearance was like burning coals of fire, and like the appearance of lamps: it went up and down among the living creatures; and the fire was bright, and out of the fire went forth lightning. And the living creatures ran and returned as the appearance of a flash of lightning.

\readingcitation{Gospel}{Matthew 9:9}
\lett{A}{t that time:} As Jesus passed forth from thence, he saw a man, named Matthew, sitting at the receipt of custom: and he saith unto him, Follow me. And he arose, and followed him. And it came to pass, as Jesus sat at meat in the house, behold, many publicans and sinners came and sat down with him and his disciples. And when the Pharisees saw it, they said unto his disciples, Why eateth your Master with publicans and sinners? But when Jesus heard that, he said unto them, They that be whole need not a physician, but they that are sick. But go ye and learn what that meaneth, I will have mercy, and not sacrifice: for I am not come to call the righteous, but sinners.

\subby{St. Michael}
\fancyhead[RE,LO]{St. Michael}
\begin{inhead}
    {First Class Double\\
29 September}
\end{inhead}

\readingcitation{Epistle}{Revelation 1:1}

\lett{I}{n those days:} God shewed the things which must shortly come to pass; and he sent and signified it by his angel unto his servant John: Who bare record of the word of God, and of the testimony of Jesus Christ, and of all things that he saw. Blessed is he that readeth, and they that hear the words of this prophecy, and keep those things which are written therein: for the time is at hand. John to the seven churches which are in Asia: Grace be unto you, and peace, from him which is, and which was, and which is to come; and from the seven Spirits which are before his throne; And from Jesus Christ, who is the faithful witness, and the first begotten of the dead, and the prince of the kings of the earth. Unto him that loved us, and washed us from our sins in his own blood.

\subby{St. Luke}
\fancyhead[RE,LO]{St. Luke}
\begin{inhead}
    {Second Class Double\\
18 October}
\end{inhead}

%Considered for BCP:
%\collect
%\lett{W}{e} beseech thee, O Lord: that thy holy Evangelist Luke may intercede for us: who for the honour of thy name continually bare in his own body the mortification of the Cross. Through.

\readingcitation{Epistle}{2 Corinthians 8:16}
\lett{B}{rethren:} I give thanks to God, which put the same earnest care into the heart of Titus for you. For indeed he accepted the exhortation; but being more forward, of his own accord he went unto you. And we have sent with him the brother, whose praise is in the gospel throughout all the churches; And not that only, but who was also chosen of the churches to travel with us with this grace, which is administered by us to the glory of the same Lord, and declaration of your ready mind: Avoiding this, that no man should blame us in this abundance which is administered by us: Providing for honest things, not only in the sight of the Lord, but also in the sight of men. And we have sent with them our brother, whom we have oftentimes proved diligent in many things, but now much more diligent, upon the great confidence which I have in you. Whether any do enquire of Titus, he is my partner and fellowhelper concerning you: or our brethren be enquired of, they are the messengers of the churches, and the glory of Christ. Wherefore shew ye to them, and before the churches, the proof of your love, and of our boasting on your behalf.

\readingcitation{Gospel}{Luke 10:1}
\lett{A}{t that time:} The Lord appointed other seventy also, and sent them two and two before his face into every city and place, whither he himself would come. Therefore said he unto them, The harvest truly is great, but the labourers are few: pray ye therefore the Lord of the harvest, that he would send forth labourers into his harvest. Go your ways: behold, I send you forth as lambs among wolves. Carry neither purse, nor scrip, nor shoes: and salute no man by the way. And into whatsoever house ye enter, first say, Peace be to this house. And if the son of peace be there, your peace shall rest upon it: if not, it shall turn to you again. And in the same house remain, eating and drinking such things as they give: for the labourer is worthy of his hire. Go not from house to house. And into whatsoever city ye enter, and they receive you, eat such things as are set before you: And heal the sick that are therein, and say unto them, The kingdom of God is come nigh unto you.

\subby{Sts. Simon and Jude}
\fancyhead[RE,LO]{Sts. Simon \& Jude}
\begin{inhead}
    {Second Class Double\\
28 October}
\end{inhead}

\collect
\lett{O}{God,} who through thy blessed Apostles Simon and Jude hast suffered us to come unto the knowledge of thy name: grant us by advancing in virtue to celebrate their everlasting glory, and by the celebration thereof to advance in virtue. Through.

\readingcitation{Epistle}{Ephesians 4:7}
\lett{B}{rethren:} Unto every one of us is given grace according to the measure of the gift of Christ. Wherefore he saith, When he ascended up on high, he led captivity captive, and gave gifts unto men. (Now that he ascended, what is it but that he also descended first into the lower parts of the earth? He that descended is the same also that ascended up far above all heavens, that he might fill all things.) And he gave some, apostles; and some, prophets; and some, evangelists; and some, pastors and teachers; For the perfecting of the saints, for the work of the ministry, for the edifying of the body of Christ: Till we all come in the unity of the faith, and of the knowledge of the Son of God, unto a perfect man, unto the measure of the stature of the fulness of Christ:

\readingcitation{Gospel}{John 15:17}
\lett{A}{t that time:} Jesus said unto his disciples: 17 These things I command you, that ye love one another. If the world hate you, ye know that it hated me before it hated you. If ye were of the world, the world would love his own: but because ye are not of the world, but I have chosen you out of the world, therefore the world hateth you. Remember the word that I said unto you, The servant is not greater than his lord. If they have persecuted me, they will also persecute you; if they have kept my saying, they will keep yours also. But all these things will they do unto you for my name's sake, because they know not him that sent me. If I had not come and spoken unto them, they had not had sin: but now they have no cloke for their sin. He that hateth me hateth my Father also. If I had not done among them the works which none other man did, they had not had sin: but now have they both seen and hated both me and my Father. But this cometh to pass, that the word might be fulfilled that is written in their law, They hated me without a cause.

\subby{All Hallows}
\fancyhead[RE,LO]{All Hallows}
\begin{inhead}
    {Second Class Double\\
1 November}
\end{inhead}

%Considered for BCP:
%\collect
%\lett{A}{lmighty} and everlasting God, who in one solemnity hast vouchsafed unto us to venerate the merits of all thy Saints: we beseech thee; that, at the intercession of so great a multitude, thou wouldest bestow upon us, who entreat thee, the abundance of thy mercy. Through.

%RE-CHECK:
\readingcitation{Epistle}{Revelation 7:2}
\lett{I}{n those days:} Behold, I, John, saw another angel ascending from the east, having the seal of the living God: and he cried with a loud voice to the four angels, to whom it was given to hurt the earth and the sea, Saying, Hurt not the earth, neither the sea, nor the trees, till we have sealed the servants of our God in their foreheads. And I heard the number of them which were sealed: and there were sealed an hundred and forty and four thousand of all the tribes of the children of Israel. Of the tribe of Juda were sealed twelve thousand. Of the tribe of Reuben were sealed twelve thousand. Of the tribe of Gad were sealed twelve thousand. Of the tribe of Aser were sealed twelve thousand. Of the tribe of Nepthalim were sealed twelve thousand. Of the tribe of Manasses were sealed twelve thousand. Of the tribe of Simeon were sealed twelve thousand. Of the tribe of Levi were sealed twelve thousand. Of the tribe of Issachar were sealed twelve thousand. Of the tribe of Zabulon were sealed twelve thousand. Of the tribe of Joseph were sealed twelve thousand. Of the tribe of Benjamin were sealed twelve thousand. After this I beheld, and, lo, a great multitude, which no man could number, of all nations, and kindreds, and people, and tongues, stood before the throne, and before the Lamb, clothed with white robes, and palms in their hands; And cried with a loud voice, saying, Salvation to our God which sitteth upon the throne, and unto the Lamb. And all the angels stood round about the throne, and about the elders and the four beasts, and fell before the throne on their faces, and worshipped God, Saying, Amen: Blessing, and glory, and wisdom, and thanksgiving, and honour, and power, and might, be unto our God for ever and ever. Amen.

\readingcitation{Gospel}{Matthew 5:1}
\lett{A}{t that time:} Jesus, seeing the multitudes, went up into a mountain: and when he was set, his disciples came unto him: And he opened his mouth, and taught them, saying, Blessed are the poor in spirit: for theirs is the kingdom of heaven. Blessed are they that mourn: for they shall be comforted. Blessed are the meek: for they shall inherit the earth. Blessed are they which do hunger and thirst after righteousness: for they shall be filled. Blessed are the merciful: for they shall obtain mercy. Blessed are the pure in heart: for they shall see God. Blessed are the peacemakers: for they shall be called the children of God. Blessed are they which are persecuted for righteousness' sake: for theirs is the kingdom of heaven. Blessed are ye, when men shall revile you, and persecute you, and shall say all manner of evil against you falsely, for my sake. Rejoice, and be exceeding glad: for great is your reward in heaven.



%\fancyhead[RO,LE]{\textit{Of Season}}
\fancyhead[RE,LO]{}

   \fancyhead[RE,LO]{Advent I}
   \subby{First Sunday of Advent}\label{AdventI}
\begin{inhead}
    {First Class Semidouble}
\end{inhead}

\collect
\lett{S}{tir} up, we beseech thee, O Lord, thy power and come: that by thy protection we may be found worthy to be set free from the dangers of our sins which beset us, and to be saved by thy deliverance. Who livest.
\begin{rubric}
    This Collect is to be repeated every day, with the other Collects in Advent, until Christmas Day.
\end{rubric}
\begin{rubric}
    \nth{2} Collect is \emph{Of St. Mary in Advent} (p. \pageref{SPMaryInAdvent}) \& \nth{3} \emph{Against the Persecutors of the Church} (p. \pageref{SPAgainst}) or \emph{For the Chief Bishop} (p. \pageref{SPChiefBishop}).
\end{rubric}
\readingcitation{Epistle}{Romans 13:8}
\lett{B}{rethren:} Know the time, that now it is high time to awake out of sleep: for now is our salvation nearer than when we believed. The night is far spent, the day is at hand: let us therefore cast off the works of darkness, and let us put on the armour of light. Let us walk honestly, as in the day; not in rioting and drunkenness, not in chambering and wantonness, not in strife and envying. But put ye on the Lord Jesus Christ.

\readingcitation{Gospel}{Luke 21:25}
\lett{A}{t that time:} Jesus said unto his disciples: There shall be signs in the sun, and in the moon, and in the stars; and upon the earth distress of nations, with perplexity; the sea and the waves roaring; Men's hearts failing them for fear, and for looking after those things which are coming on the earth: for the powers of heaven shall be shaken. And then shall they see the Son of man coming in a cloud with power and great glory. And when these things begin to come to pass, then look up, and lift up your heads; for your redemption draweth nigh. And he spake to them a parable; Behold the fig tree, and all the trees; When they now shoot forth, ye see and know of your own selves that summer is now nigh at hand. So likewise ye, when ye see these things come to pass, know ye that the kingdom of God is nigh at hand. Verily I say unto you, This generation shall not pass away, till all be fulfilled. Heaven and earth shall pass away: but my words shall not pass away.

\subby{Second Sunday of Advent}
\fancyhead[RE,LO]{Advent II}
\begin{inhead}
{Second Class Semidouble}
\end{inhead}

\collect
\lett{S}{tir} up our hearts, O Lord, to prepare the ways of thine only-begotten Son: that through his advent we may be worthy to serve thee with pure minds. Who liveth.
\readingcitation{Gospel}{Matthew 11:2}
\lett{A}{t that time:} When John had heard in the prison the works of Christ, he sent two of his disciples, And said unto him, Art thou he that should come, or do we look for another? Jesus answered and said unto them, Go and shew John again those things which ye do hear and see: The blind receive their sight, and the lame walk, the lepers are cleansed, and the deaf hear, the dead are raised up, and the poor have the gospel preached to them. And blessed is he, whosoever shall not be offended in me. And as they departed, Jesus began to say unto the multitudes concerning John, What went ye out into the wilderness to see? A reed shaken with the wind? But what went ye out for to see? A man clothed in soft raiment? behold, they that wear soft clothing are in kings' houses. But what went ye out for to see? A prophet? yea, I say unto you, and more than a prophet. For this is he, of whom it is written, Behold, I send my messenger before thy face, which shall prepare thy way before thee.

\subby{Third Sunday of Advent}
\fancyhead[RE,LO]{Advent III}
\begin{inhead}
{Second Class Semidouble}
\end{inhead}

\collect
\lett{I}{ncline} thine ear to our prayers, we beseech thee, O Lord: and by the grace of thy visitation enlighten the darkness of our minds. Who livest.
\begin{rubric}
    Additional Collects are as on \nth{1} Sunday of Advent (p. \pageref{AdventI}).
\end{rubric}
\readingcitation{Epistle}{Philippians 4:4}
\lett{B}{rethren:} Rejoice in the Lord alway: and again I say, Rejoice. Let your moderation be known unto all men. The Lord is at hand. Be careful for nothing; but in every thing by prayer and supplication with thanksgiving let your requests be made known unto God. And the peace of God, which passeth all understanding, shall keep your hearts and minds through Christ Jesus our Lord.

\readingcitation{Gospel}{John 1:19}
\lett{A}{t that time:} The Jews sent priests and Levites from Jerusalem to ask him, Who art thou? And he confessed, and denied not; but confessed, I am not the Christ. And they asked him, What then? Art thou Elias? And he saith, I am not. Art thou that prophet? And he answered, No. Then said they unto him, Who art thou? that we may give an answer to them that sent us. What sayest thou of thyself? He said, I am the voice of one crying in the wilderness, Make straight the way of the Lord, as said the prophet Esaias. And they which were sent were of the Pharisees. And they asked him, and said unto him, Why baptizest thou then, if thou be not that Christ, nor Elias, neither that prophet? John answered them, saying, I baptize with water: but there standeth one among you, whom ye know not; He it is, who coming after me is preferred before me, whose shoe's latchet I am not worthy to unloose. These things were done in Bethabara beyond Jordan, where John was baptizing.

\subby{Fourth Sunday of Advent}
\fancyhead[RE,LO]{Advent IV}
\begin{inhead}
{Second Class Semidouble}
\end{inhead}

\collect
\lett{S}{tir} up, we beseech thee, O Lord, thy power, and come: and with great might succour us; that whereas through our sins we are hindered, we may through the help of thy grace speedily obtain thy merciful pardon. Who livest.
\begin{rubric}
    Additional Collects are as on \nth{1} Sunday of Advent (p. \pageref{AdventI}).
\end{rubric}

\readingcitation{Epistle}{1 Corinthians 4:1}
\lett{B}{rethren:} Let a man so account of us, as of the ministers of Christ, and stewards of the mysteries of God. Moreover it is required in stewards, that a man be found faithful. But with me it is a very small thing that I should be judged of you, or of man's judgment: yea, I judge not mine own self. For I know nothing by myself; yet am I not hereby justified: but he that judgeth me is the Lord. Therefore judge nothing before the time, until the Lord come, who both will bring to light the hidden things of darkness, and will make manifest the counsels of the hearts: and then shall every man have praise of God.

\readingcitation{Gospel}{Luke 3:1}
\lett{I}{n} the fifteenth year of the reign of Tiberius Caesar, Pontius Pilate being governor of Judaea, and Herod being tetrarch of Galilee, and his brother Philip tetrarch of Ituraea and of the region of Trachonitis, and Lysanias the tetrarch of Abilene, Annas and Caiaphas being the high priests, the word of God came unto John the son of Zacharias in the wilderness. And he came into all the country about Jordan, preaching the baptism of repentance for the remission of sins; As it is written in the book of the words of Esaias the prophet, saying, The voice of one crying in the wilderness, Prepare ye the way of the Lord, make his paths straight. Every valley shall be filled, and every mountain and hill shall be brought low; and the crooked shall be made straight, and the rough ways shall be made smooth; And all flesh shall see the salvation of God.

\subby{Nativity of Our Lord}
\fancyhead[RE,LO]{Christmas}
\begin{inhead}
{First Class Double, Third Class Octave}
\end{inhead}
\centerline{\small{(The First Mass at Night)}}

\introit
\lett{T}{he} Lord hath said unto me: Thou art my Son, this day have I begotten thee. \textit{Ps.} Why do the heathen so furiously rage together: and why do the people imagine a vain thing?
\collect
\lett{O}{God,} who hast made this most sacred night to shine with the brightness of the true light: grant, we beseech thee; that we, who have known the mysteries of his light on earth, may also attain to the fruition of his joys in heaven. Who liveth.

\subby{St. Stephen}
\fancyhead[RE,LO]{St. Stephen}
\begin{inhead}
{26 December}\par
{Second Class Double, Simple Octave}
\end{inhead}

\collect
%Roman Collect without intercession of the saint
\lett{G}{rant} us, we beseech thee, O Lord, so to imitate that which we revere: that we may learn to love even our enemies; for as much as we celebrate the birthday of him who knew how to pray even for his persecutors to our Lord Jesus Christ thy Son. Who liveth.

\readingcitation{Epistle}{Acts 6:8}
\lett{I}{n those days:} Stephen, full of faith and power, did great wonders and miracles among the people. Then there arose certain of the synagogue, which is called the synagogue of the Libertines, and Cyrenians, and Alexandrians, and of them of Cilicia and of Asia, disputing with Stephen. And they were not able to resist the wisdom and the spirit by which he spake. When they heard these things, they were cut to the heart, and they gnashed on him with their teeth. But Stephen, being full of the Holy Ghost, looked up stedfastly into heaven, and saw the glory of God, and Jesus standing on the right hand of God, And said, Behold, I see the heavens opened, and the Son of man standing on the right hand of God. Then they cried out with a loud voice, and stopped their ears, and ran upon him with one accord, And cast him out of the city, and stoned him: and the witnesses laid down their clothes at a young man's feet, whose name was Saul. And they stoned Stephen, calling upon God, and saying, Lord Jesus, receive my spirit. And he kneeled down, and cried with a loud voice, Lord, lay not this sin to their charge. And when he had said this, he fell asleep.

\subby{St. John}
\fancyhead[RE,LO]{St. John}
\begin{inhead}
{27 December}\par
{Second Class Double, Simple Octave}
\end{inhead}

\collect
%Roman collect without intercession of the Saints
\lett{O}{Lord,} graciously illumine thy Church: that it, being enlightened by the doctrine of blessed John thy Apostle and Evangelist, may attain unto everlasting gifts. Through.
\begin{rubric}
    Commemoration is made of the Nativity from the Third Mass (p. \pageref{NativityMassIII}).
\end{rubric}

\readingcitation{Epistle}{Ecclesiasticus 15:1}
\lett{H}{e} that feareth the Lord will do good, and he that hath the knowledge of the law shall obtain her. And as a mother shall she meet him, and receive him as a wife married of a virgin. With the bread of understanding shall she feed him, and give him the water of wisdom to drink. He shall be stayed upon her, and shall not be moved; and shall rely upon her, and shall not be confounded. She shall exalt him above his neighbours, and in the midst of the congregation shall she open his mouth. He shall find joy and a crown of gladness, and she shall cause him to inherit an everlasting name.

\readingcitation{Gospel}{John 21:19}
\lett{A}{t that time:} Jesus said unto Peter: Follow me. Then Peter, turning about, seeth the disciple whom Jesus loved following; which also leaned on his breast at supper, and said, Lord, which is he that betrayeth thee? Peter seeing him saith to Jesus, Lord, and what shall this man do? Jesus saith unto him, If I will that he tarry till I come, what is that to thee? follow thou me. Then went this saying abroad among the brethren, that that disciple should not die: yet Jesus said not unto him, He shall not die; but, If I will that he tarry till I come, what is that to thee? This is the disciple which testifieth of these things, and wrote these things: and we know that his testimony is true.


\subby{Holy Innocents}
\fancyhead[RE,LO]{Holy Innocents}
\begin{inhead}
{28 December}\par
{Second Class Double, Simple Octave}
\end{inhead}

\collect
\lett{O}{God,} whose praise the Innocents, thy Martyrs did this day, not by speaking, but by dying, confess: mortify all evil vices in us; that, as our tongues do confess thy faith, so our lives may by their actions set forth the same. Through.
\begin{rubric}
    Commemoration is made of the Nativity from the Third Mass (p. \pageref{NativityMassIII}).
\end{rubric}

\subby{Sunday within the Nativity Octave}
\fancyhead[RE,LO]{Christmas I}
\begin{inhead}
{Semidouble}
\end{inhead}

\collect
\lett{A}{lmighty} and everlasting God, direct our actions according to thy good pleasure: that in the name of thy well-beloved Son we may be made worthy to abound in good works. Who liveth and reigneth with thee.

\readingcitation{Gospel}{Luke 2:33}
\lett{A}{t that time:} Joseph and Mary, the Mother of Jesus, marvelled at those things which were spoken of him. And Simeon blessed them, and said unto Mary his mother, Behold, this child is set for the fall and rising again of many in Israel; and for a sign which shall be spoken against; (Yea, a sword shall pierce through thy own soul also,) that the thoughts of many hearts may be revealed. And there was one Anna, a prophetess, the daughter of Phanuel, of the tribe of Aser: she was of a great age, and had lived with an husband seven years from her virginity; And she was a widow of about fourscore and four years, which departed not from the temple, but served God with fastings and prayers night and day. And she coming in that instant gave thanks likewise unto the Lord, and spake of him to all them that looked for redemption in Jerusalem. And when they had performed all things according to the law of the Lord, they returned into Galilee, to their own city Nazareth. And the child grew, and waxed strong in spirit, filled with wisdom: and the grace of God was upon him.

\subby{Circumcision of Our Lord}
\fancyhead[RE,LO]{Circumcision of Our Lord}
\begin{inhead}
{1 January}\par
{Second Class Double}
\end{inhead}

\collect
\lett{O}{God,} who by the virgin child-bearing of blessed Mary hast bestowed upon mankind the rewards of eternal salvation: grant, we beseech thee; that we may perceive her intercession for us, through whom we have been counted worthy to receive the author of life, Jesus Christ thy Son, our Lord. Who liveth.

%CHECK:
\readingcitation{Epistle}{Titus 2:11}
\lett{B}{eloved:} The grace of God that bringeth salvation hath appeared to all men, Teaching us that, denying ungodliness and worldly lusts, we should live soberly, righteously, and godly, in this present world; Looking for that blessed hope, and the glorious appearing of the great God and our Saviour Jesus Christ; Who gave himself for us, that he might redeem us from all iniquity, and purify unto himself a peculiar people, zealous of good works. These things speak, and exhort, and rebuke with all authority. Let no man despise thee.

\readingcitation{Gospel}{Luke 2:15}
\lett{A}{t that time:} When eight days were accomplished for the circumcising of the child, his name was called \divineName{Jesus}, which was so named of the angel before he was conceived in the womb.


\subby{Epiphany of Our Lord}\label{EpiphanyMass}
\fancyhead[RE,LO]{Epiphany}
\begin{inhead}
{First Class Double, Second Class Octave}
\end{inhead}

\collect
\lett{O}{God,} who on this day by the leading of a star didst manifest thy only-begotten Son to the Gentiles: mercifully grant; that we, which know thee now by faith, may be brought even to the contemplation of the beauty of thy majesty. Through the same.
\begin{rubric}
    This Collect to be said daily throughout the Octave.
\end{rubric}

\readingcitation{Epistle}{Isaiah 60:1}
\lett{A}{rise,} shine; for thy light is come, and the glory of the \divineName{Lord} is risen upon thee. For, behold, the darkness shall cover the earth, and gross darkness the people: but the \divineName{Lord} shall arise upon thee, and his glory shall be seen upon thee. And the Gentiles shall come to thy light, and kings to the brightness of thy rising. Lift up thine eyes round about, and see: all they gather themselves together, they come to thee: thy sons shall come from far, and thy daughters shall be nursed at thy side. Then thou shalt see, and flow together, and thine heart shall fear, and be enlarged; because the abundance of the sea shall be converted unto thee, the forces of the Gentiles shall come unto thee. The multitude of camels shall cover thee, the dromedaries of Midian and Ephah; all they from Sheba shall come: they shall bring gold and incense; and they shall shew forth the praises of the \divineName{Lord}.


\subby{Sunday within the Epiphany Octave}
\fancyhead[RE,LO]{Epiphany I}
\begin{inhead}
{Semidouble}
\end{inhead}

%CHECK:
\readingcitation{Epistle}{Romans 12:1}
\lett{B}{rethren:} I beseech you therefore, brethren, by the mercies of God, that ye present your bodies a living sacrifice, holy, acceptable unto God, which is your reasonable service. And be not conformed to this world: but be ye transformed by the renewing of your mind, that ye may prove what is that good, and acceptable, and perfect, will of God. For I say, through the grace given unto me, to every man that is among you, not to think of himself more highly than he ought to think; but to think soberly, according as God hath dealt to every man the measure of faith. For as we have many members in one body, and all members have not the same office: So we, being many, are one body in Christ, and every one members one of another: in Christ Jesus our Lord.

\readingcitation{Gospel}{Luke 2:41}
\lett{W}{hen} Jesus was twelve years old, they went up to Jerusalem after the custom of the feast. And when they had fulfilled the days, as they returned, the child Jesus tarried behind in Jerusalem; and Joseph and his mother knew not of it. But they, supposing him to have been in the company, went a day's journey; and they sought him among their kinsfolk and acquaintance. And when they found him not, they turned back again to Jerusalem, seeking him. And it came to pass, that after three days they found him in the temple, sitting in the midst of the doctors, both hearing them, and asking them questions. And all that heard him were astonished at his understanding and answers. And when they saw him, they were amazed: and his mother said unto him, Son, why hast thou thus dealt with us? behold, thy father and I have sought thee sorrowing. And he said unto them, How is it that ye sought me? wist ye not that I must be about my Father's business? And they understood not the saying which he spake unto them. And he went down with them, and came to Nazareth, and was subject unto them: but his mother kept all these sayings in her heart. And Jesus increased in wisdom and stature, and in favour with God and man.


\subby{Second Sunday after Epiphany}\label{EpiphanyII}
\fancyhead[RE,LO]{Epiphany II}
\begin{inhead}
{Semidouble}
\end{inhead}

\readingcitation{Gospel}{John 2:1}
\lett{A}{t that time:} There was a marriage in Cana of Galilee; and the mother of Jesus was there: And both Jesus was called, and his disciples, to the marriage. And when they wanted wine, the mother of Jesus saith unto him, They have no wine. Jesus saith unto her, Woman, what have I to do with thee? mine hour is not yet come. His mother saith unto the servants, Whatsoever he saith unto you, do it. And there were set there six waterpots of stone, after the manner of the purifying of the Jews, containing two or three firkins apiece. Jesus saith unto them, Fill the waterpots with water. And they filled them up to the brim. And he saith unto them, Draw out now, and bear unto the governor of the feast. And they bare it. When the ruler of the feast had tasted the water that was made wine, and knew not whence it was: (but the servants which drew the water knew;) the governor of the feast called the bridegroom, And saith unto him, Every man at the beginning doth set forth good wine; and when men have well drunk, then that which is worse: but thou hast kept the good wine until now. This beginning of miracles did Jesus in Cana of Galilee, and manifested forth his glory; and his disciples believed on him.


\subby{Fourth Sunday after Epiphany}
\fancyhead[RE,LO]{Epiphany IV}
\begin{inhead}
{Semidouble}
\end{inhead}

%CHECK:
\collect
%Is the English collect sufficiently similar?
\lett{O}{God,} who knowest us to be set in the midst of so many and great dangers, that by reason of the frailty of our nature we cannot always stand upright; Grant to us such strength and protection, as may support us in all dangers, and carry us through all temptations. Through.
\begin{rubric}
    If it be after 2 February, \nth{2} Collect is \emph{Of the Saints} (p. \pageref{SPSaints}) \& \nth{3} is \emph{ad libitum}.\par
    Otherwise, additional Collects are as in the Second Sunday after Epiphany (p. \pageref{EpiphanyII}).
\end{rubric}

\readingcitation{Epistle}{Romans 13:8}
\lett{B}{rethren:} Owe no man any thing, but to love one another: for he that loveth another hath fulfilled the law. For this, Thou shalt not commit adultery, Thou shalt not kill, Thou shalt not steal, Thou shalt not bear false witness, Thou shalt not covet; and if there be any other commandment, it is briefly comprehended in this saying, namely, Thou shalt love thy neighbour as thyself Love worketh no ill to his neighbour: therefore love is the fulfilling of the law.

\readingcitation{Gospel}{Matthew 8:23}
\lett{A}{t that time:} When Jesus was entered into a ship, his disciples followed him. And, behold, there arose a great tempest in the sea, insomuch that the ship was covered with the waves: but he was asleep. And his disciples came to him, and awoke him, saying, Lord, save us: we perish. And he saith unto them, Why are ye fearful, O ye of little faith? Then he arose, and rebuked the winds and the sea; and there was a great calm. But the men marvelled, saying, What manner of man is this, that even the winds and the sea obey him!


\subby{Fifth Sunday after Epiphany}\label{EpiphanyV}
\fancyhead[RE,LO]{Epiphany V}
\begin{inhead}
{Semidouble}
\end{inhead}
%CHECK:
\collect
\lett{O}{Lord,} we beseech thee to keep thy Church and household continually in thy true religion; that they who do lean only upon the hope of thy heavenly grace may evermore be defended by thy mighty power. Through.
\begin{rubric}
    \nth{2} Collect is \emph{Of the Saints} (p. \pageref{SPSaints}) \& \nth{3} is \emph{ad libitum}.
\end{rubric}
%CHECK:
\readingcitation{Epistle}{Colossians 3:12}
\lett{B}{rethren:} Put on therefore, as the elect of God, holy and beloved, bowels of mercies, kindness, humbleness of mind, meekness, longsuffering; Forbearing one another, and forgiving one another, if any man have a quarrel against any: even as Christ forgave you, so also do ye. And above all these things put on charity, which is the bond of perfectness. And let the peace of God rule in your hearts, to the which also ye are called in one body; and be ye thankful. Let the word of Christ dwell in you richly in all wisdom; teaching and admonishing one another in psalms and hymns and spiritual songs, singing with grace in your hearts to the Lord. And whatsoever ye do in word or deed, do all in the name of the Lord Jesus, giving thanks to God and the Father by %
Our Lord Jesus Christ.


\subby{Sixth Sunday after Epiphany}
\fancyhead[RE,LO]{Epiphany VI}
\begin{inhead}
{Semidouble}
\end{inhead}

\collect
\lett{G}{rant,} we beseech thee, almighty God: that we, ever meditating those things that be reasonable and pleasing unto thee, may both in word and deed fulfil the same. Through.
\begin{rubric}
    Additional Collects are as in the Fifth Sunday after Epiphany (p. \pageref{EpiphanyV}).
\end{rubric}

%Double-check introduction.
\readingcitation{Epistle}{1 Thessalonians 1:2}
\lett{B}{rethren:} We give thanks to God always for you all, making mention of you in our prayers; Remembering without ceasing your work of faith, and labour of love, and patience of hope in our Lord Jesus Christ, in the sight of God and our Father; Knowing, brethren beloved, your election of God. For our gospel came not unto you in word only, but also in power, and in the Holy Ghost, and in much assurance; as ye know what manner of men we were among you for your sake. And ye became followers of us, and of the Lord, having received the word in much affliction, with joy of the Holy Ghost: So that ye were ensamples to all that believe in Macedonia and Achaia. For from you sounded out the word of the Lord not only in Macedonia and Achaia, but also in every place your faith to God-ward is spread abroad; so that we need not to speak any thing. For they themselves shew of us what manner of entering in we had unto you, and how ye turned to God from idols to serve the living and true God; And to wait for his Son from heaven, whom he raised from the dead, even Jesus, which delivered us from the wrath to come.

\readingcitation{Gospel}{Matthew 13:31}
\lett{A}{t that time:} Jesus spake this parable unto the crowds: The kingdom of heaven is like to a grain of mustard seed, which a man took, and sowed in his field: Which indeed is the least of all seeds: but when it is grown, it is the greatest among herbs, and becometh a tree, so that the birds of the air come and lodge in the branches thereof. Another parable spake he unto them; The kingdom of heaven is like unto leaven, which a woman took, and hid in three measures of meal, till the whole was leavened. All these things spake Jesus unto the multitude in parables; and without a parable spake he not unto them: That it might be fulfilled which was spoken by the prophet, saying, I will open my mouth in parables; I will utter things which have been kept secret from the foundation of the world.


\subby{Septuagesima Sunday}
\fancyhead[RE,LO]{Septuagesima}
\begin{inhead}
{Second Class Semidouble}
\end{inhead}

%CHECK:
\collect
\lett{O}{Lord,} we beseech thee favourably to hear the prayers of thy people; that we, who are justly punished for our offences, may be mercifully delivered by thy goodness, for the glory of thy Name; through Jesus Christ our Saviour, who liveth and reigneth with thee and the Holy Ghost ever, one God, world without end. \textit{Amen.}
\begin{rubric}
    Additional Collects are as in the Fifth Sunday after Epiphany (p. \pageref{EpiphanyV}).
\end{rubric}

\readingcitation{Epistle}{1 Corinthians 9:24}
\lett{B}{rethren:} Know ye not that they which run in a race run all, but one receiveth the prize? So run, that ye may obtain. And every man that striveth for the mastery is temperate in all things. Now they do it to obtain a corruptible crown; but we an incorruptible. I therefore so run, not as uncertainly; so fight I, not as one that beateth the air: But I keep under my body, and bring it into subjection: lest that by any means, when I have preached to others, I myself should be a castaway. Moreover, brethren, I would not that ye should be ignorant, how that all our fathers were under the cloud, and all passed through the sea; And were all baptized unto Moses in the cloud and in the sea; And did all eat the same spiritual meat; And did all drink the same spiritual drink: for they drank of that spiritual Rock that followed them: and that Rock was Christ. But with many of them God was not well pleased.


\subby{Sexagesima Sunday}
\fancyhead[RE,LO]{Sexagesima}
\begin{inhead}
{Second Class Semidouble}
\end{inhead}

%Collect used in BCP:
%\collect
%\lett{O}{God,} who seest that we put not our trust in any thing that we do: mercifully grant; that by the protection of the Doctor of the Gentiles we may be defended against all adversity. Through.
\begin{rubric}
    Additional Collects are as in the Fifth Sunday after Epiphany (p. \pageref{EpiphanyV}).
\end{rubric}

\readingcitation{Epistle}{2 Corinthians 11:19}
\lett{B}{rethren:} Ye suffer fools gladly, seeing ye yourselves are wise. For ye suffer, if a man bring you into bondage, if a man devour you, if a man take of you, if a man exalt himself, if a man smite you on the face. I speak as concerning reproach, as though we had been weak. Howbeit whereinsoever any is bold, (I speak foolishly,) I am bold also. Are they Hebrews? so am I. Are they Israelites? so am I. Are they the seed of Abraham? so am I. Are they ministers of Christ? (I speak as a fool) I am more; in labours more abundant, in stripes above measure, in prisons more frequent, in deaths oft. Of the Jews five times received I forty stripes save one. Thrice was I beaten with rods, once was I stoned, thrice I suffered shipwreck, a night and a day I have been in the deep; In journeyings often, in perils of waters, in perils of robbers, in perils by mine own countrymen, in perils by the heathen, in perils in the city, in perils in the wilderness, in perils in the sea, in perils among false brethren; In weariness and painfulness, in watchings often, in hunger and thirst, in fastings often, in cold and nakedness. Beside those things that are without, that which cometh upon me daily, the care of all the churches. Who is weak, and I am not weak? who is offended, and I burn not? If I must needs glory, I will glory of the things which concern mine infirmities. The God and Father of our Lord Jesus Christ, which is blessed for evermore, knoweth that I lie not. In Damascus the governor under Aretas the king kept the city of the Damascenes with a garrison, desirous to apprehend me: And through a window in a basket was I let down by the wall, and escaped his hands. It is not expedient for me doubtless to glory. I will come to visions and revelations of the Lord. I knew a man in Christ above fourteen years ago, (whether in the body, I cannot tell; or whether out of the body, I cannot tell: God knoweth;) such an one caught up to the third heaven. And I knew such a man, (whether in the body, or out of the body, I cannot tell: God knoweth;) How that he was caught up into paradise, and heard unspeakable words, which it is not lawful for a man to utter. Of such an one will I glory: yet of myself I will not glory, but in mine infirmities. For though I would desire to glory, I shall not be a fool; for I will say the truth: but now I forbear, lest any man should think of me above that which he seeth me to be, or that he heareth of me. And lest I should be exalted above measure through the abundance of the revelations, there was given to me a thorn in the flesh, the messenger of Satan to buffet me, lest I should be exalted above measure. For this thing I besought the Lord thrice, that it might depart from me. And he said unto me, My grace is sufficient for thee: for my strength is made perfect in weakness. Most gladly therefore will I rather glory in my infirmities, that the power of Christ may rest upon me.


\subby{Quinquagesima Sunday}
\fancyhead[RE,LO]{Quinquagesima}
\begin{inhead}
{Second Class Semidouble}
\end{inhead}

\collect
\lett{W}{e} beseech thee, O Lord, graciously hear our prayers: that we, being loosed from the bonds of our sins, may be defended against all adversity. Through.
\begin{rubric}
    Additional Collects are as in the Fifth Sunday after Epiphany (p. \pageref{EpiphanyV}).
\end{rubric}


\subby{Ash Wednesday}
\fancyhead[RE,LO]{Ash Wednesday}
\begin{inhead}
{First Class Feria}
\end{inhead}

\collect
\lett{G}{rant,} O Lord, to thy faithful people: that they may both enter upon the venerable solemnity of these fasts with fitting devotion, and perform the same in godly quietness. Through.
\begin{rubric}
    \nth{2} Collect is of the Saints (p. \pageref{SPSaints}) \& \nth{3} for the Living and Departed (p. \pageref{SPLivingDeparted}).
\end{rubric}

\readingcitation{Epistle}{Joel 2:12}
\lett{T}{hus} saith the Lord: Turn ye even to me with all your heart, and with fasting, and with weeping, and with mourning: And rend your heart, and not your garments, and turn unto the \divineName{Lord} your God: for he is gracious and merciful, slow to anger, and of great kindness, and repenteth him of the evil. Who knoweth if he will return and repent, and leave a blessing behind him; even a meat offering and a drink offering unto the \divineName{Lord} your God? Blow the trumpet in Zion, sanctify a fast, call a solemn assembly: Gather the people, sanctify the congregation, assemble the elders, gather the children, and those that suck the breasts: let the bridegroom go forth of his chamber, and the bride out of her closet. Let the priests, the ministers of the \divineName{Lord}, weep between the porch and the altar, and let them say, Spare thy people, O \divineName{Lord}, and give not thine heritage to reproach, that the heathen should rule over them: wherefore should they say among the people, Where is their God? Then will the \divineName{Lord} be jealous for his land, and pity his people. Yea, the \divineName{Lord} will answer and say unto his people, Behold, I will send you corn, and wine, and oil, and ye shall be satisfied therewith: and I will no more make you a reproach among the heathen.


\subby{First Sunday of Lent}
\fancyhead[RE,LO]{Lent I}
\begin{inhead}
{First Class Semidouble}
\end{inhead}

\collect
\lett{O}{God,} who year by year dost purify the Church by this Lenten observance: grant unto thy family; that those things which they strive to obtain from thee by abstinence, they may by good works fulfil. Through.
\begin{rubric}
    Additional Collects are as on Ash Wednesday (p. \pageref{AshWednesdayMass}).
\end{rubric}


\subby{Second Sunday of Lent}
\fancyhead[RE,LO]{Lent II}
\begin{inhead}
{First Class Semidouble}
\end{inhead}

%CHECK:
\readingcitation{Epistle}{1 Thessalonians 4:1}
\lett{B}{rethren:} we beseech you, and exhort you by the Lord Jesus, that as ye have received of us how ye ought to walk and to please God, so ye would abound more and more. For ye know what commandments we gave you by the Lord Jesus. For this is the will of God, even your sanctification, that ye should abstain from fornication: That every one of you should know how to possess his vessel in sanctification and honour; Not in the lust of concupiscence, even as the Gentiles which know not God That no man go beyond and defraud his brother in any matter: because that the Lord is the avenger of all such, as we also have forewarned you and testified. For God hath not called us unto uncleanness, but unto holiness. He therefore that despiseth, despiseth not man, but God, who hath also given unto us his holy Spirit.

\readingcitation{Gospel}{Matthew 17:1}
\lett{A}{t that time:} Jesus taketh Peter, James, and John his brother, and bringeth them up into an high mountain apart, And was transfigured before them: and his face did shine as the sun, and his raiment was white as the light. And, behold, there appeared unto them Moses and Elias talking with him. Then answered Peter, and said unto Jesus, Lord, it is good for us to be here: if thou wilt, let us make here three tabernacles; one for thee, and one for Moses, and one for Elias. While he yet spake, behold, a bright cloud overshadowed them: and behold a voice out of the cloud, which said, This is my beloved Son, in whom I am well pleased; hear ye him. And when the disciples heard it, they fell on their face, and were sore afraid. And Jesus came and touched them, and said, Arise, and be not afraid. And when they had lifted up their eyes, they saw no man, save Jesus only. And as they came down from the mountain, Jesus charged them, saying, Tell the vision to no man, until the Son of man be risen again from the dead.

\subby{Third Sunday of Lent}
\fancyhead[RE,LO]{Lent III}
\begin{inhead}
{First Class Semidouble}
\end{inhead}

%CHECK:
\collect
\lett{W}{e} beseech thee, Almighty God, look upon the hearty desires of thy humble servants, and stretch forth the right hand of thy Majesty, to be our defence. Through.
\begin{rubric}
    Additional Collects are as on Ash Wednesday (p. \pageref{AshWednesdayMass}).
\end{rubric}

\readingcitation{Epistle}{Ephesians 5:1}
\lett{B}{rethren:} Be ye therefore followers of God, as dear children; And walk in love, as Christ also hath loved us, and hath given himself for us an offering and a sacrifice to God for a sweetsmelling savour. But fornication, and all uncleanness, or covetousness, let it not be once named among you, as becometh saints; Neither filthiness, nor foolish talking, nor jesting, which are not convenient: but rather giving of thanks. For this ye know, that no whoremonger, nor unclean person, nor covetous man, who is an idolater, hath any inheritance in the kingdom of Christ and of God. Let no man deceive you with vain words: for because of these things cometh the wrath of God upon the children of disobedience. Be not ye therefore partakers with them. For ye were sometimes darkness, but now are ye light in the Lord: walk as children of light: (For the fruit of the Spirit is in all goodness and righteousness and truth;).


\subby{Fourth Sunday of Lent}
\fancyhead[RE,LO]{Lent IV}
\begin{inhead}
{First Class Semidouble}
\end{inhead}

%CHECK:
\readingcitation{Epistle}{Galatians 4:21}
%Anglican Addition:
\lett{B}{rethren:} tell me, ye that desire to be under the law, do ye not hear the law? For it is written, that Abraham had two sons, the one by a bondmaid, the other by a freewoman. But he who was of the bondwoman was born after the flesh; but he of the freewoman was by promise. Which things are an allegory: for these are the two covenants; the one from the mount Sinai, which gendereth to bondage, which is Agar. For this Agar is mount Sinai in Arabia, and answereth to Jerusalem which now is, and is in bondage with her children. But Jerusalem which is above is free, which is the mother of us all. For it is written, Rejoice, thou barren that bearest not; break forth and cry, thou that travailest not: for the desolate hath many more children than she which hath an husband. Now we, brethren, as Isaac was, are the children of promise. But as then he that was born after the flesh persecuted him that was born after the Spirit, even so it is now. Nevertheless what saith the scripture? Cast out the bondwoman and her son: for the son of the bondwoman shall not be heir with the son of the freewoman. So then, brethren, we are not children of the bondwoman, but of the free.

\readingcitation{Gospel}{John 6:1}
\lett{A}{t that time:} Jesus went over the sea of Galilee, which is the sea of Tiberias. And a great multitude followed him, because they saw his miracles which he did on them that were diseased. And Jesus went up into a mountain, and there he sat with his disciples. And the passover, a feast of the Jews, was nigh. When Jesus then lifted up his eyes, and saw a great company come unto him, he saith unto Philip, Whence shall we buy bread, that these may eat? And this he said to prove him: for he himself knew what he would do. Philip answered him, Two hundred pennyworth of bread is not sufficient for them, that every one of them may take a little. One of his disciples, Andrew, Simon Peter's brother, saith unto him, There is a lad here, which hath five barley loaves, and two small fishes: but what are they among so many? And Jesus said, Make the men sit down. Now there was much grass in the place. So the men sat down, in number about five thousand. And Jesus took the loaves; and when he had given thanks, he distributed to the disciples, and the disciples to them that were set down; and likewise of the fishes as much as they would. When they were filled, he said unto his disciples, Gather up the fragments that remain, that nothing be lost. Therefore they gathered them together, and filled twelve baskets with the fragments of the five barley loaves, which remained over and above unto them that had eaten. Then those men, when they had seen the miracle that Jesus did, said, This is of a truth that prophet that should come into the world. When Jesus therefore perceived that they would come and take him by force, to make him a king, he departed again into a mountain himself alone.

\subby{Passion Sunday}\label{PassionSunday}
\fancyhead[RE,LO]{Passion Sunday}
\begin{inhead}
{First Class Semidouble}
\end{inhead}

%CHECK:
\readingcitation{Epistle}{Hebrews 9:11}
\lett{B}{rethren:} Christ being come an high priest of good things to come, by a greater and more perfect tabernacle, not made with hands, that is to say, not of this building; Neither by the blood of goats and calves, but by his own blood he entered in once into the holy place, having obtained eternal redemption for us. For if the blood of bulls and of goats, and the ashes of an heifer sprinkling the unclean, sanctifieth to the purifying of the flesh: How much more shall the blood of Christ, who through the eternal Spirit offered himself without spot to God, purge your conscience from dead works to serve the living God? And for this cause he is the mediator of the new testament, that by means of death, for the redemption of the transgressions that were under the first testament, they which are called might receive the promise of eternal inheritance, in Christ Jesus our Lord.


\subby{Palm Sunday}\label{PalmSunday}
\fancyhead[RE,LO]{Palm Sunday}
\begin{inhead}
{First Class Semidouble}
\end{inhead}

\collect
\lett{A}{lmighty} and everlasting God, who didst cause our Saviour to take upon him our flesh, and to suffer on the cross, that mankind might follow the example of his humility: mercifully grant; that we may both keep in mind the lessons of his patience, and be made partakers of his resurrection. Through the same.
\begin{rubric}
    This Collect only is said, even if a Commemoration has been made in the Office.
\end{rubric}
\begin{rubric}
    This Collect is to be said every day, after the Collect appointed for the day, until Good Friday.
\end{rubric}


\subby{Easter Day}
\fancyhead[RE,LO]{Easter Day}
\begin{inhead}
    {First Class Double, First Class Octave}
\end{inhead}

\collect
\lett{O}{God,} who through thine only-begotten Son hast this day overcome death, and opened unto us the gate of eternity: prosper by thine assistance our desires, which by thy preventing grace thou dost put into our minds. Through the same.
\begin{rubric}
    This Collect is to be said daily throughout the Easter Octave.
\end{rubric}

\readingcitation{Epistle}{1 Corinthians 5:6}
\lett{B}{rethren:} Purge out therefore the old leaven, that ye may be a new lump, as ye are unleavened. For even Christ our passover is sacrificed for us: Therefore let us keep the feast, not with old leaven, neither with the leaven of malice and wickedness; but with the unleavened bread of sincerity and truth.


\readingcitation{Gospel}{Mark 16:1}
\lett{A}{t that time:} Mary Magdalene, and Mary the mother of James, and Salome, bought sweet spices, that they might come and anoint him. And very early in the morning the first day of the week, they came unto the sepulchre at the rising of the sun. And they said among themselves, Who shall roll us away the stone from the door of the sepulchre? And when they looked, they saw that the stone was rolled away: for it was very great. And entering into the sepulchre, they saw a young man sitting on the right side, clothed in a long white garment; and they were affrighted. And he saith unto them, Be not affrighted: Ye seek Jesus of Nazareth, which was crucified: he is risen; he is not here: behold the place where they laid him. But go your way, tell his disciples and Peter that he goeth before you into Galilee: there shall ye see him, as he said unto you. And they went out quickly, and fled from the sepulchre; for they trembled and were amazed: neither said they any thing to any man; for they were afraid.


\subby{Easter Monday}
\fancyhead[RE,LO]{Easter Monday}
\begin{inhead}
    {First Class Double}
\end{inhead}

\collect
\lett{O}{God,} who by the paschal solemnity hast bestowed healing upon the world: prosper, we beseech thee, thy people with thy heavenly gift; that they may be worthy to attain unto perfect freedom, and may likewise be profited unto everlasting life. Through.

\readingcitation{Epistle}{Acts 10:37}
\lett{I}{n those days:} Peter standing in the midst of the people, said: Men and brethren, ye know the word which was published throughout all Jud{\ae}a: and began from Galilee, after the baptism which John preached: how God anointed Jesus of Nazareth with the Holy Ghost, and with power, who went about doing good, and healing all that were oppressed of the devil, for God was with him. And we are witnesses of all things, which he did both in the land of the Jews and in Jerusalem, whom they slew and hanged on a tree. Him God raised up the third day, and shewed him openly, not to all the people, but unto witnesses chosen before of God: even to us, who did eat and drink with him, after he rose from the dead. And he commanded us to preach unto the people, and to testify, that it is he who was ordained of God to be the judge of quick and dead. To him give all the Prophets witness, that through his name whosoever believeth in him shall receive remission of sins.

\subby{Easter Tuesday}
\fancyhead[RE,LO]{Easter Tuesday}
\begin{inhead}
    {First Class Double}
\end{inhead}

\collect
\lett{O}{God,} who dost ever multiply thy Church with new offspring; grant unto thy servants; that they may hold fast in their lives the sacrament which by faith they have received. Through.

\readingcitation{Epistle}{Acts 13:26}
\lett{I}{n those days:} Paul standing up, and beckoning with his hand to keep silence, said: Men and brethren, children of the stock of Abraham, and whosoever among you feareth God, to you is the word of this salvation sent. For they that dwell at Jerusalem, and their rulers, because they knew not Jesus, nor yet the voices of the Prophets which are read every sabbath-day, they have fulfilled them in condemning him: and though they found no cause of death in him, yet desired they Pilate that he should be slain. And when they had fulfilled all that was written of him, they took him down from the tree, and laid him in a sepulchre But God raised him from the dead: and he was seen many days of them which came up with him from Galilee to Jerusalem, who are his witnesses unto the people. And we declare unto you glad tidings, how that the promise which was made unto the fathers, God, hath fulfilled the same unto us their children, in that he hath raised up Jesus Christ, our Lord.

\subby{Low Sunday}
\fancyhead[RE,LO]{Easter I}
\begin{inhead}
    {First Class Double, Octave Day of Easter}
\end{inhead}


\collect
\lett{G}{rant,} we beseech thee, almighty God: that we who have celebrated the paschal feasts may, through thy bounty, ever keep them in our life and conversation. Through.

\readingcitation{Epistle}{1 John 1:4}
\lett{D}{early beloved:} Whatsoever is born of God overcometh the world: and this is the victory that overcometh the world, even our faith. Who is he that overcometh the world, but he that believeth that Jesus is the Son of God? This is he that came by water and blood, even Jesus Christ; not by water only, but by water and blood. And it is the Spirit that beareth witness, because the Spirit is truth. For there are three that bear record in heaven, the Father, the Word, and the Holy Ghost: and these three are one. And there are three that bear witness in earth, the spirit, and the water, and the blood: and these three agree in one. If we receive the witness of men, the witness of God is greater: for this is the witness of God which he hath testified of his Son. He that believeth on the Son of God hath the witness in himself.

\readingcitation{Gospel}{John 20:19}
\lett{A}{t that time:} The same day at evening, being the first day of the week, when the doors were shut where the disciples were assembled for fear of the Jews, came Jesus and stood in the midst, and saith unto them, Peace be unto you. And when he had so said, he shewed unto them his hands and his side. Then were the disciples glad, when they saw the Lord. Then said Jesus to them again, Peace be unto you: as my Father hath sent me, even so send I you. And when he had said this, he breathed on them, and saith unto them, Receive ye the Holy Ghost: Whose soever sins ye remit, they are remitted unto them; and whose soever sins ye retain, they are retained. But Thomas, one of the twelve, called Didymus, was not with them when Jesus came. The other disciples therefore said unto him, We have seen the Lord. But he said unto them, Except I shall see in his hands the print of the nails, and put my finger into the print of the nails, and thrust my hand into his side, I will not believe. And after eight days again his disciples were within, and Thomas with them: then came Jesus, the doors being shut, and stood in the midst, and said, Peace be unto you. Then saith he to Thomas, Reach hither thy finger, and behold my hands; and reach hither thy hand, and thrust it into my side: and be not faithless, but believing. And Thomas answered and said unto him, My Lord and my God. Jesus saith unto him, Thomas, because thou hast seen me, thou hast believed: blessed are they that have not seen, and yet have believed. And many other signs truly did Jesus in the presence of his disciples, which are not written in this book: But these are written, that ye might believe that Jesus is the Christ, the Son of God; and that believing ye might have life through his name.


\subby{Second Sunday after Easter}\label{EasterII}
\fancyhead[RE,LO]{Easter II}
\begin{inhead}
    {Semidouble}
\end{inhead}

\collect
\lett{O}{God,} who in the humility of thy Son hast raised up a fallen world: grant unto thy faithful abiding gladness: that they whom thou hast delivered from the peril of eternal death, may be brought to the fruition of everlasting joys. Through the same.
\begin{rubric}
    \nth{2} Collect is of St. Mary in Eastertide (p. \pageref{SPMaryInEaster}) \& \nth{3} Against the Persecutors of the Church (p. \pageref{SPAgainst}) or for the Chief Bishop (p. \pageref{SPChiefBishop}).
\end{rubric}


\readingcitation{Epistle}{1 Peter 2:19}
\lett{D}{early beloved:} Christ suffered for us, leaving us an example, that ye should follow his steps: Who did no sin, neither was guile found in his mouth: Who, when he was reviled, reviled not again; when he suffered, he threatened not; but committed himself to him that judgeth righteously: Who his own self bare our sins in his own body on the tree, that we, being dead to sins, should live unto righteousness: by whose stripes ye were healed. For ye were as sheep going astray; but are now returned unto the Shepherd and Bishop of your souls.


\subby{Third Sunday after Easter}
\fancyhead[RE,LO]{Easter III}
\begin{inhead}
    {Semidouble}
\end{inhead}

\readingcitation{Epistle}{1 Peter 2:11}
\lett{D}{early beloved:} I beseech you as strangers and pilgrims, abstain from fleshly lusts, which war against the soul; Having your conversation honest among the Gentiles: that, whereas they speak against you as evildoers, they may by your good works, which they shall behold, glorify God in the day of visitation. Submit yourselves to every ordinance of man for the Lord's sake: whether it be to the king, as supreme; Or unto governors, as unto them that are sent by him for the punishment of evildoers, and for the praise of them that do well. For so is the will of God, that with well doing ye may put to silence the ignorance of foolish men: As free, and not using your liberty for a cloke of maliciousness, but as the servants of God. Honour all men. Love the brotherhood. Fear God. Honour the king. Servants, be subject to your masters with all fear; not only to the good and gentle, but also to the froward. For this is thankworthy: in Christ Jesus, our Lord.

\subby{Fourth Sunday after Easter}
\fancyhead[RE,LO]{Easter IV}
\begin{inhead}
    {Semidouble}
\end{inhead}

\collect
\lett{O}{God,} who makest the minds of the faithful to be of one will: grant unto thy people, that they may love the thing which thou commandest, and desire that which thou dost promise; that amid the changes of the world our hearts may there be fixed, where true joys are to be found. Through.
\begin{rubric}
    Additional Collects are as on Easter II (p. \pageref{EasterII}).
\end{rubric}

\subby{Rogation Sunday}
\fancyhead[RE,LO]{Easter V}
\begin{inhead}
    {Semidouble}
\end{inhead}

\collect
\lett{O}{God,} from whom all good things do come, grant to us thy humble servants: that by thy inspiration we may think those things that be right; and by thy guiding may perform the same. Through.

\readingcitation{Gospel}{John 16:23}
\lett{A}{t that time:} Jesus said unto his disciples: Verily, verily, I say unto you, Whatsoever ye shall ask the Father in my name, he will give it you. Hitherto have ye asked nothing in my name: ask, and ye shall receive, that your joy may be full. These things have I spoken unto you in proverbs: but the time cometh, when I shall no more speak unto you in proverbs, but I shall shew you plainly of the Father. At that day ye shall ask in my name: and I say not unto you, that I will pray the Father for you: For the Father himself loveth you, because ye have loved me, and have believed that I came out from God. I came forth from the Father, and am come into the world: again, I leave the world, and go to the Father. His disciples said unto him, Lo, now speakest thou plainly, and speakest no proverb. Now are we sure that thou knowest all things, and needest not that any man should ask thee: by this we believe that thou camest forth from God.


\subby{Rogation Days}
\fancyhead[RE,LO]{Rogation Days}
\subsubsec{Daily Office Propers}\par\noindent
\begin{inhead}
    Monday\par
 \nth{2} Feria
\end{inhead}

\begin{inhead}
    Wednesday
\end{inhead}
%CHECK:

\collect
\lett{G}{rant,} we beseech thee, almighty God: that we, who in our affliction do put our trust in thy mercy: may ever be defended by thy protection against all adversities. Through.
\begin{rubric}
    Additional Collects are as on Easter II (p. \pageref{EasterII}).
\end{rubric}

\readingcitation{Epistle}{James 5:16}
\lett{D}{early beloved:} Confess your faults one to another, and pray one for another, that ye may be healed. The effectual fervent prayer of a righteous man availeth much. Elias was a man subject to like passions as we are, and he prayed earnestly that it might not rain: and it rained not on the earth by the space of three years and six months. And he prayed again, and the heaven gave rain, and the earth brought forth her fruit. Brethren, if any of you do err from the truth, and one convert him; Let him know, that he which converteth the sinner from the error of his way shall save a soul from death, and shall hide a multitude of sins.

\subby{Ascension Thursday}
\fancyhead[RE,LO]{Ascension Thursday}
\begin{inhead}
    {First Class Double, Third Class Octave}
\end{inhead}

\collect
\lett{G}{rant,} we beseech thee almighty God: that as we believe thy only-begotten Son our Redeemer to have this day ascended into the heavens: so we may also in mind dwell in heavenly places. Through the same.

%CHECK:
\readingcitation{Gospel}{Mark 16:14}
\lett{A}{t that time:} Jesus appeared unto the eleven as they sat at meat, and upbraided them with their unbelief and hardness of heart, because they believed not them which had seen him after he was risen. And he said unto them, Go ye into all the world, and preach the gospel to every creature. He that believeth and is baptized shall be saved; but he that believeth not shall be damned. And these signs shall follow them that believe; In my name shall they cast out devils; they shall speak with new tongues; They shall take up serpents; and if they drink any deadly thing, it shall not hurt them; they shall lay hands on the sick, and they shall recover. So then after the Lord had spoken unto them, he was received up into heaven, and sat on the right hand of God. And they went forth, and preached every where, the Lord working with them, and confirming the word with signs following.

\subby{Sunday within the Ascension Octave}
\fancyhead[RE,LO]{Sunday within Ascension}
\begin{inhead}
    {Semidouble}
\end{inhead}

\collect
\lett{A}{lmighty} and everlasting God: grant that our wills may ever be subjected unto thee; and that we may serve thy majesty in sincerity of heart. Through.
\lett{G}{rant,} we beseech thee, almighty God: that like as we do believe thy only-begotten Son our Lord Jesus Christ to have ascended into the heavens; so we may also in heart and mind thither ascend and with him continually dwell. Who liveth and reigneth with thee.

%CHECK:
\readingcitation{Epistle}{1 Peter 4:7}

\lett{D}{early beloved:} The end of all things is at hand:
%
be ye therefore sober, and watch unto prayer. And above all things have fervent charity among yourselves: for charity shall cover the multitude of sins. Use hospitality one to another without grudging. As every man hath received the gift, even so minister the same one to another, as good stewards of the manifold grace of God. If any man speak, let him speak as the oracles of God; if any man minister, let him do it as of the ability which God giveth: that God in all things may be glorified through Jesus Christ,
%
to whom be praise and dominion for ever and ever. Amen.

\subby{Whitsunday}
\fancyhead[RE,LO]{Whitsunday}
\begin{inhead}
    {First Class Double, First Class Octave}
\end{inhead}

%CHECK:
\readingcitation{Gospel}{John 14:15}
\lett{A}{t that time:} Jesus said unto his disciples: If ye love me, keep my commandments. And I will pray the Father, and he shall give you another Comforter, that he may abide with you for ever; Even the Spirit of truth; whom the world cannot receive, because it seeth him not, neither knoweth him: but ye know him; for he dwelleth with you, and shall be in you. I will not leave you comfortless: I will come to you. Yet a little while, and the world seeth me no more; but ye see me: because I live, ye shall live also. At that day ye shall know that I am in my Father, and ye in me, and I in you. He that hath my commandments, and keepeth them, he it is that loveth me: and he that loveth me shall be loved of my Father, and I will love him, and will manifest myself to him. Judas saith unto him, not Iscariot, Lord, how is it that thou wilt manifest thyself unto us, and not unto the world? Jesus answered and said unto him,
%
If a man love me, he will keep my words: and my Father will love him, and we will come unto him, and make our abode with him. He that loveth me not keepeth not my sayings: and the word which ye hear is not mine, but the Father's which sent me. These things have I spoken unto you, being yet present with you. But the Comforter, which is the Holy Ghost, whom the Father will send in my name, he shall teach you all things, and bring all things to your remembrance, whatsoever I have said unto you. Peace I leave with you, my peace I give unto you: not as the world giveth, give I unto you. Let not your heart be troubled, neither let it be afraid. Ye have heard how I said unto you, I go away, and come again unto you. If ye loved me, ye would rejoice, because I said, I go unto the Father: for my Father is greater than I. And now I have told you before it come to pass, that, when it is come to pass, ye might believe. Hereafter I will not talk much with you: for the prince of this world cometh, and hath nothing in me. But that the world may know that I love the Father; and as the Father gave me commandment, even so I do.


\subby{Whit-Monday}
\fancyhead[RE,LO]{Whit-Monday}
\begin{inhead}
    {First Class Double}
\end{inhead}

\collect
\lett{O}{God,} who didst give the Holy Spirit unto thine Apostles: grant to thy people the fulfilment of their godly petitions; and as thou hast given to them faith, so also bestow upon them peace. Through.

%CHECK:
\readingcitation{Epistle}{Acts 10:34}
%The English and Roman epistles are so similar, we make no distinction.
\lett{I}{n those days:} Peter opened his mouth, and said, Of a truth I perceive that God is no respecter of persons: But in every nation he that feareth him, and worketh righteousness, is accepted with him. The word which God sent unto the children of Israel, preaching peace by Jesus Christ: (he is Lord of all:) That word, I say, ye know, which was published throughout all Judaea, and began from Galilee, after the baptism which John preached; How God anointed Jesus of Nazareth with the Holy Ghost and with power: who went about doing good, and healing all that were oppressed of the devil; for God was with him. And we are witnesses of all things which he did both in the land of the Jews, and in Jerusalem; whom they slew and hanged on a tree: Him God raised up the third day, and shewed him openly; Not to all the people, but unto witnesses chosen before of God, even to us, who did eat and drink with him after he rose from the dead.
%
And he commanded us to preach unto the people, and to testify that it is he which was ordained of God to be the Judge of quick and dead. To him give all the prophets witness, that through his name whosoever believeth in him shall receive remission of sins. While Peter yet spake these words, the Holy Ghost fell on all them which heard the word. And they of the circumcision which believed were astonished, as many as came with Peter, because that on the Gentiles also was poured out the gift of the Holy Ghost. For they heard them speak with tongues, and magnify God. Then answered Peter, Can any man forbid water, that these should not be baptized, which have received the Holy Ghost as well as we? And he commanded them to be baptized in the name of the Lord. Then prayed they him to tarry certain days.

\subby{Whit-Tuesday}
\fancyhead[RE,LO]{Whit-Tuesday}
\begin{inhead}
    {First Class Double}
\end{inhead}


\collect
\lett{W}{e} beseech thee, O Lord, let the power of the Holy Ghost be with us: that it may both mercifully cleanse our hearts, and defend us from all adversities. Through.


\subby{Trinity Sunday}
\fancyhead[RE,LO]{Pentecost Sunday}
\begin{inhead}
    {First Class Double}
\end{inhead}
\readingcitation{Epistle}{Romans 11:33}
\lett{O}{the} depth of the riches both of the wisdom and knowledge of God! how unsearchable are his judgments, and his ways past finding out! For who hath known the mind of the Lord? or who hath been his counsellor? Or who hath first given to him, and it shall be recompensed unto him again? For of him, and through him, and to him, are all things: to whom be glory for ever. Amen.


\readingcitation{Gospel}{Matthew 28:18}
\lett{A}{t that time:} Jesus said unto his disciples: All power is given unto me in heaven and in earth. Go ye therefore, and teach all nations, baptizing them in the name of the Father, and of the Son, and of the Holy Ghost: Teaching them to observe all things whatsoever I have commanded you: and, lo, I am with you alway, even unto the end of the world.


\subby{Sunday in the Octave of Corpus Christi}
\fancyhead[RE,LO]{Pentecost I}
\centerline{\small{(First Sunday after Trinity)}}
\begin{inhead}
    {Semidouble}
\end{inhead}

\readingcitation{Epistle}{1 John 4:7}
\lett{D}{early beloved:} God is love. In this was manifested the love of God toward us, because that God sent his only begotten Son into the world, that we might live through him. Herein is love, not that we loved God, but that he loved us, and sent his Son to be the propitiation for our sins. Beloved, if God so loved us, we ought also to love one another. No man hath seen God at any time. If we love one another, God dwelleth in us, and his love is perfected in us. Hereby know we that we dwell in him, and he in us, because he hath given us of his Spirit. And we have seen and do testify that the Father sent the Son to be the Saviour of the world. Whosoever shall confess that Jesus is the Son of God, God dwelleth in him, and he in God. And we have known and believed the love that God hath to us. God is love; and he that dwelleth in love dwelleth in God, and God in him. Herein is our love made perfect, that we may have boldness in the day of judgment: because as he is, so are we in this world. There is no fear in love; but perfect love casteth out fear: because fear hath torment. He that feareth is not made perfect in love. We love him, because he first loved us. If a man say, I love God, and hateth his brother, he is a liar: for he that loveth not his brother whom he hath seen, how can he love God whom he hath not seen? And this commandment have we from him, That he who loveth God love his brother also.


\subby{Third Sunday after Pentecost}\label{PentecostIII}
\fancyhead[RE,LO]{Pentecost III}
\begin{inhead}
    {Semidouble}
\end{inhead}

\collect
\lett{O}{God,} the protector of all that trust in thee, without whom nothing is strong, nothing is holy: increase and multiply upon us thy mercy; that, thou being our ruler and guide, we may so pass through things temporal, that we finally lose not the things eternal. Through.
\begin{rubric}
    \nth{2} Collect is of the Saints (p. \pageref{SPSaints}) \& \nth{3} is \emph{ad libitum}.
\end{rubric}

%CHECK:
\readingcitation{Epistle}{1 Peter 5:5}
\lett{D}{early beloved:} All of you be subject one to another, and be clothed with humility: for God resisteth the proud, and giveth grace to the humble. %
Humble yourselves therefore under the mighty hand of God, that he may exalt you in due time: casting all your care upon him; for he careth for you. Be sober, be vigilant; because your adversary the devil, as a roaring lion, walketh about, seeking whom he may devour: whom resist stedfast in the faith, knowing that the same afflictions are accomplished in your brethren that are in the world. But the God of all grace, who hath called us unto his eternal glory by Christ Jesus, after that ye have suffered a while, make you perfect, stablish, strengthen, settle you. To him be glory and dominion for ever and ever. Amen.

\secret
\lett{R}{egard,} O Lord, the gifts of thy suppliant Church: and grant that the partaking thereof may avail for the salvation and continual sanctification of them that believe. Through.
\begin{rubric}
    \nth{2} Secret is of the Saints (p. \pageref{SPSaints}) \& \nth{3} is \emph{ad libitum}.
\end{rubric}

\postcommunion
\lett{M}{ay} the holy things, which we have received, quicken us, O Lord: that we being cleansed thereby may be made ready for thine everlasting mercy. Through.
\begin{rubric}
    \nth{2} Postcommunion is of the Saints (p. \pageref{SPSaints}) \& \nth{3} is \emph{ad libitum}.
\end{rubric}

\subby{Fourth Sunday after Pentecost}
\fancyhead[RE,LO]{Pentecost IV}
\begin{inhead}
    {Semidouble}
\end{inhead}

\collect
\lett{G}{rant,} O Lord, we beseech thee, that the course of this world may be so peaceably ordered by thy governance, that thy Church may joyfully serve thee in all godly quietness. Through.
\begin{rubric}
    Additional Collects are as on Pentecost III (p. \pageref{PentecostIII}).
\end{rubric}

\readingcitation{Gospel}{Luke 5:1}
\lett{A}{t that time:} It came to pass, that as the people pressed upon Jesus to hear the word of God, he stood by the lake of Gennesaret, and saw two ships standing by the lake: but the fishermen were gone out of them, and were washing their nets. And he entered into one of the ships, which was Simon’s, and prayed him that he would thrust out a little from the land. And he sat down, and taught the people out of the ship. Now when he had left speaking, he said unto Simon, Launch out into the deep, and let down your nets for a draught. And Simon answering said unto him, Master, we have toiled all the night, and have taken nothing: nevertheless at thy word I will let down the net. And when they had this done, they inclosed a great multitude of fishes: and their net brake. And they beckoned unto their partners, which were in the other ship, that they should come and help them. And they came, and filled both the ships, so that they began to sink. When Simon Peter saw it, he fell down at Jesus’ knees, saying, Depart from me; for I am a sinful man, O Lord. For he was astonished, and all that were with him, at the draught of the fishes which they had taken: and so was also James, and John, the sons of Zebedee, which were partners with Simon. And Jesus said unto Simon, Fear not; from henceforth thou shalt catch men. And when they had brought their ships to land, they forsook all, and followed him.

\secret
\lett{O}{Lord,} we beseech thee mercifully to receive our oblations: and graciously turn our rebel wills to thee. Through.
\begin{rubric}
    Additional Secrets are as on Pentecost III (p. \pageref{PentecostIII}).
\end{rubric}

\postcommunion
\lett{W}{e} beseech thee, O Lord: that the mysteries which we have received may purify us: and defend us by the gift which they bestow. Through.
\begin{rubric}
    Additional Postcommunions are as on Pentecost III (p. \pageref{PentecostIII}).
\end{rubric}

\subby{Fifth Sunday after Pentecost}
\fancyhead[RE,LO]{Pentecost V}
\begin{inhead}
    {Semidouble}
\end{inhead}

\collect
\lett{O}{God,} who hast prepared for them that love thee such good things as pass man's understanding: pour into our hearts such love toward thee: that we, loving thee above all things, may obtain thy promises, which exceed all that we can desire. Through.
\begin{rubric}
    Additional Collects are as on Pentecost III (p. \pageref{PentecostIII}).
\end{rubric}


\readingcitation{Gospel}{Matthew 5:20}
\lett{A}{t that time:} Jesus said unto his disciples: Except your righteousness shall exceed the righteousness of the scribes and Pharisees, ye shall in no case enter into the kingdom of heaven. Ye have heard that it was said by them of old time, Thou shalt not kill; and whosoever shall kill shall be in danger of the judgment: But I say unto you, That whosoever is angry with his brother without a cause shall be in danger of the judgment: and whosoever shall say to his brother, Raca, shall be in danger of the council: but whosoever shall say, Thou fool, shall be in danger of hell fire. Therefore if thou bring thy gift to the altar, and there rememberest that thy brother hath ought against thee; Leave there thy gift before the altar, and go thy way; first be reconciled to thy brother, and then come and offer thy gift.

\secret
\lett{B}{e} merciful, O Lord, unto our supplications: and graciously receive these oblations of thy servants and handmaids; that those things which each hath offered to the honour of thy name may be profitable unto all for their salvation. Through.
\begin{rubric}
    Additional Secrets are as on Pentecost III (p. \pageref{PentecostIII}).
\end{rubric}

\postcommunion
\lett{O}{Lord} who hast satisfied us with thy heavenly gift: grant, we beseech thee; that we may be cleansed from our secret faults, and delivered from the snares of our enemies. Through.
\begin{rubric}
    Additional Postcommunions are as on Pentecost III (p. \pageref{PentecostIII}).
\end{rubric}

\subby{Sixth Sunday after Pentecost}
\fancyhead[RE,LO]{Pentecost VI}
\begin{inhead}
    {Semidouble}
\end{inhead}

\collect
\lett{G}{od} of all power and might, who art the author of all good things: graft in our hearts the love of thy name and increase in us true religion; that, nourishing all that is good, thou mayest mercifully guard that which thou hast nourished. Through.
\begin{rubric}
    Additional Collects are as on Pentecost III (p. \pageref{PentecostIII}).
\end{rubric}

\readingcitation{Gospel}{Mark 8:1}
\lett{I}{n those days:} The multitude being very great, and having nothing to eat, Jesus called his disciples unto him, and saith unto them, I have compassion on the multitude, because they have now been with me three days, and have nothing to eat: And if I send them away fasting to their own houses, they will faint by the way: for divers of them came from far. And his disciples answered him, From whence can a man satisfy these men with bread here in the wilderness? And he asked them, How many loaves have ye? And they said, Seven. And he commanded the people to sit down on the ground: and he took the seven loaves, and gave thanks, and brake, and gave to his disciples to set before them; and they did set them before the people. And they had a few small fishes: and he blessed, and commanded to set them also before them. So they did eat, and were filled: and they took up of the broken meat that was left seven baskets. And they that had eaten were about four thousand: and he sent them away.


\secret
\lett{B}{e} merciful, O Lord, unto our supplications, and graciously receive these oblations of thy people: and that none may pray amiss or fail in his petition, grant; that what we ask faithfully, we may obtain effectually. Through.
\begin{rubric}
    Additional Secrets are as on Pentecost III (p. \pageref{PentecostIII}).
\end{rubric}

\postcommunion
\lett{W}{e} have been fulfilled with thy gifts, O Lord: grant, we beseech thee; that we may both be cleansed by their operation, and defended by their succour. Through.
\begin{rubric}
    Additional Postcommunions are as on Pentecost III (p. \pageref{PentecostIII}).
\end{rubric}

\subby{Seventh Sunday after Pentecost}
\fancyhead[RE,LO]{Pentecost VII}
\begin{inhead}
    {Semidouble}
\end{inhead}

\collect
\lett{O}{God,} whose never-failing providence ordereth all things both in heaven and earth: we humbly beseech thee; to put away from us all hurtful things, and to give us all things which be profitable for us. Through.
\begin{rubric}
    Additional Collects are as on Pentecost III (p. \pageref{PentecostIII}).
\end{rubric}

\readingcitation{Gospel}{Matthew 7:15}
\lett{A}{t that time:} Jesus said unto his disciples: Beware of false prophets, which come to you in sheep's clothing, but inwardly they are ravening wolves. Ye shall know them by their fruits. Do men gather grapes of thorns, or figs of thistles? Even so every good tree bringeth forth good fruit; but a corrupt tree bringeth forth evil fruit. A good tree cannot bring forth evil fruit, neither can a corrupt tree bring forth good fruit. Every tree that bringeth not forth good fruit is hewn down, and cast into the fire. Wherefore by their fruits ye shall know them. Not every one that saith unto me, Lord, Lord, shall enter into the kingdom of heaven; but he that doeth the will of my Father which is in heaven.

\secret
\lett{O}{God} who by the one perfect sacrifice didst ratify the divers offerings of the law: accept the sacrifice of thy bounden servants, and sanctify it with that blessing wherewith thou didst bless the gifts of Abel: that what each hath offered to the honour of thy majesty may be profitable unto all for their salvation. Through.
\begin{rubric}
    Additional Secrets are as on Pentecost III (p. \pageref{PentecostIII}).
\end{rubric}

\postcommunion
\lett{L}{et} the operation of thy healing, O Lord, both mercifully set us free from our perversities; and lead us to the things that be right. Through.
\begin{rubric}
    Additional Postcommunions are as on Pentecost III (p. \pageref{PentecostIII}).
\end{rubric}

\subby{Eighth Sunday after Pentecost}
\fancyhead[RE,LO]{Pentecost VIII}
\begin{inhead}
    {Semidouble}
\end{inhead}

\collect
\lett{G}{rant} to us, Lord, we beseech thee, the spirit to think and do always such things as be rightful: that we, who cannot exist without thee; may by thee be enabled to live according to thy will. Through.
\begin{rubric}
    Additional Collects are as on Pentecost III (p. \pageref{PentecostIII}).
\end{rubric}

\readingcitation{Gospel}{Luke 16:1}
\lett{A}{t that time:} Jesus spake this parable unto his disciples: There was a certain rich man, which had a steward; and the same was accused unto him that he had wasted his goods. And he called him, and said unto him, How is it that I hear this of thee? give an account of thy stewardship; for thou mayest be no longer steward. Then the steward said within himself, What shall I do? for my lord taketh away from me the stewardship: I cannot dig; to beg I am ashamed. I am resolved what to do, that, when I am put out of the stewardship, they may receive me into their houses. So he called every one of his lord's debtors unto him, and said unto the first, How much owest thou unto my lord? And he said, An hundred measures of oil. And he said unto him, Take thy bill, and sit down quickly, and write fifty. Then said he to another, And how much owest thou? And he said, An hundred measures of wheat. And he said unto him, Take thy bill, and write fourscore. And the lord commended the unjust steward, because he had done wisely: for the children of this world are in their generation wiser than the children of light. And I say unto you, Make to yourselves friends of the mammon of unrighteousness; that, when ye fail, they may receive you into everlasting habitations.

\secret
\lett{A}{ccept,} we beseech thee, O Lord, the gifts which of thine own bounty we offer unto thee: that, by the mighty working of thy grace, these most sacred mysteries may both sanctify our conversation in this life, and bring us to everlasting joys. Through.
\begin{rubric}
    Additional Secrets are as on Pentecost III (p. \pageref{PentecostIII}).
\end{rubric}

\postcommunion
\lett{M}{ay} this heavenly mystery be unto us, O Lord, the renewing both of mind and of body: that like as we do offer unto thee our worship, so we may perceive the benefit of the same. Through.
\begin{rubric}
    Additional Postcommunions are as on Pentecost III (p. \pageref{PentecostIII}).
\end{rubric}

\subby{Ninth Sunday after Pentecost}
\fancyhead[RE,LO]{Pentecost IX}
\begin{inhead}
{Semidouble}
\end{inhead}

\collect
\lett{L}{et} thy merciful ears, O Lord, be open to the prayers of thy humble servants: and that they may obtain their petitions; make them to ask such things as shall please thee. Through.
\begin{rubric}
    Additional Collects are as on Pentecost III (p. \pageref{PentecostIII}).
\end{rubric}

\readingcitation{Gospel}{Luke 19:41}
\lett{A}{t that time:} When Jesus was come near to Jerusalem, he beheld the city, and wept over it, Saying, If thou hadst known, even thou, at least in this thy day, the things which belong unto thy peace! but now they are hid from thine eyes. For the days shall come upon thee, that thine enemies shall cast a trench about thee, and compass thee round, and keep thee in on every side, And shall lay thee even with the ground, and thy children within thee; and they shall not leave in thee one stone upon another; because thou knewest not the time of thy visitation. And he went into the temple, and began to cast out them that sold therein, and them that bought; Saying unto them, It is written, My house is the house of prayer: but ye have made it a den of thieves. And he taught daily in the temple. But the chief priests and the scribes and the chief of the people sought to destroy him.

\secret
\lett{G}{rant} to us, we beseech thee, O Lord, worthily to frequent these mysteries: for as oft as the commemoration of this victim is celebrated, the work of our redemption is performed. Through.
\begin{rubric}
    Additional Secrets are as on Pentecost III (p. \pageref{PentecostIII}).
\end{rubric}

\postcommunion
\lett{W}{e} beseech thee, O Lord, that the communion of thy sacrament may both bestow upon us purification, and grant us purity. Through.
\begin{rubric}
    Additional Postcommunions are as on Pentecost III (p. \pageref{PentecostIII}).
\end{rubric}

\subby{Tenth Sunday after Pentecost}
\fancyhead[RE,LO]{Pentecost X}
\begin{inhead}
{Semidouble}
\end{inhead}

\collect
\lett{O}{God,} who declarest thy almighty power most chiefly in shewing mercy and pity: multiply upon us thy loving-kindness; that we, running toward thy promises, may by thee be made partakers of thy heavenly gifts. Through.
\begin{rubric}
    Additional Collects are as on Pentecost III (p. \pageref{PentecostIII}).
\end{rubric}

\readingcitation{Gospel}{Luke 18:9}
\lett{A}{t that time:} Jesus spake this parable unto certain which trusted in themselves that they were righteous, and despised others: Two men went up into the temple to pray; the one a Pharisee, and the other a publican. The Pharisee stood and prayed thus with himself, God, I thank thee, that I am not as other men are, extortioners, unjust, adulterers, or even as this publican. I fast twice in the week, I give tithes of all that I possess. And the publican, standing afar off, would not lift up so much as his eyes unto heaven, but smote upon his breast, saying, God be merciful to me a sinner. I tell you, this man went down to his house justified rather than the other: for every one that exalteth himself shall be abased; and he that humbleth himself shall be exalted.

\secret
\lett{L}{et} the sacrifices which we dedicate be offered unto thee, O Lord: which thou hast granted us so to present for the honour of thy name, that they may also be for our healing. Through.
\begin{rubric}
    Additional Secrets are as on Pentecost III (p. \pageref{PentecostIII}).
\end{rubric}

\postcommunion
\lett{W}{e} beseech thee, O Lord our God: that as thou failest not to renew us with thy divine sacraments, so thou wouldest not leave us destitute of thy gracious help. Through.
\begin{rubric}
    Additional Postcommunions are as on Pentecost III (p. \pageref{PentecostIII}).
\end{rubric}

\subby{Eleventh Sunday after Pentecost}
\fancyhead[RE,LO]{Pentecost XI}
\begin{inhead}
{Semidouble}
\end{inhead}

\collect
\lett{A}{lmighty} and everlasting God, who in the abundance of thy loving-kindness dost give more than thy suppliants either desire or deserve: pour down upon us thy mercy; forgiving us those things whereof our conscience is afraid, and giving us those things which we do not presume to ask. Through.
\begin{rubric}
    Additional Collects are as on Pentecost III (p. \pageref{PentecostIII}).
\end{rubric}

\readingcitation{Gospel}{Mark 7:31}
\lett{A}{t that time:} Jesus, departing from the coasts of Tyre and Sidon, he came unto the sea of Galilee, through the midst of the coasts of Decapolis. And they bring unto him one that was deaf, and had an impediment in his speech; and they beseech him to put his hand upon him. And he took him aside from the multitude, and put his fingers into his ears, and he spit, and touched his tongue; And looking up to heaven, he sighed, and saith unto him, Ephphatha, that is, Be opened. And straightway his ears were opened, and the string of his tongue was loosed, and he spake plain. And he charged them that they should tell no man: but the more he charged them, so much the more a great deal they published it; And were beyond measure astonished, saying, He hath done all things well: he maketh both the deaf to hear, and the dumb to speak.

\secret
\lett{L}{ook} favourably, O Lord, we beseech thee, on this our service: that the gift which we offer may be acceptable to thee, and avail for the succour of our frailty. Through.
\begin{rubric}
    Additional Secrets are as on Pentecost III (p. \pageref{PentecostIII}).
\end{rubric}

\postcommunion
\lett{W}{e} beseech thee, O Lord, that we who have received thy sacrament may in such wise feel the succour of the same: that being preserved both in body and soul, we may glory in the fulness of thy heavenly healing. Through.
\begin{rubric}
    Additional Postcommunions are as on Pentecost III (p. \pageref{PentecostIII}).
\end{rubric}

\subby{Twelfth Sunday after Pentecost}
\fancyhead[RE,LO]{Pentecost XII}
\begin{inhead}
{Semidouble}
\end{inhead}

\collect
\lett{A}{lmighty} and merciful God, of whose gift it cometh that thy faithful people do unto thee worthy and laudable service: grant, we beseech thee; that we may so faithfully serve thee in this life, that we may without stumbling run the way of thy promises. Through.
\begin{rubric}
    Additional Collects are as on Pentecost III (p. \pageref{PentecostIII}).
\end{rubric}

\readingcitation{Gospel}{Luke 10:23}
\lett{A}{t that time:} Jesus said unto his disciples: Blessed are the eyes which see the things that ye see: For I tell you, that many prophets and kings have desired to see those things which ye see, and have not seen them; and to hear those things which ye hear, and have not heard them. And, behold, a certain lawyer stood up, and tempted him, saying, Master, what shall I do to inherit eternal life? He said unto him, What is written in the law? how readest thou? And he answering said, Thou shalt love the Lord thy God with all thy heart, and with all thy soul, and with all thy strength, and with all thy mind; and thy neighbour as thyself. And he said unto him, Thou hast answered right: this do, and thou shalt live. But he, willing to justify himself, said unto Jesus, And who is my neighbour? And Jesus answering said, A certain man went down from Jerusalem to Jericho, and fell among thieves, which stripped him of his raiment, and wounded him, and departed, leaving him half dead. And by chance there came down a certain priest that way: and when he saw him, he passed by on the other side. And likewise a Levite, when he was at the place, came and looked on him, and passed by on the other side. But a certain Samaritan, as he journeyed, came where he was: and when he saw him, he had compassion on him, And went to him, and bound up his wounds, pouring in oil and wine, and set him on his own beast, and brought him to an inn, and took care of him. And on the morrow when he departed, he took out two pence, and gave them to the host, and said unto him, Take care of him; and whatsoever thou spendest more, when I come again, I will repay thee. Which now of these three, thinkest thou, was neighbour unto him that fell among the thieves? And he said, He that shewed mercy on him. Then said Jesus unto him, Go, and do thou likewise.

\secret
\lett{M}{ercifully} look, we beseech thee, O Lord, upon the sacrifices which we present on thy sacred altars: that while they do bring us pardon, they may minister unto the honour of thy name. Through.
\begin{rubric}
    Additional Secrets are as on Pentecost III (p. \pageref{PentecostIII}).
\end{rubric}

\postcommunion
\lett{Q}{uicken} us, O Lord, we pray thee, who have received these holy mysteries: that they may obtain for us both pardon and protection. Through.
\begin{rubric}
    Additional Postcommunions are as on Pentecost III (p. \pageref{PentecostIII}).
\end{rubric}

\subby{Thirteenth Sunday after Pentecost}
\fancyhead[RE,LO]{Pentecost XIII}
\begin{inhead}
{Semidouble}
\end{inhead}

\collect
\lett{A}{lmighty} and everlasting God, give unto us the increase of faith, hope, and charity: and, that we may be worthy to obtain that which thou dost promise, make us to love that which thou dost command. Through.
\begin{rubric}
    Additional Collects are as on Pentecost III (p. \pageref{PentecostIII}).
\end{rubric}

\readingcitation{Gospel}{Luke 17:11}
\lett{A}{t that time:} As Jesus went to Jerusalem, he passed through the midst of Samaria and Galilee. And as he entered into a certain village, there met him ten men that were lepers, which stood afar off: And they lifted up their voices, and said, Jesus, Master, have mercy on us. And when he saw them, he said unto them, Go shew yourselves unto the priests. And it came to pass, that, as they went, they were cleansed. And one of them, when he saw that he was healed, turned back, and with a loud voice glorified God, And fell down on his face at his feet, giving him thanks: and he was a Samaritan. And Jesus answering said, Were there not ten cleansed? but where are the nine? There are not found that returned to give glory to God, save this stranger. And he said unto him, Arise, go thy way: thy faith hath made thee whole.

\secret
\lett{B}{e} favourable, O Lord, to thy people, be favourable unto their offerings: that, being well-pleased with this oblation, thou mayest both bestow upon us pardon and grant us our petitions. Through.
\begin{rubric}
    Additional Secrets are as on Pentecost III (p. \pageref{PentecostIII}).
\end{rubric}

\postcommunion
\lett{W}{e} beseech thee, O Lord, that we who have received thy heavenly sacraments, may thereby grow in grace to the attainment of eternal redemption. Through.
\begin{rubric}
    Additional Secrets are as on Pentecost III (p. \pageref{PentecostIII}).
\end{rubric}

\subby{Fourteenth Sunday after Pentecost}
\fancyhead[RE,LO]{Pentecost XIV}
\begin{inhead}
{Semidouble}
\end{inhead}

\collect
\lett{K}{eep,} we beseech thee, O Lord, thy Church with thy perpetual mercy: and, because the frailty of man without thee cannot but fall; keep us ever by thy help from things hurtful, and lead us to those things profitable to our salvation. Through.
\begin{rubric}
    Additional Collects are as on Pentecost III (p. \pageref{PentecostIII}).
\end{rubric}

\readingcitation{Gospel}{Matthew 6:24}
\lett{A}{t that time:} Jesus said unto his disciples: No man can serve two masters: for either he will hate the one, and love the other; or else he will hold to the one, and despise the other. Ye cannot serve God and mammon. Therefore I say unto you, Take no thought for your life, what ye shall eat, or what ye shall drink; nor yet for your body, what ye shall put on. Is not the life more than meat, and the body than raiment? Behold the fowls of the air: for they sow not, neither do they reap, nor gather into barns; yet your heavenly Father feedeth them. Are ye not much better than they? Which of you by taking thought can add one cubit unto his stature? And why take ye thought for raiment? Consider the lilies of the field, how they grow; they toil not, neither do they spin: And yet I say unto you, That even Solomon in all his glory was not arrayed like one of these. Wherefore, if God so clothe the grass of the field, which to day is, and to morrow is cast into the oven, shall he not much more clothe you, O ye of little faith? Therefore take no thought, saying, What shall we eat? or, What shall we drink? or, Wherewithal shall we be clothed? (For after all these things do the Gentiles seek:) for your heavenly Father knoweth that ye have need of all these things. But seek ye first the kingdom of God, and his righteousness; and all these things shall be added unto you.

\secret
\lett{G}{rant} to us, O Lord, we beseech thee, that this saving victim may avail, both for the cleansing of our offences, and for the propitiation of thy power. Through.
\begin{rubric}
    Additional Secrets are as on Pentecost III (p. \pageref{PentecostIII}).
\end{rubric}

\postcommunion
\lett{M}{ay} thy sacraments, O God, evermore cleanse and defend us: and bring us to the attainment of eternal salvation. Through.
\begin{rubric}
    Additional Postcommunions are as on Pentecost III (p. \pageref{PentecostIII}).
\end{rubric}

\subby{Fifteenth Sunday after Pentecost}
\fancyhead[RE,LO]{Pentecost XV}
\begin{inhead}
{Semidouble}
\end{inhead}

\collect
\lett{O}{Lord,} we beseech thee, let thy continual pity cleanse and defend thy Church: and, because it cannot continue in safety without thee; govern it evermore by thy goodness. Through.
\begin{rubric}
    Additional Collects are as on Pentecost III (p. \pageref{PentecostIII}).
\end{rubric}

\readingcitation{Epistle}{Galatians 5:25}
\lett{B}{rethren:} If we live in the Spirit, let us also walk in the Spirit. Let us not be desirous of vain glory, provoking one another, envying one another. Brethren, if a man be overtaken in a fault, ye which are spiritual, restore such an one in the spirit of meekness; considering thyself, lest thou also be tempted. Bear ye one another's burdens, and so fulfil the law of Christ. For if a man think himself to be something, when he is nothing, he deceiveth himself. But let every man prove his own work, and then shall he have rejoicing in himself alone, and not in another. For every man shall bear his own burden. Let him that is taught in the word communicate unto him that teacheth in all good things. Be not deceived; God is not mocked: for whatsoever a man soweth, that shall he also reap. For he that soweth to his flesh shall of the flesh reap corruption; but he that soweth to the Spirit shall of the Spirit reap life everlasting. And let us not be weary in well doing: for in due season we shall reap, if we faint not. As we have therefore opportunity, let us do good unto all men, especially unto them who are of the household of faith.

\readingcitation{Gospel}{Luke 7:11}
\lett{A}{t that time:} Jesus went into a city called Nain; and many of his disciples went with him, and much people. Now when he came nigh to the gate of the city, behold, there was a dead man carried out, the only son of his mother, and she was a widow: and much people of the city was with her. And when the Lord saw her, he had compassion on her, and said unto her, Weep not. And he came and touched the bier: and they that bare him stood still. And he said, Young man, I say unto thee, Arise. And he that was dead sat up, and began to speak. And he delivered him to his mother. And there came a fear on all: and they glorified God, saying, That a great prophet is risen up among us; and, That God hath visited his people.

\secret
\lett{M}{ay} thy sacraments guard us, O Lord: and ever protect us against the assaults of the devil. Through.
\begin{rubric}
    Additional Secrets are as on Pentecost III (p. \pageref{PentecostIII}).
\end{rubric}

\postcommunion
\lett{W}{e} beseech thee, O Lord, that the operation of thy heavenly gift may so possess our souls and bodies: that not our own desires, but the effectual working of the same may ever prevail within us. Through.
\begin{rubric}
    Additional Postcommunions are as on Pentecost III (p. \pageref{PentecostIII}).
\end{rubric}

\subby{Sixteenth Sunday after Pentecost}
\fancyhead[RE,LO]{Pentecost XVI}
\begin{inhead}
{Semidouble}
\end{inhead}

\collect
\lett{L}{ord,} we pray thee that thy grace may always prevent and follow us: and make us continually to be given to good works. Through.
\begin{rubric}
    Additional Collects are as on Pentecost III (p. \pageref{PentecostIII}).
\end{rubric}

\readingcitation{Gospel}{Luke 14:1}
\lett{A}{t that time:} When Jesus went into the house of one of the chief Pharisees to eat bread on the sabbath day, that they watched him. And, behold, there was a certain man before him which had the dropsy. And Jesus answering spake unto the lawyers and Pharisees, saying, Is it lawful to heal on the sabbath day? And they held their peace. And he took him, and healed him, and let him go; And answered them, saying, Which of you shall have an ass or an ox fallen into a pit, and will not straightway pull him out on the sabbath day? And they could not answer him again to these things. And he put forth a parable to those which were bidden, when he marked how they chose out the chief rooms; saying unto them, When thou art bidden of any man to a wedding, sit not down in the highest room; lest a more honourable man than thou be bidden of him; And he that bade thee and him come and say to thee, Give this man place; and thou begin with shame to take the lowest room. But when thou art bidden, go and sit down in the lowest room; that when he that bade thee cometh, he may say unto thee, Friend, go up higher: then shalt thou have worship in the presence of them that sit at meat with thee. For whosoever exalteth himself shall be abased; and he that humbleth himself shall be exalted.

\secret
\lett{C}{leanse} us, O Lord, we pray thee, by the effectual working of this our sacrifice: and so accomplish in us the work of thy mercy; that we may be found worthy to be made partakers of the same. Through.
\begin{rubric}
    Additional Secrets are as on Pentecost III (p. \pageref{PentecostIII}).
\end{rubric}

\postcommunion
\lett{O}{Lord,} we beseech thee, mercifully to purify our minds and renew them with thy heavenly sacraments: that we may thereby have help also for our bodies both now and for evermore. Through.
\begin{rubric}
    Additional Postcommunions are as on Pentecost III (p. \pageref{PentecostIII}).
\end{rubric}

\subby{Seventeenth Sunday after Pentecost}
\fancyhead[RE,LO]{Pentecost XVII}
\begin{inhead}
{Semidouble}
\end{inhead}

\collect
\lett{L}{ord,} we beseech thee: grant thy people grace to avoid the infections of the devil: and with pure minds to follow thee, the only God. Through.
\begin{rubric}
    Additional Collects are as on Pentecost III (p. \pageref{PentecostIII}).
\end{rubric}

\readingcitation{Gospel}{Matthew 22:34}
\lett{A}{t that time:} The Pharisees came to Jesus: and one of them, which was a lawyer, asked him a question, tempting him, and saying, Master, which is the great commandment in the law? Jesus said unto him, Thou shalt love the Lord thy God with all thy heart, and with all thy soul, and with all thy mind. This is the first and great commandment. And the second is like unto it, Thou shalt love thy neighbour as thyself. On these two commandments hang all the law and the prophets. While the Pharisees were gathered together, Jesus asked them, Saying, What think ye of Christ? whose son is he? They say unto him, The Son of David. He saith unto them, How then doth David in spirit call him Lord, saying, The \divineName{Lord} said unto my Lord, Sit thou on my right hand, till I make thine enemies thy footstool? If David then call him Lord, how is he his son? And no man was able to answer him a word, neither durst any man from that day forth ask him any more questions.

\secret
\lett{W}{e} humbly entreat thy majesty, O Lord: that these holy mysteries which we perform may deliver us both from past and future transgressions. Through.
\begin{rubric}
    Additional Secrets are as on Pentecost III (p. \pageref{PentecostIII}).
\end{rubric}

\postcommunion
\lett{A}{lmighty} God, let thy holy mysteries both heal our vices, and bestow upon us everlasting remedies. Through.
\begin{rubric}
    Additional Postcommunions are as on Pentecost III (p. \pageref{PentecostIII}).
\end{rubric}

\subby{Eighteenth Sunday after Pentecost}
\fancyhead[RE,LO]{Pentecost XVIII}
\begin{inhead}
{Semidouble}
\end{inhead}

\collect
\lett{L}{et} the operation of thy mercy, O Lord, direct our hearts: forasmuch as without thee we are not able to please thee. Through.
\begin{rubric}
    Additional Collects are as on Pentecost III (p. \pageref{PentecostIII}).
\end{rubric}

\readingcitation{Gospel}{Matthew 9:1}
\lett{A}{t that time:} Jesus entered into a ship, and passed over, and came into his own city. And, behold, they brought to him a man sick of the palsy, lying on a bed: and Jesus seeing their faith said unto the sick of the palsy; Son, be of good cheer; thy sins be forgiven thee. And, behold, certain of the scribes said within themselves, This man blasphemeth. And Jesus knowing their thoughts said, Wherefore think ye evil in your hearts? For whether is easier, to say, Thy sins be forgiven thee; or to say, Arise, and walk? But that ye may know that the Son of man hath power on earth to forgive sins, (then saith he to the sick of the palsy,) Arise, take up thy bed, and go unto thine house. And he arose, and departed to his house. But when the multitudes saw it, they marvelled, and glorified God, which had given such power unto men.

\secret
\lett{O}{God,} who, by communion in this venerable sacrifice, hast made us partakers of the one supreme Godhead: grant, we beseech thee; that as we have the knowledge of thy truth, so by worthy conversation we may attain unto the same. Through.
\begin{rubric}
    Additional Secrets are as on Pentecost III (p. \pageref{PentecostIII}).
\end{rubric}

\postcommunion
\lett{W}{e} render thanks unto thee, O Lord, for that thou hast quickened us with thy sacred gift: humbly beseeching thy mercy; that thou wouldest make us worthy partakers of the same. Through.
\begin{rubric}
    Additional Postcommunions are as on Pentecost III (p. \pageref{PentecostIII}).
\end{rubric}


\subby{Nineteenth Sunday after Pentecost}
\fancyhead[RE,LO]{Pentecost XIX}
\begin{inhead}
{Semidouble}
\end{inhead}

\collect
\lett{O}{Almighty} and most merciful God: of thy bountiful goodness keep us from all things that may hurt us; that we, being ready both in body and soul, may cheerfully accomplish those things that thou wouldest have done. Through.
\begin{rubric}
    Additional Collects are as on Pentecost III (p. \pageref{PentecostIII}).
\end{rubric}

%CHECK:
\readingcitation{Epistle}{Ephesians 4:17}
\lett{B}{rethren:} This I say therefore, and testify in the Lord, that ye henceforth walk not as other Gentiles walk, in the vanity of their mind, Having the understanding darkened, being alienated from the life of God through the ignorance that is in them, because of the blindness of their heart: Who being past feeling have given themselves over unto lasciviousness, to work all uncleanness with greediness. But ye have not so learned Christ; If so be that ye have heard him, and have been taught by him, as the truth is in Jesus: That ye put off concerning the former conversation the old man, which is corrupt according to the deceitful lusts;
%
Be renewed in the spirit of your mind; And that ye put on the new man, which after God is created in righteousness and true holiness. Wherefore putting away lying, speak every man truth with his neighbour: for we are members one of another. Be ye angry, and sin not: let not the sun go down upon your wrath: Neither give place to the devil. Let him that stole steal no more: but rather let him labour, working with his hands the thing which is good, that he may have to give to him that needeth.
%
Let no corrupt communication proceed out of your mouth, but that which is good to the use of edifying, that it may minister grace unto the hearers. And grieve not the holy Spirit of God, whereby ye are sealed unto the day of redemption. Let all bitterness, and wrath, and anger, and clamour, and evil speaking, be put away from you, with all malice: And be ye kind one to another, tenderhearted, forgiving one another, even as God for Christ's sake hath forgiven you.

\readingcitation{Gospel}{Matthew 22:1}
\lett{A}{t that time:} Jesus spake unto the chief priests and Pharisees in parables, saying: The kingdom of heaven is like unto a certain king, which made a marriage for his son, And sent forth his servants to call them that were bidden to the wedding: and they would not come. Again, he sent forth other servants, saying, Tell them which are bidden, Behold, I have prepared my dinner: my oxen and my fatlings are killed, and all things are ready: come unto the marriage. But they made light of it, and went their ways, one to his farm, another to his merchandise: And the remnant took his servants, and entreated them spitefully, and slew them. But when the king heard thereof, he was wroth: and he sent forth his armies, and destroyed those murderers, and burned up their city. Then saith he to his servants, The wedding is ready, but they which were bidden were not worthy. Go ye therefore into the highways, and as many as ye shall find, bid to the marriage. So those servants went out into the highways, and gathered together all as many as they found, both bad and good: and the wedding was furnished with guests. And when the king came in to see the guests, he saw there a man which had not on a wedding garment: And he saith unto him, Friend, how camest thou in hither not having a wedding garment? And he was speechless. Then said the king to the servants, Bind him hand and foot, and take him away, and cast him into outer darkness; there shall be weeping and gnashing of teeth. For many are called, but few are chosen.

\secret
\lett{G}{rant,} we beseech thee, O Lord, that these gifts which we offer in the sight of thy majesty, may be for our salvation. Through.
\begin{rubric}
    Additional Secrets are as on Pentecost III (p. \pageref{PentecostIII}).
\end{rubric}

\postcommunion
\lett{M}{ay} the operation of thy healing, O Lord, both mercifully cleanse us from our perversities, and make us ever to cleave unto thy commandments. Through.
\begin{rubric}
    Additional Postcommunions are as on Pentecost III (p. \pageref{PentecostIII}).
\end{rubric}

\subby{Twentieth Sunday after Pentecost}
\fancyhead[RE,LO]{Pentecost XX}
\begin{inhead}
{Semidouble}
\end{inhead}

\collect
\lett{G}{rant,} we beseech thee, merciful Lord, to thy faithful people pardon and peace: that they may be cleansed from all their sins; and serve thee with a quiet mind. Through.
\begin{rubric}
    Additional Collects are as on Pentecost III (p. \pageref{PentecostIII}).
\end{rubric}

\readingcitation{Gospel}{John 4:46}
\lett{A}{t that time:} There was a certain nobleman, whose son was sick at Capernaum. When he heard that Jesus was come out of Judaea into Galilee, he went unto him, and besought him that he would come down, and heal his son: for he was at the point of death. Then said Jesus unto him, Except ye see signs and wonders, ye will not believe. The nobleman saith unto him, Sir, come down ere my child die. Jesus saith unto him, Go thy way; thy son liveth. And the man believed the word that Jesus had spoken unto him, and he went his way. And as he was now going down, his servants met him, and told him, saying, Thy son liveth. Then enquired he of them the hour when he began to amend. And they said unto him, Yesterday at the seventh hour the fever left him. So the father knew that it was at the same hour, in the which Jesus said unto him, Thy son liveth: and himself believed, and his whole house.

\secret
\lett{W}{e} beseech thee, O Lord, that these mysteries may bestow upon us thy heavenly healing: and purge away the vices of our hearts. Through.
\begin{rubric}
    Additional Secrets are as on Pentecost III (p. \pageref{PentecostIII}).
\end{rubric}

\postcommunion
\lett{W}{e} beseech thee, O Lord, make us ever to obey thy commandments: that we may be rendered worthy of thy sacred gifts. Through.
\begin{rubric}
    Additional Postcommunions are as on Pentecost III (p. \pageref{PentecostIII}).
\end{rubric}

%CHANGE LABEL?
\subby{Twenty-First Sunday after Pentecost}\label{trinity}
\fancyhead[RE,LO]{Pentecost XXI}
\begin{inhead}
{Semidouble}
\end{inhead}

\collect
\lett{L}{ord,} we beseech thee to keep thy household in continual godliness: that through thy protection it may be free from all adversities; and in good works be devoted to thy name. Through.
\begin{rubric}
    Additional Collects are as on Pentecost III (p. \pageref{PentecostIII}).
\end{rubric}

%CHECK:
\readingcitation{Epistle}{Ephesians 6:10}
\lett{B}{rethren:} Be strong in the Lord, and in the power of his might. Put on the whole armour of God, that ye may be able to stand against the wiles of the devil. For we wrestle not against flesh and blood, but against principalities, against powers, against the rulers of the darkness of this world, against spiritual wickedness in high places. Wherefore take unto you the whole armour of God, that ye may be able to withstand in the evil day, and having done all, to stand. Stand therefore, having your loins girt about with truth, and having on the breastplate of righteousness; And your feet shod with the preparation of the gospel of peace; Above all, taking the shield of faith, wherewith ye shall be able to quench all the fiery darts of the wicked. And take the helmet of salvation, and the sword of the Spirit, which is the word of God:
%
Praying always with all prayer and supplication in the Spirit, and watching thereunto with all perseverance and supplication for all saints; And for me, that utterance may be given unto me, that I may open my mouth boldly, to make known the mystery of the gospel, For which I am an ambassador in bonds: that therein I may speak boldly, as I ought to speak.

\readingcitation{Gospel}{Matthew 18:23}
\lett{A}{t that time:} Jesus spake unto his disciples this parable: The kingdom of heaven is likened unto a certain king, which would take account of his servants. And when he had begun to reckon, one was brought unto him, which owed him ten thousand talents. But forasmuch as he had not to pay, his lord commanded him to be sold, and his wife, and children, and all that he had, and payment to be made. The servant therefore fell down, and worshipped him, saying, Lord, have patience with me, and I will pay thee all. Then the lord of that servant was moved with compassion, and loosed him, and forgave him the debt. But the same servant went out, and found one of his fellowservants, which owed him an hundred pence: and he laid hands on him, and took him by the throat, saying, Pay me that thou owest. And his fellowservant fell down at his feet, and besought him, saying, Have patience with me, and I will pay thee all. And he would not: but went and cast him into prison, till he should pay the debt. So when his fellowservants saw what was done, they were very sorry, and came and told unto their lord all that was done. Then his lord, after that he had called him, said unto him, O thou wicked servant, I forgave thee all that debt, because thou desiredst me: Shouldest not thou also have had compassion on thy fellowservant, even as I had pity on thee? And his lord was wroth, and delivered him to the tormentors, till he should pay all that was due unto him. So likewise shall my heavenly Father do also unto you, if ye from your hearts forgive not every one his brother their trespasses.

\secret
\lett{M}{ercifully} receive, O Lord, these offerings: wherewith thou hast willed that atonement should be made unto thee, and that salvation should be restored to us by the power of thy loving kindness. Through.
\begin{rubric}
    Additional Secrets are as on Pentecost III (p. \pageref{PentecostIII}).
\end{rubric}

\postcommunion
\lett{O}{Lord,} who hast bestowed upon us the food of immortality, we beseech thee: that those things which we have taken with our lips we may seek after in purity of heart. Through.
\begin{rubric}
    Additional Postcommunions are as on Pentecost III (p. \pageref{PentecostIII}).
\end{rubric}

\subby{Twenty-Second Sunday after Pentecost}
\fancyhead[RE,LO]{Pentecost XXII}
\begin{inhead}
{Semidouble}
\end{inhead}

\collect
\lett{O}{God,} our refuge and strength, who art the author of all godliness: hear the devout prayers of thy Church, and grant; that those things which we ask faithfully we may obtain effectually. Through.
\begin{rubric}
    Additional Collects are as on Pentecost III (p. \pageref{PentecostIII}).
\end{rubric}

\readingcitation{Gospel}{Matthew 22:15}
\lett{A}{t that time:} The Pharisees went, and took counsel how they might entangle Jesus in his talk. And they sent out unto him their disciples with the Herodians, saying, Master, we know that thou art true, and teachest the way of God in truth, neither carest thou for any man: for thou regardest not the person of men. Tell us therefore, What thinkest thou? Is it lawful to give tribute unto Caesar, or not? But Jesus perceived their wickedness, and said, Why tempt ye me, ye hypocrites? Shew me the tribute money. And they brought unto him a penny. And he saith unto them, Whose is this image and superscription? They say unto him, Caesar's. Then saith he unto them, Render therefore unto Caesar the things which are Caesar's; and unto God the things that are God's.

\secret
\lett{W}{e} offer thee, O Lord, this sacrifice of praise for the increase of our dutiful service: that thou wouldest mercifully accomplish that which thou hast bestowed upon us beyond our deserving. Through.
\begin{rubric}
    Additional Secrets are as on Pentecost III (p. \pageref{PentecostIII}).
\end{rubric}

\postcommunion
\lett{W}{e} beseech thee, almighty God: that we, whom thou makest to rejoice in the partaking of heavenly things, may by thee be defended against all earthly perils. Through.
\begin{rubric}
    Additional Postcommunions are as on Pentecost III (p. \pageref{PentecostIII}).
\end{rubric}

\subby{Twenty-Third Sunday after Pentecost}
\fancyhead[RE,LO]{Pentecost XXIII}
\begin{inhead}
{Semidouble}
\end{inhead}
\begin{rubric}
	If this Sunday be hindered by the supervention of the last after Pentecost, it is anticipated on Saturday with all the privileges proper to a Sunday occurring, and in it are said \emph{Gloria in excelsis}, the Nicene Creed, Preface of the Trinity, and \emph{Ite, Missa est}.
\end{rubric}

\collect
\lett{O}{Lord,} we beseech thee absolve thy people from their offences; that through thy bountiful goodness we may be delivered from the bonds of those sins, which by our frailty we have committed. Through.
\begin{rubric}
    Additional Collects are as on Pentecost III (p. \pageref{PentecostIII}).
\end{rubric}

\readingcitation{Epistle}{Philippians 3:17}
\lett{B}{rethren:} Be followers together of me, and mark them which walk so as ye have us for an ensample. (For many walk, of whom I have told you often, and now tell you even weeping, that they are the enemies of the cross of Christ: Whose end is destruction, whose God is their belly, and whose glory is in their shame, who mind earthly things.) For our conversation is in heaven; from whence also we look for the Saviour, the Lord Jesus Christ: Who shall change our vile body, that it may be fashioned like unto his glorious body, according to the working whereby he is able even to subdue all things unto himself. Therefore, my brethren dearly beloved and longed for, my joy and crown, so stand fast in the Lord, my dearly beloved. I beseech Euodias, and beseech Syntyche, that they be of the same mind in the Lord. And I intreat thee also, true yokefellow, help those women which laboured with me in the gospel, with Clement also, and with other my fellowlabourers, whose names are in the book of life.

\readingcitation{Gospel}{Matthew 9:18}
\lett{A}{t that time:} While Jesus spake unto the multitudes, Behold, there came a certain ruler, and worshipped him, saying, My daughter is even now dead: but come and lay thy hand upon her, and she shall live. And Jesus arose, and followed him, and so did his disciples. And, behold, a woman, which was diseased with an issue of blood twelve years, came behind him, and touched the hem of his garment: For she said within herself, If I may but touch his garment, I shall be whole. But Jesus turned him about, and when he saw her, he said, Daughter, be of good comfort; thy faith hath made thee whole. And the woman was made whole from that hour. And when Jesus came into the ruler's house, and saw the minstrels and the people making a noise, He said unto them, Give place: for the maid is not dead, but sleepeth. And they laughed him to scorn. But when the people were put forth, he went in, and took her by the hand, and the maid arose. And the fame hereof went abroad into all that land.

\subby{Twenty-Fourth---and last---Sunday after Pentecost}
\fancyhead[RE,LO]{Pentecost XXIV}
\begin{inhead}
{Semidouble}
\end{inhead}

\collect
\lett{S}{tir} up, we beseech thee, O Lord, the wills of thy faithful people; that they, plenteously bringing forth the fruit of good works, may by thee be plenteously rewarded. Through.
\begin{rubric}
    Additional Collects are as on Pentecost III (p. \pageref{PentecostIII}).
\end{rubric}

\readingcitation{Epistle}{Colossians 1:9}
\lett{B}{rethren:} We do not cease to pray for you, and to desire that ye might be filled with the knowledge of his will in all wisdom and spiritual understanding; That ye might walk worthy of the Lord unto all pleasing, being fruitful in every good work, and increasing in the knowledge of God; Strengthened with all might, according to his glorious power, unto all patience and longsuffering with joyfulness; Giving thanks unto the Father, which hath made us meet to be partakers of the inheritance of the saints in light: Who hath delivered us from the power of darkness, and hath translated us into the kingdom of his dear Son: In whom we have redemption through his blood, even the forgiveness of sins.

\readingcitation{Gospel}{Matthew 24:15}
\lett{A}{t that time:} Jesus said unto his disciples: When ye shall see the abomination of desolation, spoken of by Daniel the prophet, stand in the holy place, (whoso readeth, let him understand:) Then let them which be in Judaea flee into the mountains: Let him which is on the housetop not come down to take any thing out of his house: Neither let him which is in the field return back to take his clothes. And woe unto them that are with child, and to them that give suck in those days! But pray ye that your flight be not in the winter, neither on the sabbath day: For then shall be great tribulation, such as was not since the beginning of the world to this time, no, nor ever shall be. And except those days should be shortened, there should no flesh be saved: but for the elect's sake those days shall be shortened. Then if any man shall say unto you, Lo, here is Christ, or there; believe it not. For there shall arise false Christs, and false prophets, and shall shew great signs and wonders; insomuch that, if it were possible, they shall deceive the very elect. Behold, I have told you before. Wherefore if they shall say unto you, Behold, he is in the desert; go not forth: behold, he is in the secret chambers; believe it not. For as the lightning cometh out of the east, and shineth even unto the west; so shall also the coming of the Son of man be. For wheresoever the carcase is, there will the eagles be gathered together. Immediately after the tribulation of those days shall the sun be darkened, and the moon shall not give her light, and the stars shall fall from heaven, and the powers of the heavens shall be shaken: And then shall appear the sign of the Son of man in heaven: and then shall all the tribes of the earth mourn, and they shall see the Son of man coming in the clouds of heaven with power and great glory. And he shall send his angels with a great sound of a trumpet, and they shall gather together his elect from the four winds, from one end of heaven to the other. Now learn a parable of the fig tree; When his branch is yet tender, and putteth forth leaves, ye know that summer is nigh: So likewise ye, when ye shall see all these things, know that it is near, even at the doors. Verily I say unto you, This generation shall not pass, till all these things be fulfilled. Heaven and earth shall pass away, but my words shall not pass away.

\secret
\lett{B}{e} favourable, O Lord, to our supplications: and, accepting the offerings and prayers of us thy people, turn the hearts of all to thee; that we, being delivered from earthly lusts, may turn to heavenly desires. Through.
\begin{rubric}
    Additional Secrets are as on Pentecost III (p. \pageref{PentecostIII}).
\end{rubric}

\postcommunion
\lett{G}{rant} to us, we beseech thee, O Lord: that through these sacraments, which we have received, whatsoever is corrupt in our minds may be made whole by the healing power of the same. Through.
\begin{rubric}
    Additional Postcommunions are as on Pentecost III (p. \pageref{PentecostIII}).
\end{rubric}

\begin{rubric}
    If there be more than 24 Sundays after Pentecost, then the Sundays after the \nth{23} after Pentecost are supplied from the unused Sundays after Epiphany (p. \pageref{epiphany}). Then, the \nth{24} Sunday after Pentecost.
\end{rubric}


\fancyhead[RE,LO]{}\fancyhead[RO,LE]{}
\fancyhead[C]{}\thispagestyle{empty}
\phantomsection
\addcontentsline{toc}{chapter}{Propers of the Church Year}
  \begin{center}
   \includegraphics[width=98mm, height=140mm]{Vision.eps}
   \par
   \vspace{2ex}
   	\textsc{\Huge{Propers of the\\
   	Church Year}}
   \end{center}

%\phantomsection
\addcontentsline{toc}{section}{Holy Week}
\fancyhead[C]{\LARGE Holy Week}

\subby{Holy Monday}
\fancyhead[RO,LE]{\textit{Holy Monday}}
\fancyhead[RE,LO]{}
\begin{inhead}
    {First Class Feria}
\end{inhead}

	\properantiphon{Ben.}{O Father, glorify thou me {\dag} with thine own self, with the glory which I had with thee before the world began.}

    \properantiphon{Mag.}{Thou couldest have no power {\dag} at all against me, except it were given thee from above.}
    
   \introit
   \lett{P}{lead} thou my cause, O Lord, with them that strive with me, and fight thou against them that fight against me: lay hand upon the shield and buckler, and stand up to help me, O Lord, thou strength of my salvation. \textit{Ps.} Bring forth the spear, and stop the way against them that persecute me: say unto my soul: I am thy salvation.
   
   \collect
	\lett{A}{lmighty} God, whose most dear Son went not up to joy but first he suffered pain, and entered not into glory before he was crucified; Mercifully grant that we, walking in the way of the cross, may find it none other than the way of life and peace. Through the same.
	
\begin{rubric}
    \nth{2} Collect is \emph{Against the Persecutors of the Church} (p. \SPAgainst) or \emph{for the Chief Bishop} (p. \SPChiefBishop).
\end{rubric}
	
\readingcitation{Epistle}{Isaiah 63:1}
\lett{I}{n those days:} Said Isaiah: Who is this that cometh from Edom, with dyed garments from Bozrah? this that is glorious in his apparel, travelling in the greatness of his strength? I that speak in righteousness, mighty to save. Wherefore art thou red in thine apparel, and thy garments like him that treadeth in the winefat? I have trodden the winepress alone; and of the people there was none with me: for I will tread them in mine anger, and trample them in my fury; and their blood shall be sprinkled upon my garments, and I will stain all my raiment. For the day of vengeance is in mine heart, and the year of my redeemed is come. And I looked, and there was none, to help; and I wondered that there was none to uphold: therefore mine own arm brought salvation unto me; and my fury it upheld me. And I will tread down the people in mine anger, and make them drunk in my fury, and I will bring down their strength to the earth. I will mention the lovingkindnesses of the Lord, and the praises of the Lord according to all that the Lord hath bestowed on us, and the great goodness toward the house of Israel, which he hath bestowed on them according to his mercies, and according to the multitude of his lovingkindnesses. For he said, Surely they are my people, children that will not lie: so he was their Saviour. In all their affliction he was afflicted, and the angel of his presence saved them: in his love and in his pity he redeemed them; and he bare them, and carried them all the days of old. But they rebelled, and vexed his holy Spirit; therefore he was turned to be their enemy, and he fought against them. Then he remembered the days of old, Moses, and his people, saying, Where is he that brought them up out of the sea with the shepherd of his flock? where is he that put his holy Spirit within him? that led them by the right hand of Moses with his glorious arm, dividing the water before them, to make himself an everlasting name? that led them through the deep, as an horse in the wilderness., that they should not stumble? As a beast goeth down into the valley, the spirit of the Lord caused him to rest: so didst thou lead thy people, to make thyself a glorious name. Look down from heaven, and behold from the habitation of thy holiness and of thy glory: where is thy zeal and thy strength, the sounding of thy bowels and of thy mercies toward me? are they restrained? Doubtless thou art our father, though Abraham be ignorant of us, and Israel acknowledge us not, thou, O Lord, art our father our redeemer; thy name is from everlasting. O Lord, why hast thou made us to err from thy ways, and hardened our heart from thy fear? Return for thy servants' sake, the tribes of thine inheritance. The people of thy holiness have possessed it but a little while: our adversaries have trodden down thy sanctuary. We are thine: thou never barest rule over them; they were not called by thy name.

\gradtr{Awake, O Lord, and stand up to judge my quarrel, avenge thou my cause, my God, and my Lord. ℣. Bring forth the spear, and stop the way against them that persecute me.}{O Lord, deal not with us after our sins which we have committed: nor reward us according to our wickednesses. ℣. Lord, remember not our old sins, but have mercy upon us, and that soon: for we are come to great misery. \inrub{Here genuflect.} ℣. Help us, O God of our salvation: and for the glory of thy name, O Lord, deliver us: and be merciful unto our sins, for thy name's sake. \inrub{Arise.}}

\readingcitation{Gospel}{Mark 14:1}
\lett{A}{t that time:} After two days was the feast of the passover, and of unleavened bread: and the chief priests and the scribes sought how they might take him by craft, and put him to death. But they said, Not on the feast day, lest there be an uproar of the people. And being in Bethany in the house of Simon the leper, as he sat at meat, there came a woman having an alabaster box of ointment of spikenard very precious; and she brake the box, and poured it on his head. And there were some that had indignation within themselves, and said, Why was this waste of the ointment made? for it might have been sold for more than three hundred pence, and have been given to the poor. And they murmured against her. And Jesus said, Let her alone; why trouble ye her? she hath wrought a good work on me. For ye have the poor with you always, and whensoever ye will ye may do them good: but me ye have not always. She hath done what she could: she is come aforehand to anoint my body to the burying. Verily I say unto you, Wheresoever this gospel shall be preached throughout the whole world, this also that she hath done shall be spoken of for a memorial of her. And Judas Iscariot, one of the twelve, went unto the chief priests, to betray him unto them. And when they heard it, they were glad, and promised to give him money. And he sought how he might conveniently betray him. And the first day of unleavened bread, when they killed the passover, his disciples said unto him, Where wilt thou that we go and prepare that thou mayest eat the passover? And he sendeth forth two of his disciples, and saith unto them, Go ye into the city, and there shall meet you a man bearing a pitcher of water: follow him. And wheresoever he shall go in, say ye to the good-man of the house, The Master saith, Where is the guestchamber, where I shall eat the passover with my disciples? And he will shew you a large upper room furnished and prepared: there make ready for us. And his disciples went forth, and came into the city, and found as he had said unto them: and they made ready the passover. And in the evening he cometh with the twelve. And as they sat and did eat, Jesus said, Verily I say unto you, One of you which eateth with me shall betray me. And they began to be sorrowful, and to say unto him one by one, Is it I? and another said, Is it I? And he answered and said unto them, It is one of the twelve, that dippeth with me in the dish. The Son of man indeed goeth, as it is written of him: but woe to that man by whom the Son of man is betrayed! good were it for that man if he had never been born. And as they did eat, Jesus took bread, and blessed, and brake it, and gave to them, and said, Take, eat; this is my body. And he took the cup, and when he had given thanks, he gave it to them: and they all drank of it. And he said unto them, This is my blood of the new testament, which is shed for many. Verily I say unto you, I will drink no more of the fruit of the vine, until that day that I drink it new in the kingdom of God. And when they had sung an hymn, they went out into the mount of Olives. And Jesus saith unto them, All ye shall be offended because of me this night: for it is written, I will smite the shepherd, and the sheep shall be scattered. But after that I am risen I will go before you into Galilee. But Peter said unto him, Although all shall be offended, yet will not I. And Jesus saith unto him, Verily I say unto thee, that this day even in this night, before the cock crow twice, thou shalt deny me thrice. But he spake the more vehemently, If I should die with thee, I will not deny thee in any wise. Likewise also said they all. And they came to a place which was named Gethsemane: and he saith to his disciples, Sit ye here, while I shall pray. And he taketh with him Peter and James and John, and began to be sore amazed, and to be very heavy; and saith unto them, My soul is exceeding sorrowful unto death: tarry ye here, and watch. And he went forward a little, and fell on the ground, and prayed that, if it were possible, the hour might pass from him. And he said, Abba, Father, all things are possible unto thee; take away this cup from me: nevertheless not what I will, but what thou wilt. And he cometh. and findeth them sleeping, and saith unto Peter, Simon, sleepest thou? couldest not thou watch one hour? Watch ye and pray, lest ye enter into temptation. The spirit truly is ready, but the flesh is weak. And again lie went away, and prayed, and spake the same words. And when ho returned, he found them asleep again, (for their eyes were heavy,) neither wist they what to answer him. And he cometh the third time, and saith unto them, Sleep on now, and take your rest: it is enough, the hour is come; behold, the Son of man is betrayed into the hands of sinners. Rise up, let us go; lo, he that betrayeth me is at hand. And immediately, while he yet spake, cometh Judas, one of the twelve, and with him a great multitude with swords and staves, from the chief priests and the scribes and the elders. And he that betrayed him had given them a token, saying, Whomsoever I shall kiss, that same is he; take him, and lead him away safely. And as soon as he was come, he goeth straightway to him, and saith, Master, master; and kissed him. And they laid their hands on him, and took him. And one of them that stood by drew a sword, and smote a servant of the high priest, and cut off his ear. And Jesus answered and said unto them, Are ye come out, as against a thief, with swords and with staves, to take me? I was daily with you in the temple teaching, and ye took me not: but the scriptures must be fulfilled. And they all forsook him, and fled. And there followed him a certain young man, having a linen cloth cast about his naked body; and the young men laid hold on him: and he left the linen cloth, and fled from them naked. And they led Jesus away to the high priest: and with him were assembled all the chief priests and the elders and the scribes. And Peter followed him afar off, even into the palace of the high priest: and he sat with the servants, and warmed himself at the fire. And the chief priests and all the council sought for witness against Jesus to put him to death; and found none. For many bare false witness against him, but their witness agreed not together. And there arose certain, and bare false witness against him, saying, We heard him say, I will destroy this temple that is made with hands, and within three days I will build another made without hands. But neither so did their witness agree together. And the high priest stood up in the midst, and asked Jesus, saying, Answerest thou nothing? what is it which these witness against thee? But he held his peace, and answered nothing. Again the high priest asked him, and said unto him, Art thou the Christ, the Son of the Blessed? And Jesus said, I am: and ye shall see the Son of man sitting on the right hand of power, and coming in the clouds of heaven. Then the high priest rent his clothes, and saith, What need we any further witnesses? ye have heard the blasphemy: what think ye? And they all condemned him to be guilty of death. And some began to spit on him, and to cover his face, and to buffet him, and to say unto him, Prophesy: and the servants did strike him with the palms of their hands. And as Peter was beneath in the palace, there cometh one of the maids of the high priest: and when she saw Peter warming himself, she looked upon him, and said, And thou also wast with Jesus of Nazareth. But he denied, saying, I know not, neither understand I what thou sayest. And he went out into the porch; and the cock crew. And a maid saw him again, and began to say to them that stood by, This is one of them. And he denied it again. And a little after, they that stood by said again to Peter, Surely thou art one of them: for thou art a Galilean, and thy speech agreeth thereto. But he began to curse and to swear, saying, I know not this man of whom ye speak. And the second time the cock crew. And Peter called to mind the word that Jesus said unto him, Before the cock crow twice, thou shalt deny me thrice. And when he thought thereon, he wept.

\offertory{Deliver me, O Lord, from mine enemies: for I flee unto thee to hide me, teach me to do the thing that pleaseth thee: for thou art my God.}

\secret
\lett{A}{lmighty} God, let these sacrifices cleanse us by their effectual power, that we may thereby attain in purity unto the author of the same. Through.

\begin{rubric}
    \nth{2} Secret is \emph{Against the Persecutors of the Church} (p. \SPAgainst) or \emph{for the Chief Bishop} (p. \SPChiefBishop).
\end{rubric}

\vspace{-1ex}

\communion{Let them be put to confusion and shame together, that rejoice at my trouble: let them be clothed with rebuke and dishonour that boast themselves against me.}

\vspace{-1ex}

\postcommunion
\lett{M}{ay} thy holy mysteries, O Lord, impart unto us heavenly fervour: that we may delight in their performance, and likewise in bringing forth the fruit thereof. Through.

\begin{rubric}
    \nth{2} Postcommunion is \emph{Against the Persecutors of the Church} (p. \SPAgainst) or \emph{for the Chief Bishop} (p. \SPChiefBishop).
\end{rubric}

\textsc{Priest.} Let us pray.\par
\textsc{Deacon.} Humble your heads before God.\par
\lett{H}{elp} us, O God of our salvation: and grant, that we may come with joy to the celebration of those benefits by which thou hast vouchsafed to redeem us. Through.


\subby{Holy Tuesday}
\fancyhead[RO,LE]{\textit{Holy Tuesday}}
\fancyhead[RE,LO]{}
\begin{inhead}
    {First Class Feria}
\end{inhead}

	\properantiphon{Ben.}{Now before the feast {\dag} of the Passover, when Jesus knew that his hour was come, having loved his own, he loved them unto the end.}

    \properantiphon{Mag.}{I have power {\dag} to lay down my life, and I have power to take it again.}

\vspace{-1ex}

\introit
\lett{B}{ut} it behoveth us to glory in the Cross of our Lord Jesus Christ: in whom is our salvation, life, and resurrection: by whom we are saved and set free. \textit{Ps.} God be merciful unto us, and bless us: and shew us the light of his countenance, and be merciful unto us.

\collect
	\lett{O}{Lord} God, whose blessed Son, our Saviour, gave his back to the smiters and hid not his face from shame; Grant us grace to take joyfully the sufferings of the present time, in full assurance of the glory that shall be revealed. Through the same.
	
	\begin{rubric}
    \nth{2} Collect is \emph{Against the Persecutors of the Church} (p. \SPAgainst) or \emph{for the Chief Bishop} (p. \SPChiefBishop).
	\end{rubric}
	
	\readingcitation{Epistle}{Isaiah 50:5}
	\lett{I}{n those days:} Said Isaiah: The Lord \divineName{God} hath opened mine ear, and I was not rebellious, neither turned away back. I gave my back to the smiters, and my cheeks to them that plucked off the hair: I hid not my face from shame and spitting. For the Lord God will help me; therefore shall I not be confounded: therefore have I set my face like a flint and I know that I shall not be ashamed. He is near that justifieth me; who will contend with me? let us stand together: who is mine adversary? let him come near to me. Behold, the Lord God will help me; who is he that shall condemn me? lo, they all shall wax old as a garment; the moth shall eat them up. Who is among you that feareth the Lord that obeyeth the voice of his servant, that walketh in darkness, and hath no light? let him trust in the name of the Lord, and stay upon his God. Behold, all ye that kindle a fire, that compass yourselves about with sparks: walk in the light of your fire, and in the sparks that ye have kindled. This shall ye have of mine hand; ye shall lie down in sorrow.
	
	\gradual{Nevertheless, when they were sick, I put on sackcloth, and humbled my soul with fasting: and my prayer shall turn into mine own bosom. ℣. Plead thou my cause, O Lord, with them that strive with me, and fight thou against them that fight against me: lay hand upon the shield and buckler, and stand up to help me.}
	
	\readingcitation{Gospel}{Mark 15:1}
	\lett{A}{t that time:} Straightway in the morning the chief priests held a consultation with the elders and scribes and the whole council, and bound Jesus, and carried him away, and delivered him to Pilate. And Pilate asked him, Art thou the King of the Jews? And he answering said unto him, Thou sayest it. And the chief priests accused him of many things: but he answered nothing. And Pilate asked him again, saying, Answerest thou nothing? behold how many things they witness against thee. But Jesus yet answered nothing; so that Pilate marvelled. Now at that feast he released unto them one prisoner, whomsoever they desired. And there was one named Barabbas, which lay bound with them that had made insurrection with him, who had committed murder in the insurrection. And the multitude crying aloud began to desire him to do as he had ever done unto them. But Pilate answered them, saying, Will ye that I release unto you the King of the Jews? For he knew that the chief priests had delivered him for envy. But the chief priests moved the people, that he should rather release Barabbas unto them. And Pilate answered and said again unto them, What will ye then that I shall do unto him whom ye call the King of the Jews? And they cried out again, Crucify him. Then Pilate said unto them, Why, what evil hath he done? And they cried out the more exceedingly, Crucify him, And so Pilate, willing to content the people, released Barabbas unto them, and delivered Jesus, when he had scourged him, to be crucified. And the soldiers led him away into the hall, called Pr{\ae}torium; and. they called together the whole band. And they clothed him with purple, and platted a crown of thorns, and put it about his head, and began to salute him, Hail, King of the Jews! And they smote him on the head with a reed, and. did spit upon him, and bowing their knees worshipped him. And when they had mocked. him, they took off the purple from him, and put his own clothes on him, and led him out to. crucify him. And they compel one Simon a Cyrenian, who passed by, coming out of the country, the father of Alexander and Rufus, to bear his cross. And they bring him unto the place Golgotha, which is, being interpreted, The place of a skull. And they gave him to drink wine mingled with myrrh: but he received it not. And when they had crucified him, they parted his garments, casting lots upon them, what every man should take. And it was the third hour, and they crucified him. And the superscription of his accusation was written over, \divineName{the King of the Jews}. And with him they crucify two thieves; the one on his right hand, and the other on his left. And the scripture was fulfilled, which saith, And he was numbered with the transgressors. And they that passed by railed on him, wagging their heads, and saying, Ah, thou that destroyest the temple, and buildest it in three days, save thy self, and come down from the cross. Likewise also the chief priests mocking said among them selves with the scribes, He saved others; himself he cannot save. Let Christ the King of Israel descend now from the cross, that we may see and believe. And they that were crucified with him reviled him. And when the sixth hour was come, there was darkness over the whole land until the ninth hour. And at the ninth hour Jesus cried with a loud voice, saying, Eloi, Eloi, lama sabachthani? which is, being interpreted, My God, my God, why hast thou forsaken me? And some of them that stood by, when they heard it, said, Behold, he calleth Elias. And one ran and filled a spunge full of vinegar, and put it on a reed, and gave him to drink, saying, Let alone; let us see whether Elias will come to take him down. And Jesus cried with a loud voice, and gave up the ghost. \inrub{Here genuflect, and pause awhile.} And the veil of the temple was rent in twain from the top to the bottom. And when the centurion, which stood over against him, saw that he so cried out, and gave up the ghost, he said, Truly this man was the Son of God.
	
	\offertory{Keep me, O Lord, from the hands of the ungodly: preserve me from the wicked men.}
	
	\secret
	\lett{W}{e} beseech thee, O Lord: that these sacrifices which we offer unto thee in this time of fasting, ordained for the healing of our souls, may effectually renew us. Through.

	\begin{rubric}
    \nth{2} Secret is \emph{Against the Persecutors of the Church} (p. \SPAgainst) or \emph{for the Chief Bishop} (p. \SPChiefBishop).
	\end{rubric}
	
	\communion{They that sit in the gate speak against me: and the drunkards make songs upon me: but, Lord, I make my prayer unto thee: in an acceptable time, O God, in the multitude of thy mercy.}
	
	\postcommunion
	\lett{A}{lmighty} God, let thy holy mysteries both heal our vices and bestow upon us everlasting remedies. Through.
	
	\begin{rubric}
    \nth{2} Postcommunion is \emph{Against the Persecutors of the Church} (p. \SPAgainst) or \emph{for the Chief Bishop} (p. \SPChiefBishop).
	\end{rubric}
	
	\textsc{Priest.} Let us pray.\par
	\textsc{Deacon.} Humble your heads before God.\par
	\lett{L}{et} thy merciful kindness, O God, cleanse us throughly from all the corruption of our former nature, and make us to be effectually renewed unto holiness of life. Through.


\subby{Spy Wednesday}
\fancyhead[RO,LE]{\textit{Spy Wednesday}}
\fancyhead[RE,LO]{}
\begin{inhead}
    {First Class Double}
\end{inhead}

	\properantiphon{Ben.}{Simon, sleepest thou? {\dag} couldest not thou watch with me one hour?}

    \properantiphon{Mag.}{A maid said unto Peter, {\dag} Surely thou art one of them, for thy speech betrayeth thee.}

\introit
\lett{A}{t} the name of Jesus every knee should bow, of things in heaven, and things in earth, and things under the earth: for the Lord became obedient unto death, even the death of the cross: wherefore Jesus Christ is Lord, to the glory of God the Father. \textit{Ps.} Hear my prayer, O Lord: and let my crying come unto thee.

\begin{rubric}
    After the Kyrie,
\end{rubric}
\elcol{℣. Let us pray.}{℣. Orémus.}
\elcol{℣. Let us bow the knee.}{℣. Flectámus génua.}
\elcol{℣. Arise.}{℣. Leváte.}

\collect
\lett{G}{rant,} we beseech thee, almighty God: that we, who are continually afflicted by reason of our transgressions, may be delivered by the passion of thine only-begotten Son. Who liveth.

\readingcitation{Epistle}{Isaiah 62:11}
\lett{T}{hus saith the Lord God:} Say ye to the daughter of Zion, Behold, thy salvation cometh; behold, his reward is with him. Who is this that cometh from Edom, with dyed garments from Bozrah? this that is glorious in his apparel, travelling in the greatness of his strength? I that speak in righteousness, mighty to save. Wherefore art thou red in thine apparel, and thy garments like him that treadeth in the winefat? I have trodden the winepress alone; and of the people there was none with me: for I will tread them in mine anger, and trample them in my fury; and their blood shall be sprinkled upon my garments, and I will stain all my raiment. For the day of vengeance is in mine heart, and the year of my redeemed is come. And I looked, and there was none to help; and I wondered that there was none to uphold: therefore mine own arm brought salvation unto me; and my fury, it upheld me. And I will tread down the people in mine anger, and make them drunk in my fury, and I will bring down their strength to the earth. I will mention the lovingkindnesses of the \divineName{Lord}, and the praises of the \divineName{Lord}, according to all that the \divineName{Lord} hath bestowed on us.

\gradual{Hide not thy face from thy servant, for I am in trouble: O haste thee, and hear me. ℣. Save me, O God, for the waters are come in, even unto my soul: I stick fast in the deep mire, where no ground is.}

\elcol{℣. The Lord be with you.}{℣. Dóminus vobíscum.}
\elcol{℟. And with thy spirit.}{℟. Et cum spíritu tuo.}
\elcol{℣. Let us pray.}{℣. Orémus.}

\collect
\lett{A}{ssist} us mercifully with thy help, O Lord God of our salvation; that we may enter with joy upon the meditation of those mighty acts, whereby thou hast given unto us life and immortality. Through.
	
	\begin{rubric}
    \nth{2} Collect is \emph{Against the Persecutors of the Church} (p. \SPAgainst) or \emph{for the Chief Bishop} (p. \SPChiefBishop).
	\end{rubric}

\readingcitation{Epistle}{Hebrews 9:16}
\lett{B}{rethren:} Where a testament is, there must also of necessity be the death of the testator. For a testament is of force after men are dead: otherwise it is of no strength at all while the testator liveth. Whereupon neither the first testament was dedicated without blood. For when Moses had spoken every precept to all the people according to the law, he took the blood of calves and of goats, with water, and. scarlet wool, and hyssop, and sprinkled both the book and all the people, saying, This is the blood of the testament which God hath enjoined unto you. Moreover, he sprinkled likewise with blood both the tabernacle, and all the vessels of the ministry. And almost all things are by the law purged with blood; and without shedding of blood is no remission. It was therefore necessary that the patterns of things in the heavens should be purified with these; but the heavenly things themselves with better sacrifices than these. For Christ is not entered into the holy places made with hands, which are the figures of the true; but into heaven itself, now to appear in the presence of God for us: nor yet that he should offer himself often, as the high priest entereth into the holy place every year with blood of others; for then must he often have suffered since the foundation of the world: but now once in the end of the world hath he appeared to put away sin by the sacrifice of himself. And as it is appointed unto men once to die, but after this the judgment: so Christ was once offered to bear the sins of many; and unto them that look for him shall he appear the second time without sin unto salvation.

\tract{Hear my prayer, O Lord, and let my crying come unto thee. ℣. Hide not thy face from me in the time of my trouble. ℣. Incline thine ear unto me when I call, O hear me, and that right soon. ℣. For my days are consumed away like smoke: and my bones are burnt up as it were a fire-brand. ℣. My heart is smitten down, and withered like grass: so that I forget to eat my bread. ℣. Thou shalt arise, O Lord, and have mercy upon Sion: for it is time that thou have mercy upon her.}

\readingcitation{Gospel}{Luke 22:1}
\lett{A}{t that time:} The feast of unleavened bread drew nigh, which is called the Passover. And the chief priests and scribes sought how they might kill him; for they feared the people. Then entered Satan into Judas surnamed Iscariot, being of the number of the twelve. And he went his way, and communed with the chief priests and captains, how he might betray him unto them. And they were glad, and covenanted to give him money. And he promised, and sought opportunity to betray him unto them in the absence of the multitude. Then came the day of unleavened bread, when the passover must be killed. And he sent Peter and John, saying, Go and prepare us the passover, that we may eat. And they said unto him, Where wilt thou that we prepare? And he said unto them, Behold, when ye are entered into the city, there shall a man meet you, bearing a pitcher of water; follow him into the house where he entereth in. And ye shall say unto the goodman of the house, The Master saith unto thee, Where is the guest-chamber, where I shall eat the passover with my disciples? And he shall shew you a large upper room furnished: there make ready. And they went, and found as he had said unto them: and they made ready the passover. And when the hour was come, he sat down, and the twelve apostles with him. And he said unto them, With desire I have desired to eat this passover with you before I suffer: for I say unto you, I will not any more eat thereof, until it be fulfilled in the kingdom of God. And he took the cup, and gave thanks, and said, Take this, and divide it among yourselves: for I say unto you I will not drink of the fruit of the vine, until the kingdom of God shall come. And he took bread, and gave thanks, and brake it, and gave unto them, saying, This is my body which is given for you: this do in remembrance of me. Likewise also the cup after supper, saying, This cup is the new testament in my blood, which is shed for you. But, behold, the hand of him that betrayeth me is with me on the table. And truly the Son of man goeth, as it was determined: but woe unto that man by whom he is betrayed! And they began to enquire among themselves, which of them it was that should do this thing. And there was also a strife among them, which of them should be accounted the greatest And he said unto them, The kings of the Gentiles exercise lordship over them; and they that exercise authority upon them are called benefactors. But ye shall not be so: but he that is greatest among you, let him be as the younger; and he that is chief, as he that doth serve. For whether is greater, he that sitteth at meat, or he that serveth? is not he that sitteth at meat? but I am among you as he that serveth. Ye are they which have continued with me in my temptations. And I appoint unto you a kingdom, as my Father hath appointed unto me; that ye may eat and drink at my table in my kingdom. and sit on thrones, judging the twelve tribes of Israel. And the Lord said, Simon, Simon, behold, Satan hath desired to have you, that he may sift you as wheat: but I have prayed for thee, that thy faith fail not; and when thou art converted, strengthen thy brethren. And he said unto him, Lord, I am ready to go with thee, both into prison, and to death. And he said, I tell thee, Peter, the cock shall not crow this day, before that thou shalt thrice deny that thou knowest me. And he said unto them, When I sent you without purse, and scrip, and shoes, lacked ye any thing? And they said, Nothing. Then said he unto them, But now, he that hath a purse, let him take it, and likewise his scrip: and he that hath no sword, let him sell his garment, and buy one. For I say unto you, that this that is written must yet be accomplished in me, And he was reckoned among the transgressors: for the things concerning me have an end. And they said, Lord, behold, here are two swords. And he said unto them, It is enough. And he came out, and went, as he was wont, to the mount of Olives; and his disciples also followed him. And when he was at the place, he said unto them, Pray that ye enter not into temptation. And he was withdrawn from them about a stone's cast, and kneeled down, and prayed, saying, Father, if thou be willing, remove this cup from me: nevertheless not my will, but thine, be done. And there appeared an angel unto him from heaven, strengthening him. And being in an agony he prayed more earnestly: and his sweat was as it were great drops of blood falling down to the ground. And when he rose up from prayer, and was come to his disciples, he found them sleeping for sorrow, and said unto them, Why sleep ye? rise and pray, lest ye enter into temptation. And while he yet spake behold a multitude, and he that was called Judas, one of the twelve, went before them, and drew near unto Jesus to kiss him. But Jesus said unto him. Judas, betrayest thou the Son of man with a kiss? When they which were about him saw what would follow, they said unto him, Lord, shall we smite with the sword? And one of them smote the servant of the high priest, and cut off his right ear. And Jesus answered and said, Suffer ye thus far. And he touched his ear, and healed him. Then Jesus said unto the chief priests, and captains of the temple, and the elders, which were come to him, Be ye come out, as against a thief, with swords and staves? When I was daily with you in the temple, ye stretched forth no hands against me: but this is your hour, and the power of darkness. Then took they him, and led him, and brought him into the high priest's house. And Peter followed afar off. And when they had kindled a fire in the midst of the hall, and were set down together, Peter sat down among them. But a certain maid beheld him as he sat by the fire, and earnestly looked upon him, and said, This man was also with him. And he denied him, saying, Woman, I know him not. And after a little while another saw him, and said, Thou art also of them. And Peter said, Man, I am not. And about the space of one hour after another confidently affirmed, saying, Of a truth this fellow also was with him: for he is a Galilean. And Peter said, Man, I know not what thou sayest. And immediately, while he yet spake, the cock crew. And the Lord turned, and looked upon Peter. And Peter remembered the word of the Lord, how he had said unto him, Before the cock crow, thou shalt deny me thrice. And Peter went out, and wept bitterly. And the men that held Jesus mocked him, and smote him. And when they had blindfolded him, they struck him on the face, and asked him, saying, Prophesy, who is it that smote thee? And many other things blasphemously spake they against him. And as soon as it was day the elders of the people and the chief priests and the scribes came together, and led him into their council, saying, Art thou the Christ? tell us. And he said unto them, If I tell you, ye will not believe: and if I also ask you, ye will not answer me, nor let me go. Hereafter shall the Son of man sit on the right hand of the power of God. Then said they all, Art thou then the Son of God? And he said unto them, Ye say that I am. And they said, What need we any further witness? for we ourselves have heard of his own mouth.

\offertory{Hear my prayer, O Lord, and let my crying come unto thee: hide not thy face from me.}

\secret
\lett{A}{ccept,} we beseech thee, O Lord, the gift which we offer, and vouchsafe so to work in us: that as we set forth in a mystery the passion of thy Son our Lord, so by devout affection we may attain unto the same. Through the same.

	\begin{rubric}
    \nth{2} Secret is \emph{Against the Persecutors of the Church} (p. \SPAgainst) or \emph{for the Chief Bishop} (p. \SPChiefBishop).
	\end{rubric}

\communion{I have mingled my drink with weeping: for thou hast taken me up, and cast me down: and I am withered like grass: but thou, O Lord, shalt endure for ever: thou shalt arise, and have mercy upon Sion, for it is time that thou have mercy upon her.}

\postcommunion
\lett{A}{lmighty} God, enlighten our understanding: that, as these wondrous mysteries do testify to the temporal death of thy Son, so thereby we may have assurance that thou hast bestowed upon us life eternal. Through the same.

	\begin{rubric}
    \nth{2} Postcommunion is \emph{Against the Persecutors of the Church} (p. \SPAgainst) or \emph{for the Chief Bishop} (p. \SPChiefBishop).
	\end{rubric}
	
	\textsc{Priest.} Let us pray.\par
	\textsc{Deacon.} Humble your heads before God.\par
	\lett{A}{lmighty} God, we beseech thee graciously to behold this thy family, for which our Lord Jesus Christ was contented to be betrayed, and given up into the hands of wicked men, and to suffer death upon the Cross. Who liveth.
	

\subby{Maundy Thursday}
\fancyhead[RO,LE]{\textit{Maundy Thursday}}
\fancyhead[RE,LO]{}
\begin{inhead}
    {First Class Double}
\end{inhead}

    \properantiphon{Ben.}{Now he that betrayed him {\dag} gave them a sign, saying: Whomsoever I shall kiss, that same is he; hold him fast.}

    \properantiphon{Mag.}{As they were eating, {\dag} Jesus took bread, and blessed it, and brake it, and gave it to his disciples.}
    
    \introit
    \lett{B}{ut} it behoveth us to glory in the Cross of our Lord Jesus Christ: in whom is our salvation, life, and resurrection: by whom we are saved, and set free. \textit{Ps.} God be merciful unto us, and bless us: and shew us the light of his countenance, and be merciful unto us.
    
    \begin{rubric}
    	\emph{Gloria in excelsis} is said, and after the intonation thereof the bells are rung: and, the Hymn being ended, they are silent until Holy Saturday, as is there noted.
    \end{rubric}
    
    \collect
    \lett{A}{lmighty} Father, whose dear Son, on the night before he suffered, did institute the Sacrament of his Body and Blood; Mercifully grant that we may thankfully receive the same in remembrance of him, who in these holy mysteries giveth us a pledge of life eternal; through the same thy Son Jesus Christ our Lord, who now liveth and reigneth with thee and the Holy Spirit ever, one God, world without end. \textit{Amen.}
    
    \readingcitation{Epistle}{1 Corinthians 11:17}
    \lett{B}{rethren:} In this that I declare unto you I praise you not, that ye come together not for the better, but for the worse. For first of all, when ye come together in the church, I hear that there be divisions among you; and I partly believe it. For there must be also heresies among you, that they which are approved may be made manifest among you. When ye come together therefore into one place, this is not to eat the Lord's supper. For in eating every one taketh before other his own supper: and one is hungry, and another is drunken. What? have ye not houses to eat and to drink in? or despise ye the church of God, and shame them that have not? What shall I say to you? shall I praise you in this? I praise you not. For I have received of the Lord that which also I delivered unto you, That the Lord Jesus the same night in which he was betrayed took bread: and when he had given thanks, he brake it, and said, Take, eat: this is my body, which is broken for you: this do in remembrance of me After the same manner also he took the cup, when he had supped, saying, This cup is the new testament in my blood: this do ye, as oft as ye drink it, in remembrance of me. For as often as ye eat this bread, and drink this cup, ye do shew the Lord's death till he come. Wherefore whosoever shall eat this bread, and drink this cup of the Lord, unworthily, shall be guilty of the body and blood of the Lord. But let a man examine himself, and so let him eat of that bread, and drink of that cup. For he that eateth and drinketh unworthily, eateth and drinketh damnation to himself, not discerning the Lord's body. For this cause many are weak and sickly among you, and many sleep. For if we would judge ourselves, we should not be judged. But when we are judged, we are chastened of the Lord, that we should not be condemned with the world. Wherefore, my brethren, when ye come together to eat, tarry one for another. And if any man hunger, let him eat at home; that ye come not together unto condemnation. And the rest will I set in order when I come.
    
    \gradual{Christ for us became obedient unto death, even the death of the cross. ℣. Wherefore God also hath highly exalted him: and given him a name which is above every name.}
    
    \readingcitation{Gospel}{Luke 23:1}
    \lett{A}{t that time:} The whole multitude of them arose, and led him unto Pilate. And they began to accuse him, saying, We found this fellow perverting the nation, and forbidding to give tribute to Cæsar, saying that he himself is Christ a King. And Pilate asked him, saying, Art thou the King of the Jews? And he answered him and said, Thou sayest it. Then said Pilate to the chief priests and to the people, I find no fault in this man. And they were the more fierce, saying, He stirreth up the people, teaching through-out all Jewry, beginning from Galilee to this place. When Pilate heard of Galilee, he asked whether the man were a Galilean. And as soon as he knew that he belonged unto Herod's jurisdiction, he sent him to Herod, who himself also was at Jerusalem at that time. And when Herod saw Jesus, he was exceeding glad: for he was desirous to see him of a long season, because he had heard many things of him; and he hoped to have seen some miracle done by him. Then he questioned with him in many words; but he answered him nothing. And the chief priests and scribes stood and vehemently accused him. And Herod with his men of war set him at nought, and mocked him, and arrayed him in a gorgeous robe, and sent him again to Pilate. And the same day Pilate and Herod were made friends together: for before they were at enmity between themselves. And Pilate, when he had called together the chief priests and the rulers and the people, said unto them, Ye have brought this man unto me, as one that perverted the people: and, behold, I, having examined him before you, have found no fault in this man touching those things whereof ye accuse him: no, nor yet Herod: for I sent you to him; and, lo, nothing worthy of death is done unto him. I will therefore chastise him, and release him. (For of necessity he must release one unto them at the feast.) And they cried out all at once, saying, Away with this man, and release unto us Barabbas: (who for a certain sedition made in the city, and for murder, was cast into prison.) Pilate therefore, willing to release Jesus, spake again to them. But they cried, saying, Crucify him, crucify him. And he said unto them the third time, Why, what evil hath he done? I have found no cause of death in him: I will therefore chastise him, and let him go. And they were instant with loud voices, requiring that he might be crucified. And the voices of them and of the chief priests prevailed. And Pilate gave sentence that it should be as they required. And he released unto them him that for sedition and murder was cast into prison, whom they had desired; but he delivered Jesus to their will. And as they led him away, they laid hold upon one Simon, a Cyrenian, coming out of the country, and on him they laid the cross, that he might bear it after Jesus. And there followed him a great company of people, and of women, which also bewailed and lamented him. But Jesus turning unto them said, Daughters of Jerusalem, weep not for me, but weep for yourselves, and for your children. For, behold, the days are coming, in the which they shall say, Blessed are the barren, and the wombs that never bare, and the paps which never gave suck. Then shall they begin to say to the mountains, Fall on us; and to the hills, Cover us. For if they do these things in a green tree, what shall be done in the dry? And there were also two other, malefactors, led with him to be put to death. And when they were come to the place which is called Calvary, there they crucified him, and the malefactors, one on the right hand, and the other on the left. Then said Jesus, Father, forgive them; for they know not what they do. And they parted his raiment, and cast lots. And the people stood beholding. And the rulers also with them derided him, saying, He saved others; let him save himself, if he be Christ, the chosen of God. And the soldiers also mocked him, coming to him, and offering him vinegar, and saying, If thou be the king of the Jews, save thyself. And a superscription also was written over him, in letters of Greek, and Latin, and Hebrew, \textsc{This is the King of the Jews}. And one of the malefactors which were hanged railed on him, saying, If thou be Christ, save thyself and us. But the other answering rebuked him, saying, Dost not thou fear God, seeing thou art in the same condemnation? And we indeed justly; for we receive the due reward of our deeds: but this man hath done nothing amiss. And he said unto Jesus, Lord, remember me when thou comest into thy kingdom. And Jesus said unto him, Verily I say unto thee, To day shalt thou be with me in paradise. And it was about the sixth hour, and there was a darkness over all the earth until the ninth hour. And the sun was darkened, and the veil of the temple was rent in the midst. And when Jesus had cried with a loud voice, he said, Father, into thy hands I commend my spirit: and having said thus, he gave up the ghost. \inrub{Here genuflect, and pause awhile.} Now when the centurion saw what was done, he glorified God, saying, Certainly this was a righteous man. And all the people that came together to that sight, beholding the things which were done, smote their breasts, and returned. And all his acquaintance, and the women that followed him from Galilee, stood afar off, beholding these things.
    
    \offertory{The right hand of the Lord bringeth mighty things to pass, the right of the Lord hath exalted me: I shall not die, but live, and declare the works of the Lord.}
    
    %MANUAL ADJUSTMENT:
    \clearpage
    \secret
    \lett{W}{e} beseech thee, O Lord holy, Father almighty, everlasting God: that he may render our sacrifice acceptable unto thee, who on this day commanded and taught his disciples to do this in remembrance of him, even Jesus Christ thy Son our Lord. Who liveth.
    
    \begin{rubric}
    	After the Comfortable Words, the \emph{Pax} is not given.
    \end{rubric}
    
    \begin{rubric}
    	\textsc{Note,} To-day the Priest consecrates two Hosts, of which he consumes one, and reserves the other (touched with the Blood) for the following day, on which the Sacrament is not consecrated; he also reserves some consecrated particles, in case of need, for the sick.\par
    	But be consumes the whole of the Blood, and before washing his fingers he places the reserved Host in another Chalice, which the Deacon covers with the Pall and the Paten inverted, and spreads a Veil over it, and places it in the midst of the Altar.\par
    	Then the Communion is said and the Mass finished. The Priest genuflects, however, whenever he approaches or departs from the middle of the Altar, or crosses before the Sacrament reserved in the Chalice: and when he is to say \emph{The Lord be with you}, he does not turn to the people in the middle of the Altar, lest be turn his back on the Sacrament, but from the Gospel side: and at the end he gives the blessing from the same place and does not complete the circle.
    \end{rubric}
    
    \communion{The Lord Jesus, after he had supped with his disciples, washed their feet, and said to them: Know ye what I, your Lord and Master, have done to you? I have given you an example, that ye also should so do.}
    
    \postcommunion
    \lett{O}{Lord} our God, who hast refreshed us with life-giving sustenance, we beseech thee: that we who observe this institution in the time of our mortal life, may obtain the benefit thereof in thy gift of immortality. Through.
    
    \begin{rubric}
    	\emph{Ite, Missa est} is said, and the Blessing is given, and the Gospel of St. John is read, at the beginning of which the Priest does not sign the Altar, but himself only.
    \end{rubric}
    
    \subby{Altar of Repose \& Stripping of the Altar}
    
    \begin{rubric}
    	To-day a suitable place is to be prepared in some Chapel of the Church, or Altar: and adorned as decently as may be, with veils and lights: where the Chalice with the Host, reserved as above, shall be set. When the Mass is over, torches are lighted, and a Procession is formed in the usual manner, but with another Subdeacon vested carrying the Cross. The Celebrant, vested in a white Cope, standing before the Altar, places incense in two thuribles without a blessing: then kneeling in the middle, with one he censes the Sacrament thrice: and receiving the Chalice with the Sacrament from the hand of the Deacon, who stands, and having it covered with the ends of the Veil with which his shoulders are covered, he proceeds between the same Deacon on his right, and the Subdeacon on his left, under a baldachino with two Acolytes continually censing the Sacrament, to the place prepared, where it is to be kept for the morrow. Meanwhile, during the Procession, the following hymn is sung.
    \end{rubric}
    
    \subbysub{Pange Lingua}
    
    \gregorioscore{resources/gabc/PangeLingua.gabc}
    
    
    \begin{rubric}
    	When they arrive at the place prepared, the Deacon kneels down and receives from the Priest, who stands, the Chalice with the Sacrament and places it first upon the Altar, where it is censed by the Priest kneeling, as above: then be places it in the receptacle.\par
    	Afterwards Evensong is said in the Choir without singing. Which being finished, the Priest, vested in an Alb and violet stole, with the Ministers shall lay bare the Altars, reading the following:
    \end{rubric}

    \subbysub{Psalm 22. \textit{Deus, Deus meus}}
    
        \antiphon{}{They parted my garments among them: and upon my vesture did they cast lots.}

\lett{M}{y} God, my God, look upon me; why hast thou forsaken me : and art so far from my health, and from the words of my complaint?\par
\secondline{\psanum{2}O my God, I cry in the day-time, but thou hearest not : and in the night-season also I take no rest.}
\thirdline{\psanum{3}And thou continuest holy : O thou worship of Israel.}
\psanum{4}Our fathers hoped in thee : they trusted in thee, and thou didst deliver them.\par
\psanum{5}They called upon thee, and were holpen : they put their trust in thee, and were not confounded.\par
\psanum{6}But as for me, I am a worm, and no man : a very scorn of men, and the outcast of the people.\par
\psanum{7}All they that see me laugh me to scorn : they shoot our their lips, and shake their heads, saying,\par
\psanum{8}He trusted in God, that he would deliver him : let him deliver him, if he will have him.\par
\psanum{9}But thou art he that took me out of my mother's womb : thou wast my hope, when I hanged yet upon my mother's breasts.\par
\psanum{10}I have been left unto thee ever since I was born : thou art my God, even from my mother's womb.\par
\psanum{11}O go not from me, for trouble is hard at hand : and there is none to help me.\par
\psanum{12}Many oxen are come about me : fat bulls of Basan close me in on every side.\par
\psanum{13}They gape upon me with their mouths : as it were a ramping and a roaring lion.\par
\psanum{14}I am poured out like water, and all my bones are out of joint : my heart also in the midst of my body is even like melting wax.\par
\psanum{15}My strength is dried up like a potsherd, and my tongue cleaveth to my gums : and thou shalt bring me into the dust of death.\par
\psanum{16}For many dogs are come about me : and the council of the wicked layeth siege against me.\par
\psanum{17}They pierced my hands and my feet; I may tell all my bones : they stand staring and looking upon me.\par
\psanum{18}They part my garments among them : and casts lots upon my vesture.\par
\psanum{19}But be not thou far from me, O Lord : thou art my succour, haste thee to help me.\par
\psanum{20}Deliver my soul from the sword : my darling from the power of the dog.\par
\psanum{21}Save me from the lion's mouth : thou hast heard me also from among the horns of the unicorns.\par
\psanum{22}I will declare thy Name unto my brethren : in the midst of the congregation will I praise thee.\par
\psanum{23}O praise the Lord, ye that fear him : magnify him, all ye of the seed of Jacob, and fear him, all ye seed of Israel.\par
\psanum{24}For he hath not despised, nor abhorred, the low estate of the poor : he hath not hid his face from him, but when he called unto him he heard him.\par
\psanum{25}My praise is of thee in the great congregation : my vows will I perform in the sight of them that fear him.\par
\psanum{26}The poor shall eat and be satisfied : they that seek after the Lord shall praise him; your heart shall live for ever.\par
\psanum{27}All the ends of the world shall remember themselves, and be turned unto the Lord : and all the kindreds of the nations shall worship before him.\par
\psanum{28}For the kingdom is the Lord's : and he is the Governor among the people.\par
\psanum{29}All such as be fat upon earth : have eaten and worshipped.\par
\psanum{30}All they that go down into the dust shall kneel before him : and no man hath quickened his own soul.\par
\psanum{31}My seed shall serve him : they shall be counted unto the Lord for a generation.\par
\psanum{32}They shall come, and the heavens shall declare his righteousness : unto a people that shall be born, whom the Lord hath made.\par

    \antiphon{}{They parted my garments among them: and upon my vesture did they cast lots.}
    
    \begin{rubric}
    	After the stripping of the Altars, at the appointed hour a signal is given with a board, and the Clergy come together to perform the Maundy.\par
    	The Prelate or Superior is vested with Stole and violet Cope over Amice and Alb, and in the appointed place with the assistance of the Deacon (who with the Subdeacon assists him in white vestments, as at Mass) puts incense in the thurible: then the Deacon holding the book of the Gospels before his breast, kneels before the Superior and asks a blessing: on receiving this, while two Acolytes with lighted candles stand by, and the Subdeacon holds the book, he signs it and censes it, and sings as usual the following Gospel.
    	\end{rubric}
    
    \readingcitation{Gospel}{John 13:1}
    \lett{B}{efore} the feast of the passover, when Jesus knew that his hour was come that he should depart out of this world unto the Father, having loved his own which were in the world, he loved them unto the end. And supper being ended, the devil having now put into the heart of Judas Iscariot, Simon's son, to betray him; Jesus knowing that the Father had given all things into his hands, and that he was come from God, and went to God; he riseth from supper, and laid aside his garments; and took a towel, and girded himself. After that he poureth water into a bason, and began to wash the disciples' feet, and to wipe them with the towel wherewith he was girded. Then cometh he to Simon Peter: and Peter saith unto him, Lord, dost thou wash my feet? Jesus answered and said unto him, What I do thou knowest not now; but thou shalt know hereafter. Peter saith unto him, Thou shalt never wash my feet. Jesus answered him, If I wash thee not, thou hast no part with me. Simon Peter saith unto him, Lord, not my feet only, but also my hands and my head. Jesus saith to him, He that is washed needeth not save to wash his feet, but is clean every whit: and ye are clean, but not all. For he knew who should betray him; therefore said he, Ye are not all clean. So after he had washed their feet, and had taken his garments, and was set down again, he said unto them, Know ye what I have done to you? Ye call me Master and Lord: and ye say well; for so I am. If I then, your Lord and Master, have washed your feet; ye also ought to wash one another's feet. For I have given you an example, that ye should do as I have done to you.
    
    \begin{rubric}
    	When this is finished, the Subdeacon carries the open book to be kissed by the Superior, and the Deacon censes him as usual.\par
    	Then the Superior puts off the Cope, and is girded by the Deacon and Subdeacon with a linen cloth, and thus girded, the Deacon and Subdeacon assisting, proceeds to the washing of feet. Those who are to be washed are arranged in order, and while Clerks minister a bason and water, and the Subdeacon holds the right foot of each, the Superior kneels before each, washes his foot, wipes and kisses it, the Deacon giving the towel to wipe them. Meanwhile the following are sung:
    \end{rubric}
    
    \antiphon{}{A new commandment I give unto you: that ye love one another, as I have loved you, saith the Lord. \textit{Ps. 119.} Blessed are those that are undefiled in the way: and walk in the law of the Lord. ℣. A new commandment I give unto you: that ye love one another, as I have loved you, saith the Lord.}
    
    \antiphon{}{After the Lord had risen from supper, he poured water into a bason, and began to wash his disciples' feet: this example left he unto them. \textit{Ps. 48.} Great is the Lord, and highly to be praised: in the city of our God, even upon his holy hill. ℣. After the Lord had risen from supper, he poured water into a bason, and began to wash his disciples' feet: this example left he unto them.}
    
    \antiphon{}{The Lord Jesus, after he had supped with his disciples, washed their feet, and said to them: Know ye what I, your Lord and Master, have done to you? I have given you an example, that ye also should so do. \textit{Ps. 85.} Lord, thou art become gracious unto thy land: thou hast turned away the captivity of Jacob. ℣. The Lord Jesus, after he had supped with his disciples, washed their feet, and said to them: Know ye what I, your Lord and Master, have done to you? I have given you an example, that ye also should so do.}
    
    \antiphon{}{Lord, dost thou wash my feet? Jesus answered, and said unto him: If I wash not thy feet, thou hast no part with me. ℣. Then cometh he to Simon Peter, and Peter saith unto him.}
    
    \antiphon{}{Lord, dost thou wash my feet? Jesus answered, and said unto him: If I wash not thy feet, thou hast no part with me. ℣. What I do, thou knowest not now: but thou shalt know hereafter.}
    
    \antiphon{}{Lord, dost thou wash my feet? Jesus answered, and said unto him: If I wash not thy feet, thou hast no part with me. ℣. If I, your Lord and Master, have washed your feet: how much more ought ye to wash one another's feet? \textit{Ps. 49.} O hear ye this, all ye people: ponder it with your ears, all ye that dwell in the world. ℣. If I, your Lord and Master, have washed your feet: how much more ought ye to wash one another's feet?}
    
    \antiphon{}{By this shall all men know that ye are my disciples, if ye have love one to another. ℣. Jesus said to his disciples. ℣. By this shall all men know that ye are my disciples, if ye have love one to another.}
    
    \antiphon{}{Let there abide in you faith, hope, charity, these three: but the greatest of these is charity. ℣. And now abideth faith, hope, charity, these three, but the greatest of these is charity. ℣. Let there abide in you faith, hope, charity, these three: but the greatest of these is charity.}
    
    \antiphon{}{Blessed be the holy Trinity and the undivided Unity: we will give thanks unto him, because he hath shewed his mercy upon us. ℣. Let us bless the Father and the Son with the Holy Ghost. \textit{Ps. 84.} O how amiable are thy dwellings, thou Lord of hosts! My soul hath a desire and longing to enter into the courts of the Lord. ℣. Blessed be the holy Trinity and the undivided Unity: we will give thanks unto him, because he hath shewed his mercy upon us.}
    
    \antiphon{}{Where charity and love are, there is God. ℣. The love of Christ hath gathered us together in one. ℣. Let us rejoice and be glad in him. ℣. Let us fear and love the living God. ℣. And love one another in sincerity of heart.}
    
    \antiphon{}{Where charity and love are, there is God. ℣. When therefore we are gathered together in one: ℣. Take we heed that we be not divided in mind. ℣. Let malicious quarrels, let contentions cease. ℣. And let Christ our God be in the midst of us.}
    
    \antiphon{}{Where charity and love are, there is God. ℣. So may we also with the Blessed see ℣. In glory thy countenance, O Christ our God: ℣. Joy, that is infinite and undefiled, ℣. For ever and for evermore. Amen.}
    
    \begin{rubric}
    	After the washing, the Superior, or he that hath washed the feet, washes his hands, and wipes then with another towel: then returning to the place where he was at first he puts on the Cope, and standing with head uncovered says:
    \end{rubric}
    
    \begin{rubric}
    	\emph{Our Father}, secretly until,
    \end{rubric}
    
    ℣. And lead us not into temptation.
    
    ℟. But deliver us from evil.
    
    ℣. Thou hast charged, O Lord.
    
    ℟. That we shall diligently keep thy commandments.
    
    ℣. Thou didst wash the feet of thy disciples.
    
    ℟. Despise not thou the works of thine hands.
    
    ℣. O Lord, hear my prayer.
    
    ℟. And let my cry come unto thee.
    
    ℣. The Lord be with you.
    
    ℟. And with thy spirit.
    
    \letuspray
    
    \lett{A}{ssist} us, we beseech thee, O Lord, in this our bounden service: and since thou wast pleased to wash the feet of thy disciples, despise not the works of thine own hands which thou hast bidden us to follow: that as here outward stains by us and from us are washed away; so also the secret sins of us all may be blotted out by thee. Which do thou vouchsafe to grant, who livest and reignest God: throughout all ages, world without end. \textit{Amen.}

\subby{Good Friday}
\fancyhead[RO,LE]{\textit{Good Friday}}
\fancyhead[RE,LO]{}
\begin{inhead}
    {First Class Double}
\end{inhead}

\properantiphon{Ben.}{They set up {\dag} over his head his accusation written: Jesus of Nazareth, the King of the Jews.}

    \properantiphon{Mag.}{When he had received {\dag} the vinegar, he said, It is finished; and he bowed his head, and gave up the ghost.}

\begin{rubric}
	In Choir, None being ended, the Priest and Ministers in black vestments, without lights and incense, proceed to the Altar: and pray awhile prostrate before it. Meanwhile the Acolytes spread one cloth only upon the Altar. The Priest with the Ministers, having finished their prayer, go up to the Altar, and he kisses it in the middle: then the Reader proceeds to read the Prophecy in the place where the Epistle is read, and begins it without any announcement: the Priest meanwhile reads the same in a low voice at the Epistle corner.
\end{rubric}

\readingcitation{Lesson}{Hosea 6:1}
\lett{T}{hus saith the Lord:} In their affliction they will seek me early: Come, and let us return unto the \divineName{Lord}: for he hath torn, and he will heal us; he hath smitten, and he will bind us up. After two days will he revive us: in the third day he will raise us up, and we shall live in his sight. Then shall we know, if we follow on to know the \divineName{Lord}: his going forth is prepared as the morning; and he shall come unto us as the rain, as the latter and former rain unto the earth. O Ephraim, what shall I do unto thee? O Judah, what shall I do unto thee? for your goodness is as a morning cloud, and as the early dew it goeth away. Therefore have I hewed them by the prophets; I have slain them by the words of my mouth: and thy judgments are as the light that goeth forth. For I desired mercy, and not sacrifice; and the knowledge of God more than burnt offerings.

\begin{rubric}
	\emph{Thanks be to God} is not answered here, nor after the following Lesson.
\end{rubric}

\tract{O Lord, I have heard thy speech, and was afraid: I have considered thy works, and was confounded. ℣. O Lord, revive thy work in the midst of the years: in the midst of the years make it known. ℣. In the time of confusion of my soul: in wrath, remember mercy. ℣. God came from Teman, and his Holy One from the thick woods of the mountains. ℣. His glory covered the heavens: and the earth was full of his praise.}

\begin{rubric}
	The Tract being finished, the Priest at the Epistle corner says: \emph{Let us pray}. The Deacon: \emph{Let us bow the knee}, and the Subdeacon: \emph{Arise}.
\end{rubric}

\collect
\lett{O}{God,} from whom Judas received the punishment of his guilt, and the thief the reward of his confession, grant unto us the effects of thy propitiation: that as in his passion Jesus Christ, our Lord, gave unto both the divers rewards of their merits; so he may take away the transgressions of our old nature, and bestow upon us the grace of his resurrection. Who liveth.

\begin{rubric}
	The Subdeacon in the tone of the Epistle, likewise without announcement, sings the following Lesson:
\end{rubric}

\readingcitation{Lesson}{Hebrews 10:1}
\lett{B}{rethren:} The law having a shadow of good things to come, and not the very image of the things, can never with those sacrifices which they offered year by year continually make the corners thereunto perfect. For then would they not have ceased to be offered? because that the worshippers once purged should have had no more conscience of sins. But in those sacrifices there is a remembrance again made of sins every year. For it is not possible that the blood of bulls and of goats should take away sins. Wherefore when he cometh into the world, he saith, Sacrifice and offering thou wouldest not, but a body hast thou prepared me: in burnt offerings and sacrifices for sin thou hast had no pleasure. Then said I, Lo, I come (in the volume of the book it is written of me,) to do thy will, O God. Above when he said, Sacrifice and offering and burnt offerings and offering for sin thou wouldest not, neither hadst pleasure therein; which are offered by the law; then said he, Lo, I come to do thy will, O God. He taketh away the first, that he may establish the second. By the which will we are sanctified through the offering of the body of Jesus Christ once for all. And every priest standeth daily ministering and offering oftentimes the same sacrifices, which can never take away sins: but this man, after he had offered one sacrifice for sins for ever, sat down on the right hand of God; from hence-forth expecting till his enemies be made his footstool. For by one offering he hath perfected for ever them that are sanctified. Whereof the Holy Ghost also is a witness to us: for after that he had said before, This is the covenant that I will make with them after those days, saith the Lord, I will put my laws into their hearts, and in their minds will I write them; and their sins and iniquities will I remember no more. Now where remission of these is, there is no more offering for sin. Having therefore, brethren, boldness to enter into the holiest by the blood of Jesus, by a new and living way, which he hath consecrated for us, through the veil, that is to say, his flesh; and having an high priest over the house of God; let us draw near with a true heart in full assurance of faith, having our hearts sprinkled from an evil conscience, and our bodies washed with pure water. Let us hold fast the profession of our faith without wavering; (for he is faithful that promised;) and let us consider one another to provoke unto love and to good works: not forsaking the assembling of ourselves together, as the manner of some is; but exhorting one another: and so much the more, as ye see the day approaching.

\tract{Deliver me, O Lord, from the evil man: and preserve me from the wicked man. ℣. Who imagine mischief in their hearts: and stir up strife all the day long. ℣. They have sharpened their tongues like a serpent: adders' poison is under their lips. ℣. Keep me, O Lord, from the hands of the ungodly: and preserve me from the wicked men. ℣. Who are purposed to overthrow my goings: the proud have laid a snare for me. ℣. And spread a net abroad with cords: yea, and set traps in my way. ℣. I said unto the Lord: Thou art my God: hear the voice of my prayers, O Lord. ℣. O Lord God, thou strength of my health: thou hast covered my head in the day of battle. ℣. Let not the ungodly have his desire, O Lord: let not his mischievous imagination prosper, lest they be too proud. ℣. Let the mischief of their own lips fall upon the head of them: that compass me about. ℣. The righteous also shall give thanks unto thy name: and the just shall continue in thy sight.}

\begin{rubric}
	The Tract ended, the Gospel is read upon a bare pulpit: the Celebrant meanwhile reads it in a low voice at the Epistle corner.
\end{rubric}

\readingcitation{Gospel}{John 19:1}
\lett{A}{t that time:} Pilate took Jesus, and scourged him. And the soldiers platted a crown of thorns, and put it on his head, and they put on him a purple robe, and said, Hail, King of the Jews! and they smote him with their hands. Pilate therefore went forth again, and saith unto them, Behold, I bring him forth to you, that ye may know that I find no fault in him. Then came Jesus forth, wearing the crown of thorns, and the purple robe. And Pilate saith unto them, Behold the man! When the chief priests therefore and officers saw him, they cried out, saying, Crucify him, crucify him., Pilate saith unto them, Take ye him, and crucify him: for I find no fault in him. The Jews answered him, We have a law, and by our law he ought to die, because he made himself the Son of God. When Pilate therefore heard that saying, he was the more afraid; and went again into the judgment hall, and saith unto Jesus, Whence art thou? But Jesus gave him no answer. Then saith Pilate unto him, Speakest thou not unto me? knowest thou not that I have power to crucify thee, and have power to release thee? Jesus answered, Thou couldest have no power at all against me, except it were given thee from above: therefore he that delivered me unto thee hath the greater sin. And from thenceforth Pilate sought to release him: but the Jews cried out, saying, If thou let this man go, thou art not C{\ae}sar's friend: whosoever maketh himself a king, speaketh against C{\ae}sar. When Pilate therefore heard that saying, he brought Jesus forth, and sat down in the judgment seat in a place that is called the Pavement, but in the Hebrew, Gabbatha. And it was the preparation of the passover, and about the sixth hour: and he saith unto the Jews, Behold your King! But they cried out, Away with him, away with him, crucify him. Pilate saith unto them, Shall I crucify your King? The chief priests answered, We have no king but Cæsar. Then delivered he him therefore unto them to be crucified. And they took Jesus, and led him away. And he bearing his cross went forth into a place called the place of a skull, which is called in the Hebrew Golgotha: where they crucified him, and two other with him, on either side one, and Jesus in the midst. And Pilate wrote a title and put it on the cross. And the writing was, \divineName{Jesus of Nazareth the King of the Jews}. This title then read many of the Jews: for the place where Jesus was crucified was nigh to the city; and it was written in Hebrew, and Greek, and Latin. Then said the chief priests of the Jews to Pilate, Write not, The King of the Jews; but that he said, I am King of the Jews. Pilate answered, What I have written I have written. Then the soldiers, when they had crucified Jesus, took his garments, and made four parts, to every soldier a part; and also his coat: now the coat was without seam, woven from the top throughout. They said therefore among themselves, Let us not rend it, but cast lots for it, whose it shall be: that the scripture might be fulfilled, which saith, They parted my raiment among them, and for my vesture they did cast lots. These things therefore the soldiers did. Now there stood by the cross of Jesus his mother, add his mother's sister, Mary the wife of Cleophas, and Mary Magdalene. When Jesus therefore saw his mother, and the disciple standing by, whom he loved, he saith unto his mother, Woman, behold thy son! Then saith he to the disciple, Behold thy mother! And from that hour that disciple took her unto his own home. After this, Jesus knowing that all things were now accomplished, that the scripture might be fulfilled, saith, I thirst. Now there was set a vessel full of vinegar: and they filled a spunge with vinegar, and put it upon hyssop, and put it to his mouth. When Jesus therefore had received the vinegar, he said, It is finished: and he bowed his head, and gave up the ghost. \inrub{Here genuflect, and pause awhile.} The Jews therefore, because it was the preparation, that the bodies should not remain upon the cross on the sabbath day, (for that sabbath day was an high day,) besought Pilate that their legs might be broken, and that they might be taken away. Then came the soldiers, and brake the legs of the first, and of the other which was crucified with him. But when they came to Jesus, and saw that he was dead already, the brake not his legs: but one of the soldiers with a spear pierced his side, and forthwith came there out blood and water. And he that saw it bare record, and his record is true: and he knoweth that he saith true, that ye might believe. For these things were done, that the scripture should be fulfilled, A bone of him shall not be broken. And again another scripture saith, They shall look on him whom they pierced.

\begin{rubric}
	Then the Priest, standing at the Epistle side of the Altar, begins at once with joined hands:
\end{rubric}
\par\noindent
Let us pray dearly beloved, for the holy Church of God: that our God and Lord would vouchsafe to give her peace and unity, and preserve her throughout all the world: making subject unto her principalities and powers: and grant that, leading a quiet and peaceful life, we may glorify God the Father almighty.

\elcol{℣. Let us pray.}{℣. Orémus.}
\elcol{℣. Let us bow the knee.}{℣. Flectámus génua.}
\elcol{℣. Arise.}{℣. Leváte.}

\begin{rubric}
	The Collect is sung in the ferial tone of the Collect of the Mass, with hands extended. And the same manner is observed in all that follow.
\end{rubric}

\lett{A}{lmighty} and everlasting God, who in Christ hast revealed thy glory to all nations: preserve the works of thy mercy; that thy Church, spread abroad over the whole world, may with steadfast faith persevere in the confession of thy name. Through the same.\\
\par\noindent
Let us pray also for our most blessed Father \textit{N.}: that our God and Lord, who hath chosen him unto the order of bishops, may preserve him in health and safety to his holy Church, for the governance of the holy people of God.

\elcol{℣. Let us pray.}{℣. Orémus.}
\elcol{℣. Let us bow the knee.}{℣. Flectámus génua.}
\elcol{℣. Arise.}{℣. Leváte.}
\needspace{4\baselineskip}
\lett{A}{lmighty} and everlasting God, by whose judgement all things are established: mercifully regard our prayers, and in thy goodness preserve him whom thou hast chosen to be our Bishop; that the Christian people who are governed by thine authority may under so great a Pontiff increase in the merits of their faith. Through.\\
\par\noindent
Let us pray also for all Bishops, Priests, and Deacons, for all Subdeacons, Acolytes. Exorcists, Door-keepers, Confessors, Readers, Virgins, and Widows: and for all the holy people of God.

\elcol{℣. Let us pray.}{℣. Orémus.}
\elcol{℣. Let us bow the knee.}{℣. Flectámus génua.}
\elcol{℣. Arise.}{℣. Leváte.}
\needspace{4\baselineskip}
\lett{A}{lmighty} and everlasting God, by whose Spirit the whole body of the Church is governed and sanctified: receive our supplications, which we offer before thee for all orders of the same; that by the bounty of thy grace they may faithfully serve thee in their several estates. Through . . . in the unity of the same.\\
\par\noindent
Let us pray also for our most Christian Emperor (elect) \textit{N.}: that our God and Lord would make subject unto him all the nations of the heathen, for our perpetual peace.

\elcol{℣. Let us pray.}{℣. Orémus.}
\elcol{℣. Let us bow the knee.}{℣. Flectámus génua.}
\elcol{℣. Arise.}{℣. Leváte.}
\needspace{4\baselineskip}
\lett{A}{lmighty} and everlasting God, in whose hand is the dominion of all things and the government of all kingdoms: graciously look upon the Roman Empire; that the nations which trust in their own fierceness may be overthrown by the right hand of thy power. Through.\\
\par\noindent
Let us pray also for our catechumens: that our God and Lord would open the ears of their hearts, and the gate of mercy: that, receiving in the waters of regeneration the remission of all their sins, they also may be found in Christ Jesus our Lord.

\elcol{℣. Let us pray.}{℣. Orémus.}
\elcol{℣. Let us bow the knee.}{℣. Flectámus génua.}
\elcol{℣. Arise.}{℣. Leváte.}
\needspace{4\baselineskip}
\lett{A}{lmighty} and everlasting God, who dost continually enrich thy Church with a new offspring: increase the faith and understanding of our catechumens; that they, being born again in the water of baptism, may be numbered among the sons of thine adoption. Through.\\
\par\noindent
Let us pray, dearly beloved, unto God the Father almighty, that he would purge the world from all errors: would take away diseases: drive away famine: open the prisons: loosen the chains: grant unto pilgrims a safe return: to the sick, healing: and to them that travel by sea a haven of safety.

\elcol{℣. Let us pray.}{℣. Orémus.}
\elcol{℣. Let us bow the knee.}{℣. Flectámus génua.}
\elcol{℣. Arise.}{℣. Leváte.}
\needspace{4\baselineskip}
\lett{A}{lmighty} and everlasting God, the comfort of them that mourn, the strength of them that travail: let the prayers of them that cry out of any tribulation ascend unto thee; that in their necessities all may rejoice in the succour of thy loving kindness. Through.\\
\par\noindent
Let us pray also for heretics and schismatics: that our God and Lord would deliver them from all their errors; and vouchsafe to call them back to their holy mother, the Catholic and Apostolic Church.

\elcol{℣. Let us pray.}{℣. Orémus.}
\elcol{℣. Let us bow the knee.}{℣. Flectámus génua.}
\elcol{℣. Arise.}{℣. Leváte.}
\needspace{4\baselineskip}
\lett{A}{lmighty} and everlasting God, who savest all men, and wouldest not that any should perish: look upon the souls that are deceived by the craft of the devil; that the hearts of them that are gone astray being delivered from all perversity of heresy, may turn to wisdom and come again to the unity of thy truth. Through.\\
\par\noindent
Let us pray also for the faithless Jews: that our God and Lord would take away the veil from their hearts; that they also may acknowledge Jesus Christ, our Lord.

\begin{rubric}
	\emph{Amen} is not said, nor \emph{Let us pray}, nor \emph{Let us bow the knee}, nor \emph{Arise}, but at once is said:
\end{rubric}
\needspace{4\baselineskip}
\lett{A}{lmighty} and everlasting God, who deniest not thy mercy even to the faithless Jews: graciously hear our prayers, which we offer for the blindness of this people; that they acknowledging the light of thy truth, which is Christ, may be delivered from their darkness. Through the same.\\
\par\noindent
Let us pray also for the heathen: that God almighty would take away the iniquity from their hearts; that, forsaking their idols, they may be turned unto the living and true God, and to his only Son, Jesus Christ, our God and Lord.

\elcol{℣. Let us pray.}{℣. Orémus.}
\elcol{℣. Let us bow the knee.}{℣. Flectámus génua.}
\elcol{℣. Arise.}{℣. Leváte.}
\needspace{4\baselineskip}
\lett{A}{lmighty} and everlasting God, who desirest not the death of sinners but rather that they should live: mercifully receive our prayer, and deliver them from the worship of idols; and gather them unto thy holy Church, to the praise and glory of thy name. Through.

\begin{rubric}
	Having finished the Prayers, the Priest lays aside his Chasuble and goes to the Epistle side of the Altar, where at the back corner he receives from the Deacon the Cross which has previously been made ready on the Altar. With his face to the people, he uncovers the upper portion, beginning alone \emph{Behold the wood of the Cross}, and from thence he is helped in singing the rest by the Ministers until \emph{O come, let us worship}.\par
When the Choir sing \emph{O come, let us worship}, all prostrate themselves except the Celebrant. He then advances to the front corner of the same Epistle side: and uncovering the right arm of the Cross, and raising it a little, he begins in a higher key than at first: \emph{Behold the wood of the Cross}, the others singing and worshipping as before. Then the Priest proceeds to the middle of the Altar: and uncovering the Cross entirely and raising it, for the third time he begins in a still higher key: \emph{Behold the wood of the Cross}, the others singing and worshipping as before.
\end{rubric}
Behold the wood of the Cross, whereon was hung the world's salvation.

\textit{Choir.} ℟. O come, let us worship.

\begin{rubric}
	Then the Priest alone carries the Cross to the place prepared before the Altar, and kneeling there he puts it down: then laying aside his shoes, he proceeds to worship the Cross %ADDITION TO CLARIFY MANNER OF WORSHIP:
	with the worship of 
	\begin{otherlanguage}{greek}
	προσκύνησις, 
	\end{otherlanguage}
	that is, dulia.
	\par
	And kneeling thrice before, he kisses it. This done, he returns, and puts on his shoes and the Chasuble. Afterwards the Ministers of the Altar, then the other Clerics and laics, two and two, kneeling thrice, as has been said, worship the Cross. Meantime, while the worship of the Cross is proceeding, the Reproaches are sung, together with the other things that follow, either all or in part, according as the many or few that worship the Cross require: which also the Priest sitting at the bench reads with the Ministers, as below:
\end{rubric}

\textsc{Celebrant.} O my people, what have I done unto thee? or wherein have I wearied thee? Answer me.\par
℣. Because I brought thee forth from the land of Egypt: thou hast prepared a Cross for thy Saviour.

\textsc{Deacon.} \begin{otherlanguage}{greek} Ἅγιος ὁ Θεός.\end{otherlanguage}\par
\textsc{Subdeacon.} Holy God.\par
\textsc{Deacon.} \begin{otherlanguage}{greek} Ἅγιος ἰσχυρός.\end{otherlanguage}\par
\textsc{Subdeacon.} Holy Mighty.\par
\textsc{Deacon.} \begin{otherlanguage}{greek} Ἅγιος ἀθάνατος, ἐλέησον ἡμᾶς.\end{otherlanguage}\par
\textsc{Subdeacon.} Holy Immortal, have mercy upon us.\par
\textsc{Celebrant.} Because I led thee through the desert forty years, and fed thee with manna, and brought thee into a land exceeding good: thou hast prepared a Cross for thy Saviour.\par
\textsc{Deacon.} \begin{otherlanguage}{greek} Ἅγιος ὁ Θεός.\end{otherlanguage}\par
\textsc{Subdeacon.} Holy God.\par
\textsc{Deacon.} \begin{otherlanguage}{greek} Ἅγιος ἰσχυρός.\end{otherlanguage}\par
\textsc{Subdeacon.} Holy Mighty.\par
\textsc{Deacon.} \begin{otherlanguage}{greek} Ἅγιος ἀθάνατος, ἐλέησον ἡμᾶς.\end{otherlanguage}\par
\textsc{Subdeacon.} Holy Immortal, have mercy upon us.\par
\textsc{Celebrant.} What more could I have done for thee that I have not done? I indeed did plant thee, my vineyard, exceeding fair: and thou art become very bitter unto me: for vinegar thou gavest to quench my thirst: and hast pierced with a spear the side of thy Saviour.\par
\textsc{Deacon.} \begin{otherlanguage}{greek} Ἅγιος ὁ Θεός.\end{otherlanguage}\par
\textsc{Subdeacon.} Holy God.\par
\textsc{Deacon.} \begin{otherlanguage}{greek} Ἅγιος ἰσχυρός.\end{otherlanguage}\par
\textsc{Subdeacon.} Holy Mighty.\par
\textsc{Deacon.} \begin{otherlanguage}{greek} Ἅγιος ἀθάνατος, ἐλέησον ἡμᾶς.\end{otherlanguage}\par
\textsc{Subdeacon.} Holy Immortal, have mercy upon us.\par
\textsc{Celebrant.} I did scourge Egypt with her first-born for thy sake: and thou hast scourged me and delivered me up.\par
\textsc{Deacon and Subdeacon.} O my people, what have I done unto thee, or wherein have I wearied thee? Answer me.\par 
\textsc{Celebrant.} I led thee out of Egypt, drowning Pharaoh in the Red Sea: and thou hast delivered me unto the chief priests.\par
\textsc{Deacon and Subdeacon.} O my people, what have I done unto thee, or wherein have I wearied thee? Answer me.\par
\textsc{Celebrant.} I opened the sea before thee: and thou hast opened my side with a spear.\par
\textsc{Deacon and Subdeacon.} O my people, what have I done unto thee, or wherein have I wearied thee? Answer me.\par
\textsc{Celebrant.} I went before thee in a pillar of cloud: and thou hast led me unto the judgment hall of Pilate.\par
\textsc{Deacon and Subdeacon.} O my people, what have I done unto thee, or wherein have I wearied thee? Answer me.\par
\textsc{Celebrant.} I fed thee with manna in the desert: and thou hast stricken me with blows and scourges.\par
\textsc{Deacon and Subdeacon.} O my people, what have I done unto thee, or wherein have I wearied thee? Answer me.\par
\textsc{Celebrant.} I gave thee to drink of the water of salvation from the rock: and thou hast given me gall and vinegar to drink.\par
\textsc{Deacon and Subdeacon.} O my people, what have I done unto thee, or wherein have I wearied thee? Answer me.\par
\textsc{Celebrant.} For thee I smote the kings of the Canaan-ites: and thou hast smitten my head with a reed.\par
\textsc{Deacon and Subdeacon.} O my people, what have I done unto thee, or wherein have I wearied thee? Answer me.\par
\textsc{Celebrant.} I gave thee a royal sceptre: and thou hast given unto my head a crown of thorns.\par
\textsc{Deacon and Subdeacon.} O my people, what have I done unto thee, or wherein have I wearied thee? Answer me.\par
\textsc{Celebrant.} I exalted thee with great power: and thou hast hanged me upon the gibbet of the Cross.\par
\textsc{Deacon and Subdeacon.} O my people, what have I done unto thee, or wherein have I wearied thee? Answer me.\par
\textsc{Celebrant.} We worship thy Cross, O Lord: and praise and glorify thy holy resurrection: for behold, by virtue of the tree joy hath come to the whole world. \textit{Ps.} God be merciful unto us and bless us :\par
\textsc{Deacon and Subdeacon.} And shew us the light of his countenance, and be merciful unto us.\par
\textsc{Celebrant.} We worship thy Cross, O Lord: and praise and glority thy holy resurrection: for behold, by virtue of the tree joy hath come to the whole world.\par
\textsc{Deacon and Subdeacon.} Faithful Cross, above all other, one and only noble tree: none in foliage, none in blossom, none in fruit thy peer may be.\par
Sweetest wood and sweetest iron, sweetest weight is hung on thee.\par
\textsc{Celebrant.} \textit{Hymn.} Sing, my tongue, the glorious battle, sing the ending of the fray, now above the Cross, loud triumphant lay: tell how Christ, the world's Redeemer, as a victim won the day.\par
\textsc{Deacon and Subdeacon.} Faithful Cross, above all other, one and only noble tree: none in foliage, none in blossom, none in fruit thy peer may be.\par
\textsc{Celebrant.} God in pity saw man fallen, shamed and sunk in misery, when he fell on death by tasting fruit of the forbidden tree: then another tree was chosen which the world from death should free.\par
\textsc{Deacon and Subdeacon.} Sweetest wood and sweetest iron, sweetest weight is hung on thee.\par
\textsc{Celebrant.} Thus the scheme of our salvation, was of old in order laid: that the manifold deceiver's art by art might be outweighed: and the lure the toe put forward into means of healing made.\par
\textsc{Deacon and Subdeacon.} Faithful Cross, above all other, one and only noble tree: none in foliage, none in blossom, none in fruit thy peer may be.\par
\textsc{Celebrant.} Therefore when the appointed fulness of the holy time was come, he was sent who maketh all things forth from God's eternal home; thus he came to earth, incarnate, offspring of a virgin's womb.\par
\textsc{Deacon and Subdeacon.} Sweetest wood and sweetest iron, sweetest weight is hung on thee.\par
\textsc{Celebrant.} Lo! he lies, an Infant weeping, where the narrow manger stands, while the Mother-Maid his members wraps in mean and lowly bands, and the swaddling clothes is winding round his helpless feet and hands.\par
\textsc{Deacon and Subdeacon.} Faithful Cross, above all other, one and only noble tree: none in foliage, none in blossom, none in fruit thy peer may be.\par
\textsc{Celebrant.} Thirty years among us dwelling, his appointed time fulfilled, born for this, he meets his Passion, for that this he freely willed, on the Cross the Lamb is lifted where his life-blood shall be spilled.\par
\textsc{Deacon and Subdeacon.} Sweetest wood and sweetest iron, sweetest weight is hung on thee.\par
\textsc{Celebrant.} He endured the nails, the spitting, vinegar, and spear, and reed; from that holy body broken blood and water forth proceed: earth, and stars, and sky, and ocean by that flood from stain are freed.\par
\textsc{Deacon and Subdeacon.} Faithful Cross, above all other, one and only noble tree: none in foliage, none in blossom, none in fruit thy peer may be.\par
\textsc{Celebrant.} Bend thy boughs, O tree of glory! thy relaxing sinews bend; for awhile the ancient rigour that thy birth bestowed, suspend; and the King of heavenly beauty on thy bosom gently tend!\par
\textsc{Deacon and Subdeacon.} Sweetest wood and sweetest iron, sweetest weight is hung on thee.\par
\textsc{Celebrant.} Thou alone wast counted worthy this world's ransom to uphold ; for a shipwreck'd race preparing harbour, like the ark of old; with the sacred blood anointed from the smitten Lamb that rolled.\par
\textsc{Deacon and Subdeacon.} Faithful Cross, above all other, one and only noble tree: none in foliage, none in blossom, none in fruit thy peer may be.\par
\textsc{Celebrant.} To the Trinity be glory everlasting, as is meet; equal to the Father, equal to the Son, and Paraclete: Trinal Unity, whose praises all created things repeat. Amen.\par
\textsc{Deacon and Subdeacon.} Sweetest wood and sweetest iron, sweetest weight is hung on thee.\par

\begin{rubric}
	Towards the end of the worship of the Cross candles are lighted upon the Altar: and the Deacon, receiving the Burse with the Corporal, spreads the Corporal in the usual way, and puts a Purificator by it; and when the adoration is over, reverently receives the Cross and carries it back to the Altar. Then a Procession is formed to the place where on the previous day the Sacrament bad been laid. A Subdeacon goes first with the Cross between two Acolytes bearing candlesticks with lighted candles, and then the Clergy in order, last of all the Priest with the Ministers. When they arrive at the place of the Sacrament, torches are kindled, which are not extinguished till after the consumption of the Sacrament. The Priest genuflects before the Sacrament and prays awhile: the Deacon meanwhile opens the receptacle in which lies hidden the Body of the Lord. Then the Priest, rising, puts incense into two thuribles without blessing, the Deacon presenting the boat, and then kneeling he censes the Sacrament.\par
	The Deacon takes the Chalice with the Sacrament from the receptacle and gives it into the hands of the Priest, covering it with the ends of the Veil which is around his shoulders. Then they proceed in the order in which they came: a baldachino is carried over the Sacrament: and two Acolytes with thuribles continually cense the Sacrament. Meanwhile is sung the following hymn:
\end{rubric}

\subbysub{Vexilla Regis prodeunt}

\gregorioscore{resources/gabc/ProperSeason/GoodFriday.gabc}

\begin{rubric}
	When the Priest arrives at the Altar, he places the Chalice upon it, genuflects, and again censes it, and going up takes the Host out of the Chalice and places it upon the Paten which the Deacon holds: and receiving the Paten from the hand of the Deacon, he places the sacred Host upon the Corporal, saying nothing. If he touches the Sacrament, he washes his hands in some vessel. Meanwhile the Deacon puts wine in the Chalice, and the Subdeacon water, which the Priest does not bless, nor does he say over it the usual Prayer: but receiving the Chalice from the Deacon, be places it upon the Altar, saying nothing: and the Deacon covers it with the Pall.\par
Then he puts incense in the thurible without blessing it, and censes the Oblations, Cross, and Altar in the accustomed manner, genuflecting before and after, and whenever be crosses before the Sacrament.
\end{rubric}

\begin{rubric}
	When he censes the Altar, he says:
\end{rubric}
\lett{L}{et} my prayer, O Lord, be set forth in thy sight as the incense: and let the lifting up of my hands be an evening sacrifice. Set a watch, O Lord, before my mouth, and keep the door of my lips: O let not mine heart be inclined to any evil thing, let me not be occupied in ungodly works.

\begin{rubric}
	When he gives back the thurible to the Deacon he says:
\end{rubric}

\lett{T}{he} Lord kindle in us the fire of his love and the flame of everlasting charity. Amen.

\begin{rubric}
	And he himself is not censed.
\end{rubric}

\begin{rubric}
	Afterwards he washes his hands a little beyond the Altar on the Epistle side, saying nothing: then bowing in the middle of the Altar, with hands joined, he says:
\end{rubric}

\lett{I}{n} a humble spirit and with a contrite heart may we be accepted of thee, O Lord: and so let our sacrifice be offered in thy sight this day, that it may be pleasing unto thee, O Lord God.

\begin{rubric}
	Then turning to the people at the Gospel corner, he says as usual:
\end{rubric}

\lett{P}{ray,} brethren, that my sacrifice and yours may be acceptable to God the Father almighty.

\begin{rubric}
	And he turns back by the same way, not completing the circle: and, omitting the rest, proceeds:
\end{rubric}

Let us pray: Commanded by saving precepts, and taught by divine institution, we are bold to say:

\lett{O}{ur} Father, who art in heaven: Hallowed be thy name: Thy kingdom come: Thy will be done, on earth as it is in heaven. Give us this day our daily bread: And forgive us our trespasses, as we forgive those that trespass against us. And lead us not into temptation.\par
℟. But deliver us from evil.

\begin{rubric}
	The Priest says silently \emph{Amen}. In the same voice in which he said \emph{Our Father}, and without \emph{Let us pray}, he says in the tone of the Collect of the ferial Mass:
\end{rubric}

\lett{D}{eliver} us, O Lord, we beseech thee, from all evils, past, present, and to come: and at the intercession of the blessed and glorious ever Virgin Mary Mother of God, with thy blessed Apostles Peter and Paul, and with Andrew, and all the Saints, \inrub{he does not sign himself with the Paten} favourably grant peace in our days: that by the help of thine availing mercy we may ever both be free from sin, and safe from all distress. Through the same thy Son Jesus Christ our Lord: Who liveth and reigneth with thee, in the unity of the Holy Ghost, God, throughout all ages, world without end.\par
℟. Amen.

\subbysub{Confession}
\begin{rubric}
    Then the Deacon turneth to the People and saith to those who come to receive the Holy Communion,
\end{rubric}
\lett{Y}{e} who do truly and earnestly repent you of your sins, and are in love and charity with your neighbours, and intend to lead a new life, following the commandments of God, and walking from henceforth in his holy ways; Draw near with faith, and take this holy Sacrament to your comfort; and make your humble confession to Almighty God, devoutly kneeling.
\begin{rubric}
Then shall this General Confession be made, in an audible voice, by the Priest, bowing, and all those who are minded to receive the Holy Communion, humbly kneeling.
\end{rubric}
\needspace{4\baselineskip}
\lett{A}{lmighty} God, Father of our Lord Jesus Christ, Maker of all things, Judge of all men; We acknowledge and bewail our manifold sins and wickedness, Which we, from time to time, most grievously have committed, By thought, word, and deed, Against thy Divine Majesty, Provoking most justly thy wrath and indignation against us. We do earnestly repent, And are heartily sorry for these our misdoings; The remembrance of them is grievous unto us; The burden of them is intolerable. Have mercy upon us, Have mercy upon us, most merciful Father; For thy Son our Lord Jesus Christ's sake, Forgive us all that is past; And grant that we may ever hereafter Serve and please thee in newness of life, To the honour and glory of thy Name; Through Jesus Christ our Lord. Amen.
\begin{rubric}
	Then shall the Priest stand up and, turning to the People, say,
\end{rubric}
\needspace{4\baselineskip}
\lett{A}{lmighty} God, our heavenly Father, who of his great mercy hath promised forgiveness of sins to all those who with hearty repentance and true faith turn unto him; Have mercy upon you; pardon {\ding{64}} and deliver you from all your sins; confirm and strengthen you in all goodness; and bring you to everlasting life; through Jesus Christ our Lord.

℟. Amen.
\subbysub{Comfortable Words}
\begin{rubric}
Then shall the Priest, facing the People, say,
\end{rubric}\noindent
\begin{center}
	\textsc{Hear what comfortable words our Saviour Christ saith unto all who truly turn to him.}
\end{center}
\par\noindent
Come unto me, all ye that travail and are heavy laden, and I will refresh you. \vr{Matt. 11:28}
\par\noindent
    So God loved the world, that he gave his only-begotten Son, to the end that all that believe in him should not perish, but have everlasting life. \vr{John 3:16}
    \par\noindent
    \begin{center}
		\textsc{Hear also what Saint Paul saith.}
	\end{center}
\par\noindent
    This is a true saying, and worthy of all men to be received, That Christ Jesus came into the world to save sinners. \vr{1 Tim. 1:15}

    \begin{center}
		\textsc{Hear also what Saint John saith.}
	\end{center}
    \par\noindent
    If any man sin, we have an Advocate with the Father, Jesus Christ the righteous; and he is the Propitiation for our sins. \vr{1 John 2:1-2}

\subbysub{Prayer of Humble Access}
\begin{rubric}
    The Priest alone then saith, in an audible voice,
\end{rubric}
\lett{W}{e} do not presume to come to this thy Table, O merciful Lord, trusting in our own righteousness, but in thy manifold and great mercies. We are not worthy so much as to gather up the crumbs under thy Table. But thou art the same Lord, whose property is always to have mercy: Grant us therefore, gracious Lord, so to eat the flesh of thy dear Son Jesus Christ, and to drink his blood, that our sinful bodies may be made clean by his Body, and our souls washed through his most precious Blood, and that we may evermore dwell in him, and he in us.\par
℟. Amen.

\subbysub{Communion of the Priest}

\begin{rubric}
	Then the Celebrant, having made a reverence down to the ground, places the Paten beneath the Sacrament, which he takes in his right hand and elevates so that it can be seen by the people: and immediately he divides it over the Chalice into three parts, the last of which he puts into the Chalice in the usual way, saying nothing. \emph{The Peace of the Lord} is not said, nor \emph{O Lamb of God}, nor is the kiss of peace given. Then, omitting the first two Prayers, he says only the following:
\end{rubric}

\lett{L}{et} the partaking of thy Body, O Lord Jesu Christ, which I, unworthy, presume to receive, turn not to my judgement and condemnation: but of thy goodness let it avail unto me for protection of mind and of body, that I may receive thy healing: Who livest and reignest with God the Father in the unity of the Holy Ghost, God, throughout all ages, world without end. Amen.

\begin{rubric}
	Then he genuflects and receives the Paten with the Body of Christ: and with the deepest humility and reverence says:
\end{rubric}

I will receive the bread of heaven, and call upon the name of the Lord.

\begin{rubric}
	He beats his breast, saying thrice:
\end{rubric}

Lord, I am not worthy that thou shouldest come under my roof: but speak the word only, and my soul shall be healed.

\begin{rubric}
	Then he signs himself with the Sacrament, saying:
\end{rubric}

The Body of our Lord Jesus Christ preserve my soul unto everlasting life. Amen.

\begin{rubric}
	And he reverently consumes the Body.
\end{rubric}

\begin{rubric}
	Then, omitting everything which is wont to be said before the consumption of the Blood, he at once reverently consumes the particle of the Host with the wine from the Chalice.
\end{rubric}

\subbysub{Communion of the People}
\begin{rubric}
	The Priest then, turning towards the Congregation, holds a particle over the Chalice and says:
\end{rubric}
\elcol{℣. Behold the Lamb of God. Behold him who taketh away the sins of the world.}{℣. Ecce Agnus Dei, ecce, qui tollit peccáta mundi.}
\elcol{℟. Lord, I am not worthy that thou shouldest come under my roof, but speak the word only, and my soul shall be healed.}{℟. Dómine, non sum dignus, ut intres sub tectum meum, sed tantum dic verbo, et sanábitur ánima mea.}
\elcol{℟. Lord, I am not worthy that thou shouldest come under my roof, but speak the word only, and my soul shall be healed.}{℟. Dómine, non sum dignus, ut intres sub tectum meum, sed tantum dic verbo, et sanábitur ánima mea.}
\elcol{℟. Lord, I am not worthy that thou shouldest come under my roof, but speak the word only, and my soul shall be healed.}{℟. Dómine, non sum dignus, ut intres sub tectum meum, sed tantum dic verbo, et sanábitur ánima mea.}

\begin{rubric}
	The Priest then distributes Communion, in both kinds, to the People in the usual manner.
\end{rubric}

\begin{rubric}
	After distributing the Communion to the People, he makes the ablution of his fingers in the usual manner and taken the purification. He then says, bowing in the midst of the Altar, with hands joined:
\end{rubric}

\lett{G}{rant,} O Lord, that what we have taken with our mouths we may receive with a pure heart: and from a temporal gift may it become to us an everlasting remedy.

%MANUAL ADJUSTMENT:
\clearpage
\begin{rubric}
	\emph{May thy Body} is not said, nor the Postcommunion, nor \emph{Let this my bounden duty}, nor is the blessing given: but having made a reverence to the Altar, the Priest retires with the Ministers.
\end{rubric}

\subbysub{Thanksgiving}
\begin{rubric}
    Standing with his hands joined, the Priest shall alone say the Thanksgiving, in an audible voice.
\end{rubric}
\letuspray
\lett{A}{lmighty} and everliving God, we most heartily thank thee, for that thou dost vouchsafe to feed us who have duly received these holy mysteries with the spiritual food of the most precious Body and Blood of thy Son our Saviour Jesus Christ; and dost assure us thereby of thy favour and goodness towards us; and that we are very members incorporate in the mystical body of thy Son, which is the blessed company of all faithful people; and are also heirs through hope of thy everlasting kingdom, by the merits of his most precious death and passion. And we humbly beseech thee, O heavenly Father, so to assist us with thy grace, that we may continue in that holy fellowship, and do all such good works as thou hast prepared for us to walk in; through Jesus Christ our Lord, to whom, with thee and the Holy Ghost, be all honour and glory, world without end.\par
℟. Amen.

\begin{rubric}
	Evensong is then said without singing: and the Altar is stripped.
\end{rubric}


\subby{Easter Even}
\fancyhead[RO,LE]{\textit{Easter Even}}
\fancyhead[RE,LO]{}
\begin{inhead}
    {First Class Double}
\end{inhead}

\properantiphon{Ben.}{The women, {\dag} sitting over against the sepulchre, made lamentation, weeping for the Lord.}
    
\subbysub{Blessing of the Paschal Fire}
\begin{rubric}
	At a suitable time the Altars are covered with linen cloths, but the candles remain unlit until the beginning of Mass. Meanwhile fire is struck from a flint outside the church, and from it coals are kindled. After Mattins (or None) has been said in Choir, the Priest, vested in Amice, Alb, Girdle, Stole, and violet Cope, or without a Chasuble, standing with the Ministers, with the Cross, holy water, and incense, before the door of the Church, if it can conveniently be done, or in the entry of the church itself, blesses the new fire, saying:
\end{rubric}

℣. The Lord be with you.

℟. And with thy spirit.

\letuspray
\lett{O}{God,} who through thy Son, the true corner-stone, hast bestowed upon thy faithful people the fire of thy brightness: sanct {\ding{64}} ify this new fire, now struck from the flint-stone, to be profitable to our service: and grant unto us that by this paschal feast we may be so inflamed with heavenly desires; that we may with pure hearts attain unto the feast of thy eternal brightness. Through the same Christ, our Lord. \textit{Amen.}

\letuspray
\lett{O}{Lord} God, Father almighty, who art light unfailing, and the author of all lights: bl {\ding{64}} ess this light which hath here been sanctified and blessed by thee, who hast enlightened the whole world: that we may be enkindled by that light, and enlightened with the fire of thy brightness: and like as thou didst enlighten Moses going forth from Egypt, so do thou enlighten our hearts and understandings; that we may be found worthy to attain unto life and light everlasting. Through Christ, our Lord. \textit{Amen.}

\letuspray
\lett{O}{Lord} holy, Father almighty, everlasting God: vouchsafe to assist us who bless this fire in thy name, and in the name of thy only-begotten Son, Jesus Christ, our God and Lord, and of the Holy Ghost; and assist us against the fiery darts of the enemy, and enlighten us with thy heavenly grace: Who livest and reignest with the same thy only-begotten Son and the Holy Ghost, God: throughout all ages world without end. \textit{Amen.}

\begin{rubric}
	He then blesses the five grains of incense, to be placed in the Candle, saying at once this Collect :
\end{rubric}
\lett{W}{e} beseech thee, almighty God, that the abundance of thy bles {\ding{64}} sing may come down upon this incense; and as thou, thyself unseen, art the regenerator of all things, so kindle the brightness of this night; that not only may our sacrifice, which is offered on this night, shine with the secret mingling of thy light; but that in all places wheresoever aught shall be carried from this mystery here sanctified, the wickedness of the crafts of the devil may be driven forth, and the power of thy majesty ever be present. Through Christ, our Lord. \textit{Amen.}

\subbysub{Procession of the Light of Christ}

\begin{rubric}
	Whilst be blesses the grains of incense, an Acolyte takes some of the blessed coals, and sets them in the thurible: and the Priest, when he has finished the last Prayer, puts incense from the boat into the thurible, giving the usual blessing. Then he sprinkles the aforesaid grains of incense and the fire with holy water, saying: \emph{Thou shalt purge me, O Lord}, without note and without the Psalm, and censes them thrice.

	Meanwhile all the lights in the church are extinguished, to be relighted later from the blessed fire.
	
	Then the Deacon, vested in a white Dalmatic, takes the reed upon the top of which are three candles in the shape of a triangle. The Thurifer goes first with an Acolyte bearing the five grains of incense: the Subdeacon follow's with the Cross, the Clergy in their order: then the Deacon with the reed, after him the Celebrant.
	
	When the Deacon enters the church, he lowers the reed, and an Acolyte, carrying a candle which has been lighted from the new fire, lights one of the three candles placed upon it, and the Deacon, elevating the reed, genuflects, as do all others with him except the Subdeacon bearing the Cross, and sings alone: \emph{The Light of Christ}.
\end{rubric}

℟. Thanks be to God.

\begin{rubric}
	And proceeding to the middle of the church, another candle is lighted: and having again genuflected, he sings in a higher tone: \emph{The Light of Christ}.
\end{rubric}

℟. Thanks be to God.

\begin{rubric}
	For the third time he proceeds before the Altar, where the third candle is lighted: and again, having genuflected as before, he sings still higher: \emph{The Light of Christ}.
\end{rubric}

℟. Thanks be to God.

\begin{rubric}
	The Celebrant then ascends to the Epistle corner of the Altar, and the Deacon, banding the reed with the triple candle to an Acolyte, receives the book, and asks a blessing from the Celebrant as at the Gospel, the Priest saying:
\end{rubric}
\par\noindent
The Lord be in thy heart, and on thy lips: that thou mayest worthily and fitly proclaim his paschal praise: In the name of the Father, and of the Son, {\ding{64}} and of the Holy Ghost. Amen.

\begin{rubric}
	The Deacon then goes to the desk and placing the book upon it, censes it. On the right hand of the Deacon stand the Subdeacon with the Cross, and the Thurifer: on his left an Acolyte holding the reed, and another the five grains of blessed incense to be fixed into the Candle. Then, all rising and standing, as at the Gospel, the Deacon sings:
\end{rubric}

\lett{N}{ow} let the Angelic host of heaven rejoice: let the divine mysteries rejoice: and for the victory of the mighty King let the trumpet of salvation sound forth. Let the earth also be glad, illumined by the rays of this great brightness: and enlightened by the splendour of the eternal King, let her know that she hath put away the darkness of the whole world. Let our mother the Church also rejoice, adorned with the brightness of so great a light: and let this temple resound with the triumphant voices of the peoples. Wherefore, dearly beloved brethren, who are here present in the wondrous clearness of this holy light, join with me, I beseech you, in calling upon the mercy of almighty God. That he who hath been pleased, for no merit of mine, to admit me into the number of his Levites: may pour on me the brightness of his light, and make me meet to proclaim the praises of this Candle. Through Jesus Christ his Son our Lord: Who liveth and reigneth with him in the unity of the Holy Ghost, God. Throughout all ages, world without end. \textit{Amen.}


℣. The Lord be with you.

℟. And with thy spirit.

℣. Lift up your hearts.

℟. We lift them up unto the Lord.

℣. Let us give thanks unto our Lord God.

℟. It is meet and right so to do.
\par\noindent
It is very meet, right, and our bounden duty, that we should, with the whole affection of our heart and mind, and with the service of our lips, give praise unto the invisible God, the Father almighty, and unto his only begotten Son, Jesus Christ our Lord. Who paid for us to the eternal Father the debt of Adam's transgression: and with his dear blood wiped away the reproach of our former offences. For this is the paschal feast wherein the very Lamb is slain, by whose blood the door-posts of the faithful are made holy. This is the night, wherein of old thou didst lead forth our fathers, the children of Israel, out of Egypt, and didst make them to pass as on dry land through the Red Sea. Yea, this is the night, that with the brightness of the pillar hath purged away the darkness of iniquity. This is the night, which throughout the whole world today doth separate them that believe in Christ from the wickedness of the age and from the darkness of transgressions, reneweth them unto grace, restoreth them to holiness. This is the night wherein, breaking the chains of death, Christ ascendeth from hell in triumph. For nought indeed had it profited us to be born, if it had not profited us to be redeemed. O how wonderful the condescension of thy loving kindness! O how inestimable the goodness of thy love: who to redeem a slave didst deliver up thy Son! O truly necessary sin of Adam, which by the death of Christ was done away! O happy fault, which was counted worthy to have such and so great a Redeemer! O night truly blessed, which alone was worthy to know the time and the hour wherein Christ rose again from hell! This is the night whereof it is written: And the night is as clear as the day: And, then shall my night be turned into day. The sanctifying power therefore of this night putteth to flight the deeds of wickedness, washeth away sins: restoreth innocence to the fallen, and joy to them that mourn: casteth out enmities, prepareth concord, and boweth down principalities.

\begin{rubric}
	Here the Deacon fixes the five grains of incense in the Candle, placing them in the form of a cross, in this order:
\end{rubric}

\begin{center}
	1
	
	4 2 5
	
	3
\end{center}
\par\noindent
Therefore in this night of grace accept, O holy Father, the evening sacrifice of this incense: which, by the hands of thy ministers, holy Church doth lay before thee, in the solemn offering of this Candle, made from the work of bees. But we already know the excellency of this pillar, which for the honour of God the sparkling fire doth kindle.

\begin{rubric}
	Here the Deacon lights the Candle with one of the three candles set on the reed.
\end{rubric}
\par\noindent
Which, though it be divided into parts, suffereth not loss by the borrowing of its light. For it is fed by melting wax, which the bee the mother hath wrought into the substance of this precious candle.

\begin{rubric}
	Here the lamps are lighted.
\end{rubric}
\par\noindent
O night truly blessed, which did spoil the Egyptians, and made rich the Hebrews! O night wherein heavenly things are joined unto earthly, things human unto things divine. We therefore pray thee, O Lord: that this Candle, consecrated to the honour of thy name, may continue without ceasing to vanquish the darkness of this night. That, being accepted for a savour of sweetness, it may be mingled with the lights of heaven. May the morning star find it burning. That morning star, I say, which knoweth not his going down. That star, which, rising again from hell, steadfastly giveth light to all mankind. We therefore pray thee, O Lord: that thou wouldest vouchsafe to rule, govern and preserve with thy continual protection us thy servants, the whole clergy and thy most faithful people: together with our most blessed Father \textit{N.} (and our Bishop \textit{N.}), granting us peaceful times in this out paschal joy. Through the same Jesus Christ thy Son our Lord: Who liveth and reigneth with thee in the unity of the Holy Ghost God: throughout all ages, world without end. \textit{Amen.}

\subbysub{The Prophecies}

\begin{rubric}
	The blessing of the Candle ended, the Deacon lays aside his white vestments, and puts on violet, and goes to the Celebrant: who, laying aside his Cope, puts on a Maniple and violet Chasuble. Then the Prophecies are read without title, and at the end the response \emph{Thanks be to God} is not made, and the Celebrant reads them in a low voice at the Epistle corner of the Altar.
	
	At the end of the Prophecies are said the Prayers as given below.
	
	Before, or during the reading of the Prophecies, the Priests catechize any Catechumens about to be baptized, and prepare them for Baptism.
\end{rubric}

\readingcitation{Prophecy I}{Genesis 1:1}
\lett{I}{n} the beginning God created the heaven and the earth. And the earth was without form, and void; and darkness was upon the face of the deep. And the Spirit of God moved upon the face of the waters. And God said, Let there be light: and there was light. And God saw the light, that it was good: and God divided the light from the darkness. And God called the light Day, and the darkness he called Night. And the evening and the morning were the first day.

And God said, Let there be a firmament in the midst of the waters, and let it divide the waters from the waters. And God made the firmament, and divided the waters which were under the firmament from the waters which were above the firmament: and it was so. And God called the firmament Heaven. And the evening and the morning were the second day.

And God said, Let the waters under the heaven be gathered together unto one place, and let the dry land appear: and it was so. And God called the dry land Earth; and the gathering together of the waters called he Seas: and God saw that it was good. And God said, Let the earth bring forth grass, the herb yielding seed, and the fruit tree yielding fruit after his kind, whose seed is in itself, upon the earth: and it was so. And the earth brought forth grass, and herb yielding seed after his kind, and the tree yielding fruit, whose seed was in itself, after his kind: and God saw that it was good. And the evening and the morning were the third day.

And God said, Let there be lights in the firmament of the heaven to divide the day from the night; and let them be for signs, and for seasons, and for days, and years: and let them be for lights in the firmament of the heaven to give light upon the earth: and it was so. And God made two great lights; the greater light to rule the day, and the lesser light to rule the night: he made the stars also. And God set them in the firmament of the heaven to give light upon the earth, and to rule over the day and over the night, and to divide the light from the darkness: and God saw that it was good. And the evening and the morning were the fourth day.

And God said, Let the waters bring forth abundantly the moving creature that hath life, and fowl that may fly above the earth in the open firmament of heaven. And God created great whales, and every living creature that moveth, which the waters brought forth abundantly, after their kind, and every winged fowl after his kind: and God saw that it was good. And God blessed them, saying, Be fruitful, and multiply, and fill the waters in the seas, and let fowl multiply in the earth. And the evening and the morning were the fifth day.

And God said, Let the earth bring forth the living creature after his kind, cattle, and creeping thing, and beast of the earth after his kind: and it was so. And God made the beast of the earth after his kind, and cattle after their kind, and every thing that creepeth upon the earth after his kind: and God saw that it was good.

And God said, Let us make man in our image, after our likeness: and let them have dominion over the fish of the sea, and over the fowl of the air, and over the cattle, and over all the earth, and over every creeping thing that creepeth upon the earth. So God created man in his own image, in the image of God created he him; male and female created he them. And God blessed them, and God said unto them, Be fruitful, and multiply, and replenish the earth, and subdue it: and have dominion over the fish of the sea, and over the fowl of the air, and over every living thing that moveth upon the earth.

And God said, Behold, I have given you every herb bearing seed, which is upon the face of all the earth, and every tree, in the which is the fruit of a tree yielding seed; to you it shall be for meat. And to every beast of the earth, and to every fowl of the air, and to every thing that creepeth upon the earth, wherein there is life, I have given every green herb for meat: and it was so. And God saw every thing that he had made, and, behold, it was very good. And the evening and the morning were the sixth day.

Thus the heavens and the earth were finished, and all the host of them. And on the seventh day God ended his work which he had made; and he rested on the seventh day from all his work which he had made.

\begin{rubric}
	\emph{Thanks be to God} is not said either here or after the other Prophecies.
\end{rubric}

\textsc{Priest.} Let us pray.

\textsc{Deacon.} Let us bow the knee.

\textsc{Subdeacon.} Arise.

\lett{O}{God,} who didst wonderfully create man, and hast more wonderfully redeemed him: grant us, we beseech thee, such strength of mind to withstand the enticements of sin; that we may be found worthy to attain to everlasting joys. Through Jesus Christ thy Son our Lord. Who liveth.

\readingcitation{Prophecy II}{Genesis 5:32}
\lett{A}{nd} Noah was five hundred years old: and Noah begat Shem, Ham, and Japheth. And it came to pass, when men began to multiply on the face of the earth, and daughters were born unto them, that the sons of God saw the daughters of men that they were fair; and they took them wives of all which they chose. And the Lord said, My spirit shall not always strive with man, for that he also is flesh: yet his days shall be an hundred and twenty years. There were giants in the earth in those days; and also after that, when the sons of God came in unto the daughters of men, and they bare children to them, the same became mighty men which were of old, men of renown.

And God saw that the wickedness of man was great in the earth, and that every imagination of the thoughts of his heart was only evil continually. And it repented the Lord that he had made man on the earth, and it grieved him at his heart. And the Lord said, I will destroy man whom I have created from the face of the earth; both man, and beast, and the creeping thing, and the fowls of the air; for it repenteth me that I have made them. But Noah found grace in the eyes of the Lord.

These are the generations of Noah: Noah was a just man and perfect in his generations, and Noah walked with God. And Noah begat three sons, Shem, Ham, and Japheth. The earth also was corrupt before God, and the earth was filled with violence. And God looked upon the earth, and, behold, it was corrupt; for all flesh had corrupted his way upon the earth. And God said unto Noah, The end of all flesh is come before me; for the earth is filled with violence through them; and, behold, I will destroy them with the earth. Make thee an ark of gopher wood; rooms shalt thou make in the ark, and shalt pitch it within and without with pitch. And this is the fashion which thou shalt make it of: The length of the ark shall be three hundred cubits, the breadth of it fifty cubits, and the height of it thirty cubits. A window shalt thou make to the ark, and in a cubit shalt thou finish it above; and the door of the ark shalt thou set in the side thereof; with lower, second, and third stories shalt thou make it. And, behold, I, even I, do bring a flood of waters upon the earth, to destroy all flesh, wherein is the breath of life, from under heaven; and every thing that is in the earth shall die. But with thee will I establish my covenant; and thou shalt come into the ark, thou, and thy sons, and thy wife, and thy sons' wives with thee. And of every living thing of all flesh, two of every sort shalt thou bring into the ark, to keep them alive with thee; they shall be male and female. Of fowls after their kind, and of cattle after their kind, of every creeping thing of the earth after his kind, two of every sort shall come unto thee, to keep them alive. And take thou unto thee of all food that is eaten, and thou shalt gather it to thee; and it shall be for food for thee, and for them. Thus did Noah; according to all that God commanded him, so did he. 

And Noah was six hundred years old when the flood of waters was upon the earth. The same day were all the fountains of the great deep broken up, and the windows of heaven were opened. And the rain was upon the earth forty days and forty nights. In the selfsame day entered Noah, and Shem, and Ham, and Japheth, the sons of Noah, and Noah's wife, and the three wives of his sons with them, into the ark; they, and every beast after his kind, and all the cattle after their kind, and every creeping thing that creepeth upon the earth after his kind, and every fowl after his kind, every bird of every sort. And the ark went upon the face of the waters. And the waters prevailed exceedingly upon the earth; and all the high hills, that were under the whole heaven, were covered. Fifteen cubits upward did the waters prevail; and the mountains were covered. And all flesh died that moved upon the earth, both of fowl, and of cattle, and of beast, and of every creeping thing that creepeth upon the earth. And Noah only remained alive, and they that were with him in the ark. And the waters prevailed upon the earth an hundred and fifty days.

And God remembered Noah, and every living thing, and all the cattle that was with him in the ark: and God made a wind to pass over the earth, and the waters asswaged; the fountains also of the deep and the windows of heaven were stopped, and the rain from heaven was restrained; and the waters returned from off the earth continually: and after the end of the hundred and fifty days the waters were abated.

And it came to pass at the end of forty days, that Noah opened the window of the ark which he had made: and he sent forth a raven, which went forth to and fro, until the waters were dried up from off the earth. Also he sent forth a dove from him, to see if the waters were abated from off the face of the ground; but the dove found no rest for the sole of her foot, and she returned unto him into the ark, for the waters were on the face of the whole earth: then he put forth his hand, and took her, and pulled her in unto him into the ark. And he stayed yet other seven days; and again he sent forth the dove out of the ark; and the dove came in to him in the evening; and, lo, in her mouth was an olive leaf pluckt off: so Noah knew that the waters were abated from off the earth. And he stayed yet other seven days; and sent forth the dove; which returned not again unto him any more.

And God spake unto Noah, saying, Go forth of the ark, thou, and thy wife, and thy sons, and thy sons' wives with thee. Bring forth with thee every living thing that is with thee, of all flesh, both of fowl, and of cattle, and of every creeping thing that creepeth upon the earth; that they may breed abundantly in the earth, and be fruitful, and multiply upon the earth. And Noah went forth, and his sons, and his wife, and his sons' wives with him: every beast, every creeping thing, and every fowl, and whatsoever creepeth upon the earth, after their kinds, went forth out of the ark.

And Noah builded an altar unto the Lord; and took of every clean beast, and of every clean fowl, and offered burnt offerings on the altar. And the Lord smelled a sweet savour.

\textsc{Priest.} Let us pray.

\textsc{Deacon.} Let us bow the knee.

\textsc{Subdeacon.} Arise.

\lett{O}{God,} who art strength unchangeable and light eternal: look down in mercy on the wondrous mystery of thy whole Church, and by the operation of thy continual providence accomplish in tranquillity the work of man's salvation; and let the whole world perceive and know that things cast down are being raised up, things grown old are being made new, and all things are returning to their perfection, through him from whom they took their beginning: Jesus Christ thy Son our Lord: Who liveth and reigneth with thee.

\readingcitation{Prophecy III}{Genesis 22:1}
\lett{I}{n those days:} God did tempt Abraham, and said unto him, Abraham: and he said, Behold, here I am. And he said, Take now thy son, thine only son Isaac, whom thou lovest, and get thee into the land of Moriah; and offer him there for a burnt offering upon one of the mountains which I will tell thee of.

And Abraham rose up early in the morning, and saddled his ass, and took two of his young men with him, and Isaac his son, and clave the wood for the burnt offering, and rose up, and went unto the place of which God had told him. Then on the third day Abraham lifted up his eyes, and saw the place afar off. And Abraham said unto his young men, Abide ye here with the ass; and I and the lad will go yonder and worship, and come again to you. And Abraham took the wood of the burnt offering, and laid it upon Isaac his son; and he took the fire in his hand, and a knife; and they went both of them together. And Isaac spake unto Abraham his father, and said, My father: and he said, Here am I, my son. And he said, Behold the fire and the wood: but where is the lamb for a burnt offering? And Abraham said, My son, God will provide himself a lamb for a burnt offering: so they went both of them together. And they came to the place which God had told him of; and Abraham built an altar there, and laid the wood in order, and bound Isaac his son, and laid him on the altar upon the wood. And Abraham stretched forth his hand, and took the knife to slay his son. And the angel of the Lord called unto him out of heaven, and said, Abraham, Abraham: and he said, Here am I. And he said, Lay not thine hand upon the lad, neither do thou any thing unto him: for now I know that thou fearest God, seeing thou hast not withheld thy son, thine only son from me. And Abraham lifted up his eyes, and looked, and behold behind him a ram caught in a thicket by his horns: and Abraham went and took the ram, and offered him up for a burnt offering in the stead of his son. And Abraham called the name of that place Jehovah-jireh: as it is said to this day, In the mount of the Lord it shall be seen.

And the angel of the Lord called unto Abraham out of heaven the second time, and said, By myself have I sworn, saith the Lord, for because thou hast done this thing, and hast not withheld thy son, thine only son: that in blessing I will bless thee, and in multiplying I will multiply thy seed as the stars of the heaven, and as the sand which is upon the sea shore; and thy seed shall possess the gate of his enemies; and in thy seed shall all the nations of the earth be blessed; because thou hast obeyed my voice. So Abraham returned unto his young men, and they rose up and went together to Beer-sheba; and Abraham dwelt at Beer-sheba.

\textsc{Priest.} Let us pray.

\textsc{Deacon.} Let us bow the knee.

\textsc{Subdeacon.} Arise.

\lett{O}{God,} the sovereign Father of the faithful, who dost spread abroad throughout all the world the grace of adoption, to the increasing of the sons of thy promise: and, according as thou didst swear unto thy servant Abraham, dost through this paschal sacrament make him the father of all nations; grant that thy peoples may worthily enter upon the grace of thy calling. Through.

\readingcitation{Prophecy IV}{Exodus 14:24}
\lett{I}{n those days:} It came to pass, that in the morning watch the Lord looked unto the host of the Egyptians through the pillar of fire and of the cloud, and troubled the host of the Egyptians, and took off their chariot wheels, that they drave them heavily: so that the Egyptians said, Let us flee from the face of Israel; for the Lord fighteth for them against the Egyptians.

 And the Lord said unto Moses, Stretch out thine hand over the sea, that the waters may come again upon the Egyptians, upon their chariots, and upon their horsemen. And Moses stretched forth his hand over the sea, and the sea returned to his strength when the morning appeared; and the Egyptians fled against it; and the Lord overthrew the Egyptians in the midst of the sea. And the waters returned, and covered the chariots, and the horsemen, and all the host of Pharaoh that came into the sea after them; there remained not so much as one of them. But the children of Israel walked upon dry land in the midst of the sea; and the waters were a wall unto them on their right hand, and on their left. Thus the Lord saved Israel that day out of the hand of the Egyptians; and Israel saw the Egyptians dead upon the sea shore. And Israel saw that great work which the Lord did upon the Egyptians: and the people feared the Lord, and believed the Lord, and his servant Moses.

Then sang Moses and the children of Israel this song unto the Lord, and spake, saying,

\tract{We will sing unto the Lord: for he hath triumphed gloriously: the horse and his rider hath he thrown into the sea: the Lord is my strength and song, and he is become my salvation. ℣. He is my God, and I will praise him: my father's God, and I will exalt him. ℣. The Lord is a man of war: the Lord is his name.}

\textsc{Priest.} Let us pray.

\textsc{Deacon.} Let us bow the knee.

\textsc{Subdeacon.} Arise.

\lett{O}{God,} whose miracles of old we perceive to shine forth even in our times: who by the water of regeneration dost now work for the salvation of the nations, as thou didst deliver one people from the pursuit of the Egyptians by the power of thy right hand: vouchsafe; that the fulness of the whole world may be numbered among the sons of Abraham and made partakers of the dignity of Israel. Through.

\readingcitation{Prophecy V}{Isaiah 54:17}
\lett{T}{his} is the heritage of the servants of the \divineName{Lord}: and their righteousness is of me, saith the Lord.

\secondline{Ho, every one that thirsteth, come ye to the waters,
and he that hath no money; come ye, buy, and eat;
yea, come, buy wine and milk without money and without price.
Wherefore do ye spend money for that which is not bread?
and your labour for that which satisfieth not?
hearken diligently unto me, and eat ye that which is good,
and let your soul delight itself in fatness.
Incline your ear, and come unto me:
hear, and your soul shall live;
and I will make an everlasting covenant with you,
even the sure mercies of David.
Behold, I have given him for a witness to the people,
a leader and commander to the people.
Behold, thou shalt call a nation that thou knowest not,
and nations that knew not thee shall run unto thee
because of the Lord thy God, and for the Holy One of Israel;
for he hath glorified thee.}

Seek ye the Lord while he may be found,
call ye upon him while he is near:
let the wicked forsake his way,
and the unrighteous man his thoughts:
and let him return unto the Lord,
and he will have mercy upon him;
and to our God, for he will abundantly pardon.

For my thoughts are not your thoughts,
neither are your ways my ways, saith the Lord.
For as the heavens are higher than the earth,
so are my ways higher than your ways,
and my thoughts than your thoughts.
For as the rain cometh down, and the snow from heaven,
and returneth not thither, but watereth the earth,
and maketh it bring forth and bud,
that it may give seed to the sower, and bread to the eater:
so shall my word be that goeth forth out of my mouth:
it shall not return unto me void,
but it shall accomplish that which I please,
and it shall prosper in the thing whereto I sent it, saith the Lord almighty.

\textsc{Priest.} Let us pray.

\textsc{Deacon.} Let us bow the knee.

\textsc{Subdeacon.} Arise.

\lett{A}{lmighty} and everlasting God, multiply to the honour of thy name what thou didst swear to the faith of the fathers: and spread abroad the sons of promise by holy adoption; that as the Saints of old doubted not thine oath, so thy Church may now see in abundant measure the fulfilment of the same. Through.

\readingcitation{Prophecy VI}{Baruch 3:9}

\lett{H}{ear,} O Israel, the commandments of life: give ear to understand wisdom. How happeneth it Israel, that thou art in thine enemies' land, that thou art waxen old in a strange country, that thou art defiled with the dead, That thou art counted with them that go down into the grave? Thou hast forsaken the fountain of wisdom. For if thou hadst walked in the way of God, thou shouldest have dwelled in peace for ever. Learn where is wisdom, where is strength, where is understanding; that thou mayest know also where is length of days, and life, where is the light of the eyes, and peace.

Who hath found out her place? or who hath come into her treasures? Where are the princes of the heathen become, and such as ruled the beasts upon the earth; They that had their pastime with the fowls of the air, and they that hoarded up silver and gold, wherein men trust, and made no end of their getting? For they that wrought in silver, and were so careful, and whose works are unsearchable, They are vanished and gone down to the grave, and others are come up in their steads. Young men have seen light, and dwelt upon the earth: but the way of knowledge have they not known, Nor understood the paths thereof, nor laid hold of it: their children were far off from that way. It hath not been heard of in Chanaan, neither hath it been seen in Theman. The Agarenes that seek wisdom upon earth, the merchants of Meran and of Theman, the authors of fables, and searchers out of understanding; none of these have known the way of wisdom, or remember her paths.

O Israel, how great is the house of God! and how large is the place of his possession! Great, and hath none end; high, and unmeasurable. There were the giants famous from the beginning, that were of so great stature, and so expert in war. Those did not the Lord choose, neither gave he the way of knowledge unto them: But they were destroyed, because they had no wisdom, and perished through their own foolishness. Who hath gone up into heaven, and taken her, and brought her down from the clouds? Who hath gone over the sea, and found her, and will bring her for pure gold? No man knoweth her way, nor thinketh of her path.

But he that knoweth all things knoweth her, and hath found her out with his understanding: he that prepared the earth for evermore hath filled it with fourfooted beasts: He that sendeth forth light, and it goeth, calleth it again, and it obeyeth him with fear. The stars shined in their watches, and rejoiced: when he calleth them, they say, Here we be; and so with cheerfulness they shewed light unto him that made them. This is our God, and there shall none other be accounted of in comparison of him. He hath found out all the way of knowledge, and hath given it unto Jacob his servant, and to Israel his beloved. Afterward did he shew himself upon earth, and conversed with men.

\textsc{Priest.} Let us pray.

\textsc{Deacon.} Let us bow the knee.

\textsc{Subdeacon.} Arise.

\lett{O}{God,} who by the calling of the nations dost ever multiply thy Church: mercifully grant; that they whom thou dost cleanse by the waters of baptism may be defended by thy continual protection. Through.

\readingcitation{Prophecy VII}{Ezekiel 37:1}
\lett{I}{n those days:} The hand of the \divineName{Lord} was upon me, and carried me out in the spirit of the Lord, and set me down in the midst of the valley which was full of bones, and caused me to pass by them round about: and, behold, there were very many in the open valley; and, lo, they were very dry. And he said unto me, Son of man, can these bones live? And I answered, O Lord God, thou knowest.

Again he said unto me, Prophesy upon these bones, and say unto them, O ye dry bones, hear the word of the Lord. Thus saith the Lord God unto these bones; Behold, I will cause breath to enter into you, and ye shall live: and I will lay sinews upon you, and will bring up flesh upon you, and cover you with skin, and put breath in you, and ye shall live; and ye shall know that I am the Lord.

So I prophesied as I was commanded: and as I prophesied, there was a noise, and behold a shaking, and the bones came together, bone to his bone. And when I beheld, lo, the sinews and the flesh came up upon them, and the skin covered them above: but there was no breath in them. Then said he unto me, Prophesy unto the wind, prophesy, son of man, and say to the wind, Thus saith the Lord God; Come from the four winds, O breath, and breathe upon these slain, that they may live. So I prophesied as he commanded me, and the breath came into them, and they lived, and stood up upon their feet, an exceeding great army.

Then he said unto me, Son of man, these bones are the whole house of Israel: behold, they say, Our bones are dried, and our hope is lost: we are cut off for our parts. Therefore prophesy and say unto them, Thus saith the Lord God; Behold, O my people, I will open your graves, and cause you to come up out of your graves, and bring you into the land of Israel. And ye shall know that I am the Lord, when I have opened your graves, O my people, and brought you up out of your graves, and shall put my spirit in you, and ye shall live, and I shall place you in your own land, faith the Lord almighty.

\textsc{Priest.} Let us pray.

\textsc{Deacon.} Let us bow the knee.

\textsc{Subdeacon.} Arise.

\lett{O}{God,} who dost instruct us in the pages of either Testament in the celebrating of the paschal sacrament: grant unto us so to understand thy mercy; that, through the receiving of these present gifts, we may be established in the expectation of those which are to come. Through.

\readingcitation{Prophecy VIII}{Isaiah 4:1}
\lett{I}{n that day} seven women shall take hold of one man, saying: We will eat our own bread, and wear our own apparel: only let us be called by thy name, to take away our reproach. In that day shall the branch of the Lord be beautiful and glorious, and the fruit of the earth shall be excellent and comely for them that are escaped of Israel. And it shall come to pass, that he that is left in Zion, and he that remaineth in Jerusalem, shall be called holy, even every one that is written among the living in Jerusalem: when the Lord shall have washed away the filth of the daughters of Zion, and shall have purged the blood of Jerusalem from the midst thereof by the spirit of judgment, and by the spirit of burning. And the Lord will create upon every dwelling place of mount Zion, and upon her assemblies, a cloud and smoke by day, and the shining of a flaming fire by night: for upon all the glory shall be a defence. And there shall be a tabernacle for a shadow in the daytime from the heat, and for a place of refuge, and for a covert from storm and from rain.

\tract{My well-beloved hath a vineyard in a very fruitful hill. ℣. And he fenced it, and gathered out the stones thereof: and planted it with the choicest vine, and built a tower in the midst of it. ℣. And also made a wine-press therein: for the vineyard of the Lord of hosts is the house of Israel.}

\textsc{Priest.} Let us pray.

\textsc{Deacon.} Let us bow the knee.

\textsc{Subdeacon.} Arise.

\lett{O}{God,} who by the mouth of thy holy prophets hast manifested thyself in all the children of thy Church, to be in all places of thy dominion the sower of good seed and the husbandman of thine elect branches: grant unto thy peoples, who are named thy vineyard and thy harvest field; that being purged from the filthiness of thorns and briars, they may be made to bring forth worthy fruit in abundance. Through.

\readingcitation{Prophecy IX}{Exodus 12:1}
\lett{I}{n those days:} The \divineName{Lord} spake unto Moses and Aaron in the land of Egypt, saying, This month shall be unto you the beginning of months: it shall be the first month of the year to you.

Speak ye unto all the congregation of Israel, saying, In the tenth day of this month they shall take to them every man a lamb, according to the house of their fathers, a lamb for an house: and if the household be too little for the lamb, let him and his neighbour next unto his house take it according to the number of the souls; every man according to his eating shall make your count for the lamb. Your lamb shall be without blemish, a male of the first year: ye shall take it out from the sheep, or from the goats: and ye shall keep it up until the fourteenth day of the same month: and the whole assembly of the congregation of Israel shall kill it in the evening. And they shall take of the blood, and strike it on the two side posts and on the upper door post of the houses, wherein they shall eat it. And they shall eat the flesh in that night, roast with fire, and unleavened bread; and with bitter herbs they shall eat it. Eat not of it raw, nor sodden at all with water, but roast with fire; his head with his legs, and with the purtenance thereof. And ye shall let nothing of it remain until the morning; and that which remaineth of it until the morning ye shall burn with fire.

And thus shall ye eat it; with your loins girded, your shoes on your feet, and your staff in your hand; and ye shall eat it in haste: it is the Lord’s passover.

\textsc{Priest.} Let us pray.

\textsc{Deacon.} Let us bow the knee.

\textsc{Subdeacon.} Arise.

\lett{A}{lmighty} and everlasting God, who art wondrous in the dispensation of all thy works: let thy redeemed understand that the creation of the world in the beginning was not a work more excellent than the sacrifice, in the end of time, of Christ our Passover: Who liveth and reigneth with thee.

\readingcitation{Prophecy X}{Jonah 3:1}
\lett{I}{n those days:} The word of the Lord came unto Jonah the Prophet the second time, saying: Arise, go unto Nineveh, that great city, and preach unto it the preaching that I bid thee. So Jonah arose, and went unto Nineveh, according to the word of the Lord. Now Nineveh was an exceeding great city of three days’ journey. And Jonah began to enter into the city a day’s journey, and he cried, and said, Yet forty days, and Nineveh shall be overthrown.

So the people of Nineveh believed God, and proclaimed a fast, and put on sackcloth, from the greatest of them even to the least of them. For word came unto the king of Nineveh, and he arose from his throne, and he laid his robe from him, and covered him with sackcloth, and sat in ashes. And he caused it to be proclaimed and published through Nineveh by the decree of the king and his nobles, saying, Let neither man nor beast, herd nor flock, taste any thing: let them not feed, nor drink water: but let man and beast be covered with sackcloth, and cry mightily unto God: yea, let them turn every one from his evil way, and from the violence that is in their hands. Who can tell if God will turn and repent, and turn away from his fierce anger, that we perish not?

And God saw their works, that they turned from their evil way; and he had mercy on his people, even he the Lord, our God.

\textsc{Priest.} Let us pray.

\textsc{Deacon.} Let us bow the knee.

\textsc{Subdeacon.} Arise.

\lett{O}{God,} who hast united the diversity of nations in the confession of thy name: grant us both to will, and to have power to do the things which thou commandest; that the people whom thou hast called to everlasting life may agree in unity of faith and in godliness of conversation. Through.

\readingcitation{Prophecy XI}{Deuteronomy 31:22}
\lett{I}{n those days:} Moses wrote this song, and taught it the children of Israel. And he gave Joshua the son of Nun a charge, and said, Be strong and of a good courage: for thou shalt bring the children of Israel into the land which I sware unto them: and I will be with thee.

And it came to pass, when Moses had made an end of writing the words of this law in a book, until they were finished, that Moses commanded the Levites, which bare the ark of the covenant of the Lord, saying, Take this book of the law, and put it in the side of the ark of the covenant of the Lord your God, that it may be there for a witness against thee. For I know thy rebellion, and thy stiff neck: behold, while I am yet alive with you this day, ye have been rebellious against the Lord; and how much more after my death?

Gather unto me all the elders of your tribes, and your officers, that I may speak these words in their ears, and call heaven and earth to record against them. For I know that after my death ye will utterly corrupt yourselves, and turn aside from the way which I have commanded you; and evil will befall you in the latter days; because ye will do evil in the sight of the Lord, to provoke him to anger through the work of your hands.

And Moses spake in the ears of all the congregation of Israel the words of this song, until they were ended.

\tract{Give ear, O ye heavens, and I will speak: and hear, O earth, the words of my mouth. ℣. My doctrine shall drop as the rain: my speech shall distil as the dew. ℣. As the small rain upon the tender herb, and as the showers upon the grass: because I will publish the name of the Lord: ℣. Ascribe ye greatness unto our God: he is the rock, his work is perfect, for all his ways are judgment. ℣. A God of truth and without iniquity: just and right is the Lord.}

\textsc{Priest.} Let us pray.

\textsc{Deacon.} Let us bow the knee.

\textsc{Subdeacon.} Arise.

\lett{O}{God,} the exaltation of the humble and the strength of them that do stand, who by thy holy servant Moses didst vouchsafe so to instruct thy people in the chanting of thy sacred song, that the repeating of the law by them might likewise be for our instruction: stir up thy power among all the fulness of the nations whom thou hast justified, and grant us gladness, assuaging fear; that the sins of all may be blotted out by thy forgiveness, and that which was denounced in vengeance may turn to our salvation. Through.

\readingcitation{Prophecy XII}{Daniel 3:1}
\lett{I}{n those days:} Nebuchadnezzar the king made an image of gold, whose height was threescore cubits, and the breadth thereof six cubits: he set it up in the plain of Dura, in the province of Babylon. Then Nebuchadnezzar the king sent to gather together the princes, the governors, and the captains, the judges, the treasurers, the counsellors, the sheriffs, and all the rulers of the provinces, to come to the dedication of the image which Nebuchadnezzar the king had set up. Then the princes, the governors, and captains, the judges, the treasurers, the counsellors, the sheriffs, and all the rulers of the provinces, were gathered together unto the dedication of the image that Nebuchadnezzar the king had set up; and they stood before the image that Nebuchadnezzar had set up. Then an herald cried aloud, To you it is commanded, O people, nations, and languages, that at what time ye hear the sound of the cornet, flute, harp, sackbut, psaltery, dulcimer, and all kinds of musick, ye fall down and worship the golden image that Nebuchadnezzar the king hath set up: and whoso falleth not down and worshippeth shall the same hour be cast into the midst of a burning fiery furnace. Therefore at that time, when all the people heard the sound of the cornet, flute, harp, sackbut, psaltery, and all kinds of musick, all the people, the nations, and the languages, fell down and worshipped the golden image that Nebuchadnezzar the king had set up.

Wherefore at that time certain Chaldeans came near, and accused the Jews. They spake and said to the king Nebuchadnezzar, O king, live for ever. Thou, O king, hast made a decree, that every man that shall hear the sound of the cornet, flute, harp, sackbut, psaltery, and dulcimer, and all kinds of musick, shall fall down and worship the golden image: and whoso falleth not down and worshippeth, that he should be cast into the midst of a burning fiery furnace. There are certain Jews whom thou hast set over the affairs of the province of Babylon, Shadrach, Meshach, and Abed-nego; these men, O king, have not regarded thee: they serve not thy gods, nor worship the golden image which thou hast set up.

Then Nebuchadnezzar in his rage and fury commanded to bring Shadrach, Meshach, and Abed-nego. Then they brought these men before the king. Nebuchadnezzar spake and said unto them, Is it true, O Shadrach, Meshach, and Abed-nego, do not ye serve my gods, nor worship the golden image which I have set up? Now if ye be ready that at what time ye hear the sound of the cornet, flute, harp, sackbut, psaltery, and dulcimer, and all kinds of musick, ye fall down and worship the image which I have made; well: but if ye worship not, ye shall be cast the same hour into the midst of a burning fiery furnace; and who is that God that shall deliver you out of my hands? Shadrach, Meshach, and Abed-nego, answered and said to the king, O Nebuchadnezzar, we are not careful to answer thee in this matter. If it be so, our God whom we serve is able to deliver us from the burning fiery furnace, and he will deliver us out of thine hand, O king. But if not, be it known unto thee, O king, that we will not serve thy gods, nor worship the golden image which thou hast set up.

Then was Nebuchadnezzar full of fury, and the form of his visage was changed against Shadrach, Meshach, and Abed-nego: therefore he spake, and commanded that they should heat the furnace one seven times more than it was wont to be heated. And he commanded the most mighty men that were in his army to bind Shadrach, Meshach, and Abed-nego, and to cast them into the burning fiery furnace. Then these men were bound in their coats, their hosen, and their hats, and their other garments, and were cast into the midst of the burning fiery furnace. Therefore because the king’s commandment was urgent, and the furnace exceeding hot, the flame of the fire slew those men that took up Shadrach, Meshach, and Abed-nego. And these three men, Shadrach, Meshach, and Abed-nego, fell down bound into the midst of the burning fiery furnace. And they walked in the midst of the fire, praising God, and blessing the Lord.

\letuspray
\lett{A}{lmighty} and everlasting God, the only hope of the world, who, by the preaching of thy Prophets hast declared the mysteries of these present times: mercifully increase the devotion of thy people; forasmuch as without thine inspiration thy faithful can in no wise advance in virtue. Through.

    \par\noindent
	\centerline{\rule{0.25\textwidth}{0.4pt}}
	\par\noindent

\subbysub{Blessing of the Baptismal Font}\label{EasterFontBlessing}

\begin{rubric}
	These things being finished, if the church has a baptismal Font, the Priest who is to bless the Font puts on a violet Cope, and preceded by the Cross with lights and the blessed Candle lighted, descends with the Clergy and Ministers, vested, to the Font: and meanwhile the following is sung:
\end{rubric}

\tract{Like as the hart desireth the water-brooks: so longeth my soul after thee, O God. ℣. My soul is athirst for God, yea, even for the living God: when shall I come to appear before the presence of God? ℣. My tears have been my meat day and night, while they daily say unto me: Where is now thy God?}

\begin{rubric}
	Then the Priest, before he enters for the blessing of the Font, says this Prayer near by the Font:
\end{rubric}

\elcol{℣. The Lord be with you.}{℣. Dóminus vobíscum.}
\elcol{℟. And with thy spirit.}{℟. Et cum spíritu tuo.}

\letuspray
\lett{A}{lmighty} and everlasting God, mercifully look upon the devotion of this people called to a new birth, who, like the hart, seek the fountain of thy waters: and mercifully grant; that the thirst of their faith may, by the mystery of Baptism, sanctify them in body and soul. Through.

\begin{rubric}
	Then he proceeds to the blessing of the Font, saying:
\end{rubric}

\elcol{℣. The Lord be with you.}{℣. Dóminus vobíscum.}
\elcol{℟. And with thy spirit.}{℟. Et cum spíritu tuo.}

\letuspray
\lett{A}{lmighty} and everlasting God, be present at the mysteries, be present at the sacraments of thy great goodness: and send forth the spirit of adoption for the regenerating of the new peoples whom the font of baptism doth bring forth unto thee; that what is to be done by our humble ministry, may be effectually fulfilled by thy power. Through Jesus Christ, thy Son, our Lord: Who liveth and reigneth with thee in the unity of the Holy Ghost, God.

\begin{rubric}
	Then raising his voice to the tone of a Preface, he proceeds with joined hands:
\end{rubric}

Throughout all ages, world without end.

℟. Amen.

℣. The Lord be with you.

℟. And with thy spirit.

℣. Lift up your hearts.

℟. We lift them up unto the Lord.

℣. Let us give thanks unto our Lord God.

℟. It is meet and right so to do.
\par\noindent
It is very meet, right, and our bounden duty, that we should at all times and in all places give thanks unto thee: O Lord holy Father almighty, everlasting God: Who by thy invisible power dost wondrously give effect to the work of thy sacraments: And though we be unworthy to perform mysteries so great: Yet thou dost not leave us destitute of the gifts of thy grace, but mercifully inclinest thine ears even unto these our supplications. O God, whose Spirit, in the first beginnings of the world, moved over the waters: that even then the nature of water might conceive the virtue of sanctification. O God, who didst wash away in the waters the iniquities of a sinful world, and didst even in the outpouring of the flood shew forth in a figure our regeneration: that by the mystery of this same element, there should be an end to sin and likewise a beginning of virtue. Look, O Lord, upon the face of thy Church, and multiply in her the power of thy regeneration, for thou dost by the streams of thine abundant grace make glad thy city : and openest the fount of baptism throughout the whole world for the renewing of the nations: that, by the command of thy majesty, she may receive from the Holy Ghost the grace of thine only begotten Son.

\begin{rubric}
	Here the Priest divides the water in the form of a cross with his extended hand, which wipes at once with a linen cloth, saying:
\end{rubric}
\par\noindent
May he, by the secret mingling of his divine power, make fruitful this water prepared for the regeneration of mankind: that having received sanctification, and being born again a new creature, there may come forth an offspring of heaven from the spotless womb of this divine fountain: And that all, whether sex divide them in body, or age in time, may alike be brought forth into one childhood by grace, their mother. Far hence, O Lord, at thy command let every unclean spirit depart: far hence be all the wickedness of the craft of the devil. Let no power of the enemy here have place or portion: let it not ensnare and compass us around: let it not waylay us in secret: let it not taint us with corruption.

\begin{rubric}
	He touches the water with his hand.
\end{rubric}
\par\noindent
May this holy and undefiled creature be free from every assault of the enemy, and purified by the departure of all iniquity. May it be a living fountain, water that doth regenerate, a purifying stream: that all who shall be washed in this laver of salvation may, by the operation of the Holy Ghost within them, obtain the grace of perfect purification.

\begin{rubric}
	He makes three crosses over the Font, saying:
\end{rubric}
\par\noindent
Wherefore I bless thee, O creature of water, by the living {\ding{64}} God, by the true {\ding{64}} God, by the holy {\ding{64}} God: by God who in the beginning through his word divided thee from the dry land: whose spirit moved upon thee.

\begin{rubric}
	Here he divides the water with his hand, and scatters it towards the four quarters of the world, saying:
\end{rubric}
\par\noindent
Who made thee to flow from the fountain of paradise, and commanded thee to water the whole earth with thy tour rivers. Who in the desert bestowed upon thee sweetness when thou wast bitter, that men might drink, and brought thee forth from the rock for the thirsting people. I ble {\ding{64}} ss thee also by Jesus Christ his only Son, our Lord: who in Cana of Galilee by a wondrous miracle did change thee through his power into wine. Who walked upon thee with his feet: and was baptized in thee by John in Jordan. Who brought thee forth together with blood from his side: and commanded his disciples, that believing they should be baptized in thee, saying: Go ye, teach all nations, baptizing them in the name of the Father, and of the Son, and of the Holy Ghost.

\begin{rubric}
	He changes his voice, and proceeds in the tone of the Lesson:
\end{rubric}
\par\noindent
Do thou, almighty God, mercifully assist us, who keep these thy commandments: do thou graciously breathe upon us.

\begin{rubric}
	He breathes thrice upon the water in the form of a cross, saying:
\end{rubric}
\par\noindent
Do thou with thy mouth bless this element of water: that as by nature it hath power to cleanse and wash the body, so also it may be effectual for the purifying of the soul.

\begin{rubric}
	Here the Priest lowers the Candle into the water a little way: and resuming the tone of the Preface, says:
\end{rubric}
\par\noindent
May the power of the Holy Ghost descend upon the fulness of this font.

\begin{rubric}
	Then withdrawing the Candle from the water, he again dips it more deeply, and repeats in a somewhat higher voice, \emph{May the power}. Then he again withdraws the Candle from the water, and the third time dips it to the bottom, repeating in a still higher voice: \emph{May the power}, as above.
	
	Then breathing three times on the water in this form \textsc{Ψ} he proceeds:
\end{rubric}
\par\noindent
And make the whole substance of this water to be fruitful unto regeneration.

\begin{rubric}
	Here the Candle is taken out of the water, and he proceeds:
\end{rubric}
\par\noindent
Here may the stains of every sin be blotted out: here may nature, created in thine image, be restored to the honour of its first estate, and cleansed from all the defilement of its old condition: that every man who cometh to this sacrament of regeneration may be born again unto the new childhood of true innocence.

\begin{rubric}
	He says that which follows in the reading voice:
\end{rubric}
\par\noindent
Through Jesus Christ thy Son, our Lord: Who shall come to judge the quick and the dead, and the world by fire. \textit{Amen.}

\begin{rubric}
	Then some of the water that has been blessed is sprinkled by the assisting Priests over the People.
	
	And meantime one of the ministers of the church takes some of the water in a vessel, for sprinkling in houses and other places. Which done, the Priest who is blessing the Font, pours some of the Oil of Catechumens into the water in the form of a cross, saying in an audible voice:
\end{rubric}
\par\noindent
May this font be sanctified and made fruitful by the Oil of salvation, for such as shall be born again therefrom, unto life everlasting. \textit{Amen.}

\begin{rubric}
	Then, in the same manner, he pours in some of the Chrism, saying:
\end{rubric}
\par\noindent
May this inpouring of the Chrism of our Lord Jesus Christ, and of the Holy Ghost the Paraclete, be wrought in the name of the holy Trinity. \textit{Amen.}

\begin{rubric}
	Lastly he takes both phials of the said holy Oil and the Chrism and pours in some of both together in the form of a cross, saying:
\end{rubric}
\par\noindent
May this commingling of the Chrism of sanctification, and of the Oil of unction and of the Water of baptism be likewise wrought in the name of the Fa {\ding{64}} ther, and of the S {\ding{64}} on, and of the Holy {\ding{64}} Ghost. \textit{Amen.}

\begin{rubric}
	Then he mingles the Oil itself with the water, and with his hand scatters it over the whole Font. If there be any to be baptized, he baptizes them in the usual manner.
\end{rubric}

\begin{rubric}
	Then as the Priest and Ministers return to the Altar, the Litany is sung by two Cantors, and the Choir repeat it also together, as is directed below.
\end{rubric}

    \par\noindent
	\centerline{\rule{0.25\textwidth}{0.4pt}}
	\par\noindent

\subbysub{Paschal Litany}\label{PaschalLitany}

\begin{rubric}
	But where there is no baptismal Font, when the last Prophecy with its Collect is ended, the Celebrant lays aside his Chasuble, and with the Ministers lies prostrate before the Altar: and all kneeling, the Litany is chanted by two Cantors in the midst of the Choir, and both sides of the Choir respond together.
	
	But when they come to the \emph{℣. We sinners}, the Priest and his Ministers rise, and proceed to the sacristy, where they put on white vestments for the solemn Mass, and meanwhile the candles on the Altar are lighted.
	
	At the end of the Litany, \emph{Kyrie eleison} is sung solemnly, and repeated, as is usual at Mass.
\end{rubric}

\elcoldent{Lord, have mercy upon us.}{Kyrie, eléison.}
\elcol{℟. Christ, have mercy upon us.}{℟. Christe, eléison.}
\elcoldent{Lord, have mercy upon us.}{Kyrie, eléison.}
\elcoldent{O Christ hear us.}{Christe, audi nos.}
\elcol{℟. O Christ graciously hear us.}{℟. Christe, exáudi nos.}
\elcoldent{O God the Father of heaven,}{Pater de c{\ae}lis, Deus,}
\elcol{℟. Have mercy upon us.}{℟. Miserére nobis.}
\elcoldent{O God the Son, Redeemer of the world,}{Fili, Redémptor mundi, Deus,}
\elcol{℟. Have mercy upon us.}{℟. Miserére nobis.}
\elcoldent{O God the Holy Ghost,}{Spíritus Sancte, Deus,}
\elcol{℟. Have mercy upon us.}{℟. Miserére nobis.}
\elcoldent{Holy Trinity one God,}{Sancta Trínitas, unus Deus,}
\elcol{℟. Have mercy upon us.}{℟. Miserére nobis.}
\elcoldent{Holy Mary,}{Sancta María,}
\elcol{℟. Pray for us.}{℟. Ora pro nobis.}
\elcoldent{Holy Mother of God,}{Sancta Dei Génitrix,}
\elcol{℟. Pray for us.}{℟. Ora pro nobis.}
\elcoldent{Holy Virgin of virgins,}{Sancta Virgo vírginum,}
\elcol{℟. Pray for us.}{℟. Ora pro nobis.}
\elcoldent{Holy Michael,}{Sancte Micha{\"e}l,}
\elcol{℟. Pray for us.}{℟. Ora pro nobis.}
\elcoldent{Holy Gabriel,}{Sancte Gábri{\"e}l,}
\elcol{℟. Pray for us.}{℟. Ora pro nobis.}
\elcoldent{Holy Raphael,}{Sancte Rápha{\"e}l,}
\elcol{℟. Pray for us.}{℟. Ora pro nobis.}
\elcoldent{All ye holy Angels and Archangels,}{Omnes sancti Angeli et Archángeli,}
\elcol{℟. Pray for us.}{℟. Oráte pro nobis.}
\elcoldent{All ye holy orders of blessed Spirits,}{Omnes sancti beatórum Spírituum órdines,}
\elcol{℟. Pray for us.}{℟. Oráte pro nobis.}
\elcoldent{St. John Baptist,}{Sancte Joánnes Baptísta,}
\elcol{℟. Pray for us.}{℟. Ora pro nobis.}
\elcoldent{St. Joseph,}{Sancte Josephe,}
\elcol{℟. Pray for us.}{℟. Ora pro nobis.}
\elcoldent{All ye holy Patriarchs and Prophets,}{Omnes sancti Patriárch{\ae} et Prophét{\ae},}
\elcol{℟. Pray for us.}{℟. Oráte pro nobis.}
\elcoldent{St. Peter,}{Sancte Petre,}
\elcol{℟. Pray for us.}{℟. Ora pro nobis.}
\elcoldent{St. Paul,}{Sancte Paule,}
\elcol{℟. Pray for us.}{℟. Ora pro nobis.}
\elcoldent{St. Andrew,}{Sancte Andréa,}
\elcol{℟. Pray for us.}{℟. Ora pro nobis.}
\elcoldent{St. John,}{Sancte Joánnes,}
\elcol{℟. Pray for us.}{℟. Ora pro nobis.}
\elcoldent{All ye holy Apostles and Evangelists,}{Omnes sancti Apóstoli et Evangelíst{\ae},}
\elcol{℟. Pray for us.}{℟. Oráte pro nobis.}
\elcoldent{All ye holy Disciples of the Lord,}{Omnes sancti Discípuli Dómini,}
\elcol{℟. Pray for us.}{℟. Oráte pro nobis.}
\elcoldent{St. Stephen,}{Sancte Stéphane,}
\elcol{℟. Pray for us.}{℟. Ora pro nobis.}
\elcoldent{St. Lawrence,}{Sancte Laurénti,}
\elcol{℟. Pray for us.}{℟. Ora pro nobis.}
\elcoldent{St. Vincent,}{Sancte Vincénti,}
\elcol{℟. Pray for us.}{℟. Ora pro nobis.}
\elcoldent{All ye holy Martyrs,}{Omnes sancti Mártyres,}
\elcol{℟. Pray for us.}{℟. Oráte pro nobis.}
\elcoldent{St. Sylvester,}{Sancte Silvéster,}
\elcol{℟. Pray for us.}{℟. Ora pro nobis.}
\elcoldent{St. Gregory,}{Sancte Gregóri,}
\elcol{℟. Pray for us.}{℟. Ora pro nobis.}
\elcoldent{St. Augustine,}{Sancte Augustíne,}
\elcol{℟. Pray for us.}{℟. Ora pro nobis.}
\elcoldent{All ye holy Bishops and Confessors,}{Omnes sancti Pontífices et Confessóres,}
\elcol{℟. Pray for us.}{℟. Oráte pro nobis.}
\elcoldent{All ye holy Doctors,}{Omnes sancti Doctóres,}
\elcol{℟. Pray for us.}{℟. Oráte pro nobis.}
\elcoldent{Holy St. Anthony,}{Sancte Antóni,}
\elcol{℟. Pray for us.}{℟. Ora pro nobis.}
\elcoldent{Holy St. Benedict,}{Sancte Benedícte,}
\elcol{℟. Pray for us.}{℟. Ora pro nobis.}
\elcoldent{All ye holy Priests and Levites,}{Omnes sancti Sacerdótes et Levít{\ae},}
\elcol{℟. Pray for us.}{℟. Oráte pro nobis.}
\elcoldent{All ye holy Monks and Hermits,}{Omnes sancti Mónachi et Eremít{\ae},}
\elcol{℟. Pray for us.}{℟. Oráte pro nobis.}
\elcoldent{St. Mary Magdalene,}{Sancta María Magdaléna,}
\elcol{℟. Pray for us.}{℟. Ora pro nobis.}
\elcoldent{St. Agnes,}{Sancta Agnes,}
\elcol{℟. Pray for us.}{℟. Ora pro nobis.}
\elcoldent{St. Cecilia,}{Sancta C{\ae}cília,}
\elcol{℟. Pray for us.}{℟. Ora pro nobis.}
\elcoldent{St. Agatha,}{Sancta Agatha,}
\elcol{℟. Pray for us.}{℟. Ora pro nobis.}
\elcoldent{St. Anastasia,}{Sancta Anastásia,}
\elcol{℟. Pray for us.}{℟. Ora pro nobis.}
\elcoldent{All ye holy Virgins and Widows,}{Omnes sanct{\ae} Vírgines et Vídu{\ae},}
\elcol{℟. Pray for us.}{℟. Oráte pro nobis.}
\elcoldent{All ye Saints of God,}{Omnes Sancti et Sanct{\ae} Dei,}
\elcol{℟. Intercede for us.}{℟. Intercédite pro nobis.}
\elcoldent{Be thou merciful,}{Propítius esto,}
\elcol{℟. Spare us, O Lord.}{℟. Parce nobis, Dómine.}
\elcoldent{Be thou merciful,}{Propítius esto,}
\elcol{℟. Hear us, O Lord.}{℟. Exáudi nos, Dómine.}
\elcoldent{From all evil,}{Ab omni malo,}
\elcol{℟. Deliver us, O Lord.}{℟. Líbera nos, Dómine.}
\elcoldent{From all sin,}{Ab omni peccáto,}
\elcol{℟. Deliver us, O Lord.}{℟. Líbera nos, Dómine.}
\elcoldent{From everlasting death,}{A morte perpétua,}
\elcol{℟. Deliver us, O Lord.}{℟. Líbera nos, Dómine.}
\elcoldent{By the mystery of thy holy Incarnation,}{Per mystérium sanct{\ae} Incarnatiónis tu{\ae},}
\elcol{℟. Deliver us, O Lord.}{℟. Líbera nos, Dómine.}
\elcoldent{By thine advent,}{Per Advéntum tuum,}
\elcol{℟. Deliver us, O Lord.}{℟. Líbera nos, Dómine.}
\elcoldent{By thy nativity,}{Per Nativitátem tuam,}
\elcol{℟. Deliver us, O Lord.}{℟. Líbera nos, Dómine.}
\elcoldent{By thy baptism and holy fasting,}{Per Baptísmum et sanctum Jejúnium tuum,}
\elcol{℟. Deliver us, O Lord.}{℟. Líbera nos, Dómine.}
\elcoldent{By thy cross and passion,}{Per Crucem et Passiónem tuam,}
\elcol{℟. Deliver us, O Lord.}{℟. Líbera nos, Dómine.}
\elcoldent{By thy death and burial,}{Per Mortem et sepultúram tuam,}
\elcol{℟. Deliver us, O Lord.}{℟. Líbera nos, Dómine.}
\elcoldent{By thy holy resurrection,}{Per sanctam Resurrectiónem tuam,}
\elcol{℟. Deliver us, O Lord.}{℟. Líbera nos, Dómine.}
\elcoldent{By thy wonderful ascension,}{Per admirábilem Ascensiónem tuam,}
\elcol{℟. Deliver us, O Lord.}{℟. Líbera nos, Dómine.}
\elcoldent{By the coming of the Holy Ghost the Paraclete,}{Per advéntum Spíritus Sancti Parácliti,}
\elcol{℟. Deliver us, O Lord.}{℟. Líbera nos, Dómine.}
\elcoldent{In the day of judgment,}{In die judícii,}
\elcol{℟. Deliver us, O Lord.}{℟. Líbera nos, Dómine.}
\elcoldent{We sinners,}{Peccatóres,}
\elcol{℟. Beseech thee, hear us.}{℟. Te rogámus, audi nos.}
\elcoldent{That thou wouldest spare us,}{Ut nobis parcas,}
\elcol{℟. We beseech thee, hear us.}{℟. Te rogámus, audi nos.}
\elcoldent{That thou wouldest vouchsafe to govern and preserve thy holy Church,}{Ut Ecclésiam tuam sanctam régere et conserváre dignéris,}
\elcol{℟. We beseech thee, hear us.}{℟. Te rogamus, audi nos.}
\elcoldent{That thou wouldest vouchsafe to preserve our Apostolic lord and all orders of the Church in holy religion,}{Ut Domnum Apostólicum, et omnes ecclesiásticos órdines in sancta religióne conserváre dignéris,}
\elcol{℟. We beseech thee, hear us.}{℟. Te rogamus, audi nos.}
\elcoldent{That thou wouldest vouchsafe to humble the enemies of holy Church,}{Ut inimícos sanct{\ae} Ecclési{\ae} humiliáre dignéris,}
\elcol{℟. We beseech thee, hear us.}{℟. Te rogamus, audi nos.}
\elcoldent{That thou wouldest vouchsafe to give unto Christian kings and rulers peace and true concord,}{Ut régibus et princípibus christiánis pacem et veram concórdiam donáre dignéris,}
\elcol{℟. We beseech thee, hear us.}{℟. Te rogamus, audi nos.}
\elcoldent{That thou wouldest vouchsafe to strengthen, and preserve us in thy holy service,}{Ut nosmetípsos in tuo sancto servítio confortáre et conserváre dignéris,}
\elcol{℟. We beseech thee, hear us.}{℟. Te rogamus, audi nos.}
\elcoldent{That thou wouldest vouchsafe to reward all our benefactors with everlasting blessings,}{Ut ómnibus benefactóribus nostris sempitérna bona retríbuas,}
\elcol{℟. We beseech thee, hear us.}{℟. Te rogamus, audi nos.}
\elcoldent{That thou wouldest vouchsafe to give and preserve the fruits of the earth,}{Ut fructus terr{\ae} dare et conserváre dignéris,}
\elcol{℟. We beseech thee, hear us.}{℟. Te rogamus, audi nos.}
\elcoldent{That thou wouldest vouchsafe to grant unto all the faithful departed rest eternal,}{Ut ómnibus fidélibus defúnctis réquiem {\ae}térnam donáre dignéris,}
\elcol{℟. We beseech thee, hear us.}{℟. Te rogamus, audi nos.}
\elcoldent{That thou wouldest vouchsafe graciously to hear us,}{Ut nos exáudire dignéris,}
\elcol{℟. We beseech thee, hear us.}{℟. Te rogamus, audi nos.}
\elcoldent{O Lamb of God, who takest away the sins of the world,}{Agnus Dei, qui tollis peccáta mundi,}
\elcol{℟. Spare us, O Lord, hear us.}{℟. Parce nobis, Dómine.}
\elcoldent{O Lamb of God, who takest away the sins of the world,}{Agnus Dei, qui tollis peccáta mundi,}
\elcol{℟. Graciously hear us, O Lord.}{℟. Exáudi nos, Dómine.}
\elcoldent{O Lamb of God, who takest away the sins of the world,}{Agnus Dei, qui tollis peccáta mundi,}
\elcol{℟. Have mercy upon us.}{℟. Miserére nobis.}
\elcoldent{O Christ, hear us.}{Christe, audi nos.}
\elcol{℟. O Christ, graciously hear us.}{℟. Christe, exáudi nos.}

\subbysub{Holy Mass}

\begin{rubric}
	Here the Cantors solemnly intone \emph{Kyrie, eléison}, \emph{Christe, eléison}, \emph{Kyrie, eléison}, repeating each invocation thrice.
	
	Meanwhile the Priest with the Ministers in white vestments approaches the Altar: and having said the Psalm \emph{Give sentence with me, O God}, with \emph{Glory be}, he makes the Confession: then going up, he kisses the Altar, and censes it in the usual manner: and after the Choir has finished \emph{Kyrie, eléison}, he solemnly intones \emph{Gloria in excelsis}, and the bells are rung. Then the Priest says:
\end{rubric}

\elcol{℣. The Lord be with you.}{℣. Dóminus vobíscum.}
\elcol{℟. And with thy spirit.}{℟. Et cum spíritu tuo.}
\elcol{℣. Let us pray.}{℣. Orémus.}

\collect
\lett{O}{God,} who dost enlighten this most sacred night with the glory of the Resurrection of the Lord: preserve in the new offspring of thy family the spirit of adoption, which thou hast given; that they being renewed in body and soul may render thee pure service. Through the same.

\readingcitation{Epistle}{Colossians 3:1}
\lett{B}{rethren:} If ye be risen with Christ, seek those things which are above, where Christ sitteth on the right hand of God. Set your affection on things above, not on things on the earth. For ye are dead, and your life is hid with Christ in God. When Christ, who is our life, shall appear, then shall ye also appear with him in glory.

\begin{rubric}
	The Epistle ended, the Celebrant begins: \emph{Alleluia}.
	
	And be sings it through thrice, raising his voice each time: the Choir repeats it each time in the same tone.
	
	Afterwards the Choir continues:
\end{rubric}
\par\noindent
℣. O give thanks unto the Lord, for he is gracious: because his mercy endureth for ever.

\tract{O praise the Lord, all ye heathen: praise him, all ye nations. ℣. For his merciful kindness is ever more and more towards us: and the truth of the Lord endureth for ever.}

\begin{rubric}
	At the Gospel, lights are not carried, but only incense. A blessing is asked, and other things are done as usual.
\end{rubric}

\readingcitation{Gospel}{Matthew 28:1}
\lett{I}{n} the end of the sabbath, as it began to dawn toward the first day of the week, came Mary Magdalene and the other Mary to see the sepulchre. And, behold, there was a great earthquake: for the angel of the Lord descended from heaven, and came and rolled back the stone from the door, and sat upon it. His countenance was like lightning, and his raiment white as snow: and for fear of him the keepers did shake, and became as dead men. And the angel answered and said unto the women, Fear not ye: for I know that ye seek Jesus, which was crucified. He is not here: for he is risen, as he said. Come, see the place where the Lord lay. And go quickly, and tell his disciples that he is risen from the dead; and, behold, he goeth before you into Galilee; there shall ye see him: lo, I have told you.

\begin{rubric}
	The Creed is not said, but, the Gospel ended, the Priest says: \emph{The Lord be with you}, then: \emph{Let us pray}. The Offertory Verse is not said. At the Lavabo \emph{Gloria Patri} is said.
\end{rubric}

\secret
\lett{R}{eceive,} we beseech thee O Lord, the prayers of thy people together with the offering of these sacrifices: that those things which we have begun in these paschal mysteries may by thine operation be profitable for our healing in eternity. Through.

\begin{rubric}
    	\emph{The Peace of the Lord} is said, but the Kiss of Peace is not given.
    	
    	\emph{Agnus Dei} is not said, nor the Postcommunion, but the accustomed Prayers before the Communion are said.
\end{rubric}

\begin{rubric}
	After the receiving of the Sacrament, for Vespers in the choir is sung the Antiphon: \emph{Alleluia, alleluia, alleluia}, and
\end{rubric}

\subbysub{Psalm 117. \textit{Laudate Dominum}}

\elcol{\lett{O}{praise} the \textsc{Lord}, all ye heathen : praise him, all ye nations.}{\lett{L}{aud\smash{á}te} Dóminum, omnes gentes: * laudáte eum, omnes pópuli:}
\elcol{For his merciful kindness is ever more and more toward us : and the truth of the \textsc{Lord} endureth for ever. Praise the \textsc{Lord}.}{Quóniam confirmáta est super nos misericórdia ejus: * et véritas Dómini manet in {\ae}térnum.}
\elcol{Glory be to the Father, and to the Son, *
    and to the Holy Ghost:}{Glória Patri, et Fílio, * et Spirítui Sancto.}
\elcol{As it was in the beginning, is now, and ever shall be, *
    world without end. Amen.}{Sicut erat in princípio, et nunc, et semper, * et in sǽcula s{\ae}culórum. Amen.}
    
\antiphon{}{Alleluia, alleluia, alleluia.}

\begin{rubric}
	The Chapter, Hymn, and Verse are not said, but the Celebrant begins at once to sing the Antiphon to the Magnificat:
\end{rubric}

\lett{I}{n the end of the Sabbath.}

\vspace{1ex}

\begin{rubric}
	And the Choir continues:
\end{rubric}

As it began to dawn toward the first day of the week: came Mary Magdalene and the other Mary to see the sepulchre, alleluia.

\begin{rubric}
	Then the Magnificat is sung with \emph{Gloria Patri} at the end, and the censing is done, as at other times at Evensong.
\end{rubric}

\begin{rubric}
	The Antiphon having been repeated, the Celebrant says:
\end{rubric}

\elcol{℣. The Lord be with you.}{℣. Dóminus vobíscum.}
\elcol{℟. And with thy spirit.}{℟. Et cum spíritu tuo.}

\letuspray
\lett{P}{our} forth upon us, O Lord, the Spirit of thy charity: that as thou hast fulfilled us with these paschal sacraments, so of thy goodness thou wouldest make us to be of one heart and mind. Through ... in the unity of the same.

\elcol{℣. The Lord be with you.}{℣. Dóminus vobíscum.}
\elcol{℟. And with thy spirit.}{℟. Et cum spíritu tuo.}

\begin{rubric}
	And the Deacon, turning himself to the people, sings:
\end{rubric}

    \elcol{℣. Depart in peace, alleluia, alleluia.}{℣. Ite, Missa est, allelúja, allelúja.}
    \elcol{℟. Thanks be to God, alleluia, alleluia.}{℟. Deo grátias, allelúja, allelúja.}

\begin{rubric}
	And the Priest, having said \emph{Let this my bounden duty}, gives the Blessing in the usual manner.
\end{rubric}

\phantomsection
\addcontentsline{toc}{section}{Proper of Season}
\fancyhead[C]{\LARGE Proper of Season}

%\begin{secrubric}
%	Except during Advent \& Lent, the preceding Sunday's propers may be used for a weekday Feria.
%\end{secrubric}

\subby{First Sunday of Advent}\label{AdventI}
\fancyhead[RO,LE]{\textit{Advent I}}
\fancyhead[RE,LO]{}
%\begin{inhead}
%    {First Class Semidouble}
%\end{inhead}

\subbysub{I Evensong}\label{FirstAdventEvensong}

\gregorioscore{resources/gabc/ProperSeason/FirstAdventEvensong.gabc}

℣. Drop down, ye heavens, from above, and let the skies pour down righteousness.

℟. Let the earth open, and let them bring forth salvation.\\

\properantiphon{Mag.}{Behold, the Name of the Lord {\dag} cometh from afar: and his glory filleth all the earth.}

%Creator of the stars of night,
% 
% thy people's everlasting light,
% 
% Jesu, Redeemer, save us all,
% 
% And hear thy servants when they call.
%
%2. Thou, grieving that the ancient curse
% 
% Should doom to death a universe,
% 
% Hast found the med'cine, full of grace,
% 
% To save and heal a ruin'd race.
%
%3. Thou cam'st, the Bridegroom of the bride,
% 
% As drew the world to ev'ning-tide;
% 
% Proceeding from a virgin shrine,
% 
% The spotless victim all divine.
%
%4. At whose dread name, majestic now,
% 
% All knees must bend, all hearts must bow;
% 
% And things celestial thee shall own,
% 
% And things terrestrial, Lord alone.
%
%5. O thou whose coming is with dread
% 
% To judge and doom the quick and dead,
% 
% Preserve us, while we dwell below,
% 
% From every insult of the foe.
%
%6. To God the Father, God the Son,
% 
% And God the Spirit, Three in One,
% 
% Laud, honour, might, and glory be
% 
% From age to age eternally. Amen.

\subbysub{Mattins}

\invitatoryhymn\label{FirstAdventInvitatory}

\gregorioscore{resources/gabc/ProperSeason/FirstAdventInvitatory.gabc}

\officehymn\label{FirstAdventMattins}

\gregorioscore{resources/gabc/ProperSeason/FirstAdventMattins.gabc}

℣. The voice of one crying in the wilderness, Prepare ye the way of the Lord.

℟. Make his paths straight.

\properantiphon{Ben.}{The Holy Ghost {\dag} shall come upon thee, Mary; fear not, thou shalt bear in thy womb the Son of God, alleluia.}

%A thrilling voice by Jordan rings,
%
%rebuking guilt and darksome things:
%
%vain dreams of sin and visions fly;
%
%Christ in his might shines forth on high.
%
%2. Now let each torpid soul arise,
%
%that sunk in guilt and wounded lies;

%see! the new Star's refulgent ray
%
%shall chase disease and sin away.
%
%3. The Lamb descends from heav'n above
%
%to pardon sin with freest love:
%
%for such indulgent mercy shewn
%
%with tearful joy our thanks we own.
%
%4. That when again he shines reveal'd,
%
%and trembling worlds to terror yield.
%
%He give not sin its just reward,
%
%but in his love protect and guard.
%
%5. To God the Father, God the Son,
%
%And God the Spirit, Three in One,
%
%Laud, honour, might, and glory be
%
%From age to age eternally. Amen.

\subbysub{II Evensong}

\begin{rubric}
	II Evensong as in I Evensong, except for the following Antiphon.
\end{rubric}

\properantiphon{Mag.}{Fear not, Mary, {\dag} for thou hast found favour with the Lord: behold, thou shalt conceive, and bring forth a Son, alleluia.}


\subby{Advent Ferial Office}
\begin{rubric}
	On Advent Ferias, the Invitatory Hymn is of the First Sunday of Advent. The Office Hymn and Versicle is taken from Sunday (p. \pageref{SundayMattinsWinter}). The Antiphon is of the Day.
\end{rubric}


\subby{The First Week of Advent}
\feastday{Advent I Week}

\begin{multicols}{2}
\subsubsec{Monday}
\antiphon{Ben.}{The Angel of the Lord {\dag} announced unto Mary, and she conceived by the Holy Ghost, alleluia.}

\antiphon{Mag.}{Lift up thine eyes, {\dag} O Jerusalem, and see the power of the King: behold, the Saviour cometh to loose thee from thy chain.}

\subsubsec{Tuesday}
\antiphon{Ben.}{Before they came together {\dag} Mary was found with child of the Holy Ghost, alleluia.}

\antiphon{Mag.}{Seek ye the Lord {\dag} while he may be found: call ye upon him while he is near, alleluia.}

\subsubsec{Wednesday}
\antiphon{Ben.}{Out of Sion {\dag} shall go forth the law, and the word of the Lord from Jerusalem.}

\antiphon{Mag.}{After me {\dag} cometh one mightier than I, the latchet of whose shoes I am not worthy to unloose.}

\subsubsec{Thursday}
\antiphon{Ben.}{Blessed art thou {\dag} among women, and blessed is the fruit of thy womb.}

\antiphon{Mag.}{I will wait {\dag} upon the Lord my Saviour: and I will look for him while he is near, alleluia.}

\subsubsec{Friday}
\antiphon{Ben.}{Lo, there cometh one {\dag} that is both God and Man, of the house of David, to sit upon the throne, alleluia.}

\antiphon{Mag.}{Out of Egypt {\dag} have I called my Son: he shall come to save his people.}

\subsubsec{Saturday}
\antiphon{Ben.}{Fear not, {\dag} O Sion, behold thy God cometh, alleluia.}

\end{multicols}


\subby{Second Sunday of Advent}
\feastday{Advent II}
%\begin{inhead}
%{Second Class Semidouble}
%\end{inhead}

\begin{rubric}
	The Hymns \& Versicles are of the First Sunday of Advent (p. \pageref{AdventI}), with the following Antiphons.
\end{rubric}

\properantiphon{Mag.}{Come, O Lord, in peace; {\dag} visit us with thy salvation, that we may rejoice before thee with a perfect heart.}

\properantiphon{Ben.}{He shall sit upon the throne {\dag} of David, and of his kingdom there shall be no end, alleluia.}

\properantiphon{Mag.}{Blessed art thou, {\dag} O Mary, for thou hast believed the Lord: and there shall be a performance in thee of those things which were told thee from the Lord, alleluia.}


\subby{The Second Week of Advent}
\feastday{Advent II Week}

\begin{multicols}{2}
\subsubsec{Monday}
\antiphon{Ben.}{From heaven there cometh {\dag} the Lord, the Ruler, and in his hand are honour and dominion.}

\antiphon{Mag.}{Behold the King shall come, {\dag} the Lord of the earth, and he shall take away the yoke of our captivity.}

\subsubsec{Tuesday}
\antiphon{Ben.}{The Lord shall arise upon thee, {\dag} O Jerusalem, and his glory shall be seen upon thee.}

\antiphon{Mag.}{The voice of one crying {\dag} in the wilderness, Prepare ye the way of the Lord, make straight in the desert a highway for our God.}

\subsubsec{Wednesday}
\antiphon{Ben.}{Behold I send my messenger, {\dag} and he shall prepare my way before thy face.}

\antiphon{Mag.}{Sion, {\dag} thou shalt be renewed, and shalt see thy righteous One, he that cometh unto thee.}

\subsubsec{Thursday}
\antiphon{Ben.}{Thou, O Lord, {\dag} art he that is to come, for whom we look, to save thy people.}

\antiphon{Mag.}{He that cometh after me {\dag} is preferred before me, the latchet of whose shoes I am not worthy to unloose.}

\subsubsec{Friday}
\antiphon{Ben.}{Say to them: {\dag} Ye that are of a fearful heart, be strong: for behold the Lord our God shall come.}

\antiphon{Mag.}{Sing unto the Lord {\dag} a new song; and his praise from the end of the earth.}

\subsubsec{Saturday}
\antiphon{Ben.}{The Lord shall set up an ensign {\dag} for the nations, and shall assemble the outcasts of Israel.}

\end{multicols}

\subby{Magnificat Antiphons before Christmas}\label{OAntiphons}
\fancyhead[RO,LE]{\textit{Advent Antiphons}}
\fancyhead[RE,LO]{}
\begin{rubric}
    The Magnificat Antiphon is always taken from here on these days.
\end{rubric}
%\begin{multicols}{2}
\subbysub{16 December: O Sapientia}
\lett{O}{Wisdom,} {\dag} which camest out of the mouth of the Most High, and reachest from one end to another, mightily and sweetly ordering all things: Come and teach us the way of prudence.
\subbysub{17 December: O Adonai}
\lett{O}{Adonai,} {\dag} and Leader of the house of Israel, who appearedst in the bush to Moses in a flame of fire, and gavest him the law in Sinai: Come and redeem us with an outstretched arm.
\subbysub{18 December: O Radix Jesse}
\lett{O}{Root of Jesse,} {\dag} which standest for an ensign of the people, at whom kings shall shut their mouths, unto whom the Gentiles shall seek: Come and deliver us, and tarry not.
\subbysub{19 December: O Clavis David}
\lett{O}{Key of David,} {\dag} and Sceptre of the house of Israel; that openest and no man shutteth, and shuttest and no man openeth: Come, and bring the prisoners out of the prison-house, them that sit in darkness and the shadow of death.
\subbysub{20 December: O Oriens}
\lett{O}{Day-spring,} {\dag} Brightness of the Light everlasting, and Sun of righteousness: Come and enlighten them that sit in darkness and the shadow of death.
\subbysub{21 December: O Rex gentium}
\lett{O}{King of Nations,} {\dag} and their Desire; the Cornerstone, who makest both one: Come and save mankind, whom thou formedst of clay.
\subbysub{22 December: O Emmanuel}
\lett{O}{Emmanuel,} {\dag} our King and Lawgiver, the Desire of all nations and their Salvation: Come and save us, O Lord our God.
\subbysub{23 December: O Virgo virginum}
\lett{O}{Virgin of Virgins,} {\dag} how shall this be? for neither before thee was any seen like thee, nor shall there be after. Daughters of Jerusalem, why marvel ye at me? The thing which ye behold is a divine mystery.
%\end{multicols}


\subby{Third Sunday of Advent}
\feastday{Advent III}
%\begin{inhead}
%{Second Class Semidouble}
%\end{inhead}

\begin{rubric}
	The Hymns \& Versicles are of the First Sunday of Advent (p. \pageref{AdventI}), with the following Antiphons.
\end{rubric}

\properantiphon{Mag.}{Before me {\dag} there was no God formed, neither shall there be after me: unto me every knee shall bow, and me shall every tongue confess.}

\properantiphon{Ben.}{Now when John had heard {\dag} in the prison the works of Christ, he sent two of his disciples, and said unto him: Art thou he that should come, or do we look for another?}

\properantiphon{Mag.}{Art thou he {\dag} that should come, or do we look for another? Go and shew John those things which ye do see: the blind receive their sight, the dead are raised, and the poor have the gospel preached unto them, alleluia.}


\subby{The Third Week of Advent}
\feastday{Advent III Week}

\begin{multicols}{2}
\subsubsec{Monday}
\antiphon{Ben.}{There shall come forth {\dag} a rod out of the stem of Jesse, and all the earth shall be filled with the glory of the Lord: and all flesh shall see the salvation of God.}

\antiphon{Mag.}{All generations {\dag} shall call me blessed: for God hath regarded the lowliness of his handmaiden.}

\subsubsec{Tuesday}
\antiphon{Ben.}{Thou, Bethlehem, {\dag} in the land of Juda, shalt not be the least: for out of thee shall come a Governor, that shall rule my people Israel.}

\properantiphon{Mag.}{Awake, awake, {\dag} arise, O Jerusalem: loose thyself from the bands of thy neck, O captive daughter of Sion.}

\subsubsec{Ember Wednesday}
\properantiphon{Ben.}{The Angel {\dag} Gabriel was sent to Mary, a Virgin espoused to Joseph.}

\properantiphon{Mag.}{Behold the handmaid of the Lord; {\dag} be it unto me according to thy word.}

\subsubsec{Thursday}
\antiphon{Ben.}{Be ye watchful {\dag} in your hearts, for the Lord our God is nigh at hand.}

\properantiphon{Mag.}{Rejoice ye {\dag} with Jerusalem; and be glad with her, all ye that love her, forever.}

\subsubsec{Ember Friday}
\properantiphon{Ben.}{For lo, as soon {\dag} as the voice of thy salutation sounded in mine ears, the babe leaped in my womb for joy, alleluia.}

\properantiphon{Mag.}{This is the witness {\dag} which John bare, saying: He that cometh after me is preferred before me.}

\subsubsec{Ember Saturday}
\properantiphon{Ben.}{How shall this be, {\dag} O Angel of God, seeing I know not a man? Hearken, O Virgin Mary: the Holy Ghost shall come upon thee, and the power of the Highest shall overshadow thee.}

\end{multicols}


\subby{Fourth Sunday of Advent}
\feastday{Advent IV}
\fancyhead[RE,LO]{}
%\begin{inhead}
%{Second Class Semidouble}
%\end{inhead}


\begin{rubric}
	In the Daily Office, the Hymns \& Versicles are of the First Sunday of Advent (p. \pageref{AdventI}), with the following Antiphon.
\end{rubric}

\properantiphon{Ben.}{Hail, Mary, {\dag} thou that art full of grace, the Lord is with thee; blessed art thou among women, alleluia.}


\subby{The Fourth Week of Advent}
\feastday{Advent IV Week}

\begin{multicols}{2}
\subsubsec{Monday}
\antiphon{Ben.}{Thus saith the Lord {\dag} your God: Repent ye, and turn again; for the kingdom of heaven is at hand, alleluia.}

\subsubsec{Tuesday}
\antiphon{Ben.}{Awake, awake, {\dag} put on strength, O arm of the Lord.}

\subsubsec{Wednesday}
\antiphon{Ben.}{I will place salvation {\dag} in Sion, and in Jerusalem my glory, alleluia.}

\subsubsec{Thursday}
\antiphon{Ben.}{Comfort ye, {\dag} comfort ye my people, saith your God.}

\subsubsec{Friday}
\antiphon{Ben.}{Behold all things are fulfilled {\dag} which were spoken by the Angel of the Virgin Mary.}

\end{multicols}


\subby{Christmas Eve}
\feastday{Christmas Eve}
%\begin{inhead}
%{First Class Vigil}
%\end{inhead}

\begin{rubric}
	In Mattins, the Hymn is of the First Sunday of Advent (p. \pageref{AdventI}), with the following.
\end{rubric}

℣. To-morrow the iniquity of the earth shall be blotted out. 

℟. And the Saviour of the world shall reign over us.

\properantiphon{Ben.}{The Saviour {\dag} of the world shall arise as the sun: and shall come down into the Virgin's womb, as the showers upon the grass, alleluia.}


\subby{Christmas Day}\label{ChristmasDay}
\feastday{Christmas}
%\begin{inhead}
%{First Class Double, Third Class Octave}
%\end{inhead}

\subbysub{I Evensong}\label{ChristmasEvensong}

\gregorioscore{resources/gabc/ProperSeason/ChristmasEvensong.gabc}

℣. To-morrow the iniquity of the earth shall be blotted out.

℟. And the Saviour of the world shall reign over us.

\properantiphon{Mag.}{When the sun hath risen {\dag} in the heavens, ye shall see the King of kings proceeding from the Father, as a bridegroom out of his chamber.}

\subbysub{Mattins}

\invitatoryhymn

\begin{rubric}
	The Invitatory Hymn is as in I Evensong, with the following ending.
\end{rubric}

\gregorioscore{resources/gabc/ProperSeason/ChristmasInvitatory.gabc}

\officehymn\label{ChristmasMattins}

\gregorioscore{resources/gabc/ProperSeason/ChristmasMattins.gabc}

℣. The Lord declared, alleluia.

℟. His salvation, alleluia.

\properantiphon{Ben.}{Glory {\dag} to God in the highest, and on earth peace to men of good will, alleluia, alleluia.}

\subbysub{II Evensong}

\begin{rubric}
	The Office Hymn is from I Evensong with the following Versicle and Antiphon.
\end{rubric}

℣. The Lord declared, alleluia.

℟. His salvation, alleluia.

\properantiphon{Mag.}{To-day {\dag} the Christ is born; to-day hath a Saviour appeared: to-day on earth Angels are singing, Archangels rejoicing: today the righteous exult and say. Glory to God in the highest, alleluia.}


\subby{26 December. St. Stephen}
\feastday{St. Stephen}
%\begin{inhead}
%{26 December}\par
%{Second Class Double, Simple Octave}
%\end{inhead}

\subbysub{Mattins}

\begin{rubric}
	The Invitatory Hymn is from the First Common of a Martyr Bishop (p. \pageref{CommonMartyrBishopI}).
\end{rubric}

\officehymn\label{StephenMattins}

\gregorioscore{resources/gabc/ProperSeason/StephenMattins.gabc}

℣. Devout men carried Stephen to his burial.

℟. And made great lamentation over him.

\properantiphon{Ben.}{And Stephen, {\dag} full of faith and power, did great wonders and miracles among the people.}

\subbysub{Evensong}\label{StephenEvensong}

\gregorioscore{resources/gabc/ProperSeason/StephenEvensong.gabc}

℣. Stephen saw the heavens opened.

℟. He saw, and entered in: how blessed is the man unto whom heaven stood open!

\properantiphon{Mag.}{Devout men {\dag} carried Stephen to his burial, and made great lamentation over him.}


\subby{27 December. St. John}
\feastday{St. John}
%\begin{inhead}
%{27 December}\par
%{Second Class Double, Simple Octave}
%\end{inhead}

\subbysub{Mattins}

\begin{rubric}
	The Invitatory Hymn is from the Common of Apostles (p. \pageref{CommonApostles}).
\end{rubric}

\officehymn\label{JohnMattins}

\gregorioscore{resources/gabc/ProperSeason/JohnMattins.gabc}

℣. This is that disciple which testifieth of these things.

℟. And we know that his testimony is true.

\properantiphon{Ben.}{This is the same John, {\dag} who leaned on the Lord's bosom at the last supper: the blessed Apostle, unto whom were revealed the secret things of heaven.}

\subbysub{Evensong}

\begin{rubric}
	The Office Hymn is of Mattins with the following Versicle \& Antiphon.
\end{rubric}

℣. Right worthy of honour is blessed John the Apostle.

℟. Who leaned on the Lord's bosom at the last supper.

\properantiphon{Mag.}{There went {\dag} this saying abroad among the brethren, that that disciple should not die: yet Jesus said not unto him, He shall not die; but, If I will that he tarry till I come.}


\subby{28 December. Holy Innocents}
\feastday{Holy Innocents}
%\begin{inhead}
%{28 December}\par
%{Second Class Double, Simple Octave}
%\end{inhead}

\subbysub{Mattins}

\invitatoryhymn\label{HolyInnocentsInvitatory}

\gregorioscore{resources/gabc/ProperSeason/HolyInnocentsInvitatory.gabc}

\officehymn\label{HolyInnocentsMattins}

\gregorioscore{resources/gabc/ProperSeason/HolyInnocentsMattins.gabc}

℣. Herod was exceeding wroth, and slew many children.

℟. In Bethlehem Judah, the city of David.

\properantiphon{Ben.}{These are they {\dag} which were not defiled with women; for they are virgins: and they follow the Lamb whithersoever he goeth.}

\subbysub{Evensong}

\begin{rubric}
	The Office Hymn is of Mattins with the following Versicle \& Antiphon.
\end{rubric}

℣. Beneath the altar all the Saints do cry aloud.

℟. Shall not our blood be avenged, O our God.

\properantiphon{Mag.}{Then were innocent children {\dag} slain instead of Christ by a wicked ruler; the very sucklings were put to death: spotless, they follow the Lamb himself, and say for ever: Glory be to thee, O Lord.}


\subby{Sunday within the Nativity Octave}
\feastday{Christmas I}
%\begin{inhead}
%{Semidouble}
%\end{inhead}

\begin{rubric}
	In the Daily Office, the Hymns are of Christmas Day (p. \pageref{ChristmasDay}), with the following.
\end{rubric}

℣. All the ends of the world have seen, alleluia.

℟. The salvation of our God, alleluia.

\properantiphon{Mag. \& Ben.}{While all things {\dag} were in quiet silence, and that night was in the midst of her swift course, thine Almighty Word, O Lord, leaped down out of thy royal throne, alleluia.}

\properantiphon{Mag.}{The Child Jesus {\dag} increased in wisdom and stature in the sight of God and man.}


\subby{1 January. Circumcision of Our Lord}
\feastday{Circumcision}
%\begin{inhead}
%{1 January}\par
%{Second Class Double}
%\end{inhead}

\begin{rubric}
	In the Daily Office, the Hymns are of Christmas Day (p. \pageref{ChristmasDay}), with the following.
\end{rubric}

℣. The Lord declared, alleluia.

℟. His salvation, alleluia.

\properantiphon{Mag.}{God, {\dag} for his great love wherewith he loved us, hath sent his own Son in the likeness of sinful flesh, alleluia.}\\

℣. The Word was made flesh, alleluia.

℟. And dwelt among us, alleluia.

\properantiphon{Ben.}{A great and wondrous mystery {\dag} is made known to us this day; a new thing is wrought in both natures: God is made man; that which was, remained, and that which was not, he assumed; suffering no confusion, nor yet division.}\\

℣. The Lord declared, alleluia.

℟. His salvation, alleluia.

\properantiphon{Mag.}{Great {\dag} is the mystery of the inheritance: the womb of her that knew not man is become the temple of the Godhead: by taking flesh of her, he was no way defiled: all the nations shall gather, saying: Glory be to thee, O Lord.}


\subby{2 January. Most Holy Name of Jesus}\label{MostHolyName}
\feastday{Most Holy Name}
\fancyhead[RE,LO]{2 January}
%\begin{inhead}
%{2 January}\par
%{Second Class Double}
%\end{inhead}

\subbysub{I Evensong}\label{HolyNameEvensong}

\gregorioscore{resources/gabc/ProperSeason/HolyNameEvensong.gabc}

℣. Blessed be the Name of the Lord, alleluia.

℟. From this time forth for evermore, alleluia.

\properantiphon{Mag.}{For he that is mighty {\dag} hath magnified me, and holy is his Name, alleluia.}

\subbysub{Mattins}

\invitatoryhymn\label{HolyNameInvitatory}

\gregorioscore{resources/gabc/ProperSeason/HolyNameInvitatory.gabc}

\officehymn\label{HolyNameMattins}

\gregorioscore{resources/gabc/ProperSeason/HolyNameMattins.gabc}

℣. Our help standeth in the Name of the Lord, alleluia. 

℟. Who hath made heaven and earth, alleluia.

\properantiphon{Ben.}{He gave himself {\dag} that he might deliver his people, and get him a perpetual Name, alleluia.}

\subbysub{II Evensong}

\begin{rubric}
	The Office Hymn \& Versicle are of I Evensong with the following Antiphon.
\end{rubric}

\properantiphon{Mag.}{Thou shalt call {\dag} his Name Jesus, for he shall save his people from their sins, alleluia.}

%MANUAL ADJUSTMENT:
\clearpage
\subby{Second Sunday after Nativity}
\feastday{Christmas II}
\fancyhead[RE,LO]{}
%\begin{inhead}
%{The Sunday between Circumcision \& Epiphany, if there be one.}\par
%{Semidouble}
%\end{inhead}

\properantiphon{Mag.}{Behold, the Angel of the Lord {\dag} appeareth in a dream to Joseph, saying: Arise, and take the young Child and his Mother, and flee into Egypt, and be thou there until I bring thee word; for it shall come to pass, that Herod will seek the young Child, to destroy him.}

\properantiphon{Ben.}{But when Herod was dead, {\dag} behold an Angel of the Lord appeareth in a dream to Joseph in Egypt, saying: Arise, and take the young Child and his Mother, and go into the land of Israel; for they are dead which sought the young Child's life.}

\properantiphon{Mag.}{When Joseph heard {\dag} that Archelaus did reign in Jud{\ae}a in the room of his father Herod, he was afraid to go thither; and being warned in a dream, he turned aside into the parts of Galilee: and he came and dwelt in a city which was called Nazareth; that it might be fulfilled which was spoken by the prophets, He shall be called a Nazarene.}


\subby{2 January. Octave Day of St. Stephen}
\fancyhead[RO,LE]{\textit{Stephen Octave Day}}
\fancyhead[RE,LO]{2 January}
%\begin{inhead}
%    {Simple\\
%2 January}
%\end{inhead}

\begin{rubric}
	The propers are from the Feast Day, except for that which followeth.
\end{rubric}

\collect
\lett{A}{lmighty} and everlasting God, who, in the blood of the blessed Levite Stephen, didst consecrate the first-fruits of the Martyrs: grant, we beseech thee; that he may ever intercede for us, who prayed even for his persecutors to our Lord Jesus Christ, thy Son. Who with thee liveth.


\subby{3 January. Octave Day of St. John}
\fancyhead[RO,LE]{\textit{John Octave Day}}
\fancyhead[RE,LO]{3 January}
%\begin{inhead}
%    {Simple\\
%3 January}
%\end{inhead}

\begin{rubric}
	The propers are from the Feast Day, except for that which followeth.
\end{rubric}

\antiphon{Mag.}{This is the same John {\dag} who learned on the Lord's bosom at the last supper: the blessed Apostle, unto whom were revealed the secret things of heaven.}

\collect
\lett{M}{erciful} Lord, we beseech thee to cast thy bright beams of light upon thy Church: that it being enlightened by the doctrine of thy blessed Apostle and Evangelist Saint John may so walk in the light of thy truth, that it may at length attain to the light of everlasting life. Through.

\secret
\lett{R}{eceive,} O Lord, the gifts which we offer unto thee on the solemnity of him, in whose advocacy we trust for deliverance. Through.

\postcommunion
\lett{O}{God,} who hast refreshed us with heavenly meat and drink, we humbly beseech thee: that we may be defended by the prayers of him, in whose memory we have received the same. Through.
\needspace{4\baselineskip}
\begin{rubric}
	If the Postcommunion of this Mass shall have been already said for some other Saint, the Postcommunion for St. John shall be that of his Feast before the Latin Gate, 6 May, as followeth.
\end{rubric}
\lett{W}{e} beseech thee, O Lord: that we, who have been refreshed with heavenly bread, may be nourished unto life eternal. Through.



\subby{4 January. Octave Day of Holy Innocents}
\fancyhead[RO,LE]{\textit{Innocents Octave Day}}
\fancyhead[RE,LO]{4 January}
%\begin{inhead}
%    {Simple\\
%4 January}
%\end{inhead}

\begin{rubric}
	The propers are from the Feast Day, except for that which followeth.\par
	\textsc{Note,} The \emph{Gloria in excelsis} is said.
\end{rubric}

℣. Herod was exceeding wroth, and slew many children.

℟. In Bethlehem Judah, the city of David.\\

\antiphon{Mag.}{These are they {\dag} which were not defiled with women; for they are virgins: and they follow the Lamb whithersoever he goeth.}

\begin{rubric}
	If a Commemoration only is made of the Octave, the Prayers (of the Feast) will be as followeth.
\end{rubric}

\collect
\lett{O}{almighty} God, who out of the mouths of babes and sucklings hast ordained strength, and madest infants to glorify thee by their deaths: mortify and kill all vices in us; and so strengthen us by thy grace, that by the innocency of our lives, and constancy of our faith even unto death, we may glorify thy holy name. Through.

\begin{rubric}
	Commemoration is made of St. Titus (p. \pageref{TitusCollect}).
\end{rubric}

\secret
\lett{M}{ay} the devout prayers of thy Saints never fail us, O Lord: that they may render our gifts acceptable unto thee, and ever obtain for us thy pardon. Through.

\begin{rubric}
	Commemoration is made of St. Titus (p. \pageref{TitusSecret}).
\end{rubric}

\postcommunion
\lett{W}{e} beseech thee, O Lord: that the gifts which we have offered and received may, through the prayers of the Saints, effectually avail for our succour both in this life, and in that which is to come. Through.

\begin{rubric}
	Commemoration is made of St. Titus (p. \pageref{TitusPostcommunion}).
\end{rubric}

\begin{rubric}
	\textsc{Note,} In Lent, either the Ferial Day's propers or Ash Wednesday's propers may be said.
\end{rubric}

\begin{rubric}
	The Office Hymn \& Versicles of the Feast of the Epiphany of Our Lord shall be used for the foregoing Sundays and Ferial Days until I Evensong of Septuagesima, exclusive.
\end{rubric}


\subby{4 January. St. Titus}\label{Titus}
\fancyhead[RO,LE]{\textit{Titus}}
\fancyhead[RE,LO]{4 January}
%\begin{inhead}
%    {Simple\\
%4 January}
%\end{inhead}

\begin{rubric}
	The propers are from the First Common of a Confessor Bishop (p. \pageref{CommonConfessorBishopI}), except for the Gospel \& Prayers as followeth.
\end{rubric}

\collect\label{TitusCollect}
\lett{O}{God,} who didst adorn blessed Titus, thy Confessor and Bishop, with apostolic virtues: grant, through his merits and intercession; that, living justly and godly in this world, we may be found worthy to attain unto the heavenly country. Through.

\readingcitation{Gospel}{Luke 10:1}
\lett{A}{t that time:} The Lord appointed seventy also: and sent them two and two before his face into every city and place, whither he himself would come. Therefore said he unto them, The harvest truly is great, but the labourers are few: pray ye therefore the Lord of the harvest, that he would send forth labourers into his harvest. Go your ways: behold, I send you forth as lambs among wolves. Carry neither purse, nor scrip, nor shoes: and salute no man by the way. And into whatsoever house ye enter, first say, Peace be to this house. And if the son of peace be there, your peace shall rest upon it: if not, it shall turn to you again. And in the same house remain, eating and drinking such things as they give: for the labourer is worthy of his hire. Go not from house to house. And into whatsoever city ye enter, and they receive you, eat such things as are set before you: And heal the sick that are therein, and say unto them, The kingdom of God is come nigh unto you.

\secret\label{TitusSecret}
\lett{W}{e} beseech thee, O Lord, that we remembering with gladness the merits of thy Saints, may in all places feel the succour of their intercession. Through.

\postcommunion\label{TitusPostcommunion}
\lett{G}{rant,} we beseech thee, almighty God: that we, shewing forth our thankfulness for the gifts which we have received, may, at the intercession of blessed Titus, thy Confessor and Bishop, obtain yet more abundant mercies. Through.


\subby{5 January. Vigil of Epiphany}
\fancyhead[RO,LE]{\textit{Epiphany Vigil}}
\fancyhead[RE,LO]{5 January}
%\begin{inhead}
%    {Second Class Vigil\\
%5 January}
%\end{inhead}

\begin{rubric}
	The Daily Office propers are as on the Feast of the Circumcision, except that which followeth.
\end{rubric}

\begin{rubric}
	\textsc{Note,} Commemoration is made of St. Telesphorus with the Prayers of the next Mass.
\end{rubric}

\antiphon{Mag.}{The Child Jesus {\dag} increased in wisdom and stature in the sight of God and man.}

\antiphon{Ben.}{While all things {\dag} were in quiet silence, and that night was in the midst of her swift course, thine Almighty Word, O Lord, leaped down out of thy royal throne, alleluia.}

\introit
\lett{W}{hile} all things were in in quiet silence, and night was in the midst of her swift course, thine almighty Word, O Lord, leaped down from heaven out of thy royal throne. \textit{Ps.} The Lord is King, and hath put on put glorious apparel: the Lord hath put on his apparel, and girded himself with strength.

\collect
\lett{A}{lmighty} and everlasting God, direct our actions according to thy good pleasure: that in the name of thy well-beloved Son we may be made worthy to abound in good works. Who liveth and reigneth with thee.

\begin{rubric}
	Commemoration of St. Telesphorus (p. \pageref{TelesphorusCollect}) and for St. Mary after Christmas (p. \SPMaryChristmas) is said.
\end{rubric}

\readingcitation{Epistle}{Galatians 4:1}
\lett{B}{rethren:} The heir, as long as he is a child, differeth nothing from a servant, though he be lord of all; But is under tutors and governors until the time appointed of the father. Even so we, when we were children, were in bondage under the elements of the world: But when the fulness of the time was come, God sent forth his Son, made of a woman, made under the law, To redeem them that were under the law, that we might receive the adoption of sons. And because ye are sons, God hath sent forth the Spirit of his Son into your hearts, crying, Abba, Father. Wherefore thou art no more a servant, but a son; and if a son, then an heir of God through Christ.

\gradall{Thou art fairer than the children of men: full of grace are thy lips. ℣. My heart is inditing of a good matter, I speak of the things which I have made unto the King: my tongue is the pen of a ready writer.}{Alleluia, alleluia. ℣. The Lord is King, and hath put on glorious apparel: the Lord hath put on his apparel, and girded himself with strength. Alleluia.}

\readingcitation{Gospel}{Matthew 2:19}
\lett{A}{t that time:} When Herod was dead, behold, an angel of the Lord appeareth in a dream to Joseph in Egypt, Saying, Arise, and take the young child and his mother, and go into the land of Israel: for they are dead which sought the young child's life. And he arose, and took the young child and his mother, and came into the land of Israel. But when he heard that Archelaus did reign in Judaea in the room of his father Herod, he was afraid to go thither: notwithstanding, being warned of God in a dream, he turned aside into the parts of Galilee: And he came and dwelt in a city called Nazareth: that it might be fulfilled which was spoken by the prophets, He shall be called a Nazarene.

\offertory{God hath made the round world so sure: that it cannot be moved: ever since the world began, hath thy seat, O God, been prepared, thou art from everlasting.}

\secret
\lett{G}{rant,} we beseech thee, almighty God: that the gift which we offer in the sight of thy majesty, may obtain for us grace to serve thee with all godliness, and bring us in the end to everlasting felicity. Through.

\begin{rubric}
	Commemoration of St. Telesphorus (p. \pageref{TelesphorusSecret}) and for St. Mary after Christmas (p. \SPMaryChristmas) is said.
\end{rubric}

\communion{Take the young Child and his Mother and go into the land of Israel: for they are ead which sought the young Child's life.}

\postcommunion
\lett{M}{ay} the operation of this mystery, Lord, avail for the cleansing of our sins, and for the fulfilment of our godly desires. Through.

\begin{rubric}
	Commemoration of St. Telesphorus (p. \pageref{TelesphorusPostcommunion}) and for St. Mary after Christmas (p. \SPMaryChristmas) is said.
\end{rubric}


\subby{5 January. St. Telesphorus}
\fancyhead[RO,LE]{\textit{Telesphorus}}
\fancyhead[RE,LO]{5 January}
%\begin{inhead}
%    {Memorial\\
%5 January}
%\end{inhead}

\begin{rubric}
	The propers are from the Second Common of a Martyr Bishop (p. \pageref{CommonMartyrBishopII}), with the Prayers below.
\end{rubric}

\collect\label{TelesphorusCollect}
\lett{O}{God,} who makest us glad with the yearly solemnity of blessed Telesphorus, thy Martyr and Bishop: mercifully grant; that, as we now celebrate his birthday so we may likewise rejoice in his protection. (Through.)

\secret\label{TelesphorusSecret}
\lett{S}{anctify,} O Lord, the gifts which we dedicate to thee: that at the intercession of blessed Telesphorus, thy Martyr and Bishop, they may obtain for us thy gracious favour. (Through.)

\postcommunion\label{TelesphorusPostcommunion}
\lett{W}{e} beseech thee, O Lord our God, that like as we, whom thou hast refreshed by the partaking of thy sacred gift, do offer unto thee our worship: so by the intercession of blessed Telesphorus thy Martyr and Bishop, we may perceive the benefit of the same. (Through.)\\


\subby{Epiphany of Our Lord}\label{EpiphanyDay}
\feastday{Epiphany}
\fancyhead[RE,LO]{}
%\begin{inhead}
%{First Class Double, Second Class Octave}
%\end{inhead}

\subbysub{I Evensong}\label{EpiphanyEvensong}

\gregorioscore{resources/gabc/ProperSeason/EpiphanyEvensong.gabc}

℣. The kings of Tharsis and of the isles shall give presents.

℟. The kings of Arabia and Saba shall bring gifts.

\properantiphon{Mag.}{The wise men, {\dag} beholding the star, said one to another, This is the sign of a mighty King; forth fare we, and let us seek him: and let us offer him gifts, gold, incense, and myrrh, alleluia.}

\subbysub{Mattins}

\begin{rubric}
	The Invitatory Hymn is as in I Evensong.
\end{rubric}

\officehymn\label{EpiphanyMattins}

\gregorioscore{resources/gabc/ProperSeason/EpiphanyMattins.gabc}

	℣. O worship the Lord, alleluia.

	℟. All ye Angels of his, alleluia.
	
	\properantiphon{Ben.}{To-day {\dag} the Church is joined to her heavenly Bridegroom; because in Jordan Christ hath washed away her offences: the wise men with their offerings hasten to the royal marriage, and the guests are regaled with water made wine, alleluia.}

\subbysub{II Evensong}

\begin{rubric}
	The Office Hymn \& Versicle are of I Evensong, with the following Antiphon.
\end{rubric}

\properantiphon{Mag.}{Now do we celebrate {\dag} a holy day adorned by three miracles: to-day a star led the wise men to the manger; to-day water was made wine at the wedding feast; to-day Christ vouchsafed to be baptized of John in Jordan that he might save us, alleluia.}


\begin{rubric}
	Within the Octave of the Epiphany, no Feast is kept except an occurrent I Double, on which Commemoration of the Octave is always made.
	
	During the Octave all is said as on the Feast, except the Ants. on Benedictus and Magnificat appointed for each day.
	
	The Office of Sunday within the Octave is said as set forth below, on whatever day in the Octave it may fall.
	
	After Sunday, the days in the Octave continue according to their number, as if Sunday had not intervened. For example, if Sunday should fall on the third day in the Octave, on the following day the Ants. on Benedictus and Magnificat of the third day in the Octave are said, and so thenceforth the others.
\end{rubric}


\subby{Within the Octave of the Epiphany}
\feastday{Within Epiphany Octave}

%MANUAL ADJUSTMENT:
%\begin{multicols}{2}
\subsubsec{Day II, Semidouble}
\antiphon{Ben.}{From the east {\dag} there came wise men to Bethlehem, to worship the Lord: and when they had opened their treasures, they presented unto him precious gifts: gold as to a mighty King, incense as to the true God, and myrrh to foreshew his burial, alleluia.}

\antiphon{Mag.}{When the wise men saw the star {\dag} they rejoiced with exceeding great joy: and when they were come into the house, they presented unto the Lord gold, frankincense, and myrrh, alleluia.}

\subsubsec{Day III, Semidouble}
\antiphon{Ben.}{Three are the gifts {\dag} which the wise men presented unto the Lord: gold, frankincense, and myrrh, to the Son of God, to the mighty King, alleluia.}

\antiphon{Mag.}{Light of light, {\dag} thou, O Christ, hast appeared, unto whom the wise men present their gifts, alleluia, alleluia, alleluia.}
%\end{multicols}
%\clearpage
%MANUAL ADJUSTMENT:
%\begin{multicols}{2}
\subsubsec{Day IV, Semidouble}
\antiphon{Ben.}{We have seen his star {\dag} in the East, and are come with gifts to worship the Lord.}

\antiphon{Mag.}{Herod inquired {\dag} of the wise men, What sign have ye seen, concerning him that is born King? We have seen a star shining, the brightness whereof enlightens the world.}

\subsubsec{Day V, Semidouble}
\antiphon{Ben.}{Many nations {\dag} shall come from afar, bearing their gifts, alleluia.}

\antiphon{Mag.}{All they from Sheba {\dag} shall come: they shall bring gold and incense, alleluia, alleluia.}

\subsubsec{Day VI, Semidouble}
\antiphon{Ben.}{They that despised thee {\dag} shall come unto thee, and shall bow themselves down at the soles of thy feet.}

\antiphon{Mag.}{The wise men, being warned {\dag} in dreams by an Angel, departed into their own country another way.}

\subsubsec{Saturday within the Octave}
℣. We have seen his star in the east, alleluia.

℟. And are come with gifts to worship the Lord, alleluia.

\antiphon{Mag.}{The Child Jesus {\dag} tarried behind in Jerusalem; and Joseph and his mother knew not of it, supposing him to have been in the company; and they sought him among their kinsfolk and acquaintance.}

%\end{multicols}

\begin{rubric}
	When the Octave Day of the Epiphany falls on Sunday, the Office of the Sunday in the Octave is said on the preceding Saturday, and I Evensong of Sunday is said on Friday with Commemoration of the preceding day in the Octave.
	
	But if a I Double occur on this Saturday, the Office of Sunday is anticipated on the nearest day on which only the Offce of the Octave is to be said, and in the Office of the Feast Commemoration is made of the occurrent day within the Octave. In the Office of Sunday all is said as noted below.
\end{rubric}


\subby{Sunday within the Octave of the Epiphany}
\feastday{Epiphany I}
%\begin{inhead}
%{Semidouble}
%\end{inhead}

\begin{rubric}
	The Office Hymns are from the Feast of the Epiphany (p. \pageref{EpiphanyDay}), except for the following.
\end{rubric}

℣. All they from Sheba shall come, alleluia.

℟. They shall bring gold and incense, alleluia.

\properantiphon{Ben.}{The Child Jesus {\dag} tarried behind in Jerusalem; and Joseph and his mother knew not of it, supposing him to have been in the company; and they sought him among their kinsfolk and acquaintance.}\\

%MANUAL ADJUSTMENT:
\clearpage

℣. We have seen his star in the east, alleluia.

℟. And are come with gifts to worship the Lord, alleluia.

\properantiphon{Mag.}{Son, {\dag} why hast thou thus dealt with us? behold, thy father and I have sought thee sorrowing. How is it that ye sought me? wist ye not that I must be about my Father's business?}


\subby{Baptism of Our Lord}

\centerline{\small{(Octave Day of the Epiphany)}}

\feastday{Baptism of Our Lord}
%\begin{inhead}
%{Greater Double}
%\end{inhead}

\begin{rubric}
	The Daily Office propers are as in the Feast of the Epiphany (p. \pageref{EpiphanyDay}).
\end{rubric}


\begin{rubric}
	From the Octave of Epiphany until the Saturday before the I Sunday in Lent, when the Office is of the Feria, ail is said as in the Daily Hymns except that which is appointed as Proper. The Collect is that of the preceding Sunday.
\end{rubric}


\subby{Second Sunday after Epiphany}\label{EpiphanyII}
\feastday{Epiphany II}
%\begin{inhead}
%{Semidouble}
%\end{inhead}

\begin{rubric}
	The Office Hymns \& Versicles are from the Daily Hymns (p. \pageref{DailyHymns}), with the following Antiphons.
\end{rubric}

\properantiphon{Mag.}{God hath holpen {\dag} his servant Israel, as he promised to Abraham and to his seed: and hath exalted the humble for ever and ever.}

\properantiphon{Ben.}{There was a marriage {\dag} in Cana of Galilee, and Jesus was there with Mary his Mother.}

\properantiphon{Mag.}{And when they wanted wine, {\dag} Jesus bade them to fill the water-pots with water; and it was turned into wine, alleluia.}


\subby{Third Sunday after Epiphany}\label{epiphany}
\feastday{Epiphany III}
%\begin{inhead}
%{Semidouble}
%\end{inhead}

\properantiphon{Mag.}{God hath holpen {\dag} his servant Israel, as he promised to Abraham and to his seed: and hath exalted the humble for ever and ever.}

\properantiphon{Ben.}{When Jesus {\dag} was come down from the mountain, behold there came a leper and worshipped him, saying, Lord, if thou wilt, thou canst make me clean. And Jesus put forth his hand and touched him, saying: I will; be thou clean.}

\properantiphon{Mag.}{Lord, if thou wilt, {\dag} thou canst make me clean: then saith Jesus, I will; be thou clean.}


\subby{Fourth Sunday after Epiphany}
\feastday{Epiphany IV}
%\begin{inhead}
%{Semidouble}
%\end{inhead}

\properantiphon{Mag.}{God hath holpen {\dag} his servant Israel, as he promised to Abraham and to his seed: and hath exalted the humble for ever and ever.}

\properantiphon{Ben.}{And when Jesus {\dag} was entered into a ship, behold, there arose a great tempest in the sea: and his disciples awoke him, saying: Lord, save us, we perish.}

\properantiphon{Mag.}{Save us, Lord, {\dag} we perish: rebuke the winds and the sea, O God, and make a great calm.}


\subby{Fifth Sunday after Epiphany}\label{EpiphanyV}
\feastday{Epiphany V}
%\begin{inhead}
%{Semidouble}
%\end{inhead}

\properantiphon{Mag.}{God hath holpen {\dag} his servant Israel, as he promised to Abraham and to his seed: and hath exalted the humble for ever and ever.}

\properantiphon{Ben.}{Sir, {\dag} didst not thou sow good seed in thy field? from whence then hath it tares? He said unto them, An enemy hath done this.}

\properantiphon{Mag.}{Gather ye together {\dag} first the tares, and bind them in bundles to burn them: but gather the wheat into my barn, saith the Lord.}


\subby{Sixth Sunday after Epiphany}
\feastday{Epiphany VI}
%\begin{inhead}
%{Semidouble}
%\end{inhead}

\properantiphon{Mag.}{God hath holpen {\dag} his servant Israel, as he promised to Abraham and to his seed: and hath exalted the humble for ever and ever.}

\properantiphon{Ben.}{The kingdom of heaven {\dag} is like unto a grain of mustard seed, which is the least of all seeds; but when it is grown, it is the greatest among herbs.}

\properantiphon{Mag.}{The kingdom of heaven {\dag} is like unto leaven, which a woman took and hid in three measures of meal, until the whole was leavened.}

%MANUAL ADJUSTMENT:
\vspace{-1.5ex}

\subby{Septuagesima Sunday}
\fancyhead[RO,LE]{\textit{Septuagesima}}
\fancyhead[RE,LO]{Sunday}
%\begin{inhead}
%    {Second Class Semidouble}
%\end{inhead}

\properantiphon{Mag.}{The Lord said {\dag} unto Adam, Of the tree which is in the midst of the garden thou shalt not eat: for in the day that thou eatest thereof, thou shalt surely die.}

\properantiphon{Ben.}{The kingdom {\dag} of heaven is like unto a man that is an householder, which went out early in the morning to hire labourers into his vineyard, saith the Lord.}

\properantiphon{Mag.}{The householder {\dag} saith unto the labourers, Why stand ye here all the day idle? They say unto him, Because no man hath hired us. Go ye also into the vineyard, and whatsoever is right, I will give you.}

%MANUAL ADJUSTMENT:
\vspace{-0.5ex}

\subby{Septuagesimatide Ferial Office}
\begin{rubric}
On Ferias from Septuagesima Sunday until Ash Wednesday, exclusive, when the Office is not of the Feria, it is always commemorated at Mattins and Evensong on Double Feasts, even of the I Class, and on days within Octaves.
\end{rubric}
\begin{rubric}
	The \emph{Benedictus} Antiphons are said on Ferias as set forth in the Psalter, through Tuesday after Quinquagesima.
\end{rubric}
\begin{rubric}
	\textsc{Note,} No notice is taken of an occurrent Vigil, either in the Office of a Feast, or of a day within an Octave, or of a Feria. The Office of St. Mary on Saturday is not said.
\end{rubric}

%MANUAL ADJUSTMENT:
\vspace{-1.5ex}

\subby{Antiphons within Septuagesima Week}
\fancyhead[RO,LE]{Septuagesima Week}
\fancyhead[RE,LO]{}
%\begin{multicols}{2}

\subsubsec{Septuagesima Monday, Feria}
\antiphon{Mag.}{These last {\dag} have wrought but one hour, and thou hast made them equal unto us, which have borne the burden and heat of the day.}


\subsubsec{Septuagesima Tuesday, Feria}
\antiphon{Mag.}{The householder answered and said, {\dag} Friend, I do thee no wrong: didst not thou agree with me for a penny? Take that thine is, and go thy way.}


\subsubsec{Septuagesima Wednesday, Feria}
\antiphon{Mag.}{Take that thine is {\dag} and go thy way: because I am good, saith the Lord.}


\subsubsec{Septuagesima Thursday, Feria}
\antiphon{Mag.}{Is it not lawful {\dag} for me to do what I will? Is thine eye evil, because I am good? saith the Lord.}


\subsubsec{Septuagesima Friday, Feria}
\antiphon{Mag.}{So the last {\dag} shall be first, and the first shall be last: for many are called, but few are chosen.}

%\end{multicols}


\subby{Sexagesima Sunday}
\feastday{Sexagesima}
%\begin{inhead}
%{Second Class Semidouble}
%\end{inhead}

\properantiphon{Mag.}{The Lord {\dag} said unto Noah: The end of all flesh is come before me: make thee an ark of gopher wood, that therein the seed of all flesh may be saved.}

\properantiphon{Ben.}{When much people {\dag} were gathered together to Jesus, and were come to him out of every city, he spake by a parable: A sower went out to sow his seed.}

\properantiphon{Mag.}{Unto you it is given {\dag} to know the mysteries of the kingdom of heaven: but to others in parables, said Jesus to his disciples.}


\subby{Antiphons within Sexagesima Week}
\fancyhead[RO,LE]{Sexagesima Week}
\fancyhead[RE,LO]{}
%\begin{multicols}{2}

\subsubsec{Sexagesima Monday, Feria}
\antiphon{Mag.}{If ye seek {\dag} the summit of true honour, hasten to yon heavenly country with what speed ye may.}


\subsubsec{Sexagesima Tuesday, Feria}
\antiphon{Mag.}{The seed {\dag} is the word of God, but Christ is the Sower: every one that heareth him shall abide for ever.}


\subsubsec{Sexagesima Wednesday, Feria}
\antiphon{Mag.}{But that {\dag} which fell on the good ground are they which in an honest and good heart receive the word; and bring forth fruit with patience.}


\subsubsec{Sexagesima Thursday, Feria}
\antiphon{Mag.}{Some seed {\dag} fell on good ground, and brought forth fruit; some an hundred-fold, and some sixtyfold.}


\subsubsec{Sexagesima Friday, Feria}
\antiphon{Mag.}{They who keep the word of God {\dag} with an honest and perfect heart bring forth fruit with patience.}
%\end{multicols}


\subby{Quinquagesima Sunday}
\feastday{Quinquagesima}
%\begin{inhead}
%{Second Class Semidouble}
%\end{inhead}

\properantiphon{Mag.}{Mighty Abraham, {\dag} the father of our faith, offered a burnt offering upon the altar, instead of his son.}

\properantiphon{Ben.}{Behold, we go up {\dag} to Jerusalem, and all things that are written concerning the Son of man shall be accomplished: for he shall be delivered unto the Gentiles, and shall be mocked, and spitted on; and they shall scourge him, and put him to death; and the third day he shall rise again.}

\properantiphon{Mag.}{And Jesus stood, {\dag} and commanded him to be brought unto him, and asked him, saying: What wilt thou that I shall do unto thee? Lord, that I may receive my sight. And Jesus said unto him, Receive thy sight, thy faith hath saved thee. And straightway he received his sight, and followed him, glorifying God.}


\subby{Antiphons within Quinquagesima Week}
\fancyhead[RO,LE]{Quinquagesima Week}
\fancyhead[RE,LO]{}
%\begin{multicols}{2}

\subsubsec{Quinquagesima Monday, Feria}
\antiphon{Mag.}{And they which went before {\dag} rebuked him that he should hold his peace: but he cried so much the more, Have mercy on me, thou Son of David.}


\subsubsec{Quinquagesima Tuesday, Feria}
\antiphon{Mag.}{Have mercy on me, {\dag} thou Son of David. What wilt thou that I shall do unto thee? Lord, that I may receive my sight.}

%\end{multicols}


\begin{rubric}
	If on the following Ash Wednesday there occur a I or II Double, it is transferred to the first unhindered day. Greater and lesser Doubles and Memorials are only commemorated on the other days of Lent.
\end{rubric}


\subby{Ash Wednesday}
\fancyhead[RO,LE]{\textit{Ash Wednesday}}
\fancyhead[RE,LO]{}
%\begin{inhead}
%    {Greater Privileged Feria}
%\end{inhead}

\begin{rubric}
	On this day all Octaves cease until Holy Sabbath. On this and other Ferias through None of the following Saturday all is said as in the Psalter throughout the Year except that which is appointed as proper.
\end{rubric}

\begin{rubric}
	The Daily Office propers are of the Ferial Day, except the following Antiphons.
\end{rubric}

\properantiphon{Ben.}{When ye fast, {\dag} be not, as the hypocrites, of a sad countenance.}

\properantiphon{Mag.}{Lay up for yourselves {\dag} treasures in heaven, where neither moth nor rust doth corrupt.}



\subby{Thursday after Ash Wednesday}
\fancyhead[RO,LE]{\textit{Ash Wednesday}}
\fancyhead[RE,LO]{Thursday}
%\begin{inhead}
%    {Greater Feria}
%\end{inhead}

\begin{rubric}
	The Daily Office propers are of the Ferial Day, except the following Antiphons and Collects.
\end{rubric}

\begin{paracol}{2}

\sloppy

\begin{inhead}
	Mattins
\end{inhead}

\antiphon{Ben.}{Lord, my servant {\dag} lieth at home sick of the palsy, grievously tormented. Verily I say unto thee, I will come and heal him.}

\lett{O}{God,} who art offended by sin, and reconciled by penitence: mercifully regard the prayers of thy suppliant people, and turn away the scourge of thy wrath, which for our sins we have justly deserved. Through.

\switchcolumn

\begin{inhead}
	II Evensong
\end{inhead}

\antiphon{Mag.}{Lord, I am not worthy {\dag} that thou shouldest enter under my roof: but speak the word only and my servant shall be healed.}

\lett{S}{pare} us, O Lord, spare thy people: that they who are justly chastised by thy scourges, may be relieved by thy tender mercy. Through.

\fussy

\end{paracol}

\introit
\lett{W}{hen} I called upon the Lord, he heard my voice from the battle that was against me, yea, even God, that endureth for ever, shall bring them down: O cast thy burden upon the Lord, and he shall nourish thee. \textit{Ps.} Hear my prayer, O God, and hide not thyself from my petition: take heed unto me and hear me.

\collect
\lett{O}{God,} who art wroth with them that sin against thee, and sparest them that are penitent: mercifully look upon the prayers of thy people which call upon thee; and turn away the scourges of thy wrath which for our sins we justly deserve. Through.

\begin{rubric}
    \nth{2} Collect is \emph{Of Saints} (p. \SPSaints) \& \nth{3} \emph{Of the Living and Departed} (p. \SPLiving).
\end{rubric}

\readingcitation{Epistle}{Isaiah 38:1}
\lett{I}{n those days:} Was Hezekiah sick unto death. And Isaiah the prophet the son of Amoz came unto him, and said unto him, Thus saith the \divineName{Lord}, Set thine house in order: for thou shalt die, and not live. Then Hezekiah turned his face toward the wall, and prayed unto the \divineName{Lord}, And said, Remember now, O \divineName{Lord}, I beseech thee, how I have walked before thee in truth and with a perfect heart, and have done that which is good in thy sight. And Hezekiah wept sore. Then came the word of the \divineName{Lord} to Isaiah, saying, Go, and say to Hezekiah, Thus saith the \divineName{Lord}, the God of David thy father, I have heard thy prayer, I have seen thy tears: behold, I will add unto thy days fifteen years. And I will deliver thee and this city out of the hand of the king of Assyria: and I will defend this city, saith the Lord almighty.

\gradual{O cast thy burden upon the Lord, and he shall nourish thee. ℣. When I called upon the Lord, he heard my voice from the battle that was against me.}

\readingcitation{Gospel}{Matthew 8:5}
\lett{A}{t that time:} When Jesus was entered into Capernaum, there came unto him a centurion, beseeching him, And saying, Lord, my servant lieth at home sick of the palsy, grievously tormented. And Jesus saith unto him, I will come and heal him. The centurion answered and said, Lord, I am not worthy that thou shouldest come under my roof: but speak the word only, and my servant shall be healed. For I am a man under authority, having soldiers under me: and I say to this man, Go, and he goeth; and to another, Come, and he cometh; and to my servant, Do this, and he doeth it. When Jesus heard it, he marvelled, and said to them that followed, Verily I say unto you, I have not found so great faith, no, not in Israel. And I say unto you, That many shall come from the east and west, and shall sit down with Abraham, and Isaac, and Jacob, in the kingdom of heaven. But the children of the kingdom shall be cast out into outer darkness: there shall be weeping and gnashing of teeth. And Jesus said unto the centurion, Go thy way; and as thou hast believed, so be it done unto thee. And his servant was healed in the selfsame hour.

\offertory{Unto thee, O Lord, will I lift up my soul: my God, I have put my trust in thee, O let me not be confounded: neither let mine enemies triumph over me: for all they that hope in thee shall not be ashamed.}

\secret
\lett{W}{e} beseech thee, O Lord, favourably to regard these our sacrifices: that they may be profitable for our devotion and set forward our salvation. Through.

\begin{rubric}
    \nth{2} Secret is \emph{Of Saints} (p. \pageref{SPSaints}) \& \nth{3} \emph{Of the Living and Departed} (p. \pageref{SPLivingDeparted}).
\end{rubric}

\communion{Thou shalt be pleased with the sacrifice of righteousness, with the burnt-offerings and oblations upon thine altar, O Lord.}

\postcommunion
\lett{W}{e} humbly beseech thee, almighty God: that, as we have received the blessing of this heavenly gift; so it may be made to us thy sacrament, and avail to our salvation. Through.

\begin{rubric}
    \nth{2} Postcommunion is \emph{Of Saints} (p. \SPSaints) \& \nth{3} \emph{Of the Living and Departed} (p. \pageref{SPLivingDeparted}).
\end{rubric}

\textsc{Priest.} Let us pray.\par
\textsc{Deacon.} Humble your heads before God.\par
\begin{rubric}
    The Priest then prays the following:
\end{rubric}
\lett{S}{pare,} O Lord, spare thy people: that they who are justly chastised by thy scourges, may by thy mercy be relieved. Through.



\subby{Friday after Ash Wednesday}
\fancyhead[RO,LE]{\textit{Ash Wednesday}}
\fancyhead[RE,LO]{Friday}
%\begin{inhead}
%    {Greater Feria}
%\end{inhead}

\begin{rubric}
	The Daily Office propers are of the Ferial Day, except the following Antiphons and Collects.
\end{rubric}

\begin{paracol}{2}

\sloppy

\begin{inhead}
	Mattins
\end{inhead}

\antiphon{Ben.}{When thou doest thine alms, {\dag} let not thy left hand know what thy right. hand doeth.}

\lett{W}{e} beseech thee, O Lord, to accompany with thy bounteous favour the fast upon which we have entered: that the observance which we shew forth in our bodies, we may be able also to practise with sincerity of heart. Through.

\switchcolumn

\begin{inhead}
	II Evensong
\end{inhead}

\antiphon{Mag.}{But thou, when thou prayest, {\dag} enter into thy closet, and when thou hast shut thy door, pray to thy Father.}

\lett{D}{efend} thy people, O Lord, and mercifully cleanse them from all their sins: for no adversity will harm them over whom iniquity hath no dominion. Through.

\fussy

\end{paracol}

\introit
\lett{T}{he} Lord heard, and had mercy upon me: the Lord became my helper. \textit{Ps.} I will magnify thee, O Lord, for thou hast set me up: and not made my foes to triumph over me.

\collect
\lett{W}{e} beseech thee, O Lord, to further with thy gracious favour the fasts which we have begun: that as we keep this observance in the flesh, so we may have strength to perform the same in singleness of heart. Through.

\begin{rubric}
    \nth{2} Collect is \emph{Of Saints} (p. \pageref{SPSaints}) \& \nth{3} \emph{Of the Living and Departed} (p. \pageref{SPLivingDeparted}).
\end{rubric}

\readingcitation{Epistle}{Isaiah 58:1}
\lett{T}{hus saith the Lord God:} Cry aloud, spare not, lift up thy voice like a trumpet, and shew my people their transgression, and the house of Jacob their sins. Yet they seek me daily, and delight to know my ways, as a nation that did righteousness, and forsook not the ordinance of their God: they ask of me the ordinances of justice; they take delight in approaching to God. Wherefore have we fasted, say they, and thou seest not? wherefore have we afflicted our soul, and thou takest no knowledge? Behold, in the day of your fast ye find pleasure, and exact all your labours. Behold, ye fast for strife and debate, and to smite with the fist of wickedness: ye shall not fast as ye do this day, to make your voice to be heard on high. Is it such a fast that I have chosen? a day for a man to afflict his soul? is it to bow down his head as a bulrush, and to spread sackcloth and ashes under him? wilt thou call this a fast, and an acceptable day to the \divineName{Lord}? Is not this the fast that I have chosen? to loose the bands of wickedness, to undo the heavy burdens, and to let the oppressed go free, and that ye break every yoke? Is it not to deal thy bread to the hungry, and that thou bring the poor that are cast out to thy house? when thou seest the naked, that thou cover him; and that thou hide not thyself from thine own flesh? Then shall thy light break forth as the morning, and thine health shall spring forth speedily: and thy righteousness shall go before thee; the glory of the \divineName{Lord} shall be thy reward. Then shalt thou call, and the \divineName{Lord} shall answer; thou shalt cry, and he shall say, Here I am. For I the \divineName{Lord} thy God am merciful.

\gradual{One thing have I desired of the Lord, which I will require, even that I may dwell in the house of the Lord. ℣. To behold the fair beauty of the Lord, and to hide me in his holy temple.}

\tract{O Lord, deal not with us after our sins: nor reward us according to our wickednesses. ℣. O Lord, remember not our old sins: but have mercy upon us, and that soon, for we are come to great misery. \inrub{Here genuflect.} ℣. Help us, O God of our salvation: and for the glory of thy name, O Lord, deliver us: and be merciful unto our sins, for thy name's sake.}

\readingcitation{Gospel}{Matthew 5:43}
\lett{A}{t that time:} Jesus said unto his disciples: Ye have heard that it hath been said, Thou shalt love thy neighbour, and hate thine enemy. But I say unto you, Love your enemies, bless them that curse you, do good to them that hate you, and pray for them which despitefully use you, and persecute you; That ye may be the children of your Father which is in heaven: for he maketh his sun to rise on the evil and on the good, and sendeth rain on the just and on the unjust. For if ye love them which love you, what reward have ye? do not even the publicans the same? And if ye salute your brethren only, what do ye more than others? do not even the publicans so? Be ye therefore perfect, even as your Father which is in heaven is perfect. Take heed that ye do not your alms before men, to be seen of them: otherwise ye have no reward of your Father which is in heaven. Therefore when thou doest thine alms, do not sound a trumpet before thee, as the hypocrites do in the synagogues and in the streets, that they may have glory of men. Verily I say unto you, They have their reward. But when thou doest alms, let not thy left hand know what thy right hand doeth: That thine alms may be in secret: and thy Father which seeth in secret himself shall reward thee openly.

\offertory{Quicken me, O Lord, according to thy word: that I may know thy testimonies.}

\secret
\lett{G}{rant,} we beseech thee, O Lord, that this sacrifice of Lenten observance which we offer: may both render our souls acceptable unto thee, and make us more readily to serve thee in continence. Through.

\begin{rubric}
    \nth{2} Secret is \emph{Of Saints} (p. \pageref{SPSaints}) \& \nth{3} \emph{Of the Living and Departed} (p. \pageref{SPLivingDeparted}).
\end{rubric}

\communion{Serve the Lord in fear, and rejoice unto him with reverence: lay hold on discipline, lest ye perish from the right way.}

\postcommunion
\lett{P}{our} forth upon us, O Lord, the Spirit of thy charity: that as thou hast fulfilled us with one heavenly bread, so of thy goodness thou wouldest make us to be of one heart and mind. Through . . . in the unity of the same.

\begin{rubric}
    \nth{2} Postcommunion is \emph{Of Saints} (p. \pageref{SPSaints}) \& \nth{3} \emph{Of the Living and Departed} (p. \pageref{SPLivingDeparted}).
\end{rubric}

\textsc{Priest.} Let us pray.\par
\textsc{Deacon.} Humble your heads before God.\par
\begin{rubric}
    The Priest then prays the following:
\end{rubric}
\lett{D}{efend,} O Lord, thy people, and mercifully cleanse them from all their sins: that no adversity may harm them, over whom iniquity hath no dominion. Through.


\subby{Saturday after Ash Wednesday}
\fancyhead[RO,LE]{\textit{Ash Wednesday}}
\fancyhead[RE,LO]{Saturday}
%\begin{inhead}
%    {Greater Feria}
%\end{inhead}

\begin{rubric}
	In the Daily Office, the propers are of the Ferial Day, except the following Antiphon.
\end{rubric}

\antiphon{Ben.}{Yet they seek me daily, {\dag} and delight to know my ways.}

\introit
\lett{T}{he} Lord heard, and had mercy upon me: the Lord became my helper. \textit{Ps.} I will magnify thee, O Lord, for thou hast set me up: and not made my foes to triumph over me.

\collect
\lett{A}{ssist} us, O Lord, in these our supplications: and grant; that like as this solemn fast hath been ordained for the safety and healing of our bodies and our souls, so we may with devout observance celebrate the same. Through.

\begin{rubric}
    \nth{2} Collect is \emph{Of Saints} (p. \pageref{SPSaints}) \& \nth{3} \emph{Of the Living and Departed} (p. \pageref{SPLivingDeparted}).
\end{rubric}

\readingcitation{Epistle}{Isaiah 58:9}
\lett{T}{hus saith the Lord God:} If thou take away from the midst of thee the yoke, the putting forth of the finger, and speaking vanity; And if thou draw out thy soul to the hungry, and satisfy the afflicted soul; then shall thy light rise in obscurity, and thy darkness be as the noonday: And the \divineName{Lord} shall guide thee continually, and satisfy thy soul in drought, and make fat thy bones: and thou shalt be like a watered garden, and like a spring of water, whose waters fail not. And they that shall be of thee shall build the old waste places: thou shalt raise up the foundations of many generations; and thou shalt be called, The repairer of the breach, The restorer of paths to dwell in. If thou turn away thy foot from the sabbath, from doing thy pleasure on my holy day; and call the sabbath a delight, the holy of the \divineName{Lord}, honourable; and shalt honour him, not doing thine own ways, nor finding thine own pleasure, nor speaking thine own words: Then shalt thou delight thyself in the \divineName{Lord}; and I will cause thee to ride upon the high places of the earth, and feed thee with the heritage of Jacob thy father: for the mouth of the \divineName{Lord} hath spoken it.

\gradual{One thing have I desired of the Lord, which I will require, even that I may dwell in the house of the Lord. ℣. To behold the fair beauty of the Lord, and to visit his temple.}


\readingcitation{Gospel}{Mark 6:47}
\lett{A}{t that time:} When even was come, the ship was in the midst of the sea, and Jesus alone on the land. And he saw them toiling in rowing; for the wind was contrary unto them: and about the fourth watch of the night he cometh unto them, walking upon the sea, and would have passed by them. But when they saw him walking upon the sea, they supposed it had been a spirit, and cried out: For they all saw him, and were troubled. And immediately he talked with them, and saith unto them, Be of good cheer: it is I; be not afraid. And he went up unto them into the ship; and the wind ceased: and they were sore amazed in themselves beyond measure, and wondered. For they considered not the miracle of the loaves: for their heart was hardened. And when they had passed over, they came into the land of Gennesaret, and drew to the shore. And when they were come out of the ship, straightway they knew him, And ran through that whole region round about, and began to carry about in beds those that were sick, where they heard he was. And whithersoever he entered, into villages, or cities, or country, they laid the sick in the streets, and besought him that they might touch if it were but the border of his garment: and as many as touched him were made whole.

\offertory{Quicken me, O Lord, according to thy word: that I may know thy testimonies.}

\secret
\lett{A}{ccept,} O Lord, the sacrifice which thou hast ordained to be a worthy propitiation unto thee: and grant, we beseech thee, that we being cleansed by the operation of the same, may offer unto thee the acceptable devotion of our hearts. Through.

\begin{rubric}
    \nth{2} Secret is \emph{Of Saints} (p. \pageref{SPSaints}) \& \nth{3} \emph{Of the Living and Departed} (p. \pageref{SPLivingDeparted}).
\end{rubric}

\communion{Serve the Lord in fear, and rejoice unto him with reverence: lay hold on discipline, lest ye perish from the right wav.}

\postcommunion
\lett{O}{Lord,} who hast quickened us with the gift of heavenly life: we beseech thee, that those things which in this present life are to us a mystery, may be our succour unto life eternal. Through.

\begin{rubric}
    \nth{2} Postcommunion is \emph{Of Saints} (p. \pageref{SPSaints}) \& \nth{3} \emph{Of the Living and Departed} (p. \pageref{SPLivingDeparted}).
\end{rubric}

\textsc{Priest.} Let us pray.\par
\textsc{Deacon.} Humble your heads before God.\par
\begin{rubric}
    The Priest then prays the following:
\end{rubric}
\lett{M}{ay} thy faithful people, O God, be strengthened by thy gifts: that they receiving the same may seek them the more, and seeking them may obtain them everlastingly. Through.


\subby{First Sunday of Lent}\label{LentI}
\fancyhead[RO,LE]{\textit{Lent I}}
\fancyhead[RE,LO]{}
%\begin{inhead}
%    {First Class Semidouble}
%\end{inhead}

\subbysub{I Evensong}\label{FirstLentEvensong}

\gregorioscore{resources/gabc/ProperSeason/FirstLentEvensong.gabc}

℣. God shall give his Angels charge over thee.

℟. To keep thee in all thy ways.

\properantiphon{Mag.}{Then shalt thou call, {\dag} and the Lord shall answer: thou shalt cry, and he shall say, Here am I.}

\subbysub{Mattins}

%MANUAL ADJUSTMENT:
%\invitatoryhymn\label{FirstLentInvitatory}

\begin{inhead}
	Invitatory Hymn\label{FirstLentInvitatory}
\end{inhead}

\gregorioscore{resources/gabc/ProperSeason/FirstLentInvitatory.gabc}

\officehymn\label{FirstLentMattins}

\gregorioscore{resources/gabc/ProperSeason/FirstLentMattins.gabc}

℣. God shall give his Angels charge over thee.

℟. To keep thee in all thy ways.

\properantiphon{Ben.}{Then was Jesus {\dag} led up of the Spirit into the wilderness to be tempted of the devil: and when he had fasted forty days and forty nights, he was afterward an hungred.}

\subbysub{II Evensong}

\begin{rubric}
	The Office Hymn \& Versicle are of I Evensong, with the following Antiphon.
\end{rubric}

\properantiphon{Mag.}{Behold, now {\dag} is the accepted time; behold, now is the day of salvation: let us therefore in all things approve ourselves as the ministers of God, in much patience, in watchings, in fastings, and by love unfeigned.}


\subby{Lenten Ferial Office}
\begin{rubric}
	On Lenten Ferias, the Hymns and Versicles are taken from the First Sunday of Lent. The Antiphon and Collect are of the Day.
\end{rubric}

\subby{Ember Wednesday in Lent}
\feastday{Lenten Emberday}
\fancyhead[RE,LO]{Ember Wednesday}

\properantiphon{Ben.}{This crooked {\dag} and perverse generation seeketh after a sign: and there shall no sign be given it, but the sign of the prophet Jonas.}

\properantiphon{Mag.}{For as Jonas {\dag} was three days and three nights in the whale's belly, so shall the Son of Man be in the heart of the earth.}


\subby{Ember Friday in Lent}
\feastday{Lenten Emberday}
\fancyhead[RE,LO]{Ember Friday}
%\bcpfeast{Ember Friday in Lent}{Ember Friday}{Lenten Emberday}

\properantiphon{Ben.}{An Angel of the Lord {\dag} went down from heaven, and troubled the water, and one was healed.}

\properantiphon{Mag.}{He that made me whole, {\dag} the same commanded me: Take up thy bed, and walk in peace.}

\subby{Ember Saturday in Lent}
\feastday{Lenten Emberday}
\fancyhead[RE,LO]{Ember Saturday}
%\bcpfeast{Ember Saturday in Lent}{Ember Saturday}{Lenten Emberday}

\properantiphon{Ben.}{And Jesus taketh his disciples, {\dag} and goeth up into a mountain, and was transfigured before them.}


\subby{Second Sunday of Lent}
\fancyhead[RO,LE]{\textit{Lent II}}
\fancyhead[RE,LO]{}
%\begin{inhead}
%    {First Class Semidouble}
%\end{inhead}

\begin{rubric}
	The Daily Office propers are as in the First Sunday of Lent (p. \pageref{LentI}), except for the following Antiphons.
\end{rubric}

\properantiphon{Mag.}{Tell the vision {\dag} which ye have seen to no man, until the Son of Man be risen again from the dead.}

\properantiphon{Ben.}{Jesus went thence, {\dag} and departed into the coasts of Tyre and Sidon: and behold, a woman of Canaan came out of the same coasts, and cried, saying: Have mercy on me, thou Son of David.}

\properantiphon{Mag.}{Jesus said {\dag} unto the woman of Canaan, It is not meet to take the children's bread, and to cast it to dogs. Truth, Lord; yet the dogs eat of the crumbs which fall from their master's table. Then Jesus answered, O woman, great is thy faith; be it unto thee even as thou wilt.}


\subby{Third Sunday of Lent}
\fancyhead[RO,LE]{\textit{Lent III}}
\fancyhead[RE,LO]{}
%\begin{inhead}
%    {First Class Semidouble}
%\end{inhead}

\begin{rubric}
	The Daily Office propers are as in the First Sunday of Lent (p. \pageref{LentI}), except for the following Antiphons.
\end{rubric}

\properantiphon{Mag.}{But the father {\dag} said to his servants, Bring forth the best robe and put it on him; and put a ring on his hand, and shoes on his feet.}

\properantiphon{Ben.}{When a strong man armed {\dag} keepeth his palace, all his goods are in peace.}

\properantiphon{Mag.}{A certain woman {\dag} of the company lifted up her voice and cried, Blessed is the womb that bare thee, and the paps which thou hast sucked. But Jesus answered, Yea, rather, blessed are they that hear the word of God, and keep it.}


\subby{Fourth Sunday of Lent}
\fancyhead[RO,LE]{\textit{Lent IV}}
\fancyhead[RE,LO]{}
%\begin{inhead}
%    {First Class Semidouble}
%\end{inhead}

\begin{rubric}
	The Daily Office propers are as in the First Sunday of Lent (p. \pageref{LentI}), except for the following Antiphons.
\end{rubric}

\properantiphon{Mag.}{Woman, {\dag} hath no man condemned thee? No man, Lord. Neither do I condemn thee; go and sin no more.}

\properantiphon{Ben.}{When Jesus {\dag} lifted up his eyes, and saw a great company come unto him, he saith unto Philip: Whence shall we buy bread that these may eat? And this he said to prove him; for he himself knew what he would do.}

\properantiphon{Mag.}{And Jesus {\dag} went up into a mountain, and there he sat with his disciples.}


\subby{Passion Sunday}
\fancyhead[RO,LE]{\textit{Passion Sunday}}
\fancyhead[RE,LO]{}
%\begin{inhead}
%    {First Class Semidouble}
%\end{inhead}

\subbysub{I Evensong}\label{PassionSundayEvensong}

\gregorioscore{resources/gabc/ProperSeason/PassionSundayEvensong1.gabc}

\begin{rubric}
	During the following stanza, all kneel.
\end{rubric}

\gregorioscore{resources/gabc/ProperSeason/PassionSundayEvensongPenultimateVerse.gabc}

\begin{rubric}
	\textsc{Note,} The following stanza is never changed.
\end{rubric}

\gregorioscore{resources/gabc/ProperSeason/PassionSundayEvensongUltimateVerse.gabc}

℣. Deliver me, O Lord, from the evil man.

℟. And preserve me from the wicked man.

\properantiphon{Mag.}{I am one {\dag} that bear witness of myself, and the Father that sent me beareth witness of me.}

\subbysub{Mattins}

%Tune from Lauds
\invitatoryhymn\label{PassionSundayInvitatory}

\gregorioscore{resources/gabc/ProperSeason/PassionSundayInvitatory.gabc}

\officehymn\label{PassionSundayMattins}

\gregorioscore{resources/gabc/ProperSeason/PassionSundayMattins.gabc}

℣. Deliver me from mine enemies, O God.

℟. Defend me from them that rise up against me.

\properantiphon{Ben.}{Jesus said {\dag} unto the multitude of the Jews, and to the chief priests: He that is of God heareth God's words; ye therefore hear them not, because ye are not of God.}

\subbysub{II Evensong}

\begin{rubric}
	The Office Hymn \& Versicle are of I Evensong, with the following Antiphon.
\end{rubric}

\properantiphon{Mag.}{Your father Abraham {\dag} rejoiced to see my day: and he saw it and was glad.}


\subby{Passiontide Ferial Office}
\begin{rubric}
	On Passiontide Ferias, the Hymns and Versicles are taken from Passion Sunday. The Antiphon and Collect are of the Day.
\end{rubric}


\subby{Palm Sunday}
\fancyhead[RO,LE]{\textit{Palm Sunday}}
\fancyhead[RE,LO]{}
%\begin{inhead}
%    {First Class Semidouble}
%\end{inhead}

\begin{rubric}
	The Daily Office propers are as in Passion Sunday, except for the following Antiphons.
\end{rubric}

\properantiphon{Mag.}{Righteous Father, {\dag} the world hath not known thee: but I have known thee, because thou hast sent me.}

\properantiphon{Ben.}{The multitudes {\dag} which came together for the feast day cried unto the Lord: Blessed is he that cometh in the Name of the Lord. Hosanna in the highest.}

\properantiphon{Mag.}{For it is written, {\dag} I will smite the shepherd, and the sheep shall be scattered. But after I am risen again, I will go before you into Galilee: there shall ye see me, saith the Lord.}


\subby{Easter Day}\label{EasterDay}
\fancyhead[RO,LE]{\textit{Easter Day}}
\fancyhead[RE,LO]{}
%\begin{inhead}
%    {First Class Double, First Class Octave}
%\end{inhead}

\subbysub{I Evensong}

\begin{rubric}
	If I Evensong of Easter Day be not said during the service of Easter Even, then it is said as usual with the following Antiphon, omitting the Hymn and Versicle.
\end{rubric}

\properantiphon{Mag.}{In the end of the sabbath, {\dag} as it began to dawn toward the first day of the week, came Mary Magdalene and the other Mary to see the sepulchre, alleluia.}

\subbysub{Mattins}

\invitatoryhymn\label{EasterInvitatory}

\gregorioscore{resources/gabc/ProperSeason/EasterInvitatory.gabc}

\officehymn\label{EasterMattins}

\gregorioscore{resources/gabc/ProperSeason/EasterMattins.gabc}

℣. This is the day which the Lord hath made, alleluia.

℟. We will rejoice and be glad in it, alleluia.

\properantiphon{Ben.}{And very early in the morning {\dag} the first day of the week, they came unto the sepulchre at the rising of the sun, alleluia.}

\subbysub{II Evensong}\label{EasterEvensong}

\gregorioscore{resources/gabc/ProperSeason/EasterEvensong.gabc}

℣. This is the day which the Lord hath made, alleluia.

℟. We will rejoice and be glad in it, alleluia.

\properantiphon{Mag.}{And when they looked, {\dag} they saw that the stone was rolled away: for it was very great, alleluia.}

\begin{rubric}
	During Easter Week, in the Daily Office, the Office Hymns \& Versicles are from Easter Day.
\end{rubric}


\subby{Antiphons within Easter Week}
\fancyhead[RO,LE]{Easter Week}
\fancyhead[RE,LO]{}
\begin{multicols}{2}

\subbysub{Easter Monday}
\properantiphon{Ben.}{Jesus himself {\dag} drew near to his disciples in the way, and went with them: but their eyes were holden that they should not know him: and he rebuked them, saying, O fools and slow of heart to believe all that the prophets have spoken, alleluia.}

\properantiphon{Mag.}{What manner of communications {\dag} are these that ye have one to another, as ye walk, and are sad? alleluia, alleluia.}

\subbysub{Easter Tuesday}
\properantiphon{Ben.}{Jesus stood {\dag} in the midst of his disciples and said unto them, Peace be unto you, alleluia, alleluia.}

\properantiphon{Mag.}{Behold my hands {\dag} and my feet, that it is I myself, alleluia.}

\subbysub{Easter Wednesday}
\properantiphon{Ben.}{Cast the net {\dag} on the right side of the ship, and ye shall find, alleluia.}

\properantiphon{Mag.}{Jesus saith unto his disciples, {\dag} Bring of the fish which ye have now caught. Simon Peter went up, and drew the net to land, full of great fishes, alleluia.}

\subbysub{Easter Thursday}
\properantiphon{Ben.}{Mary stood {\dag} at the sepulchre weeping, and seeth two Angels in white, sitting, and the napkin that was about the head of Jesus, alleluia.}

\properantiphon{Mag.}{They have taken {\dag} away my Lord, and I know not where they have laid him. If thou have borne him hence, tell me, alleluia: and I will take him away, alleluia.}

\subbysub{Easter Friday}
\properantiphon{Ben.}{The eleven disciples, {\dag} when they saw the Lord in Galilee, worshipped him, alleluia.}

\properantiphon{Mag.}{All power {\dag} is given unto me in heaven and in earth, alleluia.}

\subbysub{Easter Saturday}
\properantiphon{Ben.}{They ran both together, {\dag} and the other disciple did outrun Peter, and came first to the sepulchre, alleluia.}

\properantiphon{Mag.}{The same day {\dag} at evening, being the first day of the week, when the doors were shut where the disciples were assembled, came Jesus and stood in the midst, and saith unto them, Peace be unto you, alleluia.}
\end{multicols}


\subby{Low Sunday}
\feastday{Easter I}
%\begin{inhead}
%    {First Class Double, Octave Day of Easter}
%\end{inhead}

\begin{rubric}
	In the Daily Office, the Hymns are of Easter Day (p. \pageref{EasterDay}), with the Versicles and Antiphons as followeth.
\end{rubric}

℣. Abide with us, Lord, alleluia. 

℟. For it is toward evening, alleluia.

\properantiphon{Mag.}{The same day {\dag} at evening, being the first day of the week, when the doors were shut where the disciples were assembled, came Jesus and stood in the midst, and saith unto them, Peace be unto you, alleluia.}\\

%MANUAL ADJUSTMENT:
\clearpage
℣. In thy resurrection, O Christ, alleluia.

℟. Let heaven and earth rejoice, alleluia.

\properantiphon{Ben.}{The same day {\dag} at evening, being the first day of the week, when the doors were shut where the disciples were assembled, came Jesus and stood in the midst, and saith unto them, Peace be unto you, alleluia.}\\

℣. Abide with us, Lord, alleluia. 

℟. For it is toward evening, alleluia.

\properantiphon{Mag.}{After eight days, {\dag} when the doors were shut, the Lord entered, and said unto them, Peace be unto you, alleluia, alleluia.}


\subby{Eastertide Ferial Office}
\begin{rubric}
	On Ferias in Eastertide, the Hymns are as in Easter Day (p. \pageref{EasterDay}). The Versicles are as followeth. The Antiphon is proper.
\end{rubric}

\begin{inhead}
	Mattins
\end{inhead}

℣. In thy resurrection, O Christ, alleluia.

℟. Let heaven and earth rejoice, alleluia.

\begin{inhead}
	Evensong
\end{inhead}

℣. Abide with us, Lord, alleluia. 

℟. For it is toward evening, alleluia.

\subby{The First Week after the Easter Octave}
\feastday{Easter I Week}

\begin{multicols}{2}
\subbysub{Monday}
\antiphon{Ben.}{When Jesus was risen {\dag} early the first day of the week, he appeared first to Mary Magdalene, out of whom he had cast seven devils, alleluia.}

\antiphon{Mag.}{Peace be unto you, it is I, {\dag} alleluia: be not afraid, alleluia.}

\subbysub{Tuesday}
\antiphon{Ben.}{I go before you {\dag} into Galilee: there shall ye see me, as I said unto you, alleluia, alleluia.}

\antiphon{Mag.}{Reach hither thy finger, {\dag} and examine the print of the nails, alleluia: and be not faithless, but believing, alleluia.}

\subbysub{Wednesday}
\antiphon{Ben.}{I am the true vine, {\dag} alleluia: and ye are the true branches, alleluia}

\antiphon{Mag.}{Because thou hast seen me, {\dag} Thomas, thou hast believed: blessed are they that have not seen, and yet have believed, alleluia.}

\subbysub{Thursday}
\antiphon{Ben.}{My heart burns {\dag} within me; I desire to behold my Lord: I seek, and find not where they have laid him, alleluia, alleluia.}

\antiphon{Mag.}{I did put my finger {\dag} into the print of the nails, and thrust my hand into his side, and said, My Lord and my God, alleluia.}

\subbysub{Friday}
\antiphon{Ben.}{There came unto the tomb {\dag} Mary Magdalene and the other Mary, to see the sepulchre, alleluia.}

\begin{rubric}
	At Evensong, unless a Double Feast or some Octave occur on the following day, the Office is of Saint Mary, and this Office is said on the following Saturday. The same order is observed on the other Saturdays through the Saturday before Easter V.
\end{rubric}
	
\end{multicols}


\subby{Second Sunday after Easter}\label{EasterII}
\feastday{Easter II}
%\begin{inhead}
%    {Semidouble}
%\end{inhead}

\begin{rubric}
	In the Daily Office, the Hymns are of Easter Day (p. \pageref{EasterDay}), with the following.
\end{rubric}

℣. Abide with us, Lord, alleluia. 

℟. For it is toward evening, alleluia.

\properantiphon{Mag.}{I am {\dag} the Shepherd of the sheep: I am the way, the truth, and the life: I am the Good Shepherd; and I know my sheep, and am known of mine, alleluia, alleluia.}\\

℣. In thy resurrection, O Christ, alleluia.

℟. Let heaven and earth rejoice, alleluia.

\properantiphon{Ben.}{I am {\dag} the Shepherd of the sheep: I am the way, the truth, and the life: I am the Good Shepherd; and I know my sheep, and am known of mine, alleluia, alleluia.}\\

℣. Abide with us, Lord, alleluia. 

℟. For it is toward evening, alleluia.

\properantiphon{Mag.}{I am {\dag} the Good Shepherd, who feed my sheep; and I lay down my life for my sheep, alleluia.}


\subby{The Second Week after the Easter Octave}
\feastday{Easter II Week}

\begin{multicols}{2}
\subsubsec{Monday}
\antiphon{Ben.}{Go ye into the world, {\dag} alleluia, and teach all nations, alleluia.}

\antiphon{Mag.}{The Good Shepherd {\dag} layeth down his life for the sheep, alleluia.}

\subsubsec{Tuesday}
\antiphon{Ben.}{Go ye into the world {\dag} and teach all nations, baptizing them in the Name of the Father, and of the Son, and of the Holy Ghost, alleluia.}

\antiphon{Mag.}{He that is an hireling, {\dag} whose own the sheep are not, seeth the wolf coming, and leaveth the sheep, and fleeth: and the wolf catcheth them, and, scattereth the sheep, alleluia.}

\subsubsec{Wednesday}
\antiphon{Ben.}{Go unto my brethren {\dag} and say unto them, alleluia: that they go into Galilee, alleluia: there shall they see me, alleluia, alleluia, alleluia.}

\antiphon{Mag.}{As the Father knoweth me, {\dag} even so know I the Father: and I lay down my life for the sheep, alleluia.}

\subsubsec{Thursday}
\antiphon{Ben.}{Art thou only {\dag} a stranger, and hast not heard concerning Jesus, how they delivered him to be condemned to death? alleluia.}

\antiphon{Mag.}{Other sheep I have, {\dag} which are not of this fold: them also I must bring, and they shall hear my voice; and there shall be one fold and one Shepherd, alleluia.}

\subsubsec{Friday}
\antiphon{Ben.}{Ought not Christ {\dag} to have suffered these things, and to enter into his glory? alleluia.}

\end{multicols}


\subby{Patronage of St. Joseph}
\feastday{Patronage of St. Joseph}
%\begin{inhead}
%    {Wednesday after the Second Sunday after Easter}\par
%    {First Class Double, Simple Octave}
%\end{inhead}

\subbysub{I Evensong}\label{PatronageEvensong}

\gregorioscore{resources/gabc/ProperSeason/JosephEvensong.gabc}

℣. I do give praise unto thy name, alleluia.

℟. For thou art my defender and helper, alleluia.

%From the Breviarium Monasticum, using KJV and Coverdale translation:
\properantiphon{Mag.}{When Mary {\dag} the mother of Jesus was espoused to Joseph, before they came together, she was found with child of the Holy Ghost, alleluia.}

\subbysub{Mattins}

\begin{rubric}
	The Invitatory Hymn is of I Evensong.
\end{rubric}

\officehymn\label{PatronageMattins}

\gregorioscore{resources/gabc/ProperSeason/JosephMattins.gabc}

℣. Thou hast given me the defence of thy salvation, alleluia.

℟. And thy right hand also shall hold me up, alleluia.

\properantiphon{Ben.}{Joseph, thou son of David, {\dag} fear not to take unto thee Mary thy wife: for that which is conceived in her is of the Holy Ghost, alleluia.}

\subbysub{II Evensong}

\begin{rubric}
	The Office Hymn is of I Evensong, with the following Versicle \& Antiphon.
\end{rubric}

℣. I sat down under his shadow with great delight, alleluia.

℟. And his fruit was sweet to my taste, alleluia.

\properantiphon{Mag.}{Son, {\dag} why hast thou thus dealt with us? Behold, thy father and I have sought thee sorrowing, alleluia.}


\subby{Third Sunday after Easter}
\feastday{Easter III}
%\begin{inhead}
%    {Semidouble}
%\end{inhead}

\begin{rubric}
	In the Daily Office, the Hymns are of Easter Day (p. \pageref{EasterDay}), with the following.
\end{rubric}

℣. Abide with us, Lord, alleluia. 

℟. For it is toward evening, alleluia.

\properantiphon{Mag.}{A little while, {\dag} and ye shall not see me, saith the Lord: and again, a little while, and ye shall see me, because I go to the Father, alleluia, alleluia.}\\

℣. In thy resurrection, O Christ, alleluia.

℟. Let heaven and earth rejoice, alleluia.

\properantiphon{Ben.}{A little while, {\dag} and ye shall not see me, saith the Lord: and again, a little while, and ye shall see me, because I go to the Father, alleluia, alleluia.}\\

℣. Abide with us, Lord, alleluia. 

℟. For it is toward evening, alleluia.

\properantiphon{Mag.}{Verily I say unto you, {\dag} that ye shall weep and lament, but the world shall rejoice: and ye shall be sorrowful, but your sorrow shall be turned into joy, alleluia.}

\subby{The Third Week after the Easter Octave}
\feastday{Easter III Week}

\begin{multicols}{2}
\subsubsec{Monday}
\antiphon{Ben.}{And beginning {\dag} at Moses and all the prophets, he expounded unto them the scriptures concerning himself, alleluia.}

\antiphon{Mag.}{Your sorrow {\dag} shall be turned into joy, alleluia: and your joy no man taketh from you, alleluia, alleluia.}

\subsubsec{Tuesday}
\antiphon{Ben.}{And they constrained him, {\dag} saying, Abide with us, O Lord, for it is toward evening, alleluia.}

\antiphon{Mag.}{Sorrow {\dag} hath filled your heart: and your joy no man taketh from you, alleluia, alleluia.}

\subsubsec{Wednesday}
\antiphon{Ben.}{Abide with us, {\dag} for it is toward evening, and the day is far spent, alleluia.}

\antiphon{Mag.}{Your sorrow, {\dag} alleluia, shall be turned into joy, alleluia.}

\subsubsec{Thursday}
\antiphon{Ben.}{And he went in {\dag} to tarry with them: and it came to pass, as he sat at meat with them, he took bread, and blessed it, and brake, and gave to them, alleluia, alleluia.}

\antiphon{Mag.}{Verily, verily, I say unto you, {\dag} I will see you again, and your heart shall rejoice: and your joy no man taketh from you, alleluia.}

\subsubsec{Friday}
\antiphon{Ben.}{They knew the Lord Jesus, {\dag} alleluia, in breaking of bread, alleluia.}

\end{multicols}


\subby{Fourth Sunday after Easter}
\feastday{Easter IV}
%\begin{inhead}
%    {Semidouble}
%\end{inhead}

\begin{rubric}
	In the Daily Office, the Hymns are of Easter Day (p. \pageref{EasterDay}), with the following.
\end{rubric}

℣. Abide with us, Lord, alleluia. 

℟. For it is toward evening, alleluia.

\properantiphon{Mag.}{Now I go my way {\dag} to him that sent me; and none of you asketh me, Wither goest thou? alleluia, alleluia.}\\

℣. In thy resurrection, O Christ, alleluia.

℟. Let heaven and earth rejoice, alleluia.

\properantiphon{Ben.}{Now I go my way {\dag} to him that sent me; and none of you asketh me, Wither goest thou? alleluia, alleluia.}\\

℣. Abide with us, Lord, alleluia. 

℟. For it is toward evening, alleluia.

\properantiphon{Mag.}{Now I go my way {\dag} to him that sent me: but because I have said these things unto you, sorrow hath filled your heart, alleluia.}


\subby{The Fourth Week after the Easter Octave}
\feastday{Easter IV Week}

\begin{multicols}{2}
\subsubsec{Monday}
\antiphon{Ben.}{Did not our heart {\dag} burn within us concerning Jesus, while he talked with us by the way? alleluia.}

\antiphon{Mag.}{I tell you the truth; {\dag} it is expedient for you that I go away: for if I go not away, the Comforter will not come unto you, alleluia.}

\subsubsec{Tuesday}
\antiphon{Ben.}{Peace be unto you, it is I, {\dag} alleluia: be not afraid, alleluia.}

\antiphon{Mag.}{When the Comforter, {\dag} the Spirit of truth, is come, he will convince the world of sin, and of righteousness, and of judgement, alleluia.}

\subsubsec{Wednesday}
\antiphon{Ben.}{A spirit {\dag} hath not flesh and bones, as ye see me have: believe ye therefore, alleluia.}

\antiphon{Mag.}{I have yet {\dag} many things to say unto you, but ye cannot bear them now: howbeit, when he, the Spirit of truth, is come, he will guide you into all truth, alleluia.}

\subsubsec{Thursday}
\antiphon{Ben.}{The disciples {\dag} offered the Lord a piece of broiled fish, and of an honeycomb, alleluia, alleluia.}

\antiphon{Mag.}{For he shall not speak of himself; {\dag} but whatsoever he shall hear, that shall he speak: and he will shew you things to come, alleluia.}

\subsubsec{Friday}
\antiphon{Ben.}{These are the words {\dag} which I spake unto you, while I was yet with you, alleluia, alleluia.}

\end{multicols}


\subby{Rogation Sunday}
\feastday{Easter V}
%\begin{inhead}
%    {Semidouble}
%\end{inhead}

\begin{rubric}
	In the Daily Office, the Hymns are of Easter Day (p. \pageref{EasterDay}), with the following.
\end{rubric}

℣. Abide with us, Lord, alleluia. 

℟. For it is toward evening, alleluia.

\properantiphon{Mag.}{Hitherto {\dag} have ye asked nothing in my Name ask, and ye shall receive, alleluia.}\\

℣. In thy resurrection, O Christ, alleluia.

℟. Let heaven and earth rejoice, alleluia.

\properantiphon{Ben.}{Hitherto {\dag} have ye asked nothing in my Name ask, and ye shall receive, alleluia.}\\

℣. Abide with us, Lord, alleluia. 

℟. For it is toward evening, alleluia.

\properantiphon{Mag.}{Ask, {\dag} and ye shall receive, that your joy may be full: for the Father himself loveth you, because ye have loved me, and have believed, alleluia.}


\subby{Rogation Monday}
\feastday{Rogation Monday}
%\begin{inhead}
%    {Greater Feria}
%\end{inhead}

\begin{rubric}
	In the Daily Office, the Hymns are of Easter Day (p. \pageref{EasterDay}), with the following.
\end{rubric}

℣. In thy resurrection, O Christ, alleluia.

℟. Let heaven and earth rejoice, alleluia.

\properantiphon{Ben.}{Ask, {\dag} and ye shall receive; seek, and ye shall find; knock, and it shall be opened unto you, alleluia.}\\

℣. Abide with us, Lord, alleluia. 

℟. For it is toward evening, alleluia.

\properantiphon{Mag.}{For the Father {\dag} himself loveth you, because you have loved me, and have believed, alleluia.}


\subby{Rogation Tuesday}
\feastday{Rogation Tuesday}
%\begin{inhead}
%    {Feria}
%\end{inhead}

\begin{rubric}
	The Hymns are of Easter Day (p. \pageref{EasterDay}), with the following.
\end{rubric}

℣. In thy resurrection, O Christ, alleluia.

℟. Let heaven and earth rejoice, alleluia.

\properantiphon{Ben.}{It behoved Christ to suffer, {\dag} and to rise again from the dead, alleluia.}\\

℣. Abide with us, Lord, alleluia. 

℟. For it is toward evening, alleluia.

\properantiphon{Mag.}{I came forth from the Father, {\dag} and am come into the world: again I leave the world, and go to the Father, alleluia.}


\subby{Vigil of Ascension}
\feastday{Ascension Vigil}
%\begin{inhead}
%    {Vigil}
%\end{inhead}

\begin{rubric}
	The Hymns are of Easter Day (p. \pageref{EasterDay}), with the following.
\end{rubric}

℣. In thy resurrection, O Christ, alleluia.

℟. Let heaven and earth rejoice, alleluia.

\properantiphon{Ben.}{Father, the hour is come: {\dag} glorify thy Son with the glory which I had with thee before the world was, alleluia.}

\introit
\lett{W}{ith} a voice of singing declare ye this, and let it be heard, alleluia: utter it even unto the end of the earth: the Lord hath delivered his people, alleluia, alleluia. \textit{Ps.} O be joyful in God, all ye lands: sing praises unto his name, make his praise to be glorious.

\collect
\lett{O}{Lord,} from whom all good things do come, grant to us thy humble servants: that by thy holy inspiration we may think those things that be good; and by thy merciful guiding may perform the same. Through.

\lett{G}{rant,} we beseech thee, almighty God: that we, who in our affliction do put our trust in thy mercy; may ever be defended by thy protection against all adversities. (Through.)

\lett{G}{rant,} we beseech thee, O Lord God, that we thy servants may enjoy perpetual health of mind and of body: and, at the glorious intercession of blessed Mary ever Virgin, be delivered from present sadness, and rejoice in everlasting gladness. Through.

\readingcitation{Epistle}{Ephesians 4:7}
\lett{B}{rethren:} Unto every one of us is given grace according to the measure of the gift of Christ. Wherefore he saith, When he ascended up on high, he led captivity captive, and gave gifts unto men. (Now that he ascended, what is it but that he also descended first into the lower parts of the earth? He that descended is the same also that ascended up far above all heavens, that he might fill all things.) And he gave some, apostles; and some, prophets; and some, evangelists; and some, pastors and teachers; For the perfecting of the saints, for the work of the ministry, for the edifying of the body of Christ: Till we all come in the unity of the faith, and of the knowledge of the Son of God, unto a perfect man, unto the measure of the stature of the fulness of Christ.

\alleluia{Alleluia, alleluia. ℣. Christ is risen, and hath shewed light unto us, whom he hath redeemed with his blood. Alleluia. ℣. I came forth from the Father, and am come into the world: again I leave the world, and go to the Father. Alleluia.}

\readingcitation{Gospel}{John 17:1}
\lett{A}{t that time:} Jesus lifted up his eyes to heaven, and said: Father, the hour is come; glorify thy Son, that thy Son also may glorify thee: As thou hast given him power over all flesh, that he should give eternal life to as many as thou hast given him. And this is life eternal, that they might know thee the only true God, and Jesus Christ, whom thou hast sent. I have glorified thee on the earth: I have finished the work which thou gavest me to do. And now, O Father, glorify thou me with thine own self with the glory which I had with thee before the world was. I have manifested thy name unto the men which thou gavest me out of the world: thine they were, and thou gavest them me; and they have kept thy word. Now they have known that all things whatsoever thou hast given me are of thee. For I have given unto them the words which thou gavest me; and they have received them, and have known surely that I came out from thee, and they have believed that thou didst send me. I pray for them: I pray not for the world, but for them which thou hast given me; for they are thine. And all mine are thine, and thine are mine; and I am glorified in them. And now I am no more in the world, but these are in the world, and I come to thee.

\offertory{O praise the Lord our God, ye people, and make the voice of his praise to be heard: who holdeth our soul in life, and suffereth not out feet to slip: praised be the Lord, who hath not cast out my prayer, nor turned his mercy from me, alleluia.}

\secret
\lett{R}{eceive,} O Lord, the prayers of thy faithful people, together with the offering of these sacrifices: that through these observances of our bounden devotion, we may attain unto heavenly glory. Through.

\lett{W}{e} beseech thee, O Lord, that these our oblations may both loose the bonds of our iniquity, and obtain for us the gifts of thy loving-kindness. (Through.)

\lett{T}{hrough} thy mercy, O Lord, and the intercession of blessed Mary ever Virgin, may this oblation avail for our prosperity and peace both now and ever. Through.

\communion{O sing unto the Lord, alleluia: sing unto the Lord, and praise his name: be telling of his salvation from day to day, alleluia, alleluia.}

\postcommunion
\lett{G}{rant} unto us, O Lord, that we who have been fulfilled with the strength of thy heavenly table: may both desire those things which be right and obtain those things which we desire. Through.

\lett{W}{e} beseech thee, O Lord, to further with thy gracious favour these our supplications: that we, receiving thy gifts in our tribulation, may increase in thy love by the consolation of the same. (Through.)

\lett{G}{rant,} we beseech thee, O Lord: that we who have received these aids to our salvation may at all times and in all places be protected by the advocacy of blessed Mary ever Virgin; in whose honour we have made these offerings to thy majesty. Through.


\subby{Ascension Thursday}\label{Ascension}
\feastday{Ascension Thursday}
%\begin{inhead}
%    {First Class Double, Third Class Octave}
%\end{inhead}

\subbysub{I Evensong}

\begin{rubric}
	The Office Hymn is of Mattins, with the following Versicle \& Antiphon.
\end{rubric}

℣. God is gone up with a merry noise, alleluia.

℟. And the Lord with the sound of the trump, alleluia.

\properantiphon{Mag.}{Father, {\dag} I have manifested thy Name unto the men which thou gavest me: and now I pray for them, not for the world, because I come to thee, alleluia.}

\subbysub{Mattins}

\invitatoryhymn\label{AscensionInvitatory}

\gregorioscore{resources/gabc/ProperSeason/AscensionInvitatory.gabc}

\officehymn\label{AscensionMattins}

\gregorioscore{resources/gabc/ProperSeason/AscensionMattins.gabc}

℣. The Lord hath prepared, alleluia.

℟. His seat in heaven, alleluia.

\properantiphon{Ben.}{I ascend {\dag} unto my Father, and your Father: and to my God, and your God, alleluia.}

\subbysub{II Evensong}

\begin{rubric}
	The Office Hymn is of Mattins, with the following.
\end{rubric}

℣. God is gone up with a merry noise, alleluia.

℟. And the Lord with the sound of the trump, alleluia.

\properantiphon{Mag.}{O King of glory, {\dag} thou Lord of Sabaoth, who triumphing to-day hast ascended above all heavens, leave us not comfortless; but send on us the promise of the Father, even the Spirit of truth, alleluia.}


\begin{rubric}
	During the Octave of the Ascension, the Office is said daily or Commemoration made. All is said as on the Feast.
\end{rubric}


\subsec{Sunday within the Ascension Octave}
\feastday{Ascension Sunday}
%\begin{inhead}
%    {Semidouble}
%\end{inhead}

\begin{rubric}
	The Hymns are of Ascension Thursday (p. \pageref{Ascension}), with the following Versicles \& Antiphons.
\end{rubric}

℣. The Lord hath prepared, alleluia.

℟. His seat in heaven, alleluia.

\properantiphon{Mag.}{When the Comforter is come, {\dag} whom I will send unto you from the Father, even the Spirit of truth, which proceedeth from the Father, he shall testify of me, alleluia.}\\

℣. When Christ ascended up on high, alleluia.

℟. He led captivity captive, alleluia.

\properantiphon{Ben.}{When the Comforter is come, {\dag} whom I will send unto you from the Father, even the Spirit of truth, which proceedeth from the Father, he shall testify of me, alleluia.}\\

℣. The Lord hath prepared, alleluia.

℟. His seat in heaven, alleluia.

\properantiphon{Mag.}{These things have I told you, {\dag} that when the time shall come, ye may remember that I told you of them, alleluia.}


\subby{Vigil of Whitsunday}
\feastday{Whitsunday Vigil}
%\begin{inhead}
%    {First Class Vigil}
%\end{inhead}

\begin{rubric}
	From this day through Trinity Sunday, if a I Double or II Double occur, it is transferred to a day after Trinity Sunday. Other Feasts and Memorials are commemorated, except during the Triduum of Whitsunday. No notice is taken of other Octaves during this time.
\end{rubric}

\begin{rubric}
	At Mattins and the Hours, all is said as on the preceding Sunday, without Commemoration of the Ascension.
\end{rubric}

\begin{rubric}
	Mattins (or None) having been said in Choir, the Priest and Ministers, clad in vestments of violet colour, go up to the Altar and make a reverence, and the Priest kisses it in the middle. Then the Prophecies are read without title, the candles of the Altar remaining unlighted until the beginning of the Mass, as on Holy Saturday. The Priest reads them in a low voice at the Epistle corner of the Altar. At the end of the Prophecies the Collects are said without \emph{Let us bow the knee.}
\end{rubric}

\readingcitation{Prophecy I}{Genesis 22:1}
\lett{I}{n those days:} God did tempt Abraham, and said unto him, Abraham: and he said, Behold, here I am. And he said, Take now thy son, thine only son Isaac, whom thou lovest, and get thee into the land of Moriah; and offer him there for a burnt offering upon one of the mountains which I will tell thee of. And Abraham rose up early in the morning, and saddled his ass, and took two of his young men with him, and Isaac his son, and clave the wood for the burnt offering, and rose up, and went unto the place of which God had told him. Then on the third day Abraham lifted up his eyes, and saw the place afar off. And Abraham said unto his young men, Abide ye here with the ass; and I and the lad will go yonder and worship, and come again to you. And Abraham took the wood of the burnt offering, and laid it upon Isaac his son; and he took the fire in his hand, and a knife; and they went both of them together. And Isaac spake unto Abraham his father, and said, My father: and he said, Here am I, my son. And he said, Behold the fire and the wood: but where is the lamb for a burnt offering? And Abraham said, My son, God will provide himself a lamb for a burnt offering: so they went both of them together. And they came to the place which God had told him of; and Abraham built an altar there, and laid the wood in order, and bound Isaac his son, and laid him on the altar upon the wood. And Abraham stretched forth his hand, and took the knife to slay his son. And the angel of the \divineName{Lord} called unto him out of heaven, and said, Abraham, Abraham: and he said, Here am I. And he said, Lay not thine hand upon the lad, neither do thou any thing unto him: for now I know that thou fearest God, seeing thou hast not withheld thy son, thine only son from me. And Abraham lifted up his eyes, and looked, and behold behind him a ram caught in a thicket by his horns: and Abraham went and took the ram, and offered him up for a burnt offering in the stead of his son. And Abraham called the name of that place Jehovahjireh: as it is said to this day, In the mount of the \divineName{Lord} it shall be seen. And the angel of the \divineName{Lord} called unto Abraham out of heaven the second time, And said, By myself have I sworn, saith the \divineName{Lord}, for because thou hast done this thing, and hast not withheld thy son, thine only son: That in blessing I will bless thee, and in multiplying I will multiply thy seed as the stars of the heaven, and as the sand which is upon the sea shore; and thy seed shall possess the gate of his enemies; And in thy seed shall all the nations of the earth be blessed; because thou hast obeyed my voice. So Abraham returned unto his young men, and they rose up and went together to Beersheba; and Abraham dwelt at Beersheba.

\begin{rubric}
	After this and the other Prophecies the Response \emph{Thanks be to God} is not made: then the Priest says:
\end{rubric}
\letuspray
\lett{O}{God,} who in the deed of Abraham thy servant hast given a pattern of obedience to mankind: grant us so to conquer the perversity of our desires, that we may in all things fulfil the righteousness of thy commandments. Through.

\readingcitation{Prophecy II}{Exodus 14:24}
\lett{I}{n those days:} It came to pass, that in the morning watch the \divineName{Lord} looked unto the host of the Egyptians through the pillar of fire and of the cloud, and troubled the host of the Egyptians, And took off their chariot wheels, that they drave them heavily: so that the Egyptians said, Let us flee from the face of Israel; for the \divineName{Lord} fighteth for them against the Egyptians. And the \divineName{Lord} said unto Moses, Stretch out thine hand over the sea, that the waters may come again upon the Egyptians, upon their chariots, and upon their horsemen. And Moses stretched forth his hand over the sea, and the sea returned to his strength when the morning appeared; and the Egyptians fled against it; and the \divineName{Lord} overthrew the Egyptians in the midst of the sea. And the waters returned, and covered the chariots, and the horsemen, and all the host of Pharaoh that came into the sea after them; there remained not so much as one of them. But the children of Israel walked upon dry land in the midst of the sea; and the waters were a wall unto them on their right hand, and on their left. Thus the \divineName{Lord} saved Israel that day out of the hand of the Egyptians; and Israel saw the Egyptians dead upon the sea shore. And Israel saw that great work which the \divineName{Lord} did upon the Egyptians: and the people feared the \divineName{Lord}, and believed the \divineName{Lord}, and his servant Moses. Then sang Moses and the children of Israel this song unto the \divineName{Lord}, and spake, saying:

\tract{I will sing unto the \divineName{Lord}, for he hath triumphed gloriously: the horse and his rider hath he thrown into the sea. The \divineName{Lord} is my strength and song, and he is become my salvation. ℣. He is my God, and I will prepare him an habitation; my father's God, and I will exalt him. ℣. The \divineName{Lord} is a man of war: the \divineName{Lord} is his name.}

\letuspray
\lett{O}{God,} who by the light of thy new Covenant hast made manifest thy wonders wrought in former times, shewing in the Red Sea a pattern of the sacred font, and in the deliverance of thy people from bondage in Egypt foreshadowing the sacraments of thy Christian people: grant that all nations, being admitted by the merit of their faith to the privilege of Israel, may be regenerated by the partaking of thy Holy Spirit. Through . . . in the unity of the same Holy Ghost.

\readingcitation{Prophecy III}{Deuteronomy 31:22}
\lett{I}{n those days:} Moses wrote this song, and taught it the children of Israel. And he gave Joshua the son of Nun a charge, and said, Be strong and of a good courage: for thou shalt bring the children of Israel into the land which I sware unto them: and I will be with thee. And it came to pass, when Moses had made an end of writing the words of this law in a book, until they were finished, That Moses commanded the Levites, which bare the ark of the covenant of the \divineName{Lord}, saying, Take this book of the law, and put it in the side of the ark of the covenant of the \divineName{Lord} your God, that it may be there for a witness against thee. For I know thy rebellion, and thy stiff neck: behold, while I am yet alive with you this day, ye have been rebellious against the \divineName{Lord}; and how much more after my death? Gather unto me all the elders of your tribes, and your officers, that I may speak these words in their ears, and call heaven and earth to record against them. For I know that after my death ye will utterly corrupt yourselves, and turn aside from the way which I have commanded you; and evil will befall you in the latter days; because ye will do evil in the sight of the \divineName{Lord}, to provoke him to anger through the work of your hands. And Moses spake in the ears of all the congregation of Israel the words of this song, until they were ended.

\tract{Give ear, O ye heavens, and I will speak; and hear, O earth, the words of my mouth. ℣. My doctrine shall drop as the rain, my speech shall distil as the dew. ℣. As the small rain upon the tender herb, and as the showers upon the grass: Because I will publish the name of the \divineName{Lord}. ℣. Ascribe ye greatness unto our God. He is the Rock, his work is perfect: for all his ways are judgment. ℣. A God of truth and without iniquity, just and right is he.}

\letuspray
\lett{O}{God,} the glory of the faithful and the life of the just, who through Moses thy servant hast instructed us also in the chanting of thy sacred song: accomplish in all nations the work of thy mercy, granting them felicity, and delivering them from terror; that those things which were uttered for punishment may be turned into an everlasting remedy. Through.

\readingcitation{Prophecy IV}{Isaiah 4:1}
\lett{I}{n that day} seven women shall take hold of one man, saying, We will eat our own bread, and wear our own apparel: only let us be called by thy name, to take away our reproach. In that day shall the branch of the \divineName{Lord} be beautiful and glorious, and the fruit of the earth shall be excellent and comely for them that are escaped of Israel. And it shall come to pass, that he that is left in Zion, and he that remaineth in Jerusalem, shall be called holy, even every one that is written among the living in Jerusalem: When the Lord shall have washed away the filth of the daughters of Zion, and shall have purged the blood of Jerusalem from the midst thereof by the spirit of judgment, and by the spirit of burning. And the \divineName{Lord} will create upon every dwelling place of mount Zion, and upon her assemblies, a cloud and smoke by day, and the shining of a flaming fire by night: for upon all the glory shall be a defence. And there shall be a tabernacle for a shadow in the daytime from the heat, and for a place of refuge, and for a covert from storm and from rain.

\tract{My wellbeloved hath a vineyard in a very fruitful hill. ℣. And he fenced it, and gathered out the stones thereof, and planted it with the choicest vine, and built a tower in the midst of it. ℣. And also made a winepress therein: for the vineyard of the \divineName{Lord} of hosts is the house of Israel.}

\letuspray
\lett{A}{lmighty} and everlasting God, who through thy only Son hast revealed thyself to be the husbandman of thy Church, who dost mercifully purge every branch that bringeth forth fruit in the true vine, even the same thy Christ, to the intent that it may bring forth more fruit: let not the thorns of sin prevail against thy faithful people whom by the Font of baptism thou hast brought like a vine out of Egypt; that being fortified by the sanctifying power of thy Spirit, they may be enriched with everlasting fruit. Through the same . . . in the unity of the same.

\readingcitation{Prophecy V}{Baruch 3:9}
\lett{H}{ear,} O Israel, the commandments of life: give ear to understand wisdom. How happeneth it, O Israel, that thou art in thine enemies' land, that thou art waxen old in a strange country, that thou art defiled with the dead, that thou art counted with them that go down into the grave? Thou hast forsaken the fountain of wisdom. For if thou hadst walked in the way of God, thou shouldest have dwelled in peace for ever. Learn where is wisdom, where is strength, where is understanding; that thou mayest know also where is length of days, and life, where is the light of the eyes, and peace. Who hath found out her place? and who hath come into her treasuries? Where are the princes of the heathen, and such as ruled the beasts that are upon the earth; they that had their pastime with the fowls of the air, and they that hoarded up silver and gold, wherein men trust; and of whose getting there is no end? For they that wrought in silver, and were so careful, and whose works are past finding out, they are vanished and gone down to the grave, and others are come up in their steads. Younger men have seen the light, and dwelt upon the earth: but the way of knowledge have they not known, neither understood they the paths thereof: neither have their children laid hold of it: they are far off from their way. It hath not been heard of in Canaan, neither hath it been seen in Teman. The sons also of Agar that seek understanding, which are in the land, the merchants of Merran and Teman, and the authors of fables, and the searchers out of understanding; none of these have known the way of wisdom, or remembered her paths. O Israel, how great is the house of God! and how large is the place of his possession! great, and hath none end; high, and unmeasurable. There were the giants born that were famous of old, great of stature, and expert in war. These did not God choose, neither gave he the way of knowledge unto them: so they perished, because they had no wisdom, they perished through their own foolishness. Who hath gone up into heaven, and taken her, and brought her down from the clouds? Who hath gone over the sea, and found her, and will bring her for choice gold? There is none that knoweth her way, nor any that comprehendeth her path. But he that knoweth all things knoweth her, he found her out with his understanding: he that prepared the earth for evermore hath filled it with fourfooted beasts: he that sendeth forth the light, and it goeth; he called it, and it obeyed him with fear: and the stars shined in their watches, and were glad: when he called them, they said, Here we be; they shined with gladness unto him that made them. This is our God, and there shall none other be accounted of in comparison of him. He hath found out all the way of knowledge, and hath given it unto Jacob his servant, and to Israel that is beloved of him. Afterward did she appear upon earth, and was conversant with men.

\letuspray
\lett{O}{God,} who by the mouths of the Prophets hast commanded us to leave things temporal, and to strive after things eternal: grant unto thy servants; that we, knowing the things which thou commandest, may by thy heavenly inspiration be enabled to perform the same. Through.

\readingcitation{Prophecy VI}{Ezekiel 37:}
\lett{I}{n those days:} The hand of the \divineName{Lord} was upon me, and carried me out in the spirit of the \divineName{Lord}, and set me down in the midst of the valley which was full of bones, And caused me to pass by them round about: and, behold, there were very many in the open valley; and, lo, they were very dry. And he said unto me, Son of man, can these bones live? And I answered, O Lord \divineName{God}, thou knowest. Again he said unto me, Prophesy upon these bones, and say unto them, O ye dry bones, hear the word of the \divineName{Lord}. Thus saith the Lord \divineName{God} unto these bones; Behold, I will cause breath to enter into you, and ye shall live: And I will lay sinews upon you, and will bring up flesh upon you, and cover you with skin, and put breath in you, and ye shall live; and ye shall know that I am the \divineName{Lord}. So I prophesied as I was commanded: and as I prophesied, there was a noise, and behold a shaking, and the bones came together, bone to his bone. And when I beheld, lo, the sinews and the flesh came up upon them, and the skin covered them above: but there was no breath in them. Then said he unto me, Prophesy unto the wind, prophesy, son of man, and say to the wind, Thus saith the Lord \divineName{God}; Come from the four winds, O breath, and breathe upon these slain, that they may live. So I prophesied as he commanded me, and the breath came into them, and they lived, and stood up upon their feet, an exceeding great army. Then he said unto me, Son of man, these bones are the whole house of Israel: behold, they say, Our bones are dried, and our hope is lost: we are cut off for our parts. Therefore prophesy and say unto them, Thus saith the Lord \divineName{God}; Behold, O my people, I will open your graves, and cause you to come up out of your graves, and bring you into the land of Israel. And ye shall know that I am the \divineName{Lord}, when I have opened your graves, O my people, and brought you up out of your graves, And shall put my spirit in you, and ye shall live, and I shall place you in your own land.

\letuspray
\lett{O}{Lord} God of hosts, who restorest those things that are broken down, and preservest those things that thou restorest: increase the peoples that shall be regenerated in the sanctification of thy name; that all who are washed in holy baptism may ever be guided by thy inspiration. Through.

    \par\noindent
	\centerline{\rule{0.25\textwidth}{0.4pt}}
	\par\noindent

\begin{rubric}
	These ended, the Celebrant receives a violet Cope, and, while he descends to the Font, the following is sung:
\end{rubric}

\tract{Like as the hart desireth the water-brooks: so longeth my soul after thee, O God. ℣. My soul is athirst for God, yea, even for the living God: when shall I come to appear before the presence of God? ℣. My tears have been my meat day and night, while they daily say unto me: Where is now thy God?}

\begin{rubric}
	Then the Priest, before he enters for the blessing of the Font, says near the Font:
\end{rubric}

\elcol{℣. The Lord be with you.}{℣. Dóminus vobíscum.}
\elcol{℟. And with thy spirit.}{℟. Et cum spíritu tuo.}
\elcol{℣. Let us pray.}{℣. Orémus.}

\needspace{4\baselineskip}
\lett{G}{rant,} we beseech thee, almighty God: that we, who observe the solemnity of the gift of the Holy Ghost, being inflamed with heavenly desires, may thirst after the fountain of life. Through ... in the unity of the same.

\begin{rubric}
	Then he proceeds to the blessing of the Font, as on Holy Saturday.
\end{rubric}

    \par\noindent
	\centerline{\rule{0.25\textwidth}{0.4pt}}
	\par\noindent

\begin{rubric}
	But when there are no Fonts, the sixth Prophecy with its Collect being finished, the Celebrant lays aside the Chasuble, and prostrates himself before the Altar with the Ministers; and all kneeling, the Litany is sung by two Cantors in the middle of the choir, both choirs responding together. When they come to \emph{℣. We sinners beseech thee,} the Priest and Ministers rise and go to the Sacristy, where they put on red vestments, while the lights around the Altar are lighted.
\end{rubric}

\begin{rubric}
	At the end of the Litany \emph{Kyrie, eléison} is sung solemnly for the Mass, and repeated as usual. Meanwhile the Priest with the Ministers proceeds to the Altar, and makes the Confession, Prayer of Humble Access, and Ten Commandments (and/or Summary of the Law); then ascending, he kisses it, and censes it as usual. \emph{Kyrie, eléison} being ended, he begins solemnly \emph{Glory be to God on high}, and the bells are rung.
\end{rubric}

\begin{rubric}
	Afterwards the Priest says:
\end{rubric}

\elcol{℣. The Lord be with you.}{℣. Dóminus vobíscum.}
\elcol{℟. And with thy spirit.}{℟. Et cum spíritu tuo.}
\elcol{℣. Let us pray.}{℣. Orémus.}

\collect
\lett{G}{rant,} we beseech thee, almighty God: that the splendour of thy brightness may shine forth upon us; and that the light of thy light may, by the illumination of the Holy Spirit, strengthen the hearts of them who through thy grace are born again. Through ... in the unity of the same.

\begin{rubric}
	This Collect only is said, even if a Commemoration has been made in the Office.
\end{rubric}

%MANUAL ADJUSTMENT:
\vspace{-1ex}

\readingcitation{Epistle}{Acts 19:1}
\lett{I}{n those days:} It came to pass, that, while Apollos was at Corinth, Paul having passed through the upper coasts came to Ephesus: and finding certain disciples, He said unto them, Have ye received the Holy Ghost since ye believed? And they said unto him, We have not so much as heard whether there be any Holy Ghost. And he said unto them, Unto what then were ye baptized? And they said, Unto John's baptism. Then said Paul, John verily baptized with the baptism of repentance, saying unto the people, that they should believe on him which should come after him, that is, on Christ Jesus. When they heard this, they were baptized in the name of the Lord Jesus. And when Paul had laid his hands upon them, the Holy Ghost came on them; and they spake with tongues, and prophesied. And all the men were about twelve. And he went into the synagogue, and spake boldly for the space of three months, disputing and persuading the things concerning the kingdom of God.

\alleluia{Alleluia. ℣. O give thanks unto the Lord, for he is gracious: and his mercy endureth for ever.}

\begin{rubric}
	\emph{Alleluia} is not repeated, but there immediately follows:
\end{rubric}

\tract{O praise the Lord, all ye heathen: praise him, all ye nations. ℣. For his merciful kindness is ever more and more towards us: and the truth of the Lord endureth for ever.}

\begin{rubric}
	At the Gospel, lights are not carried, but incense only.
\end{rubric}

\readingcitation{Gospel}{John 14:15}
\lett{A}{t that time:} Jesus said unto his disciples: If ye love me, keep my commandments. And I will pray the Father, and he shall give you another Comforter, that he may abide with you for ever; Even the Spirit of truth; whom the world cannot receive, because it seeth him not, neither knoweth him: but ye know him; for he dwelleth with you, and shall be in you. I will not leave you comfortless: I will come to you. Yet a little while, and the world seeth me no more; but ye see me: because I live, ye shall live also. At that day ye shall know that I am in my Father, and ye in me, and I in you. He that hath my commandments, and keepeth them, he it is that loveth me: and he that loveth me shall be loved of my Father, and I will love him, and will manifest myself to him.

\offertory{O send forth thy Spirit, and they shall be made, and thou shalt renew the face of the earth: the glorious majesty of the Lord shall endure for ever, alleluia.}

\secret
\lett{S}{anctify,} we beseech thee, O Lord, the gifts which we offer: and cleanse our hearts by the enlightening of the Holy Spirit. Through ... in the unity of the same Holy Ghost.

\communion{In the last day of the feast, Jesus said: He that believeth on me, out of his belly shall flow rivers of living water: but this spake he of the Spirit, which they that believe on him should receive, alleluia, alleluia.}

\postcommunion
\lett{P}{our} thy Holy Spirit upon us, O Lord, and cleanse our hearts: that they may be made fruitful by the inward sprinkling of his dew. Through... in the unity of the same Holy Spirit.


\subby{Whitsunday}\label{Whitsunday}
\feastday{Whitsunday}
%\begin{inhead}
%    {First Class Double, First Class Octave}
%\end{inhead}

\subbysub{I Evensong}\label{WhitsundayEvensong}

\gregorioscore{resources/gabc/ProperSeason/WhitsundayEvensong.gabc}

℣. They were all filled with the Holy Ghost, alleluia.

℟. And began to speak, alleluia.

\properantiphon{Mag.}{I will not leave you {\dag} comfortless, alleluia: I go; but I will come to you, alleluia: and your heart shall rejoice, alleluia.}

\subbysub{Mattins}

\invitatoryhymn\label{WhitsundayInvitatory}

\gregorioscore{resources/gabc/ProperSeason/WhitsundayInvitatory.gabc}

\officehymn\label{WhitsundayMattins}

\gregorioscore{resources/gabc/ProperSeason/WhitsundayMattins.gabc}

℣. They were all filled with the Holy Ghost, alleluia.

℟. And began to speak, alleluia.

\properantiphon{Ben.}{Receive ye {\dag} the Holy Ghost: whosesoever sins ye remit, they are remitted unto them, alleluia.}

\subbysub{II Evensong}

\begin{rubric}
	The Office Hymn is of I Evensong, with the following Versicle \& Antiphon.
\end{rubric}

℣. The Apostles did speak with other tongues, alleluia.

℟. The wonderful works of God, alleluia.

\properantiphon{Mag.}{To-day {\dag} are fulfilled the days of Pentecost, alleluia: to-day the Holy Spirit appeared in fire to the disciples, and bestowed upon them his manifold graces: sending them into all the world, to preach the gospel, and to testify: He that believeth and is baptised shall be saved, alleluia.}


\subby{Whit-Monday}
\feastday{Whit-Monday}
%\begin{inhead}
%    {First Class Double}
%\end{inhead}

\begin{rubric}
	In the Daily Office, the Hymns \& Versicles are of Whitsunday (p. \pageref{Whitsunday}), with the following Antiphons.
\end{rubric}

\properantiphon{Ben.}{God so loved the world, {\dag} that he gave his only-begotten Son, that whosoever believeth in him should not perish, but have everlasting life, alleluia.}

\properantiphon{Mag.}{If a man love me, {\dag} he will keep my saying, and my Father will love him: and we will come unto him, and make our abode with him, alleluia.}


\subby{Whit-Tuesday}
\feastday{Whit-Tuesday}
%\begin{inhead}
%    {First Class Double}
%\end{inhead}

\begin{rubric}
	In the Daily Office, the Hymns \& Versicles are of Whitsunday (p. \pageref{Whitsunday}), with the following Antiphons.
\end{rubric}

\properantiphon{Ben.}{I am the door, {\dag} saith the Lord: by me if any man enter in, he shall be saved, and shall find pasture, alleluia.}

\properantiphon{Mag.}{Peace {\dag} I leave with you, my peace I give unto you: not as the world giveth, give I unto you, alleluia.}


\subby{Ember Wednesday in Whitsuntide}
\feastday{Whitsun Emberday}
\fancyhead[RE,LO]{Ember Wednesday}
%\begin{inhead}
%    {Semidouble}
%\end{inhead}

\begin{rubric}
	In the Daily Office, the Hymns \& Versicles are of Whitsunday (p. \pageref{Whitsunday}), with the following Antiphons.
\end{rubric}

\properantiphon{Ben.}{I am the living bread, {\dag} saith the Lord, which came down from heaven, alleluia, alleluia.}

\properantiphon{Mag.}{I am the living bread, {\dag} which came down from heaven: if any man eat of this bread, he shall live for ever: and the bread that I will give is my flesh, which I will give for the life of the world, alleluia.}


\subby{Whit-Thursday}
\feastday{Whit-Thursday}
\fancyhead[RE,LO]{}
%\begin{inhead}
%    {Semidouble}
%\end{inhead}


\begin{rubric}
	In the Daily Office, the Hymns \& Versicles are of Whitsunday (p. \pageref{Whitsunday}), with the following Antiphons.
\end{rubric}

\properantiphon{Ben.}{Jesus called unto him {\dag} his twelve disciples, and gave them power and authority over all devils, and to cure diseases: and he sent them to preach the kingdom of God, and to heal the sick, alleluia.}

\properantiphon{Mag.}{The Spirit {\dag} which proceedeth from the Father, alleluia: he shall glorify me, alleluia, alleluia.}


\subby{Ember Friday in Whitsuntide}
\feastday{Whitsun Emberday}
\fancyhead[RE,LO]{Ember Friday}
%\begin{inhead}
%    {Semidouble}
%\end{inhead}

\begin{rubric}
	In the Daily Office, the Hymns \& Versicles are of Whitsunday (p. \pageref{Whitsunday}), with the following Antiphons.
\end{rubric}

\properantiphon{Ben.}{Jesus said: {\dag} But that ye may know that the Son of man hath power on earth to forgive sins, (he saith unto the sick of the palsy) I say unto thee, Arise, take up thy bed, and go thy way into thine house, alleuia.}

\properantiphon{Mag.}{But the Comforter, {\dag} which is the Holy Ghost, whom the Father will send in my Name, he will teach you all things, and bring all things to your remembrance, whatsoever I have said unto you, alleluia.}


\subby{Ember Saturday in Whitsuntide}
\feastday{Whitsun Emberday}
\fancyhead[RE,LO]{Ember Saturday}
%\begin{inhead}
%    {Semidouble}
%\end{inhead}

\begin{rubric}
	In the Daily Office, the Hymns \& Versicles are of Whitsunday (p. \pageref{Whitsunday}), with the following Antiphons.
\end{rubric}

\properantiphon{Ben.}{The love of God is shed {\dag} abroad in our hearts by his Spirit which dwelleth in us, alleluia.}

%MANUAL ADJUSTMENT:
\clearpage
\subby{Trinity Sunday}
\feastday{Trinity Sunday}
\fancyhead[RE,LO]{}
%\begin{inhead}
%    {First Class Double}
%\end{inhead}

\subbysub{I Evensong}

\begin{rubric}
	The Office Hymn is of Saturday in Winter (p. \pageref{SaturdayEvensongWinter}), with the following.
\end{rubric}

℣. Blessed art thou, O Lord, in the firmament of heaven.

℟. And to be praised, and glorified, and magnified for ever.

\properantiphon{Mag.}{Thanks, O God, {\dag} be unto thee, thanks be unto thee, one and very Trinity, one and supreme Deity, holy and onely Unity.}

\subbysub{Mattins}

\invitatoryhymn\label{TrinityInvitatory}

\gregorioscore{resources/gabc/ProperSeason/TrinityInvitatory.gabc}

\officehymn\label{TrinityMattins}

\gregorioscore{resources/gabc/ProperSeason/TrinityMattins.gabc}

℣. Let us bless the Father, and the Son, and the Holy Ghost.

℟. Praise him, and magnify him for ever.

\properantiphon{Ben.}{Blessed be {\dag} the holy Creator and Governor of all things, the holy and undivided Trinity, both now and ever, and to endless ages of ages.}

\subbysub{II Evensong}

\begin{rubric}
	The Office Hymn is of Saturday in Winter (p. \pageref{SaturdayEvensongWinter}), with the following.
\end{rubric}

℣. Blessed art thou, O Lord, in the firmament of heaven.

℟. And to be praised, and glorified, and magnified for ever.

\properantiphon{Mag.}{Thee, O God, {\dag} the Father unbegotten, thee, O only-begotten Son; thee, O Holy Spirit, the Paraclete; holy and undivided Trinity: with our whole heart and mouth we confess thee, we praise thee and bless thee: to thee be glory for ever and ever.}


\subby{The Most Holy Body of Christ}\label{CorpusChristi}
\feastday{Corpus Christi}
%\begin{inhead}
%	{Thursday after Trinity Sunday}\par
%    {First Class Double, Second Class Octave}
%\end{inhead}

\subbysub{I Evensong}\label{CorpusChristiEvensong}

\gregorioscore{resources/gabc/ProperSeason/CorpusChristiEvensong1.gabc}

\begin{rubric}
	If the Office is recited in the presence of the exposed Sacrament, the following stanza is said kneeling:
\end{rubric}

\gregorioscore{resources/gabc/ProperSeason/CorpusChristiEvensong2.gabc}

℣. Thou gavest them bread from heaven, alleluia.

℟. Containing within itself all sweetness, alleluia.

\properantiphon{Mag.}{O how sweet {\dag} is thy Spirit, O Lord, who, that thou mightest shew thy kindness unto thy children, givest them that sweetest bread from heaven, fillest the hungry with good things, and sendest the rich and scornful empty away.}

\subbysub{Mattins}

\invitatoryhymn\label{CorpusChristiInvitatory}

\gregorioscore{resources/gabc/ProperSeason/CorpusChristiInvitatory.gabc}

\officehymn\label{CorpusChristiMattins}

\gregorioscore{resources/gabc/ProperSeason/CorpusChristiMattins.gabc}

℣. He maketh peace in thy borders, alleluia.

℟. And filleth thee with the flour of wheat, alleluia.

\properantiphon{Ben.}{I am {\dag} the living bread, which came down from heaven: if any man eat of this bread, he shall live for ever, alleluia.}

\subbysub{II Evensong}

\begin{rubric}
	The Office Hymn \& Versicle are of I Evensong, with the following Antiphon.
\end{rubric}

\properantiphon{Mag.}{O sacred banquet, {\dag} wherein Christ is received, the memory of his Passion is renewed; the soul with grace is filled, and a pledge of future glory is bestowed, alleluia.}


\subby{Sunday in the Octave of Corpus Christi}
\feastday{Trinity I}
\centerline{\small{(First Sunday after Trinity)}}
%\begin{inhead}
%    {Semidouble}
%\end{inhead}

\begin{rubric}
	The Hymns are of the Feast of the Holy Body of Christ (p. \pageref{CorpusChristi}), with the following Versicles \& Antiphons.
\end{rubric}

℣. He fed them with the finest wheat flour, alleluia.

℟. And with honey out of the stony rock did he satisfy them, alleluia.

\properantiphon{Mag.}{The child Samuel {\dag} ministered unto the Lord before Eli, and the word of the Lord was precious in his sight.}\\

%MANUAL ADJUSTMENT:
\clearpage
℣. He gave them bread from heaven, alleluia.

℟. So man did eat angels' food, alleluia.

\properantiphon{Ben.}{Father Abraham, {\dag} have mercy on me, and send Lazarus, that he may dip the tip of his finger in water, and cool my tongue.}\\

℣. He fed them with the finest wheat flour, alleluia.

℟. And with honey out of the stony rock did he satisfy them, alleluia.

\properantiphon{Mag.}{Son, remember {\dag} that thou in thy lifetime receivedst thy good things, and likewise Lazarus evil things.}


\subby{Compassion of Our Lord Jesus Christ}
\feastday{Divine Compassion}
%\begin{inhead}
%	{Friday after Trinity I}\par
%    {Second Class Double, Simple Octave}
%\end{inhead}

\subbysub{I Evensong}\label{CompassionEvensong}

\gregorioscore{resources/gabc/ProperSeason/CompassionEvensong.gabc}

℣. I am come to send fire on the earth.

℟. And what will I, if it be already kindled?

\properantiphon{Mag.}{Thy rebuke {\dag} hath broken my heart; I am full of heaviness: I looked for some to have pity on me, but there was no man; neither found I any to comfort me.}

%MANUAL ADJUSTMENT:
\clearpage
\subbysub{Mattins}

\invitatoryhymn\label{CompassionInvitatory}

\gregorioscore{resources/gabc/ProperSeason/CompassionInvitatory.gabc}

\officehymn\label{CompassionMattins}

\gregorioscore{resources/gabc/ProperSeason/CompassionMattins.gabc}

℣. Surely he hath borne our griefs.

℟. And carried our sorrows.

\properantiphon{Ben.}{He was wounded {\dag} for our transgressions he was bruised for our iniquities: the chastisement of our peace was upon him, and with his stripes we are healed.}

\subbysub{II Evensong}

\begin{rubric}
	The Office Hymn is as in I Evensong, with the following Versicle \& Antiphon.
\end{rubric}

℣. With joy shall ye draw water.

℟. Out of the wells of salvation.

\properantiphon{Mag.}{But when they came to Jesus, {\dag} and saw that he was dead already, they brake not his legs: but one of the soldiers with a spear pierced his side, and forthwith came there out blood and water.}


\subby{Sundays after Trinity}
\feastday{Trinitytide}

\begin{multicols}{2}
\subsubsec{Second Sunday after Trinity}
\properantiphon{Mag.}{And all Israel knew {\dag} from Dan even to Beersheba, that Samuel was established to be a prophet of the Lord.}

\properantiphon{Ben.}{A certain man {\dag} made a great supper, and bade many; and sent his servants at supper time to say to them that were bidden, Come; for all things are now ready, alleluia.}

\properantiphon{Mag.}{Go out quickly {\dag} into the streets and lanes of the city, and compel them to come in; the poor and the maimed, the halt and the blind, that my house may be filled, alleluia.}

\subsubsec{Third Sunday after Trinity}
\properantiphon{Mag.}{So David prevailed {\dag} over the Philistine with a sling and a stone, in the Name of the Lord.}

\properantiphon{Ben.}{What man of you, {\dag} having an hundred sheep, if he lose one of them, doth not leave the ninety and nine in the wilderness, and go after that which is lost, until he find it? Alleluia.}

\properantiphon{Mag.}{What woman, {\dag} having ten pieces of silver, if she lose one piece, doth not light a candle, and sweep the house, and seek diligently until she find it?}

\subsubsec{Fourth Sunday after Trinity}
\properantiphon{Mag.}{Ye mountains {\dag} of Gilboa, let there be neither dew nor rain upon you; for there the shield of the mighty is vilely cast away, the shield of Saul, as though he had not been anointed with oil. How are the mighty fallen in the midst of the battle! Jonathan was slain upon thy high places! Saul and Jonathan were lovely and exceeding pleasant in their lives, and in their death they were not divided.}

\properantiphon{Ben.}{Be ye therefore {\dag} merciful, as your Father also is merciful, saith the Lord.}

\properantiphon{Mag.}{Judge not, {\dag} that ye be not judged: for with what judgment ye judge, ye shall be judged, saith the Lord.}

\subsubsec{Fifth Sunday after Trinity}
\properantiphon{Mag.}{O Lord, I beseech thee, {\dag} do away the iniquity of thy servant, for I have done very foolishly.}

\properantiphon{Ben.}{And Jesus entered {\dag} into a ship, and sat down, and taught the people, alleluia.}

\properantiphon{Mag.}{Master, {\dag} we have toiled all the night, and have taken nothing: nevertheless, at thy word I will let down the net.}

\subsubsec{Sixth Sunday after Trinity}
\properantiphon{Mag.}{Zadok the priest {\dag} and Nathan the prophet anointed Solomon king in Gihon: and the people came up rejoicing and said, Let the king live for ever.}

\properantiphon{Ben.}{Ye have heard that it was said {\dag} by them of old time, Thou shalt not kill, and whosoever shall kill, shall be in danger of the judgment.}

\properantiphon{Mag.}{If thou bring {\dag} thy gift to the altar, and there rememberest that thy brother hath ought against thee; leave there thy gift before the altar and go thy way; first be reconciled to thy brother, and then come and offer thy gift, alleluia.}

\subsubsec{Seventh Sunday after Trinity}
\properantiphon{Mag.}{Thou hast heard, O Lord, {\dag} the supplication of thy servant, that I might build a temple to thy Name.}

\properantiphon{Ben.}{The multitude being very great, {\dag} and having nothing to eat, Jesus called his disciples unto him, and saith unto them, I have compassion on the multitude, because they have now been with me three days, and have nothing to eat, alleluia.}

\properantiphon{Mag.}{I have compassion {\dag} on the multitude, because they have now been with me three days, and have nothing to eat: and if I send them away fasting, they will faint by the way, alleluia.}

\subsubsec{Eighth Sunday after Trinity}
\properantiphon{Mag.}{When the Lord took up Elijah {\dag} by a whirlwind into heaven, Elisha cried, saying: My father, the chariot of Israel and the horses thereof.}

\properantiphon{Ben.}{Beware {\dag} of false prophets, which come to you in sheep's clothing, but inwardly they are ravening wolves. Ye shall know them by their fruits, alleluia.}

\properantiphon{Mag.}{A good tree {\dag} cannot bring forth evil fruit, neither can a corrupt tree bring forth good fruit: every tree that bringeth not forth good fruit is hewn down, and cast into the fire, alleluia.}

\subsubsec{Ninth Sunday after Trinity}
\properantiphon{Mag.}{Jehoash {\dag} did that which was right in the sight of the Lord all his days, wherein Jehoiada the priest instructed him.}

\properantiphon{Ben.}{The lord said {\dag} unto the steward, How is it that I hear this of thee? Give an account of thy stewardship, alleluia.}

\properantiphon{Mag.}{What shall I do? {\dag} for my lord taketh away from me the stewardship: I cannot dig; to beg I am ashamed: I am resolved what to do, that, when I am put out of the stewardship, they may receive me into their houses.}

\subsubsec{Tenth Sunday after Trinity}
\properantiphon{Mag.}{I beseech thee, O Lord, {\dag} remember now how I have walked before thee in truth and with a perfect heart, and have done that which is good in thy sight.}

\properantiphon{Ben.}{When the Lord was come near to Jerusalem, {\dag} he beheld the city, and wept over it, saying: If thou hadst known, even thou! For the days shall come upon thee, that thine enemies shall cast a trench about thee, that thine enemies shall cast a trench about thee, and compass thee round, and keep thee in on every side, and shall lay thee even with the ground; because thou knewest not the time of thy visitation, alleluia.}

\properantiphon{Mag.}{Is it not written, {\dag} Mine house shall be called an house of prayer for all people? but ye have made it a den of robbers. And he was daily with them, teaching int he temple.}

\subsubsec{Eleventh Sunday after Trinity}
\properantiphon{Ben.}{And the publican, {\dag} standing afar off, would not lift up so much as his eyes unto heaven, but smote upon his breast, saying, God be merciful to me a sinner.}

\properantiphon{Mag.}{This man went down {\dag} to his house justified rather than the other: for every one that exalteth himself shall be abased; and he that humbleth himself shall be exalted.}

\subsubsec{Twelfth Sunday after Trinity}
\properantiphon{Ben.}{When the Lord had passed {\dag} through the coasts of Tyre, he made the deaf to hear and the dumb to speak.}

\properantiphon{Mag.}{He hath done all things well: {\dag} he maketh both the deaf to hear, and the dumb to speak.}

\subsubsec{Thirteenth Sunday after Trinity}
\properantiphon{Ben.}{Master, {\dag} what shall I do to inherit eternal life? He said unto him, What is written in the law? how readest thou? Thou shalt love the Lord thy God with all thy heart, alleluia.}

\properantiphon{Mag.}{A certain man {\dag} went down from Jerusalem to Jericho, and fell among thieves, which stripped him of his raiment, and wounded him, and departed, leaving him half dead.}

\subsubsec{Fourteenth Sunday after Trinity}
\properantiphon{Ben.}{As Jesus passed through {\dag} a certain village, there met him ten men that were lepers, which stood afar off: and they lifted up their voices, and said, Jesus, Master, have mercy on us.}

\properantiphon{Mag.}{And one of them, {\dag} when he saw that he was healed, turned back, and with a loud voice glorified God, alleluia.}

\subsubsec{Fifteenth Sunday after Trinity}
\properantiphon{Ben.}{Be not therefore anxious, {\dag} saying, What shall we eat? or What shall we drink? for your heavenly Father knoweth that ye have need of all these things, alleluia.}

\properantiphon{Mag.}{Seek ye first {\dag} the kingdom of God, and his righteousness; and all these things shall be added unto you, alleluia.}

\subsubsec{Sixteenth Sunday after Trinity}
\properantiphon{Ben.}{Jesus went into a city {\dag} called Nain; and behold, there was a dead man carried out, the only son of his mother.}

\properantiphon{Mag.}{A great prophet {\dag} is risen up among us: and God hath visited his people.}

\subsubsec{Seventeenth Sunday after Trinity}
\properantiphon{Ben.}{And Jesus went {\dag} into the house of one of the the chief Pharisees to eat bread on the sabbath day, behold there was a certain man before him which had the dropsy; and he took him, and healed him, and let him go.}

\properantiphon{Mag.}{When thou art bidden to {\dag} a wedding, sit down in the lowest place; that he that bade thee may say unto thee, Friend, go up higher; then shalt thou have worship in the presence of them that sit at meat with thee, alleluia.}

\subsubsec{Eighteenth Sunday after Trinity}
\properantiphon{Ben.}{Master, {\dag} which is the great commandment in the law? Jesus said unto him, Thou shalt love the Lord thy God with all thy heart, alleluia.}

\properantiphon{Mag.}{What think ye of Christ? {\dag} whose son is he? They say unto him, The son of David. Jesus saith unto them, How then doth David in spirit call him Lord, saying, The Lord said unto my Lord, Sit thou on my right hand?}

\subsubsec{Nineteenth Sunday after Trinity}
\properantiphon{Ben.}{The Lord said {\dag} unto the sick of the palsy, Son, be of good cheer, thy sins be forgiven thee.}

\properantiphon{Mag.}{The sick of the palsy {\dag} therefore took up his bed whereon he lay, glorifying God: and all the people, when they saw it, give praise unto God.}

\subsubsec{Twentieth Sunday after Trinity}
\properantiphon{Ben.}{Tell them which are bidden, {\dag} Behold, I have prepared my dinner; come unto the marriage, alleluia.}

\properantiphon{Mag.}{And when the king came in {\dag} to see the guests, he saw there a man which had not on a wedding garment: and he said unto him, Friend, how camest thou in hither not having a wedding garment?}

\subsubsec{Twenty-First Sunday after Trinity}
\properantiphon{Ben.}{There was a certain nobleman, {\dag} whose son was sick at Capernaum: when he heard that Jesus was come out of Jud{\ae}a into Galilee, he besought him that he would heal his son.}

\properantiphon{Mag.}{So the father knew {\dag} that it was at the same hour in the which Jesus said, Thy son liveth: and himself believed, and his whole house.}

\subsubsec{Twenty-Second Sunday after Trinity}
\properantiphon{Ben.}{Then said the lord {\dag} unto the servant, Pay me that thou owest. The servant therefore fell down and worshipped him, saying, Lord, have patience with me, and I will pay thee all.}

\properantiphon{Mag.}{Thou wicked servant, {\dag} I forgave thee all that debt, because thou desiredst me: shouldest not thou also have had compassion on thy fellow-servant, even as I had pity on thee, alleluia.}

\subsubsec{Twenty-Third Sunday after Trinity}
\properantiphon{Ben.}{Master, {\dag} we know that thou art true, and teachest the way of God in truth, alleluia.}

\properantiphon{Mag.}{Render therefore {\dag} unto Caesar the things which are Caesar's: and unto God the things that are God's, alleluia.}

\subsubsec{Twenty-Fourth Sunday after Trinity}
\properantiphon{Ben.}{For she said within herself, {\dag} If I may but touch the hem of his garment, I shall be whole.}

\properantiphon{Mag.}{But Jesus turned him about, {\dag} and when he saw her, he said, Daughter, be of good comfort; thy faith hath made thee whole, alleluia.}

\subsubsec{Sunday Next before Advent}
\properantiphon{Ben.}{When Jesus then {\dag} lifted up his eyes, and saw a great company come unto him, he saith unto Philip, Whence shall we buy bread, that these may eat? And this he said to prove him: for he himself knew what he would do.}

\properantiphon{Mag.}{Then those men, {\dag} when they had seen the miracle that Jesus did, said among themselves, This is of a truth that prophet that should come into the world.}

\end{multicols}


%After Trinity VIII

\subby{Magnificat Antiphons for Sundays in Trinitytide}
\fancyhead[RO,LE]{Trinitytide Antiphons}
\fancyhead[RE,LO]{}
\begin{multicols}{2}
\begin{inhead}
Sunday within 29 July - 4 August
\end{inhead}\par\noindent
Wisdom hath builded her house, {\dag} she hath hewn out her seven pillars: she hath subdued the nations, and in her own might hath she trodden under the necks of the proud and lofty.

\begin{inhead}
Sunday within 5 - 11 August
\end{inhead}\par\noindent
I dwell {\dag} in high places, and my throne is in a cloudy pillar.
\begin{inhead}
Sunday within 12 - 18 August
\end{inhead}\par\noindent
All wisdom {\dag} cometh from the Lord, and is with him for ever; and is before the ages.

\begin{inhead}
Sunday within 19 - 25 August
\end{inhead}\par\noindent
Wisdom crieth {\dag} aloud in the broad places: Whosoever loveth wisdom, let him turn in hither, and he shall find her; and when he hath found her, happy is he if he hold her fast.
\begin{inhead}
Sunday within 26 - 28 August
\end{inhead}\par\noindent
My son {\dag} keep thy father's commandment, and forsake not the law of thy mother: bind them continually upon thine heart.
\begin{inhead}
Sunday within 29 August - 4 September
\end{inhead}\par\noindent
Now when Job had heard {\dag} the words of the messengers, he endured with patience, saying, Shall we receive good at the Lord's hand, and shall we not receive evil also? In all this Job sinned not with his lips, neither charged God foolishly.
\begin{inhead}
Sunday within 5 - 11 September
\end{inhead}\par\noindent
In all this {\dag} Job sinned not with his lips, neither charged God foolishly.
\begin{inhead}
Sunday within 12 - 18 September
\end{inhead}\par\noindent
Remember not, {\dag} Lord, our offences, nor the offences of our forefathers; neither take thou vengeance of our sins.
\begin{inhead}
Sunday within 19 - 25 September
\end{inhead}\par\noindent
Adonai, {\dag} Lord God Almighty, great and wonderful, who hast given salvation by the hand of a woman; hear, we beseech thee, the prayers of thy servants.
\begin{inhead}
Sunday within 26 - 27 September
\end{inhead}\par\noindent
O Lord, {\dag} the King Almighty, the whole world is in thy power, and there is no man that can gainsay thee.
\begin{inhead}
Sunday within 28 September - 4 October
\end{inhead}\par\noindent
The Lord open your hearts {\dag} in his law and commandments; and may the Lord our God send you peace.
\begin{inhead}
Sunday within 5 - 11 October
\end{inhead}\par\noindent
The sun shone {\dag} upon the shields of gold, and the mountains glistered therewith: and yet the forces of the heathen were discomfited.
\begin{inhead}
Sunday within 12 - 18 October
\end{inhead}\par\noindent
But Israel mourned Judas, {\dag} and made great lamentation for him, saying: How art thou fallen, valiant in battle, that didst deliver the people of the Lord.
\begin{inhead}
Sunday within 19 - 25 October
\end{inhead}\par\noindent
God be gracious unto you, {\dag} and hear your prayers, and be at one with you: and the Lord our God never forsake you in time of trouble.
\begin{inhead}
Sunday within 26 - 28 October
\end{inhead}\par\noindent
Thine is the power, O Lord, {\dag} and thine is the kingdom: thou art high above all nations: give peace in our time, O Lord our God.
\begin{inhead}
Sunday within 29 October - 4 November
\end{inhead}\par\noindent
I saw the Lord also, {\dag} sitting upon a throne, high and lifted up, and his train filled the temple: the whole earth was full of the majesty of his glory.

\begin{inhead}
Sunday on 5 November
\end{inhead}\par\noindent
Behold, O Lord, {\dag} how the city is become desolate that was full of precious treasure; she sitteth sorrowful that was great among the nations: none can comfort her but only thou, O our God.

\begin{inhead}
Sunday within 6 - 12 November
\end{inhead}\par\noindent
Encompass us, {\dag} O Lord, with thine impregnable wall: and with the arms of thy power defend us alway, O our God.
\begin{inhead}
Sunday within 13 - 19 November
\end{inhead}\par\noindent
Thou that upholdest {\dag} the throne of the heavens and beholdest the depths, O Lord, King of kings; that weighest the mountains and holdest the earth in the hollow of thine hand: give ear, O Lord, unto us, in the midst of our groanings.

\begin{inhead}
Sunday within 20 - 26 November
\end{inhead}\par\noindent
I have set watchmen {\dag} upon thy walls, O Jerusalem, which shall never hold their peace day nor night, praising the Name of the Lord.
\end{multicols}


\phantomsection
\addcontentsline{toc}{section}{Proper of Saints}
\fancyhead[C]{\LARGE Proper of Saints}


\subby{Vigil of St. Andrew}
\feastday{St. Andrew Vigil}
\fancyhead[RE,LO]{29 November}
\begin{inhead}
    {Vigil\\
29 November}
\end{inhead}

\begin{rubric}
	If to-day be Saturday, the anticipated Vigil of St. Andrew is kept, but the commemoration of St. Saturninus is omitted, in which case the \nth{2} is of St. Mary (p. \SPMaryAdvent) and the \nth{3} against the persecutors of the Church (p. \SPAgainst) or for the Chief Bishop (p. \SPChiefBishop).
\end{rubric}

\introit
\lett{T}{he} Lord, walking by the sea of Galilee, saw two brethren, Peter and Andrew, and he called them, saying: Follow me: I will make you fishers of men. \textit{Ps.} The heavens declare the glory of God: and the firmament sheweth his handy-work.

\collect
\lett{W}{e} beseech thee, almighty God: that the blessed Apostle Andrew, whose festival we prevent, may implore thy help for us; that we, being absolved from our offences, may likewise be delivered from all dangers. Through.

\begin{rubric}
	In Advent, \nth{2} Collect of the Feria, \nth{3} of St. Saturninus (p. \pageref{SaturninusCollect}).
\end{rubric}
\begin{rubric}
	Outside of Advent, \nth{2} Collect of St. Saturninus (p. \pageref{SaturninusCollect}), \nth{3} of St. Mary  in Eastertide (p. \SPMaryEaster).
\end{rubric}

\readingcitation{Epistle}{Ecclesiasticus 44:22}
%RV:
\lett{I}{n} Isaac did the Lord establish, for Abraham his father's sake, The blessing of all men, and the covenant: And he made it rest upon the head of Jacob; He acknowledged him in his blessings, And gave to him by inheritance, And divided his portions; Among twelve tribes did he part them. And he brought out of him a man of mercy, Which found favour in the sight of all flesh; A man beloved of God and men, even Moses, Whose memorial is blessed. He made him like to the glory of the saints, And magnified him in the fears of his enemies. By his words he caused the wonders to cease; He glorified him in the sight of kings; He gave him commandment for his people, And shewed him part of his glory. He sanctified him in his faithfulness and meekness; He chose him out of all flesh. He made him to hear his voice, And led him into the thick darkness, And gave him commandments face to face, Even the law of life and knowledge, That he might teach Jacob the covenant, And Israel his judgments. He exalted Aaron, a holy man like unto him, Even his brother, of the tribe of Levi. He established for him an everlasting covenant, And gave him the priesthood of the people; He beautified him with comely ornaments, And girded him about with a robe of glory.

%Traditional Text:
%\lett{T}{he} blessing of the Lord was upon the head of the righteous. Therefore did the Lord give him an heritage, and divided his portions among the twelve tribes: and he found favour in the sight of all flesh. And he magnified him so that his enemies stood in fear of him, and by his words he caused the wonders to cease. He made him glorious in the sight of kings, and gave him a commandment for his people, and shewed him his glory. He sanctified him in his faithfulness and meekness, and chose him out of all men. And he gave him commandments before his face, even the law of life and knowledge, and made him to be exalted. An everlasting covenant he made with him, and girded him about with the girdle of righteousness: and the Lord crowned him with the crown of glory.

\gradual{Right honourable are thy friends, O God: right well is their princedom established. ℣. If I tell them: they are more in number than the sand.}

\readingcitation{Gospel}{John 1:35}
\lett{A}{t that time:} John stood, and two of his disciples; And looking upon Jesus as he walked, he saith, Behold the Lamb of God! And the two disciples heard him speak, and they followed Jesus. Then Jesus turned, and saw them following, and saith unto them, What seek ye? They said unto him, Rabbi, (which is to say, being interpreted, Master,) where dwellest thou? He saith unto them, Come and see. They came and saw where he dwelt, and abode with him that day: for it was about the tenth hour. One of the two which heard John speak, and followed him, was Andrew, Simon Peter's brother. He first findeth his own brother Simon, and saith unto him, We have found the Messias, which is, being interpreted, the Christ. And he brought him to Jesus. And when Jesus beheld him, he said, Thou art Simon the son of Jona: thou shalt be called Cephas, which is by interpretation, A stone. The day following Jesus would go forth into Galilee, and findeth Philip, and saith unto him, Follow me. Now Philip was of Bethsaida, the city of Andrew and Peter. Philip findeth Nathanael, and saith unto him, We have found him, of whom Moses in the law, and the prophets, did write, Jesus of Nazareth, the son of Joseph. And Nathanael said unto him, Can there any good thing come out of Nazareth? Philip saith unto him, Come and see. Jesus saw Nathanael coming to him, and saith of him, Behold an Israelite indeed, in whom is no guile! Nathanael saith unto him, Whence knowest thou me? Jesus answered and said unto him, Before that Philip called thee, when thou wast under the fig tree, I saw thee. Nathanael answered and saith unto him, Rabbi, thou art the Son of God; thou art the King of Israel. Jesus answered and said unto him, Because I said unto thee, I saw thee under the fig tree, believest thou? thou shalt see greater things than these. And he saith unto him, Verily, verily, I say unto you, Hereafter ye shall see heaven open, and the angels of God ascending and descending upon the Son of man.

\offertory{Thou hast crowned him with glory and worship: thou hast made him to have dominion of the works of thy hands, O Lord.}

\secret
\lett{W}{e} offer, O Lord, this gift to be hallowed unto thee: whereby, recalling the festival of thy blessed Apostle Andrew, we likewise implore the purification of our souls. Through.

\begin{rubric}
	In Advent, \nth{2} Collect of the Feria, \nth{3} of St. Saturninus (p. \pageref{SaturninusSecret}).
\end{rubric}
\begin{rubric}
	Outside of Advent, \nth{2} Collect of St. Saturninus (p. \pageref{SaturninusSecret}), \nth{3} of St. Mary  in Eastertide (p. \SPMaryEaster).
\end{rubric}

\communion{Andrew saith to Simon his brother: We have found the Messias, which is called the Christ: and he brought him to Jesus.}

\postcommunion
\lett{O}{Lord,} who hast bestowed on us these sacraments, we humbly beseech thee: that, at the intercession of thy blessed Apostle Andrew, the mysteries which we offer for his venerable passion may be profitable for our healing. Through.

\begin{rubric}
	In Advent, \nth{2} Collect of the Feria, \nth{3} of St. Saturninus (p. \pageref{SaturninusPostcommunion}).
\end{rubric}
\begin{rubric}
	Outside of Advent, \nth{2} Collect of St. Saturninus (p. \pageref{SaturninusPostcommunion}), \nth{3} of St. Mary  in Eastertide (p. \SPMaryEaster).
\end{rubric}


\subby{St. Saturninus}
\feastday{{St. Saturninus}}
\fancyhead[RE,LO]{29 November}
\begin{inhead}
    {Memorial\\
29 November}
\end{inhead}

\begin{rubric}
	The propers are from the Second Common of a Martyr not a Bishop (p. \pageref{CommonMartyrNotBishopII}), except the following Prayers.
\end{rubric}

\collect\label{SaturninusCollect}
\lett{O}{God,} who vouchsafest unto us to rejoice in the birthday of thy blessed Martyr Saturninus: grant, we pray thee; that we may be succoured by his merits. Through.

\secret\label{SaturninusSecret}
\lett{S}{anctify,} O Lord, we beseech thee, the offerings which we dedicate unto thee: and at the intercession of blessed Saturninus, thy Martyr, for their sake graciously regard us. Through.

\postcommunion\label{SaturninusPostcommunion}
\lett{W}{e} beseech thee, O Lord, that we, being sanctified by the receiving of thy sacrament: may, at the intercession of thy Saints, be thereby rendered acceptable unto thee. Through.

%MANUAL ADJUSTMENT:
\clearpage
\bcpfeast{30 November. St. Andrew}{St. Andrew}{30 November}
%\supplement{30 November}{St. Andrew}{}

\begin{secrubric}
	The Hymns and Versicles are of the Common of Apostles (p. \pageref{CommonApostles}), with the Antiphons as followeth.
\end{secrubric}

\properantiphon{Mag.}{One of the two {\dag} which followed the Lord was Andrew, Simon Peter's brother, alleluia.}

\properantiphon{Ben.}{Yield up to us {\dag} a man so righteous, restore to us a man so holy: destroy not a man so dear to God, righteous, dutiful, and gentle.}

\properantiphon{Mag.}{When blessed Andrew {\dag} came to the place where the Cross had been prepared, he cried out and said: O goodly Cross, so long desired, and now made ready for my eager spirit; fearless and joyful do I come to thee: therefore do thou also receive me gladly, as his disciple, who did hang upon thee.}

%MANUAL ADJUSTMENT:
\vspace{-1.75ex}

\subby{St. Peter Chrysologus}
\feastday{{St. Peter Chrysologus}}
\fancyhead[RE,LO]{2 December}
\begin{inhead}
    {Double\\
2 December}
\end{inhead}

\begin{rubric}
	The propers are from the Common of Doctors (p. \pageref{CommonDoctors}), except for the following.
\end{rubric}

\begin{rubric}
	\textsc{Note,} Commemoration is made of St. Bibiana and of the Advent Feria.
\end{rubric}

%MANUAL ADJUSTMENT:
\vspace{-1.5ex}

\collect
\lett{O}{God,} who by divine foreshewing wast pleased to choose blessed Peter Chrysologus thy illustrious Doctor to rule and instruct thy Church: grant, we beseech thee; that, as we have had him for a Doctor of life on earth, so we may be found worthy to have him for an intercessor in heaven. Through.

\gradual{Behold a great priest, who in his days pleased God. ℣. There was none found like unto him, who kept the law of the Most High.}

\communion{Lord. thou deliveredst unto me five talents: behold, I have gained beside them five talents more. Well done, thou good and faithful servant, thou hast been faithful over a few things, I will make thee ruler over many things, enter thou into the joy of thy lord.}

%MANUAL ADJUSTMENT:
\vspace{-1.75ex}

\subby{St. Bibiana}
\feastday{{St. Bibiana}}
\fancyhead[RE,LO]{2 December}
\begin{inhead}
    {Memorial\\
2 December}
\end{inhead}

\begin{rubric}
	The propers are from the Second Common of a Virgin Martyr (p. \pageref{CommonVirginMartyrII}), except for the following.
\end{rubric}

%MANUAL ADJUSTMENT:
\vspace{-1ex}

%MANUAL ADJUSTMENT:
\subsubsec{Collect}
\lett{O}{God,} the giver of all good gifts, who in thy handmaid Bibiana didst unite the palm of martyrdom with the flower of virginity: unite by her intercession our hearts in charity with thee; that all perils being done away, we may attain unto everlasting rewards. Through.

\secret
\lett{G}{raciously} receive, O Lord, through the merits of blessed Bibiana, thy Virgin and Martyr, the sacrifices which we offer unto thee: and grant that they may avail for our continual help. Through.

\postcommunion
\lett{O}{Lord} our God, who hast fulfilled us with the bounty of thy heavenly gift: we beseech thee, at the intercession of blessed Bibiana, thy Virgin and Martyr, we may ever live by the partaking of the same. Through.


\subby{St. Barbara}
\feastday{{St. Barbara}}
\fancyhead[RE,LO]{4 December}
\begin{inhead}
    {Memorial\\
4 December}
\end{inhead}

\begin{rubric}
	The propers are from the First Common of a Virgin Martyr (p. \pageref{CommonVirginMartyrI}).
\end{rubric}


\subby{St. Sabbas of Jud{\ae}a}
\feastday{{St. Sabbas}}
\fancyhead[RE,LO]{5 December}
\begin{inhead}
    {Memorial\\
5 December}
\end{inhead}

\begin{rubric}
	The propers are from the Common of an Abbot (p. \pageref{CommonAbbots}).
\end{rubric}

\begin{rubric}
	If today be Saturday, the Mass is of St. Mary on Saturday with \nth{2} Collect of the Feria \& \nth{3} of St. Sabbas.
\end{rubric}


\subby{St. Nicholas}
\feastday{{St. Nicholas}}
\fancyhead[RE,LO]{6 December}
\begin{inhead}
    {Double\\
6 December}
\end{inhead}

\begin{rubric}
	The Daily Office propers are from the First Common of a Confessor Bishop (p. \pageref{CommonConfessorBishopI}).
\end{rubric}

\introit
\lett{T}{he} Lord hath established a covenant of peace with him, and made him a prince: that he should have the dignity of the priesthood for ever. \textit{Ps.} Lord, remember David: and all his trouble.

\collect
\lett{O}{God,} who didst adorn thy blessed Bishop Nicholas with innumerable miracles: grant, we beseech thee; that by his merits and prayers we may be delivered from the fires of hell. Through.

\begin{rubric}
	Commemoration is of the Feria, unless it be Saturday.
\end{rubric}

\readingcitation{Epistle}{Hebrews 13:7}
\lett{B}{rethren:} Remember them which have the rule over you, who have spoken unto you the word of God: whose faith follow, considering the end of their conversation. Jesus Christ the same yesterday, and to day, and for ever. Be not carried about with divers and strange doctrines. For it is a good thing that the heart be established with grace; not with meats, which have not profited them that have been occupied therein. We have an altar, whereof they have no right to eat which serve the tabernacle. For the bodies of those beasts, whose blood is brought into the sanctuary by the high priest for sin, are burned without the camp. Wherefore Jesus also, that he might sanctify the people with his own blood, suffered without the gate. Let us go forth therefore unto him without the camp, bearing his reproach. For here have we no continuing city, but we seek one to come. By him therefore let us offer the sacrifice of praise to God continually, that is, the fruit of our lips giving thanks to his name. But to do good and to communicate forget not: for with such sacrifices God is well pleased. Obey them that have the rule over you, and submit yourselves: for they watch for your souls, as they that must give account.

\gradall{I have found David my servant, with my holy oil have I anointed him: my hand shall hold him fast, and my arm shall strengthen him. ℣. The enemy shall not be able to do him violence, the son of wickedness shall not hurt him.}{Alleluia, alleluia. ℣. The righteous shall flourish like a palm-tree: and shall spread abroad like a cedar in Libanus. Alleluia.}

\readingcitation{Gospel}{Matthew 25:14}
\lett{A}{t that time:} Jesus spake this parable to his disciples: A man travelling into a far country, who called his own servants, and delivered unto them his goods. And unto one he gave five talents, to another two, and to another one; to every man according to his several ability; and straightway took his journey. Then he that had received the five talents went and traded with the same, and made them other five talents. And likewise he that had received two, he also gained other two. But he that had received one went and digged in the earth, and hid his lord's money. After a long time the lord of those servants cometh, and reckoneth with them. And so he that had received five talents came and brought other five talents, saying, Lord, thou deliveredst unto me five talents: behold, I have gained beside them five talents more. His lord said unto him, Well done, thou good and faithful servant: thou hast been faithful over a few things, I will make thee ruler over many things: enter thou into the joy of thy lord. He also that had received two talents came and said, Lord, thou deliveredst unto me two talents: behold, I have gained two other talents beside them. His lord said unto him, Well done, good and faithful servant; thou hast been faithful over a few things, I will make thee ruler over many things: enter thou into the joy of thy lord.

\offertory{My truth and my mercy shall be with him: and in my name shall his horn be exalted.}

\secret
\lett{S}{anctify,} we bescech thee, O Lord God, these gifts which we offer on the solemnity of thy holy Bishop Nicholas: that our life may ever thereby be directed both in prosperity and in adversity. Through.

\begin{rubric}
	Commemoration is of the Feria, unless it be Saturday.
\end{rubric}

\communion{I have sworn once by my holiness: His seed shall endure for ever, and his seat is like as the sun before me, he shall stand fast for evermore as the moon, and as the faithful witness in heaven.}

\postcommunion
\lett{M}{ay} the sacrifices which we have received, O Lord, for the solemnity of thy holy Bishop Nicholas, preserve us by their everlasting protection. Through.

\begin{rubric}
	Commemoration is of the Feria, unless it be Saturday.
\end{rubric}


\subby{St. Ambrose}
\feastday{{St. Ambrose}}
\fancyhead[RE,LO]{7 December}
\begin{inhead}
    {Greater Double\\
7 December}
\end{inhead}

\begin{rubric}
	The Daily Office propers are from the Common of a Confessor Bishop (p. \pageref{CommonConfessorBishopI}).
\end{rubric}

\introit
\lett{I}{n} the midst of the Church he opened his mouth: and the Lord filled him with the spirit of wisdom and of understanding: he clothed him with a robe of glory. \textit{Ps.} It is a good thing to give thanks unto the Lord: and to sing praises unto thy name, O most Highest.

\collect
\lett{O}{God,} who didst give blessed Ambrose unto thy people to be a minister of everlasting salvation: grant, we beseech thee; that as we have learned of him the doctrine of life on earth, so we may be found worthy to have him for our advocate in heaven. Through.

\begin{rubric}
	Commemoration is of the Feria.
\end{rubric}

\begin{rubric}
	The Epistle is from the Common of Doctors (p. \pageref{CommonDoctors}).
\end{rubric}

\gradall{Behold a great priest, who in his days pleased God. ℣. There was none foun like unto him, who kept the law of the Most High.}{Alleluia, alleluia. ℣. The Lord sware, and will not repent: Thou art a priest for ever after the order of Melchisedech.}

%MANUAL ADJUSTMENT:
\clearpage
\begin{rubric}
{In Septuagesimatide or Lent, replacing the Alleluia:}
\end{rubric}\par\noindent
\tract{Blessed is the man that feareth the Lord: he hath great delight in his commandments. ℣. His seed shall be mighty upon earth: the generation of the faithful shall be blessed. ℣. Riches and plenteousness shall be in his house: and his righteousness endureth for ever.}

\begin{rubric}
{In Eastertide, replacing the Lesser Alleluia:}
\end{rubric}\par\noindent
\alleluia{Alleluia, alleluia. ℣. The Lord loved him, and adorned him: and clothed him with a robe of glory. Alleluia. ℣. The righteous shall grow as the lily and flourish for ever before the Lord. Alleluia.}

\begin{rubric}
	The Gospel is from the Common of Doctors (p. \pageref{CommonDoctors}).
\end{rubric}

\offertory{My truth and my mercy shall be with him: and in my name shall his horn be exalted.}

\secret
\lett{A}{lmighty} and everlasting God, grant, that the gifts which we present unto thy majesty, may through the intercession of blessed Ambrose, thy Confessor and Bishop, be profitable unto us for everlasting salvation. Through.

\begin{rubric}
	Commemoration is of the Feria.
\end{rubric}

\communion{I have sworn once by my holiness: His seed shall endure for ever, and his seat is like as the sun before me, he shall stand fast for evermore as the moon, and as the faithful witness in heaven.}

\postcommunion
\lett{G}{rant,} we beseech thee, almighty God: that we, receiving the sacraments of our salvation, may ever be aided by the prayer of blessed Ambrose thy Confessor and Bishop; in whose honour we have made these offerings unto thy majesty. Through.

\begin{rubric}
	Commemoration is of the Feria. In Lent, the Last Gospel is of the Feria.
\end{rubric}

%\begin{rubric}
%	If today be Saturday, the Mass is of the anticipated Vigil of the Conception of the Blessed Virgin Mary.
%\end{rubric}

%MANUAL ADJUSTMENT:
\clearpage

\bcpfeast{8 December. Conception of the Blessed Virgin Mary}{Conception B.V.M.}{8 December}
%\supplement{8 December}{Conception}{of the Blessed Virgin Mary}

\begin{secrubric}
	The Hymns are from the Common of the Blessed Virgin Mary (p. \pageref{CommonBVM}) with the following Versicle and Antiphons.
\end{secrubric}

℣. To-day is the Conception of the holy Virgin Mary.

℟. Whose glorious life illumineth all the churches.

\properantiphon{Mag.}{All generations {\dag} shall call me blessed: for he that is mighty hath magnified me, alleluia.}\\

\properantiphon{Ben.}{The Lord God said {\dag} unto the serpent, I will put enmity between thee and the woman, and between thy seed and her seed; it shall bruise thy head, alleluia.}\\

\properantiphon{Mag.}{Let us celebrate {\dag} the worshipful Conception of the blessed and glorious Virgin Mary, whose lowliness the Lord regarded when at the word of an Angel she conceived the world's Redeemer, alleluia.}


\subby{St. Melchiades}
\feastday{{St. Melchiades}}
\fancyhead[RE,LO]{10 December}
\begin{inhead}
    {Memorial\\
10 December}
\end{inhead}

\begin{rubric}
	The propers are from the First Common of a Martyr Bishop (p. \pageref{CommonMartyrBishopI}).
\end{rubric}


\subby{St. Damasus}
\feastday{{St. Damasus}}
\fancyhead[RE,LO]{11 December}
\begin{inhead}
    {Memorial\\
11 December}
\end{inhead}

\introit
\lett{L}{et} thy priests, O Lord, be clothed with righteousness, and let thy Saints sing with joyfulness: for thy servant David's sake turn not away the presence of thine Anointed. \textit{Ps.} Lord, remember David: and all his trouble.

\collect
\lett{G}{raciously} hear our prayers, O Lord: and at the intercession of blessed Damasus, thy Confessor and Bishop, mercifully grant us pardon and peace. Through.

\begin{rubric}
	Commemoration is of the Octave \& Feria.
\end{rubric}

\begin{rubric}
	The Epistle is from the Second Common of a Confessor Bishop (p. \pageref{CommonConfessorBishopII}).
\end{rubric}

\gradall{Behold a great priest who in his days pleased God. ℣. There was none found like unto him, who kept the law of the Most High.}{Alleluia, alleluia. ℣. Thou art a priest for ever after the order of Melchisedech. Alleluia.}

\begin{rubric}
	The Gospel is from the Second Common of a Confessor Bishop (p. \pageref{CommonConfessorBishopII}).
\end{rubric}

\communion{I have found David my servant, with my holy oil have I anointed him: my hand shall hold him fast, and my arm shall strengthen him.}

\secret
\lett{G}{rant,} O Lord, that like as thy dedicated people do acknowledge that in tribulation they have been succoured by the merits of thy Saints: so this oblation, which they offer unto thee in honour of the same, may be acceptable in thy sight. Through.

\begin{rubric}
	Commemoration is of the Octave \& Feria.
\end{rubric}

\communion{Lord, thou deliveredst unto me five talents, behold I have gained beside them five talents more. Well done, thou good and faithful servant, thou hast been faithful over a few things, I will make thee ruler over many things, enter thou into the joy of thy Lord.}

\postcommunion
\lett{G}{rant,} we beseech thee, O Lord, that thy faithful people may ever rejoice in the veneration of thy saints: and be defended by their perpetual supplication. Through.

\begin{rubric}
	Commemoration is of the Octave \& Feria.
\end{rubric}


\subby{St. Lucy}
\feastday{{St. Lucy}}
\fancyhead[RE,LO]{13 December}
\begin{inhead}
    {Greater Double\\
13 December}
\end{inhead}

\begin{rubric}
	The Daily Office propers are from the Common of Virgins (p. \pageref{CommonVirginOnlyI}), except that which followeth.
\end{rubric}

\antiphon{Mag.}{In thy patience {\dag}} thou didst possess thy soul, O Lucy, spouse of Christ: thou didst hate the things which are in the world, and thou shinest among the Angels: resisting unto blood, thou didst vanquish the enemy.

\antiphon{Ben.}{A pillar art thou {\dag} that may not be moved, O Lucy, spouse of Christ: and all the people are waiting until thou receive the crown of life, alleluia.}\\

℣. Full of grace are thy lips.

℟. Because God hath blessed thee for ever.

\antiphon{Mag.}{With such gravity {\dag} was she endued by the Holy Spirit, that the virgin of the Lord remained unmoved.}

\introit
\lett{T}{hou} hast loved righteousness, and hated iniquity: wherefore God, even thy God, hath anointed thee with the oil of gladness above thy fellows. \textit{Ps.} My heart is inditing of a good matter: I speak of the things which I have made unto the King.

\collect
\lett{G}{raciously} hear us, O God of our salvation: that, like as we do rejoice in the festival of blessed Lucy thy Virgin; so we may be instructed in all godly and devout affection. Through.

\begin{rubric}
	Commemoration is of the Octave \& Feria.
\end{rubric}

\begin{rubric}
	The Epistle is from the First Common of a Virgin (p. \pageref{CommonVirginOnlyI}).
\end{rubric}

\gradall{Thou hast loved righteousness, and hated iniquity. ℣. Wherefore God, even thy God, hath anointed thee with the oil of gladness.}{Alleluia, alleluia. ℣. Full of grace are thy lips: because God hath blessed thee for ever. Alleluia.}

\begin{rubric}
	In Votive Masses after Septuagesima, the Tract, and in Eastertide, the Alleluia, is from the First Common of a Virgin (p. \pageref{CommonVirginOnlyI}).
\end{rubric}

\readingcitation{Gospel}{Matthew 13:44}
\lett{A}{t that time:} Jesus spake this parable unto his disciples: The kingdom of heaven is like unto treasure hid in a field; the which when a man hath found, he hideth, and for joy thereof goeth and selleth all that he hath, and buyeth that field. Again, the kingdom of heaven is like unto a merchant man, seeking goodly pearls: Who, when he had found one pearl of great price, went and sold all that he had, and bought it. Again, the kingdom of heaven is like unto a net, that was cast into the sea, and gathered of every kind: Which, when it was full, they drew to shore, and sat down, and gathered the good into vessels, but cast the bad away. So shall it be at the end of the world: the angels shall come forth, and sever the wicked from among the just, And shall cast them into the furnace of fire: there shall be wailing and gnashing of teeth. Jesus saith unto them, Have ye understood all these things? They say unto him, Yea, Lord. Then said he unto them, Therefore every scribe which is instructed unto the kingdom of heaven is like unto a man that is an householder, which bringeth forth out of his treasure things new and old.

\offertory{The Virgins that be her fellows shall be brought unto the King: they that bear her company shall be brought unto thee with joy and gladness: and shall enter into the palace of the Lord the King.}

\secret
\lett{G}{rant,} O Lord, that like as thy dedicated people do acknowledge that in tribulation they have been succoured by the merits of thy Saints: so this oblation, which they offer unto thee in honour of the same, may be acceptable in thy sight. Through.

\begin{rubric}
	Commemoration is of the Octave \& Feria.
\end{rubric}

\communion{Princes have persecuted me without a cause, but my heart standeth in awe of thy word: I am as glad of thy word, as one that findeth great spoils.}

\postcommunion
\lett{O}{Lord,} who hast satisfied thy family with sacred gifts: we beseech thee; that we may at all times be comforted by the intercession of her whose festival we celebrate. Through.

\begin{rubric}
	Commemoration is of the Octave \& Feria.
\end{rubric}

\subby{St. Herman of Alaska}
\feastday{{St. Herman}}
\fancyhead[RE,LO]{13 December}
\begin{inhead}
    {Memorial\\
13 December}
\end{inhead}

\begin{rubric}
	The propers are from the First Common of a Confessor not Bishop (p. \pageref{CommonConfessorNotBishopI}).
\end{rubric}


\subby{Day VII within the Octave of the Conception}
\feastday{{Conception Octave}}
\fancyhead[RE,LO]{14 December}
\begin{inhead}
    {Semidouble\\
14 December}
\end{inhead}

\begin{rubric}
	Of the Octave, as on December 9. But if Ember Wednesday fall on this day, of the Ember Day, with a Commemoration of the Octave. The Mass is with the \nth{3} Prayer of the Holy Ghost.
\end{rubric}

\begin{rubric}
	If an Ember Day occur on any of the following Feasts, the Commemoration of the Feria is omitted.
\end{rubric}


\subby{Octave Day of the Conception of the Blessed Virgin Mary}
\feastday{{Conception Octave Day}}
\fancyhead[RE,LO]{15 December}
\begin{inhead}
    {Greater Double\\
15 December}
\end{inhead}

\begin{rubric}
	Hymn, Versicle, Antiphon, \& Mass as on the Feast, with Commemoration of the Feria.
\end{rubric}


\subby{St. Eusebius of Vercelli}
\feastday{{St. Eusebius}}
\fancyhead[RE,LO]{16 December}
\begin{inhead}
    {Memorial\\
16 December}
\end{inhead}

\begin{rubric}
	The propers are from the Second Common of a Martyr Bishop (p. \pageref{CommonMartyrBishopII}), with commemoration of the Feria (and the Collect of St. Mary in Mass).
\end{rubric}


\subby{Expectation of the Blessed Virgin Mary}
\feastday{{Expectation}}
\fancyhead[RE,LO]{18 December}
\begin{inhead}
    {Double\\
18 December}
\end{inhead}

%Is this the same as the Mass of the BVM from Advent to Christmas except for the Gradual?
%\subsubsec{Daily Office Propers}\par\noindent
%%From the Dominican Breviary (https://archive.org/details/BrevariumIuxtaRitumSacriOrdinisPraedicatorum/black-and-white%20version/Breviarium_Cormier_1909_prior_cropped%20%285.18x6.5%29%20pg%201-1577%20%28vol.%201%29/page/n826/mode/1up):
\par\noindent
\textit{Opening Sentence.} Send ye the lamb to the ruler of the land from Sela to the wilderness, unto the mount of the daughter of Zion.\vr{Is. 16:1}\par

\begin{rubric}
	\textsc{Note,} The Magnificat Antiphons are of the `O' Antiphons (p. \pageref{OAntiphons}).
\end{rubric}

\begin{paracol}{2}[]
\sloppy
\begin{inhead}
	I Evensong
\end{inhead}

℣. Hail Mary, full of grace.

℟. The Lord is with thee.

\begin{rubric}
	The Office Hymns are from the First Sunday of Advent (p. \pageref{AdventI}).
\end{rubric}

\switchcolumn
	
\begin{inhead}
	Mattins
\end{inhead}

℣. The Holy Ghost shall come upon thee.

℟. And the power of the Highest shall overshadow thee.\\

\antiphon{Ben.}{He shall sit upon the throne of David, {\dag} and of his kingdom, for ever.}

\fussy

\end{paracol}

\begin{rubric}
	II Evensong as in I Evensong.
\end{rubric}


%%From Bute's Roman Breviary translation (https://archive.org/details/theromanbreviary01unknuoft/page/665/mode/1up):
%O Antiphons will always take precedence.
%\antiphon{Mag.}{The Holy Ghost shall come upon thee, O Mary, * fear not; thou shalt bear in thy womb the Son of God. Alleluia.}\par

%\antiphon{Mag.}{O maiden of maidens, how shall this be, since neither before nor henceforth hath there been, nor shall be such another? Daughters of Jerusalem, why look ye curiously upon me? What ye see is a mystery of God.}\par

\begin{rubric}
	The Mass propers are from the Common of the Blessed Virgin Mary (p. \pageref{CommonBVM}), except for the Collect as followeth.
\end{rubric}

\collect
\lett{O}{God,} who wast pleased that thy Word should take flesh of the womb of the Blessed Virgin Mary at the message of an Angel: grant to thy humble servants; that we who believe her to be truly the Mother of God may be aided by her intercession in thy sight. Through the same.

%MANUAL ADJUSTMENT:
\clearpage

\subby{Vigil of St. Thomas}
\feastday{{St. Thomas Vigil}}
\fancyhead[RE,LO]{20 December}
\begin{inhead}
    {Vigil\\
20 December}
\end{inhead}

\begin{rubric}
	If today be Sunday, in the Ember Saturday Mass, Commemoration is made of the anticipated Vigil of St. Thomas and the last Gospel of the Vigil is read at the end, the \nth{3} Collect is of St. Mary.
\end{rubric}

\begin{rubric}
	The Mass propers are from the Vigil of Apostles (p. \pageref{CommonVigilApostles}), with Commemoration of the Feria of Advent and the \nth{3} Collect of St. Mary. But if an Ember Day occur, Commemoration is made of the Vigil in the Mass of the Feria.
\end{rubric}

\bcpfeast{21 December. St. Thomas}{St. Thomas}{21 December}
%\supplement{21 December}{St. Thomas}{}

\begin{secrubric}
	The Office Hymn and Versicle are from the Common of Apostles (p. \pageref{CommonApostles}), with the following Antiphon.
\end{secrubric}

\properantiphon{Mag. \& Ben.}{Because thou hast seen me, {\dag} Thomas, thou hast believed: blessed are they that have not seen, and yet have believed, alleluia.}


\subby{St. Paul the First Hermit}
\feastday{{St. Paul Hermit}}
\fancyhead[RE,LO]{10 January}
\begin{inhead}
    {Memorial\\
10 January}
\end{inhead}

\begin{rubric}
	The propers are from the Common of Abbots (p. \pageref{CommonAbbots}), except that which followeth.
\end{rubric}

\introit
\lett{T}{he} just shall flourish like a palm-tree: and shall spread abroad like a cedar in Libanus: planted in the house of the Lord: in the courts of the house of our God. \textit{Ps.} It is a good thing to give thanks unto the Lord: and to sing praises unto thy name, O Most Highest.

\collect
\lett{O}{God,} who makest us glad with the yearly solemnity of blessed Paul, thy Confessor: mercifully grant; that, as we now celebrate his birthday, so we may follow the example of his life. Through.

\gradall{The just shall flourish like a palm-tree: and shall spread abroad like a cedar in Libanus in the house of the Lord. ℣. To tell of thy loving-kindness early in the morning, and of thy truth in the night-season.}{Alleluia, alleluia. ℣. The just shall grow as the lily: and flourish for ever before the Lord. Alleluia.}

\readingcitation{Gospel}{Matthew 11:25}
\lett{A}{t that time:} Jesus answered and said: I thank thee, O Father, Lord of heaven and earth, because thou hast hid these things from the wise and prudent, and hast revealed them unto babes. Even so, Father: for so it seemed good in thy sight. All things are delivered unto me of my Father: and no man knoweth the Son, but the Father; neither knoweth any man the Father, save the Son, and he to whomsoever the Son will reveal him. Come unto me, all ye that labour and are heavy laden, and I will give you rest. Take my yoke upon you, and learn of me; for I am meek and lowly in heart: and ye shall find rest unto your souls. For my yoke is easy, and my burden is light.

\offertory{The just shall rejoice in thy strength, O Lord: exceeding glad shall he be of thy salvation: thou hast given him his heart's desire.}

\secret
\lett{G}{rant,} we beseech thee, O Lord, that we who, trusting in this our sacrifice of praise, do offer it before thee to the honour of thy Saints: may by the same be delivered from all evils both in this life and that which is to come. Through.

\communion{The just shall rejoice in the Lord, and put his trust in him: and all they that are true of heart shall be glad.}

\postcommunion
\lett{O}{Lord,} our God, who hast refreshed us with heavenly meat and drink, we humbly beseech thee: that we may be defended by the prayers of him in whose memory we have received the same. Through.


\subby{St. Hyginus of Rome}
\feastday{{St. Hyginus}}
\fancyhead[RE,LO]{11 January}
\begin{inhead}
    {Memorial\\
11 January}
\end{inhead}

\begin{rubric}
	The propers are from the First Common of a Martyr not a Bishop (p. \pageref{CommonMartyrNotBishopI}).
\end{rubric}


\subby{St. Benedict Biscop}
\feastday{{St. Benedict Biscop}}
\fancyhead[RE,LO]{12 January}
\begin{inhead}
    {Memorial\\
12 January}
\end{inhead}

\begin{rubric}
	The propers are from the Common of Abbots (p. \pageref{CommonAbbots}), except for that which followeth.
\end{rubric}

\collect
\lett{O}{God,} by whose gift the blessed Abbot Benedict left all things that he might be made perfect: grant unto all those who have entered upon the path of evangelical perfection, that they may neither look back nor linger in the way; but hastening to thee without stumbling, may lay hold upon eternal life. Through.

\secret
\lett{W}{e} beseech thee, O Lord, that thy holy Abbot Benedict, may may intercede for us: that this sacrifice which we offer and present upon thy holy altar may be profitable unto us for our salvation. Through.

\postcommunion
\lett{L}{et} thy sacrament, O Lord, which we have now received and the prayers of the blessed Abbot Benedict, effectually defend us: that we may both imitate the example of his conversion, and receive the succour of his intercession. Through.


\subby{St. Hilary}
\feastday{{St. Hilary}}
\fancyhead[RE,LO]{14 January}
\begin{inhead}
    {Double\\
14 January}
\end{inhead}


\begin{rubric}
	The propers are from the Common of Doctors (p. \pageref{CommonDoctors}).\par
	\textsc{Note,} Commemoration is made of St. Felix, with the Prayers from the Mass below.
\end{rubric}


\subby{St. Felix}
\feastday{{St. Felix}}
\fancyhead[RE,LO]{14 January}
\begin{inhead}
    {Memorial\\
14 January}
\end{inhead}

\begin{rubric}
	The propers are from the Second Common of a Martyr not a Bishop (p. \pageref{CommonMartyrNotBishopII}), except that which followeth.\par
\end{rubric}

\collect
\lett{G}{rant,} we beseech thee, almighty God, that the examples of thy Saints may provoke us to a better life: that as we celebrate their festival so we may imitate their actions. Through.

\secret
\lett{W}{e} beseech thee, O Lord, mercifully to accept this our sacrifice, which we offer unto thee, pleading the merits of blessed Felix, thy Martyr: that the same may avail for our perpetual succour. Through.

\postcommunion
\lett{O}{Lord,} who hast fulfilled us with saving mysteries: we beseech thee that we may be aided by the prayers of blessed Felix thy Martyr, whose festival we celebrate. Through.


%https://archive.org/details/missale-monasticum-propers/mode/2up?view=theater

\subby{St. Maurus}
\feastday{{St. Maurus}}
\fancyhead[RE,LO]{15 January}
\begin{inhead}
    {Greater Double\\
15 January}
\end{inhead}

\subbysub{I Evensong}

\begin{multicols}{2}
\begin{hangparas}{1.25em}{1}
Defender, leader true, thine own companions deem'd

Thy splendour half divine, since all were less esteem'd:

For thy most worthy deeds, Maurus, accept the lays

Wherewith we celebrate thy praise.\\

Born of a noble stock, great honour was his due,

But palaces he spurn'd, and from the world withdrew;

Delights he trampl'd down, estates and robes unpric'd,

To undergo the yoke of Christ.\\

The holy Abbot's grace, before his eyes display'd,

By deeds of equal worth he eagerly portray'd;

The pattern of the life monastic shone in truth

From every action of the youth.\\

Sternly, with sackcloth rough, self-mastery he wrought,

And by the curb of law, unbroken silence sought;

The ever-watchful nights in fervent prayer he spent;

Whole days of fasting underwent\\

Right speedily he flew to do the father's hest,

Dry-shod, the waters deep with fearless feet he press'd;

And safely he return'd with Placidus, set free

Like Peter walking on the sea.\\

To thee, O Trinity, high praise and honour be,

Whose countenance desir'd the heaven-dwellers see:

Grant that the Holy Rule may be our pathway plain

The prize of Maurus to attain.  Amen.\\
\end{hangparas}

℣. The Lord loved him and adorned him.

℟. He clothed him with a robe of glory.
\end{multicols}

\antiphon{Mag.}{O most blessed of men! {\dag} who, rejecting this world, bore the yoke of Holy Rule from tender years so lovingly; and being made obedient even unto death, he denied himself, that he might wholly cling to Christ his Master, alleluia.}

\subbysub{Mattins}

\begin{rubric}
	The Office Hymn \& Versicle are from the First Common of a Confessor not a Bishop (p. \pageref{CommonConfessorNotBishopI}), with the Antiphon from I Evensong.
\end{rubric}

\subbysub{II Evensong}

\begin{rubric}
	The Office Hymn is of I Evensong, with the Versicle \& Antiphon as followeth.
\end{rubric}

℣. The Lord guided the righteous in right paths.

℟. And shewed him the kingdom of God.\\

\antiphon{Mag.}{To-day holy Maurus, {\dag} lying upon a goat-skin, died happily before the altar; to-day the first-begotten disciple of blessed Benedict, through the guiding of the Holy Rule, came up to Christ, rising untroubled, accompanied by choirs of Angels; today the obedient man, telling his victories, was worthy to be crowned by the Lord, alleluia.}

\introit
\lett{T}{hy} way is in the sea, and thy paths in the great waters: and thy footsteps are not known. Thou leddest thy people like sheep. \textit{Ps.} The waters saw thee, O God, the waters saw thee, and were afraid: the depths also were troubled.

\collect
\lett{O}{God,} who for a pattern of obedience didst cause blessed Maurus to walk dry-shod upon the waters: grant that we may both follow perfectly the example of his virtues, and also be worthy to share in his reward. Through.

\readingcitation{Epistle}{Ecclesiasticus 51:13}
%RV:
\lett{W}{hen} I was yet young, Or ever I went abroad, I sought wisdom openly in my prayer. Before the temple I asked for her, And I will seek her out even to the end. From her flower as from the ripening grape my heart delighted in her: My foot trod in uprightness, From my youth I tracked her out. I bowed down mine ear a little, and received her, And found for myself much instruction. I profited in her: Unto him that giveth me wisdom I will give glory. For I purposed to practice her, And I was zealous for that which is good; And I shall never be put to shame. My soul hath wrestled in her, And in my doing I was exact: I spread forth my hands to the heaven above, And bewailed my ignorances of her. I set my soul aright unto her, And in pureness I found her. I gat me a heart joined with her from the beginning: Therefore shall I not be forsaken. My inward part also was troubled to seek her: Therefore have I gotten a good possession. The Lord gave me a tongue for my reward; And I will praise him therewith.
\par
Draw near unto me, ye unlearned, And lodge in the house of instruction. Say, wherefore are ye lacking in these things, And your souls are very thirsty? I opened my mouth, and spake, Get her for yourselves without money. Put your neck under the yoke, And let your soul receive instruction: She is hard at hand to find. Behold with your eyes, How that I laboured but a little, And found for myself much rest. Get you instruction with a great sum of silver, And gain much gold by her. May your soul rejoice in his mercy, And may ye not be put to shame in praising him. Work your work before the time cometh, And in his time he will give you your reward. 

\gradall{I will bring thy seed from the east, and gather thee from the west. ℣. When thou passest through the waters, I will be with thee; and through the rivers, they shall not overflow thee.}{Alleluia, alleluia. ℣. An obedient man shall speak of victory; Wealth and riches shall be in his house. Alleluia.}

\begin{rubric}
	In Septuagesimatide \& Lent, the Alleluia is omitted and the Tract is said instead.
\end{rubric}

\tract{Blessed is the man that feareth the Lord: he hath great delight in his commandments. ℣. His seed shall be mighty upon earth: the generation of the faithful shall be blessed. ℣. Riches and plenteousness shall be in his house: and his righteousness endureth for ever.}

\readingcitation{Gospel}{Matthew 14:28}
\lett{A}{t that time:} Peter answered Jesus and said, Lord, if it be thou, bid me come unto thee on the water. And he said, Come. And when Peter was come down out of the ship, he walked on the water, to go to Jesus. But when he saw the wind boisterous, he was afraid; and beginning to sink, he cried, saying, Lord, save me. And immediately Jesus stretched forth his hand, and caught him, and said unto him, O thou of little faith, wherefore didst thou doubt? And when they were come into the ship, the wind ceased. Then they that were in the ship came and worshipped him, saying, Of a truth thou art the Son of God.

\offertory{And the places that have been desolate for ages shall be built in thee: thou shalt raise up the foundations of generation and generation: turning the paths into rest. And I will feed thee with the inheritance of thy Father.}

%CHECK TRANSLATION:
\secret
\lett{M}{ay} the sacrifices we offer ascend unto thee, O Lord, as an odour of sweetness; and may the intercession of blessed Maurus, Abbot, intervene for us, that thy propitious power may descend upon us. Through.

\communion{I have chosen you, and ordained you, that ye should go and bring forth fruit, and that your fruit should remain: that whatsoever ye shall ask of the Father in my name, he may give it you.}

%CHECK TRANSLATION:
\postcommunion
\lett{W}{e} implore thy mercy, O Lord our God, that having received the pledges of our salvation and giving thanks for thy help; thou may accept our offerings unto our salvation; sending upon us thy grace to support us; that the celestial blessing, which thou hast brought us by the patronage of blessed Maurus, Abbot, may be perfected by continuing in imitation of him. Through.


\subby{St. Marcellus}
\feastday{{St. Marcellus}}
\fancyhead[RE,LO]{16 January}
\begin{inhead}
    {Memorial\\
16 January}
\end{inhead}

\begin{rubric}
	The propers are from the Second Common of a Martyr Bishop (p. \pageref{CommonMartyrBishopII}), except for that which followeth.
\end{rubric}

\introit
\lett{T}{he} Lord hath established a covenant of peace with him, and made him a prince: that he should have the dignity of the priesthood for ever. \textit{Ps.} Lord, remember David: and all his trouble.

\collect
\lett{O}{Lord,} we beseech thee favourably to hear the prayers of thy people: that, as we do rejoice in the passion of blessed Marcellus thy Martyr and Bishop, so we may be succoured by his merits. Through.

\gradall{I have found David my servant, with my holy oil have I anointed him: my hand shall hold him fast, and my arm shall strengthen him. ℣. The enemy shall not be able to do him violence, the son of wickedness shall not hurt him.}{Alleluia, alleluia. ℣. Thou art a priest for ever, after the order of Melchisedech. Alleluia.}

\offertory{My truth and my mercy shall be with him: and in my name shall his horn be exalted.}

\secret
\lett{R}{eceive,} O Lord, we beseech thee, the gifts which we duly offer: and by the pleading of the merits of blessed Marcellus thy Martyr and Bishop, grant that they may avail to set forward our salvation. Through.

\communion{Lord, thou deliveredst unto me five talents: behold, I have gained beside them five talents more. Well done, thou good and faithful servant, thou hast been faithful over a few things, I will make thee ruler over many things, enter thou into the joy of thy lord.}

\postcommunion
\lett{O}{Lord,} who hast satisfied thy family with sacred gifts: we beseech thee; that we may at all times be comforted by the intercession of him whose festival we celebrate. Through.


\subby{St. Anthony}
\feastday{{St. Anthony}}
\fancyhead[RE,LO]{17 January}
\begin{inhead}
    {Double\\
17 January}
\end{inhead}

\begin{rubric}
	The propers are from the Common of Abbots (p. \pageref{CommonAbbots}), except for the Gospel which is from the First Common of a Confessor not a Bishop (p. \pageref{CommonConfessorNotBishopI}).
\end{rubric}


\subby{Chair of St. Peter at Rome}
\feastday{{Chair St. Peter Rome}}
\fancyhead[RE,LO]{18 January}
\begin{inhead}
    {Greater Double\\
18 January}
\end{inhead}

\begin{rubric}
	The Daily Office propers are from the Chair of St. Peter at Antioch (p. \pageref{CathedraAntioch}).
\end{rubric}

\introit
\lett{T}{he} Lord hath established a covenant of peace with him, and made him a prince: that he should have the dignity of the priesthood for ever. \textit{Ps.} Lord, remember David: and all his trouble.

\collect
\lett{O}{God,} who didst bestow upon thy blessed Apostle Peter the keys of the kingdom of heaven, and didst appoint unto him the high priesthood of binding and loosing: vouchsafe; that by the help of his intercession we may be delivered from the bonds of our iniquities. Who livest and reignest.

\lett{O}{God,} who by the preaching of the blessed Apostle Paul didst teach the multitude of the Gentiles: grant to us, we beseech thee; that we who celebrate his commemoration may know him to be our advocate with thee. (Through.)

\begin{rubric}
    Commemoration is made of St. Prisca of Rome (p. \pageref{PriscaCollect}).
\end{rubric}

\readingcitation{Epistle}{1 Peter 1:1}
\lett{P}{eter} Peter, an apostle of Jesus Christ, to the strangers scattered throughout Pontus, Galatia, Cappadocia, Asia, and Bithynia, elect according to the foreknowledge of God the Father, through sanctification of the Spirit, unto obedience and sprinkling of the blood of Jesus Christ: Grace unto you, and peace, be multiplied. Blessed be the God and Father of our Lord Jesus Christ, which according to his abundant mercy hath begotten us again unto a lively hope by the resurrection of Jesus Christ from the dead, to an inheritance incorruptible, and undefiled, and that fadeth not away, reserved in heaven for you, who are kept by the power of God through faith unto salvation ready to be revealed in the last time. Wherein ye greatly rejoice, though now for a season, if need be, ye are in heaviness through manifold temptations: that the trial of your faith, being much more precious than of gold that perisheth, though it be tried with fire, might be found unto praise and honour and glory at the appearing of Jesus Christ our Lord.

\gradall{Let them exalt him in the congregation of the people: and praise him in the seat of the elders. ℣. O that men would praise the Lord for his goodness, and declare the wonders that he doeth for the children of men.}{Alleluia, alleluia. ℣. Thou art Peter, and upon this rock I will build my Church. Alleluia.}

\begin{rubric}
	In Septuagesimatide \& Lent, the Alleluia is replaced with the following.
\end{rubric}

\tract{Thou art Peter, and upon this rock I will build my Church. ℣. And the gates of hell shall not prevail against it: and I will give unto thee the keys of the kingdom of heaven. ℣. Whatsoever thou shalt bind on earth shall be bound in heaven. ℣. And whatsoever thou shalt loose on earth shall be loosed in heaven.}

\begin{rubric}
	In Eastertide, the Alleluia is replaced with the following.
\end{rubric}

\alleluia{Alleluia, alleluia. ℣. O that men would praise the Lord for his goodness, and declare the wonders that he doeth for the children of men. Alleluia. ℣. Thou art Peter, and upon this rock I will build my Church. Alleluia.}

\readingcitation{Gospel}{Matthew 16:13}
\lett{A}{t that time:} When Jesus came into the coasts of C{\ae}sarea Philippi, he asked his disciples, saying, Whom do men say that I the Son of man am? And they said, Some say that thou art John the Baptist: some, Elias; and others, Jeremias, or one of the prophets. He saith unto them, But whom say ye that I am? And Simon Peter answered and said, Thou art the Christ, the Son of the living God. And Jesus answered and said unto him, Blessed art thou, Simon Bar-jona: for flesh and blood hath not revealed it unto thee, but my Father which is in heaven. And I say also unto thee, That thou art Peter, and upon this rock I will build my church; and the gates of hell shall not prevail against it. And I will give unto thee the keys of the kingdom of heaven: and whatsoever thou shalt bind on earth shall be bound in heaven: and whatsoever thou shalt loose on earth shall be loosed in heaven.

\offertory{Thou art Peter, and upon this rock I will build my Church: and the gates of hell shall not prevail against it: and I will give unto thee the keys of the kingdom of heaven.}

\secret
\lett{W}{e} beseech thee, O Lord, that the intercession of blessed Peter the Apostle may commend unto thee the prayers and sacrifices of thy Church: that those things which we celebrate for his glory may avail for our pardon. Through.
\needspace{4\baselineskip}
\lett{S}{anctify,} O Lord, through the prayers of thine Apostle Paul, the gifts of thy people: that those things, which by thine institution are pleasing unto thee, may be made more pleasing by his prayer and advocacy. (Through.)

\begin{rubric}
    Commemoration is made of St. Prisca of Rome (p. \pageref{PriscaSecret}).
\end{rubric}

%MANUAL ADJUSTMENT:
\vspace{-1.5ex}

\communion{Thou art Peter, and upon this rock I will build my Church.}

%MANUAL ADJUSTMENT:
\vspace{-1ex}

\postcommunion
\lett{M}{ay} the gift, O Lord, which we have offered, make us to rejoice: that as we proclaim thy wonders in thine Apostle Peter; so through him we may receive the abundance of thy loving-kindness. Through.
\needspace{4\baselineskip}
\lett{O}{Lord,} who hast sanctified us with this saving mystery: we beseech thee; that he whom thou hast given to be our advocate and guide may never fail in prayer for us. (Through.)

\begin{rubric}
    Commemoration is made of St. Prisca of Rome (p. \pageref{PriscaPostcommunion}).
\end{rubric}

%MANUAL ADJUSTMENT:
\vspace{-1ex}

\subby{St. Prisca of Rome}
\feastday{{St. Prisca Rome}}
\fancyhead[RE,LO]{18 January}
\begin{inhead}
    {Memorial\\
18 January}
\end{inhead}

\begin{rubric}
	The propers are from the Second Common of a Virgin Martyr (p. \pageref{CommonVirginMartyrII}), except for that which followeth.
\end{rubric}

%MANUAL ADJUSTMENT:
\vspace{-1ex}

\collect\label{PriscaCollect}
\lett{G}{rant,} we beseech thee, almighty God: that we who celebrate the heavenly birthday of blessed Prisca, thy Virgin and Martyr; may both rejoice in her yearly solemnity, and profit by the example of so great a faith. Through.

%MANUAL ADJUSTMENT:
\vspace{-1ex}

\secret\label{PriscaSecret}
\lett{W}{e} beseech thee, O Lord, that this sacrifice which we offer in remembrance of the birthday of thy Saints may both loose us from the bonds of our iniquity, and obtain for us the gifts of thy mercy. Through.

%MANUAL ADJUSTMENT:
\subsubsec{Postcommunion}\label{PriscaPostcommunion}
%\postcommunion\label{PriscaPostcommunion}
\lett{O}{Lord,} who hast fulfilled us with saving mysteries: we beseech thee that we may be aided by the prayers of her whose festival we celebrate. Through.

%MANUAL ADJUSTMENT:
\clearpage
\subby{Sts. Marius, Martha, Audifax, \& Abachum}
\feastday{{Sts. Marius \&c.}}
\fancyhead[RE,LO]{19 January}
\begin{inhead}
    {Memorial\\
19 January}
\end{inhead}

\introit
\lett{L}{et} the righteous be glad and rejoice before God: let them also be merry and joyful. \textit{Ps.} Let God arise, and let his enemies be scattered: let them also that hate him flee before him.

\collect
\lett{G}{raciously} hear thy people, O Lord, who call upon thee with the advocacy of thy Saints: and grant us both to rejoice in peace in this temporal life; and to find succour unto life everlasting. Through.

\begin{rubric}
    Commemoration is made of St. Mark of Ephesus (p. \pageref{EphesianCollect}) and St. Mary in Epiphanytide (p. \SPMaryChristmas).
\end{rubric}

\begin{rubric}
	The Epistle is from the Third Common of Many Martyrs (p. \pageref{CommonMartyrsIII}).
\end{rubric}

\gradall{The souls of the just are in the hand of God: and there shall no torment of malice touch them. ℣. In the sight of the unwise they seemed to die: but they are in peace.}{Alleluia, alleluia. ℣. Our God is wonderful in his Saints. Alleluia.}

\begin{rubric}
	In Septuagesimatide \& Lent, the Alleluia is replaced with the Tract from the Third Common of Many Martyrs (p. \pageref{CommonMartyrsIII}).
\end{rubric}

\begin{rubric}
	The Gospel is from the first additional Gospel of the Third Common of Many Martyrs (p. \pageref{Matthew243}).
\end{rubric}

\offertory{Our soul is escaped, even as a bird out of the snare of the fowler: the snare is broken, and we are delivered.}

\secret
\lett{R}{egard,} O Lord, the prayers and oblations of thy faithful people: that they may be acceptable unto thee for the festival of thy Saints, and bestow on us the succour of thy mercy. Through.

\begin{rubric}
    Commemoration is made of St. Mark of Ephesus (p. \pageref{EphesianCollect}) and St. Mary in Epiphanytide (p. \SPMaryChristmas).
\end{rubric}

\communion{I say unto you, my friends: Be not afraid of them that persecute you.}

\postcommunion
\lett{G}{rant,} we beseech thee, that the intercession of thy Saints may make us acceptable unto thee: that those things which we perform in this temporal celebration we may receive unto eternal salvation. Through.

\begin{rubric}
    Commemoration is made of St. Mark of Ephesus (p. \pageref{EphesianCollect}) and St. Mary in Epiphanytide (p. \SPMaryChristmas).
\end{rubric}

%MANUAL ADJUSTMENT:
\vspace{-2ex}

%PROVISIONAL:
\subby{St. Mark of Ephesus}
\feastday{{St. Mark Ephesus}}
\fancyhead[RE,LO]{19 January}
\begin{inhead}
    {Memorial\\
19 January}
\end{inhead}

\begin{rubric}
	The propers are from the Second Common of a Confessor Bishop (p. \pageref{CommonConfessorBishopII}), except for that which followeth.
\end{rubric}

%Through thee shone uncreated light

%Of the life given to the least,

%And like a giant of great might

%Thou didst advance from the East.\\

%Thou didst illumine the whole world

%With the rays of thy true words,

%O divine Mark, Bishop of Christ,

%Pray for us for that same light.\\

%Thine el'quent lips and honey'd tongue

%Became the mouthpiece of his grace;

%Thy sacred tongue was shown to be

%A pen of wisdom, grace, and peace.\\

%All glory, Lord, to thee we pay

%For the Ephesian which thou gav'st,

%All glory, as is ever meet,

%To Father and to Paraclete. Amen.

\par\noindent
\textit{Yet to be approved.}

\vspace{-1ex}

\begin{tcolorbox}[
  colback=white,           % background color
  colframe=black,          % frame color
  fonttitle=\itshape,     % optional font for title
  title style={left=2mm},  % optional title left alignment
  attach title to upper={},
  enhanced,
  boxed title style={size=small,colframe=black!50!white,colback=white},
]

\antiphon{Mag.}{O blessed Mark {\dag} clothed with invincible armour, thou didst cast down rebellious pride.}

\antiphon{Ben.}{Thou didst serve {\dag} as the instrument of the Paraclete, and shone forth as the champion of Orthodoxy.}

\antiphon{Mag.}{Rejoice, Mark, {\dag} the boast of the Orthodox and joy to all true Catholics!}

\collect
\label{EphesianCollect}
\lett{O}{Lord} Jesu Christ, the only Shepherd and Bishop of our souls; we beseech thee to keep us in thy truth, preserve us by thy grace, and govern us according to thy loving-kindness. That as we imitate the good example of blessed Mark thy Confessor and Bishop, we may grow in love for thy Word and service of thy Church. Who with.

\secret
\label{EphesianSecret}
\lett{R}{eceive,} O most holy Father, these gifts now offered unto thee in memory of Saint Mark. That as he was a bulwark and defender of the Catholic Faith, so we may, by these gifts, be given the same firm love of thee. Through.

\postcommunion
\label{EphesianPostcommunion}
\lett{A}{bide} in us, O divine Comforter, who have here received thine holy Sacraments. That just as thy faithful servant Mark did contend for the unity of the Body of Christ, so we may be cleansed and knit together by the same. Who with the Father, from whom alone thou dost proceed, and the Son, who sendeth thy power upon us, liveth and reigneth, God, world without end. Amen.
\end{tcolorbox}

%MANUAL ADJUSTMENT:
\clearpage
\subby{Sts. Fabian \& Sebastian}
\feastday{{Sts. Fabian \& Sebastian}}
\fancyhead[RE,LO]{20 January}
\begin{inhead}
    {Double\\
20 January}
\end{inhead}

\begin{rubric}
	The Daily Office propers are from the First Common of Many Martyrs (p. \pageref{CommonMartyrsI}).
\end{rubric}

\introit
\lett{L}{et} the sorrowful sighing of the prisoners, O Lord, come before thee, reward thou our neighbours seven-fold into their bosom: avenge thou the blood of thy Saints that is shed. \textit{Ps.} O God, the heathen are come into thine inheritance: thy holy temple have they defiled: and made Jerusalem an heap of stones.

\collect
\lett{A}{lmighty} God, mercifully look upon our infirmities: that whereas we are oppressed by the burden of our sins, the glorious intercession of thy blessed Martyrs Fabian and Sebastian may be our succour and defence. Through.

\begin{rubric}
	The Epistle is the fifth additional Epistle of the Third Common of Many Martyrs (p. \pageref{Hebrews1133}).
\end{rubric}

\gradall{God is glorious in his holy ones: fearful in praises, doing wonders. ℣. Thy right hand, O Lord, is become glorious in power: thy right hand hath dashed in pieces the enemy.}{Alleluia, alleluia. ℣. Thy Saints give thanks unto thee, O Lord: they shew the glory of thy kingdom. Alleluia.}

\begin{rubric}
	In Septuagesimatide \& Lent, the Alleluia is replaced with the following.
\end{rubric}

\tract{They that sow in tears, shall reap in joy. ℣. He that now goeth on his way weeping, and beareth forth good seed. ℣. Shall doubtless come again with joy, and bring his sheaves with him.}

\readingcitation{Gospel}{Luke 6:17}
\lett{A}{t that time:} Jesus came down from the mountain, and stood in the plain, and the company of his disciples, and a great multitude of people out of all Jud{\ae}a and Jerusalem, and from the sea coast of Tyre and Sidon, which came to hear him, and to be healed of their diseases; and they that were vexed with unclean spirits: and they were healed. And the whole multitude sought to touch him: for there went virtue out of him, and healed them all. And he lifted up his eyes on his disciples, and said, Blessed be ye poor: for yours is the kingdom of God. Blessed are ye that hunger now: for ye shall be filled. Blessed are ye that weep now: for ye shall laugh. Blessed are ye, when men shall hate you, and when they shall separate you from their company, and shall reproach you, and cast out your name as evil, for the Son of man’s sake. Rejoice ye in that day, and leap for joy: for, behold, your reward is great in heaven.

\offertory{Be glad, O ye righteous, and rejoice in the Lord: and be joyful, all ye that are true of heart.}

\secret
\lett{W}{e} beseech thee, O Lord, mercifully to accept this our sacrifice which we offer unto thee, pleading the merits of thy blessed Martyrs Fabian and Sebastian: that the same may avail for our perpetual succour. Through.

\communion{A multitude of sick folk, and they that were vexed with unclean spirits came to him: for there went virtue out of him, and healed them all.}

\postcommunion
\lett{W}{e} beseech thee, O Lord our God, that like as we whom thou hast refreshed by the partaking of thy sacred gift do offer unto thee our worship: so, by the intercession of thy holy Martyrs Fabian and Sebastian, we may perceive the benefit of the same. Through.


\subby{St. Agnes of Rome}
\feastday{{St. Agnes Rome}}
\fancyhead[RE,LO]{21 January}
\begin{inhead}
    {Greater Double\\
21 January}
\end{inhead}

\begin{rubric}
	The Daily Office propers are from the Common of a Virgin (p. \pageref{CommonVirginOnlyI}), except for that which followeth.
\end{rubric}

\antiphon{Mag.}{Blessed Agnes, {\dag} in the midst of the flames, stretched out her hands and prayed: I call on thee, O Father transcendent, august, and dread; for by thy holy Son's protection I have escaped the threats of an impious tyrant, and passed unscathed through the foulness of fleshly pollution: and behold, I come to thee, whom I have loved, whom I have sought, whom I have alway desired.}

\antiphon{Ben.}{Lo, that which I desired, {\dag} now I see; that for which I hoped, I now possess; I am united in heaven unto him, whom on earth I loved with a perfect devotion.}\\

℣. Full of grace are thy lips.

℟. Because God hath blessed thee for ever.

\antiphon{Mag.}{While the blessed Agnes {\dag} was standing in the midst of the flames, she stretched out her hands, and prayed unto God: Almighty Lord, worthy of all adoration, fear, and worship: I bless thy holy Name, and I glorify thee for ever and ever.}

\introit
\lett{T}{he} ungodly wait for me to destroy me: O Lord, I will consider thy testimonies: I see that all things come to an end: but thy commandment is exceeding broad. \textit{Ps.} Blessed are those that are undefiled in the way: and walk in the law of the Lord.

\collect
\lett{A}{lmighty} and everlasting God, who dost choose the weak things of the world to confound those things that are strong: mercifully grant; that we who keep the feast of blessed Agnes thy Virgin and Martyr may feel the succour of her intercession in thy sight. Through.

\readingcitation{Epistle}{Ecclesiasticus 51:1}
%RV:
\lett{I}{will} give thanks unto thee, O Lord, O King, And will praise thee, O God my Saviour: I do give thanks unto thy name: For thou wast my protector and helper, And didst deliver my body out of destruction, And out of the snare of a slanderous tongue, From lips that forge lies, And wast my helper before them that stood by; And didst deliver me, according to the abundance of thy mercy, and greatness of thy name, From the gnashings of teeth ready to devour, Out of the hand of such as sought my life, Out of the manifold afflictions which I had; From the choking of a fire on every side, And out of the midst of fire which I kindled not; Out of the depth of the belly of the grave, And from an unclean tongue, And from lying words, The slander of an unrighteous tongue unto the king. My soul drew near even unto death, And my life was near to the grave beneath. They compassed me on every side, And there was none to help me. I was looking for the succour of men, And it was not. And I remembered thy mercy, O Lord, And thy working which hath been from everlasting, How thou deliverest them that wait for thee, And savest them out of the hand of the enemies, O Lord our God.

%\lett{I}{will} give thanks unto thee, O Lord, O King, And will praise thee, O God my Saviour: I do give thanks unto thy name: For thou wast my protector and helper, And didst deliver my body out of destruction, And out of the snare of a slanderous tongue, From lips that forge lies, And wast my helper before them that stood by;  And didst deliver me, according to the abundance of thy mercy, and greatness of thy name, From the gnashings of teeth ready to devour, Out of the hand of such as sought my life, Out of the manifold afflictions which I had; From the choking of a fire on every side, And out of the midst of fire which I kindled not; Out of the depth of the belly of the grave, And from an unclean tongue, And from lying words, The slander of an unrighteous tongue unto the king. My soul drew near even unto death, And my life was near to the grave beneath. And I remembered thy mercy, O Lord, And thy working which hath been from everlasting, How thou deliverest them that wait for thee, And savest them out of the hand of the enemies, O Lord our God.

\gradall{Full of grace are thy lips: because God hath blessed thee for ever. ℣. Because of the word of truth, of meekness, and righteousness: and thy right hand shall teach thee terrible things.}{Alleluia, alleluia. ℣. The five wise virgins took oil in their vessels with their lamps: and at midnight there was a cry made: Behold, the bridegroom cometh: go ye out to meet Christ the Lord. Alleluia.}

\begin{rubric}
{In Septuagesimatide or Lent, replacing the Alleluia:}
\end{rubric}\par\noindent
\tract{Come, Spouse of Christ, receive the crown which the Lord hath prepared for thee for ever: for love of whom thou didst shed thy blood. ℣. Thou hast loved righteousness and hated iniquity: wherefore God, even thy God, hath anointed thee with the oil of gladness above thy fellows. ℣. In thy comeliness and in thy beauty go forth, ride prosperously, and reign.}

\readingcitation{Gospel}{Matthew 25:1}
\lett{A}{t that time:} Jesus spake this parable unto his disciples: The kingdom of heaven shall be likened unto ten virgins, which took their lamps, and went forth to meet the bridegroom. And five of them were wise, and five were foolish. They that were foolish took their lamps, and took no oil with them:  But the wise took oil in their vessels with their lamps. While the bridegroom tarried, they all slumbered and slept. And at midnight there was a cry made, Behold, the bridegroom cometh; go ye out to meet him. Then all those virgins arose, and trimmed their lamps. And the foolish said unto the wise, Give us of your oil; for our lamps are gone out. But the wise answered, saying, Not so; lest there be not enough for us and you: but go ye rather to them that sell, and buy for yourselves. And while they went to buy, the bridegroom came; and they that were ready went in with him to the marriage: and the door was shut. Afterward came also the other virgins, saying, Lord, Lord, open to us. But he answered and said, Verily I say unto you, I know you not. Watch therefore, for ye know neither the day nor the hour wherein the Son of man cometh.

\offertory{The Virgins that be her fellows shall be brought unto the King: they that bear her company shall be brought unto thee with joy and gladness: and shall enter into the palace of the Lord the King.}

\secret
\lett{O}{Lord,} mercifully regard the sacrifices which we offer unto thee: and at the intercession of blessed Agnes, thy Virgin and Martyr, absolve us from the bonds of our sins. Through.

\communion{The five wise virgins took oil in their vessels with their lamps: and at midnight there was a cry made: Behold, the bridegroom cometh: go ye out to meet Christ the Lord.}

\postcommunion
\lett{O}{Lord,} our God, who hast refreshed us with heavenly meat and drink, we humbly beseech thee: that we may be defended by the prayers of her in whose memory we have received the same. Through.


\subby{Sts. Vincent \& Anastasius}
\feastday{{St. Vincent \& Anastasius}}
\fancyhead[RE,LO]{22 January}
\begin{inhead}
    {Memorial\\
22 January}
\end{inhead}

%New propers needed for addition of Anastasius

\begin{rubric}
	The propers are from the First Common of Many Martyrs (p. \pageref{CommonMartyrsI}).
\end{rubric}

%\begin{rubric}
%	The propers are from the Second Common of a Martyr not a Bishop (p. \pageref{CommonMartyrNotBishopII}), except for that which followeth.
%\end{rubric}

%\collect
%\lett{A}{ssist} us, O Lord, in our supplications: that we who perceive the guilt of our iniquities, may be delivered through the intercession of thy blessed Martyr Vincent. Through.

%CHECK TRANSLATION:
%\secret
%\lett{W}{e} beseech thee, O Lord, to accept our service and prayers, cleanse us by thy heavenly mysteries, and graciously hear us. Through.

%CHECK TRANSLATION:
%\postcommunion
%\lett{W}{e} beseech thee, almighty God; that, by the intercession of thy blessed Martyr Vincent, we who have received heavenly food may be strengthened against all adversity by the same. Through.


\subby{St. Emerentiana}
\feastday{{St. Emerentiana}}
\fancyhead[RE,LO]{23 January}
\begin{inhead}
    {Memorial\\
23 January}
\end{inhead}

\begin{rubric}
	The propers are from the Second Common of a Virgin and Martyr (p. \pageref{CommonVirginMartyrII}).
\end{rubric}


\subby{St. Timothy}
\feastday{{St. Timothy}}
\fancyhead[RE,LO]{24 January}
\begin{inhead}
    {Double\\
24 January}
\end{inhead}

\begin{rubric}
	The propers come from the First Common of a Bishop and Martyr (p. \pageref{CommonMartyrBishopI}), except for the Epistle, as below.
\end{rubric}

\readingcitation{Epistle}{1 Timothy 6:11}
\lett{D}{early beloved:} Follow after righteousness, godliness, faith, love, patience, meekness. Fight the good fight of faith, lay hold on eternal life, whereunto thou art also called, and hast professed a good profession before many witnesses. I give thee charge in the sight of God, who quickeneth all things, and before Christ Jesus, who before Pontius Pilate witnessed a good confession; That thou keep this commandment without spot, unrebukeable, until the appearing of our Lord Jesus Christ: Which in his times he shall shew, who is the blessed and only Potentate, the King of kings, and Lord of lords; Who only hath immortality, dwelling in the light which no man can approach unto; whom no man hath seen, nor can see: to whom be honour and power everlasting. Amen.

\bcpfeast{25 January. Conversion of St. Paul the Apostle}{St. Paul's Conversion}{25 January}
%\supplement{25 January}{Conversion}{of St. Paul the Apostle}

\subbysub{I Evensong}\label{PaulEvensong}

\gregorioscore{resources/gabc/ProperTime/PaulEvensong.gabc}

℣. Thou art a chosen vessel, holy Apostle Paul.

℟. A preacher of the truth throughout all the world.

\properantiphon{Mag.}{Go forth, Ananias, {\dag} and enquire for Saul, for behold, he prayeth: for he is a chosen vessel unto me, to bear my Name before the Gentiles, and Kings, and the children of Israel.}

\subbysub{Mattins}

\begin{rubric}
	The Invitatory Hymn is as in I Evensong.
\end{rubric}

\begin{rubric}
The Office Hymn is from the Common of the Apostles (p. \pageref{CommonApostles}), with the following Versicle \& Antiphon.
\end{rubric}

℣. Thou art a chosen vessel, holy Apostle Paul.

℟. A preacher of the truth throughout all the world.

\properantiphon{Ben.}{Ye which have followed me {\dag} shall sit upon twelve thrones, judging the twelve tribes of Israel, saith the Lord.}

\subbysub{II Evensong}

\begin{rubric}
	The Office Hymn \& Versicle are as in I Evensong, with the following Antiphon.
\end{rubric}

\properantiphon{Mag.}{O holy Apostle Paul, {\dag} thou preacher of the truth and Doctor of the Gentiles, intercede for us unto God, who hath chosen thee.}


\subby{St. Polycarp of Smyrna}
\feastday{{St. Polycarp Smyrna}}
\fancyhead[RE,LO]{26 January}
\begin{inhead}
    {Memorial\\
26 January}
\end{inhead}

\introit
\lett{O}{ye} priests of the Lord, bless ye the Lord: O ye holy and humble men of heart, bless ye the Lord. \textit{Ps.} O all ye works of the Lord, bless ye the Lord: praise him, and magnify him for ever.

\collect
\lett{O}{God,} who makest us glad with the yearly solemnity of blessed Polycarp, thy Martyr and Bishop: mercifully grant; that, as we now cebrate his birthday, so we may likewise rejoice in his protection. Through.

\readingcitation{Epistle}{1 John 3:10}
\lett{D}{early beloved:} Whosoever doeth not righteousness is not of God, neither he that loveth not his brother. For this is the message that ye heard from the beginning, that we should love one another. Not as Cain, who was of that wicked one, and slew his brother. And wherefore slew he him? Because his own works were evil, and his brother's righteous. Marvel not, my brethren, if the world hate you. We know that we have passed from death unto life, because we love the brethren. He that loveth not his brother abideth in death. Whosoever hateth his brother is a murderer: and ye know that no murderer hath eternal life abiding in him. Hereby perceive we the love of God, because he laid down his life for us: and we ought to lay down our lives for the brethren.

\gradall{Thou hast crowned him with glory and worship. ℣. Thou hast made him to have dominion of the works of thy hands, O Lord.}{Alleluia, alleluia. ℣. This is a priest whom the Lord hath crowned. Alleluia.}

\begin{rubric}
{In Septuagesimatide or Lent, replacing the Alleluia:}
\end{rubric}\par\noindent
\tract{Blessed is the man that feareth the Lord: he hath great delight in his commandments. ℣. His seed shall be mighty upon earth: the generation of the faithful shall be blessed. ℣. Riches and plenteousness shall be in his house: and his righteousness endureth for ever.}

\readingcitation{Gospel}{Matthew 10:26}
\lett{A}{t that time:} Jesus said to his disciples: There is nothing covered, that shall not be revealed; and hid, that shall not be known. What I tell you in darkness, that speak ye in light: and what ye hear in the ear, that preach ye upon the housetops. And fear not them which kill the body, but are not able to kill the soul: but rather fear him which is able to destroy both soul and body in hell. Are not two sparrows sold for a farthing? and one of them shall not fall on the ground without your Father. But the very hairs of your head are all numbered. Fear ye not therefore, ye are of more value than many sparrows. Whosoever therefore shall confess me before men, him will I confess also before my Father which is in heaven.

\offertory{I have found David my servant, with my holy oil have I anointed him: my hand shall hold him fast, and my arm shall strengthen him.}

\secret
\lett{S}{anctify,} O Lord, the gifts which we dedicate to thee: that at the intercession of blessed Polycarp, thy Martyr and Bishop, they may obtain for us thy gracious favour. Through.

\communion{Thou hast set, O Lord, a crown of pure gold upon his head.}

\postcommunion
\lett{W}{e} beseech thee, O Lord our God, that like as we, whom thou hast refreshed by the partaking of thy sacred gift, do offer unto thee our worship: so, by the intercession of blessed Polycarp, thy Martyr and Bishop, we may perceive the benefit of the same. Through.


\subby{St. John Chrysostom}
\feastday{{St. John Chrysostom}}
\fancyhead[RE,LO]{27 January}
\begin{inhead}
    {Double\\
27 January}
\end{inhead}

\begin{rubric}
	The propers are from the Common of Doctors (p. \pageref{CommonDoctors}), except for that which followeth.
\end{rubric}

\collect
\lett{M}{ultiply,} we beseech thee, O Lord, thy Church with thy heavenly grace: even as thou didst vouchsafe to enlighten her with the glorious merits and doctrine of blessed John Chrysostom, thy Confessor and Bishop. Through.

\gradall{Behold a great priest, who in his days pleased God. ℣. There was none found like unto him, who kept the law of the Most High.}{Alleluia, alleluia. ℣. Blessed is the man that endureth temptation: for when he is tried, he shall receive the crown of life, alleluia.}

\begin{rubric}
{In Septuagesimatide or Lent, replacing the Alleluia:}
\end{rubric}\par\noindent
\tract{Blessed is the man that feareth the Lord: he hath great delight in his commandments. ℣. His seed shall be mighty upon earth: the generation of the faithful shall be blessed. ℣. Riches and plenteousness shall be in his house: and his righteousness remaineth for ever.}

\secret
\lett{M}{ay} the devout prayers of Saint John Chrysostom thy Bishop and Doctor, never fail to succour us, O Lord: that they may render our oblations acceptable in thy sight; and may ever obtain for us thy merciful pardon. Through.

\postcommunion
\lett{W}{e} beseech thee, O Lord, that blessed John Chrysostom, thy Bishop and illustrious Doctor; may stand before thee as our advocate: that these thy sacrifices may avail for our salvation. Through.


\subby{St. Cyril of Alexandria}
\feastday{{St. Cyril Alexandria}}
\fancyhead[RE,LO]{28 January}
\begin{inhead}
    {Double\\
28 January}
\end{inhead}

\begin{rubric}
	The propers are from the Common of Doctors (p. \pageref{CommonDoctors}), except for that which followeth.
\end{rubric}

\collect
\lett{O}{God,} who didst make blessed Cyril, thy Confessor and Bishop, an invincible defender of the divine Motherhood of the most blessed Virgin Mary: grant, by his intercession; that we, who believe her to be indeed the Mother of God, may be saved through her maternal protection. Through the same.

\begin{rubric}
    Commemoration is made of St. Agnes (p. \pageref{AgnesCollectII}).
\end{rubric}

\secret
\lett{A}{lmighty} God, graciously look upon our gifts: and at the intercession of blessed Cyril vouchsafe; that we may be found meet worthily to receive in our hearts thine only-begotten Son, Jesus Christ our Lord, co-eternal with thee in thy glory. Who liveth and reigneth with thee.

\begin{rubric}
    Commemoration is made of St. Agnes (p. \pageref{AgnesSecretII}).
\end{rubric}

\postcommunion
\lett{O}{Lord,} who hast refreshed us with divine mysteries, we humbly beseech thee: that, being aided by the example and merits of the blessed Bishop Cyril, we may be enabled worthily to serve the most holy Mother of thine only-begotten Son. Who liveth and reigneth with thee.

\begin{rubric}
    Commemoration is made of St. Agnes (p. \pageref{AgnesPostcommunionII}).
\end{rubric}


\subby{The Second Feast of St. Agnes}
\feastday{{Second St. Agnes}}
\fancyhead[RE,LO]{28 January}
\begin{inhead}
    {Memorial\\
28 January}
\end{inhead}

\begin{rubric}
	The propers are from the First Common of a Virgin (p. \pageref{CommonVirginOnlyI}), except for that which followeth.
\end{rubric}

\introit
\lett{A}{ll} the rich among the people shall make their supplication before thee: the Virgins that be her fellows shall be brought unto the King: they that bear her company shall be brought unto thee with joy and gladness. \textit{Ps.} My heart is inditing of a good matter: I speak of the things which I have made unto the King.

\collect\label{AgnesCollectII}
\lett{O}{God,} who makest us glad with the yearly solemnity of blessed Agnes, thy Virgin and Martyr: grant, we beseech thee; that as we venerate her in our service, so we may follow the example of her godly conversation. Through.

\readingcitation{Gospel}{Matthew 13:44}
\lett{A}{t that time:} Jesus spake this parable unto his disciples: The kingdom of heaven is like unto treasure hid in a field; the which when a man hath found, he hideth, and for joy thereof goeth and selleth all that he hath, and buyeth that field. Again, the kingdom of heaven is like unto a merchant man, seeking goodly pearls: Who, when he had found one pearl of great price, went and sold all that he had, and bought it. Again, the kingdom of heaven is like unto a net, that was cast into the sea, and gathered of every kind: Which, when it was full, they drew to shore, and sat down, and gathered the good into vessels, but cast the bad away. So shall it be at the end of the world: the angels shall come forth, and sever the wicked from among the just, And shall cast them into the furnace of fire: there shall be wailing and gnashing of teeth. Jesus saith unto them, Have ye understood all these things? They say unto him, Yea, Lord. Then said he unto them, Therefore every scribe which is instructed unto the kingdom of heaven is like unto a man that is an householder, which bringeth forth out of his treasure things new and old.

\offertory{Full of grace are thy lips, because God hath blessed thee for ever and ever.}

\secret\label{AgnesSecretII}
\lett{L}{et} thy plenteous benediction, we beseech thee, O Lord, come down upon these sacrifices: that it may mercifully work out our sanctification, and make us to rejoice in the solemnity of thy Martyrs. Through.

\communion{The kingdom of heaven is like unto a merchant man, seeking goodly pearls: who when he had found one pearl of great price, gave all that he had, and bought it.}

\postcommunion\label{AgnesPostcommunionII}
\lett{G}{rant,} we beseech thee, O Lord, that the sacrament which we have received in our observance of this yearly festival: may bestow on us thy healing; both in this temporal life and unto life eternal. Through.


\subby{St. Martina}
\feastday{{St. Martina}}
\fancyhead[RE,LO]{30 January}
\begin{inhead}
    {Memorial\\
30 January}
\end{inhead}

\begin{rubric}
	The propers are from the First Common of a Virgin and Martyr (p.\pageref{CommonVirginMartyrI}).
\end{rubric}


\subby{St. Ignatius of Antioch}
\feastday{{St. Ignatius Antioch}}
\fancyhead[RE,LO]{1 February}
\begin{inhead}
    {Double\\
1 February}
\end{inhead}

\begin{rubric}
	The Daily Office propers are from the First Common of Martyr Bishops (p. \pageref{CommonMartyrBishopI}).
\end{rubric}

\introit
\lett{B}{ut} God forbid that I should glory, save in the Cross of our Lord Jesus Christ: by whom the world is crucified unto me, and I unto the world. \textit{Ps.} Lord, remember David: and all his trouble.

\collect
\lett{A}{lmighty} God, mercifully look upon our infirmities; that whereas we are oppressed by the burden of our sins, the glorious intercession of blessed Ignatius thy Martyr and Bishop may be our succour and defence. Through.

\begin{rubric}
    Commemoration is made of St. Bridget of Ireland, from the First Common of a Virgin (p. \pageref{CommonVirginOnlyI}).
\end{rubric}

\readingcitation{Epistle}{Romans 8:35}
\lett{B}{rethren:} Who shall separate us from the love of Christ? shall tribulation, or distress, or persecution, or famine, or nakedness, or peril, or sword? As it is written, For thy sake we are killed all the day long; we are accounted as sheep for the slaughter. Nay, in all these things we are more than conquerors through him that loved us. For I am persuaded, that neither death, nor life, nor angels, nor principalities, nor powers, nor things present, nor things to come, Nor height, nor depth, nor any other creature, shall be able to separate us from the love of God, which is in Christ Jesus our Lord.

\gradall{Behold a great priest, who in his days pleased God. ℣. There was none found like unto him, who kept the law of the Most High.}{Alleluia, alleluia. ℣. I am crucifed with Christ: I live, yet not I, but Christ liveth in me. Alleluia.}

\begin{rubric}
{In Septuagesimatide or Lent, replacing the Alleluia:}
\end{rubric}\par\noindent
\tract{Thou hast given him his heart's desire: and hast not denied him the request of his lips. ℣. For thou hast prevented him with the blessings of goodness. ℣. Thou hast set a crown of pure gold upon his head.}

\readingcitation{Gospel}{John 12:24}
\lett{A}{t that time:} Jesus said unto his disciples: Verily, verily, I say unto you, Except a corn of wheat fall into the ground and die, it abideth alone: but if it die, it bringeth forth much fruit. He that loveth his life shall lose it; and he that hateth his life in this world shall keep it unto life eternal. If any man serve me, let him follow me; and where I am, there shall also my servant be: if any man serve me, him will my Father honour.

\offertory{Thou hast crowned him with glory and worship: and hast made him to have dominion of the works of thy hands, O Lord.}

\secret
\lett{W}{e} beseech thee, O Lord, mercifully to accept this our sacrifice which we offer unto thee, pleading the merits of blessed Ignatius thy Martyr and Bishop: that the same may avail for our perpetual succour. Through.

\begin{rubric}
    Commemoration is made of St. Bridget of Ireland, from the First Common of a Virgin (p. \pageref{CommonVirginOnlyI}).
\end{rubric}

\communion{I am the wheat of Christ: let me be ground by the teeth of beasts, that I may be found pure bread.}

\postcommunion
\lett{W}{e} beseech thee, O Lord, our God, that as we whom thou hast refreshed by the partaking of thy sacred gift offer unto thee our worship: so by the intercession of blessed Ignatius thy Martyr and Bishop, we may perceive the benefit of the same. Through.

\begin{rubric}
    Commemoration is made of St. Bridget of Ireland, from the First Common of a Virgin (p. \pageref{CommonVirginOnlyI}).
\end{rubric}


%%CHECK MISSAL:
\subby{St. Bridget of Ireland}
\feastday{{St. Bridget}}
\fancyhead[RE,LO]{1 February}
\begin{inhead}
    {Memorial\\
1 February}
\end{inhead}

\begin{rubric}
	The propers are from the First Common of a Virgin (p. \pageref{CommonVirginOnlyI}).
\end{rubric}

%From Common(?):
%\collect\label{BridgetCollect}
%\lett{H}{ear} us, O God our Saviour: and grant that as we rejoice in the festival of blessed Bridget, thy Virgin, so we may be brought to perfect holiness by the love of piety and devotion. Through.

\bcpfeast{2 February. Purification of the Blessed Virgin Mary}{Purification B.V.M.}{2 February}
%\supplement{2 February}{Purification}{of the Blessed Virgin Mary}

\begin{secrubric}
	The Office Hymn and Versicle are from the Common of the Blessed Virgin Mary (p. \pageref{CommonBVM}), except for the Evensong Versicles and the Antiphons as followeth.
\end{secrubric}

℣. It was revealed unto Simeon by the Holy Ghost. 

℟. That he should not see death before he had seen the Lord's Christ.

\properantiphon{Mag.}{The ancient {\dag} carried the Infant, but the Infant governed the ancient: he whom a Virgin bare, and after bearing, remained virgin, the same was worshipped by her who bare him.}\\

\properantiphon{Ben.}{And when the parents {\dag} brought in the Child Jesus, then Simeon took him up in his arms, and blessed God, saying: Lord, now lettest thou thy servant depart in peace.}\\

\properantiphon{Mag.}{To-day {\dag} the blessed Virgin Mary presented the Child Jesus in the temple; and Simeon, filled with the Holy Spirit, received him into his arms, and blessed God for ever.}


\subby{St. Blaise}
\feastday{{St. Blaise}}
\fancyhead[RE,LO]{3 February}
\begin{inhead}
    {Memorial\\
3 February}
\end{inhead}

\begin{rubric}
	The propers are from the Second Common of a Martyr Bishop (p. \pageref{CommonMartyrBishopII}).
\end{rubric}


\subby{St. Joseph of Aleppo}
\feastday{{St. Joseph Aleppo}}
\fancyhead[RE,LO]{4 February}
\begin{inhead}
    {Memorial\\
4 February}
\end{inhead}

\begin{rubric}
	The propers are from the First Common of a Martyr not a Bishop (p. \pageref{CommonMartyrNotBishopI}).
\end{rubric}

\begin{rubric}
	Commemoration is made of the New Martyrs of Russia, from the Prayers in the First Common of Many Martyrs (p. \pageref{CommonMartyrsI}).
\end{rubric}

\subby{The New Martyrs of Russia}
\feastday{{New Russian Martyrs}}
\fancyhead[RE,LO]{4 February}
\begin{inhead}
    {Memorial\\
4 February}
\end{inhead}

\begin{rubric}
	The propers are from the First Common of Many Martyrs (p. \pageref{CommonMartyrsI}).
\end{rubric}


\subby{St. Agatha}
\feastday{{St. Agatha}}
\fancyhead[RE,LO]{5 February}
\begin{inhead}
    {Greater Double\\
5 February}
\end{inhead}

\begin{rubric}
	The Office Hymns and Versicles are from the First Common of a Virgin Martyr (p. \pageref{CommonVirginMartyrI}), except for the II Evensong Versicle as below.
\end{rubric}

\antiphon{Mag.}{The blessed Agatha, {\dag} standing in the midst of the prison, with outstretched hands entreated the Lord: O Lord Jesus Christ, my gracious Master, I give thanks unto thee, who hast enabled me to overcome the torments of the executioner: bid me now, O Lord, joyfully to enter into thine unfading glory.}\\

\antiphon{Ben.}{The multitude {\dag} of the heathen, fleeing to the tomb of the virgin, took thence her veil to defend them from the fire: that the Lord might shew himself a deliverer from the burning, for the merits of Agatha his blessed Martyr.}\\

℣. Full of grace are thy lips.

℟. Because God hath blessed thee for ever.

\antiphon{Mag.}{The blessed Agatha, {\dag} standing in the midst of the prison, with outstretched hands entreated the Lord: O Lord Jesus Christ, my gracious Master, I give thanks unto thee, who hast enabled me to overcome the torments of the executioner: bid me now, O Lord, joyfully to enter into thine unfading glory.}

\introit
\lett{R}{ejoice} we all in the Lord, keeping feast day in honour of blessed Agatha, the Virgin and Martyr: in whose passion the Angels rejoice, and glorify the Son of God. \textit{Ps.} My heart is inditing of a good matter: I speak of the things which I have made unto the King.

\collect
\lett{O}{God,} who among the manifold works of thy power hast bestowed even upon the weakness of women the victory of martyrdom: mercifully grant; that we, who celebrate the birthday of blessed Agatha, thy Virgin and Martyr, may by her example be drawn nearer unto thee. Through.

\readingcitation{Epistle}{1 Corinthians 1:26}
\lett{B}{rethren:} Ye see your calling: how that not many wise men after the flesh, not many mighty, not many noble, are called: But God hath chosen the foolish things of the world to confound the wise; and God hath chosen the weak things of the world to confound the things which are mighty; And base things of the world, and things which are despised, hath God chosen, yea, and things which are not, to bring to nought things that are: That no flesh should glory in his presence. But of him are ye in Christ Jesus, who of God is made unto us wisdom, and righteousness, and sanctification, and redemption: That, according as it is written, He that glorieth, let him glory in the Lord.

\gradall{God shall help her with his countenance: God is in the midst of her, therefore shall she not be removed. ℣. The rivers of the flood thereof shall make glad the city of God: the holy place of the tabernacle of the Most Highest.}{Alleluia, alleluia. ℣. I will speak of thy testimonies even before kings: and will not be ashamed. Alleluia.}

\begin{rubric}
{In Septuagesimatide or Lent, replacing the Alleluia:}
\end{rubric}\par\noindent
\tract{They that sow in tears, shall reap in joy. ℣. They that now go on their way weeping, and bear forth good seed. ℣. Shall doubtless come again with joy, and bring their sheaves with them.}

\readingcitation{Gospel}{Matthew 19:3}
\lett{A}{t that time:} The Pharisees came unto Jesus, tempting him, and saying unto him: Is it lawful for a man to put away his wife for every cause? And he answered and said unto them, Have ye not read, that he which made them at the beginning made them male and female, And said, For this cause shall a man leave father and mother, and shall cleave to his wife: and they twain shall be one flesh? Wherefore they are no more twain, but one flesh. What therefore God hath joined together, let not man put asunder. They say unto him, Why did Moses then command to give a writing of divorcement, and to put her away? He saith unto them, Moses because of the hardness of your hearts suffered you to put away your wives: but from the beginning it was not so. And I say unto you, Whosoever shall put away his wife, except it be for fornication, and shall marry another, committeth adultery: and whoso marrieth her which is put away doth commit adultery. His disciples say unto him, If the case of the man be so with his wife, it is not good to marry. But he said unto them, All men cannot receive this saying, save they to whom it is given. For there are some eunuchs, which were so born from their mother's womb: and there are some eunuchs, which were made eunuchs of men: and there be eunuchs, which have made themselves eunuchs for the kingdom of heaven's sake. He that is able to receive it, let him receive it.

\offertory{The Virgins that be her fellows shall be brought unto the King: they that bear her company shall be brought unto thee.}

%MANUAL ADJUSTMENT:
\vspace{-1ex}

\secret
\lett{R}{eceive,} O Lord, the gifts which we offer on the solemnity of blessed Agatha, thy Virgin and Martyr, through whose advocacy we trust to be delivered. Through.

\communion{He who deigned to heal my every wound, and to restore my breast unto my body, on him do I call, the living God.}

%MANUAL ADJUSTMENT:
\vspace{-1ex}

\postcommunion
\lett{M}{ay} the mysteries which we have received be for our succour, O Lord: and at the intercession of blessed Agatha, thy Virgin and Martyr, cause us to rejoice in thy continual protection. Through.


\subby{St. Dorothea}
\feastday{{St. Dorothea}}
\fancyhead[RE,LO]{6 February}
\begin{inhead}
    {Memorial\\
6 February}
\end{inhead}

\begin{rubric}
	The propers are from the Second Common of a Virgin Martyr (p. \pageref{CommonVirginMartyrII}), except for that which followeth.
\end{rubric}

%MANUAL ADJUSTMENT:
\vspace{-1.5ex}

\collect
\lett{W}{e} beseech thee, O Lord, that as blessed Dorothea, thy Virgin and Martyr, was ever found pleasing unto thee, both by the merit of her chastity, and by her confession of thy power: so she may implore for us thy pardon. Through.

%MANUAL ADJUSTMENT:
\vspace{-1.5ex}

\secret
\lett{G}{raciously} receive, O Lord, through the merits of bressed Dorotnea thy Virgin and Martyr, the sacrifices which we offer unto thee and grant that they may avail for our continual help. Through.

\postcommunion
\lett{O}{Lord} our God, who hast fulfilled us with the bounty of thy heavenly gift: we beseech thee that, at the intercession of blessed Dorothea thy Virgin and Martyr, we may ever live by the partaking of the same. Through.

%MANUAL ADJUSTMENT:
\clearpage
\subby{St. Romuald}
\feastday{{St. Romuald}}
\fancyhead[RE,LO]{7 February}
\begin{inhead}
    {Double\\
7 February}
\end{inhead}

\begin{rubric}
	The propers are from the Common of Abbots (p. \pageref{CommonAbbots}).
\end{rubric}


\subby{St. Apollonia of Alexandria}
\feastday{{St. Apollonia}}
\fancyhead[RE,LO]{9 February}
\begin{inhead}
    {Memorial\\
9 February}
\end{inhead}

\begin{rubric}
	The propers are from the First Common of a Virgin Martyr (p. \pageref{CommonVirginMartyrI}), except for that which followeth.
\end{rubric}

\collect
\lett{O}{God} who among the manifold works of thy power hast bestowed even upon the weakness of women the victory of martyrdom mercifully grant; that we, who celebrate the birthday of blessed Apollonia thy Virgin and Martyr, may by her example be drawn nearer unto thee. Through.

\secret
\lett{R}{eceive,} O Lord, the gifts which we offer on the solemnity of blessed Apollonia thy Virgin and Martyr: through whose advocacy we trust to be delivered. Through.

\postcommunion
\lett{M}{ay} the mysteries which we have received be for our succour, O Lord: and at the intercession of blessed Apollonia thy Virgin and Martyr, cause us to rejoice in thy continual protection. Through.


\bcpfeast{10 February. St. Scholastica}{St. Scholastica}{10 February}

\subbysub{I Evensong}\label{ScholasticaEvensong}

\gregorioscore{resources/gabc/ProperTime/ScholasticaEvensong.gabc}

℣. Who is this that flieth as a cloud? 

℟. And as a dove to her windows?

\antiphon{Mag.}{Let all the multitude {\dag} of the faithful exult in the glory of the gracious virgin Scholastica: and chiefly let the company of virgins be joyful, celebrating her Solemnity; for she besought the Lord, pouring forth her tears, and of him received greater power, because her love was greater.}

\subbysub{Mattins}

\invitatoryhymn\label{ScholasticaInvitatory}

\gregorioscore{resources/gabc/ProperTime/ScholasticaInvitatory.gabc}

\officehymn\label{ScholasticaMattins}

\gregorioscore{resources/gabc/ProperTime/ScholasticaMattins.gabc}

℣. Behold, thou art fair, my love.  

℟. Behold, thou art fair; thou hast dove's eyes.

\antiphon{Ben.}{O how illustrious {\dag} are the merits of blessed Scholastica! O how great the power of her tears! through which the renowned virgin, out of sunny clearness, drew down from the air a mighty flood of rain.}

\subbysub{II Evensong}

\begin{rubric}
	The Office Hymn \& Versicle are of I Evensong, with the following Antiphon.
\end{rubric}

\antiphon{Mag.}{To-day {\dag} the holy virgin Scholastica, in the likeness of a dove, went forth with all gladness to the heavenly places: to-day she was found worthy to enjoy for ever the bliss of celestial life beside her brother.}


\subby{Pope St. Gregory II}
\feastday{{Pope St. Gregory II}}
\fancyhead[RE,LO]{11 February}
\begin{inhead}
    {Memorial\\
11 February}
\end{inhead}

\begin{rubric}
	The propers are from the First Common of a Confessor Bishop (p. \pageref{CommonConfessorBishopI}).
\end{rubric}


\subby{St. Valentine}
\feastday{{St. Valentine}}
\fancyhead[RE,LO]{14 February}
\begin{inhead}
    {Memorial\\
14 February}
\end{inhead}

\begin{rubric}
	The propers are from the First Common of a Martyr not a Bishop (p. \pageref{CommonMartyrNotBishopI}), except for that which followeth.
\end{rubric}

\collect
\lett{G}{rant,} we beseech thee, almighty God: that we who observe the birthday of blessed Valentine thy Martyr may by his intercession be delivered from all evils that beset us. Through.

\secret
\lett{R}{eceive,} O Lord, we beseech thee, the gifts which we duly offer: and by the pleading of the merits of blessed Valentine thy Martyr, grant; that they may avail to set forward our salvation. Through.

\postcommunion
\lett{M}{ay} this heavenly mystery, O Lord, renew us in soul and body: that as we offer unto thee our worship, so by the intercession of blessed Valentine thy Martyr, we may perceive the benefit of the same. Through.


\subby{Sts. Faustinus and Jovita}
\feastday{{Sts. Faustinus \& Jovita}}
\fancyhead[RE,LO]{15 February}
\begin{inhead}
    {Memorial\\
15 February}
\end{inhead}

\begin{rubric}
	The propers are from the Third Common of Many Martyrs (p. \pageref{CommonMartyrsIII}), except for that which followeth.
\end{rubric}

\collect
\lett{O}{God,} who makest us glad with the yearly solemnity of thy holy Martyrs Faustinus and Jovita: mercifully grant; that as we rejoice in their merits, so we may be enkindled by their example. Through.

\secret
\lett{A}{ssist} us mercifully, O Lord, in these our supplications which we make before thee in remembrance of thy Saints: that we who trust not in our own rightneousness may be succoured by the merits of them that have found favour in thy sight. Through.

\postcommunion
\lett{O}{Lord,} who hast fulfilled us with saving mysteries, we beseech thee: that we may be aided by the prayers of those whose festival we celebrate. Through.


\subby{St. Simeon of Jerusalem}
\feastday{{St. Simeon}}
\fancyhead[RE,LO]{18 February}
\begin{inhead}
    {Memorial\\
18 February}
\end{inhead}

\begin{rubric}
	The propers are from the First Common of a Martyr Bishop (p. \pageref{CommonMartyrBishopI}).
\end{rubric}

\bcpfeast{22 February. Chair of St. Peter at Antioch}{Chair St. Peter Antioch}{22 February}\label{CathedraAntioch}
%\supplement{22 February}{Chair}{of St. Peter at Antioch}

\subbysub{I Evensong}

\gregorioscore{resources/gabc/ProperTime/PeterAntiochEvensong.gabc}\label{PeterAntiochEvensong}
    
    ℣. Thou art Peter.

	℟. And upon this rock I will build my Church.
	
\properantiphon{Mag.}{Thou art the shepherd of the sheep, {\dag} Prince of the Apostles, unto thee were given the keys of the kingdom of heaven.}

\subbysub{Mattins}

\begin{rubric}
	The Invitatory Hymn is as in I Evensong.
\end{rubric}

\officehymn\label{PeterAntiochMattins}

\gregorioscore{resources/gabc/ProperTime/PeterAntiochMattins.gabc}

    ℣. Thou art Peter.

	℟. And upon this rock I will build my Church.
	
	\properantiphon{Ben.}{Whatsoever {\dag} thou shalt bind on earth shall be bound in heaven: and whatsoever thou shalt loose on earth shall be loosed in heaven, saith the Lord unto Simon Peter.}

\subbysub{II Evensong}

\begin{rubric}
	II Evensong as in I Evensong, except for the Antiphon, as followeth.
\end{rubric}

\properantiphon{Mag.}{Whilst he was chief Bishop, he feared nothing on earth, but ascended gloriously to the heavenly kingdoms.}


\subby{Vigil of St. Matthias}
\feastday{{St. Matthias Vigil}}
\fancyhead[RE,LO]{23 February}
\begin{inhead}
    {Vigil\\
23 February}
\end{inhead}

\begin{rubric}
	The propers are from the Common of Vigils of the Apostles (p. \pageref{CommonVigilApostles}).
\end{rubric}


\bcpfeast{24 February. St. Matthias}{St. Matthias}{24 February}
%\supplement{24 February}{St. Matthias}{}

\begin{secrubric}
	The Daily Office propers are of the Common of Apostles (p. \pageref{CommonApostles}).
\end{secrubric}


\subby{St. Walburga of Heidenheim}
\feastday{{St. Walburga}}
\fancyhead[RE,LO]{25 February}
\begin{inhead}
    {Memorial\\
25 February (26 February in a Leap Year)}
\end{inhead}

\begin{rubric}
	The propers are from the Common of a Virgin (p. \pageref{CommonVirginOnlyI}), except for that which followeth.
\end{rubric}

%From Latin Diurnale Monasticum \& Missale Romanum (https://www.google.com/books/edition/Missale_Romanum/p2HPgCPUBPsC?hl=en&gbpv=1&dq=%22Walb%C3%BArgae%22&pg=RA3-PA126&printsec=frontcover):
\collect
\lett{O}{God,} who among the manifold gifts of thy grace dost also work great wonders even in the weaker sex: mercifully grant that we may feel the help of the intercession of blessed Walburga, thy Virgin, whose example of chastity doth enlighten us, and whose glory in miracles doth gladden us. Through.
%\lett{D}{eus,} qui inter innúmera gráti{\ae} tu{\ae} dona, étiam in sexu frágili tua operáris magnália: {\dag} concéde propítius; ut beát{\ae} Walbúrg{\ae}, Vírginis tu{\ae}, apud misericórdiam tuam patrocínia sentiámus, * cujus non solum castitátis illustrámur exémplis, verum étiam miraculórum glória jucundámur. Per Dóminum.

\begin{rubric}
	In Lent, Commemoration of the Feria.
\end{rubric}

\secret
\lett{W}{e} beseech thee, O Lord, that the prayer of blessed Walburga thy Virgin may render this sacrifice, which we devoutly offer unto thee in her venerable commemoration, acceptable unto thy Majesty. Through.
%\lett{S}{acrifícium} oblátum, quǽsumus Dómine, beát{\ae} Walbúrg{\ae} Vírginis tu{\ae} orátio Majestáti tu{\ae} reddat accéptum: pro cujus veneránda commemoratióne tibi devóto offértur obséquio. Per Dóminum.

\begin{rubric}
	In Lent, Commemoration of the Feria.
\end{rubric}

\postcommunion
\lett{G}{rant,} O Lord, that by the intercession of blessed Walburga, thy Virgin, we may obtain the grace of thy blessing: that as we proclaim her venerable glory, she may aid us in all our necessities. Through.
%\lett{B}{enedictiónis,} Dómine, grátiam, intercedénte beáta Walbúrga Vírgine tua, cónsequi mereámur: ut, cujus venerándam glóriam pr{\ae}dicámus; ejus in ómnibus nostris necessitátibus auxílium sentiámus. Per Dóminum.

\begin{rubric}
	In Lent, Commemoration \& Last Gospel of the Feria.
\end{rubric}


\subby{Pope St. Alexander of Alexandria}
\feastday{{Pope St. Alexander Alexandria}}
\fancyhead[RE,LO]{26 February}
\begin{inhead}
    {Memorial\\
26 February (27 February in a Leap Year)}
\end{inhead}

\begin{rubric}
	The propers are from the First Common of a Confessor Bishop (p. \pageref{CommonConfessorBishopI}).
\end{rubric}
\begin{rubric}
	In Lent, Commemoration \& Last Gospel of the Feria.
\end{rubric}


\subby{St. Raphael of Brooklyn}
\feastday{{St. Raphael Brooklyn}}
\fancyhead[RE,LO]{27 February}
\begin{inhead}
    {Greater Double\\
27 February (28 February in a Leap Year)}
\end{inhead}

\begin{rubric}
	In Lent, Commemoration only.
\end{rubric}

\begin{rubric}
	The propers are from the First Common of a Confessor Bishop (p. \pageref{CommonConfessorBishopI}).
\end{rubric}
\begin{rubric}
	In Lent, Commemoration \& Last Gospel of the Feria.
\end{rubric}


\subby{St. David of Wales}
\feastday{{St. David Wales}}
\fancyhead[RE,LO]{1 March}
\begin{inhead}
    {Memorial\\
1 March}
\end{inhead}

\begin{rubric}
	The propers are from the First Common of a Confessor Bishop (p. \pageref{CommonConfessorBishopI}), except for that which followeth.
\end{rubric}

\collect
\lett{G}{rant} to us, almighty God: that the loving intercession of blessed David, thy Confessor and Bishop, may protect us; that while we celebrate his festival we may imitate his steadfastness in the defence of the Catholic faith. Through.

\begin{rubric}
	In Lent, Commemoration of the Feria.
\end{rubric}

\secret
\lett{W}{e} beseech thee, O Lord: that we, remembering with gladness the merits of thy Saints, may in all places feel the succour of their intercession. Through.

\begin{rubric}
	In Lent, Commemoration of the Feria.
\end{rubric}

\postcommunion
\lett{G}{rant,} we beseech thee, almighty God: that we, shewing forth our thankfulness for the gifts which we have received, may, at the intercession of blessed David, thy Confessor and Bishop, obtain yet more abundant mercies. Through.

\begin{rubric}
	In Lent, Commemoration \& Last Gospel of the Feria.
\end{rubric}


\subby{St. Chad}
\feastday{{St. Chad}}
\fancyhead[RE,LO]{2 March}
\begin{inhead}
    {Memorial\\
2 March}
\end{inhead}

\begin{rubric}
	The propers are from the Second Common of a Confessor Bishop (p. \pageref{CommonConfessorBishopII}), except for that which followeth.
\end{rubric}

\collect
\lett{A}{lmighty} and everlasting God, who on this day dost gladden us by the festival of blessed Chad, thy Confessor and Bishop: we humbly beseech thy mercy; that we, who devoutly observe and venerate his festival, may by his loving advocacy obtain the reward of everlasting life. Through.

\begin{rubric}
	In Lent, Commemoration of the Feria.
\end{rubric}

\secret
\lett{W}{e} beseech thee, O Lord, mercifully to have respect unto our supplications: and at the intercession of blessed Chad thy Confessor and Bishop on our behalf, grant that we who minister thy heavenly sacraments may be free from all sin; that through thy purifying grace we may be cleansed by those same mysteries which we serve. Through.

\begin{rubric}
	In Lent, Commemoration of the Feria.
\end{rubric}

\postcommunion
\lett{G}{rant,} we beseech thee, O Lord our God: that, at the intercession of blessed Chad, thy Confessor and Bishop, we who have tasted of holy things, being cleansed by divine mysteries, may attain to the fulness of this heavenly sacrament. Through.

\begin{rubric}
	In Lent, Commemoration \& Last Gospel of the Feria.
\end{rubric}


\subby{Pope St. Lucius I}
\feastday{{Pope St. Lucius I}}
\fancyhead[RE,LO]{4 March}
\begin{inhead}
    {Memorial\\
4 March}
\end{inhead}

\begin{rubric}
	The propers are from the Second Common of a Martyr Bishop (p. \pageref{CommonMartyrBishopII}), except for that which followeth.
\end{rubric}

\collect
\lett{O}{God,} who makest us glad with the yearly solemnity of blessed Lucius thy Martyr and Bishop: mercifully grant; that, as we now celebrate his birthday, so we may likewise rejoice in his protection. Through.

\begin{rubric}
	In Lent, Commemoration of the Feria.
\end{rubric}

\secret
\lett{W}{e} beseech thee, O Lord, mercifully to accept this our sacrifice, which we offer unto thee, pleading the merits of blessed Lucius, thy Martyr and Bishop: that the same may avail for our perpetual succour. Through.

\begin{rubric}
	In Lent, Commemoration of the Feria.
\end{rubric}

\postcommunion
\lett{W}{e} beseech thee, O Lord our God, that like as we whom thou hast refreshed by the partaking of thy sacred gift do offer unto thee our worship: so, by the intercession of blessed Lucius thy Martyr and Bishop, we may perceive the benefit of the same. Through.

\begin{rubric}
	In Lent, Commemoration \& Last Gospel of the Feria.
\end{rubric}


\subby{Sts. Perpetua and Felicitas}
\feastday{{Sts. Perpetua \& Felicitas}}
\fancyhead[RE,LO]{6 March}
\begin{inhead}
    {Double\\
6 March}
\end{inhead}

\begin{rubric}
	In Lent, Commemoration only.
\end{rubric}

\begin{rubric}
	The propers are from the Common of a Martyr not a Virgin (p. \pageref{CommonMartyrNotVirgin}), except for that which followeth.
\end{rubric}

%Felicity changed to Felicitas
\collect
\lett{G}{rant,} we beseech thee, O Lord our God, that we may at all times so devoutly honour the triumphs of thy holy Martyrs Perpetua and Felicitas: that, although we cannot worthily shew forth their praises, yet we may continually honour them with lowly service. Through.

\begin{rubric}
	In Lent, Commemoration of the Feria.
\end{rubric}

\secret
\lett{O}{Lord,} we beseech thee, look down upon these gifts, which we offer on thine altars on this festival of thy holy Martyrs Perpetua and Felicitas: that as by these blessed mysteries thou hast bestowed glory upon them, so likewise of thy bounty thou wouldest vouchsafe to us thy pardon. Through.

\begin{rubric}
	In Lent, Commemoration of the Feria.
\end{rubric}

\postcommunion
\lett{O}{Lord,} who hast fulfilled us with mystic gifts and joys: grant, we beseech thee; that by the intercession of thy holy Martyrs, Perpetua and Felicitas, we may spiritually attain to those things which we temporally perform. Through.

\begin{rubric}
	In Lent, Commemoration \& Last Gospel of the Feria.
\end{rubric}


\subby{St. Gregory of Nyssa}
\feastday{{St. Gregory Nyssa}}
\fancyhead[RE,LO]{9 March}
\begin{inhead}
    {Memorial\\
9 March}
\end{inhead}

\begin{rubric}
	The propers are from the Common of Doctors (p. \pageref{CommonDoctors}).
\end{rubric}
\begin{rubric}
	In Lent, Commemoration \& Last Gospel of the Feria.
\end{rubric}


\subby{The Forty Holy Martyrs}
\feastday{{40 Holy Martyrs}}
\fancyhead[RE,LO]{10 March}
\begin{inhead}
    {Memorial\\
10 March}
\end{inhead}

\introit
\lett{T}{he} just cry, and the Lord heareth them: and delivereth them out of all their troubles. \textit{Ps.} I will alway give thanks unto the Lord; his praise shall ever be in my mouth.

\collect
\lett{G}{rant} we beseech thee, almighty God: that, like as we have known thy glorious Martyrs to be constant in their confession, so we may perceive their loving intercession for us with thee. Through.

\begin{rubric}
	In Lent, Commemoration of the Feria.
\end{rubric}

\begin{rubric}
	The Epistle is the fifth additional Epistle of the Third Common of Many Martyrs (p. \pageref{Hebrews1133}).
\end{rubric}

\gradtr{Behold, how good and joyful a thing it is, brethren, to dwell together in unity! ℣. It is like the precious ointment upon the head, that ran down unto the beard, even unto Aaron's beard.}{They that sow in tears, shall reap in joy. ℣. They that now go on their way weeping, and bear forth good seed. ℣. Shall doubtless come again with joy, and bring their sheaves with them.}

\begin{rubric}
	The Gospel is from the Second Common of Many Martyrs (p. \pageref{CommonMartyrsII}).
\end{rubric}

\offertory{Be glad, O ye righteous, and rejoice in the Lord: and be joyful, all ye that are true of heart.}

\secret
\lett{R}{egard,} O Lord, the prayers and oblations of thy faithful people: that they may be acceptable unto thee for the festival of thy Saints, and bestow on us the succour of thy mercy. Through.

\begin{rubric}
	In Lent, Commemoration of the Feria.
\end{rubric}

\communion{Whosoever shall do the will of my Father which is in heaven: the same is my brother, and sister, and mother, saith the Lord.}

\postcommunion
\lett{G}{rant,} O Lord, we beseech thee, that the intercession of thy Saints may make us acceptable unto thee: that those things which we perform in this temporal celebration we may receive unto eternal salvation. Through.

\begin{rubric}
	In Lent, Commemoration \& Last Gospel of the Feria.
\end{rubric}

%MANUAL ADJUSTMENT:
\clearpage
\bcpfeast{12 March. Pope St. Gregory the Great}{Pope St. Gregory}{12 March}
%\supplement{12 March}{St. Gregory}{Pope, Confessor, \& Martyr}

\subbysub{I Evensong}\label{GregoryEvensong}

\gregorioscore{resources/gabc/ProperTime/GregoryEvensong.gabc}

    ℣. The Lord loved him and adorned him.

	℟. He clothed him with a robe of glory.

\properantiphon{Mag.}{O Teacher right excellent, {\dag} O light of Holy Church, O blessed Gregory, lover of the divine law: intercede for us unto the Son of God.}

\subbysub{Mattins}\label{GregoryMattins}

\begin{rubric}
	The Invitatory Hymn is as in I Evensong.
\end{rubric}

\gregorioscore{resources/gabc/ProperTime/GregoryMattins.gabc}

    ℣. The Lord guided the righteous in right paths.

	℟. And shewed him the kingdom of God.

\properantiphon{Ben.}{Gregory, {\dag} when he looked upon the youthful Angles, said: They have the countenance of Angels, and such as these should be of the fellowship of the Angels in heaven.}

\subbysub{II Evensong}

\begin{rubric}
	The Office Hymn \& Antiphon are of I Evensong, with the Versicle from Mattins.
\end{rubric}


%MANUAL ADJUSTMENT:
\clearpage
\subby{St. Patrick}
\feastday{{St. Patrick}}
\fancyhead[RE,LO]{17 March}
\begin{inhead}
    {Double\\
17 March}
\end{inhead}

\begin{rubric}
	In Lent, Commemoration only.
\end{rubric}

\begin{rubric}
	The propers are from the First Common of a Confessor Bishop (p. \pageref{CommonConfessorBishopI}), except for that which followeth.
\end{rubric}

\collect
\lett{O}{God,} who for the preaching of thy glory unto the Gentiles wast pleased to send forth blessed Patrick, thy Confessor and Bishop: grant by his merits and intercession; that we may through thy mercy be enabled to accomplish those things which thou commandest us to do. Through.

\begin{rubric}
    Commemoration is made of St. Joseph of Arimathea (p. \pageref{ArimatheanCollect}).
\end{rubric}

\begin{rubric}
	In Lent, Commemoration \& Last Gospel of the Feria.
\end{rubric}

\secret
\lett{W}{e} beseech thee, O Lord, that we remembering with gladness the merits of thy Saints, may in all places feel the succour of their intercession. Through.

\begin{rubric}
    Commemoration is made of St. Joseph of Arimathea, from the Second Common of a Confessor not a Bishop (p. \pageref{CommonConfessorNotBishopII})
\end{rubric}

\begin{rubric}
	In Lent, Commemoration \& Last Gospel of the Feria.
\end{rubric}

\postcommunion
\lett{G}{rant,} we beseech thee, almighty God: that we, shewing forth our thankfulness for the gifts which we have received, may at the intercession of blessed Patrick, thy Confessor and Bishop, obtain yet more abundant mercies. Through.

\begin{rubric}
    Commemoration is made of St. Joseph of Arimathea, from the Second Common of a Confessor not a Bishop (p. \pageref{CommonConfessorNotBishopII})
\end{rubric}

\begin{rubric}
	In Lent, Commemoration \& Last Gospel of the Feria.
\end{rubric}


%MANUAL ADJUSTMENT:
\clearpage
\subby{St. Joseph of Arimathea}
\feastday{{St. Joseph Arimathea}}
\fancyhead[RE,LO]{17 March}
\begin{inhead}
    {Memorial\\
17 March}
\end{inhead}

\begin{rubric}
	The propers are from the Second Common of a Confessor not a Bishop (p. \pageref{CommonConfessorNotBishopII}), except for that which followeth.
\end{rubric}
\begin{rubric}
	In Lent, Commemoration \& Last Gospel of the Feria.
\end{rubric}

    ℣. Then took they the body of Jesus.

	℟. And wound it in linen clothes with the spices.

\antiphon{Mag. \& Ben.}{Joseph of Arimathea, {\dag} being a disciple of Jesus, but secretly for fear of the Jews, besought Pilate that he might take away the body of Jesus: and Pilate gave him leave.}
 
\collect\label{ArimatheanCollect}
\lett{O}{God,} who didst give such grace unto thy servant Joseph that he boldly craved the body of Jesus, and with great reverence laid him in the rock-hewn sepulchre: grant, we beseech thee, that we likewise may be so emboldened for thee, as to do works meet for thy Kingdom. Through the same.


\subby{St. Cyril of Jerusalem}
\feastday{{St. Cyril Jerusalem}}
\fancyhead[RE,LO]{18 March}
\begin{inhead}
    {Double\\
18 March}
\end{inhead}

\begin{rubric}
	In Lent, Commemoration only.
\end{rubric}

\begin{rubric}
	The Daily propers are from the Common of Doctors (p. \pageref{CommonDoctors}), except for that which followeth.
\end{rubric}

\introit
\lett{I}{n} the midst of the Church he opened his mouth: and the Lord filled him with the spirit of wisdom and of understanding: he clothed him with a robe of glory. \textit{Ps.} It is a good thing to give thanks unto the Lord: and to sing praises unto thy name, O most Highest.

\collect\label{CyrilCollect}
\lett{G}{rant} to us, we beseech thee, almighty God, at the intercession of the blessed Bishop Cyril: so to know thee, the only true God, and Jesus Christ whom thou hast sent; that we may be found worthy to be numbered for evermore among the sheep who hear his voice. Through the same.

\begin{rubric}
	Commemoration is made of St. Edward (p. \pageref{EdwardCollect}).
\end{rubric}
\begin{rubric}
	In Lent, Commemoration of the Feria.
\end{rubric}

\begin{rubric}
	The Epistle is from the additional Epistle of the Common of Doctors (p. \pageref{CommonDoctors}).
\end{rubric}

\gradtr{The mouth of the righteous is exercised in wisdom, and his tongue will be talking of judgement. ℣. The law of his God is in his heart: and his goings shall not slide.}{Blessed is the man that feareth the Lord: he hath great delight in his commandments. ℣. His seed shall be mighty upon earth: the generation of the faithful shall be blessed. ℣. Riches and plenteousness shall be in his house: and his righteousness endureth for ever.}

\readingcitation{Gospel}{Matthew 10:23}
\lett{A}{t that time:} Jesus said unto his disciples: When they persecute you in this city, flee ye into another: for verily I say unto you, Ye shall not have gone over the cities of Israel, till the Son of man be come. The disciple is not above his master, nor the servant above his lord. It is enough for the disciple that he be as his master, and the servant as his lord. If they have called the master of the house Beelzebub, how much more shall they call them of his household? Fear them not therefore: for there is nothing covered, that shall not be revealed; and hid, that shall not be known. What I tell you in darkness, that speak ye in light: and what ye hear in the ear, that preach ye upon the housetops. And fear not them which kill the body, but are not able to kill the soul: but rather fear him which is able to destroy both soul and body in hell.

\offertory{The righteous shall flourish like a palm-tree: and shall spread abroad like a cedar in Libanus.}

\secret\label{CyrilSecret}
\lett{L}{ook} down, O Lord, upon the spotless victim which we offer unto thee: and grant; that by the merits of thy blessed Bishop and Confessor Cyril we may endeavour ourselves to receive it with clean hearts. Through.

\begin{rubric}
	Commemoration is made of St. Edward (p. \pageref{EdwardSecret}).
\end{rubric}
\begin{rubric}
	In Lent, Commemoration of the Feria.
\end{rubric}

\communion{A faithful and wise servant, whom the Lord hath made ruler over his household: to give them their portion of meat in due season.}

\postcommunion\label{CyrilPostcommunion}
\lett{O}{Lord} Jesu Christ, may the sacraments of thy Body and Blood, which we have received: sanctify our minds and hearts through the prayers of the blessed Bishop Cyril; that we may be worthy to be made partakers of the divine nature. Who livest.

\begin{rubric}
	Commemoration is made of St. Edward (p. \pageref{EdwardPostcommunion}).
\end{rubric}
\begin{rubric}
	In Lent, Commemoration \& Last Gospel of the Feria.
\end{rubric}


\subby{St. Edward}
\feastday{{St. Edward}}
\fancyhead[RE,LO]{18 March}
\begin{inhead}
    {Memorial\\
18 March}
\end{inhead}

\begin{rubric}
	The propers are from the First Common of a Martyr not a Bishop (p. \pageref{CommonMartyrNotBishopI}), except for the following.
\end{rubric}

\collect\label{EdwardCollect}
\lett{O}{God,} the triumphant ruler of an everlasting kingdom, mercifully behold this thy family who celebrate the memory of blessed Edward, thy King and Martyr; and, by his merits and intercession, vouchsafe; that they, who glory in his triumph, may also attain unto his rewards. Through.

\begin{rubric}
	Commemoration is made of St. Cyril (p. \pageref{CyrilCollect}).
\end{rubric}
\begin{rubric}
	In Lent, Commemoration of the Feria.
\end{rubric}

\secret\label{EdwardSecret}
\lett{G}{rant,} O Lord, that this our bounden service may be acceptable in thy sight: that these our oblations may, by the prayers of him on whose solemnity they are offered, be made profitable unto our salvation. Through.
\begin{rubric}
	Commemoration is made of St. Cyril (p. \pageref{CyrilSecret}).
\end{rubric}
\begin{rubric}
	In Lent, Commemoration of the Feria.
\end{rubric}

\postcommunion\label{EdwardPostcommunion}
\lett{W}{e} beseech thee, O Lord our God, that like as we, whom thou hast refreshed by the partaking of thy sacred gift, do offer unto thee our worship: so by the intercession of blessed Edward thy Martyr, we may perceive the benefit of the same. Through.

\begin{rubric}
	Commemoration is made of St. Cyril (p. \pageref{CyrilPostcommunion}).
\end{rubric}
\begin{rubric}
	In Lent, Commemoration \& Last Gospel of the Feria.
\end{rubric}

%MANUAL ADJUSTMENT:
\clearpage
\bcpfeast{19 March. St. Joseph, Spouse of the Blessed Virgin Mary}{St. Joseph Spouse}{19 March}
%\supplement{19 March}{St. Joseph}{Spouse of the B.V.M.}

\subbysub{I Evensong}\label{JosephEvensong}

\gregorioscore{resources/gabc/ProperTime/JosephEvensong.gabc}

    ℣. He made him lord of his house.

	℟. And ruler of all his substance.
	
	\properantiphon{Mag.}{Then Joseph, {\dag} being raised from sleep, did as the Angel of the Lord had bidden him, and took unto him his wife.}
	
\subbysub{Mattins}
	
	\invitatoryhymn\label{JosephInvitatory}
	
\gregorioscore{resources/gabc/ProperTime/JosephInvitatory.gabc}
	
	\officehymn\label{JosephMattins}
	
\gregorioscore{resources/gabc/ProperTime/JosephMattins.gabc}

    ℣. The mouth of the righteous is exercised in wisdom.

	℟. And his tongue will be talking of judgement.

	\properantiphon{Ben.}{Jesus himself {\dag} began to be about thirty years of age, being, as was supposed, the son of Joseph.}

\subbysub{II Evensong}

\begin{rubric}
	The Office Hymn is of I Evensong, with the following Versicle \& Antiphon.
\end{rubric}

    ℣. Riches and plenteousness shall be in his house.

	℟. And his righteousness endureth for ever.
	
	\properantiphon{Mag.}{Behold a faithful {\dag} and wise servant, whom his Lord hath made ruler over his household.}
	

\subby{St. Cuthbert}
\feastday{{St. Cuthbert}}
\fancyhead[RE,LO]{20 March}
\begin{inhead}
    {Double\\
20 March}
\end{inhead}

\begin{rubric}
	In Lent, Commemoration only.
\end{rubric}

\begin{rubric}
	The propers are from the Second Common of a Bishop Confessor (p. \pageref{CommonConfessorBishopII}), except for that which followeth.
\end{rubric}

\collect
\lett{O}{God,} who dost make thy Saints glorious by the inestimable gift of thy grace: grant, we beseech thee; that at the intercession of blessed Cuthbert, thy Confessor and Bishop, we may be found worthy to attain to the perfection of all virtue. Through.

\begin{rubric}
	In Lent, Commemoration of the Feria.
\end{rubric}

\secret
\lett{A}{ccept,} we beseech thee, O Lord, the sacrifice of man's redemption: and at the intercession of blessed Cuthbert, thy Confessor and Bishop, mercifully grant us health of mind and of body. Through.

\begin{rubric}
	In Lent, Commemoration of the Feria.
\end{rubric}

\postcommunion
\lett{W}{e} beseech thee, O Lord, that thy holy things which we have received may protect us by their power: and at the intercession of blessed Cuthbert, thy Confessor and Bishop, whose life shone forth in glory, guard us in peace and holiness. Through.

\begin{rubric}
	In Lent, Commemoration \& Last Gospel of the Feria.
\end{rubric}

\bcpfeast{21 March. St. Benedict}{St. Benedict}{21 March}
%\supplement{21 March}{St. Benedict}{}

\subbysub{I Evensong}

	\gregorioscore{resources/gabc/ProperTime/BenedictEvensong.gabc}\label{BenedictEvensong}

    ℣. The Lord loved him and adorned him.

	℟. He clothed him with a robe of glory.
	
	\properantiphon{Mag.}{Let all the multitude {\dag} of the faithful exult in the glory of the gracious Father Benedict: and chiefly let the company of monks be joyful, celebrating his festival on earth, in whose goodly fellowship the Saints rejoice in heaven.}

\subbysub{Mattins}
	
	\invitatoryhymn
	
	\gregorioscore{resources/gabc/ProperTime/BenedictMattins.gabc}\label{BenedictInvitatory}
	
	\officehymn

	\gregorioscore{resources/gabc/ProperTime/BenedictMattins.gabc}\label{BenedictMattins}

    ℣. The Lord guided the righteous in right paths.

	℟. And shewed him the kingdom of God.
	
	\properantiphon{Ben.}{Benedict, {\dag} thou father and guide of monks, thou most holy Confessor of the Lord, intercede for us all and for our salvation.}

\subbysub{II Evensong}

\begin{rubric}
	The Office Hymn is of I Evensong, with the Versicle from Mattins and the following Antiphon.
\end{rubric}

\properantiphon{Mag.}{To-day {\dag} holy Benedict, while his disciples beheld it, ascending by way of the East on a straight pathway to heaven: to-day with hands uplifted, he died between the words of his supplication: to-day he was received into glory by the Angels.}


\subby{St. Gabriel}
\feastday{St. Gabriel Archangel}
\fancyhead[RE,LO]{24 March}
\begin{inhead}
    {Greater Double\\
24 March}
\end{inhead}
\par\noindent

\begin{rubric}
	Commemoration only is made, unless a Votive Mass of St. Gabriel is said outside of Lent or as a patronal Feast Day.
\end{rubric}

\begin{rubric}
	I Evensong: Psalms 9 \& 11,15. Daniel 8:15-19. Revelation 8:1-5.\par
	Mattins: Psalms 30, 34, \& 77. Daniel 9:20-26. Luke 1:5-17\par
	II Evensong: Psalms 97, 99, \& 103. Isaiah 7:10-14. Luke 1:26-38.
\end{rubric}

%AOB:
\textit{Opening Sentence.} I saw another angel come down from heaven, having great power; and the earth was lightened with his glory.

%MANUAL ADJUSTMENT:
\clearpage

\begin{paracol}{2}[]
\sloppy
\begin{inhead}
	I Evensong
\end{inhead}
\begin{hangparas}{1.25em}{1}

Christ, the fair glory of the holy Angels,

Thou who hast made us, thou who o'er us rulest,

Grant of thy mercy unto us thy servants

Steps up to heaven.\\

Send thy Archangel, Gabriel, the mighty,

Herald of heaven; may he from us mortals

Spurn the old serpent, watching o'er the temples

Where thou art worshipp'd.\\

May the blest Mother of our God and Saviour,

May the assembly of the Saints in glory,

May the celestial companies of Angels

Ever assist us.\\

This he vouchsafe us, God for ever blessed

Father eternal, Son, and Holy Spirit,

Whose is the glory which through all creation.

Ever resoundeth. Amen.\\
\end{hangparas}

    ℣. An Angel stood at the altar of the temple.

	℟. Having in his hand a golden censer.
	
	\properantiphon{Mag.}{The Angel Gabriel {\dag} came in unto Mary and said, Hail, thou that art full of grace, the Lord is with thee; blessed art thou among women.}
	
	\switchcolumn
	
\begin{inhead}
	Mattins
\end{inhead}
\begin{hangparas}{1.25em}{1}
O Christ, Redeemer of us all,

Protect thy servants when they call,

And hear with reconciling care

The blessed Virgin's holy prayer.\\

Be ever present, Angel high

Whose name `God's might' doth signify:

To all the weak new strength impart,

And solace to the sad of heart.\\

And ye, O ever blissful throng

Of heav'nly Spirits, guardians strong,

Our past and present ills dispel,

From future peril shield us well.\\

From lands wherein thy faithful dwell

Drive far away the infidel;

So we to Christ due hymns of praise

Henceforth with eager hearts may raise.\\

To thee, O Father, born of none,

And thee, O sole-begotten Son,

One with the Holy Paraclete,

Be glory ever, as is meet. Amen.\\
\end{hangparas}

    ℣. An Angel stood at the altar of the temple.

	℟. Having in his hand a golden censer.
	
	\properantiphon{Ben.}{The Angel Gabriel {\dag} descended to Zacharias, and said unto him: Thy wife shall bear thee a son, and thou shalt call his name John; and many shall rejoice at his birth; for he shall go before the face of the Lord, to prepare his ways.}

\fussy
\end{paracol}

\begin{rubric}
	In II Evensong, the Office Hymn is of I Evensong, with the following Versicle \& Antiphon.
\end{rubric}

    ℣. In the presence of the Angels I will sing praise unto thee, O my God.

	℟. I will worship toward thy holy temple, and praise thy Name.
	
	\properantiphon{Mag.}{The Archangel Gabriel {\dag} said unto Mary: With God nothing shall be impossible. And Mary said, Behold the handmaid of the Lord: be it unto me according to thy word. And the Angel departed from her.}

\introit
\lett{O}{praise} the Lord, ye Angels of his, ye that excel in strength: ye that fulfil his commandment, and hearken unto the voice of his words. \textit{Ps.} Praise the Lord, O my soul: and all that is within me praise his holy name.

\collect
\lett{O}{God,} who from the company of Angels didst choose the Archangel Gabriel to proclaim the mystery of thine Incarnation: mercifully grant; that we who celebrate his festival (commemoration) on earth, may perceive his advocacy in heaven. Who livest.

\begin{rubric}
	Commemoration is made of the Feria.
\end{rubric}

\readingcitation{Epistle}{Daniel 9:21}
\lett{I}{n those days:} Behold the man Gabriel, whom I had seen in the vision at the beginning, being caused to fly swiftly, touched me about the time of the evening oblation. And he informed me, and talked with me, and said, O Daniel, I am now come forth to give thee skill and understanding. At the beginning of thy supplications the commandment came forth, and I am come to shew thee; for thou art greatly beloved: therefore understand the matter, and consider the vision. Seventy weeks are determined upon thy people and upon thy holy city, to finish the transgression, and to make an end of sins, and to make reconciliation for iniquity, and to bring in everlasting righteousness, and to seal up the vision and prophecy, and to anoint the most Holy. Know therefore and understand, that from the going forth of the commandment to restore and to build Jerusalem unto the Messiah the Prince shall be seven weeks, and threescore and two weeks: the street shall be built again, and the wall, even in troublous times. And after threescore and two weeks shall Messiah be cut off, but not for himself: and the people of the prince that shall come shall destroy the city and the sanctuary; and the end thereof shall be with a flood, and unto the end of the war desolations are determined.

\gradtr{O praise the Lord, ye Angels of his, ye that excel in strength, ye that fulfil his commandment. ℣. Praise the Lord, O my soul, and all that is within me praise his holy name.}{Hail, Mary, full of grace; the Lord is with thee. ℣. Blessed art thou among women: and blessed is the fruit of thy womb. ℣. Behold, thou shalt conceive, and bring forth a Son, and shalt call his name Emmanuel. ℣. The Holy Ghost shall come upon thee, and the power of the Highest shall overshadow thee. ℣. Therefore also that Holy Thing which shall be born of thee, shall be called the Son of God.}

\begin{rubric}
	In Votive Masses before Septuagesima or after Pentecost, the Gradual as above, but in place of the Tract is said:
\end{rubric}

\alleluia{Alleluia, alleluia. ℣. O praise the Lord, all ye his hosts: ye servants of his that do his pleasure. Alleluia.}

\begin{rubric}
	In Votive Masses in Eastertide, replacing the Gradual \& Tract:
\end{rubric}

\alleluia{Alleluia, alleluia. ℣. Who maketh his Angels spirits: and his ministers a flaming fire. Alleluia. ℣. Hail, Mary, full of grace; the Lord is with thee: blessed art thou among women. Alleluia.}

\readingcitation{Gospel}{Luke 1:26}
\lett{A}{t that time:} The Angel Gabriel was sent from God unto a city of Galilee named Nazareth, to a virgin espoused to a man whose name was Joseph, of the house of David; and the virgin's name was Mary. And the angel came in unto her, and said, Hail, thou that art highly favoured, the Lord is with thee: blessed art thou among women. And when she saw him, she was troubled at his saying, and cast in her mind what manner of salutation this should be. And the angel said unto her, Fear not, Mary: for thou hast found favour with God. And, behold, thou shalt conceive in thy womb, and bring forth a son, and shalt call his name \divineName{Jesus}. He shall be great, and shall be called the Son of the Highest: and the Lord God shall give unto him the throne of his father David: and he shall reign over the house of Jacob for ever; and of his kingdom there shall be no end. Then said Mary unto the angel, How shall this be, seeing I know not a man? And the angel answered and said unto her, The Holy Ghost shall come upon thee, and the power of the Highest shall overshadow thee: therefore also that holy thing which shall be born of thee shall be called the Son of God. And, behold, thy cousin Elisabeth, she hath also conceived a son in her old age: and this is the sixth month with her, who was called barren. For with God nothing shall be impossible. And Mary said, Behold the handmaid of the Lord; be it unto me according to thy word.

\offertory{An Angel stood at the altar of the temple, having a golden censer in his hand; and there was given unto him much incense: and the smoke of the incense ascended up before God.}

\secret
\lett{M}{ay} the offering of our service, and the prayer of blessed Gabriel the Archangel, be accepted in thy sight, O Lord: that he, who is venerated by us on earth, may be an advocate for us with thee in heaven. Through.

\begin{rubric}
	Commemoration of the Feria.
\end{rubric}

\communion{O all ye Angels of the Lord, bless ye the Lord: sing ye praises, and magnify him above all for ever.}

\postcommunion
\lett{H}{aving} received the mysteries of thy Body and Blood, we entreat thy loving kindness, O Lord our God: that, as we have known thine Incarnation by the message of Gabriel, so we may through his help obtain the benefits of that Incarnation. Who livest.

\begin{rubric}
	Commemoration \& Last Gospel of the Feria.
\end{rubric}


\bcpfeast{25 March. Annunciation of the Blessed Virgin Mary}{Annunciation}{25 March}
%\supplement{25 March}{Annunciation}{of the Blessed Virgin Mary}

\begin{secrubric}
	The Hymns are from the Common of the Blessed Virgin Mary (p. \pageref{CommonBVM}), with the following Versicles \& Antiphons.
\end{secrubric}

    ℣. Hail Mary, full of grace.

	℟. The Lord is with thee.
	
	\properantiphon{Mag.}{The Holy Ghost {\dag} shall come upon thee, Mary: and the power of the Highest shall overshadow thee.}\\
	
	℣. Hail Mary, full of grace.

	℟. The Lord is with thee.

	\properantiphon{Ben.}{How shall this be, {\dag} O Angel of God, seeing I know not a man? Hearken, O Virgin Mary: The Holy Ghost shall come upon thee, and the power of the Highest shall overshadow thee.}\\
	
	℣. Hail Mary, full of grace.

	℟. The Lord is with thee.

	\properantiphon{Mag.}{The Angel Gabriel {\dag} spake unto Mary, saying: Hail, thou that art full of grace, the Lord is with thee; blessed art thou among women.}


\subby{St. John Damascene}
\feastday{{St. John Damascene}}
\fancyhead[RE,LO]{27 March}
\begin{inhead}
    {Double\\
27 March}
\end{inhead}

\begin{rubric}
	Commemoration only.
\end{rubric}

\begin{rubric}
	The Daily Office propers are from the Common of Doctors (p. \pageref{CommonDoctors}), except for that which followeth.
\end{rubric}

\introit
\lett{T}{hou} hast holden me by my right hand: thou shalt guide me with thy counsel, and after that receive me with glory. \textit{Ps.} Truly God is loving unto Israel, even unto such as are of a clean heart!

\collect
\lett{A}{lmighty} and everlasting God, who, for the defence of the veneration of sacred images, didst endue blessed John with heavenly doctrine and wondrous strength of spirit: grant unto us by his intercession and example; that we may imitate the virtues and perceive the advocacy of those whose images we venerate. Through.

\begin{rubric}
	Commemoration of the Feria.
\end{rubric}

\readingcitation{Epistle}{Wisdom 10:10}
%RV (unedited):
\lett{W}{isdom} guided him in straight paths; She shewed him God’s kingdom, and gave him knowledge of holy things; She prospered him in his toils, and multiplied the fruits of his labour; When in their covetousness men dealt hardly with him, She stood by him and made him rich; She guarded him from enemies, And from those that lay in wait she kept him safe, And over his sore conflict she watched as judge, That he might know that godliness is more powerful than all. When a righteous man was sold, wisdom forsook him not, But from sin she delivered him; She went down with him into a dungeon, And in bonds she left him not, Till she brought him the sceptre of a kingdom, And authority over those that dealt tyrannously with him; She shewed them also to be false that had mockingly accused him, And gave him eternal glory. Wisdom delivered a holy people and a blameless seed from a nation of oppressors. She entered into the soul of a servant of the Lord, And withstood terrible kings in wonders and signs. She rendered unto holy men a reward of their toils.

\gradtr{It is God that girdeth me with strength of war: and maketh my way perfect, ℣. He teacheth mine hands to fight: and mine arms shall break even a bow of steel.}{I will follow upon mine enemies, and overtake them. ℣. I will smite them that they shall not be able to stand: but fall under my feet. ℣. For this cause will I give thanks unto thee, O Lord, among the Gentiles, and sing praises unto thy name.}

\readingcitation{Gospel}{Luke 6:6}
\lett{A}{t that time:} It came to pass on another sabbath, that Jesus entered into the synagogue and taught. And there was a man whose right hand was withered. And the scribes and Pharisees watched him, whether he would heal on the sabbath day; that they might find an accusation against him. But he knew their thoughts, and said to the man which had the withered hand, Rise up, and stand forth in the midst. And he arose and stood forth. Then said Jesus unto them, I will ask you one thing; Is it lawful on the sabbath days to do good, or to do evil? to save life, or to destroy it? And looking round about upon them all, he said unto the man, Stretch forth thy hand. And he did so: and his hand was restored whole as the other. And they were filled with madness; and communed one with another what they might do to Jesus.

\offertory{There is hope of a tree, if it be cut down, that it will sprout again, and that the tender branch thereof will not cease.}

\secret
\lett{O}{Lord,} let the devout intercession of blessed John and of the Saints, who through his labours are set forth in the temples for our veneration, avail to render the gifts which we offer acceptable in thy sight. Through.

\begin{rubric}
	Commemoration of the Feria.
\end{rubric}

\communion{The arms of the ungodly shall be broken, and the Lord upholdeth the righteous.}

\postcommunion
\lett{W}{e} beseech thee, O Lord, that the gifts which we have received may shield us with heavenly armour: and that the advocacy of blessed John, together with the united intercession of the Saints, the veneration of whose images in the Church be victoriously upheld, be our defence. Through.

\begin{rubric}
	Commemoration \& Last Gospel of the Feria.
\end{rubric}


\subby{St. John Climacus}
\feastday{{St. John Climacus}}
\fancyhead[RE,LO]{30 March}
\begin{inhead}
    {Memorial\\
30 March}
\end{inhead}

\begin{rubric}
	The propers are from the Common of Abbots (p. \pageref{CommonAbbots}).
\end{rubric}
\begin{rubric}
	Commemoration \& Last Gospel of the Feria.
\end{rubric}


\subby{St. Innocent of Alaska}
\feastday{{St. Innocent Alaska}}
\fancyhead[RE,LO]{31 March}
\begin{inhead}
    {Memorial\\
31 March}
\end{inhead}

\begin{rubric}
	The propers are from the Second Common of a Confessor Bishop (p. \pageref{CommonConfessorBishopII}).
\end{rubric}
\begin{rubric}
	Commemoration \& Last Gospel of the Feria.
\end{rubric}

%MANUAL ADJUSTMENT:
\clearpage

\subby{St. Isidore of Seville}
\feastday{{St. Isidore Seville}}
\fancyhead[RE,LO]{4 April}
\begin{inhead}
    {Double\\
4 April}
\end{inhead}

\begin{rubric}
	In Lent or Easter Week, Commemoration only.
\end{rubric}

\begin{rubric}
	The propers are from the Common of Doctors (p. \pageref{CommonDoctors}).
\end{rubric}
\begin{rubric}
	In Lent, Commemoration \& Last Gospel of the Feria.
\end{rubric}


\subby{St. Tikhon of Moscow}
\feastday{{St. Tikhon Moscow}}
\fancyhead[RE,LO]{7 April}
\begin{inhead}
    {\nth{1} Double, Common Octave\\
7 April}
\end{inhead}

\begin{rubric}
	The propers are from the Second Common of a Bishop Confessor (p. \pageref{CommonConfessorBishopII}).
\end{rubric}
\begin{rubric}
	In Lent, Commemoration \& Last Gospel of the Feria.
\end{rubric}


\subby{St. Leo the Great}
\feastday{{Pope St. Leo the Great}}
\fancyhead[RE,LO]{11 April}
\begin{inhead}
    {Double\\
11 April}
\end{inhead}

\begin{rubric}
	In Lent or Easter Week, Commemoration only.
\end{rubric}

\begin{rubric}
	The Daily Office propers are from the Common of Doctors (p. \pageref{CommonDoctors}).
\end{rubric}

\introit
\lett{I}{n} the midst of the Church he opened his mouth: and the Lord filled him with the spirit of wisdom and of understanding: he clothed him with a robe of glory. (Alleluia, alleluia.) \textit{Ps.} It is  a good thing to give thanks unto the Lord: and to sing praises unto thy name, O Most Highest.

\collect
\lett{W}{e} beseech thee, O Lord, graciously to hear the prayers which we offer unto thee on the solemnity of blessed Leo, thy Confessor and Bishop: that, like as he was found worthy to do thee faithful service, so by his merits and intercession we may be absolved from all our sins. Through.

\begin{rubric}
	In Lent, Commemoration of the Feria.
\end{rubric}

\begin{rubric}
	The Epistle is the additional Epistle of the Common of Doctors (p. \pageref{Ecclesiasticus395}).
\end{rubric}

\gradtr{The mouth of the righteous is exercised in wisdom, and his tongue will be talking of judgement. ℣. The law of his God is in his heart: and his goings shall not slide.}{Blessed is the man that feareth the Lord: he hath great delight in his commandments. ℣. His seed shall be mighty upon earth: the generation of the faithful shall be blessed. ℣. Riches and plenteousness shall be in his house: and his righteousness endureth for ever.}

\begin{rubric}
	In Eastertide, the following Alleluia Verse replaces the Gradual \& Tract.
\end{rubric}

\alleluia{Alleluia, alleluia. ℣. The Lord loved him, and adorned him: he clothed him with a robe of glory. Alleluia. ℣. The righteous shall grow as the lily: and flourish for ever before the Lord. Alleluia.}

\readingcitation{Gospel}{Matthew 16:13}
\lett{A}{t that time:} When Jesus came into the coasts of Caesarea Philippi, he asked his disciples, saying, Whom do men say that I the Son of man am? And they said, Some say that thou art John the Baptist: some, Elias; and others, Jeremias, or one of the prophets. He saith unto them, But whom say ye that I am? And Simon Peter answered and said, Thou art the Christ, the Son of the living God. And Jesus answered and said unto him, Blessed art thou, Simon Barjona: for flesh and blood hath not revealed it unto thee, but my Father which is in heaven. And I say also unto thee, That thou art Peter, and upon this rock I will build my church; and the gates of hell shall not prevail against it. And I will give unto thee the keys of the kingdom of heaven: and whatsoever thou shalt bind on earth shall be bound in heaven: and whatsoever thou shalt loose on earth shall be loosed in heaven.

\offertory{I have found David my servant, with my holy oil have I anointed him: my hand shall hold him fast, and my arm shall strengthen him. (Alleluia.)}


\secret
\lett{W}{e} beseech thee, O Lord, that our devout observance of the yearly solemnity of Saint Leo, thy Confessor and Bishop, may render us acceptable unto thy loving kindness: that this service of propitiation, which we duly offer, may be profitable unto him for the reward of blessed-ness, and obtain for us the gifts of thy grace. Through.

\begin{rubric}
	In Lent, Commemoration of the Feria.
\end{rubric}

\communion{Blessed is the servant, whom the lord when he cometh shall find watching: verily I say unto you, that he shall make him ruler over all his goods. (Alleluia.)}

\postcommunion
\lett{O}{God,} who rewardest the souls of them that put their trust in thee: vouchsafe; that we who keep the solemn festival of blessed Leo, thy Confessor and Bishop, may by his prayers obtain thy merciful pardon. Through.

\begin{rubric}
	In Lent, Commemoration \& Last Gospel of the Feria.
\end{rubric}


\subby{St. Hermengild}
\feastday{{St. Hermengild}}
\fancyhead[RE,LO]{13 April}
\begin{inhead}
    {Memorial\\
13 April}
\end{inhead}

\begin{rubric}
	In Lent, the propers from the First Common of a Martyr not a Bishop (p. \pageref{CommonMartyrNotBishopI}), with the following. And Commemoration \& Last Gospel of the Feria.
\end{rubric}

\begin{rubric}
	In Eastertide, the propers from the Common of Martyrs in Eastertide (p. \pageref{CommonMartyrsEaster}), with the following.
\end{rubric}

\collect
\lett{O}{God,} who didst teach blessed Hermengild, thy Martyr, to lay down an earthly for a heavenly kingdom: grant to us, we beseech thee; by his example to despise things temporal, and to seek after things eternal. Through.

\readingcitation{Gospel}{Luke 14:26}
\lett{A}{t that time:} Jesus said unto the multitudes: If any man come to me, and hate not his father, and mother, and wife, and children, and brethren, and sisters, yea, and his own life also, he cannot be my disciple. And whosoever doth not bear his cross, and come after me, cannot be my disciple. For which of you, intending to build a tower, sitteth not down first, and counteth the cost, whether he have sufficient to finish it? Lest haply, after he hath laid the foundation, and is not able to finish it, all that behold it begin to mock him, Saying, This man began to build, and was not able to finish. Or what king, going to make war against another king, sitteth not down first, and consulteth whether he be able with ten thousand to meet him that cometh against him with twenty thousand? Or else, while the other is yet a great way off, he sendeth an ambassage, and desireth conditions of peace. So likewise, whosoever he be of you that forsaketh not all that he hath, he cannot be my disciple.

\secret
\lett{W}{e} beseech thee, O Lord, to accept our prayers and oblations: and graciously hearken unto us, whom thou dost cleanse by thy heavenly mysteries. Through.

\postcommunion
\lett{G}{rant,} we beseech thee, O Lord our God: that like as we in this life do gladly honour the memory of thy Saints; so we may rejoice to behold them for ever. Through.


\subby{St. Justin}
\feastday{{St. Justin}}
\fancyhead[RE,LO]{14 April}
\begin{inhead}
    {Memorial\\
14 April}
\end{inhead}

\introit
\lett{T}{he} proud have digged pits for me, which are not after thy law: I will speak of thy testimonies also, even before kings, and will not be ashamed. (Alleluia, alleluia.) \textit{Ps.} Blessed are those that are undefiled in the way: and walk in the law of the Lord.

\collect
\lett{O}{God,} who through the foolishness of the Cross didst wondrously teach blessed Justin Martyr the excellent knowledge of Jesus Christ: grant to us through his intercession; that we, driving away the errors that beset us, may attain unto steadfastness of faith. Through the same.

\begin{rubric}
	Commemoration is made of St. Tiburtius, Valerian, \& Maximus (p. \pageref{TiburtiusCollect}).
\end{rubric}
\begin{rubric}
	In Lent, Commemoration of the Feria.
\end{rubric}

\readingcitation{Epistle}{1 Corinthians 1:18}
\lett{B}{rethren:} The preaching of the cross is to them that perish foolishness; but unto us which are saved it is the power of God. For it is written, I will destroy the wisdom of the wise, and will bring to nothing the understanding of the prudent. Where is the wise? where is the scribe? where is the disputer of this world? hath not God made foolish the wisdom of this world? For after that in the wisdom of God the world by wisdom knew not God, it pleased God by the foolishness of preaching to save them that believe. For the Jews require a sign, and the Greeks seek after wisdom: But we preach Christ crucified, unto the Jews a stumblingblock, and unto the Greeks foolishness; But unto them which are called, both Jews and Greeks, Christ the power of God, and the wisdom of God. Because the foolishness of God is wiser than men; and the weakness of God is stronger than men. For ye see your calling, brethren, how that not many wise men after the flesh, not many mighty, not many noble, are called: But God hath chosen the foolish things of the world to confound the wise; and God hath chosen the weak things of the world to confound the things which are mighty; And base things of the world, and things which are despised, hath God chosen, yea, and things which are not, to bring to nought things that are: That no flesh should glory in his presence. But of him are ye in Christ Jesus, who of God is made unto us wisdom, and righteousness, and sanctification, and redemption.

\gradtr{The wisdom of this world is foolishness with God, for it is written: The Lord knoweth the thoughts of the wise, that they are vain. ℣. I will destroy the wisdom of the wise, and will bring to nothing the understanding of the prudent.}{I determined not to know any thing among you, save Jesus Christ, and him crucified. ℣. We speak the wisdom of God in a mystery, even the hidden wisdom, which God ordained before the world unto our glory. ℣. Which none of the princes of this world knew. For had they known it, they would not have crucified the Lord of glory.}

\alleluia{Alleluia, alleluia. ℣. The wisdom of this world is foolishness with God: for it is written: The Lord knoweth the thoughts of the wise, that they are vain. Alleluia. ℣. Yea doubtless, and I count all things but loss for the excellency of the knowledge of Christ Jesus my Lord. Alleluia.}

\begin{rubric}
	In Votive Masses before Septuagesima or after Pentecost, the Gradual is said, as above, but, the Tract being omitted, there is added
\end{rubric}

\alleluia{Alleluia, alleluia. ℣. Yea doubtless, and I count all things but loss for the excellency of the knowledge of Christ Jesus my Lord. Alleluia.}

\readingcitation{Gospel}{Luke 12:2}
\lett{A}{t that time:} Jesus said to his disciples: There is nothing covered, that shall not be revealed; neither hid, that shall not be known. Therefore whatsoever ye have spoken in darkness shall be heard in the light; and that which ye have spoken in the ear in closets shall be proclaimed upon the housetops. And I say unto you my friends, Be not afraid of them that kill the body, and after that have no more that they can do. But I will forewarn you whom ye shall fear: Fear him, which after he hath killed hath power to cast into hell; yea, I say unto you, Fear him. Are not five sparrows sold for two farthings, and not one of them is forgotten before God? But even the very hairs of your head are all numbered. Fear not therefore: ye are of more value than many sparrows. Also I say unto you, Whosoever shall confess me before men, him shall the Son of man also confess before the angels of God.


\secret
\lett{O}{Lord} God, graciously receive our gifts: the wondrous mystery whereof the holy Martyr Justin did manfully defend against the slanders of the ungodly. Through.

\begin{rubric}
	Commemoration is made of St. Tiburtius, Valerian, \& Maximus (p. \pageref{TiburtiusSecret}).
\end{rubric}
\begin{rubric}
	In Lent, Commemoration of the Feria.
\end{rubric}

\communion{There is laid up for me a crown of righteousness, which the Lord, the righteous judge, shall give me at that day. (Alleluia.)}

\postcommunion
\lett{O}{Lord,} who hast refreshed us with heavenly food, we humbly beseech thee: that following the counsels of blessed Justin, thy Martyr; we may ever continue in thanksgiving for the gifts which we have received. Through.

\begin{rubric}
	Commemoration is made of St. Tiburtius, Valerian, \& Maximus (p. \pageref{TiburtiusPostcommunion}).
\end{rubric}
\begin{rubric}
	In Lent, Commemoration \& Last Gospel of the Feria.
\end{rubric}


\subby{St. Tiburtius, Valerian, and Maximus}
\feastday{{Sts. Tiburtius \&c.}}
\fancyhead[RE,LO]{14 April}
\begin{inhead}
    {Memorial\\
14 April}
\end{inhead}

\begin{rubric}
	Outside of Eastertide, the propers are from the Second Common of Many Martyrs (p. \pageref{CommonMartyrsII}), with the Prayers following.
\end{rubric}

\begin{rubric}
	In Eastertide, the propers are from the Common of Many Martyrs in Eastertide (p. \pageref{CommonMartyrsEaster}), with the Epistle and Gospel from the Common of a Martyr in Eastertide (p. \pageref{CommonMartyrEaster}), with the Prayers following.
\end{rubric}

\collect\label{TiburtiusCollect}
\lett{G}{rant,} we beseech thee, almighty God: that we, almighty God: that we, who celebrate the festival of thy holy Martyrs, Tiburtius, Valerian, and Maximus; may also imitate their virtues. Through.

\secret\label{TiburtiusSecret}
\lett{W}{e} beseech thee, O Lord, that this oblation, which we offer unto thee in commemoration of the birthday of thy holy Martyrs: may both loose the bonds of our iniquity, and obtain for us the gifts of thy mercy. Through.

\postcommunion\label{TiburtiusPostcommunion}
\lett{O}{Lord,} who hast satisfied us with this sacred gift, we humbly beseech thee: that the mysteries, which we celebrate in this service of our bounden duty, we may know to be an increase of thy salvation. Though.


\subby{St. Anicetus}
\feastday{{St. Anicetus}}
\fancyhead[RE,LO]{17 April}
\begin{inhead}
    {Memorial\\
17 April}
\end{inhead}


\begin{rubric}
	Outside of Eastertide, the propers are from the Second Common of a Martyr Bishop outside of Eastertide (p. \pageref{CommonMartyrBishopII}).
\end{rubric}

\begin{rubric}
	In Eastertide, the propers are from the Common of a Martyr in Eastertide (p. \pageref{CommonMartyrEaster}), using the Second Prayers for a Martyr Bishop, with the additional Gospel from the Common of Many Martyrs in Eastertide (p. \pageref{John1620}).
\end{rubric}


\subby{Sts. Soter and Caius}
\feastday{{Sts. Soter \& Caius}}
\fancyhead[RE,LO]{22 April}
\begin{inhead}
    {Memorial\\
22 April}
\end{inhead}

\begin{rubric}
	Outside of Eastertide, the propers are from the First Common of Many Martyrs outside of Eastertide (p. \pageref{CommonMartyrsI}).
\end{rubric}

\begin{rubric}
	In Eastertide, the propers are from the Common of Many Martyrs in Eastertide (p. \pageref{CommonMartyrsEaster}), with the additional Epistle (p. \pageref{Revelation191}).
\end{rubric}

\bcpfeast{23 April. St. George}{St. George}{23 April}
%\supplement{23 April}{St. George}{}

\begin{secrubric}
	Outside of Eastertide, the Daily Office propers are from the First Common of a Martyr not a Bishop out of Eastertide (p. \pageref{CommonMartyrNotBishopI}).
\end{secrubric}

\begin{rubric}
	In Eastertide, the Daily Office propers are from the Common of a Martyr in Eastertide (p. \pageref{CommonMartyrEaster}).
\end{rubric}


\bcpfeast{25 April. St. Mark}{St. Mark}{25 April}
%\supplement{25 April}{St. Mark}{}

\begin{secrubric}
	Outside of Eastertide, the Daily Office propers are from the Common of Apostles (p. \pageref{CommonApostles}).
\end{secrubric}

\begin{rubric}
	In Eastertide, the Daily Office propers are from the Common of Apostles in Eastertide (p. \pageref{CommonApostlesEaster}).
\end{rubric}


\subby{Sts. Cletus and Marcellinus}
\feastday{{Sts. Cletus \& Marcellinus}}
\fancyhead[RE,LO]{26 April}
\begin{inhead}
    {Memorial\\
26 April}
\end{inhead}

\begin{rubric}
	Outside of Eastertide, propers are from the First Common of Many Martyrs out of Eastertide (p. \pageref{CommonMartyrsI}), with the following Prayers.
\end{rubric}

\begin{rubric}
	In Eastertide, the propers are from the Common of Many Martyrs in Eastertide (p. \pageref{CommonMartyrsEaster}), with the following Prayers.
\end{rubric}

\collect
\lett{O}{Lord,} let the precious confession of thy blessed Martyrs and Bishops Cletus and Marcellinus be our defence: and let their loving intercession be a continual defence. Through.

\secret
\lett{A}{ssist} us mercifully, O Lord, in these our supplications, which we make before thee in remembrance of thy Saints: that we who trust not in our own righteousness may be succoured by the merits of them that have found favour in thy sight. Through.

\postcommunion
\lett{O}{Lord,} who hast fulfilled us with saving mysteries: we beseech thee that we may be aided by the prayers of those whose festival we celebrate. Through.


\subby{St. Vitalis}
\feastday{{St. Vitalis}}
\fancyhead[RE,LO]{28 April}
\begin{inhead}
    {Memorial\\
28 April}
\end{inhead}

\begin{rubric}
	Outside of Eastertide, propers are from the First Common of a Martyr not a Bishop out of Eastertide (p. \pageref{CommonMartyrNotBishopI}).
\end{rubric}

\begin{rubric}
	In Eastertide, the propers are from the Common of a Martyr in Eastertide (p. \pageref{CommonMartyrEaster}), using the first Prayers for a Martyr not a Bishop
\end{rubric}

\bcpfeast{1 May. Sts. Philip \& James}{Sts. Philip \& James}{1 May}
%\supplement{1 May}{Sts. Philip \& James}{}

\begin{secrubric}
	The Daily Office propers are from the Common of Apostles in Eastertide (p. \pageref{CommonApostlesEaster}), except for the following.
\end{secrubric}

\properantiphon{Mag.}{Let not your heart {\dag} be troubled, neither let it be afraid; ye believe in God, believe also in me: in my Father's house are many mansions, alleluia, alleluia.}\\

\properantiphon{Ben.}{I am the way, {\dag} the truth, and the life: no man cometh unto the Father, but by me, alleluia.}\\

	℣. Right dear in the sight of the Lord, alleluia.

	℟. Is the death of his Saints, alleluia.

\properantiphon{Mag.}{If ye abide in me, {\dag} and my words abide in you, ye shall ask what ye will, and it shall be done unto you, alleluia, alleluia, alleluia.}


\subby{St. Athanasius}
\feastday{{St. Athanasius}}
\fancyhead[RE,LO]{2 May}
\begin{inhead}
    {Double\\
2 May}
\end{inhead}

\begin{rubric}
	The Daily Office propers are from the Common of Doctors (p. \pageref{CommonDoctors}).
\end{rubric}

\introit
\lett{I}{n} the midst of the Church he opened his mouth: and the Lord filled him with the spirit of wisdom and of understanding: he clothed him with a robe of glory, alleluia, alleluia. \textit{Ps.} It is a good thing to give thanks unto the Lord: and to sing praise unto thy name, O Most Highest.

\collect
\lett{W}{e} beseech thee, O Lord, graciously to hear the prayers which we offer unto thee on the solemnity of blessed Athanasius, thy Confessor and Bishop: that, like as he was found worthy to do thee faithful service, so by his merits and intercession we may be absolved from all our sins. Through.

\readingcitation{Epistle}{2 Corinthians 4:5}
\lett{B}{rethren:} We preach not ourselves, but Christ Jesus the Lord; and ourselves your servants for Jesus' sake. For God, who commanded the light to shine out of darkness, hath shined in our hearts, to give the light of the knowledge of the glory of God in the face of Jesus Christ. But we have this treasure in earthen vessels, that the excellency of the power may be of God, and not of us. We are troubled on every side, yet not distressed; we are perplexed, but not in despair; Persecuted, but not forsaken; cast down, but not destroyed; Always bearing about in the body the dying of the Lord Jesus, that the life also of Jesus might be made manifest in our body. For we which live are alway delivered unto death for Jesus' sake, that the life also of Jesus might be made manifest in our mortal flesh. So then death worketh in us, but life in you. We having the same spirit of faith, according as it is written, I believed, and therefore have I spoken; we also believe, and therefore speak; Knowing that he which raised up the Lord Jesus shall raise up us also by Jesus, and shall present us with you.

\alleluia{Alleluia, alleluia. ℣. Thou art a priest for ever, after the order of Melchisedech. Alleluia. ℣. Blessed is the man that endureth temptation: for when he is tried, he shall receive the crown of life. Alleluia.}

\begin{rubric}
	Outside of Eastertide, the following Gradual \& Lesser Alleluia replaces the Greater Alleluia.
\end{rubric}

\gradall{Behold a great priest who in his days pleased God. ℣. There was none found like unto him who kept the law of the Most High.}{Alleluia, alleluia. ℣. Blessed is the man that endureth temptation: for when he is tried, he shall receive the crown of life. Alleluia.}

\readingcitation{Gospel}{Matthew 10:23}
\lett{A}{t that time:} Jesus said unto his disciples: When they persecute you in this city, flee ye into another: for verily I say unto you, Ye shall not have gone over the cities of Israel, till the Son of man be come. The disciple is not above his master, nor the servant above his lord. It is enough for the disciple that he be as his master, and the servant as his lord. If they have called the master of the house Beelzebub, how much more shall they call them of his household? Fear them not therefore: for there is nothing covered, that shall not be revealed; and hid, that shall not be known. What I tell you in darkness, that speak ye in light: and what ye hear in the ear, that preach ye upon the housetops. And fear not them which kill the body, but are not able to kill the soul: but rather fear him which is able to destroy both soul and body in hell.

\offertory{I have found David my servant, with my holy oil have I anointed him: my hand shall hold him fast, and my arm shall strengthen him, alleluia.}

\secret
\lett{W}{e} beseech thee, O Lord, that our devout observance of the yearly solemnity of Saint Athanasius, thy Confessor and Bishop, may render us acceptable unto thy loving kindness: that this service of propitiation, which we duly offer, may be profitable unto him for the reward of blessed-ness, and obtain for us the gifts of thy grace. Through.

\communion{What I tell you in darkness that speak ye in light, saith the Lord: and what ye hear in the ear, that preach ye upon the housetops, alleluia.}

\postcommunion
\lett{O}{God,} who rewardest the souls of them that put their trust in thee: vouchsafe; that we who keep the solemn festival of blessed Athanasius, thy Confessor and Bishop, may by his prayers obtain thy merciful pardon. Through.

\bcpfeast{3 May. Invention of the Holy Cross}{Invention of the Cross}{3 May}\label{InventionoftheCross}
%\supplement{3 May}{Invention}{of the Holy Cross}

\subbysub{I Evensong}\label{InventionEvensong}

\gregorioscore{resources/gabc/ProperTime/InventionEvensong1.gabc}

\begin{rubric}
	During the following stanza all kneel, and the last stanza is never changed.
\end{rubric}

\gregorioscore{resources/gabc/ProperTime/InventionEvensong2.gabc}

	℣. This sign of the Cross shall be in heaven, alleluia.
	
	℟. When the Lord shall come to judgement, alleluia.

\properantiphon{Mag.}{O Cross, {\dag} surpassing all the stars in splendour, world renowned, exceeding dear unto the hearts of men, holier than all things: thou only wert counted worthy to uphold the world's random. Sweet the wood, sweet the iron, bearing so sweet a burden: bring aid to this congregation, who are here assembled to celebrate thy praises, alleluia, alleluia.}

\subbysub{Mattins}

\invitatoryhymn\label{InventionInvitatory}

\gregorioscore{resources/gabc/ProperTime/InventionInvitatory.gabc}

\officehymn\label{InventionMattins}

\gregorioscore{resources/gabc/ProperTime/InventionMattins.gabc}

	℣. We adore thee, O Christ, and we bless thee, alleluia.
	
	℟. Because by thy Cross thou hast redeemed the world, alleluia.

\properantiphon{Ben.}{Thou alone {\dag} excellest in stature all the cedars of Lebanon: for on thee the Life of the world was hanged, on thee was Christ victorious, and death over death did for ever triumph, alleluia.}

\subbysub{II Evensong}

\begin{rubric}
	II Evensong as in I Evensong, with the following Versicle \& Antiphon.
\end{rubric}

	℣. This sign of the Cross shall be in heaven, alleluia.
	
	℟. When the Lord shall come to judgement, alleluia.

\properantiphon{Mag.}{He the holy Cross endured, {\dag} Burst the gates of hell in twain: Begirt with might and majesty, On Easter morn he rose again, alleluia.}


\subby{Sts. Alexander, Eventius, Theodulus, and Juvenalis}
\feastday{{Sts. Alexander \&c.}}
\fancyhead[RE,LO]{3 May}
\begin{inhead}
    {Memorial\\
3 May}
\end{inhead}

\begin{rubric}
	The propers are from the Common of Many Martyrs in Eastertide (p. \pageref{CommonMartyrsEaster}).
\end{rubric}


\subby{St. Monica}
\feastday{{St. Monica}}
\fancyhead[RE,LO]{4 May}
\begin{inhead}
    {Memorial\\
4 May}
\end{inhead}

\begin{rubric}
	The propers are from the Common of Neither Virgin nor Martyr (p. \pageref{CommonNeitherVirginMartyr}), except for the following.
\end{rubric}

\collect
\lett{O}{God,} the comforter of them that mourn, and the salvation of them that hope in thee, who didst mercifully receive the loving tears of blessed Monica for the conversion of Augustine her son: grant to us by the intercession of them both; that we may bewail our sins, and obtain the pardon of thy grace. Through.

\begin{rubric}
	The Epistle is the additional Epistle from the same Common (p. \pageref{1Timothy53}).
\end{rubric}

\readingcitation{Gospel}{Luke 7:11}
\lett{A}{t that time:} Jesus went into a city called Nain: and many of his disciples went with him, and much people. Now when he came nigh to the gate of the city, behold, there was a dead man carried out, the only son of his mother, and she was a widow: and much people of the city was with her. And when the Lord saw her, he had compassion on her, and said unto her, Weep not. And he came and touched the bier: and they that bare him stood still. And he said, Young man, I say unto thee, Arise. And he that was dead sat up, and began to speak. And he delivered him to his mother. And there came a fear on all: and they glorified God, saying, That a great prophet is risen up among us; and, That God hath visited his people.

\secret
\lett{G}{rant,} O Lord, that like as thy faithful people do acknowledge that in tribulation they have been succoured by the merits of thy Saints: so this oblation, which they here do offer unto thee in honour of the same, may be acceptable in thy sight. Through.

\postcommunion
\lett{O}{Lord,} who satisfied thy family with sacred gifts: we beseech thee that we may at all times be comforted by the intercession of her whose festival we celebrate. Through.


\bcpfeast{6 May. St. John before the Latin Gate}{St. John Latin Gate}{6 May}
%\supplement{6 May}{St. John}{before the Latin Gate}

\begin{secrubric}
	The Daily Office propers are from the Common of Apostles in Easter (p. \pageref{CommonApostlesEaster}), with the following Antiphon for both Evensongs.
\end{secrubric}

\properantiphon{Mag.}{John the Apostle, {\dag} being cast into  a caldron of boiling oil, by virtue of protecting grace came forth unharmed, alleluia.}


\subby{St. Alexis Toth}
\feastday{{St. Alexis}}
\fancyhead[RE,LO]{7 May}
\begin{inhead}
    {Memorial\\
7 May}
\end{inhead}

%CHECK:
\begin{rubric}
	The propers are from the First Common of a Confessor not a Bishop (p. \pageref{CommonConfessorNotBishopI}).
\end{rubric}


\subby{The Apparition of St. Michael the Archangel}
\feastday{{St. Michael}}
\fancyhead[RE,LO]{8 May}
\begin{inhead}
    {Greater Double\\
8 May}
\end{inhead}

\introit
\lett{O}{praise} the Lord, all ye Angels of his: ye that excel in strength ye that fulfil his commandment, and hearken unto the voice of his words. \textit{Ps.} Praise the Lord, O my soul; and all that is within me praise his holy name

\collect
\lett{O}{everlasting} God, who hast ordained and constituted the services of Angels and men in a wonderful order; Mercifully grant that, as thy holy Angels always do thee service in heaven, so, by thy appointment, they may succour and defend us on earth. Through.

\readingcitation{Epistle}{Revelation 1:1}
\lett{I}{n those days:} God shewed the things which must shortly come to pass, and he sent and. signified it by his Angel unto his servant John, Who bare record of the word of God, and of the testimony of Jesus Christ, and of all things that he saw. Blessed is he that readeth, and they that hear the words of this prophecy, and keep those things which are written therein: for the time is at hand. John to the seven churches which are in Asia: Grace be unto you, and peace, from him which is, and which was, and which is to come; and from the seven Spirits which are before his throne; And from Jesus Christ, who is the faithful witness, and the first begotten of the dead, and the prince of the kings of the earth. Unto him that loved us, and washed us from our sins in his own blood.

\alleluia{Alleluia, alleluia. ℣. Holy Archangel Michael, defend us in the battle; that we perish not in the dreadful judgement. Alleluia. ℣. The sea was shaken, and the earth trembled, when the Archangel Michael came down from heaven. Alleluia.}

\readingcitation{Gospel}{Matthew 18:1}
\lett{A}{t that time:} The disciples came unto Jesus, saying, Who is the greatest in the kingdom of heaven? And Jesus called a little child unto him, and set him in the midst of them, And said, Verily I say unto you, Except ye be converted, and become as little children, ye shall not enter into the kingdom of heaven. Whosoever therefore shall humble himself as this little child, the same is greatest in the kingdom of heaven. And whoso shall receive one such little child in my name receiveth me. But whoso shall offend one of these little ones which believe in me, it were better for him that a millstone were hanged about his neck, and that he were drowned in the depth of the sea. Woe unto the world because of offences! for it must needs be that offences come; but woe to that man by whom the offence cometh! Wherefore if thy hand or thy foot offend thee, cut them off, and cast them from thee: it is better for thee to enter into life halt or maimed, rather than having two hands or two feet to be cast into everlasting fire. And if thine eye offend thee, pluck it out, and cast it from thee: it is better for thee to enter into life with one eye, rather than having two eyes to be cast into hell fire. Take heed that ye despise not one of these little ones; for I say unto you, That in heaven their angels do always behold the face of my Father which is in heaven.

\secret
\lett{W}{e} offer thee, O Lord, sacrifices of praise, humbly beseeching thee: that by the prayers of the Angels interceding for us, thou wouldest both graciously accept the same, and grant that they may avail to our salvation. Through.

\communion{O all ye Angels of the Lord, bless ye the Lord: sing ye praises, and magnify him above all for ever.}

\postcommunion
\lett{O}{Lord,} forasmuch as we put our trust in the intercession of thy blessed Archangel Michael; we humbly beseech thee; that those things which we touch with our lips we may likewise receive into our hearts. Through.

\begin{rubric}
	If it be a Sunday, the Gospel of this Feast is read at the end.
\end{rubric}


\subby{St. Boniface IV}
\feastday{{St. Boniface IV}}
\fancyhead[RE,LO]{8 May}
\begin{inhead}
    {Memorial\\
8 May}
\end{inhead}

%CHECK:
\begin{rubric}
	The propers are from the First Common of a Confessor Bishop (p. \pageref{CommonConfessorBishopI}).
\end{rubric}


\subby{St. Gregory Nazianzen}
\feastday{{St. Gregory Nazianzen}}
\fancyhead[RE,LO]{9 May}
\begin{inhead}
    {Double\\
9 May}
\end{inhead}

\begin{rubric}
	The propers are from the Common of Doctors (p. \pageref{CommonDoctors}), using the additional Epistle.
\end{rubric}


\subby{Sts. Gordian \& Epimachus}
\feastday{{St. Gordian \& Epimachus}}
\fancyhead[RE,LO]{10 May}
\begin{inhead}
    {Memorial\\
10 May}
\end{inhead}

\begin{rubric}
	The propers are from the Common of Many Martyrs in Eastertide (p. \pageref{CommonMartyrsEaster}), with the additional Epistle (p. \pageref{Revelation191}) and the Prayers as followeth.
\end{rubric}

\collect
\lett{G}{rant,} we beseech thee, almighty God: that we, who celebrate the festival of thy blessed Martyrs Gordian and Epimachus, may be aided by their intercession with thee. Through.

\secret
\lett{W}{e} beseech thee, O Lord, mercifully to accept this our sacrifice, which we offer unto thee, pleading the merits of thy blessed Martyrs Gordian and Epimachus: that the same may effectually avail for our perpetual succour. Through.

\postcommunion
\lett{W}{e} beseech thee, almighty God: that we, who have received this heavenly food, may at the intercession of thy blessed Martyrs, Gordian and Epimachus, be thereby defended against all adversities. Through.


\subby{Sts. Nereus, Achilles, Domitilla, and Pancras}
\feastday{{Sts. Nereus \&c.}}
\fancyhead[RE,LO]{12 May}
\begin{inhead}
    {Memorial\\
12 May}
\end{inhead}

\introit
\lett{B}{ehold,} the eye of the Lord is upon them that fear him, and put their trust in his mercy, alleluia: to deliver their soul from death: for he is our help and our shield, alleluia, alleluia. \textit{Ps.} Rejoice in the Lord, O ye righteous: for it becometh well the just to be thankful.

\collect
\lett{W}{e} beseech thee, O Lord, that the blessed solemnity of thy Martyrs, Nereus, Achilles, Domitilla, and Pancras, may ever protect us: and render us worthy of thy service. Through.

\begin{rubric}
	The Epistle is from the Common of a Martyr in Eastertide (p. \pageref{CommonMartyrEaster}).
\end{rubric}

\alleluia{Alleluia, alleluia. ℣. This is the true brotherhood which overcame the wickedness of the world: which followed Christ, gaining heaven's glorious realms. Alleluia. ℣. The noble army of Martyrs praise thee, O Lord. Alleluia.}

\readingcitation{Gospel}{John 4:46}
\lett{A}{t that time:} There was a certain nobleman, whose son was sick at Capernaum. When he heard that Jesus was come out of Judaea into Galilee, he went unto him, and besought him that he would come down, and heal his son: for he was at the point of death. Then said Jesus unto him, Except ye see signs and wonders, ye will not believe. The nobleman saith unto him, Sir, come down ere my child die. Jesus saith unto him, Go thy way; thy son liveth. And the man believed the word that Jesus had spoken unto him, and he went his way. And as he was now going down, his servants met him, and told him, saying, Thy son liveth. Then enquired he of them the hour when he began to amend. And they said unto him, Yesterday at the seventh hour the fever left him. So the father knew that it was at the same hour, in the which Jesus said unto him, Thy son liveth: and himself believed, and his whole house.

\offertory{O Lord, the very heavens shall praise thy wondrous works: and thy truth in the congregation of the saints, alleluia, alleluia.}

\secret
\lett{W}{e} beseech thee, O Lord, that the confession of thy holy Martyrs Nereus, Achilles, Domitilla and Pancras may be pleasing unto thee: that it may both commend our gifts, and ever implore for us thy pardon. Through.

\communion{Rejoice in the Lord, O ye righteous, alleluia: for it becometh well the just to be thankful, alleluia.}

\postcommunion
\lett{W}{e} beseech thee, O Lord: that by the prayers of thy blessed Martyrs Nereus, Achilles, Domitilla and Pancras, the holy sacraments which we have received may profit us to the increase of thy merciful pardon. Though.


\subby{St. Boniface of Tarsus}
\feastday{{St. Boniface Tarsus}}
\fancyhead[RE,LO]{14 May}
\begin{inhead}
    {Memorial\\
14 May}
\end{inhead}

\begin{rubric}
	The propers are from the Common of a Martyr in Eastertide (p. \pageref{CommonMartyrEaster}), with the Prayers as followeth..
\end{rubric}

\collect
\lett{G}{rant,} we beseech thee, almighty God: that we, who celebrate the festival of blessed Boniface thy Martyr, may be aided by his intercession with thee. Through.

\secret
\lett{W}{e} beseech thee, O Lord, to accept our prayers and oblations: and graciously hearken unto us whom thou dost cleanse by thy heavenly mysteries. Through.

\postcommunion
\lett{W}{e} beseech thee, O Lord our God, that like as we, whom thou hast refreshed by the partaking of thy sacred gift, do offer unto thee our worship: so by the intercession of blessed Boniface, thy Martyr, we may perceive the benefits of the same. Through.


\subby{St. Venantius}
\feastday{{St. Venantius}}
\fancyhead[RE,LO]{18 May}
\begin{inhead}
    {Double\\
18 May}
\end{inhead}

\begin{rubric}
	The propers are from the Common of a Martyr in Eastertide (p. \pageref{CommonMartyrEaster}), with the Prayers as followeth.
\end{rubric}

\collect
\lett{O}{God,} who hast hallowed this day by the triumph of blessed Venantius, thy Martyr: graciously hear the prayers of thy people, and grant; that we who venerate his merits may imitate the constancy of his faith. Through.

\secret
\lett{A}{lmighty} God, let the merits of blessed Venantius render this oblation acceptable unto thee: that we, being succoured by his assistance, may be made partakers of his glory. Through.

\postcommunion
\lett{W}{e} have received, O Lord, the sacraments of everlasting life, humbly beseeching thee: that through the prayers of blessed Venantius thy Martyr on our behalf, they may obtain for us both pardon and grace. Through.


\subby{St. Pudentiana}
\feastday{{St. Pudentiana}}
\fancyhead[RE,LO]{19 May}
\begin{inhead}
    {Memorial\\
19 May}
\end{inhead}

\begin{rubric}
	The propers are from the First Common of a Virgin (p. \pageref{CommonVirginOnlyI}).
\end{rubric}


\subby{St. Romanus}
\feastday{{St. Romanus}}
\fancyhead[RE,LO]{22 May}
\begin{inhead}
    {Memorial\\
22 May}
\end{inhead}

\begin{rubric}
	The propers are from the Common of Abbots (p. \pageref{CommonAbbots}).
\end{rubric}


\subby{St. Vincent of Lerins}
\feastday{{St. Vincent Lerins}}
\fancyhead[RE,LO]{24 May}
\begin{inhead}
    {Memorial\\
24 May}
\end{inhead}

\begin{rubric}
	The propers are from the First Common of a Confessor not a Bishop (p. \pageref{CommonConfessorNotBishopI}).
\end{rubric}


\subby{Pope St. Urban I}
\feastday{{Pope St. Urban I}}
\fancyhead[RE,LO]{25 May}
\begin{inhead}
    {Memorial\\
25 May}
\end{inhead}

\begin{rubric}
	The propers are from the First Common of a Martyr Bishop (p. \pageref{CommonMartyrBishopI}).
\end{rubric}


\subby{St. Augustine of Canterbury}
\feastday{{St. Augustine Canterbury}}
\fancyhead[RE,LO]{26 May}
\begin{inhead}
    {Double\\
26 May}
\end{inhead}

\begin{rubric}
	The Daily Office propers are from the First Common of a Confessor Bishop (p. \pageref{CommonConfessorBishopI}).
\end{rubric}

\introit
\lett{L}{et} thy priests, O Lord, be clothed with righteousness: and let thy saints sing with joyfulness: for thy servant David's sake, turn not away the presence of thine Anointed. Alleluia, alleluia. \textit{Ps.} Lord, remember David: and all his trouble.

\collect
\lett{O}{God,} who didst give the blessed Bishop Augustine to be the first Teacher of the English people: grant us, we beseech thee; that we who proclaim his merits on earth may perceive his intercession in heaven. Through.

\begin{rubric}
	Commemoration of St. Eleutherius (p. \pageref{EleutheriusCollect}).
\end{rubric}

\readingcitation{Epistle}{Hebrews 7:23}
\lett{B}{rethren:} They were many priests, because they were not suffered to continue by reason of death: But this man, because he continueth ever, hath an unchangeable priesthood. Wherefore he is able also to save them to the uttermost that come unto God by him, seeing he ever liveth to make intercession for them. For such an high priest became us, who is holy, harmless, undefiled, separate from sinners, and made higher than the heavens; Who needeth not daily, as those high priests, to offer up sacrifice, first for his own sins, and then for the people's: for this he did once, when he offered up himself, Jesus Christ our Lord.

\alleluia{Alleluia, alleluia. ℣. The Lord sware, and will not repent: Thou art a priest for ever after the order of Melchisedech. Alleluia. ℣. The Lord loved him, and adorned him: he clothed him with a robe of glory. Alleluia.}

\readingcitation{Gospel}{Luke 10:1}
\lett{A}{t that time:} The Lord appointed other seventy also, and sent them two and two before his face into every city and place, whither he himself would come. Therefore said he unto them, The harvest truly is great, but the labourers are few: pray ye therefore the Lord of the harvest, that he would send forth labourers into his harvest. Go your ways: behold, I send you forth as lambs among wolves. Carry neither purse, nor scrip, nor shoes: and salute no man by the way. And into whatsoever house ye enter, first say, Peace be to this house. And if the son of peace be there, your peace shall rest upon it: if not, it shall turn to you again. And in the same house remain, eating and drinking such things as they give: for the labourer is worthy of his hire. Go not from house to house. And into whatsoever city ye enter, and they receive you, eat such things as are set before you: And heal the sick that are therein, and say unto them, The kingdom of God is come nigh unto you.

\offertory{My truth and my mercy shall be with him: and in my name shall his horn be exalted. Alleluia.}

\secret
\lett{W}{e} beseech thee, O Lord, that the gifts which we offer may be acceptable unto thee: whereby we venerate the merits of blessed Augustine, thy Confessor and Bishop, and likewise call to remembrance the pledges of our life and freedom. Through.

\begin{rubric}
	Commemoration of St. Eleutherius (p. \pageref{EleutheriusSecret}).
\end{rubric}

\communion{Blessed is the servant, whom the lord when he cometh shall find watching: verily I say unto you, He shall make him ruler over all his goods. Alleluia.}

\postcommunion
\lett{W}{e} beseech thee, O Lord, that we may be nourished by thy holy things, which we have received for the solemnity of blessed Augustine, thy Confessor and Bishop: wherewith we may continually be satisfied, and evermore desire to be filled. Through.

\begin{rubric}
	Commemoration of St. Eleutherius (p. \pageref{EleutheriusPostcommunion}).
\end{rubric}


\subby{Pope St. Eleutherius}
\feastday{{Pope St. Eleutherius}}
\fancyhead[RE,LO]{26 May}
\begin{inhead}
    {Memorial\\
26 May}
\end{inhead}

\begin{rubric}
	The propers are from the Common of a Martyr in Eastertide (p. \pageref{CommonMartyrEaster}), with the following Prayers.
\end{rubric}

\collect\label{EleutheriusCollect}
\lett{A}{lmighty} God, mercifully look upon our infirmities: that, whereas we are oppressed by the burden of our sins, the glorious intercession of blessed Eleutherius, thy Martyr and Bishop, may be our succour and defence. Through.

\secret\label{EleutheriusSecret}
\lett{S}{anctify,} O Lord, the gifts which we dedicate to thee: that at the intercession of blessed Eleutherius, thy Martyr and Bishop, they may obtain for us thy gracious favour. Through.

\postcommunion\label{EleutheriusPostcommunion}
\lett{M}{ay} this communion, O Lord, cleanse us from guilt: and, at the intercession of blessed Eleutherius, thy Martyr and Bishop, make us partakers of thy heavenly healing. Through.


\subby{St. Bede the Venerable}
\feastday{{St. Bede the Venerable}}
\fancyhead[RE,LO]{27 May}
\begin{inhead}
    {Double\\
27 May}
\end{inhead}

\begin{rubric}
	The propers are from the Common of Doctors (p. \pageref{CommonDoctors}), with the following Prayers.
\end{rubric}

\collect
\lett{O}{God,} who dost illumine thy Church with the learning of blessed Bede, thy Confessor and Doctor: mercifully grant unto thy servants; that they may ever be enlightened by his wisdom, and succoured by his merits. Through.

\begin{rubric}
	Commemoration of Pope St. John I from the Second Collect of the Common of a Martyr in Eastertide (p. \pageref{CommonMartyrEaster}).
\end{rubric}

\secret
\lett{M}{ay} the devout prayers of thy Confessor and Doctor, Saint Bede, never fail to succour us, O Lord: that they may render our oblations acceptable in thy sight; and may ever obtain for us thy merciful pardon. Through.

\begin{rubric}
	Commemoration of Pope St. John I from the Second Secret of the Common of a Martyr in Eastertide (p. \pageref{CommonMartyrEaster}).
\end{rubric}

\postcommunion
\lett{W}{e} beseech thee, O Lord, that blessed Bede, thy Confessor and illustrious Doctor, may stand before thee as our advocate: that these thy sacrifices may avail for our salvation. Through.

\begin{rubric}
	Commemoration of Pope St. John I from the Second Postcommunion of the Common of a Martyr in Eastertide (p. \pageref{CommonMartyrEaster}).
\end{rubric}


\subby{Pope St. John I}
\feastday{{Pope St. John I}}
\fancyhead[RE,LO]{27 May}
\begin{inhead}
    {Memorial\\
27 May}
\end{inhead}

\begin{rubric}
	The propers are from the Common of a Martyr in Eastertide (p. \pageref{CommonMartyrEaster}), using the Second Prayers.
\end{rubric}


\subby{Pope St. Felix I}
\feastday{{Pope St. Felix I}}
\fancyhead[RE,LO]{30 May}
\begin{inhead}
    {Memorial\\
30 May}
\end{inhead}

\begin{rubric}
	In Eastertide, the propers are from Common of a Martyr in Eastertide (p. \pageref{CommonMartyrEaster}), with the First Collect and the Second Secret and Postcommunion.
\end{rubric}

\begin{rubric}
	Outside of Eastertide, the propers are from the First Common of a Martyr Bishop (p. \pageref{CommonMartyrBishopI}).
\end{rubric}


\subby{St. Petronilla}
\feastday{{St. Petronilla}}
\fancyhead[RE,LO]{31 May}
\begin{inhead}
    {Memorial\\
31 May}
\end{inhead}

\begin{rubric}
	The propers are from the Second Common of a Virgin (p. \pageref{CommonVirginOnlyII}).
\end{rubric}


\subby{Sts. Marcellinus, Peter, \& Erasmus}
\feastday{{Sts. Marcellinus \&c.}}
\fancyhead[RE,LO]{2 June}
\begin{inhead}
    {Memorial\\
2 June}
\end{inhead}

\begin{rubric}
	Outside of Eastertide, the propers are as followeth.
\end{rubric}

\begin{rubric}
	Within Eastertide, the propers are from the Common of Many Martyrs in Eastertide (p. \pageref{CommonMartyrsEaster}) with the Prayers and Epistle as followeth.
\end{rubric}

\introit
\lett{T}{he} righteous cry, and the Lord heareth them: and delivereth them out of all their troubles. \textit{Ps.} I will alway give thanks unto the Lord: his praise shall ever be in my mouth.

\collect
\lett{O}{God,} who makest us glad with the yearly solemnity of thy blessed Martyrs, Marcellinus, Peter, and Erasmus: grant, we beseech thee; that, as we rejoice in their merits, so we may be enkindled by their example. Through.

\begin{rubric}
	The Epistle is from the third additional Epistle of the Third Common of Many Martyrs (p. \pageref{Romans818}).
\end{rubric}

\gradall{The righteous cry, and the Lord heareth them: and delivereth them out of all their troubles. ℣. The Lord is nigh unto them that are of a contrite heart: and will save such as be of an humble spirit.}{Alleluia, alleluia. ℣. I have chosen you out of the world that ye should go and bring forth fruit; and that your fruit should remain. Alleluia.}

\begin{rubric}
	In Eastertide, the following Alleluia Verse replaces the Gradual \& Lesser Alleluia.
\end{rubric}

\alleluia{Alleluia, alleluia. ℣. I have chosen you out of the world, that ye should go and bring forth fruit; and that your fruit should remain. Alleluia. ℣. Right dear in the sight of the Lord is the death of his saints. Alleluia.}

\begin{rubric}
	The Gospel is from the First Common of Many Martyrs (p. \pageref{CommonMartyrsI}).
\end{rubric}

\offertory{Be glad, ye righteous, and rejoice in the Lord: and be joyful, all ye that are true of heart.}

\secret
\lett{W}{e} beseech thee, O Lord, that this sacrifice which we offer in remembrance of the birthday of thy holy Martyrs: may both loose the bonds of our iniquity, and obtain for us the gifts of thy mercy. Through.

\communion{The souls of the just are in the hand of God, and there shall no torment of malice touch them: in the sight of the unwise they seemed to die, but they are in peace.}

\postcommunion
\lett{O}{Lord,} who hast satisfied us with this sacred gift, we humbly beseech thee: that in the mysteries which we celebrate in this service of our bounden duty, we may perceive the increase of thy saving grace. Through.


\subby{St. Boniface}
\feastday{{St. Boniface}}
\fancyhead[RE,LO]{5 June}
\begin{inhead}
    {Double\\
5 June}
\end{inhead}

\begin{rubric}
	The Daily Office propers are from the First Common of a Martyr Bishop (p. \pageref{CommonMartyrBishopI}).
\end{rubric}

\introit
\lett{I}{will} rejoice in Jerusalem and joy in my people: and the voice of weeping shall be no more heard in her, nor the voice of crying. Mine elect shall not labour in vain, nor bring forth for trouble: for they are the seed of the blessed of the Lord, and their offspring with them. (Alleluia, alleluia.) \textit{Ps.} We have heard with our ears, O God: our fathers have told us, what thou hast done in their time of old.

\collect
\lett{O}{God,} who by the zeal of blessed Boniface, thy Martyr and Bishop, didst vouchsafe to call a multitude of peoples to the knowledge of thy name: mercifully grant; that we who celebrate his festival may likewise perceive his advocacy. Through.

\readingcitation{Epistle}{Ecclesiasticus 44:1}
\lett{L}{et} us now praise famous men, And our fathers that begat us. The Lord manifested in them great glory, Even his mighty power from the beginning. Such as did bear rule in their kingdoms, And were men renowned for their power, Giving counsel by their understanding, Such as have brought tidings in prophecies: Leaders of the people by their counsels, And by their understanding men of learning for the people; Wise were their words in their instruction: Such as sought out musical tunes, And set forth verses in writing: Rich men furnished with ability, Living peaceably in their habitations: All these were honoured in their generations, And were a glory in their days. There be of them, that have left a name behind them, To declare their praises. And some there be, which have no memorial; Who are perished as though they had not been, And are become as though they had not been born; And their children after them. But these were men of mercy, Whose righteous deeds have not been forgotten. With their seed shall remain continually a good inheritance; Their children are within the covenants. Their seed standeth fast, And their children for their sakes. Their seed shall remain for ever, And their glory shall not be blotted out. Their bodies were buried in peace, And their name liveth to all generations. Peoples will declare their wisdom, And the congregation telleth out their praise.

\gradall{Rejoice, inasmuch as ye are partakers of Christ's sufferings, that, when his glory shall be revealed, ye may be glad also with exceeding joy. ℣. If ye be reproached for the name of Christ, happy are ye: for the spirit of glory and of God resteth upon you.}{Alleluia, alleluia. ℣. I will extend peace to her like a river, and glory like a flowing stream. Alleluia.}

\begin{rubric}
	In Eastertide, the following Alleluia Verse replaces the Gradual \& Lesser Alleluia.
\end{rubric}

\alleluia{Alleluia, alleluia. ℣. Rejoice ye with Jerusalem, and be glad with her, all ye that love the Lord. Alleluia. ℣. When ye see this, your heart shall rejoice: and the hand of the Lord shall be known toward his servants. Alleluia.}

\readingcitation{Gospel}{Matthew 5:1}
\lett{A}{t that time:} Jesus, seeing the multitudes, went up into a mountain: and when he was set, his disciples came unto him: And he opened his mouth, and taught them, saying, Blessed are the poor in spirit: for theirs is the kingdom of heaven. Blessed are they that mourn: for they shall be comforted. Blessed are the meek: for they shall inherit the earth. Blessed are they which do hunger and thirst after righteousness: for they shall be filled. Blessed are the merciful: for they shall obtain mercy. Blessed are the pure in heart: for they shall see God. Blessed are the peacemakers: for they shall be called the children of God. Blessed are they which are persecuted for righteousness' sake: for theirs is the kingdom of heaven. Blessed are ye, when men shall revile you, and persecute you, and shall say all manner of evil against you falsely, for my sake. Rejoice, and be exceeding glad: for great is your reward in heaven.

\offertory{I will thank the Lord for giving me warning: I have set God always before me, for he is on my right hand, therefore I shall not fall. (Alleluia.)}

\secret
\lett{L}{et} thy plenteous benediction, we beseech thee, O Lord, come down upon these sacrifices: that it may mercifully work out our sanctification; and make us to rejoice in the solemnity of holy Boniface, thy Martyr and Bishop. Through.

\communion{To him that overcometh will I grant to sit with me in my throne: even as I also overcame, and am set down with my Father in his throne. (Alleluia.)}

\postcommunion
\lett{O}{Lord,} who hast sanctified us with this saving mystery: we beseech thee; that holy Boniface thy Martyr and Bishop, whom thou hast given to be our advocate and guide, may never fail devoutly to pray for us. Through.


\subby{St. Columba of Iona}
\feastday{{St. Columba}}
\fancyhead[RE,LO]{9 June}
\begin{inhead}
    {Double\\
9 June}
\end{inhead}

\begin{rubric}
	The propers are from the Common of Abbots (p. \pageref{CommonAbbots}), with Commemoration of Sts. Primus \& Felician.
\end{rubric}


\subby{Sts. Primus \& Felician}\label{PrimusFelician}
\feastday{{Sts. Primus \& Felician}}
\fancyhead[RE,LO]{9 June}
\begin{inhead}
    {Memorial\\
9 June}
\end{inhead}

\begin{rubric}
	In Eastertide, the propers are from the Common of Many Martyrs in Eastertide (p. \pageref{CommonMartyrsEaster}, with the Prayers, Greater Alleluia, and Gospel as followeth.
\end{rubric}

\begin{rubric}
	Outside of Eastertide, the propers are as followeth.
\end{rubric}

\introit
\lett{L}{et} the people tell of the wisdom of the Saints, and let the church shew forth their praise: their names shall live for evermore. \textit{Ps.} Rejoice in the Lord, O ye righteous: for it becometh well the just to be thankful.

\collect
\lett{O}{Lord,} we beseech thee, make us ever to rejoice in the festival of thy holy Martyrs Primus and Felician: that by their prayers we may perceive the gifts of thy protection. Through.

\begin{rubric}
	The Epistle is from the Second Common of Many Martyrs (p. \pageref{CommonMartyrsII}).
\end{rubric}

\gradall{O Lord, the very heavens shall praise thy wondrous works: and thy truth in the congregation of the saints. ℣. My song shall be alway of the loving-kindness of the Lord: from one generation to another.}{Alleluia, alleluia. ℣. This is the true brotherhood, which overcame the wickedness of the world: which followed Christ, gaining heaven's glorious realms. Alleluia.}

\begin{rubric}
	In Eastertide, the following Alleluia Verse replaces the Gradual \& Lesser Alleluia.
\end{rubric}

\alleluia{Alleluia, alleluia. ℣. This is the true brotherhood which overcame the wickedness of the world: which followed Christ, gaining heaven's glorious realms. Alleluia. ℣. The noble army of Martyrs praise thee, O Lord. Alleluia.}

\readingcitation{Gospel}{Matthew 11:25}
\lett{A}{t that time:} Jesus answered and said: I thank thee, O Father, Lord of heaven and earth, because thou hast hid these things from the wise and prudent, and hast revealed them unto babes. Even so, Father: for so it seemed good in thy sight. All things are delivered unto me of my Father: and no man knoweth the Son, but the Father; neither knoweth any man the Father, save the Son, and he to whomsoever the Son will reveal him. Come unto me, all ye that labour and are heavy laden, and I will give you rest. Take my yoke upon you, and learn of me; for I am meek and lowly in heart: and ye shall find rest unto your souls. For my yoke is easy, and my burden is light.

\offertory{God is wonderful in his holy ones: even the God of Israel, he will give strength and power unto his people: blessed be God, alleluia.}

\secret
\lett{W}{e} beseech thee, O Lord, that the oblation to be consecrated on the day of the precious death of thy Martyrs may be acceptable unto thee: for the cleansing of our sins, and for the commending to thee of the prayers of thy servants. Through.

\communion{I have chosen you out of the world, that ye should go and bring forth fruit, and that your fruit should remain.}

\postcommunion
\lett{W}{e} beseech thee, almighty God: that the solemnity of thy holy Martyrs Primus and Felician, which we have celebrated with heavenly mysteries, may obtain for us thy merciful pardon. Through.


\subby{St. Margaret of Scotland}
\feastday{{St. Margaret Scotland}}
\fancyhead[RE,LO]{10 June}
\begin{inhead}
    {Memorial\\
10 June}
\end{inhead}

\begin{rubric}
	The propers are from the Common of a Woman Neither Virgin Nor Martyr (p. \pageref{CommonNeitherVirginMartyr}), with the Prayers as followeth.
\end{rubric}

\collect
\lett{O}{God,} who didst render the blessed Queen Margaret wondrous by reason of her eminent charity towards the poor: grant; that, by her intercession and example, thy charity may continually increase in our hearts. Through.

\secret
\lett{G}{rant,} O Lord, that like as thy faithful people do acknowledge that in tribulation they have been succoured by the merits of thy Saints: so this oblation, which they here do offer unto thee in honour of the same, may be acceptable in thy sight. Through.

\postcommunion
\lett{O}{Lord,} who hast satisfied thy family with sacred gifts: we beseech thee that we may at all times be comforted by the intercession of her whose festival we celebrate. Through.


\bcpfeast{11 June. St. Barnabas}{St. Barnabas}{11 June}
%\supplement{11 June}{St. Barnabas}{}

\begin{secrubric}
	The Daily Office propers are from the Common of Apostles (p. \pageref{CommonApostles}).
\end{secrubric}

\begin{secrubric}
	In Eastertide, the Daily Office propers are from the Common of Apostles in Eastertide (p. \pageref{CommonApostlesEaster}).
\end{secrubric}


\subby{Sts. Basilides, Cyrinus, Nabor, \& Nazarius}
\feastday{{Sts. Basilides \&c.}}
\fancyhead[RE,LO]{12 June}
\begin{inhead}
    {Memorial\\
12 June}
\end{inhead}

\begin{rubric}
	In Eastertide, the propers are from the Common of Many Martyrs in Eastertide (p. \pageref{CommonMartyrsEaster}), with the following Prayers.
\end{rubric}

\begin{rubric}
	Outside of Eastertide, the propers are as followeth.
\end{rubric}

\introit
\lett{L}{et} the sorrowful sighing of the prisoners, O Lord, come before thee: reward thou our neighbours sevenfold into their bosom: avenge thou the blood of thy Saints, that is shed. \textit{Ps.} O God, the heathen are come into thine inheritance: thy holy temple have they defiled: and made Jerusalem an heap of stones.

\collect
\lett{W}{e} beseech thee, O Lord, that the observance of the birthday of thy holy Martyrs Basilides, Cyrinus, Nabor, and Nazarius may shine brightly upon us: that the everlasting excellence which they have obtained may increase by the fruits of our devotion. Through.

\begin{rubric}
	The Epistle is from the Third Common of Many Martyrs (p. \pageref{Hebrews1032}).
\end{rubric}

\gradall{Avenge thou, O Lord, the blood of thy Saints that is shed. ℣. The dead bodies of thy servants have they given to be meat unto the fowls of the air: and the flesh of thy saints unto the beasts of the land}{Alleluia, alleluia. ℣. The bodies of the Saints are buried in peace: but their name liveth for evermore. Alleluia.}

\begin{rubric}
	The Gospel is from the first additional Gospel from the Third Common of Many Martyrs (p. \pageref{Matthew243}).
\end{rubric}

\offertory{Let the Saints be joyful with glory, let them rejoice in their beds: let the praises of God be in their mouth.}

\secret
\lett{W}{e} solemnly offer unto thee, O Lord, these sacrifices, rehearsing thy wondrous works, in honour of the blood of thy Saints Basilides, Cyrinus, Nabor, and Nazarius: whereby their glorious victory was made perfect. Through.

\communion{The dead bodies of thy servants, O Lord, have they given to be meat unto the fowls of the air, and the flesh of thy saints unto the beasts of the land: according to the greatness of thy power, preserve thou those that are appointed to die.}

\postcommunion
\lett{G}{rant,} we beseech thee, O Lord: that we, ever celebrating the festival of thy holy Martyrs Basilides, Cyrinus, Nabor, and Nazarius; may continually perceive their advocacy. Through.


\subby{St. Basil the Great}
\feastday{{St. Basil the Great}}
\fancyhead[RE,LO]{14 June}
\begin{inhead}
    {Greater Double\\
14 June}
\end{inhead}

\begin{rubric}
	The Daily Office propers are from the Common of Doctors (p. \pageref{CommonDoctors}).
\end{rubric}

\introit
\lett{I}{n} the midst of the Church he opened his mouth: and the Lord filled him with the spirit of wisdom and of understanding: he clothed him with a robe of glory. \textit{Ps.} It is a good thing to give thanks unto the Lord: and to sing praises unto thy name, O Most Highest.

\collect
\lett{W}{e} beseech thee, O Lord, graciously to hear the prayers which we offer unto thee on the solemnity of blessed Basil, thy Confessor and Bishop: that, like as he was found worthy to do thee faithful service, so by his merits and intercession we may be absolved from all our sins. Through.

\begin{rubric}
	The Epistle is from the Common of Doctors (p. \pageref{2Timothy41}).
\end{rubric}

\gradall{The mouth of the righteous is exercised in wisdom, and his tongue will be talking of judgment. ℣. The law of his God is in his heart: and his goings shall not slide.}{Alleluia, alleluia. ℣. I have found David my servant, with my holy oil have I anointed him. Alleluia.}

\readingcitation{Gospel}{Luke 14:26}
\lett{A}{t that time:} Jesus said unto the multitudes: If any man come to me, and hate not his father, and mother, and wife, and children, and brethren, and sisters, yea, and his own life also, he cannot be my disciple. And whosoever doth not bear his cross, and come after me, cannot be my disciple. For which of you, intending to build a tower, sitteth not down first, and counteth the cost, whether he have sufficient to finish it? Lest haply, after he hath laid the foundation, and is not able to finish it, all that behold it begin to mock him, Saying, This man began to build, and was not able to finish. Or what king, going to make war against another king, sitteth not down first, and consulteth whether he be able with ten thousand to meet him that cometh against him with twenty thousand? Or else, while the other is yet a great way off, he sendeth an ambassage, and desireth conditions of peace. So likewise, whosoever he be of you that forsaketh not all that he hath, he cannot be my disciple. Salt is good: but if the salt have lost his savour, wherewith shall it be seasoned? It is neither fit for the land, nor yet for the dunghill; but men cast it out. He that hath ears to hear, let him hear.

\offertory{My truth and my mercy shall be with him: and in my name shall his horn be exalted.}

\secret
\lett{W}{e} beseech thee, O Lord, that our devout observance of the yearly solemnity of Saint Basil, thy Confessor and Bishop, may render us acceptable unto thy loving kindness: that this service of propitiation, which we duly offer, may be profitable unto him for the reward of blessedness, and obtain for us the gifts of thy grace. Through.

\communion{A faithful and wise servant, whom the lord hath made ruler over his household: to give them their portion of meat in due season.}

\postcommunion
\lett{O}{God,} who rewardest the souls of them that put their trust in thee: vouchsafe; that we who keep the solemn festival of blessed Basil, thy Confessor and Bishop, may by his prayers obtain thy merciful pardon. Through.


\subby{Sts. Vitus, Modestus, \& Crescentia}
\feastday{{Sts. Vitus \&c.}}
\fancyhead[RE,LO]{15 June}
\begin{inhead}
    {Memorial\\
15 June}
\end{inhead}

\begin{rubric}
	In Eastertide, the propers are from the Common of Many Martyrs in Eastertide (p. \pageref{CommonMartyrsEaster}), with Prayers and Gospel as followeth.
\end{rubric}

\begin{rubric}
	Outside of Eastertide, the propers are as followeth.
\end{rubric}

\introit
\lett{G}{reat} are the troubles of the righteous, but the Lord delivereth him out of all: the Lord keepeth all his bones: so that not one of them is broken. \textit{Ps.} I will alway give thanks unto the Lord: his praise shall ever be in my mouth.

\collect
\lett{G}{rant,} O Lord, we beseech thee, that, at the intercession of thy holy Martyrs, Vitus, Modestus, and Crescentia, thy Church may learn not to be highminded, but to grow in such lowliness as is acceptable unto thee: that, despising things evil, she may with bounteous charity perform those things that be right. Through.

\begin{rubric}
	The Epistle is from First Common of Many Martyrs (p. \pageref{Wisdom31}).
\end{rubric}

\gradall{Let the Saints be joyful with glory: let them rejoice in their beds. ℣. O sing unto the Lord a new song: let the congregation of saints praise him.}{Alleluia, alleluia. ℣. Thy Saints shall give thanks unto thee, O Lord: they shew the glory of thy kingdom. Alleluia.}

\readingcitation{Gospel}{Luke 10:16}
\lett{A}{t that time:} Jesus said unto his disciples: He that heareth you heareth me; and he that despiseth you despiseth me; and he that despiseth me despiseth him that sent me. And the seventy returned again with joy, saying, Lord, even the devils are subject unto us through thy name. And he said unto them, I beheld Satan as lightning fall from heaven. Behold, I give unto you power to tread on serpents and scorpions, and over all the power of the enemy: and nothing shall by any means hurt you. Notwithstanding in this rejoice not, that the spirits are subject unto you; but rather rejoice, because your names are written in heaven.

\secret
\lett{O}{Lord,} forasmuch as the gifts which we offer for thy Saints do shew forth the glory of thy divine power: so may they achieve in us the effects of thy salvation. Through.

\communion{The souls of the just are in the hand of God, and there shall no torment of malice touch them: in the sight of the unwise they seemed to die: but they are in peace.}

\postcommunion
\lett{O}{Lord,} who hast fulfilled us with thy solemn benediction: we beseech thee; that through the intercession of thy holy Martyrs Vitus, Modestus, and Crescentia, the healing of this sacrament may be profitable both for our bodies and our souls. Through.


\subby{St. Ephrem the Syrian}
\feastday{{St. Ephrem Syrian}}
\fancyhead[RE,LO]{18 June}
\begin{inhead}
    {Double\\
18 June}
\end{inhead}

\begin{rubric}
	The propers are from the Common of Doctors (p. \pageref{CommonDoctors}), with the Prayers as followeth.
\end{rubric}

\begin{rubric}
	\textsc{Note,} The Creed is said.
\end{rubric}

%MANUAL ADJUSTMENT:
\vspace{-1ex}

\collect
\lett{O}{God,} who wast pleased to adorn thy Church with the wondrous learning and glorious merits of the life of blessed Ephram, thy Confessor and Doctor: we humbly entreat thee; that at his intercession thou wouldest defend her with thy continual power against the snares of error and wickedness. Through.

\begin{rubric}
	Commemoration of Sts. Marcus \& Marcellianus (p. \pageref{MarcusMarcellianusCollect}).
\end{rubric}

%MANUAL ADJUSTMENT:
\vspace{-1ex}

\secret
\lett{M}{ay} the devout prayers of Saint Ephrem, thy Confessor and Doctor, never fail to succour us, O Lord: that they may render our oblations acceptable in thy sight; and may ever obtain for us thy merciful pardon. Through.

\begin{rubric}
	Commemoration of Sts. Marcus \& Marcellianus (p. \pageref{MarcusMarcellianusSecret}).
\end{rubric}

%MANUAL ADJUSTMENT:
\vspace{-1ex}

\postcommunion
\lett{W}{e} beseech thee, O Lord, that blessed Ephrem, thy Confessor and illustrious Doctor, may stand before thee as our advocate: that these thy sacrifices may avail for our salvation. Through.

\begin{rubric}
	Commemoration of Sts. Marcus \& Marcellianus (p. \pageref{MarcusMarcellianusPostcommunion}).
\end{rubric}

%MANUAL ADJUSTMENT:
\clearpage

\subby{Sts. Marcus \& Marcellianus}
\feastday{{Sts. Marcus \&c.}}
\fancyhead[RE,LO]{18 June}
\begin{inhead}
    {Memorial\\
18 June}
\end{inhead}

\begin{rubric}
	In Eastertide, the propers are from the Common of Many Martyrs in Eastertide (p. \pageref{CommonMartyrsEaster}), with the Prayers and Gospel as followeth.
\end{rubric}

\begin{rubric}
	Outside Eastertide, the propers are as followeth.
\end{rubric}

\introit
\lett{B}{ut} the salvation of the righteous cometh of the Lord: who is also their strength in the time of trouble. \textit{Ps.} Fret not thyself because of the ungodly: neither be thou envious against the evil-doers.

\collect\label{MarcusMarcellianusCollect}
\lett{G}{rant,} we beseech thee, almighty God: that we who devoutly celebrate the birthday of thy holy Martyrs Marcus and Marcellianus; may, by their intercession, be delivered from all evils that beset us. Through.

\begin{rubric}
	The Epistle is from the second additional Epistle from the Third Common of Many Martyrs (p. \pageref{Romans51}).
\end{rubric}

\gradall{The souls of the just are in the hand of God: and there shall no torment of malice touch them. ℣. In the eyes of the unwise they seemed to die: but they are in peace.}{Alleluia, alleluia. ℣. This is the true brotherhood which no conflict could sunder: they who, shedding their blood, followed the Lord. Alleluía.}

\begin{rubric}
	In Eastertide, the following Alleluia Verse replaces the Gradual \& Lesser Alleluia.
\end{rubric}

\alleluia{Alleluia, alleluia. ℣. This is the true brotherhood which no conflict could sunder: they who, shedding their blood, followed the Lord. Alleluia.  ℣. The noble army of Martyrs praise thee, O Lord. Alleluia.}

\begin{rubric}
	The Gospel is from the fourth additional Gospel from the Third Common of Many Martyrs (p. \pageref{Luke1147}).
\end{rubric}

\offertory{Our soul is escaped even as a bird out of the snare of the fowler: the snare is broken, and we are delivered.}

\secret\label{MarcusMarcellianusSecret}
\lett{S}{anctify,} O Lord, the gifts which we dedicate unto thee: that, at the intercession of thy holy Martyrs Marcus and Marcellianus, they may obtain for us thy gracious favour. Through.

\communion{Verily I say unto you, whatsoever ye have done unto one of the least of mine, ye have done it unto me: come, ye blessed of my Father, inherit the kingdom prepared for you from the foundation of the world.}

\postcommunion\label{MarcusMarcellianusPostcommunion}
\lett{O}{Lord,} who hast satisfied us with the gift of thy salvation, we humbly beseech thee: that as we joyfully taste thereof, so at the intercession of thy holy Martyrs Marcus and Marcellianus, we may effectually be renewed by the same. Through.


\subby{Sts. Gervase \& Protase}
\feastday{{Sts. Gervase \& Protase}}
\fancyhead[RE,LO]{19 June}
\begin{inhead}
    {Memorial\\
19 June}
\end{inhead}

\begin{rubric}
	In Eastertide, the propers are from the Common of Many Martyrs in Eastertide (p. \pageref{CommonMartyrsEaster}), with the Prayers as followeth and the Greater Alleluia from the Feast of Sts. Primus \& Felician (p. \pageref{PrimusFelician}).
\end{rubric}

\begin{rubric}
	Outside of Eastertide, the propers are as followeth.
\end{rubric}

\introit
\lett{T}{he} Lord shall speak peace unto his people: and to his saints, that they turn not again. \textit{Ps.} Lord, thou art become gracious unto thy land, thou hast turned away the captivity of Jacob.

\collect
\lett{O}{God,} who makest us glad with the yearly solemnity of thy holy Martyrs Gervasius and Protasius: mercitully grant; that as we do rejoice in their merits, so we may be enkindled by their example. Through.

\begin{rubric}
	The Epistle is from the second additional Epistle from the Second Common of a Martyr not a Bishop (p. \pageref{1Peter413}).
\end{rubric}

\gradall{God is glorious in his holy ones: fearful in praises, doing wonders. ℣. Thy right hand, O Lord, is become glorious in power: thy right hand hath dashed in pieces the enemy.}{Alleluia, alleluia. ℣. This is the true brotherhood, which overcame the wickedness of the world: which followed Christ, gaining heaven's glorious realms. Alleluia.}

\begin{rubric}
	The Gospel is from the Second Common of Many Martyrs (p. \pageref{CommonMartyrsII}).
\end{rubric}

\offertory{Be glad, O ye righteous, and rejoice in the Lord: and be joyful, all ye that are true of heart.}

\secret
\lett{W}{e} beseech thee, O Lord, mercifully to accept these our oblations: that, at the intercession of thy holy Martyrs Gervasius and Protasius, we may be defended against all adversities. Through.

\communion{The dead bodies of thy servants, O Lord, have they given to be meat unto the fowls of the air, and the flesh of thy Saints unto the beasts of the land: according to the greatness of thy power, preserve thou those that are appointed to die.}

\postcommunion
\lett{M}{ay} this communion, O Lord, cleanse us from guilt: and at the intercession of thy holy Martyrs Gervasius and Protasius, make us partakers of thy heavenly healing. Through.


\subby{St. Silverius}
\feastday{{St. Silverius}}
\fancyhead[RE,LO]{20 June}
\begin{inhead}
    {Memorial\\
20 June}
\end{inhead}

\begin{rubric}
	The propers are from the First Common of a Martyr Bishop (p. \pageref{CommonMartyrBishopI}), with the Epistle as followeth.
\end{rubric}

\readingcitation{Epistle}{Jude 17}
\lett{D}{early beloved:} Remember ye the words which were spoken before of the apostles of our Lord Jesus Christ; How that they told you there should be mockers in the last time, who should walk after their own ungodly lusts. These be they who separate themselves, sensual, having not the Spirit. But ye, beloved, building up yourselves on your most holy faith, praying in the Holy Ghost, Keep yourselves in the love of God, looking for the mercy of our Lord Jesus Christ unto eternal life.


\subby{St. Alban}
\feastday{{St. Alban}}
\fancyhead[RE,LO]{22 June}
\begin{inhead}
    {Double\\
22 June}
\end{inhead}

\begin{rubric}
	The Daily Office propers are from the First Common of a Martyr not a Bishop (p. \pageref{CommonMartyrNotBishopI}).
\end{rubric}

\introit
\lett{T}{he} just shall rejoice in thy strength, O Lord: exceeding glad shall he be of thy salvation: thou hast given him his heart's desire.\textit{Ps.} For thou hast prevented him with the blessings of goodness: thou hast set a crown of pure gold upon his head.

\collect
\lett{O}{God,} who hast hallowed this day by the martyrdom of blessed Alban: grant, we beseech thee; that as year by year we rejoice to pay him honour, so we may be defended by his continual help. Through.

\begin{rubric}
	Commemoration of St. Paulinus (p. \pageref{PaulinusCollect}).
\end{rubric}

\readingcitation{Epistle}{Wisdom 10:10}
\lett{W}{isdom} guided the righteous man in straight paths; She shewed him God’s kingdom, and gave him knowledge of holy things; She prospered him in his toils, and multiplied the fruits of his labour; When in their covetousness men dealt hardly with him, She stood by him and made him rich; She guarded him from enemies, And from those that lay in wait she kept him safe, And over his sore conflict she watched as judge, That he might know that godliness is more powerful than all. When a righteous man was sold, wisdom forsook him not, But from sin she delivered him; She went down with him into a dungeon, And in bonds she left him not, Till she brought him the sceptre of a kingdom, And authority over those that dealt tyrannously with him; She shewed them also to be false that had mockingly accused him, And gave him eternal glory.

\gradall{Blessed is the man that feareth the Lord: he hath great delight in his commandments. ℣. His seed shall be mighty upon earth: the generation of the faithful shall be blessed.}{Alleluia, alleluia. ℣. Thou hast set, O Lord, a crown of pure gold upon his head. Alleluia.}

\readingcitation{Gospel}{Matthew 16:24}
\lett{A}{t that time:} Jesus said unto his disciples: If any man will come after me, let him deny himself, and take up his cross, and follow me. For whosoever will save his life shall lose it: and whosoever will lose his life for my sake shall find it. For what is a man profited, if he shall gain the whole world, and lose his own soul? or what shall a man give in exchange for his soul? For the Son of man shall come in the glory of his Father with his angels; and then he shall reward every man according to his works.

\offertory{Thou hast crowned him with glory and worship, and hast made him to have dominion of the works of thy hands, O Lord.}

\secret
\lett{W}{e} beseech thee, O Lord, that like as in the veneration of blessed Alban, thy Martyr, we do shew forth thy wonders: so through this bounden service of propitiation he may be a faithful intercessor for us in the sight of thy mercy. Through.

\begin{rubric}
	Commemoration of St. Paulinus (p. \pageref{PaulinusSecret}).
\end{rubric}

\communion{If any man will come after me, let him deny himself, and take up his cross and follow me.}

\postcommunion
\lett{L}{et} blessed Alban, thy Martyr, we beseech thee, O Lord, ever implore thy holy majesty: that these thy sacraments may cleanse us from guilt, and preserve in us the fervour of thy charity. Through.

\begin{rubric}
	Commemoration of St. Paulinus (p. \pageref{PaulinusPostcommunion}).
\end{rubric}


\subby{St. Paulinus}
\feastday{{St. Paulinus}}
\fancyhead[RE,LO]{22 June}
\begin{inhead}
    {Memorial\\
22 June}
\end{inhead}

\introit
\lett{L}{et} thy priests, O Lord, be clothed with righteousness, and let thy saints sing with joyfulness: for thy servant David's sake, turn not away the presence of thine anointed. \textit{Ps.} Lord, remember David, and all his trouble.

\collect\label{PaulinusCollect}
\lett{O}{God,} who to them that forsake all things in this world for thee hast promised a hundredfold in the world to come, and life eternal: mercifully grant; that, following in the footsteps of thy holy Bishop, Paulinus, we may be enabled to despise earthly things, and to seek only after things heavenly: Who livest.

\begin{rubric}
	If to-day be Saturday, a Commemoration is made of the anticipated Vigil of St. John Baptist, as below.
\end{rubric}

\readingcitation{Epistle}{2 Corinthians 8:9}
\lett{B}{rethren:} Ye know the grace of our Lord Jesus Christ, that, though he was rich, yet for your sakes he became poor, that ye through his poverty might be rich. And herein I give my advice: for this is expedient for you, who have begun before, not only to do, but also to be forward a year ago. Now therefore perform the doing of it; that as there was a readiness to will, so there may be a performance also out of that which ye have. For if there be first a willing mind, it is accepted according to that a man hath, and not according to that he hath not. For I mean not that other men be eased, and ye burdened: But by an equality, that now at this time your abundance may be a supply for their want, that their abundance also may be a supply for your want: that there may be equality: As it is written, He that had gathered much had nothing over; and he that had gathered little had no lack.

\gradall{Behold a great priest who in his days pleased God. ℣. There was none found like unto him, who kept the law of the Most High.}{Alleluia, alleluia. ℣. Thou art a priest for ever after the order of Melchisedech. Alleluia.}

\begin{rubric}
	The Gospel is from the Second Common of a Confessor not a Bishop (p. \pageref{CommonConfessorNotBishopII}).
\end{rubric}

\offertory{I have found David my servant, with my holy oil have I anointed him: my hand shall hold him fast, and my arm shall strengthen him.}

\secret\label{PaulinusSecret}
\lett{G}{rant,} O Lord, that, after the example of thy holy Bishop Paulinus, we may join the sacrifice of perfect charity to the oblation of the altar: and by zeal in well-doing be found worthy of thine everlasting mercy. Through.

\communion{A faithful and wise servant whom the lord hath made ruler over his household, to give them their portion of meat in due season.}

\postcommunion\label{PaulinusPostcommunion}
\lett{G}{rant} us, O Lord, through these holy things that love of piety and humility which thy holy Bishop Paulinus drew from this divine fountain: and by his intercession graciously pour forth on all who entreat thee the riches of thy grace. Through.


\subby{Vigil of St. John Baptist}
\feastday{{St. John Baptist Vigil}}
\fancyhead[RE,LO]{23 June}
\begin{inhead}
    {Vigil\\
23 June}
\end{inhead}

\introit
\lett{F}{ear} not, Zacharias, thy prayer is heard: and thy wife Elisabeth shall bear thee a son, and thou shalt call his name John: he shall be great in the sight of the Lord: and shall be filled with the Holy Ghost, even from his mother's womb: and many shall rejoice at his birth. \textit{Ps.} The king shall rejoice in thy strength, O Lord: exceeding glad shall he be of thy salvation.

\collect
\lett{G}{rant,} we beseech thee, almighty God: that thy family may walk in the way of salvation; and, following the teachings of blessed John the Forerunner, may attain in safety unto him whom he fore-told, Jesus Christ thy Son, our Lord: Who liveth.

\begin{rubric}
The \nth{2} is of St. Mary and the \nth{3} against the persecutors of the Church or for the Chief Bishop.
\end{rubric}

\readingcitation{Epistle}{Jeremiah 1:4}
\lett{I}{n those days:} The word of the \divineName{Lord} came unto me, saying, Before I formed thee in the belly I knew thee; and before thou camest forth out of the womb I sanctified thee, and I ordained thee a prophet unto the nations. Then said I, Ah, Lord \divineName{God}! behold, I cannot speak: for I am a child. But the \divineName{Lord} said unto me, Say not, I am a child: for thou shalt go to all that I shall send thee, and whatsoever I command thee thou shalt speak. Be not afraid of their faces: for I am with thee to deliver thee, saith the \divineName{Lord}. Then the \divineName{Lord} put forth his hand, and touched my mouth. And the \divineName{Lord} said unto me, Behold, I have put my words in thy mouth. See, I have this day set thee over the nations and over the kingdoms, to root out, and to pull down, and to destroy, and to throw down, to build, and to plant: saith the Lord almighty.

\gradual{There was a man sent from God, whose name was John. ℣. The same came to bear witness of the light, to make ready a people prepared for the Lord.}

\readingcitation{Gospel}{Luke 1:5}
\lett{T}{here} was in the days of Herod, the king of Judaea, a certain priest named Zacharias, of the course of Abia: and his wife was of the daughters of Aaron, and her name was Elisabeth. And they were both righteous before God, walking in all the commandments and ordinances of the Lord blameless. And they had no child, because that Elisabeth was barren, and they both were now well stricken in years. And it came to pass, that while he executed the priest's office before God in the order of his course, According to the custom of the priest's office, his lot was to burn incense when he went into the temple of the Lord. And the whole multitude of the people were praying without at the time of incense. And there appeared unto him an angel of the Lord standing on the right side of the altar of incense. And when Zacharias saw him, he was troubled, and fear fell upon him. But the angel said unto him, Fear not, Zacharias: for thy prayer is heard; and thy wife Elisabeth shall bear thee a son, and thou shalt call his name John. And thou shalt have joy and gladness; and many shall rejoice at his birth. For he shall be great in the sight of the Lord, and shall drink neither wine nor strong drink; and he shall be filled with the Holy Ghost, even from his mother's womb. And many of the children of Israel shall he turn to the Lord their God. And he shall go before him in the spirit and power of Elias, to turn the hearts of the fathers to the children, and the disobedient to the wisdom of the just; to make ready a people prepared for the Lord.

\offertory{Thou hast crowned him with glory and worship: thou hast made him to have dominion of the works of thy hands, O Lord.}

\secret
\lett{S}{anctify,} O Lord, the gifts which we ofter: and, at the intercession of blessed John Baptist, cleanse us thereby from the defilements of our iniquities. Through.

\begin{rubric}
The \nth{2} is of St. Mary and the \nth{3} against the persecutors of the Church or for the Chief Bishop.
\end{rubric}

\communion{His honour is great in thy salvation: glory and great worship shalt thou lay upon him, O Lord.}

\postcommunion
\lett{O}{Lord,} let the glorious prayer of blessed John Baptist ever be with us: and may he implore for us the mercy of him, whose coming he fore-told, Jesus Christ thy Son our Lord: Who liveth.

\begin{rubric}
The \nth{2} is of St. Mary and the \nth{3} against the persecutors of the Church or for the Chief Bishop.
\end{rubric}


\subby{St. Etheldreda}
\feastday{{St. Etheldreda}}
\fancyhead[RE,LO]{23 June}
\begin{inhead}
    {Memorial\\
23 June}
\end{inhead}

\begin{rubric}
	The propers are from the First Common of a Virgin (p. \pageref{CommonVirginOnlyI}), with the Collect as followeth.
\end{rubric}

\collect
\lett{O}{God,} who makest us glad with the yearly solemnity of blessed Etheldreda, thy Virgin: mercifully grant; that as we are enlightened by the example of her purity, so we may be succoured by her prayers. Through.


\bcpfeast{24 June. St. John Baptist}{St. John Baptist}{24 June}
%\supplement{24 June}{St. John Baptist}{}

\subbysub{I Evensong}\label{JohnBaptistEvensong}

\gregorioscore{resources/gabc/ProperTime/JohnBaptistEvensong.gabc}

	℣. There was a man sent from God.
	
	℟. Whose name was John.

\properantiphon{Mag.}{When Zacharias {\dag} went into the temple, there appeared unto him the Angel Gabriel, standing on the right side of the altar of incense.}

\subbysub{Mattins}

\invitatoryhymn\label{JohnBaptistInvitatory}

\gregorioscore{resources/gabc/ProperTime/JohnBaptistInvitatory.gabc}

\officehymn\label{JohnBaptistMattins}

\gregorioscore{resources/gabc/ProperTime/JohnBaptistMattins.gabc}

	℣. This child shall be great in the sight of the Lord.
	
	℟. For the hand of the Lord is with him.

\properantiphon{Ben.}{The mouth of Zacharias {\dag} was opened, and he prophesied, saying: Blessed be the God of Israel.}

\subbysub{II Evensong}

\begin{rubric}
	The Office Hymn is of I Evensong, with the following Versicle \& Antiphon.
\end{rubric}

	℣. This child shall be great in the sight of the Lord.
	
	℟. For the hand of the Lord is with him.

\properantiphon{Mag.}{The child {\dag} that is born unto us is more than a prophet; for this is he of whom the Saviour saith: Among them that are born of women, there hath not risen a greater than John the Baptist.}


\subby{Sts. John \& Paul}
\feastday{{Sts. John \& Paul}}
\fancyhead[RE,LO]{26 June}
\begin{inhead}
    {Double\\
26 June}
\end{inhead}

\begin{rubric}
	The Daily Office propers are from the First Common of Many Martyrs (p. \pageref{CommonMartyrsI}).
\end{rubric}

\introit
\lett{G}{reat} are the troubles of the righteous, but the Lord delivereth him out of all: the Lord keepeth all his bones: so that not one of them is broken. \textit{Ps.} I will alway give thanks unto the Lord: his praise shall ever be in my mouth.

\collect
\lett{W}{e} beseech thee, almighty God: that, as their faith and passion did cause blessed John and Paul to be brothers indeed; so this day's festival may bestow upon us a twofold gladness in their glory. Through.

\begin{rubric}
	Commemoration of the Octave of St. John Baptist, as followeth.
\end{rubric}

\lett{A}{lmighty} God, by whose providence thy servant John Baptist was wonderfully born, and sent to prepare the way of thy Son our Saviour by preaching repentance; Make us so to follow his doctrine and holy life, that we may truly repent according to his preaching; and after his example constantly speak the truth, boldly rebuke vice, and patiently suffer for the truth's sake. Through the same.

\readingcitation{Epistle}{Ecclesiasticus 44:10}
\lett{T}{hese} were men of mercy, Whose righteous deeds have not been forgotten. With their seed shall remain continually a good inheritance; Their children are within the covenants. Their seed standeth fast, And their children for their sakes. Their seed shall remain for ever, And their glory shall not be blotted out. Their bodies were buried in peace, And their name liveth to all generations. Peoples will declare their wisdom, And the congregation telleth out their praise. 

\gradall{Behold, how good and joyful a thing it is, brethren, to dwell together in unity! ℣. It is like the precious ointment upon the head, that ran down unto the beard, even Aaron's beard.}{Alleluia, alleluia. ℣. This is the true brotherhood, which overcame the wickedness of the world: which followed Christ, gaining heaven's glorious realms. Alleluia.}

\begin{rubric}
	The Gospel is from the Third Common of Many Martyrs (p. \pageref{CommonMartyrsIII}).
\end{rubric}

\offertory{They that love thy name shall be joyful in thee, for thou, Lord, wilt give thy blessing unto the righteous: and with thy favourable kindness wilt thou defend him as with a shield, O Lord.}

\secret
\lett{W}{e} beseech thee, O Lord, mercifully to accept this our sacrifice which we offer unto thee, pleading the merits of thy holy Martyrs John and Paul: that the same may avail for our perpetual succour. Through.

\begin{rubric}
	Commemoration of the Octave of St. John Baptist, as followeth.
\end{rubric}

\lett{W}{e} set upon thine altars, O Lord, these gifts: celebrating with due honour the nativity of him, who sang of the coming and proclaimed the presence of the Saviour of the world, Jesus Christ thy Son our Lord. Who liveth.

\communion{Though they be punished in the sight of men, God proved them: as gold in the furnace hath he tried them, and received them as a burnt-offering.}

\postcommunion
\lett{G}{rant,} we beseech thee, O Lord: that we who have received these heavenly sacraments, in celebration of the festival of thy holy Martyrs John and Paul; may attain in everlasting joys the fulfilment of our service in this life. Through.

\begin{rubric}
	Commemoration of the Octave of St. John Baptist, as followeth.
\end{rubric}

\lett{L}{et} thy Church, O God, rejoice at the birth of blessed John Baptist: through whom she hath known the author of her new birth, Jesus Christ thy Son our Lord. Who liveth.

%MANUAL ADJUSTMENT:
\clearpage

\subby{Vigil of St. Peter}
\feastday{{St. Peter Vigil}}
\fancyhead[RE,LO]{28 June}
\begin{inhead}
    {Vigil\\
28 June}
\end{inhead}

\introit
\lett{T}{he} Lord saith unto Peter: When thou wast young, thou girdedst thyself, and walkedst whither thou wouldest: but when thou shalt be old, thou shalt stretch forth thy hands, and another shall gird thee, and carry thee whither thou wouldest not: this spake he, signifying by what death he should glorify God. \textit{Ps.} The heavens declare the glory of God: and the firmament sheweth his handy-work.

\collect
\lett{G}{rant,} we beseech thee, almighty God: that as thou hast stablished us on the rock of the apostolic confession; so thou wouldest not suffer us to be troubled by any adversities. Through.

\readingcitation{Epistle}{Acts 3:1}
\lett{I}{n those days:} Peter and John went up together into the temple at the hour of prayer, being the ninth hour. And a certain man lame from his mother's womb was carried, whom they laid daily at the gate of the temple which is called Beautiful, to ask alms of them that entered into the temple; Who seeing Peter and John about to go into the temple asked an alms. And Peter, fastening his eyes upon him with John, said, Look on us. And he gave heed unto them, expecting to receive something of them. Then Peter said, Silver and gold have I none; but such as I have give I thee: In the name of Jesus Christ of Nazareth rise up and walk. And he took him by the right hand, and lifted him up: and immediately his feet and ankle bones received strength. And he leaping up stood, and walked, and entered with them into the temple, walking, and leaping, and praising God. And all the people saw him walking and praising God: And they knew that it was he which sat for alms at the Beautiful gate of the temple: and they were filled with wonder and amazement at that which had happened unto him.

\gradual{Their sound is gone out into all lands: and their words into the ends of the world. ℣. The heavens declare the glory of God: and the firmament sheweth his handy-work.}

%RV, due to textual variant issue:
\readingcitation{Gospel}{John 21:15}
\lett{A}{t that time:} Jesus said unto Simon Peter: Simon, son of John, lovest thou me more than these? He saith unto him, Yea, Lord; thou knowest that I love thee. He saith unto him, Feed my lambs. He saith to him again a second time, Simon, son of John, lovest thou me? He saith unto him, Yea, Lord; thou knowest that I love thee. He saith unto him, Tend my sheep. He saith unto him the third time, Simon, son of John, lovest thou me? Peter was grieved because he said unto him the third time, Lovest thou me? And he said unto him, Lord, thou knowest all things; thou knowest that I love thee. Jesus saith unto him, Feed my sheep. Verily, verily, I say unto thee, When thou wast young, thou girdedst thyself, and walkedst whither thou wouldest: but when thou shalt be old, thou shalt stretch forth thy hands, and another shall gird thee, and carry thee whither thou wouldest not. Now this he spake, signifying by what manner of death he should glorify God. 

\offertory{Right honourable are thy friends unto me, O God: right well is their princedom established.}

\secret
\lett{S}{anctify,} the gift of thy people, we beseech thee, O Lord, by the intercession of thine Apostles: and cleanse us from the defilements of our iniquities. Through.

\communion{Simon, son of Jonas, lovest thou me more than these? Lord, thou knowest all things: thou knowest, Lord, that I love thee.}

\postcommunion
\lett{O}{Lord,} who hast satisfied us with heavenly food: defend us by the intercession of thine Apostles against all adversity. Through.



\subby{Pope St. Leo II}
\feastday{{Pope St. Leo II}}
\fancyhead[RE,LO]{28 June}
\begin{inhead}
    {Memorial\\
28 June}
\end{inhead}

\begin{rubric}
	The propers are from the First Common of a Confessor Bishop (p. \pageref{CommonConfessorBishopI}).
\end{rubric}


\bcpfeast{29 June. St. Peter}{St. Peter}{29 June}
%\supplement{29 June}{St. Peter}{}

\subbysub{I Evensong}\label{PeterEvensong}

\gregorioscore{resources/gabc/ProperTime/PeterEvensong.gabc}

	℣. Their sound is gone out into all lands.
	
	℟. And their words into the ends of the world.

\properantiphon{Mag.}{Thou art the shepherd of the sheep, {\dag} O chief of the Apostles: unto thee were given the keys of the kingdom of heaven.}

\subbysub{Mattins}

\begin{rubric}
	The Invitatory Hymn is of the Common of Apostles (p. \pageref{CommonApostles}).
\end{rubric}

\officehymn\label{PeterMattins}

\gregorioscore{resources/gabc/ProperTime/PeterMattins.gabc}

	℣. They declared the work of God.
		
	℟. And wisely considered of his doing.

\properantiphon{Ben.}{Whatsoever {\dag} thou shalt bind on earth shall be bound in heaven: and whatsoever thou shalt loose on earth shall be loosed in heaven: saith the Lord unto Simon Peter.}

\subbysub{II Evensong}

\begin{rubric}
	The Office Hymn is of I Evensong, with the following Versicle \& Antiphon.
\end{rubric}

	℣. They declared the work of God.
		
	℟. And wisely considered of his doing.

\properantiphon{Mag.}{To-day {\dag} Simon Peter ascended the gibbet of the cross, alleluia: to-day the Key-Bearer of the Kingdom joyfully departed to Christ: to-day Paul the Apostle, the light of the whole world, bowed down his head; and for the Name of Christ, received the crown of martyrdom, alleluia.}


\bcpfeast{30 June. Commemoration of St. Paul}{St. Paul Commemoration}{30 June}
%\supplement{30 June}{Commemoration of St. Paul}{}

\subbysub{I Evensong}\label{PaulEvensong}

\gregorioscore{resources/gabc/ProperTime/PaulEvensong.gabc}

	℣. Thou art a chosen vessel, holy Apostle Paul.
		
	℟. A preacher of the truth throughout all the world.

\properantiphon{Mag.}{O holy Apostle Paul, {\dag} thou preacher of the truth and Doctor of the Gentiles, intercede for us unto God, who hath chosen thee.}

\subbysub{Mattins}

\begin{rubric}
	The Invitatory Hymn is as in I Evensong.
\end{rubric}

\begin{rubric}
	The Office Hymn is as in the Common of Apostles (p. \pageref{CommonApostles}).
\end{rubric}

\begin{rubric}
	In Eastertide, the Office Hymn is as in the Common of Apostles in Eastertide (p. \pageref{CommonApostlesEaster}).
\end{rubric}

	℣. Thou art a chosen vessel, holy Apostle Paul.
		
	℟. A preacher of the truth throughout all the world.

\properantiphon{Ben.}{Ye which have followed me {\dag} shall sit upon twelve thrones, judging the twelve tribes of Israel, saith the Lord.}

\subbysub{II Evensong}

\begin{rubric}
	II Evensong is as in I Evensong.
\end{rubric}

%MANUAL ADJUSTMENT:
\clearpage
\bcpfeast{1 July. Most Precious Blood of Our Lord Jesus Christ}{Precious Blood}{1 July}
%\supplement{1 July}{Most Precious Blood}{of Our Lord Jesus Christ}

\subbysub{I Evensong}

\gregorioscore{resources/gabc/ProperTime/PreciousBloodEvensong.gabc}\label{PreciousBloodEvensong}

	℣. Thou hast redeemed us, O Lord, by thy Blood.
		
	℟. And hast made us unto our God a kingdom.

\properantiphon{Mag.}{But ye are come {\dag} unto mount Sion and unto the city of the living God, the heavenly Jerusalem, and to Jesus, the mediator of the new covenant, and to the Blood of sprinkling, that speaketh better things than that of Abel.}

\subbysub{Mattins}

\invitatoryhymn

\gregorioscore{resources/gabc/ProperTime/PreciousBloodInvitatory.gabc}\label{PreciousBloodInvitatory}

\officehymn

\gregorioscore{resources/gabc/ProperTime/PreciousBloodMattins.gabc}\label{PreciousBloodMattins}

	℣. Being justified by the Blood of Christ.
		
	℟. We shall be saved from wrath through him.

\properantiphon{Ben.}{The Blood of the Lamb {\dag} shall be to you for a token, saith the Lord: and when I see the Blood, I will pass over you, and the plague shall not be upon you to destroy you.}

\subbysub{II Evensong}

\begin{rubric}
	II Evensong is as in I Evensong, except the following Antiphon.
\end{rubric}

\properantiphon{Mag.}{And this day {\dag} shall be unto you for a memorial: and ye shall keep it a feast to the Lord throughout your generations by an ordinance for ever.}


\bcpfeast{2 July. Visitation of the Blessed Virgin Mary}{Visitation}{2 July}
%\supplement{2 July}{Visitation}{of the B.V.M.}

\begin{rubric}
	The Hymns are as in the Common of the Blessed Virgin Mary (p. \pageref{CommonBVM}), with the following Versicles \& Antiphons.
\end{rubric}

	℣. Blessed art thou amongst women.
		
	℟. And blessed is the fruit of thy womb.

\properantiphon{Mag.}{Blessed art thou, {\dag} O Mary, for thou hast believed: and there shall be a performance in thee of those things which were told thee from the Lord, alleluia.}\\

	℣. Blessed art thou amongst women.
		
	℟. And blessed is the fruit of thy womb.

\properantiphon{Ben.}{When Elisabeth {\dag} heard the salutation of Mary, she spake out with a loud voice, and said: Whence is this to me, that the mother of my Lord should come to me? Alleluia.}\\

	℣. Blessed art thou amongst women.
		
	℟. And blessed is the fruit of thy womb.

\properantiphon{Mag.}{All generations {\dag} shall call me blessed: for God hath regarded the lowliness of his handmaiden, alleluia.}


\subby{St. John of San Francisco}
\feastday{{St. John San Francisco}}
\fancyhead[RE,LO]{2 July}
\begin{inhead}
    {Memorial\\
2 July}
\end{inhead}

\begin{rubric}
	The propers are from the Second Common of a Confessor Bishop (p. \pageref{CommonConfessorBishopII}).
\end{rubric}


\subby{Sts. Processus \& Martinian}
\feastday{{St. Processus}}
\fancyhead[RE,LO]{2 July}
\begin{inhead}
    {Memorial\\
2 July}
\end{inhead}

\begin{rubric}
	The propers are from the Second Common of a Confessor Bishop (p. \pageref{CommonConfessorBishopII}).
\end{rubric}


\subby{St. Iren{\ae}us of Lyon}
\feastday{{St. Iren{\ae}us}}
\fancyhead[RE,LO]{3 July}
\begin{inhead}
    {Double\\
3 July}
\end{inhead}

\begin{rubric}
	The propers are from the First Common of a Martyr Bishop (p. \pageref{CommonMartyrBishopI}), except for the following Collect.
\end{rubric}

\collect
\lett{O}{God,} who gavest grace to blessed Iren{\ae}us thy Martyr and Bishop to overcome false doctrine by the truth of his teaching, and favourably to establish peace in thy Church: give thy people, we beseech thee, constancy in holy religion, and grant us thy peace all the days of our life. Through.


%Last Day of Eastertide Possible---
\subby{Sts. Cyril \& Methodius}
\feastday{{St. Cyril}}
\fancyhead[RE,LO]{7 July}
\begin{inhead}
    {Double\\
7 July}
\end{inhead}

\begin{rubric}
	The propers are from the Second Common of a Confessor Bishop (p. \pageref{CommonConfessorBishopII}), except for the following.
\end{rubric}

\collect
\lett{A}{lmighty} and everlasting God, who by thy blessed Confessors and Bishops, Cyril and Methodius, didst suffer the peoples of Slavonia to come to the knowledge of thy name: vouchsafe; that we, who glory in their festival may be joined unto their fellowship. Through.

\readingcitation{Gospel}{Luke 10:1}
\lett{A}{t that time:} the Lord appointed other seventy also: and sent them two and two before his face into every city and place, whither he himself would come. Therefore said he unto them, The harvest truly is great, but the labourers are few: pray ye therefore the Lord of the harvest, that he would send forth labourers into his harvest. Go your ways: behold, I send you forth as lambs among wolves. Carry neither purse, nor scrip, nor shoes: and salute no man by the way. And into whatsoever house ye enter, first say, Peace be to this house. And if the son of peace be there, your peace shall rest upon it: if not, it shall turn to you again. And in the same house remain, eating and drinking such things as they give: for the labourer is worthy of his hire. Go not from house to house. And into whatsoever city ye enter, and they receive you, eat such things as are set before you: And heal the sick that are therein, and say unto them, The kingdom of God is come nigh unto you.

\offertory{God is wonderful in his holy ones: even the God of Israel, he will give strength and power unto his people: blessed be God.}

\secret
\lett{W}{e} beseech thee, O Lord, to have respect unto the prayers and oblations of thy faithful people: that they may be acceptable unto thee on this festival of thy Saints, and effectually bestow on us the assistance of thy mercy. Through.

\communion{What I tell you in darkness, that speak ye in light, saith the Lord: and what ye hear in the ear, that preach ye upon the house-tops.}

\postcommunion
\lett{W}{e} beseech thee, almighty God: that as thou dost vouchsafe to bestow upon us heavenly gifts, so, at the intercession of thy Saints Cyril and Methodius, thou wouldest grant unto us to despise things earthly. Through.


\subby{Seven Holy Brothers, with Sts. Rufina \& Secunda}
\feastday{{7 Holy Brothers}}
\fancyhead[RE,LO]{10 July}
\begin{inhead}
    {Memorial\\
10 July}
\end{inhead}

\introit
\lett{P}{raise} the Lord, ye servants, O praise the name of the Lord: who maketh the barren woman to keep house, and to be a joyful mother of children. \textit{Ps.} Blessed be the name of the Lord: from this time forth for evermore.

\collect
\lett{G}{rant,} we beseech thee, almighty God: that, like as we have known thy glorious Martyrs to be constant in their confession, so we may perceive their loving intercession. Through.

\begin{rubric}
	The Epistle is from the Common of Neither Virgin Nor Martyr (p. \pageref{CommonNeitherVirginMartyr}).
\end{rubric}

\gradall{Our soul is escaped, even as a bird out of the snare of the fowler. ℣. The snare is broken, and we are delivered : our help is in the name of the Lord, who hath made heaven and earth.}{Alleluia, alleluia. ℣. This is the true brotherhood, which overcame the wickedness of the world: which followed Christ, gaining heaven's glorious realms. Alleluia.}

\readingcitation{Gospel}{Matthew 12:46}
\lett{A}{t that time:} While Jesus spake to the multitudes, behold his Mother and brethren stood without, desiring to speak with him. Then one said unto him, Behold, thy mother and thy brethren stand without, desiring to speak with thee. But he answered and said unto him that told him, Who is my mother? and who are my brethren? And he stretched forth his hand toward his disciples, and said, Behold my mother and my brethren! For whosoever shall do the will of my Father which is in heaven, the same is my brother, and sister, and mother.

\offertory{Our soul is escaped, even as a bird out of the snare of the fowler: the snare is broken, and we are delivered.}

\secret
\lett{W}{e} beseech thee, O Lord, mercifully to have respect unto these our sacrifices: that through the intercession of thy Saints, they may increase our devotion, and set forward our salvation. Through.

\communion{Whosoever shall do the will of my Father which is in heaven: the same is my brother, and sister, and mother, saith the Lord.}

\postcommunion
\lett{W}{e} beseech thee, almighty God: that as in these holy mysteries we have received the pledge of our salvation, so at the intercession of thy Saints, we may be brought unto the fulfilment of the same. Through.


\subby{St. Joseph of Damascus (Priest) \& Companions}
\feastday{{St. Joseph Damascene \&c.}}
\fancyhead[RE,LO]{10 July}
\begin{inhead}
    {Memorial\\
10 July}
\end{inhead}

\begin{rubric}
	The propers are from the Third Common of Many Martyrs (p. \pageref{CommonMartyrsIII}).
\end{rubric}


\subby{Pope St. Pius I}
\feastday{{Pope St. Pius I}}
\fancyhead[RE,LO]{11 July}
\begin{inhead}
    {Memorial\\
11 July}
\end{inhead}

\begin{rubric}
	The propers are from the First Common of a Martyr Bishop (p. \pageref{CommonMartyrBishopI}).
\end{rubric}


\subby{Sts. Nabor \& Felix}
\feastday{{Sts. Nabor \& Felix}}
\fancyhead[RE,LO]{12 July}
\begin{inhead}
    {Memorial\\
12 July}
\end{inhead}

\begin{rubric}
	The propers come from the Third Common of Many Martyrs (p. \pageref{CommonMartyrsIII}), except for the following.
\end{rubric}

\collect
\lett{G}{rant,} we beseech thee, O Lord: that as the birthday of thy holy Martyrs, Nabor and Felix, faileth not to return for our observance; so they may continually assist us by their prayers. Through.

\secret
\lett{W}{e} beseech thee, O Lord, that the gifts of thy people may, by the prayers of thy holy Martyrs, Nabor and Felix, be made acceptable unto thee: that, as they are offered for their triumph to thy name, so by their merits they may be rendered worthy in thy sight. Through.

\postcommunion
\lett{O}{Lord,} who on the birthday of thy Saints hast quickened us with the gift of thy sacrament: we beseech thee; that, as by thy grace we now receive the comfort of thy blessings, so we may be brought to the everlasting fruition of the same. Through.


\subby{St. Anacletus}
\feastday{{St. Anacletus}}
\fancyhead[RE,LO]{13 July}
\begin{inhead}
    {Memorial\\
13 July}
\end{inhead}

\begin{rubric}
	The propers are from the Second Common of a Martyr Bishop (p. \pageref{CommonMartyrBishopII}), with the Gospel from the First Common of a Martyr Bishop (p. \pageref{CommonMartyrBishopI}).
\end{rubric}


\subby{St. Vladimir of Kiev}
\feastday{{St. Vladimir}}
\fancyhead[RE,LO]{15 July}
\begin{inhead}
    {Double\\
15 July}
\end{inhead}

\begin{rubric}
	The propers are from the First Common of a Confessor not a Bishop (p. \pageref{CommonConfessorNotBishopI}).
\end{rubric}


\subby{Sts. Nicholas \& Habib Khasha}
\feastday{Sts. Nicholas \& Habib Khasha}
\fancyhead[RE,LO]{16 July}
\begin{inhead}
    {Memorial\\
16 July}
\end{inhead}

\begin{rubric}
	The propers are from the Second Common of Many Martyrs (p. \pageref{CommonMartyrsII}).
\end{rubric}


%\subby{Translation of St. Swithun}
%\feastday{{St. Swithun}}
%\fancyhead[RE,LO]{15 July}
%\begin{inhead}
%    {Memorial\\
%15 July}
%\end{inhead}


%\subby{Our Lady of Einsiedeln}
%\feastday{{Einsiedeln}}
%\fancyhead[RE,LO]{16 July}
%\begin{inhead}
%    {Greater Double\\
%16 July}
%\end{inhead}

\subby{St. Alexius}
\feastday{{St. Alexius}}
\fancyhead[RE,LO]{17 July}
\begin{inhead}
    {Memorial\\
17 July}
\end{inhead}

\begin{rubric}
	The propers are from the First Common of a Confessor not a Bishop (p. \pageref{CommonConfessorNotBishopI}), except for the following.
\end{rubric}

\readingcitation{Epistle}{1 Timothy 6:6}
\lett{D}{early beloved:} Godliness with contentment is great gain. For we brought nothing into this world, and it is certain we can carry nothing out. And having food and raiment let us be therewith content. But they that will be rich fall into temptation and a snare, and into many foolish and hurtful lusts, which drown men in destruction and perdition. For the love of money is the root of all evil: which while some coveted after, they have erred from the faith, and pierced themselves through with many sorrows. But thou, O man of God, flee these things; and follow after righteousness, godliness, faith, love, patience, meekness. Fight the good fight of faith, lay hold on eternal life.

\readingcitation{Gospel}{Matthew 19:27}
\lett{A}{t that time:} Peter said unto Jesus: Behold, we have forsaken all, and followed thee; what shall we have therefore? And Jesus said unto them, Verily I say unto you, That ye which have followed me, in the regeneration when the Son of man shall sit in the throne of his glory, ye also shall sit upon twelve thrones, judging the twelve tribes of Israel. And every one that hath forsaken houses, or brethren, or sisters, or father, or mother, or wife, or children, or lands, for my name's sake, shall receive an hundredfold, and shall inherit everlasting life.


\subby{Translation of St. Raphael of Brooklyn}
\feastday{{St. Raphael Brooklyn Translation}}
\fancyhead[RE,LO]{18 July}
\begin{inhead}
    {Double\\
18 July}
\end{inhead}

\begin{rubric}
	The propers are from the First Common of a Confessor Bishop (p. \pageref{CommonConfessorBishopI}).
\end{rubric}


\subby{St. Sergius of Radonezh}
\feastday{{St. Sergius Radonezh}}
\fancyhead[RE,LO]{18 July}
\begin{inhead}
    {Memorial\\
18 July}
\end{inhead}

\begin{rubric}
	The propers are from the Common of Abbots (p. \pageref{CommonAbbots}).
\end{rubric}


%The rubrics from the English Missal are confusing here, so I use the Anglican Missal for rubrics and translation of the Prayers.
\subby{Sts. Symphorosa \& Her Seven Sons}
\feastday{{Sts. Symphorosa \&c.}}
\fancyhead[RE,LO]{18 July}
\begin{inhead}
    {Memorial\\
18 July}
\end{inhead}

\introit
\lett{T}{he} just cry, and the Lord heareth them: and delivereth them out of all their troubles. \textit{Ps.} I will alway give thanks unto the Lord; his praise shall ever be in my mouth.

\collect
\lett{O}{God,} who vouchsafest unto us to keep the heavenly birthday of thy holy Martyrs Symphorosa and her sons: grant, we beseech thee; that we may rejoice in the perpetual felicity of their friendship. Through.

\begin{rubric}
	The Epistle is the fifth additional Epistle of the Third Common of Many Martyrs (p. \pageref{Hebrews1133}).
\end{rubric}

\gradall{Behold, how good and joyful a thing it is, brethren, to dwell together in unity! ℣. It is like the precious ointment upon the head, that ran down unto the beard, even unto Aaron's beard.}{Alleluia, alleluia. ℣. This is the true brotherhood which overcame the wickedness of the world: which followed Christ, gaining heaven's glorious realms. Alleluia.}

\begin{rubric}
	The Gospel is from the Third Common of Many Martyrs (p. \pageref{CommonMartyrsIII}).
\end{rubric}

\offertory{Be glad, O ye righteous, and rejoice in the Lord: and be joyful, all ye that are true of heart.}

\secret
\lett{W}{e} beseech thee, O Lord, that the gifts which we offer unto thee of our bounden duty and service may be acceptable unto thee for the honour of thy Saints: and by thy mercy profitable unto us for our salvation. Through.

\begin{rubric}
	In Lent, Commemoration of the Feria.
\end{rubric}

\communion{Whosoever shall do the will of my Father which is in heaven: the same is my brother, and sister, and mother, saith the Lord.}

\postcommunion
\lett{G}{rant,} we beseech thee, almighty God: that, at the intercession of thy holy Martyrs Symphorosa and her sons, we, who with our outward lips have partaken of this Sacrament, may inwardly receive the same in purity of heart. Through.


\subby{St. Seraphim of Sarov}
\feastday{{St. Seraphim}}
\fancyhead[RE,LO]{19 July}
\begin{inhead}
    {Memorial\\
19 July}
\end{inhead}

\begin{rubric}
	The propers are from the First Common of a Confessor not a Bishop (p. \pageref{CommonConfessorNotBishopI}).
\end{rubric}


\subby{St. Elias the Prophet}
\feastday{{St. Elias}}
\fancyhead[RE,LO]{20 July}
\begin{inhead}
    {Double\\
20 July}
\end{inhead}\par\noindent
%Elijah changed to Elias \& Elisha to Eliseus.
%From the Carmelite Breviary and American Missal:
%Also from the `Saints of Mount Carmel' (https://archive.org/details/saintsofmountcar0000john/page/252/mode/2up):
%%From the Carmelite Breviary (1930s), translation from the KJV:
\textit{Opening Sentence.} Then stood up Elias the prophet as fire, and his word burned like a lamp, alleluia.%\vr{Ecclus. 48:1}\par

\begin{paracol}{2}[]
\sloppy
\begin{inhead}
	I Evensong
\end{inhead}
\begin{hangparas}{1.25em}{1}
The lofty peaks of Carmel

With tuneful praises ring,

The anthems of Elias

'Tis our delight to sing.\\

%Order changed to Churches:
The glory of our Churches,

Our leader, prop, and stay,

From east to west his offspring

Increaseth day by day.\\

When sorely press'd with famine,

A raven serv'd him bread,

With meal and cruse unfailing,

The widow'd hearth was fed.\\

The boy from death delivered

Is to his home restor'd,

And light so much desir'd,

In radiant flood is pour'd.\\

Behold the Heaven closeth,

To open at his voice,

And copious welcome showers

The thirsty lands rejoice.\\

To Father, Son, and Spirit,

Be equal power and praise,

All glory and dominion

Henceforth for endless days. Amen.\\
\end{hangparas}

    ℣. By the word of the Lord he shut up the heaven.

	℟. And he brought down fire from heaven thrice.

\antiphon{Mag.}{Behold, I {\dag}  will send you Elias the prophet before the coming of the great and dreadful day of the Lord: And he shall turn the heart of the fathers to the children, and the heart of the children to their fathers}\par\noindent

\switchcolumn

\begin{inhead}
	Mattins Invitatory Hymn
\end{inhead}
\begin{hangparas}{1.25em}{1}
Creator of the universe, our grateful hearts rejoice,

Exalting thee in this, the mighty Thesbite of thy choice;\\

Whose zeal for thy great glory, enkindled as a flame,

Defied the impious Prophets, and slew them in thy name.\\

His holocaust ascendeth, for thy fire may never fail,

While scorn and deep derision mock the clamorous priests of Baal.\\

The baneful rage of Jezabel he flieth in his dread,

And sleeping `neath a juniper is by an angel fed.\\

Empowered by the vision, with new courage he hath trod

Unto the heights of Horeb, unto the Mount of God.\\

O virtue of this bread divine, imparted by the Lord!

Full forty days he fasteth in the strength it doth afford.\\

All glory be to thee, O Father, Son, and Paraclete,

O undivided Trinity, thy praise all hearts repeat. Amen.
\end{hangparas}

\fussy
\end{paracol}


\begin{paracol}{2}[]
\sloppy

\begin{inhead}
	Mattins Office Hymn
\end{inhead}
\begin{hangparas}{1.25em}{1}
Come, blest companions, let our joy resounding

Extol to heav'n the leader of our line.

`Tis meet the mem'ry of his deeds abounding

Should waken ceaseless canticles divine.\\

He knows the gentle breathing of the Spirit

Cloth'd in the whistling murmur of the air,

By God's command the chastisements they merit

Proud Jezebel and Ahab justly share.\\

The caverns green of Carmel form his dwelling,

With leathern tunic is he rudely clad,

To impious Ahaziah his foretelling

Gives portent of a dissolution sad.\\

Twice at his pray'r the fire from Heav'n descending

Consumeth trembling soldiers in its flame,

The flowing wat'rs met with his mantle rending,

Dry shod he passeth safely through the same.\\

O Father, let thy help and thy protection

Be o'er thy children as they humbly plead,

Entreat the Spirit, by his sweet election,

To multiply his graces in their need.\\

O unbegotten Father, we adore thee,

O Son begotten, rev'rence be to thee,

O glorious Spirit, bow we low before thee,

Thou simple undivided Trinity. Amen.\\
\end{hangparas}

    ℣. Elias was covered with the whirlwind.

	℟. And his spirit was filled up in Eliseus.

\antiphon{Ben.}{Elias was a man {\dag} subject to like passions as we are, and he prayed earnestly that it might not rain: and it rained not on the earth by the space of three years and six months. And he prayed again, and the heaven gave rain, and the earth brought forth her fruit.}\par\noindent

\switchcolumn

\begin{inhead}
	II Evensong
\end{inhead}
\begin{hangparas}{1.25em}{1}
%Order changed to Churches
Thou prop of our Churches, thou pride of our race,

Let thy praises resound far and near,

Let sea and the land and the air give them place,

Rehearsing in gladness thy glory and grace,

Till the earth and the heav'ns give ear.\\

O sun of the heavens, how lovely thy rays!

What power thy wonders unfold,

How fruitful in merits the length of thy days,

Commission'd by God in His manifold ways,

For noble endeavours of old!\\

To regions celestial, in power and might,

Triumphant thy chariot speeds;

Uplifted by angels to marvellous height,

While shining in splendour and dazzling with light,

Thou guidest the fiery steeds.\\

As witness to men of his Sonship divine,

With Jesus thy glory we view;

The Father hath call'd thee on Tabor to shine,

Companion to Moses, and with him to sign

A testament faithful and true.\\

Protect us we pray, ‘neath thy powerful shield,

Incline to our aid from above,

Let thy fost'ring guidance be ever reveal'd

To thy children of Carmel, whose bosoms are seal'd

With the strength of thy fath'rly love\\

All power, dominion, all glory and praise,

Be given to Father and Son,

To thee, Holy Spirit, for numberless days

Our homage eternal we equally raise

All glory to God, Three in One. Amen.\\
\end{hangparas}

    ℣. Blessed are they that saw thee.

	℟. And were honoured with thy friendship.
	
%Based on KJV with changes based off the Latin.
\antiphon{Mag.}{And Elias took {\dag} his mantle, and wrapped it together, and smote the waters of the Jordan, and they were divided hither and thither, so that he and Eliseus went over on dry ground. As they still went on, behold, there appeared a chariot of fire, and horses of fire, and parted them both asunder; and Elias went up by a whirlwind into heaven, and Eliseus saw him no more.}

\fussy
\end{paracol}

%Tulit Elías {\dag} pállium suum, et invólvit illud, et percússit aquas jordanis, qua divisa sunt in utrámque partem, et transiérunt ipse et Eliséus per siccum. Qui incedentes, ecce, currus igneus et equi igner diviserunt utrúmque : et ascéndit Elías per túrbinem in c{\ae}lum, Eliséus autem non vidit eum ámplius.

\introit
\lett{I}{have} been very zealous for the Lord God of hosts: for the children of Israel have forsaken thy covenant, thrown down thine altars, and slain thy prophets with the sword; and I, even I only, am left; and they seek my life, to take it away. \textit{Ps.} I will magnify thee, O Lord, for thou hast set me up: and not made my foes to triumph over me.
\collect
\lett{G}{rant,} we beech thee, Almighty God: that we who believe that thou didst marvellously lift up thy Prophet Elias in a fiery chariot, while yet in this life; may at his intercession, while still alive, be raised to spiritual heights and rejoice in the resurrection of the just. Through.
\begin{rubric}
	Commemoration of St. Margaret of Antioch, from the Second Common of a Virgin Martyr (p. \pageref{CommonVirginMartyrII}).
\end{rubric}

\readingcitation{Epistle}{Ecclesiasticus 48:1}
%RV:
\lett{T}{here} arose Elijah the prophet as fire, And his word burned like a torch:  Who brought a famine upon them, And by his zeal made them few in number. By the word of the Lord he shut up the heaven: Thrice did he thus bring down fire. How wast thou glorified, O Elijah, in thy wondrous deeds! And who shall glory like unto thee? Who did raise up a dead man from death, And from the place of the dead, by the word of the Most High: Who brought down kings to destruction, And honourable men from their bed: Who heard rebuke in Sinai, And judgements of vengeance in Horeb: Who anointed kings for retribution, And prophets to succeed after him: Who was taken up in a tempest of fire, In a chariot of fiery horses: Who was recorded for reproofs in their seasons, To pacify anger, before it brake forth into wrath; To turn the heart of the father unto the son, And to restore the tribes of Jacob. 

%\lett{T}{hen} stood up Elias the prophet as fire, and his word burned like a lamp. He brought a sore famine upon them, and by his zeal he diminished their number. By the word of the Lord he shut up the heaven, and also three times brought down fire. O Elias, how wast thou honoured in thy wondrous deeds! and who may glory like unto thee! Who didst raise up a dead man from death, and his soul from the place of the dead, by the word of the most High: Who broughtest kings to destruction, and honorable men from their bed: Who heardest the rebuke of the Lord in Sinai, and in Horeb the judgment of vengeance: Who annointedst kings to take revenge, and prophets to succeed after him: Who was taken up in a whirlwind of fire, and in a chariot of fiery horses: Who wast ordained for reproofs in their times, to pacify the wrath of the Lord's judgment, before it brake forth into fury, and to turn the heart of the father unto the son, and to restore the tribes of Jacob.

\gradall{It came to pass, when the Lord would take up Elias into heaven by a whirlwind, that Elias went with Eliseus from Gilgal. ℣. And it came to pass, as they still went on, that, behold, there appeared a chariot of fire, and horses of fire, and parted them both asunder; and Elias went up by a whirlwind into heaven.}{Alleluia, alleluia. ℣. Elias, while he was full of zeal for the law, was taken up into heaven. Alleluia.}

\readingcitation{Gospel}{Luke 9:28}
\lett{A}{t that time:} It came to pass about an eight days after these sayings, Jesus took Peter and John and James, and went up into a mountain to pray. And as he prayed, the fashion of his countenance was altered, and his raiment was white and glistering. And, behold, there talked with him two men, which were Moses and Elias: Who appeared in glory, and spake of his decease which he should accomplish at Jerusalem. But Peter and they that were with him were heavy with sleep: and when they were awake, they saw his glory, and the two men that stood with him. And it came to pass, as they departed from him, Peter said unto Jesus, Master, it is good for us to be here: and let us make three tabernacles; one for thee, and one for Moses, and one for Elias: not knowing what he said. While he thus spake, there came a cloud, and overshadowed them: and they feared as they entered into the cloud. And there came a voice out of the cloud, saying, This is my beloved Son: hear him. And when the voice was past, Jesus was found alone. And they kept it close, and told no man in those days any of those things which they had seen.
\offertory{Elias was a man subject to like passions as we are: and he prayed earnestly that it might not rain, and it rained not on the earth for the space of three years and six months: and he prayed again, and the heaven gave rain, and the earth brought forth her fruit.}

\secret
\lett{W}{e} offer unto thee, O Lord, this sacrifice of praise in honour of thy Prophet Elias: that, as thou didst accept his burnt-offering, so thou wouldest vouchsafe to accept our own sacrifice; that through it we may be made worthy to attain unto everlasting gladness. Through.
\begin{rubric}
	Commemoration of St. Margaret of Antioch, from the Second Common of a Virgin Martyr (p. \pageref{CommonVirginMartyrII}).
\end{rubric}

\subbysub{Preface}
\lett{A}{nd} that we, with glad hearts, should praise, bless, and glorify thee in the Solemnity of blessed Elias thy Prophet: Who at thy word stood up like fire; shut the heavens and raised the dead; smote tyrants and slew the blasphemous; and laid the foundation of the monastic profession: Who, being fed with meat and drink served by Angels, walked in the strength of that food, even to the holy mountain: Who, being raised up from earth in a fiery chariot, will return to us as the Forerunner of the second Coming of Jesus Christ our Lord. Therefore.

\communion{Behold, I will send you Elias the Prophet before the coming of the great and terrible day of the Lord: and he shall turn the hearts of the fathers to the children, and the hearts of the children to their fathers.}

\postcommunion
\lett{O}{God,} who by thy holy Angel didst give meat and drink to blessed Elias thy Prophet: grant, by his intercession, that what we have received of thy heavenly table, we may keep undefiled in purity of mind. Through.
\begin{rubric}
	Commemoration of St. Margaret of Antioch, from the Second Common of a Virgin Martyr (p. \pageref{CommonVirginMartyrII}).
\end{rubric}


\subby{St. Margaret of Antioch}
\feastday{{St. Margaret Antioch}}
\fancyhead[RE,LO]{20 July}
\begin{inhead}
    {Memorial\\
20 July}
\end{inhead}

\begin{rubric}
	The propers are from the Second Common of a Virgin Martyr (p. \pageref{CommonVirginMartyrII}).
\end{rubric}


\subby{St. Praxedes}
\feastday{{St. Praxedes}}
\fancyhead[RE,LO]{21 July}
\begin{inhead}
    {Memorial\\
21 July}
\end{inhead}

\begin{rubric}
	The Prayers and Epistle are from the Second Common of a Virgin (p. \pageref{CommonVirginOnlyII}), with the following.
\end{rubric}

\introit
\lett{I}{will} speak of thy testimonies even before kings, and will not be ashamed: and my delight shall be in thy commandments, which I have loved exceedingly. \textit{Ps.} Blessed are those that are undefiled in the way: and walk in the law of the Lord.

\gradall{Thou hast loved righteousness, and hated iniquity. ℣. Wherefore God, even thy God, hath anointed thee with the oil of gladness.}{Alleluia, alleluia. ℣. In thy comeliness and in thy beauty, go forth, proceed prosperously, and reign. Alleluia.}

\readingcitation{Gospel}{Matthew 13:44}
\lett{A}{t that time:} Jesus spake this parable unto his disciples: The kingdom of heaven is like unto treasure hid in a field; the which when a man hath found, he hideth, and for joy thereof goeth and selleth all that he hath, and buyeth that field. Again, the kingdom of heaven is like unto a merchant man, seeking goodly pearls: Who, when he had found one pearl of great price, went and sold all that he had, and bought it. Again, the kingdom of heaven is like unto a net, that was cast into the sea, and gathered of every kind: Which, when it was full, they drew to shore, and sat down, and gathered the good into vessels, but cast the bad away. So shall it be at the end of the world: the angels shall come forth, and sever the wicked from among the just, And shall cast them into the furnace of fire: there shall be wailing and gnashing of teeth. Jesus saith unto them, Have ye understood all these things? They say unto him, Yea, Lord. Then said he unto them, Therefore every scribe which is instructed unto the kingdom of heaven is like unto a man that is an householder, which bringeth forth out of his treasure things new and old.

\offertory{Full of grace are thy lips, because God hath blessed thee for ever and ever.}

\communion{The kingdom of heaven is like unto a merchant man, seeking goodly pearls: who when he had found one pearl of great price, gave all that he had, and bought it.}


\bcpfeast{22 July. St. Mary Magdalene, Penitent}{St. Mary Magdalene}{22 July}
%\supplement{22 July}{St. Mary Magdalene}{}

\subbysub{I Evensong}\label{MaryMagdaleneEvensong}

\gregorioscore{resources/gabc/ProperTime/MaryMagdaleneEvensong.gabc}

    ℣. Mary hath chosen that good part.

	℟. Which shall not be taken away from her.

	\properantiphon{Mag.}{A woman {\dag} in the city which was a sinner, when she knew that Jesus sat at meat in the house of Simon the leper, brought an alabaster box of ointment, and stood behind at the feet of Jesus, and began to wash his feet with tears, and did wipe them with the hairs of her head, and kissed his feet, and anointed them with the ointment.}

\subbysub{Mattins}

\invitatoryhymn\label{MaryMagdaleneInvitatory}

\gregorioscore{resources/gabc/ProperTime/MaryMagdaleneInvitatory.gabc}

\officehymn\label{MaryMagdaleneMattins}

\gregorioscore{resources/gabc/ProperTime/MaryMagdaleneMattins.gabc}

    ℣. Her sins, which are many, are forgiven.

	℟. For she loved much.
	
	\properantiphon{Ben.}{Mary therefore {\dag} anointed the feet of Jesus, and wiped them with her hair, and the house was filled with the odour of the ointment.}

\subbysub{II Evensong}

\begin{rubric}
	The Office Hymn is of I Evensong with the following Versicle \& Antiphon.
\end{rubric}

    ℣. God hath chosen her and preferred her.

	℟. He hath made her to dwell in his tabernacle.
	
	\properantiphon{Mag.}{A woman {\dag} in the city which was a sinner, brought an alabaster box of ointment, and stood at the Lord's feet, and began to wash his feet with tears, and did wipe them with the hairs of her head.}


\subby{St. John Cassian}
\feastday{{St. John}}
\fancyhead[RE,LO]{23 July}
\begin{inhead}
    {Memorial\\
23 July}
\end{inhead}

\begin{rubric}
	The propers are from the Common of Abbots (p. \pageref{CommonAbbots}).
\end{rubric}


\subby{St. Apollinaris of Ravenna}
\feastday{{St. Apollinaris Ravenna}}
\fancyhead[RE,LO]{23 July}
\begin{inhead}
    {Memorial\\
23 July}
\end{inhead}

\introit
\lett{O}{ye} priests of the Lord, bless ye the Lord: O ye holy and humble men of heart, bless ye the Lord. \textit{Cant.} O all ye works of the Lord, bless ye the Lord: praise him and magnify him for ever. 

\collect
\lett{O}{God,} the rewarder of faithful souls, who hast hallowed this day by the martyrdom of blessed Apollinaris, thy Priest: grant, we beseech thee, unto us thy servants; that we who keep his solemn festival may by his prayers obtain thy pardon. Through.

\readingcitation{Epistle}{1 Peter 5:1}
\lett{D}{early beloved:} The elders which are among you I exhort, who am also an elder, and a witness of the sufferings of Christ, and also a partaker of the glory that shall be revealed: Feed the flock of God which is among you, taking the oversight thereof, not by constraint, but willingly; not for filthy lucre, but of a ready mind; Neither as being lords over God's heritage, but being ensamples to the flock. And when the chief Shepherd shall appear, ye shall receive a crown of glory that fadeth not away. Likewise, ye younger, submit yourselves unto the elder. Yea, all of you be subject one to another, and be clothed with humility: for God resisteth the proud, and giveth grace to the humble. Humble yourselves therefore under the mighty hand of God, that he may exalt you in due time: Casting all your care upon him; for he careth for you. Be sober, be vigilant; because your adversary the devil, as a roaring lion, walketh about, seeking whom he may devour: Whom resist stedfast in the faith, knowing that the same afflictions are accomplished in your brethren that are in the world. But the God of all grace, who hath called us unto his eternal glory by Christ Jesus, after that ye have suffered a while, make you perfect, stablish, strengthen, settle you. To him be glory and dominion for ever and ever. Amen.

\gradall{I have found David my servant, with my holy oil have I anointed him: my hand shall hold him fast, and my arm shall strengthen him. ℣. The enemy shall not be able to do him violence, the son of wickedness shall not hurt him.}{Alleluia, alleluia. ℣. The Lord sware, and will not repent: Thou art a priest for ever, after the order of Melchisedech. Alleluia.}

\readingcitation{Gospel}{Luke 22:24}
\lett{A}{t that time:} There was a strife among the disciples, which of them should be accounted the greatest. And he said unto them, The kings of the Gentiles exercise lordship over them; and they that exercise authority upon them are called benefactors. But ye shall not be so: but he that is greatest among you, let him be as the younger; and he that is chief, as he that doth serve. For whether is greater, he that sitteth at meat, or he that serveth? is not he that sitteth at meat? but I am among you as he that serveth. Ye are they which have continued with me in my temptations. And I appoint unto you a kingdom, as my Father hath appointed unto me; That ye may eat and drink at my table in my kingdom, and sit on thrones judging the twelve tribes of Israel.

\offertory{My truth and my mercy shall be with him: and in my name shall his horn be exalted.}

\secret
\lett{G}{raciously} look, O Lord, upon these gifts: which in remembrance of blessed Apollinaris, thy Priest and Martyr, we lay before thee, and offer up for our offences. Through.

\communion{Lord, thou deliveredst unto me five talents, behold, I have gained beside them five talents more. Well done, thou good and faithful servant, thou hast been faithful over a few things, I will make thee ruler over many things, enter thou into the joy of thy Lord.}

\postcommunion
\lett{W}{e} who have received thy holy things, beseech thee, O Lord, that the protection of blessed Apollinaris may continually defend us: forasmuch as thou failest not to look with mercy on those to whom thou dost grant the succour of his assistance. Through.


\subby{St. Liborius}
\feastday{{St. Liborius}}
\fancyhead[RE,LO]{23 July}
\begin{inhead}
    {Memorial\\
23 July}
\end{inhead}

\begin{rubric}
	The propers are from the First Common of a Confessor Bishop (p. \pageref{CommonConfessorBishopI}).
\end{rubric}


\subby{Vigil of St. James}
\feastday{{St. James Vigil}}
\fancyhead[RE,LO]{23 July}
\begin{inhead}
    {Vigil\\
24 July}
\end{inhead}

\begin{rubric}
	The propers are from the Common of Vigils of the Apostles (p. \pageref{CommonVigilApostles}), with commemoration of St. Christina from the Second Common of a Virgin Martyr (p. \pageref{CommonVirginMartyrII}), and the \nth{3} Collect of St. Mary in Eastertide (p. \SPMaryEaster).
\end{rubric}


\subby{St. Christina}
\feastday{{St. Christina}}
\fancyhead[RE,LO]{24 July}
\begin{inhead}
    {Memorial\\
24 July}
\end{inhead}


\bcpfeast{25 July. St. James}{St. James}{25 July}
%\supplement{25 July}{St. James}{}

\begin{secrubric}
	The Daily Office propers are of the Common of Apostles (p. \pageref{CommonApostles}).
\end{secrubric}


\subby{St. Christopher}
\feastday{{St. Christopher}}
\fancyhead[RE,LO]{25 July}
\begin{inhead}
    {Memorial\\
25 July}
\end{inhead}

\begin{rubric}
	The propers are from the First Common of a Martyr not a Bishop (p. \pageref{CommonMartyrNotBishopI}).
\end{rubric}


\bcpfeast{26 July. St. Anne}{St. Anne}{26 July}
%\supplement{26 July}{St. Anne}{}

\subbysub{I Evensong}\label{AnneEvensong}

\gregorioscore{resources/gabc/ProperTime/AnneEvensong.gabc}

    ℣. Be glad, O ye righteous and rejoice in the Lord.
    
	℟. And be joyful, all ye that are true of heart.

\properantiphon{Mag.}{The kingdom of heaven {\dag} is likened unto a merchantman seeking goodly pearls; who, when he had found one pearl of great price, went and sold all that he had, and bought it.}

\subbysub{Mattins}

\begin{rubric}
	The Invitatory Hymn is as in I Evensong.
\end{rubric}

\officehymn\label{AnneMattins}

\gregorioscore{resources/gabc/ProperTime/AnneMattins.gabc}

    ℣. Let the Saints be joyful with glory.
    
	℟. Let them rejoice in their beds.
	
\properantiphon{Ben.}{Give her {\dag} of the fruit of her hands; and let her own works praise her in the gates.}

\subbysub{II Evensong}\label{AnneEvensongII}

\gregorioscore{resources/gabc/ProperTime/AnneEvensongII.gabc}

    ℣. Let the Saints be joyful with glory.
    
	℟. Let them rejoice in their beds.

\properantiphon{Mag.}{She stretched out {\dag} her hand to the poor; yea, she reacheth forth her hands to the needy; she eateth not the bread of idleness.}


\subby{St. Pantaleon}
\feastday{{St. Pantaleon}}
\fancyhead[RE,LO]{27 July}
\begin{inhead}
    {Memorial\\
27 July}
\end{inhead}

\begin{rubric}
	The propers are from the Second Common of a Martyr not a Bishop (p. \pageref{CommonMartyrNotBishopII}).
\end{rubric}


\subby{Sts. Nazarius, Celsus, Victor, \& Innocent}
\feastday{{Sts. Nazarius \&c.}}
\fancyhead[RE,LO]{28 July}
\begin{inhead}
    {Memorial\\
28 July}
\end{inhead}

\begin{rubric}
	The propers are from the First Common of Many Martyrs (p. \pageref{CommonMartyrsI}), except for the following.
\end{rubric}

\collect
\lett{O}{Lord,} let the blessed confession of thy Saints, Nazarius, Celsus, Victor, and Innocent, strengthen us: and obtain for our frailty the succour of thy goodness. Through.

\begin{rubric}
	The Epistle is the first additional Epistle of the Third Common of Many Martyrs (p. \pageref{Wisdom1017}).
\end{rubric}

\secret
\lett{G}{rant} to us, almighty God: that we, who present these gifts unto thee to the honour of thy Saints, Nazarius, Celsus, Victor, and Innocent, may propitiate thee by the offering, and be quickened by the receiving of the same. Through.

\postcommunion
\lett{G}{rant,} O Lord, we beseech thee, that the intercession of thy Saints Nazarius, Celsus, Victor, and Innocent, may so make us acceptable unto thee: that those things which we perform in this temporal celebration we may receive unto eternal salvation. Through.


\subby{St. Martha of Bethany}
\feastday{{St. Martha}}
\fancyhead[RE,LO]{29 July}
\begin{inhead}
    {Double\\
29 July}
\end{inhead}

\begin{rubric}
	The propers are from the First Common of a Virgin (p. \pageref{CommonVirginOnlyI}), except for the following, with a Commemoration of St. Felix etc. as followeth.
\end{rubric}

\readingcitation{Gospel}{Luke 10:38}
\lett{A}{t that time:} Jesus entered into a certain village: and a certain woman named Martha received him into her house. And she had a sister called Mary, which also sat at Jesus' feet, and heard his word. But Martha was cumbered about much serving, and came to him, and said, Lord, dost thou not care that my sister hath left me to serve alone? bid her therefore that she help me. And Jesus answered and said unto her, Martha, Martha, thou art careful and troubled about many things: But one thing is needful: and Mary hath chosen that good part, which shall not be taken away from her.


\subby{Sts. Felix II, Simplicius, Faustinus, \& Beatrice}
\feastday{{St. Felix II \&c.}}
\fancyhead[RE,LO]{29 July}
\begin{inhead}
    {Memorial\\
29 July}
\end{inhead}

\begin{rubric}
	The propers are from the Second Common of Many Martyrs (p. \pageref{CommonMartyrsII}), except for the following.
\end{rubric}

\collect
\lett{G}{rant,} we beseech thee, O Lord: that as thy Christian people rejoice in the temporal solemnity of thy Martyrs, Felix, Simplicius, Faustinus, and Beatrice, so they may eternally enjoy the same; that those things which they devoutly celebrate, they may effectually obtain. Through.

\secret
\lett{W}{e} offer our sacrifices unto thee, O Lord, in commemoration of thy Holy Martyrs Felix, Simplicius, Faustinus, and Beatrice: humbly entreating thee; that they may bestow on us both pardon and salvation. Through.

\postcommunion
\lett{G}{rant,} we beseech thee, almighty God: that the solemnity of thy holy Martyrs Felix, Simplicius, Faustinus, and Beatrice: which we have celebrated in these heavenly mysteries, may obtain for us merciful pardon. Through.


\subby{Sts. Abdon \& Sennen}
\feastday{{Sts. Abdon \& Sennen}}
\fancyhead[RE,LO]{30 July}
\begin{inhead}
    {Memorial\\
30 July}
\end{inhead}

\introit
\lett{L}{et} the sorrowful sighing of the prisoners, O Lord, come before thee: reward thou our neighbours sevenfold in their bosom: avenge thou the blood of thy Saints that is shed. \textit{Ps.} O God, the heathen are come into thine inheritance: thy holy temple have they defiled: and made Jerusalem an heap of stones.

\collect
\lett{O}{God,} who on thy Saints Abdon and Sennen didst bestow the bounteous gift of grace to attain unto thy glory: grant unto thy servants the remission of their sins; that, by the intercession of the merits of thy Saints, they may be found worthy to be delivered from all adversities. Through.

\begin{rubric}
	The Epistle is the fourth additional Epistle of the Third Common of Many Martyrs (p. \pageref{2Corinthians64}).
\end{rubric}

\gradall{God is glorious in his holy ones, fearful in praises, doing wonders. ℣. Thy right hand, O Lord, is become glorious in power: thy right hand hath dashed in pieces the enemy.}{Alleluia, alleluia. ℣. The souls of the just are in the hand of God, and there shall no torment of malice touch them. Alleluia.}

\readingcitation{Gospel}{Matthew 5:1}
\lett{A}{t that time:} Jesus, seeing the multitudes, went up into a mountain, and when he was set, his disciples came unto him: And he opened his mouth, and taught them, saying, Blessed are the poor in spirit: for theirs is the kingdom of heaven. Blessed are they that mourn: for they shall be comforted. Blessed are the meek: for they shall inherit the earth. Blessed are they which do hunger and thirst after righteousness: for they shall be filled. Blessed are the merciful: for they shall obtain mercy. Blessed are the pure in heart: for they shall see God. Blessed are the peacemakers: for they shall be called the children of God. Blessed are they which are persecuted for righteousness' sake: for theirs is the kingdom of heaven. Blessed are ye, when men shall revile you, and persecute you, and shall say all manner of evil against you falsely, for my sake. Rejoice, and be exceeding glad: for great is your reward in heaven: 

\offertory{God is wonderful in his holy ones: even the God of Israel, he will give strength and power unto his people: blessed be God.}

\secret
\lett{W}{e} beseech thee, O Lord, that the sacrifice, which we offer in remembrance of the birthday of thy holy Martyrs, may both loose the bonds of our iniquity, and obtain for us the gifts of thy mercy. Through.

\communion{The dead bodies of thy servants, O Lord, have they given to be meat unto the fowls of the air, and the flesh of thy Saints unto the beasts of the land: according to the greatness of thy power preserve thou those that are appointed to die.}

\postcommunion
\lett{M}{ay} the operation of this mystery, O Lord, both cleanse our sins: and, at the intercession of thy holy Martyrs Abdon and Sennen, obtain the fulfilment of our rightful desires. Through.


\bcpfeast{1 August. Chains of St. Peter}{Lammas}{1 August}
%\supplement{1 August}{Chains of St. Peter}{}

\subbysub{I Evensong}\label{LammasEvensong}

\gregorioscore{resources/gabc/ProperTime/LammasEvensong.gabc}

    ℣. Thou art Peter.

	℟. And upon this rock I will build my Church.

\properantiphon{Mag.}{Thou art the shepherd of the sheep, {\dag} O chief of the Apostles: unto thee were given the keys of the kingdom of heaven.}

\subbysub{Mattins}

\invitatoryhymn\label{LammasInvitatory}

\gregorioscore{resources/gabc/ProperTime/LammasInvitatory.gabc}

\officehymn

\begin{rubric}
	The Office Hymn is from the Feast of the Chair of St. Peter at Antioch, 22 February (p. \pageref{PeterAntiochMattins}), with the following Versicle \& Antiphon.
\end{rubric}

    ℣. Thou art Peter.

	℟. And upon this rock I will build my Church.

\properantiphon{Ben.}{Whatsoever {\dag} thou shalt bind on earth shall be bound in heaven: and whatsoever thou shalt loose on earth shall be loosed in heaven: saith the Lord unto Simon Peter.}

\subbysub{II Evensong}

\begin{secrubric}
	The Office Hymn \& Versicle are as in I Evensong, with the following Antiphon.
\end{secrubric}

\properantiphon{Mag.}{Loosen these earthly fetters {\dag} at the command of God, O Peter; and let the kingdom of heaven be opened unto the blessed.}


\subby{Holy Maccabees}
\feastday{{Holy Maccabees}}
\fancyhead[RE,LO]{1 August}
\begin{inhead}
    {Memorial\\
1 August}
\end{inhead}

\introit
\lett{T}{he} just cry, and the Lord heareth them: and delivereth them out of all their troubles. \textit{Ps.} I will alway give thanks unto the Lord; his praise shall ever be in my mouth.

\collect
\lett{O}{Lord,} let the crown of the brethren, thy Martyrs, cause us to rejoice: that we may thereby be strengthened and increased in our faith; and comforted by their manifold intercession. Through.

\begin{rubric}
	The Epistle is the fifth additional Epistle of the Third Common of Many Martyrs (p. \pageref{Hebrews1133}).
\end{rubric}

\gradall{Behold, how good and joyful a thing it is, brethren, to dwell together in unity! ℣. It is like the precious ointment upon the head, that ran down unto the beard, even unto Aaron's beard.}{Alleluia, alleluia. ℣. This is the true brotherhood, which overcame the wickedness of the world: which followed Christ, gaining heaven's glorious realms. Alleluia.}

\begin{rubric}
	The Gospel is from the Third Common of Many Martyrs (p. \pageref{CommonMartyrsIII}).
\end{rubric}

\offertory{Let the Saints be joyful with glory: let them rejoice in their beds: let the praises of God be in their mouth, alleluia.}

\secret
\lett{G}{rant,} O Lord, that we may with devout hearts celebrate thy mysteries in honour of thy holy Martyrs: and thereby obtain an increase both of protection and joy. Through.

\communion{And I say unto you, my friends: Be not afraid of them that persecute you.}

\postcommunion
\lett{G}{rant,} we beseech thee, almighty God: that growing in virtue we may follow the faith of them whose memory we recall by the partaking of this sacrament. Through.


\subby{Pope St. Stephen}
\feastday{{Pope St. Stephen}}
\fancyhead[RE,LO]{2 August}
\begin{inhead}
    {Memorial\\
2 August}
\end{inhead}

\begin{rubric}
	The propers are from the Second Common for Confessor Bishops (p. \pageref{CommonConfessorBishopII}), except for that which followeth.
\end{rubric}

\introit
\lett{I}{will} deck her priests with health, and her saints shall rejoice and sing. \textit{Ps.} Lord, remember David: and all his trouble.

\readingcitation{Epistle}{Acts 20:17}
\lett{I}{n those days:} From Miletus Paul sent to Ephesus, and called the elders of the church. And when they were come to him, he said unto them, Ye know, from the first day that I came into Asia, after what manner I have been with you at all seasons, serving the Lord with all humility of mind, and with many tears, and temptations, which befell me by the lying in wait of the Jews: and how I kept back nothing that was profitable unto you, but have shewed you, and have taught you publickly, and from house to house, testifying both to the Jews, and also to the Greeks, repentance toward God, and faith toward our Lord Jesus Christ.

\gradall{Behold, a great priest, who in his days pleased God. ℣. There was none found like unto him, who kept the law of the Most High.}{Alleluia, alleluia. ℣. Thou art a priest for ever, after the order of Melchisedech. Alleluia.}

\communion{Lord, thou deliveredst unto me five talents: behold, I have gained beside them five talents more. Well done, thou good and faithful servant, thou hast been faithful over a few things. I will make thee ruler over many things, enter thou into the joy of thy Lord.}


\subby{The Invention of St. Stephen}
\feastday{{Invention of St. Stephen}}
\fancyhead[RE,LO]{3 August}
\begin{inhead}
    {Double\\
3 August}
\end{inhead}
\begin{rubric}
	The Office \& Mass as on 26 December, except for the following.
\end{rubric}

\begin{rubric}
	\textsc{Note,} The Creed is not said.
\end{rubric}

\collect
\lett{G}{rant} us, we beseech thee, O Lord, so to imitate that which we revere: that we may learn to love even our enemies; forasmuch as we celebrate the Finding of him, who was able to pray even for his persecutors to our Lord Jesus Christ thy Son. Who liveth.

\readingcitation{Epistle}{Acts 6:8}
\lett{I}{n those days:} Stephen, full of faith and power, did great wonders and miracles among the people. Then there arose certain of the synagogue, which is called the synagogue of the Libertines, and Cyrenians, and Alexandrians, and of them of Cilicia and of Asia, disputing with Stephen. And they were not able to resist the wisdom and the spirit by which he spake. When they heard these things, they were cut to the heart, and they gnashed on him with their teeth. But he, being full of the Holy Ghost, looked up stedfastly into heaven, and saw the glory of God, and Jesus standing on the right hand of God, and said, Behold, I see the heavens opened, and the Son of man standing on the right hand of God. Then they cried out with a loud voice, and stopped their ears, and ran upon him with one accord, and cast him out of the city, and stoned him: and the witnesses laid down their clothes at a young man’s feet, whose name was Saul. And they stoned Stephen, calling upon God, and saying, Lord Jesus, receive my spirit. And he kneeled down, and cried with a loud voice, Lord, lay not this sin to their charge. And when he had said this, he fell asleep in the Lord.


\subby{Dedication of Our Lady of the Snows}
\feastday{{Our Lady of the Snows}}
\fancyhead[RE,LO]{5 August}
\begin{inhead}
    {Greater Double\\
5 August}
\end{inhead}
\begin{rubric}
	The Office and Mass are from the Common of the Blessed Virgin Mary (p. \pageref{CommonBVM}), with Commemoration of St. Oswald from the First Common of a Martyr not a Bishop (p. \pageref{CommonMartyrNotBishopI}).
\end{rubric}

\begin{rubric}
	\textsc{Note,} The Creed is said and the Preface of the B.V.M. \emph{And that in the Festivity.}
\end{rubric}

%\collect
%\lett{A}{lmighty} and everlasting God, who by the martyrdom of the blessed King Oswald hast hallowed this day with holy joy and gladness: grant unto our hearts the increase of thy charity; that we, who honour his glorious battle for the faith, may imitate his constancy even unto death. Through.

%\readingcitation{Epistle}{Wisdom 4:7}
%RV:
%\lett{B}{ut} a righteous man, though he die before his time, shall be at rest. (For honourable old age is not that which standeth in length of time, Nor is its measure given by number of years: But understanding is gray hairs unto men, And an unspotted life is ripe old age.) Being found well-pleasing unto God he was beloved of him, And while living among sinners he was translated: He was caught away, lest †† wickedness should change his understanding, Or guile deceive his soul. (For the bewitching of naughtiness bedimmeth the things which are good, And the giddy whirl of desire perverteth an innocent mind.) Being made perfect in a little while, he fulfilled long years; For his soul was pleasing unto the Lord: Therefore hasted he out of the midst of wickedness. But as for the peoples, seeing and understanding not, Neither laying this to heart, That grace and mercy are with his chosen, And that he visiteth his holy ones.

%\lett{T}{hough} the righteous be prevented with death, yet shall he be in rest. For honourable age is not that which standeth in length of time, nor that is measured by number of years. But wisdom is the gray hair unto men, and an unspotted life is old age. He pleased God, and was beloved of him: so that living among sinners he was translated. Yea speedily was he taken away, lest that wickedness should alter his understanding, or deceit beguile his soul. For the bewitching of naughtiness doth obscure things that are honest; and the wandering of concupiscence doth undermine the simple mind. He, being made perfect in a short time, fulfilled a long time: For his soul pleased the Lord: therefore hasted he to take him away from among the wicked. This the people saw, and understood it not, neither laid they up this in their minds, That his grace and mercy is with his saints, and that he hath respect unto his chosen.

%\readingcitation{Gospel}{Matthew 16:24}
%\lett{A}{t that time:} Jesus said unto his disciples: If any man will come after me, let him deny himself, and take up his cross, and follow me. For whosoever will save his life shall lose it: and whosoever will lose his life for my sake shall find it. For what is a man profited, if he shall gain the whole world, and lose his own soul? or what shall a man give in exchange for his soul? For the Son of man shall come in the glory of his Father with his angels; and then he shall reward every man according to his works.


\subby{St. Oswald}
\feastday{{St. Oswald}}
\fancyhead[RE,LO]{3 August}
\begin{inhead}
    {Double\\
3 August}
\end{inhead}

\begin{rubric}
	The propers are from the First Common of a Martyr not a Bishop (p. \pageref{CommonMartyrNotBishopI}), except for the following.
\end{rubric}

\collect
\lett{A}{lmighty} and everlasting God, who by the martyrdom of the blessed King Oswald hast hallowed this day with holy joy and gladness: grant unto our hearts the increase of thy charity; that we, who honour his glorious battle for the faith, may imitate his constancy even unto death. Through.

%RV:
\readingcitation{Epistle}{Wisdom 4:7}
\lett{B}{ut} a righteous man, though he die before his time, shall be at rest. (For honourable old age is not that which standeth in length of time, Nor is its measure given by number of years: But understanding is gray hairs unto men, And an unspotted life is ripe old age.) Being found well-pleasing unto God he was beloved of him, And while living among sinners he was translated: He was caught away, lest wickedness should change his understanding, Or guile deceive his soul. (For the bewitching of naughtiness bedimmeth the things which are good, And the giddy whirl of desire perverteth an innocent mind.) Being made perfect in a little while, he fulfilled long years; For his soul was pleasing unto the Lord: Therefore hasted he out of the midst of wickedness. But as for the peoples, seeing and understanding not, Neither laying this to heart, That grace and mercy are with his chosen, And that he visiteth his holy ones.

\readingcitation{Gospel}{Matthew 16:24}
\lett{A}{t that time:} Jesus said unto his disciples: If any man will come after me, let him deny himself, and take up his cross, and follow me. For whosoever will save his life shall lose it: and whosoever will lose his life for my sake shall find it. For what is a man profited, if he shall gain the whole world, and lose his own soul? or what shall a man give in exchange for his soul? For the Son of man shall come in the glory of his Father with his angels; and then he shall reward every man according to his works.


\bcpfeast{6 August. Transfiguration of Our Lord Jesus Christ}{Transfiguration}{6 August}
%\supplement{6 August}{Transfiguration}{of Our Lord Jesus Christ}

\subbysub{I Evensong}\label{TransfigurationEvensong}

\gregorioscore{resources/gabc/ProperTime/TransfigurationEvensong.gabc}

    ℣. Glorious didst thou appear in the sight of the Lord. 
    
	℟. Because the Lord hath clothed thee with majesty.


\properantiphon{Mag.}{Christ Jesus, {\dag} the brightness of the Father and the express image of his person, who upholdeth all things by the word of his power, while he was by himself purging away our sins, vouchsafed on this day to shew himself in glory upon an high mountain.}

\subbysub{Mattins}

\begin{rubric}
	The Invitatory Hymn is as in I Evensong.
\end{rubric}

\officehymn\label{TransfigurationMattins}

\gregorioscore{resources/gabc/ProperTime/TransfigurationMattins.gabc}

    ℣. A crown of gold is upon his head.
    
	℟. A visible sign of holiness, glory, and honour.

\properantiphon{Ben.}{And behold a voice {\dag} out of the cloud, which said, This is my beloved Son, in whom I am well pleased; hear ye him, alleluia.}

\subbysub{II Evensong}

\begin{rubric}
	The Office Hymn \& Versicle are of I Evensong, with the following Antiphon.
\end{rubric}

\properantiphon{Mag.}{And when the disciples heard it, {\dag} they fell on their face, and were sore afraid: and Jesus came and touched them, and said, Arise, and be not afraid, alleluia.}


\subby{Sts. Sixtus II, Felicissimus, \& Agapitus}
\feastday{{Sts. Sixtus II \&c.}}
\fancyhead[RE,LO]{6 August}
\begin{inhead}
    {Memorial\\
6 August}
\end{inhead}

\begin{rubric}
	The propers are from the Second Common of Many Martyrs (p. \pageref{CommonMartyrsII}).
\end{rubric}

%MANUAL ADJUSTMENT:
\clearpage
\subby{St. Donatus}
\feastday{{St. Donatus}}
\fancyhead[RE,LO]{7 August}
\begin{inhead}
    {Memorial\\
7 August}
\end{inhead}

\introit
\lett{O}{ye} priest of the Lord, bless ye the Lord: O ye holy and humble men of heart, bless ye the Lord. \textit{Cant.} O all ye works of the Lord, bless ye the Lord: praise him and magnify him for ever.

\collect
\lett{O}{God,} the glory of thy priests: grant, we beseech thee; that we may perceive the succour of thy holy Martyr and Bishop Donatus, whose festival we celebrate. Through.

\begin{rubric}
	The Epistle is the first additional Epistle from the Second Common of a Martyr not a Bishop (p. \pageref{James12}).
\end{rubric}

\gradall{The mouth of the just is exercised in wisdom, and his tongue will be talking of judgment. ℣. The law of his God is in his heart: and his goings shall not slide.}{Alleluia, alleluia. ℣. The just shall not be moved, for the Lord strengtheneth his hand. Alleluia.}

\begin{rubric}
	The Gospel is the second additional Gospel from the Second Common of a Confessor Bishop (p. \pageref{Mark1333}).
\end{rubric}

\offertory{I have found David my servant, with my holy oil have I anointed him: my hand shall hold him fast, and my arm shall strengthen him.}

\secret
\lett{G}{rant,} we beseech thee, O Lord: that, as by these gifts which we offer to the praise of thy name, we render honour to thy holy Martyr and Bishop Donatus, so by his intercession we may receive the reward of this our bounden service. Through.

\communion{A faithful and wise servant, whom the lord hath made ruler over his household: to give them their portion of meat in due season.}

\postcommunion
\lett{A}{lmighty} and merciful God, who makest us alike partakers and ministers of thy sacraments: grant, we beseech thee; that at the intercession of blessed Donatus, thy Martyr and Bishop, we, sharing in his faith and worthily serving thee, may be profited by the same. Through.


\subby{Sts. Cyriacus, Largus, \& Smaragdus}
\feastday{{Sts. Cyriacus \&c.}}
\fancyhead[RE,LO]{8 August}
\begin{inhead}
    {Memorial\\
8 August}
\end{inhead}

\introit
\lett{O}{fear} the Lord, ye that are his saints, for they that fear him lack nothing: the lions do lack, and suffer hunger: but they who seek the Lord shall want no manner of thing that is good. \textit{Ps.} I will alway give thanks unto the Lord: his praise shall ever be in my mouth.

\lett{O}{God,} who makest us glad with the yearly solemnity of thy holy Martyrs Cyriacus, Largus, and Smaragdus: mercifully grant; that as we now celebrate their birthday, so we may imitate their constancy in suffering. Through.

\begin{rubric}
	If today be Saturday, Commemoration is made of the anticipated Vigil of St. Lawrence, as on the following day; the \nth{3} Collect of St. Mary, and the Last Gospel of the Vigil.
\end{rubric}

\readingcitation{Epistle}{1 Thessalonians 2:13}
\lett{B}{rethren:} We thank God without ceasing, because, when ye received the word of God which ye heard of us, ye received it not as the word of men, but as it is in truth, the word of God, which effectually worketh also in you that believe. For ye, brethren, became followers of the churches of God which in Jud{\ae}a are in Christ Jesus: for ye also have suffered like things of your own countrymen, even as they have of the Jews: who both killed the Lord Jesus, and their own prophets, and have persecuted us; and they please not God, and are contrary to all men: forbidding us to speak to the Gentiles that they might be saved, to fill up their sins alway: for the wrath is come upon them to the uttermost.

\gradall{O fear the Lord, ye that are his saints: for they that fear him lack nothing. ℣. But they who seek the Lord, shall want no manner of thing that is good.}{Alleluia, alleluia. ℣. The righteous shall shine, and run to and fro like sparks among the stubble for ever. Alleluia.}

\readingcitation{Gospel}{Mark 16:15}
\lett{A}{t that time:} Jesus said to his disciples: Go ye into all the world, and preach the gospel to every creature. He that believeth and is baptized shall be saved; but he that believeth not shall be damned. And these signs shall follow them that believe; In my name shall they cast out devils; they shall speak with new tongues; they shall take up serpents; and if they drink any deadly thing, it shall not hurt them; they shall lay hands on the sick, and they shall recover.

\offertory{Be glad, O ye righteous, and rejoice in the Lord: and be joyful, all ye that are true of heart.}

\secret
\lett{G}{rant,} O Lord, that this our bounden service may be acceptable in thy sight: that these our oblations may, by the prayers of those on whose solemnity they are offered, be made profitable unto our salvation. Through.

\communion{These signs shall follow them that believe in me: they shall cast out devils: they shall lay hands on the sick, and they shall recover.}

\postcommunion
\lett{W}{e} beseech thee, O Lord our God, that like as we, whom thou hast refreshed by the partaking of thy sacred gift, do offer unto thee our worship: so by the intercession of thy holy Martyrs, Cyriacus, Largus and Smaragdus, we may perceive the benefit of the same. Through.


\subby{Vigil of St. Lawrence}
\feastday{{St. Lawrence Vigil}}
\fancyhead[RE,LO]{9 August}
\begin{inhead}
    {Vigil\\
9 August}
\end{inhead}

\introit
\lett{H}{e} hath dispersed abroad, and given to the poor: his righteousness remaineth for ever: his horn shall be exalted with honour. \textit{Ps.} Blessed is the man that feareth the Lord: he hath great delight in his commandments.

\begin{rubric}
	\emph{Gloria in excelsis} is not said.
\end{rubric}

\collect
\lett{A}{ssist} us, O Lord, in these our supplications: and at the intercession of thy blessed Martyr Lawrence, whose festival we prevent, graciously bestow upon us thy perpetual mercy. Through.

\begin{rubric}
	Commemoration of St. Romanus (p. \pageref{RomanusCollect}).
\end{rubric}

%RV:
\readingcitation{Epistle}{Ecclesiasticus 51:1}
\lett{I}{will} give thanks unto thee, O Lord, O King, And will praise thee, O God my Saviour: I do give thanks unto thy name: For thou wast my protector and helper, And didst deliver my body out of destruction, And out of the snare of a slanderous tongue, From lips that forge lies, And wast my helper before them that stood by; And didst deliver me, according to the abundance of thy mercy, and greatness of thy name, From the gnashings of teeth ready to devour, Out of the hand of such as sought my life, Out of the manifold afflictions which I had; From the choking of a fire on every side, And out of the midst of fire which I kindled not; Out of the depth of the belly of the grave, And from an unclean tongue, And from lying words, The slander of an unrighteous tongue unto the king. My soul drew near even unto death, And my life was near to the grave beneath. They compassed me on every side, And there was none to help me. I was looking for the succour of men, And it was not. And I remembered thy mercy, O Lord, And thy working which hath been from everlasting, How thou deliverest them that wait for thee, And savest them out of the hand of the enemies, O Lord our God.

\gradual{He hath dispersed abroad, and given to the poor: and his righteousness remaineth for ever. ℣. His seed shall be mighty upon earth: the generation of the faithful shall be blessed.}

\readingcitation{Gospel}{Matthew 16:24}
\lett{A}{t that time:} Jesus said unto his disciples: If any man will come after me, let him deny himself, and take up his cross, and follow me. For whosoever will save his life shall lose it: and whosoever will lose his life for my sake shall find it. For what is a man profited, if he shall gain the whole world, and lose his own soul? or what shall a man give in exchange for his soul? For the Son of man shall come in the glory of his Father with his angels; and then he shall reward every man according to his works.

\offertory{My prayer is pure: and therefore I ask that a place be given to my voice in heaven: for my witness is in heaven, and my record is on high: let my prayer ascend to the Lord.}

\secret
\lett{O}{Lord,} mercifully regard the sacrifices which we offer unto thee: and at the intercession of blessed Lawrence, thy Martyr, absolve us from the bonds of our sins. Through.

\begin{rubric}
	Commemoration of St. Romanus (p. \pageref{RomanusSecret}).
\end{rubric}

\communion{He that will come after me, let him deny himself, and take up his cross, and follow me.}

\postcommunion
\lett{G}{rant,} we beseech thee, O Lord, our God: that like as we in this life do gladly honour the memory of blessed Lawrence, thy Martyr; so we may rejoice to behold him for ever. Through.

\begin{rubric}
	Commemoration of St. Romanus (p. \pageref{RomanusPostcommunion}).
\end{rubric}


\subby{St. Romanus}
\feastday{{St. Romanus}}
\fancyhead[RE,LO]{9 August}
\begin{inhead}
    {Memorial\\
9 August}
\end{inhead}

\begin{rubric}
	The propers are from Second Common of a Martyr not a Bishop (p. \pageref{CommonMartyrNotBishopII}), except for the following.
\end{rubric}

\collect\label{RomanusCollect}
\lett{G}{rant,} we beseech thee, almighty God: that, at the intercession of blessed Romanus, thy Martyr, we may both be delivered from all adversities which may happen to the body, and from all evil thoughts which may assault and hurt the soul. Through.

\secret\label{RomanusSecret}
\lett{W}{e} beseech thee, O Lord, to accept our prayers and oblations: and graciously hearken unto us, whom thou dost cleanse by thy heavenly mysteries. Through.

\postcommunion\label{RomanusPostcommunion}
\lett{W}{e} beseech thee, almighty God: that we, who have received this heavenly food, may, at the intercession of blessed Romanus, thy Martyr, be thereby defended against all adversities. Through.


\bcpfeast{10 August. St. Lawrence}{St. Lawrence}{10 August}
%\supplement{10 August}{St. Lawrence}{}

\begin{rubric}
	The Daily Office propers are of the First Common of a Martyr Bishop (p. \pageref{CommonMartyrBishopI}), except for the following.
\end{rubric}

\properantiphon{Mag.}{Laurence the Deacon {\dag} hath wrought a pious work; who by the sign of the Cross enlightened the blind: and distributed to the poor the Church's treasures.}\\

    ℣. He hath dispersed abroad and given to the poor.
    
	℟. And his righteousness remaineth for ever.

\properantiphon{Ben.}{On the iron grate, {\dag} O God, I denied thee not; and when fire was kindled beneath me, O Christ, I confessed thee: thou hast proved my heart and hast visited me in the night season; thou hast tried me with fire, and hast found no wickedness in me.}\\

    ℣. Laurence the Deacon hath wrought a pious work.
    
	℟. Who by the sign of the Cross enlightened the blind.

\properantiphon{Mag.}{Blessed Laurence, {\dag} when laid and burning on the iron grating, spake to the impious tyrant, saying: The feast is ready, turn and eat; but the Church's treasures, which thou claimest, have been garnered up in heaven, by the hands of the poor and needy.}


\subby{Sts. Tiburtius \& Susanna}
\feastday{{Sts. Tiburtius \& Susanna}}
\fancyhead[RE,LO]{11 August}
\begin{inhead}
    {Memorial\\
11 August}
\end{inhead}

\begin{rubric}
	The propers are from the Third Common of Many Martyrs (p. \pageref{CommonMartyrsIII}), except for the following.
\end{rubric}

\collect
\lett{O}{Lord,} let the protection of thy holy Martyrs Tiburtius and Susanna continually defend us: forasmuch as thou failest not to look with mercy on those to whom thou dost grant the succour of their assistance. Through.

\begin{rubric}
	The Epistle is from the third additional epistle of the Third Common of Many Martyrs (p. \pageref{Romans818}).
\end{rubric}

\secret
\lett{A}{ssist,} O Lord, the prayers of thy people, assist their oblations: that those things which are offered in these sacred mysteries may, by the intercession of thy Saints, be acceptable unto thee. Through.

\postcommunion
\lett{O}{Lord,} through whom we have received the pledge of eternal redemption: we beseech thee that at the intercession of thy holy Martyrs, it may avail for our succour both in this life and that which is to come. Through.


\subby{St. Maximus of Constantinople}
\feastday{{St. Maximus Constantinople}}
\fancyhead[RE,LO]{13 August}
\begin{inhead}
    {Double\\
13 August}
\end{inhead}

\begin{rubric}
	The propers are from the Second Common of a Confessor not a Bishop (p. \pageref{CommonConfessorNotBishopII}), with Commemoration of Sts. Hippolytus \& Cassian.
\end{rubric}

%MANUAL ADJUSTMENT:
\clearpage
\subby{Sts. Hippolytus \& Cassian}
\feastday{{Sts. Hippolytus \& Cassian}}
\fancyhead[RE,LO]{13 August}
\begin{inhead}
    {Memorial\\
13 August}
\end{inhead}

\begin{rubric}
	The propers are from the Third Common of Many Martyrs (p. \pageref{CommonMartyrsIII}), except for the following.
\end{rubric}

\collect
\lett{G}{rant,} we beseech thee, almighty God: that the venerable solemnity of thy blessed Martyrs, Hippolytus and Cassian, may increase our devotion and set forward our salvation. Through.

\secret
\lett{R}{egard,} O Lord, the gifts of thy people, which we offer on the festival of thy Saints: and let this confession of thy truth be profitable for our salvation. Through.

\postcommunion
\lett{M}{ay} the communion of thy sacraments, O Lord, which we have received, avail for our salvation: stablish us in the light of thy truth. Through.

\begin{rubric}
	If today be Saturday, the anticipated Vigil of the Assumption of the B.V. Mary is kept, as is noted on the following day, but with Commemoration of Sts. Hippolytus and Cassian, instead of St. Eusebius.
\end{rubric}


%The following is from Sarum, to match the BCP. It's mostly similar to the English Missal's text.
\subby{Vigil of the Assumption of the Blessed Virgin Mary}
\feastday{{Assumption Vigil}}
\fancyhead[RE,LO]{14 August}
\begin{inhead}
    {Vigil\\
14 August}
\end{inhead}

\introit
\lett{H}{ail,} Holy Mother, who didst bring forth the King who ruleth over heaven and earth for ever and ever. \textit{Gospel.} Blessed art thou among women, and blessed is the fruit of thy womb.

\collect
\lett{O}{God,} who didst vouchsafe to choose the virgin womb of blessed Mary wherein to make thy dwelling: grant, we beseech thee; that, being defended by her protection, we may by thee be enabled to attain with gladness to her festival. Who livest.

\begin{rubric}
	Commemoration of St. Eusebius (p. \pageref{EusebiusCollect}).
\end{rubric}

\begin{rubric}
	\nth{3} Collect of the Holy Ghost (p. \SPHolyGhost).
\end{rubric}

\readingcitation{Epistle}{Ecclesiasticus 24:9}
%RV:
\lett{H}{e} created me from the beginning before the world; And to the end I shall not fail. In the holy tabernacle I ministered before him; And so was I established in Sion. In the beloved city likewise he gave me rest; And in Jerusalem was my authority. And I took root in a people that was glorified, Even in the portion of the Lord's own inheritance. 

%\lett{H}{e} created me from the beginning before the world, and I shall never fail. In the holy tabernacle I served before him; and so was I established in Sion. Likewise in the beloved city he gave me rest, and in Jerusalem was my power. And I took root in an honourable people, even in the portion of the Lord's inheritance, and my abode is in the full congregation of the saints.

\gradual{Blessed and venerable art thou, O Virgin Mary: who without spot wast found the Mother of the Saviour. ℣. Virgin, Mother of God, he whom the whole world containeth not, being made man lay hid in thy womb.}

\readingcitation{Gospel}{Luke 11:27}
\lett{A}{t that time:} As Jesus spake to the multitudes, a certain woman of the company lifted up her voice, and said unto him, Blessed is the womb that bare thee, and the paps which thou hast sucked. But he said, Yea rather, blessed are they that hear the word of God, and keep it.

\offertory{Happy art thou, O holy Virgin Mary, and most worthy of all praise, for from thee sprang the Sun of Righteousness, Christ our God. ℣. Blessed art thou, Virgin Mary, who didst bear the Lord, didst give birth to the Creator of the world, Who made thee, and ever remainest Virgin.}

\secret
\lett{O}{Lord,} who didst translate the Mother of God from this present life, to the intent that she might faithfully intercede before thee for our sins: grant that her prayers may render these our oblations acceptable in the sight of thy mercy. Through the same.

\begin{rubric}
	Commemoration of St. Eusebius (p. \pageref{EusebiusSecret}).
\end{rubric}

\begin{rubric}
	\nth{3} Secret of the Holy Ghost (p. \SPHolyGhost).
\end{rubric}

\communion{Gentle Mother of God, succour all that pray; we, too, with them, humbly entreat that by the aid of thy prayers we may sing praises unto the Trinity.}

%While the English Missal's text and the Sarum's are very similar, the Sarum is used because of its distinctive use of `requiem'.
\postcommunion
\lett{W}{e} beseech thee, O merciful God, to strengthen our frailty, that we who keep the requiem of the Holy Virgin Mother of God, may by her intercession rise again from our iniquities. Through.

\begin{rubric}
	Commemoration of St. Eusebius (p. \pageref{EusebiusPostcommunion}).
\end{rubric}

\begin{rubric}
	\nth{3} Postcommunion of the Holy Ghost (p. \SPHolyGhost).
\end{rubric}


\subby{St. Eusebius}
\feastday{{St. Eusebius}}
\fancyhead[RE,LO]{14 August}
\begin{inhead}
    {Memorial\\
14 August}
\end{inhead}

\introit
\lett{T}{he} just shall flourish like a palm-tree: and shall spread abroad like a cedar in Libanus: planted in the house of the Lord: in the courts of the house of our God. \textit{Ps.} It is a good thing to give thanks unto the Lord: and to sing praises unto thy name, O Most Highest.

\collect\label{EusebiusCollect}
\lett{O}{God,} who makest us glad with the yearly solemnity of blessed Eusebius thy Confessor: mercifully grant; that we, who celebrate his birthday, may by his example, be drawn nearer unto thee. Through.

\begin{rubric}
	The Epistle is additional Epistle from the Second Common of a Confessor not a Bishop (p. \pageref{Philippians37}).
\end{rubric}

\gradall{The just shall flourish like a palm-tree: and shall spread abroad like a cedar in Libanus in the house of the Lord. ℣. To tell of thy loving-kindness early in the morning, and of thy truth in the night-season.}{Alleluia, alleluia. ℣. The just shall grow as the lily: and flourish for ever before the Lord. Alleluia.}

\readingcitation{Gospel}{Matthew 11:25}
\lett{A}{t that time:} Jesus answered and said: I thank thee, O Father, Lord of heaven and earth, because thou hast hid these things from the wise and prudent, and hast revealed them unto babes. Even so, Father: for so it seemed good in thy sight. All things are delivered unto me of my Father: and no man knoweth the Son, but the Father; neither knoweth any man the Father, save the Son, and he to whomsoever the Son will reveal him. Come unto me, all ye that labour and are heavy laden, and I will give you rest. Take my yoke upon you, and learn of me; for I am meek and lowly in heart: and ye shall find rest unto your souls. For my yoke is easy, and my burden is light.

\offertory{The just shall rejoice in strength, O Lord: exceeding glad shall he be of thy salvation: thou hast given him his heart’s desire.}

\secret\label{EusebiusSecret}
\lett{G}{rant,} we beseech thee, O Lord, that we who, trusting in this our sacrifice of praise, do offer it before thee to the honour of thy Saints: may by the same be delivered from all evils both in this life and in that which is to come. Through.

\communion{The just shall rejoice in the Lord, and put his trust in him: and all they that are true of heart shall be glad.}

\postcommunion\label{EusebiusPostcommunion}
\lett{O}{Lord,} our God, who hast refreshed us with heavenly meat and drink, we humbly beseech thee: that we may be defended by the prayers of him in whose memory we have received the same. Through.


\bcpfeast{15 August. Assumption of the Blessed Virgin Mary}{Assumption}{15 August}
%\supplement{15 August}{Assumption}{of the Blessed Virgin Mary}

\begin{rubric}
	The Daily Office propers are of the Common of the Blessed Virgin Mary (p. \pageref{CommonBVM}), except for the following.
\end{rubric}

    ℣. The holy Mother of God is exalted.
    
	℟. Above choirs of Angels to the heavenly kingdom.

\properantiphon{Mag.}{O most prudent Virgin, {\dag} whither goest thou, shining resplendent like the glowing dawn? Daughter of Sion, thou art all comely and beautiful, fair as the moon, clear as the sun.}\\

    ℣. The holy Mother of God is exalted.
    
	℟. Above choirs of Angels to the heavenly kingdom.

\properantiphon{Ben.}{Who is she {\dag} that riseth up as the morning, fair as the moon, clear as the sun, and terrible as an army with banners?}

\begin{rubric}
	II Evensong is as in I Evensong, except the following Antiphon.
\end{rubric}

\properantiphon{Mag.}{On this day {\dag} the Virgin Mary went up into heaven: rejoice, because she reigneth with Christ for ever and ever.}


\bcpfeast{16 August. St. Joachim}{St. Joachim}{16 August}
%\supplement{16 August}{St. Joachim}{}

\begin{rubric}
	The Daily Office propers are of the First Common of a Confessor not a Bishop (p. \pageref{CommonConfessorNotBishopI}), except for the following Versicle \& Antiphon used for I Evensong and Mattins.
\end{rubric}

    ℣. His seed shall be mighty upon earth.
    
	℟. The generation of the faithful shall be blessed.

\properantiphon{Mag. \& Ben.}{Let us praise a man famous {\dag} in his generation, with whom the Lord established the blessing of all nations, and the covenant, and made it rest upon his head.}

%MANUAL ADJUSTMENT:
\clearpage
\subby{Octave Day of St. Lawrence}
\feastday{{St. Lawrence Octave Day}}
\fancyhead[RE,LO]{17 August}
\begin{inhead}
    {Simple\\
17 August}
\end{inhead}

\introit
\lett{T}{hou} hast proved and visited mine heart, O Lord, in the night-season: thou hast tried me with fire, and hast found no wickedness in me. \textit{Ps.} Hear the right, O Lord: consider my complaint.

\collect
\lett{S}{tir} up, O Lord, in thy Church that Spirit whom the blessed Levite Lawrence served: that we, being filled with the same, may study to love that which he loved, and to perform in deed that which he taught. Through . . . in the unity of the same Holy Spirit.

\readingcitation{Epistle}{2 Corinthians 9:6}
\lett{B}{rethren:} He which soweth sparingly shall reap also sparingly; and he which soweth bountifully shall reap also bountifully. Every man according as he purposeth in his heart, so let him give; not grudgingly, or of necessity: for God loveth a cheerful giver. And God is able to make all grace abound toward you; that ye, always having all sufficiency in all things, may abound to every good work: (As it is written, He hath dispersed abroad; he hath given to the poor: his righteousness remaineth for ever. Now he that ministereth seed to the sower both minister bread for your food, and multiply your seed sown, and increase the fruits of your righteousness;)

\gradall{Thou hast crowned him with glory and worship, O Lord. ℣. And hast made him to have dominion of the works of thy hands.}{Alleluia, alleluia. ℣. The Levite Lawrence wrought a good work: who by the sign of the cross gave light to the blind. Alleluia.}

\readingcitation{Gospel}{John 12:24}
\lett{A}{t that time:} Jesus said unto his disciples: Verily, verily, I say unto you, Except a corn of wheat fall into the ground and die, it abideth alone: but if it die, it bringeth forth much fruit. He that loveth his life shall lose it; and he that hateth his life in this world shall keep it unto life eternal. If any man serve me, let him follow me; and where I am, there shall also my servant be: if any man serve me, him will my Father honour.

\offertory{The just shall rejoice in thy strength, O Lord: exceeding glad shall he be of thy salvation: thou hast given him his heart's desire.}

\secret
\lett{O}{Lord,} we beseech thee, let the holy prayer of blessed Lawrence commend the sacrifices unto thee: that it may be rendered acceptable by the merits of him in whose honour it is solemnly offered forth. Through.

\communion{He that will come after me, let him deny himself, and take up his cross, and follow me.}

\postcommunion
\lett{W}{e} humbly beseech thee, almighty God: that we, whom thou hast fulfilled with heavenly gifts, may, at the intercession of blessed Lawrence thy Martyr, be defended by thy continual protection. Through.

%Found in the supplement of the 1920 Missale Romanum.
\subby{St. Helen}
\feastday{{St. Helen}}
\fancyhead[RE,LO]{18 August}
\begin{inhead}
    {Memorial\\
18 August}
\end{inhead}

\begin{rubric}
	\textsc{Note,} When celebrated within the Octave of the Assumption, this Memorial is merely commemorated at Mass.
\end{rubric}

\introit
\lett{B}{ut} God forbid that I should glory, save in the Cross of our Lord Jesus Christ: by whom the world is crucified unto me, and I unto the world. \textit{Ps.} Thy rod and thy staff comfort me.

\collect
\lett{O}{Lord} Jesu Christ, who didst reveal unto blessed Helen the place where thy Cross lay hid, that through her thou mightest enrich thy Church with this precious treasure: grant unto us at her intercession; that by the ransom of the life-giving tree we may attain unto the rewards of everlasting life. Who livest.

\readingcitation{Epistle}{Proverbs 31:10}
\lett{W}{ho} can find a virtuous woman? for her price is far above rubies. The heart of her husband doth safely trust in her, so that he shall have no need of spoil. She will do him good and not evil all the days of her life. She seeketh wool, and flax, and worketh willingly with her hands. She is like the merchants' ships; she bringeth her food from afar. She riseth also while it is yet night, and giveth meat to her household, and a portion to her maidens. She considereth a field, and buyeth it: with the fruit of her hands she planteth a vineyard. She girdeth her loins with strength, and strengtheneth her arms. She perceiveth that her merchandise is good: her candle goeth not out by night. She layeth her hands to the spindle, and her hands hold the distaff. She stretcheth out her hand to the poor; yea, she reacheth forth her hands to the needy. She is not afraid of the snow for her household: for all her household are clothed with scarlet. She maketh herself coverings of tapestry; her clothing is silk and purple. Her husband is known in the gates, when he sitteth among the elders of the land. She maketh fine linen, and selleth it; and delivereth girdles unto the merchant. Strength and honour are her clothing; and she shall rejoice in time to come. She openeth her mouth with wisdom; and in her tongue is the law of kindness. She looketh well to the ways of her household, and eateth not the bread of idleness. Her children arise up, and call her blessed; her husband also, and he praiseth her. Many daughters have done virtuously, but thou excellest them all. Favour is deceitful, and beauty is vain: but a woman that feareth the \divineName{Lord}, she shall be praised. Give her of the fruit of her hands; and let her own works praise her in the gates.

\gradall{Like as the rich also among the people shall make their supplication before thee. Kings' daughters were among thy honourable women. ℣. She shall be brought unto the King in raiment of needle-work: the virgins that be her fellows shall bear her company, and shall be brought unto thee. With joy and gladness shall they be brought: and shall enter into the King's palace.}{Alleluia, alleluia. ℣. He hath dispersed abroad, and given to the poor : and his righteousness remaineth for ever. Alleluia.}

\readingcitation{Gospel}{Matthew 13:44}
\lett{A}{t that time:} Jesus spake this parable to his disciples: The kingdom of heaven is like unto treasure hid in a field; the which when a man hath found, he hideth, and for joy thereof goeth and selleth all that he hath, and buyeth that field. Again, the kingdom of heaven is like unto a merchant man, seeking goodly pearls: Who, when he had found one pearl of great price, went and sold all that he had, and bought it. Again, the kingdom of heaven is like unto a net, that was cast into the sea, and gathered of every kind: Which, when it was full, they drew to shore, and sat down, and gathered the good into vessels, but cast the bad away. So shall it be at the end of the world: the angels shall come forth, and sever the wicked from among the just, And shall cast them into the furnace of fire: there shall be wailing and gnashing of teeth. Jesus saith unto them, Have ye understood all these things? They say unto him, Yea, Lord. Then said he unto them, Therefore every scribe which is instructed unto the kingdom of heaven is like unto a man that is an householder, which bringeth forth out of his treasure things new and old.

\offertory{For I determined not to know any thing, save Jesus Christ, and him crucified.}

\secret
\lett{T}{hrough} these sacred mysteries vouchsafe unto us, O Lord: that as in thy mercy thou didst grant unto blessed Helen ever to carry thy Son crucified in her heart; so we may likewise continually hear him in our hearts. Who liveth and reigneth with thee.

\communion{I will go up to the palm tree, I will take hold of the boughs thereof.}

\postcommunion
\lett{G}{rant} unto us, O merciful God: that we who have been refreshed by the benefits of thy life-giving Cross on earth; may through the intercession of blessed Helen attain unto the eternal fruition of the same in heaven. Who livest.


\subby{St. Agapitus}
\feastday{{St. Agapitus}}
\fancyhead[RE,LO]{18 August}
\begin{inhead}
    {Memorial\\
18 August}
\end{inhead}

\begin{rubric}
	The propers come from the Second Common of a Martyr not a Bishop (p. \pageref{CommonMartyrNotBishopII}), except for that which followeth.
\end{rubric}

\begin{rubric}
	\textsc{Note,} When celebrated within the Octave of the Assumption, this Memorial is merely commemorated at Mass.
\end{rubric}

\collect
\lett{O}{Lord,} let thy Church trust with gladness in the advocacy of thy blessed Martyr Agapitus: that by his glorious prayers she may continue in devotion, and abide in safety. Through.

\readingcitation{Gospel}{John 12:24}
\lett{A}{t that time:} Jesus said unto his disciples: Verily, verily, I say unto you, Except a corn of wheat fall into the ground and die, it abideth alone: but if it die, it bringeth forth much fruit. He that loveth his life shall lose it; and he that hateth his life in this world shall keep it unto life eternal. If any man serve me, let him follow me; and where I am, there shall also my servant be: if any man serve me, him will my Father honour.

\secret
\lett{R}{eceive,} O Lord, the gifts which we offer on the solemnity of him, through whose advocacy we trust to be delivered. Through.

\postcommunion
\lett{O}{Lord,} who hast satisfied thy family with sacred gifts: we beseech thee; that we may at all times be comforted by the intercession of him whose festival we celebrate. Through.


\subby{Octave Day of the Assumption of the Blessed Virgin Mary}
\feastday{{Assumption Octave}}
\fancyhead[RE,LO]{22 August}
\begin{inhead}
    {Greater Double\\
22 August}
\end{inhead}
\begin{rubric}
	Mass as on the Feast with commemoration of Sts. Timothy, Hippolytus, \& Symphorian, as in the following Mass.
\end{rubric}
\begin{rubric}
	If today be Saturday, the anticipated Vigil of St. Bartholomew is kept, as is noted on the following day.
\end{rubric}

\subby{Sts. Timothy, Hippolytus, \& Symphorian}
\feastday{{Sts. Timothy \&c.}}
\fancyhead[RE,LO]{22 August}
\begin{inhead}
    {Memorial\\
22 August}
\end{inhead}

\begin{rubric}
	The propers come from the Third Common of Many Martyrs (p. \pageref{CommonMartyrsIII}), except for that which followeth.
\end{rubric}

\collect
\lett{W}{e} beseech thee, O Lord, graciously to impart unto us thy help: and, at the intercession of thy blessed Martyrs Timothy, Hippolytus, and Symphorian, stretch forth upon us the right hand of thy mercy. Through.

\secret
\lett{G}{rant,} O Lord, that like as thy dedicated people do acknowledge that in tribulation they have been succoured by the merits of thy Saints: so this oblation, which they offer unto thee in honour of the same, may be acceptable in thy sight. Through.

\postcommunion
\lett{O}{Lord} our God, who hast fulfilled us with the bounty of thy heavenly gift: we beseech thee, that, at the intercession of thy holy Martyrs, Timothy, Hippolytus, and Symphorian, we may ever live by the partaking of the same. Through.


\subby{Vigil of St. Bartholomew}
\feastday{{St. Bartholomew Vigil}}
\fancyhead[RE,LO]{23 August}
\begin{inhead}
    {Vigil\\
23 August}
\end{inhead}

\begin{rubric}
	The propers come from the Common of the Vigil of Apostles (p. \pageref{CommonVigilApostles}).\par
	\textsc{Note,} \nth{2} Collect of St. Mary in Eastertide (p. \SPMaryEaster) \& \nth{3} against the Persecutors of the Church (p. \SPAgainst) or for the Chief Bishop (p. \SPChiefBishop).
\end{rubric}


\bcpfeast{24 August. St. Bartholomew}{St. Bartholomew}{24 August}
%\supplement{24 August}{St. Bartholomew}{}

\begin{secrubric}
	The Daily Office propers are from the Common of Apostles (p. \pageref{CommonApostles}), with the following Antiphons.
\end{secrubric}

\properantiphon{Mag.}{They will deliver you up {\dag} to the councils, and they will scourge you in their synagogues; and ye shall be brought before governors and kings for my sake, for a testimony against them and the Gentiles.}

\properantiphon{Ben.}{Ye which have forsaken all, {\dag} and followed me, shall receive an hundredfold, and shall inherit everlasting life.}

\properantiphon{Mag.}{Be ye valiant in warfare, {\dag} and contend ye with the old serpent: and ye shall receive an eternal kingdom, alleluia.}


\subby{St. Zephyrinus}
\feastday{{St. Zephyrinus}}
\fancyhead[RE,LO]{26 August}
\begin{inhead}
    {Memorial\\
26 August}
\end{inhead}

\begin{rubric}
	The propers come from the Second Common of a Martyr Bishop out of Eastertide (p. \pageref{CommonMartyrBishopII}), except for that which followeth.
\end{rubric}

\collect
\lett{G}{rant,} we beseech thee, almighty God: that we who rejoice in the merits of blessed Zephyrinus, thy Martyr and Bishop, may be instructed by his example. Through.

\secret
\lett{S}{anctify,} O Lord, the gifts which we dedicate to thee: that at the intercession of blessed Zephyrinus thy Martyr and Bishop they may obtain for us thy gracious favour. Through.

\postcommunion
\lett{M}{ay} this communion, O Lord, cleanse us from guilt: and, at the intercession of blessed Zephyrinus, thy Martyr and Bishop, make us partakers of thy heavenly healing. Through.

%MANUAL ADJUSTMENT:
\clearpage
\subby{St. Augustine of Hippo}
\feastday{{St. Augustine Hippo}}
\fancyhead[RE,LO]{28 August}
\begin{inhead}
    {Greater Double\\
28 August}
\end{inhead}
\begin{rubric}
	The propers are from the Common of Doctors (p. \pageref{CommonDoctors}), except for that which followeth.
\end{rubric}

\introit
\lett{I}{n} the midst of the Church he opened his mouth: and the Lord filled him with the spirit of wisdom and of understanding: he clothed him with a robe of glory. \textit{Ps.} It is a good thing to give thanks unto the Lord: and to sing praises unto thy name, O most Highest.

\collect
\lett{A}{ssist} us, almighty God, in these our supplications: that as thou dost suffer us to put our trust and confidence in thy mercy, so, at the intercession of blessed Augustine thy Confessor and Bishop, thou wouldest graciously vouchsafe unto us the wonted effects of thy compassion. Through.

\begin{rubric}
    Commemoration is made of St. Hermes (p. \pageref{HermesCollect}).
\end{rubric}

\gradall{The mouth of the righteous is exercised in wisdom, and his tongue will be talking judgment. ℣. The law of his God is in his heart: and his goings shall not slide.}{Alleluia, alleluia. ℣. I have found David my servant, with my holy oil have I anointed him. Alleluia.}

\offertory{The righteous shall flourish like a palm-tree: and shall spread abroad like a cedar in Libanus.}

\secret
\lett{M}{ay} the devout prayers of thy Bishop and Doctor, Saint Augustine, never fail to succour us, O Lord: that they may render our oblations acceptable in thy sight; and may ever obtain for us thy merciful pardon. Through.

\begin{rubric}
    Commemoration is made of St. Hermes (p. \pageref{HermesSecret}).
\end{rubric}

\communion{A faithful and wise servant, whom the lord hath made ruler over his household: to give them their portion of meat in due season.}

\postcommunion
\lett{W}{e} beseech thee, O Lord, that blessed Augustine, thy Bishop and illustrious Doctor, may stand before thee as our advocate: that these thy sacrifices may avail for our salvation. Through.

\begin{rubric}
    Commemoration is made of St. Hermes (p. \pageref{HermesPostcommunion}).
\end{rubric}


\subby{St. Hermes}
\feastday{{St. Hermes}}
\fancyhead[RE,LO]{28 August}
\begin{inhead}
    {Memorial\\
28 August}
\end{inhead}

\begin{rubric}
	The propers are from the Second Common of a Martyr not a Bishop (p. \pageref{CommonMartyrNotBishopII}), except for that which followeth.
\end{rubric}

\collect\label{HermesCollect}
\lett{O}{God,} who didst strengthen blessed Hermes thy Martyr with the virtue of constancy in his passion: Grant unto us by his example; to despise for love of thee the prosperity of the world, and to fear none of its adversities. Through.

\secret\label{HermesSecret}
\lett{W}{e} offer unto thee, O Lord, this sacrifice of praise in commemoration of thy Saints: grant, we beseech thee; that as it hath bestowed glory on them, so may it avail for our salvation. Through.

\postcommunion\label{HermesPostcommunion}
\lett{O}{Lord,} who hast fulfilled us with thy heavenly benediction, we beseech thy mercy: that, at the intercession of thy blessed Martyr Hermes, this service of our lowliness may avail for the comforting of our souls. Through.


\bcpfeast{29 August. Decollation of St. John Baptist}{Decollation}{29 August}
%\supplement{29 August}{Beheading}{of St. John Baptist}

\begin{rubric}
	The Hymns and Versicle are from the Common of Apostles (p. \pageref{CommonApostles}), with the Antiphons as followeth.
\end{rubric}

\properantiphon{Mag. \& Ben.}{Herod sent {\dag} an executioner, and commanded him to behead John in the prison: and when his disciples heard of it, they came, and took up his body, and laid it in a tomb.}

\properantiphon{Mag.}{The king, the unbeliever, {\dag} sent his detestable servants, and commanded them to behead John the Baptist.}


\subby{St. Sabina}
\feastday{{St. Sabina}}
\fancyhead[RE,LO]{29 August}
\begin{inhead}
    {Memorial\\
29 August}
\end{inhead}

\begin{rubric}
	The propers are from the Common of a Martyr not a Virgin (p. \pageref{CommonMartyrNotVirgin}), except for that which followeth.
\end{rubric}


\subby{Sts. Felix \& Adauctus}
\feastday{{Sts. Felix \& Adauctus}}
\fancyhead[RE,LO]{30 August}
\begin{inhead}
    {Memorial\\
30 August}
\end{inhead}

\introit
\lett{L}{et} the people tell of the wisdom of the Saints, and let the church shew forth their praise: their names shall live for evermore. \textit{Ps.} Rejoice in the Lord, O ye righteous: for it becometh well the just to be thankful.

\collect
\lett{O}{Lord,} we humbly entreat thy majesty: that, like as thou dost continually gladden us with the commemoration of thy Saints; so thou wouldest evermore defend us with their supplication. Through.

\begin{rubric}
	The Epistle is the third additional Epistle from the Third Common of Many Martyrs (p. \pageref{Wisdom1017}).
\end{rubric}

\gradall{The souls of the just are in the hand of God, and there shall no torment of malice touch them. ℣. In the eyes of the unwise they seemed to die: but they are in peace.}{Alleluia, alleluia. ℣. The righteous shall shine forth, and run to and fro like sparks among the stubble for ever. Alleluia.}

\begin{rubric}
	The Gospel is the fifth additional Gospel from the Third Common of Many Martyrs (p. \pageref{Luke1016}).
\end{rubric}

\offertory{Be glad, O ye righteous, and rejoice in the Lord: and be joyful, all ye that are true of heart.}

\secret
\lett{L}{ook} down, O Lord, upon the sacrifices of thy people: that, as with devout hearts they celebrate them to the honour of thy Saints, so they may perceive them to be profitable to their salvation. Through.

\communion{What I tell you in darkness, that speak ye in light, saith the Lord: and what ye hear in the ear, that preach ye upon the housetops.}

\postcommunion
\lett{W}{e} beseech thee, O Lord: that we, being filled with the sacred gifts; may at the intercession of thy Saints ever continue in thanksgiving for the same. Through.


\subby{St. Aidan of Lindisfarne}
\feastday{{St. Aidan}}
\fancyhead[RE,LO]{31 August}
\begin{inhead}
    {Memorial\\
31 August}
\end{inhead}

\begin{rubric}
	The propers are from the First Common of a Confessor Bishop (p. \pageref{CommonConfessorBishopI}).
\end{rubric}


\subby{St. Giles}
\feastday{{St. Giles}}
\fancyhead[RE,LO]{1 September}
\begin{inhead}
    {Memorial\\
1 September}
\end{inhead}

\begin{rubric}
	The propers are from the Common of Abbots (p. \pageref{CommonAbbots}), except for that which followeth.\par
	Commemoration is made of the Twelve Brothers.
\end{rubric}


\subby{Twelve Holy Brothers}
\feastday{{12 Holy Brothers}}
\fancyhead[RE,LO]{1 September}
\begin{inhead}
    {Memorial\\
1 September}
\end{inhead}

\introit
\lett{T}{he} just cry, and the Lord heareth them: and delivereth them out of all their troubles. \textit{Ps.} I will alway give thanks unto the Lord; his praise shall ever be in my mouth.

\collect
\lett{O}{Lord,} let the crown of the brethren thy Martyrs, cause us to rejoice: that we may thereby be strengthened and increased in our faith; and comforted by their manifold intercession. Through.

\begin{rubric}
	The Epistle is the fifth additional Epistle of the Third Common of Many Martyrs (p. \pageref{Hebrews1133}).
\end{rubric}

\gradall{Behold, how good and joyful a thing it is, brethren, to dwell together in unity! ℣. It is like the precious ointment upon the head, that ran down unto the beard, even unto Aaron's beard.}{Alleluia, alleluia. ℣. This is the true brotherhood which overcame the wickedness of the world: which followed Christ, gaining heaven's glorious realms. Alleluia.}

\begin{rubric}
	The Gospel is from the Third Common of Many Martyrs (p. \pageref{CommonMartyrsII}).
\end{rubric}

\offertory{Be glad, O ye righteous, and rejoice in the Lord: and be joyful, all ye that are true of heart.}

\secret
\lett{G}{rant,} O Lord, that we may with devout hearts celebrate thy mysteries in honour of thy holy Martyrs: and thereby obtain an increase both of protection and joy. Through.

\communion{Whosoever shall do the will of my Father which is in heaven: the same is my brother, and sister, and mother, saith the Lord.}

\postcommunion
\lett{G}{rant,} we beseech thee, almighty God: that growing in virtue we may follow the faith of them whose memory we recall by the partaking of this sacrament. Through.


\subby{St. Stephen of Hungary}
\feastday{{St. Stephen Hungary}}
\fancyhead[RE,LO]{2 September}
\begin{inhead}
	{Memorial\\
		2 September}
\end{inhead}

\begin{rubric}
	The propers are from the First Common of Confessor not Bishop (p. \pageref{CommonConfessorNotBishopI}), except for the following.
\end{rubric}

\collect
\lett{G}{rant,} we beseech thee, almighty God unto thy Church: that as thy blessed Confessor Stephen, while he reigned on earth, did spread abroad her faith, so she may be found worthy to have him for her glorious defender in the heavens. Through.

\begin{rubric}
	The Gospel is the additional Gospel from the Second Common of a Confessor not a Bishop (p. \pageref{Luke1912}).
\end{rubric}

\secret
\lett{A}{lmighty} God, look upon the sacrifices which we offer: and grant; that we, who celebrate the mysteries of the Passion of the Lord, may imitate that which we perform. Through the same.

\postcommunion
\lett{G}{rant,} we beseech thee, almighty God: that, as thy blessed Confessor Stephen, for the propagation of thy faith, was counted worthy to pass from an earthly kingdom to the glory of the heavenly realm; so we may imitate his faith with due devotion. Through.


\subby{St. Gorazde of Prague}
\feastday{{St. Gorazde}}
\fancyhead[RE,LO]{4 September}
\begin{inhead}
	{Double\\
		4 September}
\end{inhead}

\begin{rubric}
	The propers are from the Second Common of a Martyr Bishop (p. \pageref{CommonMartyrBishopII}.
\end{rubric}


\bcpfeast{8 September. Nativity of the Blessed Virgin Mary}{Nativity B.V.M.}{8 September}
%\supplement{8 September}{Nativity}{of the Blessed Virgin Mary}

\begin{rubric}
	The Hymns are from the Common of the Blessed Virgin Mary (p. \pageref{CommonBVM}), with the following Versicles \& Antiphons.
\end{rubric}

    ℣. To-day is the Nativity of the holy Virgin Mary.
    
	℟. Whose glorious life illumineth all the churches.

\properantiphon{Mag.}{Let us celebrate {\dag} the worshipful Nativity of the blessed and glorious Virgin Mary: for she hath obtained the dignity of Motherhood, and yet lost not her maiden purity.}\\

    ℣. To-day is the Nativity of the holy Virgin Mary.
    
	℟. Whose glorious life illumineth all the churches.

\properantiphon{Ben.}{To-day let us celebrate {\dag} with due solemnity the Nativity of God's Mother, the ever Virgin Mary: from whom the Son of the Highest proceeded, alleluia.}\\

    ℣. To-day is the Nativity of the holy Virgin Mary.
    
	℟. Whose glorious life illumineth all the churches.

\properantiphon{Mag.}{Thy Nativity, {\dag} O Virgin Mother of God, hath proclaimed joyful tidings unto all the world: for out of thee hath arisen the Sun of righteousness, even Christ our God: who, taking away the curse, hath bestowed a blessing; and confounding death, hath given unto us life everlasting.}


\subby{St. Hadrian}
\feastday{{St. Hadrian}}
\fancyhead[RE,LO]{8 September}
\begin{inhead}
	{Memorial\\
		8 September}
\end{inhead}

\begin{rubric}
	The propers are from the First Common of a Martyr not a Bishop (p. \pageref{CommonMartyrNotBishopI}).
\end{rubric}

%MANUAL ADJUSTMENT:
\clearpage
\subby{St. Gorgonius}
\feastday{{St. Gorgonius}}
\fancyhead[RE,LO]{9 September}
\begin{inhead}
	{Memorial\\
		9 September}
\end{inhead}

\begin{rubric}
	The propers are from the Second Common of a Martyr not a Bishop (p. \pageref{CommonMartyrNotBishopII}), except for the following.
\end{rubric}

\collect
\lett{L}{et} thy Saint Gorgonius, O Lord, gladden us by his intercession; and make us to rejoice in this holy solemnity. Through.

\secret
\lett{O}{Lord,} let thy holy Martyr Gorgonius so intercede for us: that this oblation of our service may be acceptable unto thee. Through.

\postcommunion
\lett{G}{rant,} O God, that thy household may be quickened and refreshed by thine eternal goodness: and in thy Martyr Gorgonius be continually nourished with the sweet savour of Christ, thy Son. Who liveth and reigneth with thee.


\subby{Sts. Protus \& Hyacinth}
\feastday{{Sts. Protus \& Hyacinth}}
\fancyhead[RE,LO]{11 September}
\begin{inhead}
	{Memorial\\
		11 September}
\end{inhead}

\begin{rubric}
	The propers are from the Third Common of Many Martyrs out of Eastertide (p. \pageref{CommonMartyrsIII}), except for the following.
\end{rubric}

\collect
\lett{O}{Lord,} let the meritorious confession of thy blessed Martyrs, Protus and Hyacinth, comfort us; and may their loving intercession ever defend us. Through.

\secret
\lett{G}{rant,} we beseech thee, O Lord, that the oblation of our bounden service which we ofier unto thee for the commemoration of thy holy Martyrs, Protus and Hyacinth; may effectually avail for our healing unto everlasting salvation. Through.

\postcommunion
\lett{W}{e} beseech thee, O Lord, that by the supplication of thy blessed Martyrs, Protus and Hyacinth: thy holy mysteries which we have received may avail for our cleansing. Through.


\subby{Most Holy Name of Mary}
\feastday{{Name of Mary}}
\fancyhead[RE,LO]{12 September}
\begin{inhead}
	{Greater Double\\
		12 September}
\end{inhead}

\begin{rubric}
	The Daily Office propers are from the Common of the Blessed Virgin Mary (p. \pageref{CommonBVM}), except for the antiphon for I Evensong, as followeth.
\end{rubric}

%Last clause translated from the Diurnale Monasticum:
\antiphon{Mag.}{O holy Mary, {\dag} help thou the suffering, strengthen the faint-hearted, comfort the sorrowful; pray for the people, entreat for the clergy, intercede for all womankind vowed unto God: may all acknowledge the help of thy prayer, who celebrate the commemoration of thine holy Name.}\par\noindent

\introit
\lett{A}{ll} the rich among the people shall make their supplication before thee: the Virgins that be her fellows shall be brought unto the King: they that bear her company shall be brought unto thee with joy and gladness. \textit{Ps.} My heart is inditing of a good matter: I speak of the things which I have made unto the King.

\collect
\lett{G}{rant,} we beseech thee, almighty God: that thy faithful people, who rejoice in the Name and protection of the most Holy Virgin Mary; may by her loving intercession be delivered from all evils upon earth, and be found worthy to attain unto everlasting joys in heaven. Through.

\readingcitation{Epistle}{Ecclesiasticus 24:17}
%KJV:
%\lett{I}{as} the vine brought forth pleasant savour, and my flowers are the fruit of honour and riches. I am the mother of fair love, and fear, and knowledge, and holy hope: I therefore, being eternal, am given to all my children which are named of him. Come unto me, all ye that be desirous of me, and fill yourselves with my fruits. For my memorial is sweeter than honey, and mine inheritance than the honeycomb. They that eat me shall yet be hungry, and they that drink me shall yet be thirsty. He that obeyeth me shall never be confounded, and they that work by me shall not do amiss.
%RV:
\lett{A}{s} the vine I put forth grace; And my flowers are the fruit of glory and riches. Come unto me, ye that are desirous of me, And be ye filled with my produce. For my memorial is sweeter than honey, And mine inheritance than the honeycomb. They that eat me shall yet be hungry; And they that drink me shall yet be thirsty. He that obeyeth me shall not be ashamed; And they that work in me shall not do amiss.

\gradall{Blessed and venerable art thou, O Virgin Mary: who without spot wast found the Mother of the Saviour. ℣. Virgin, Mother of God, he whom the world containeth not, being made man lay hid in thy womb.}{Alleluia, alleluia. ℣. After child-birth, O Virgin, thou didst remain inviolate: Mother of God, intercede for us. Alleluia.}

\readingcitation{Gospel}{Luke 1:26}
\lett{A}{t that time:} The angel Gabriel was sent from God unto a city of Galilee, named Nazareth, to a virgin espoused to a man whose name was Joseph, of the house of David; and the virgin’s name was Mary. And the angel came in unto her, and said, Hail, thou that art highly favoured, the Lord is with thee: blessed art thou among women. And when she saw him, she was troubled at his saying, and cast in her mind what manner of salutation this should be. And the angel said unto her, Fear not, Mary: for thou hast found favour with God. And, behold, thou shalt conceive in thy womb, and bring forth a son, and shalt call his name \divineName{Jesus}. He shall be great, and shall be called the Son of the Highest: and the Lord God shall give unto him the throne of his father David: and he shall reign over the house of Jacob for ever; and of his kingdom there shall be no end. Then said Mary unto the angel, How shall this be, seeing I know not a man? And the angel answered and said unto her, The Holy Ghost shall come upon thee, and the power of the Highest shall overshadow thee: therefore also that holy thing which shall be born of thee shall be called the Son of God. And, behold, thy cousin Elisabeth, she hath also conceived a son in her old age: and this is the sixth month with her, who was called barren. For with God nothing shall be impossible. And Mary said, Behold the handmaid of the Lord; be it unto me according to thy word.

\offertory{Hail, Mary, full of grace; the Lord is with thee: blessed art thou among women, and blessed is the fruit of thy womb.}

\secret
\lett{T}{hrough} thy mercy, O Lord, and the intercession of blessed Mary ever Virgin, may this oblation avail for our prosperity and peace, both now and for ever. Through.

\communion{Blessed is the womb of the Virgin Mary, that bore the Son of the everlasting Father.}

\postcommunion
\lett{G}{rant,} we beseech thee, O Lord: that we, who have received these aids to our salvation, may at all times and in all places be protected through the advocacy of blessed Maty ever Virgin; in whose honour we have made these offerings to thy majesty. Through.


\bcpfeast{14 September. Exaltation of the Holy Cross}{Exaltation of the Cross}{14 September}
%\supplement{14 September}{Exaltation}{of the Holy Cross}

\begin{rubric}
	The Daily Office propers are as in the Feast of the Invention of the Holy Cross, 3 May (p. \pageref{InventionoftheCross}), omitting the \emph{Alleluia}. If II Evensong be said, the following Antiphon is said.
\end{rubric}

\antiphon{Mag.}{O Cross {\dag} exceeding blessed, which alone wast counted worthy to bear the Lord, the King of heaven, alleluia.}


\bcpfeast{Ember Wednesday in Autumn}{Autumnal Ember Wednesday}{}

\properantiphon{Ben.}{This kind {\dag} of evil spirit can come forth by nothing but by prayer and fasting.}


\bcpfeast{Ember Friday in Autumn}{Autumnal Ember Friday}{}

\properantiphon{Ben.}{A woman {\dag} in the city, which was a sinner, stood at the Lord's feet behind him, and began to wash his feet with tears, and did wipe them with the hairs of her head, and kissed his feet, and anointed them with the ointment.}


\bcpfeast{Ember Saturday in Autumn}{Autumnal Ember Saturday}{}

\properantiphon{Ben.}{Give light, O Lord, {\dag} to them that sit in darkness: and gudie our feet into the way of peace, thou God of Israel.}


\bcpfeast{15 September. Seven Sorrows of the Blessed Virgin Mary}{Seven Sorrows}{15 September}
%\supplement{15 September}{Seven Sorrows}{of the Blessed Virgin Mary}

\subbysub{I Evensong}\label{SevenSorrowsEvensong}

\gregorioscore{resources/gabc/ProperTime/SevenSorrowsEvensong.gabc}

    ℣. Pray for us, O Queen of Martyrs.

	℟. Who didst stand by the Cross of Jesus.

\properantiphon{Mag.}{Look not {\dag} upon me, because I am black, because the sun hath looked upon me: my mother's children were angry with me.}

\subbysub{Mattins}

\invitatoryhymn\label{SevenSorrowsInvitatory}

\gregorioscore{resources/gabc/ProperTime/SevenSorrowsInvitatory.gabc}

\officehymn\label{SevenSorrowsMattins}

\gregorioscore{resources/gabc/ProperTime/SevenSorrowsMattins.gabc}

    ℣. O Virgin Mary, by thy sorrow's might.

	℟. Make us rejoice in heaven's kingdom bright.

\properantiphon{Ben.}{Come ye, {\dag} and let us go up to the mountain of the Lord, and see if there be any sorrow like unto my sorrow.}

\subbysub{II Evensong}

\begin{rubric}
	The Office Hymn \& Versicle are as in I Evensong, with the following Antiphon.
\end{rubric}

\properantiphon{Mag.}{Sorrow oppresseth me, {\dag} my face is swollen with weeping, and on my eyelids is the shadow of death.}


\subby{St. Nicomedes}
\feastday{{St. Nicomedes}}
\fancyhead[RE,LO]{16 September}
\begin{inhead}
	{Memorial\\
		15 September}
\end{inhead}

\begin{rubric}
	The propers are from the First Common of a Martyr not a Bishop (p. \pageref{CommonMartyrNotBishopI}), except for the following.
\end{rubric}

\collect
\lett{A}{ssist,} O Lord, thy people: that as they do profit by the glorious merits of blessed Nicomedes thy Martyr, so his advocacy may at all times succour them to the obtaining of thy mercy. Through.

\secret
\lett{M}{ercifully} receive, O Lord, the gifts which we offer: and let the prayer of the blessed Martyr Nicomedes commend them unto thy majesty. Through.

\postcommunion
\lett{C}{leanse} us, O Lord, by the sacraments which we have received: that through the intercession of blessed Nicomedes, thy Martyr, they may set us free from all our offences. Through.


\subby{Sts. Cornelius \& Cyprian}
\feastday{{Sts. Cornelius \& Cyprian}}
\fancyhead[RE,LO]{16 September}
\begin{inhead}
	{Double\\
		16 September}
\end{inhead}

\begin{rubric}
	The propers are from the First Common of Many Martyrs out of Eastertide (p. \pageref{CommonMartyrsI}), except for the following.
\end{rubric}

\collect
\lett{P}{rotect} us, O Lord, we beseech thee, who observe the feast of thy blessed Martyrs and Bishops Cornelius and Cyprian: and grant that by their meritorious supplication we may find favour in thy sight. Through.

\begin{rubric}
    Commemoration is made of St. Euphemia, Lucy, and Geminian (p. \pageref{EuphemiaCollect}).
\end{rubric}

\secret
\lett{A}{ssist} us mercifully, O Lord, in these our supplications which we make before thee in remembrance of thy Saints: that we who trust not in our own righteousness may be succoured by the merits of them that have found favour in thy sight. Through.

\begin{rubric}
    Commemoration is made of St. Euphemia, Lucy, and Geminian (p. \pageref{EuphemiaSecret}).
\end{rubric}

\postcommunion
\lett{O}{Lord,} who hast fulfilled us with saving mysteries, we beseech thee: that we may be aided by the prayers of those whose festival we celebrate. Through.

\begin{rubric}
    Commemoration is made of St. Euphemia, Lucy, and Geminian (p. \pageref{EuphemiaPostcommunion}).
\end{rubric}


\subby{Sts. Euphemia, Lucy, \& Geminian}
\feastday{{Sts. Euphemia \&c.}}
\fancyhead[RE,LO]{16 September}
\begin{inhead}
	{Memorial\\
		16 September}
\end{inhead}

\begin{rubric}
	The propers are from the First Common of Many Martyrs (p. \pageref{CommonMartyrsI}), except for the following.
\end{rubric}

\collect\label{EuphemiaCollect}
\lett{G}{rant,} O Lord, that our prayers in this time of our rejoicing may be brought to good effect: that as with yearly service we recall the day of the passion of thy holy Martyrs, Euphemia, Lucy, and Geminian, so we may imitate the steadfastness of their faith. Through.

\begin{rubric}
	The Gospel is from the Second Common of Many Martyrs (p. \pageref{CommonMartyrsII}).
\end{rubric}

\secret\label{EuphemiaSecret}
\lett{G}{raciously} hearken, we beseech thee, O Lord, unto the prayers of thy people: and make us to rejoice in the intercession of those, whose festival thou dost suffer us to celebrate. Through.

\postcommunion\label{EuphemiaPostcommunion}
\lett{O}{Lord,} graciously hear our prayers: that we, who solemnly observe the feast of thy holy Martyrs, Euphemia, Lucy, and Geminian, may be succoured by their continual help. Through.


\subby{Sts. Januarius \& Companions}
\feastday{{Sts. Januarius \&c.}}
\fancyhead[RE,LO]{19 September}
\begin{inhead}
	{Double\\
		19 September}
\end{inhead}

\begin{rubric}
	The propers are from the Third Common of Many Martyrs (p. \pageref{CommonMartyrsIII}), except for the first additional Gospel of the same Common (p. \pageref{Matthew243}).
\end{rubric}
\begin{rubric}
	\textsc{Note,} Commemoration is make of St. Theodore of Canterbury.
\end{rubric}


\subby{St. Theodore of Canterbury}
\feastday{{St. Theodore Canterbury}}
\fancyhead[RE,LO]{19 September}
\begin{inhead}
	{Memorial\\
		19 September}
\end{inhead}

\begin{rubric}
	The propers are from the First Common of a Confessor Bishop (p. \pageref{CommonConfessorBishopI}).
\end{rubric}

\begin{rubric}
	If today be Saturday, the anticipated Vigil of St. Matthew is kept.
\end{rubric}


\subby{Sts. Eustace \& Companions}
\feastday{{Sts. Eustace \&c.}}
\fancyhead[RE,LO]{20 September}
\begin{inhead}
	{Double\\
		20 September}
\end{inhead}

\begin{rubric}
	The propers are from the Second Common of Many Martyrs (p. \pageref{CommonMartyrsII}).
\end{rubric}


\subby{Vigil of St. Matthew}
\feastday{{St. Matthew Vigil}}
\fancyhead[RE,LO]{20 September}
\begin{inhead}
	{Vigil\\
		20 September}
\end{inhead}

\begin{rubric}
	The propers are from the Common of Vigils of the Apostles (p. \pageref{CommonVigilApostles}) (with the \nth{2} Prayers of St. Mary in Eastertide (p. \SPMaryEaster) and the \nth{3} against the Persecutors of the Church (p. \SPAgainst) or for the Chief Bishop (p. \SPChiefBishop)), except for the following Gospel.
\end{rubric}

\readingcitation{Gospel}{Luke 5:27}
\lett{A}{t that time:} Jesus saw a publican, named Levi, sitting at the receipt of custom: and he said unto him, Follow me. And he left all, rose up, and followed him. And Levi made him a great feast in his own house: and there was a great company of publicans and of others that sat down with them. But their scribes and Pharisees murmured against his disciples, saying, Why do ye eat and drink with publicans and sinners? And Jesus answering said unto them, They that are whole need not a physician; but they that are sick. I came not to call the righteous, but sinners to repentance.


\bcpfeast{21 September. St. Matthew}{St. Matthew}{21 September}

\begin{secrubric}
	The Daily Office propers are of the Common of Apostles (p. \pageref{CommonApostles}).
\end{secrubric}


\subby{Sts. Maurice \& Companions}
\feastday{{Sts. Maurice \&c.}}
\fancyhead[RE,LO]{22 September}
\begin{inhead}
	{Memorial\\
		22 September}
\end{inhead}

\begin{rubric}
	The propers are from the First Common of Many Martyrs (p. \pageref{CommonMartyrsI}), except for the following.
\end{rubric}

\collect
\lett{G}{rant,} we beseech thee, almighty God: that the solemn festival of thy holy Martyrs, Maurice and his Companions, may in such wise gladden us: that, as we do lean upon their advocacy, so we may glory in their heavenly birth. Through.

\begin{rubric}
	The Epistle is the sixth additional Epistle from the Third Common of Many Martyrs (p. \pageref{Revelation713}).
\end{rubric}

\secret
\lett{R}{egard,} we beseech thee, O Lord, the gifts which we offer unto thee in commemoration of thy holy Martyrs, Maurice and his Companions: and grant; that through the intercession of them, for whose sake they are acceptable unto thee, they may be profitable unto us for evermore. Through.

\postcommunion
\lett{O}{Lord,} who hast refreshed us with the gladness of thy heavenly sacraments: we humbly entreat thee; that we may be protected by the succour of them in whose triumphs we glory. Through.


\subby{Pope St. Linus}
\feastday{{Pope St. Linus}}
\fancyhead[RE,LO]{23 September}
\begin{inhead}
	{Memorial\\
		23 September}
\end{inhead}

\begin{rubric}
	The propers are from the First Common of a Martyr Bishop (p. \pageref{CommonMartyrBishopI}), except for the following.
\end{rubric}

\collect
\lett{O}{God,} who makest us glad with the yearly solemnity of blessed Linus thy Martyr and Bishop: mercifully grant; that, as we now celebrate his birthday, so we may likewise rejoice in his protection. Through.

\begin{rubric}
    Commemoration is made of St. Thecla (p. \pageref{TheclaCollect}).
\end{rubric}

\secret
\lett{S}{anctify,} O Lord, the gifts which we dedicate to thee: that at the intercession of blessed Linus, thy Martyr and Bishop, they may obtain for us thy gracious favour. Through.

\begin{rubric}
    Commemoration is made of St. Thecla (p. \pageref{TheclaSecret}).
\end{rubric}

\postcommunion
\lett{M}{ay} this communion, O Lord, cleanse us from guilt: and, at the intercession of blessed Linus, thy Martyr and Bishop, make us partakers of thy heavenly healing. Through.

\begin{rubric}
    Commemoration is made of St. Thecla (p. \pageref{TheclaPostcommunion}).
\end{rubric}


\subby{St. Thecla}
\feastday{{St. Thecla}}
\fancyhead[RE,LO]{23 September}
\begin{inhead}
	{Memorial\\
		23 September}
\end{inhead}

\begin{rubric}
	The propers are from the First Common of a Virgin Martyr (p. \pageref{CommonVirginMartyrI}), except for the following.
\end{rubric}

\collect\label{TheclaCollect}
\lett{G}{rant,} we beseech thee, almighty God: that we, who celebrate the birthday of blessed Thecla, thy Virgin and Martyr; may both rejoice in her yearly solemnity, and likewise profit by the example of her faith. Through.

\secret\label{TheclaSecret}
\lett{R}{eceive,} O Lord, the gifts which we offer on the solemnity of blessed Thecla, thy Virgin and Martyr: through whose advocacy we trust to be delivered. Through.

\postcommunion\label{TheclaPostcommunion}
\lett{M}{ay} the mysteries which we have received be for our succour, O Lord: and at the intercession of blessed Thecla, thy Virgin and Martyr, cause us to rejoice in thy continual protection. Through.


\subby{Sts. Cyprian \& Justina}
\feastday{{Sts. Cyprian \& Justina}}
\fancyhead[RE,LO]{26 September}
\begin{inhead}
	{Memorial\\
		26 September}
\end{inhead}

\begin{rubric}
	The propers are from the Third Common of Many Martyrs (p. \pageref{CommonMartyrsIII}), except for the following.
\end{rubric}

\collect
\lett{O}{Lord,} let the protection of thy blessed Martyrs, Cyprian and Justina, continually defend us: forasmuch as thou failest not to look with mercy on those to whom thou dost grant the succour of thy assistance. Through.

\secret
\lett{W}{e} beseech thee, O Lord, that the gifts which we offer unto thee of our bounden duty and service may be acceptable unto thee for the honour of thy Just ones: and by thy mercy profitable unto us for our salvation. Through.

\postcommunion
\lett{G}{rant} to us, we beseech thee, O Lord: at the intercession of thy holy Martyrs Cyprian and Justina; that those things which we touch with our mouths we may receive in purity of heart. Through.


\subby{Sts. Cosmas \& Damian}
\feastday{{Sts. Cosmas \& Damian}}
\fancyhead[RE,LO]{27 September}
\begin{inhead}
	{Memorial\\
		27 September}
\end{inhead}

\introit
\lett{L}{et} the people tell of the wisdom of the Saints, and let the Church shew forth their praise: their names shall live for evermore. \textit{Ps.} Rejoice in the Lord, O ye righteous: for it becometh well the just to be thankful.

\collect
\lett{G}{rant,} we beseech thee, almighty God: that we, who observe the heavenly birthday of thy holy Martyrs, Cosmas and Damian, may by their intercession be delivered from all evils that beset us. Through.

\begin{rubric}
	The Epistle is from the Second Common of Many Martyrs (p. \pageref{CommonMartyrsII}).
\end{rubric}

\gradall{The righteous cry, and the Lord heareth them: and delivereth them out of all their troubles. ℣. The Lord is nigh unto them that are of a contrite heart: and will save such as be of an humble spirit.}{Alleluia, alleluia. ℣. This is the true brotherhood, which overcame the wickedness of the world: which followed Christ, gaining heaven’s glorious realms. Alleluia.}

\begin{rubric}
	The Gospel is from the Second Common of Many Martyrs (p. \pageref{CommonMartyrsII}).
\end{rubric}

\offertory{All they that love thy Name shall be joyful in thee: for thou, Lord, wilt give thy blessing unto the righteous: and with thy favourable kindness, O Lord, wilt thou defend us as with a shield.}

\secret
\lett{M}{ay} the devout prayers of thy Saints never fail us, O Lord: that they may both render our oblations acceptable, and ever obtain for us thy merciful pardon. Through.

\communion{The dead bodies of thy servants, O Lord, have they given to be meat unto the fowls of the air, and the flesh of thy Saints unto the beasts of the land: according to the greatness of thy power preserve thou those that are appointed to die.}

\postcommunion
\lett{W}{e} beseech thee, O Lord, that thy people may be protected by the prayers which thy Saints do offer, and by this heavenly banquet, whereof thou hast suffered them to be partakers. Through.


\subby{St. Wenceslaus}
\feastday{{St. Wenceslaus}}
\fancyhead[RE,LO]{28 September}
\begin{inhead}
	{Memorial\\
		28 September}
\end{inhead}

\begin{rubric}
	The propers are from the First Common of a Martyr not a Bishop (p. \pageref{CommonMartyrNotBishopI}), except for the following.
\end{rubric}

\collect
\lett{O}{God,} who through the victory of martyrdom didst translate blessed Wenceslas from an earthly principality to the glory of heaven: defend us through his prayers from all adversity; and grant us likewise to rejoice in his fellowship. Through.


\bcpfeast{29 September. St. Michael}{St. Michael}{29 September}\label{SaintMichael}
%\supplement{29 September}{St. Michael}{}

\subbysub{I Evensong}\label{MichaelEvensong}

\gregorioscore{resources/gabc/ProperTime/MichaelEvensong.gabc}

    ℣. An Angel stood at the altar of the temple.

	℟. Having in his hand a golden censer.

\properantiphon{Mag.}{While John beheld {\dag} the sacred mystery, the Archangel Michael sounded the trumpet: Forgive us, O Lord our God, that openest the book, and loosest the seals thereof, alleluia.}

\subbysub{Mattins}

\begin{rubric}
	The Invitatory Hymn is as in I Evensong.
\end{rubric}

\officehymn\label{MichaelMattins}

\gregorioscore{resources/gabc/ProperTime/MichaelMattins.gabc}

    ℣. An Angel stood at the altar of the temple.

	℟. Having in his hand a golden censer.

\properantiphon{Ben.}{There was silence in heaven {\dag} while the dragon waged war: and Michael fought against him, and had the victory, alleluia.}

\subbysub{II Evensong}

\begin{rubric}
	The Office Hymn is as in I Evensong, with the following Versicle \& Antiphon.
\end{rubric}

    ℣. In the presence of the Angels I will sing praise unto thee, O my God.

	℟. I will worship toward thy holy temple, and praise thy Name.

\properantiphon{Mag.}{O Prince most glorious, {\dag} Michael the Archangel, keep us in remembrance: here and everywhere, always, entreat the Son of God for us, alleluia, alleluia.}

%MANUAL ADJUSTMENT:
\clearpage
\subby{St. Jerome}
\feastday{{St. Jerome}}
\fancyhead[RE,LO]{30 September}
\begin{inhead}
	{Greater Double\\
		30 September}
\end{inhead}

\begin{rubric}
	The propers are from the First Common of Many Martyrs out of Eastertide (p. \pageref{CommonMartyrsI}), except for the following.\par
	\textsc{Note,} The Creed is said.
\end{rubric}

\collect
\lett{O}{God,} who for the exposition of the sacred Scriptures didst bestow upon thy Church blessed Jerome, thy Confessor and most illustrious Doctor: grant, we beseech thee; that, by the intercession of his merits, we may through thine assistance be enabled to perform those things which he taught both in word and deed. Through.

\begin{rubric}
    Commemoration is made of St. Gregory of Armenia (p. \pageref{ArmeniaCollect}).
\end{rubric}

\secret
\lett{G}{rant} us, we beseech thee, O Lord, through these heavenly gifts to serve thee in freedorn of spirit: that the gifts which we offer may, through the mediation of blessed Jerome, thy Confessor, work in us both healing and glory. Through.

\begin{rubric}
    Commemoration is made of St. Gregory of Armenia (p. \pageref{ArmeniaSecret}).
\end{rubric}

\postcommunion
\lett{W}{e} beseech thee, O Lord, that we whom thou hast fulfilled with heavenly nourishment: may through the mediation of blessed Jerome, thy Confessor, be found worthy to obtain the grace of thy loving-kindness. Through.

\begin{rubric}
    Commemoration is made of St. Gregory of Armenia (p. \pageref{ArmeniaPostcommunion}).
\end{rubric}


\subby{St. Gregory of Armenia}
\feastday{{St. Gregory Armenia}}
\fancyhead[RE,LO]{30 September}
\begin{inhead}
	{Memorial\\
		30 September}
\end{inhead}

\begin{rubric}
	The propers are from the First Common of a Confessor Bishop (p. \pageref{CommonConfessorBishopI}), except for the following.
\end{rubric}

%Propers changed to indicate that he is a Confessor not a Martyr.
\collect\label{ArmeniaCollect}
\lett{O}{God,} who, through thy blessed Bishop and Confessor Gregory, didst grant that the people and King of Armenia to receive the light of true faith: likewise grant unto thy Church to rejoice in such mighty triumphs, and, by the merits and intercession of the same, to be succoured before thee. Through.
%Deus, qui per beátum Gregórium Pontíficem et Mártyrem tuum Arméniæ gentis pópulum régémque veræ fídei lucem recípere tribuísti: da Ecclésiæ tuæ de tantis gaudére triúmphis, et apud te méritis ejusdem et précibus adjuvári. Per Dóminum.

\secret\label{ArmeniaSecret}
\lett{G}{raciously} receive, O Lord, the sacrifices dedicated unto thee, by the merits of thy blessed Bishop and Confessor Gregory: and grant that they may avail us unto everlasting help. Through.
%Hóstias tibi, Dómine, beáti Gregórii Mártyris tui atque Pontíficis dicátas méritis, benígnus assúme: et ad perpétuum nobis tríbue proveníre subsídium. Per Dóminum.


\postcommunion\label{ArmeniaPostcommunion}
\lett{G}{rant,} we beseech thee, O Lord our God; that we, having been refreshed by participation in the holy gifts, may experience the fruit of the intercession of blessed Gregory thy Bishop and Confessor, whose festival we now celebrate. Through.
%Refécti participatióne múneris sacri, quǽsumus, Dómine, Deus noster: ut, cujus exséquimur cultum, intercedénte beáto Gregório Mártyre tuo atque Pontífice, sentiámus efféctum. Per Dóminum nostrum.


\subby{St. Remigius}
\feastday{{St. Remigius}}
\fancyhead[RE,LO]{1 October}
\begin{inhead}
	{Memorial\\
		1 October}
\end{inhead}

\begin{rubric}
	The propers are from the First Common of a Confessor Bishop (p. \pageref{CommonConfessorBishopI}).
\end{rubric}


\bcpfeast{2 October. Holy Guardian Angels}{Guardian Angels}{2 October}
%\supplement{2 October}{Holy Guardian Angels}{}

\subbysub{I Evensong}\label{GuardianAngelsEvensong}

\gregorioscore{resources/gabc/ProperTime/GuardianAngelsEvensong.gabc}

    ℣. In the presence of the Angels I will sing praise unto thee, O my God.

	℟. I will worship toward thy holy temple, and praise thy Name.
	
\properantiphon{Mag.}{They are all {\dag} ministering spirits, sent forth to minister to them who shall be heirs of salvation.}

\subbysub{Mattins}

\begin{rubric}
	The Invitatory Hymn is as in I Evensong.
\end{rubric}

\officehymn\label{GuardianAngelsMattins}

\gregorioscore{resources/gabc/ProperTime/GuardianAngelsMattins.gabc}

    ℣. In the presence of the Angels I will sing praise unto thee, O my God.

	℟. I will worship toward thy holy temple, and praise thy Name.

\properantiphon{Ben.}{And the Angel {\dag} that talked with me came again, and waked me, as a man that is wakened out of his sleep.}

\subbysub{II Evensong}

\begin{rubric}
	The Office Hymn \& Versicle is as in I Evensong, with the following Antiphon.
\end{rubric}

\properantiphon{Mag.}{Ye holy Angels, {\dag} our Guardians, defend us in time of battle, lest we perish in the dreadful day of judgment.}


\subby{Sts. Placidus \& Companions}
\feastday{{Sts. Placidus \&c.}}
\fancyhead[RE,LO]{5 October}
\begin{inhead}
	{Double\\
		5 October}
\end{inhead}

\begin{rubric}
	The propers are from the Third Common of Many Martyrs (p. \pageref{CommonMartyrsIII}), with the Prayers from the Second Common of Many Martyrs (p. \pageref{CommonMartyrsII}).
\end{rubric}


\bcpfeast{7 October. Holy Rosary of the Blessed Virgin Mary}{Holy Rosary}{7 October}
%\supplement{7 October}{Holy Rosary}{of the Blessed Virgin Mary}

\subbysub{I Evensong}\label{HolyRosaryEvensong}

\gregorioscore{resources/gabc/ProperTime/HolyRosaryEvensong.gabc}

    ℣. Queen of the holy Rosary, pray for us.

	℟. That we may be made worthy of the promises of Christ.

\properantiphon{Mag.}{Blessed art thou, {\dag} O Virgin Mary, Mother of God, for thou hast believed the Lord: for those things are accomplished in thee which were told thee: intercede for us unto the Lord our God.}

\subbysub{Mattins}

\invitatoryhymn\label{HolyRosaryInvitatory}

\gregorioscore{resources/gabc/ProperTime/HolyRosaryInvitatory.gabc}

\officehymn\label{HolyRosaryMattins}

\gregorioscore{resources/gabc/ProperTime/HolyRosaryMattins.gabc}

    ℣. God hath chosen her and preferred her.

	℟. He hath made her to dwell in his tabernacle.

\properantiphon{Ben.}{To-day let us celebrate {\dag} devoutly the Solemnity of the holy Rosary of Mary, Mother of God, that she may intercede for us unto Jesus Christ the Lord.}

\subbysub{II Evensong}\label{HolyRosaryEvensongII}

\gregorioscore{resources/gabc/ProperTime/HolyRosaryEvensongII.gabc}

    ℣. Queen of the holy Rosary, pray for us.

	℟. That we may be made worthy of the promises of Christ.

\properantiphon{Mag.}{O blessed Mother {\dag} and spotless Virgin, thou glorious Queen of the world, may all acknowledge the help of thy prayer who celebrate the solemnity of thy holy Rosary.}


\subby{St. Mark of Rome}
\feastday{{St. Mark Rome}}
\fancyhead[RE,LO]{7 October}
\begin{inhead}
	{Memorial\\
		7 October}
\end{inhead}

\begin{rubric}
	The propers are from the Second Common of a Confessor Bishop (p. \pageref{CommonConfessorBishopII}), except for the following.
\end{rubric}

\collect
\lett{G}{raciously} hear our prayers, O Lord: and at the intercession of blessed Mark, thy Confessor and Bishop, mercifully grant us pardon and peace. Through.

\secret
\lett{G}{rant,} O Lord, that like as thy dedicated people do acknowledge that in tribulation they have been succoured by the merits of thy Saints: so this oblation, which they offer unto thee in honour of the same, may be acceptable in thy sight. Through.

\postcommunion
\lett{G}{rant,} we beseech thee, O Lord, that thy faithful people may ever rejoice in the veneration of thy Saints: and be defended by their perpetual supplication. Through.


\subby{Sts. Sergius, Bacchus, Marcellus, \& Apuleius}
\feastday{{Sts. Sergius \&c.}}
\fancyhead[RE,LO]{7 October}
\begin{inhead}
	{Memorial\\
		7 October}
\end{inhead}

\begin{rubric}
The propers are from the Second Common of Many Martyrs (p. \pageref{CommonMartyrsII}), except for the following.
\end{rubric}

\collect\label{SergiusCollect}
\lett{M}{ay} the blessed merits of thy holy Martyrs Sergius, Bacchus, Marcellus, and Apuleius uphold us, O Lord: and ever make us fervent in thy love. Through.

\secret\label{SergiusSecret}
\lett{W}{e} beseech thee, O Lord, that, through the meritorious supplication of thy Saints: this sacrifice, which we offer, may obtain for us the favour of thy majesty. Through.

\postcommunion\label{SergiusPostcommunion}
\lett{G}{rant,} O Lord, that, being strengthened by the sacraments which we have received: we may, at the intercession of thy holy Martyrs, Sergius, Bacchus, Marcellus, and Apuleius, fight against all iniquity, and be defended by thy heavenly armour. Through.


\subby{Sts. Denys, Rusticus, \& Eleutherius}
\feastday{{Sts. Denys \&c.}}
\fancyhead[RE,LO]{9 October}
\begin{inhead}
	{Memorial\\
		9 October}
\end{inhead}

\introit
\lett{L}{et} the people tell of the wisdom of the Saints, and let the Church shew forth their praise: their names shall live for evermore. \textit{Ps.} Rejoice in the Lord, O ye righteous: for it becometh well the just to be thankful.

\collect
\lett{O}{God,} who on this day didst strengthen blessed Denys, thy Martyr and Bishop, with the virtue of constancy in his passion, and didst vouchsafe to join unto him Rusticus and Eleutherius for the preaching of thy glory to the Gentiles: grant us, we beseech thee; by their example, to despise for love of thee the prosperity of the world, and to fear none of its adversities. Through.

\readingcitation{Epistle}{Acts 17:22}
\lett{I}{n those days:} Paul stood in the midst of Mars' hill, and said: Ye men of Athens, I perceive that in all things ye are too superstitious. For as I passed by, and beheld your devotions, I found an altar with this inscription, \textsc{To the Unknown God}. Whom therefore ye ignorantly worship, him declare I unto you. God that made the world and all things therein, seeing that he is Lord of heaven and earth, dwelleth not in temples made with hands; Neither is worshipped with men's hands, as though he needed any thing, seeing he giveth to all life, and breath, and all things; And hath made of one blood all nations of men for to dwell on all the face of the earth, and hath determined the times before appointed, and the bounds of their habitation; That they should seek the Lord, if haply they might feel after him, and find him, though he be not far from every one of us: For in him we live, and move, and have our being; as certain also of your own poets have said, For we are also his offspring. Forasmuch then as we are the offspring of God, we ought not to think that the Godhead is like unto gold, or silver, or stone, graven by art and man's device. And the times of this ignorance God winked at; but now commandeth all men every where to repent: Because he hath appointed a day, in the which he will judge the world in righteousness by that man whom he hath ordained; whereof he hath given assurance unto all men, in that he hath raised him from the dead. And when they heard of the resurrection of the dead, some mocked: and others said, We will hear thee again of this matter. So Paul departed from among them. Howbeit certain men clave unto him, and believed: among the which was Dionysius the Areopagite, and a woman named Damaris, and others with them.

\gradall{Our soul is escaped even as a bird out of the snare of the fowler. ℣. The snare is broken, and we are delivered: our help is in the name of the Lord, who hath made heaven and earth.}{Alleluia, alleluia. ℣. Let the righteous be glad and rejoice before God: let them also be merry and joyful. Alleluia.}

\readingcitation{Gospel}{Luke 12:1}
\lett{A}{t that time:} Jesus said unto his disciples: Beware ye of the leaven of the Pharisees, which is hypocrisy. For there is nothing covered, that shall not be revealed; neither hid, that shall not be known. Therefore whatsoever ye have spoken in darkness shall be heard in the light; and that which ye have spoken in the ear in closets shall be proclaimed upon the housetops. And I say unto you my friends, Be not afraid of them that kill the body, and after that have no more that they can do. But I will forewarn you whom ye shall fear: Fear him, which after he hath killed hath power to cast into hell; yea, I say unto you, Fear him. Are not five sparrows sold for two farthings, and not one of them is forgotten before God? But even the very hairs of your head are all numbered. Fear not therefore: ye are of more value than many sparrows. Also I say unto you, Whosoever shall confess me before men, him shall the Son of man also confess before the angels of God.

\offertory{Let the Saints be joyful with glory, let them rejoice in their beds: let the praises of God be in their mouth, alleluia.}

\secret
\lett{G}{raciously} receive, we beseech thee, O Lord, the oblations of thy people, for the honour of thy Saints: and sanctify us by their intercession. Through.

\communion{And I say unto you, my friends: Be not afraid of them that persecute you.}

\postcommunion
\lett{W}{e} beseech thee, O Lord that we, who have received thy sacraments: may at the intercession of thy blessed Martyrs Denys, Rusticus, and Eleutherius, be profited thereby unto the increase of eternal redemption. Through.


\bcpfeast{11 October. Motherhood of the Blessed Virgin Mary}{Motherhood}{11 October}

\subbysub{I Evensong}

\begin{rubric}
	The Office Hymn is from the Common of the Blessed Virgin Mary (p. \pageref{CommonBVM}), with the following Versicle \& Antiphon.
\end{rubric}

    ℣. Blessed art thou amongst women.

	℟. And blessed is the fruit of thy womb.
	
	%Original translation from the Breviarium Monasticum:
	\antiphon{Mag.}{Let us celebrate with joy {\dag} the Motherhood of the blessed Mary ever Virgin.}
%Cum jucunditáte Maternitátem beáta Maria semper Virginis celebrémus.

\subbysub{Mattins}

\invitatoryhymn\label{MotherhoodInvitatory}

\gregorioscore{resources/gabc/ProperTime/MotherhoodInvitatory.gabc}

\officehymn\label{MotherhoodMattins}

\gregorioscore{resources/gabc/ProperTime/MotherhoodMattins.gabc}

    ℣. The root of Jesse hath budded forth: a star hath arisen out of Jacob.

	℟. The Virgin hath brought forth the Saviour: we praise thee, O our God.
	
	%Original translation of the last clause.
	\antiphon{Ben.}{O holy Mary, {\dag} help thou the suffering, strengthen the faint-hearted, comfort the sorrowful; pray for the people, entreat for the clergy, intercede for all womankind vowed unto God: may all experience thy help, who celebrate thine admirable motherhood.}
	
\subbysub{II Evensong}

\begin{rubric}
	The Office Hymn is from the Common of the Blessed Virgin Mary (p. \pageref{CommonBVM}), with the following Versicle \& Antiphon.
\end{rubric}

    ℣. Blessed art thou amongst women.

	℟. And blessed is the fruit of thy womb.
	
	%Original translation from the Breviarium Monasticum:
	\antiphon{Mag.}{Thy motherhood, O Virgin Birthgiver of God, announces joy unto the whole world: for from thee the Sun of Righteousness hath arisen, Christ our God.}


\subby{St. Wilfred}
\feastday{{St. Wilfred}}
\fancyhead[RE,LO]{12 October}
\begin{inhead}
	{Memorial\\
		12 October}
\end{inhead}

\begin{rubric}
	The propers are from the Second Common of a Confessor Bishop (p. \pageref{CommonConfessorBishopII}), except for the following.
\end{rubric}

\collect
\lett{O}{God,} by whose grace the blessed Bishop Wilfrid did wondrously shine forth with the glorious tokens of his merits: mercifully grant unto us; that we, who by his doctrine are taught to seek after things heavenly, may ever be defended by his advocacy. Through.

\secret
\lett{P}{urify,} we beseech thee, almighty God, the hearts of thy family by the enlightening of thy Holy Spirit: that through the intercession of blessed Wilfrid, thy Confessor and Bishop, these gifts of devotion may be rendered acceptable unto thee. Through . . . in the unity of the same.

\postcommunion
\lett{O}{Lord,} who hast fulfilled us with the food of eternal redemption, we humbly entreat thy mercy: that by the intercession of blessed Wilfrid, thy Confessor and Bishop, we may receive the gifts of everlasting salvation. Through.

%MANUAL ADJUSTMENT:
\clearpage
\subby{St. Edward}
\feastday{{St. Edward}}
\fancyhead[RE,LO]{13 October}
\begin{inhead}
	{Memorial\\
		13 October}
\end{inhead}

\begin{rubric}
	The propers are from the First Common of a Confessor not a Bishop (p. \pageref{CommonConfessorNotBishopI}), except for the following.
\end{rubric}

\collect
\lett{O}{God,} who didst bestow upon thy blessed Confessor King Edward the crown of everlasting glory: grant us, we beseech thee; so to venerate him on earth, that we may be enabled to reign with him in heaven. Through.


\subby{St. Callistus}
\feastday{{St. Callistus}}
\fancyhead[RE,LO]{14 October}
\begin{inhead}
	{Double\\
		14 October}
\end{inhead}

\introit
\lett{O}{ye} priests of the Lord, bless ye the Lord: O ye holy and humble men of heart, bless ye the Lord. \textit{Cant.} O all ye works of the Lord, bless ye the Lord: praise him, and magnify him for ever.

\collect
\lett{O}{God,} who seest that we fail by reason of our infirmity: through the examples of thy Saints mercifully restore us to the love of thee. Through.

\readingcitation{Epistle}{Hebrews 5:1}
\lett{B}{rethren:} Every high priest taken from among men is ordained for men in things pertaining to God, that he may offer both gifts and sacrifices for sins: Who can have compassion on the ignorant, and on them that are out of the way; for that he himself also is compassed with infirmity. And by reason hereof he ought, as for the people, so also for himself, to offer for sins. And no man taketh this honour unto himself, but he that is called of God, as was Aaron.

\gradall{I have found David my servant, with my holy oil have I anointed him: my hand shall hold him fast, and my arm shall strengthen him. ℣. The enemy shall not be able to do him violence: the son of wickedness shall not hurt him.}{Alleluia, alleluia. ℣. The Lord loved him, and adorned him: he clothed him with a robe of glory. Alleluia.}

\begin{rubric}
	The Gospel is of the Second Common of a Martyr not a Bishop (p. \pageref{CommonMartyrNotBishopII}).
\end{rubric}

\offertory{My truth and my mercy shall be with him: and in my name shall his horn be exalted.}

\secret
\lett{O}{Lord,} let this mystical oblation be profitable unto us: that it may both deliver us from our offences, and stablish us in everlasting salvation. Through.

\communion{Blessed is the servant whom the lord, when he cometh, shall find watching: verily I say unto you, that he shall make him ruler over all his goods.}

\postcommunion
\lett{W}{e} beseech thee, almighty God: that the gifts which we have hallowed may cleanse us from our iniquities, and bring forth in us the fruit of godly living. Through.


\bcpfeast{15 October. Our Lady of Walsingham}{Walsingham}{15 October}
%\supplement{15 October}{Our Lady of Walsingham}{}

%\gregorioscore{resources/gabc/ProperTime/WalsinghamEvensong.gabc}
	
\begin{paracol}{2}[]
\sloppy
\begin{inhead}
	I Evensong\label{WalsinghamEvensong}
\end{inhead}
\begin{hangparas}{1.25em}{1}
Hail Mary, ever blessed,

Of Walsingham the Queen!

Through vision of Richeldis,

Thy favours there were seen.

When England was thy dowry,

There pilgrims bowed the knee.

At morn and noon and even,

They knelt to honour thee.\\

Hail Mary, ever blessed.

Thy children still delight

To tell abroad thy praises,

Thy miracles, thy might.

Still pilgrim feet are treading

Along the holy way.

Hostess of England's Nazareth,

Receive us home today!\\

Hail Mary, ever blessed.

The wells of water pure

Which mark thy holy places

Are signs that God doth cure

For sick of soul and body.

E’er since Richeldis' day,

They spring in benediction

Beside the Pilgrims' Way.\\

Hail Mary, ever blessed.

Thy name is great indeed;

For Jesus Christ our Saviour

Was in thy womb conceived.

Thy name be ever praised,

Increasing in this place,

And loud the angel's greeting:

`Hail Mary, full of grace!'\\

All glory, laud, and honour be

O Jesu, Virgin-born

Our worship is alone to thee

With Father and the Ghost.

Thou art our Mediator true,

Appealing for our sin;

Send unto us thine own Spirit,

Who makes us thy true kin. Amen.
\end{hangparas}

    ℣. O Mary, Queen of Heaven and of England, praise ye the Lord.

	℟. And bring us unto the majesty of thy Son.
	
	\properantiphon{Mag.}{O all ye Saints, {\dag} come unto the Holy Place of Walsingham, and bless ye the Lord.}

\switchcolumn

\begin{inhead}
	Mattins\label{WalsinghamMattins}
\end{inhead}

\begin{rubric}
		\textsc{Note,} The Invitatory Hymn is the \emph{Salve Regina} (p. \pageref{SalveReginaEnglish}).
	\end{rubric}

\begin{hangparas}{1.25em}{1}
When Christ was born in Bethlehem,\par
The Angels worshipped and adored;\par
And kings and shepherds bowed the knee\par
To him their Saviour and their Lord.\\

For thirty years in Nazareth,\par
His Mother worshipped at the shrine\par
Of him who came in human form\par
To save the world, your Lord and mine.\\

For thirty years in Nazareth\par
He lived a hidden life of prayer;\par
And then went forth to work and die,\par
And every human sorrow share.\\

A thousand years have passed away;\par
Another Nazareth must rise\par
In England’s fair and lovely land,\par
Beneath those distant northern skies.\\

So Mary, God's own Mother blest,\par
Seeks out Richeldis, lady fair;\par
And bids her build at Walsingham\par
A second Nazareth of prayer.\\

Richeldis hastens to obey;\par
And soon the builder’s work is done:\par
The Shrine of England’s Nazareth\par
Proclaims the Mother and the Son.\\

But times did come when wicked hands\par
Were laid upon that holy place;\par
And men who knew not Mary's love\par
Did Walsingham's fair Shrine deface.\\

Four hundred years again have passed;\par
Once more the Mother of our Lord\par
Is shrined and honoured where her Son\par
Upon his Altar is adored.\\

So when you come to Walsingham,\par
Remember this, and mark it well:\par
It is to Nazareth you come,\par
Where Jesus, Mary, Joseph dwell.\\

All glory, laud, and honour be\par
To Jesu, Virgin-born, to thee;\par
All glory, as is ever meet,\par
To Father and to Paraclete. Amen.
\end{hangparas}

℣. Blessed is the holy Virgin Mary, and most worthy of all praise.

℟. Through her hath arisen the Sun of Justice, Christ our God, by whom we are saved and redeemed.

\properantiphon{Ben.}{Wicked men in their lust defied the Church, {\dag} and razed the Shrine of Our Lady.}
\end{paracol}

\begin{rubric}
	In II Evensong, the Office Hymn \& Versicle are of I Evensong, with the following Antiphon.
\end{rubric}

\properantiphon{Mag.}{O God, raise thou up again devotion to thy daughter: {\dag} Our blessed Lady of Walsingham.}


\bcpfeast{18 October. St. Luke}{St. Luke}{18 October}
%\supplement{18 October}{St. Luke}{}

\begin{secrubric}
	The Daily Office propers are of the Common of Apostles (p. \pageref{CommonApostles}).
\end{secrubric}


\subby{St. Frideswide}
\feastday{{St. Frideswide}}
\fancyhead[RE,LO]{19 October}
\begin{inhead}
	{Memorial\\
		19 October}
\end{inhead}

\begin{rubric}
	The propers are from the First Common of a Virgin (p. \pageref{CommonVirginOnlyI}), except for the following.
\end{rubric}

\collect
\lett{A}{lmighty} and everlasting God, the author of virtue and lover of virginity: grant us, we beseech thee; that, like as thy Virgin Frideswide was pleasing unto thee by the merit of chastity of life; so through her merits we may find favour in thy sight. Through.

\secret
\lett{W}{e} offer thee, O Lord, our prayers and gifts with gladness in honour of Saint Frideswide thy Virgin: that we may be enabled worthily to perform the same, and to obtain thine everlasting healing. Through.

\postcommunion
\lett{W}{e} beseech thee, O Lord, that the mysteries which we have received may be profitable unto many: through the intercession of thy blessed Virgin Frideswide, may they both deliver us from our sins, and obtain for us thy gracious protection. Through.


\subby{St. Hilarion}
\feastday{{St. Hilarion}}
\fancyhead[RE,LO]{21 October}
\begin{inhead}
	{Memorial\\
		21 October}
\end{inhead}

\begin{rubric}
	The propers are from the Common of Abbots (p. \pageref{CommonAbbots}).\par
	\textsc{Note,} Commemoration is made of Sts. Ursula \& Companions.
\end{rubric}


\subby{Sts. Ursula \& Companions}
\feastday{{Sts. Ursula \&c.}}
\fancyhead[RE,LO]{21 October}
\begin{inhead}
	{Memorial\\
		21 October}
\end{inhead}

\begin{rubric}
	The propers are from the First Common of a Virgin Martyr (p. \pageref{CommonVirginMartyrI}), with the Prayers from the Common of Many Virgin Martyrs (p. \pageref{CommonVirginMartyrs}).
\end{rubric}


\bcpfeast{24 October. St. Raphael}{St. Raphael}{24 October}
%\supplement{24 October}{St. Raphael}{}

\begin{rubric}
	The Hymns \& Versicles are of Michaelmas Day (p. \pageref{SaintMichael}), with the following Antiphons.
\end{rubric}

\properantiphon{Mag. \& Ben.}{I am Raphael {\dag} the Angel, that stand before the Holy One: wherefore bless ye God for ever, and confess ye all his great and wonderful doings.}

\properantiphon{Mag.}{O Prince most glorious, {\dag} Raphael the Archangel, keep us in remembrance: here and everywhere, always, entreat the Son of God for us.}


\subby{Sts. Chrysanthus \& Daria}
\feastday{{Sts. Chrysanthus \& Daria}}
\fancyhead[RE,LO]{25 October}
\begin{inhead}
	{Memorial\\
		25 October}
\end{inhead}

\begin{rubric}
	The propers are from the First Common of Many Martyrs (p. \pageref{CommonMartyrsI}), except for the following.
\end{rubric}

\collect
\lett{W}{e} beseech thee, O Lord, that the prayers of thy blessed Martyrs Chrysanthus and Daria may assist us: that, as we venerate them with our homage, so we may ever feel their loving help. Through.

\begin{rubric}
	The Epistle is the fourth additional Epistle of the Third Common of Many Martyrs (p. \pageref{2Corinthians64}).
\end{rubric}

\begin{rubric}
	The Gospel is the fourth additional Gospel of the Third Common of Many Martyrs (p. \pageref{Luke1147}).
\end{rubric}

\secret
\lett{W}{e} beseech thee, O Lord, that the sacrifice of thy people, which they solemnly offer on the birthday of thy holy Martyrs Chrysanthus and Daria, may be pleasing unto thee. Through.

\postcommunion
\lett{O}{Lord,} who hast fulfilled us with mystic gifts and joys: grant, we beseech thee; that by the intercession of thy holy Martyrs Chrysanthus and Daria, we may spiritually attain to those things which we temporally perform. Through.

%MANUAL ADJUSTMENT:
\clearpage
\subby{St. Evaristus}
\feastday{{St. Evaristus}}
\fancyhead[RE,LO]{26 October}
\begin{inhead}
	{Memorial\\
		26 October}
\end{inhead}

\begin{rubric}
	The propers are from the First Common of a Martyr Bishop (p. \pageref{CommonMartyrBishopI}).
\end{rubric}

\begin{rubric}
	If to-day be Saturday, the anticipated Vigil of Sts. Simon and Jude is kept, as is noted on the following day: but the \nth{2} Collect is of St. Evaristus, and the \nth{3} of St. Mary in Eastertide (p. \SPMaryEaster).
\end{rubric}

\subby{Vigil of Sts. Simon \& Jude}
\feastday{{Sts. Simon \& Jude Vigil}}
\fancyhead[RE,LO]{27 October}
\begin{inhead}
	{Vigil\\
		27 October}
\end{inhead}

\introit
\lett{L}{et} the sorrowful sighing of the prisoners, O Lord, come before thee: reward thou our neighbours sevenfold into their bosom: avenge thou the blood of thy saints that is shed. \textit{Ps.} O God, the heathen are come into thine inheritance: thy holy temple have they defiled: and made Jerusalem an heap of stones.

\begin{rubric}
	\emph{Gloria in excelsis} is not said.
\end{rubric}

\collect
\lett{G}{rant,} we beseech thee, almighty God: that as we prevent the glorious birthday of thine Apostles Simon and Jude; so they may prevent us in the sight of thy majesty, for the obtaining of thy blessings. Through.

\begin{rubric}
	\nth{2} Collect of St. Mary in Eastertide (p. \SPMaryEaster). \nth{3} against the Persecutors of the Church (\SPAgainst) or for the Chief Bishop (p. \SPChiefBishop).
\end{rubric}

\readingcitation{Epistle}{1 Corinthians 4:9}
\lett{B}{rethren:} We are made a spectacle unto the world, and to Angels, and to men. We are fools for Christ's sake, but ye are wise in Christ; we are weak, but ye are strong; ye are honourable, but we are despised. Even unto this present hour we both hunger, and thirst, and are naked, and are buffeted, and have no certain dwellingplace; And labour, working with our own hands: being reviled, we bless; being persecuted, we suffer it: Being defamed, we intreat: we are made as the filth of the world, and are the offscouring of all things unto this day. I write not these things to shame you, but as my beloved sons I warn you.

\gradual{Avenge thou, O Lord, the blood of thy Saints that is shed. ℣. The dead bodies of thy servants, O Lord, have they given to be meat unto the fowls of the air: and the flesh of thy Saints unto the beasts of the land.}

\readingcitation{Gospel}{John 15:1}
\lett{A}{t that time:} Jesus said unto his disciples: I am the true vine, and my Father is the husbandman. Every branch in me that beareth not fruit he taketh away: and every branch that beareth fruit, he purgeth it, that it may bring forth more fruit. Now ye are clean through the word which I have spoken unto you. Abide in me, and I in you. As the branch cannot bear fruit of itself, except it abide in the vine; no more can ye, except ye abide in me. I am the vine, ye are the branches: He that abideth in me, and I in him, the same bringeth forth much fruit: for without me ye can do nothing. If a man abide not in me, he is cast forth as a branch, and is withered; and men gather them, and cast them into the fire, and they are burned. If ye abide in me, and my words abide in you, ye shall ask what ye will, and it shall be done unto you.

\offertory{Let the Saints be joyful with glory: let them rejoice in their beds: let the praises of God be in their mouth.}

\secret
\lett{W}{e} humbly beseech thee, O Lord: that, although our consciences be burdened by reason of our sins: the gifts, wherewith we do prevent the festival of thy holy Apostles Simon and Jude, may through their merits be rendered acceptable in thy sight. Through.

\begin{rubric}
	\nth{2} Secret of St. Mary in Eastertide (p. \SPMaryEaster). \nth{3} against the Persecutors of the Church (\SPAgainst) or for the Chief Bishop (p. \SPChiefBishop).
\end{rubric}

\communion{The dead bodies of thy servants, O Lord, have they given to be meat unto the fowls of the air, and the flesh of thy Saints unto the beasts of the land: according to the greatness of thy power, preserve thou those that are appointed to die.}

\postcommunion
\lett{W}{e} who have received thy sacrament, humbly entreat thee, O Lord: that, by the intercession of thy blessed Apostles Simon and Jude, those things which we temporally perform we may receive unto life eternal. Through.

\begin{rubric}
	\nth{2} Postcommunion of St. Mary in Eastertide (p. \SPMaryEaster). \nth{3} against the Persecutors of the Church (\SPAgainst) or for the Chief Bishop (p. \SPChiefBishop).
\end{rubric}


\bcpfeast{28 October. Sts. Simon and Jude}{Sts. Simon \& Jude}{28 October}
%\supplement{28 October}{Sts. Simon and Jude}{}

\begin{secrubric}
	The Daily Office propers are of the Common of Apostles (p. \pageref{CommonApostles}).
\end{secrubric}


\bcpfeast{Christ the King Sunday}{Christ the King}{}
%\supplement{}{Christ the King}{Sunday}

\subbysub{I Evensong}\label{ChristTheKingEvensong}

\gregorioscore{resources/gabc/ProperTime/ChristTheKingEvensong.gabc}

℣. All power is given unto me.

℟. In heaven and in earth.

\properantiphon{Mag.}{The Lord God shall give unto him {\dag} the throne of his father David: and he shall reign over the house of Jacob for ever; and of his kingdom there shall be no end, alleluia.}

\subbysub{Mattins}

\invitatoryhymn\label{ChristTheKingInvitatory}

\gregorioscore{resources/gabc/ProperTime/ChristTheKingInvitatory.gabc}

\officehymn\label{ChristTheKingMattins}

\gregorioscore{resources/gabc/ProperTime/ChristTheKingMattins.gabc}

℣. His kingdom shall be increased.

℟. And of his peace there shall be no end.

\properantiphon{Ben.}{Unto God and his Father {\dag} hath he made us to be a kingdom, who is the first begotten of the dead, and the Prince of the kings of the earth, alleluia.}

\subbysub{II Evensong}

\begin{rubric}
	The Office Hymn is as in I Evensong, with the following Versicle \& Antiphon.
\end{rubric}

℣. His kingdom shall be increased.

℟. And of his peace there shall be no end.

\properantiphon{Mag.}{He hath on his vesture {\dag} and on his thigh a name written, King of kings, and Lord of lords: to him be glory and dominion for ever and ever.}


\subby{Vigil of All Hallows}
\feastday{{Hallowe'en}}
\fancyhead[RE,LO]{31 October}
\begin{inhead}
	{Vigil\\
		31 October}
\end{inhead}

\introit
\lett{T}{he} Saints and judges of the nations, and have dominion over the people: and the Lord their God shall reign for ever. \textit{Ps.} Rejoice in the Lord, O ye righteous: for it becometh well the just to be thankful.

\collect
\lett{O}{Lord,} our God, multiply upon us thy grace: and grant, that by our holy profession we may follow after the gladness of them, whose glorious festival we prevent. Through.

\begin{rubric}
	\nth{2} Collect of the Holy Ghost (p. \SPHolyGhost). \nth{3} Collect against the Persecutors of the Church (p. \SPAgainst) or for the Chief Bishop (p. \SPChiefBishop).
\end{rubric}

%RV:
\readingcitation{Epistle}{Revelation 5:6}
\lett{I}{n those days:} Lo, I, John, saw in the midst of the throne and of the four living creatures, and in the midst of the elders, a Lamb standing, as though it had been slain, having seven horns, and seven eyes, which are the seven Spirits of God, sent forth into all the earth. And he came, and he taketh it out of the right hand of him that sat on the throne. And when he had taken the book, the four living creatures and the four and twenty elders fell down before the Lamb, having each one a harp, and golden bowls full of incense, which are the prayers of the saints. And they sing a new song, saying, Worthy art thou to take the book, and to open the seals thereof: for thou wast slain, and didst purchase unto God with thy blood men of every tribe, and tongue, and people, and nation,  and madest them to be unto our God a kingdom and priests; and they reign upon the earth. And I saw, and I heard a voice of many angels round about the throne and the living creatures and the elders; and the number of them was ten thousand times ten thousand, and thousands of thousands; saying with a great voice, Worthy is the Lamb that hath been slain to receive the power, and riches, and wisdom, and might, and honour, and glory, and blessing. 

\gradual{Let the Saints be joyful with glory: let them rejoice in their beds. ℣. O sing unto the Lord a new song: let the congregation of saints praise him.}

\readingcitation{Gospel}{Luke 6:17}
\lett{A}{t that time:} Jesus came down from the mountain, and stood in the plain, and the company of his disciples, and a great multitude of people out of all Judaea and Jerusalem, and from the sea coast of Tyre and Sidon, which came to hear him, and to be healed of their diseases; And they that were vexed with unclean spirits: and they were healed. And the whole multitude sought to touch him: for there went virtue out of him, and healed them all. And he lifted up his eyes on his disciples, and said, Blessed be ye poor: for yours is the kingdom of God. Blessed are ye that hunger now: for ye shall be filled. Blessed are ye that weep now: for ye shall laugh. Blessed are ye, when men shall hate you, and when they shall separate you from their company, and shall reproach you, and cast out your name as evil, for the Son of man's sake. Rejoice ye in that day, and leap for joy: for, behold, your reward is great in heaven.

\offertory{Let the Saints be joyful with glory, let them rejoice in their beds: let the praises of God be in their mouth.}

\secret
\lett{W}{e} set upon thine altars, O Lord, the gifts which we offer: grant, we beseech thee; that they may be profitable for our salvation through the supplication of all thy Saints, whose coming festival we prevent. Through.

\begin{rubric}
	\nth{2} Secret of the Holy Ghost (p. \SPHolyGhost). \nth{3} Collect against the Persecutors of the Church (p. \SPAgainst) or for the Chief Bishop (p. \SPChiefBishop).
\end{rubric}

\communion{The souls of the righteous are in the hand of God, and there shall no torment of malice touch them: in the sight of the unwise they seemed to die, but they are in peace.}

\postcommunion
\lett{W}{e} beseech thee, O Lord: that, having accomplished with joy the mysteries of the coming feast; we may be aided by the prayers of them in whose memory they are shewn forth. Through.

\begin{rubric}
	\nth{2} Postcommunion of the Holy Ghost (p. \SPHolyGhost). \nth{3} Collect against the Persecutors of the Church (p. \SPAgainst) or for the Chief Bishop (p. \SPChiefBishop).
\end{rubric}


\bcpfeast{1 November. All Hallows}{All Hallows}{1 November}
%\supplement{1 November}{All Hallows}{}

\subbysub{I Evensong}\label{AllHallowsEvensong}

\gregorioscore{resources/gabc/ProperTime/AllHallowsEvensong.gabc}

℣. Be glad, O ye righteous, and rejoice in the Lord.

℟. And be joyful, all ye that are true of heart.

\properantiphon{Mag.}{O ye Angels and Archangels, {\dag} O ye Thrones and Dominions, Principalities and Powers, Virtues of the heavens, Cherubim and Seraphim, O ye Patriarchs and Prophets, ye holy Doctors of the law, O all ye Apostles, ye Martyrs of Christ, ye holy Confessors, ye Virgins of the Lord, ye holy Hermits, and all Saints: offer for us your intercessions.}

\subbysub{Mattins}

\begin{rubric}
	The Invitatory Hymn is as in I Evensong.
\end{rubric}

\officehymn\label{AllHallowsMattins}

\gregorioscore{resources/gabc/ProperTime/AllHallowsMattins.gabc}

℣. Let the Saints be joyful with glory.

℟. Let them rejoice in their beds.

\properantiphon{Ben.}{The glorious company {\dag} of the Apostles praise thee; the goodly fellowship of the Prophets praise thee; the white-robed army of Martyrs praise thee; with one heart and voice do all the elect acknowledge thee: O blessed Trinity, one only God.}

\subbysub{II Evensong}

\begin{rubric}
	The Office Hymn is as in I Evensong, with the following Versicle \& Antiphon.
\end{rubric}

℣. Let the Saints be joyful with glory.

℟. Let them rejoice in their beds.

\properantiphon{Mag.}{O how glorious {\dag} is the kingdom wherein all the Saints rejoice with Christ; arrayed in white robes, they follow the Lamb whithersoever he goeth.}

\begin{rubric}
	The Office is of All Hallows through its Octave, unless there occur a Sunday or Feast Day (Double or higher), any concurring Memorial being commemorated.
\end{rubric}


\bcpfeast{2 November. All Souls}{All Souls}{2 November}

\begin{rubric}
	The Office of the Day is the Office of the Dead.
\end{rubric}


\subby{St. Winifred}
\feastday{{St. Winifred}}
\fancyhead[RE,LO]{3 November}
\begin{inhead}
	{Memorial\\
		3 November}
\end{inhead}

\begin{rubric}
	The propers are from the First Common of a Virgin Martyr (p. \pageref{CommonVirginMartyrI}).
\end{rubric}


\subby{Sts. Vitalis \& Agricola}
\feastday{{Sts. Vitalis \& Agricola}}
\fancyhead[RE,LO]{4 November}
\begin{inhead}
	{Memorial\\
		4 November}
\end{inhead}

\begin{rubric}
	The propers are from the Second Common of Many Martyrs (p. \pageref{CommonMartyrsII}), except for the following.
\end{rubric}

\collect
\lett{G}{rant,} we beseech thee, almighty God: that we, who devoutly celebrate the festival of thy holy Martyrs Vitalis and Agricola, may be aided by their intercession with thee. Through.

\begin{rubric}
	The Epistle is from the Third Common of Many Martyrs (p. \pageref{CommonMartyrsIII}).
\end{rubric}

\begin{rubric}
	The Gospel is from the Second Common of a Martyr Bishop (p. \pageref{CommonMartyrBishopII}).
\end{rubric}

\secret
\lett{W}{e} beseech thee, O Lord, mercifully to accept these our oblations: that, at the intercession of thy holy Martyrs Vitalis and Agricola, we may be defended against all adversities. Through.

\postcommunion
\lett{M}{ay} this communion, O Lord, cleanse us from guilt: and at the intercession of thy holy Martyrs, Vitalis and Agricola, make us partakers of thy heavenly healing. Through.


\subby{St. Elizabeth}
\feastday{{St. Elizabeth}}
\fancyhead[RE,LO]{5 November}
\begin{inhead}
	{Memorial\\
		5 November}
\end{inhead}

\begin{rubric}
	The propers are from the Common of Neither Virgin nor Martyr (p. \pageref{CommonNeitherVirginMartyr}).
\end{rubric}


\subby{St. Willibrord}
\feastday{{St. Willibrord}}
\fancyhead[RE,LO]{7 November}
\begin{inhead}
	{Memorial\\
		7 November}
\end{inhead}

\begin{rubric}
	The propers are from the First Common of a Confessor Bishop (p. \pageref{CommonConfessorBishopI}).
\end{rubric}


\subby{Octave Day of All Hallows}
\feastday{{All Hallows Octave Day}}
\fancyhead[RE,LO]{8 November}
\begin{inhead}
	{Greater Double\\
		8 November}
\end{inhead}

\begin{rubric}
	The propers as on the Feast, with Commemoration of the Four Crowned Martyrs.
\end{rubric}


\subby{Four Crowned Martyrs}
\feastday{{Four Crowned Martyrs}}
\fancyhead[RE,LO]{8 November}
\begin{inhead}
	{Memorial\\
		8 November}
\end{inhead}

\begin{rubric}
	The propers are from the First Common of Many Martyrs (p. \pageref{CommonMartyrsI}), except for the following.
\end{rubric}

\begin{rubric}
	\textsc{Note,} The Epistle and Gospel are of All Hallows' Day.
\end{rubric}

\collect
\lett{G}{rant,} we beseech thee, almighty God: that, like as we have known thy glorious Martyrs to be constant in their confession, so we may perceive their loving intercession for us with thee. Through.

\secret
\lett{L}{et} thy plenteous benediction descend, O Lord: and both render our gifts acceptable unto thee at the intercession of thy holy Martyrs, and make them to us a sacrament of redemption. Through.

\postcommunion
\lett{O}{Lord,} who hast refreshed us with the gladness of thy heavenly sacraments: we humbly entreat thee; that we may be protected by the succour of them in whose triumphs we glory. Through.


\subby{Dedication of the Basilica of St. Saviour}
\feastday{{St. Saviour Basilica}}
\fancyhead[RE,LO]{9 November}
\begin{inhead}
	{Greater Double\\
		9 November}
\end{inhead}

\begin{rubric}
	The propers are from the Common of a Dedication of a Church (p. \pageref{CommonDedication}).
\end{rubric}

\begin{rubric}
	\textsc{Note,} The Creed is said, and Commemoration is made of St. Theodore Tyro.
\end{rubric}


\subby{St. Theodore Tyro}
\feastday{{St. Theodore Tyro}}
\fancyhead[RE,LO]{9 November}
\begin{inhead}
	{Memorial\\
		9 November}
\end{inhead}

\begin{rubric}
	The propers are from the Second Common of a Martyr not a Bishop (p. \pageref{CommonMartyrNotBishopII}), except for the following.
\end{rubric}

\collect
\lett{O}{God,} who dost encompass and protect us with the glorious confession of blessed Theodore, thy Martyr: grant us both to profit by his example, and to be supported by his intercession. Through.

\secret
\lett{A}{ccept,} O Lord, the prayers of thy faithful people, and the oblations of their sacrifices: and at the intercession of blessed Theodore, thy Martyr, may we through these offices of godly devotion enter into heavenly glory. Through.

\postcommunion
\lett{G}{rant} to us, we beseech thee, O Lord: at the intercession of blessed Theodore thy Martyr; that those things which we touch with our lips we may receive in purity of heart. Through.


\subby{Sts. Tryphon, Respicius, \& Nympha}
\feastday{{Sts. Tryphon \&c.}}
\fancyhead[RE,LO]{10 November}
\begin{inhead}
	{Memorial\\
		10 November}
\end{inhead}

\introit
\lett{T}{he} righteous cry, and the Lord heareth them: and delivereth them out of all their troubles. \textit{Ps.} I will alway give thanks unto the Lord: his praise shall ever be in my mouth.

\collect
\lett{G}{rant,} we beseech thee, O Lord, that we, ever keeping the feast of thy holy Martyrs Tryphon, Respicius, and Nympha: may through their intercession enjoy the benefit of thy protection. Through.

\begin{rubric}
	The Epistle is the third additional Epistle from the Third Common of Many Martyrs (p. \pageref{Romans818}).
\end{rubric}

\gradall{Avenge, O Lord, the blood of thy Saints that is shed. ℣. The dead bodies of thy servants have they given to be meat unto the fowls of the air: the flesh of thy Saints unto the beasts of the land.}{Alleluia, alleluia. ℣. Right dear in the sight of the Lord is the death of his Saints. Alleluia.}

\begin{rubric}
	The Gospel is from the Third Common of Many Martyrs (p. \pageref{CommonMartyrsIII}).
\end{rubric}

\offertory{Be glad, O ye righteous, and rejoice in the Lord: and be joyful, all ye that are true of heart.}

\secret
\lett{W}{e} beseech thee, O Lord, that the gifts which we offer unto thee of our bounden duty and service may be acceptable unto thee for honour of thy Just ones: and by thy mercy profitable unto us for our salvation. Through.

\communion{Whosoever shall do the will of my Father who is in heaven: the same is my brother, and sister, and mother, saith the Lord.}

\postcommunion
\lett{G}{rant} to us, we beseech thee, O Lord: at the intercession of thy holy Martyrs Tryphon, Respicius and Nympha; that those things which we touch with our lips we may receive in purity of heart. Through.


\subby{St. Martin of Tours}
\feastday{{St. Martin Tours}}
\fancyhead[RE,LO]{11 November}
\begin{inhead}
	{Greater Double\\
		11 November}
\end{inhead}

\begin{rubric}
	The Hymns and Versicles are from the First Common of a Confessor Bishop (p. \pageref{CommonConfessorBishopI}), with the following Antiphons.
\end{rubric}

\antiphon{Mag. \& Ben.}{O blessed Martin, {\dag} whose righteous soul possesseth Paradise; whereat the Angels triumph, the Archangels are jubilant: thee the choir of Saints proclaimeth, the throng of Virgins inviteth, saying, Abide with us for ever.}

\antiphon{Mag.}{O blessed Bishop, {\dag} who with all his heart loved Christ the King, and feared not the dominion of princes! O holy soul, which, although withheld from the persecutor's sword, lacked not the palm of martyrdom!}

\introit
\lett{T}{he} Lord hath established a covenant of peace with him, and made him a prince: that he should have the dignity of the priesthood for ever. \textit{Ps.} Lord, remember David: and all his trouble.

\collect
\lett{O}{God,} who seest that we stand not in our own strength: mercifully grant; that, by the intercession of blessed Martin, thy Confessor and Bishop, we may be defended against all adversities. Through.

\begin{rubric}
    Commemoration is made of St. Mennas, from the First Common of a Martyr not a Bishop (p. \pageref{CommonMartyrNotBishopI}).
\end{rubric}

\readingcitation{Epistle}{Ecclesiasticus 44:16}
\lett{E}{noch} pleased the Lord, and was translated, Being an example of repentance to all generations. Noah was found perfect and righteous; In the season of wrath he was taken in exchange for the world; Therefore was there left a remnant unto the earth, When the flood came. Everlasting covenants were made with him, That all flesh should no more be blotted out by a flood. Abraham was a great father of a multitude of nations; And there was none found like him in glory; Who kept the law of the Most High, And was taken into covenant with him: In his flesh he established the covenant; And when he was proved, he was found faithful. Therefore he assured him by an oath, That the nations should be blessed in his seed; That he would multiply him as the dust of the earth, And exalt his seed as the stars, And cause them to inherit from sea to sea, And from the River unto the utmost part of the earth. In Isaac also did he establish likewise, for Abraham his father’s sake, The blessing of all men, and the covenant: And he made it rest upon the head of Jacob; He acknowledged him in his blessings, And gave to him by inheritance, And divided his portions; Among twelve tribes did he part them. And he brought out of him a man of mercy, Which found favour in the sight of all flesh.

\gradall{Behold a great priest, who in his days pleased God. ℣. There was none found like unto him, who kept the law of the Most High.}{Alleluia, alleluia. ℣. The blessed man Saint Martin, Bishop of the city of Tours, entered into rest: whom Angels and Archangels, Thrones, Dominations, and Virtues received. Alleluia.}

\readingcitation{Gospel}{Luke 11:33}
\lett{A}{t that time:} Jesus said unto his disciples: No man, when he hath lighted a candle, putteth it in a secret place, neither under a bushel, but on a candlestick, that they which come in may see the light. The light of the body is the eye: therefore when thine eye is single, thy whole body also is full of light; but when thine eye is evil, thy body also is full of darkness. Take heed therefore that the light which is in thee be not darkness. If thy whole body therefore be full of light, having no part dark, the whole shall be full of light, as when the bright shining of a candle doth give thee light.

\offertory{My truth and my mercy shall be with him: and in my name shall his horn be exalted.}

\secret
\lett{S}{anctify,} we beseech thee, O Lord God, these gifts which we offer on the solemnity of thy holy Bishop Martin: that our life may ever thereby be directed both in prosperity and adversity. Through.

\begin{rubric}
    Commemoration is made of St. Mennas, from the First Common of a Martyr not a Bishop (p. \pageref{CommonMartyrNotBishopI}).
\end{rubric}

\communion{Blessed is the servant, whom the lord when he cometh shall find watching: verily I say unto you, that he shall make him ruler over all his goods.}

\postcommunion
\lett{G}{rant,} we beseech thee, O Lord our God: that through the intercession of them, on whose festival they are offered, these sacraments may avail for our salvation. Through.

\begin{rubric}
    Commemoration is made of St. Mennas, from the First Common of a Martyr not a Bishop (p. \pageref{CommonMartyrNotBishopI}).
\end{rubric}


\subby{St. Mennas}
\feastday{{St. Mennas}}
\fancyhead[RE,LO]{11 November}
\begin{inhead}
	{Memorial\\
		11 November}
\end{inhead}

\begin{rubric}
	The propers are from the Second Common of a Martyr not a Bishop (p. \pageref{CommonMartyrNotBishopII}), with the Prayers from the First Common of a Martyr not a Bishop (p. \pageref{CommonMartyrNotBishopI}).
\end{rubric}


\subby{Pope St. Martin I}
\feastday{{Pope St. Martin I}}
\fancyhead[RE,LO]{12 November}
\begin{inhead}
	{Memorial\\
		12 November}
\end{inhead}

\begin{rubric}
	The propers are from the Second Common of a Martyr Bishop (p. \pageref{CommonMartyrBishopII}), with the second additional Epistle from the Second Common of a Martyr not a Bishop (p. \pageref{1Peter413}) and the Gospel from the First Common of a Martyr Bishop (p. \pageref{CommonMartyrBishopI}).
\end{rubric}


\subby{St. Britius of Tours}
\feastday{{St. Britius Tours}}
\fancyhead[RE,LO]{13 November}
\begin{inhead}
	{Memorial\\
		13 November}
\end{inhead}

\begin{rubric}
	The propers are from the First Common of a Confessor Bishop (p. \pageref{CommonConfessorBishopI}).
\end{rubric}


\subby{St. Gregory Palamas}
\feastday{{St. Gregory Palamas}}
\fancyhead[RE,LO]{14 November}
\begin{inhead}
	{Memorial\\
		14 November}
\end{inhead}

\begin{rubric}
	The propers are from the First Common of a Bishop Confessor (p. \pageref{CommonConfessorBishopI}).
\end{rubric}


\subby{St. Gregory the Wonder-worker}
\feastday{{St. Gregory Wonder-worker}}
\fancyhead[RE,LO]{17 November}
\begin{inhead}
	{Memorial\\
		17 November}
\end{inhead}

\begin{rubric}
	The propers are from the First Common of a Bishop Confessor (p. \pageref{CommonConfessorBishopI}), except for the following.
\end{rubric}

\readingcitation{Gospel}{Mark 11:22}
\lett{A}{t that time:} Jesus answering his disciples, saith unto them: Have faith in God. For verily I say unto you, That whosoever shall say unto this mountain, Be thou removed, and be thou cast into the sea; and shall not doubt in his heart, but shall believe that those things which he saith shall come to pass; he shall have whatsoever he saith. Therefore I say unto you, What things soever ye desire, when ye pray, believe that ye receive them, and ye shall have them.


\subby{St. Hilda of Whitby}
\feastday{{St. Hilda Whitby}}
\fancyhead[RE,LO]{17 November}
\begin{inhead}
	{Memorial\\
		17 November}
\end{inhead}

\begin{rubric}
	The propers are from the First Common of a Virgin (p. \pageref{CommonVirginOnlyI}), except for the following.
\end{rubric}

\collect
\lett{G}{rant,} we beseech thee, almighty God: that we, who rejoice in the yearly solemnity of blessed Hilda thy Virgin, may by her intercession be led from our old nature to newness of life. Through.

\gradall{Thou hast loved justice and hated iniquity. ℣. Wherefore God, even thy God, hath anointed thee with the oil of gladness.}{Alleluia, alleluia. ℣. For I am jealous over you with godly jealousy: for I have espoused you to one husband, that I may present you as a chaste virgin to Christ. Alleluia.}

\secret
\lett{W}{e} bring unto thee, O Lord, our offerings for the sacrifice: that we, being reconciled to thy mercy through the merits of blessed Hilda thy Virgin, may be made thereby a living sacrifice acceptable unto thee. Through.

\postcommunion
\lett{O}{Lord,} through whom we have received the benediction of this heavenly banquet, we entreat thee: that as unto us it is thy sacrament, so at the intercession of blessed Hilda thy Virgin, it may effectually avail for our salvation. Through.


\subby{Dedication of the Basilica of the Holy Apostles Peter \& Paul}
\feastday{{Basilica of Peter \& Paul}}
\fancyhead[RE,LO]{18 November}
\begin{inhead}
	{Greater Double\\
		18 November}
\end{inhead}

\begin{rubric}
	The propers are from the Dedication of a Church (p. \pageref{CommonDedication}).
\end{rubric}

\begin{rubric}
	\textsc{Note,} The Creed is said.
\end{rubric}


\subby{Pope St. Pontianus}
\feastday{{Pope St. Pontianus}}
\fancyhead[RE,LO]{19 November}
\begin{inhead}
	{Memorial\\
		19 November}
\end{inhead}

\begin{rubric}
	The propers are from the First Common of a Martyr Bishop (p. \pageref{CommonMartyrBishopI}), except for the Gospel from the Second Common of a Martyr not a Bishop (p. \pageref{CommonMartyrNotBishopII}).
\end{rubric}


\subby{St. Edmund}
\feastday{{St. Edmund}}
\fancyhead[RE,LO]{20 November}
\begin{inhead}
	{Double\\
		20 November}
\end{inhead}

\begin{rubric}
	The propers are from the First Common of a Martyr not a Bishop (p. \pageref{CommonMartyrNotBishopI}), except for the following.
\end{rubric}

\collect
\lett{O}{God} of unspeakable mercy, who didst give grace to the most blessed King Edmund to overcome the enemy by dying for thy name: mercifully grant unto this thy family; that, by his intercession they may be found worthy to vanquish and destroy in themselves the temptations of their ancient foe. Through.

\secret
\lett{W}{e} beseech thee, almighty God, mercifully regard this sacrifice of our redemption: and at the intercession of blessed Edmund, thy King and Martyr, graciously accept it for this thy family. Through.

\postcommunion
\lett{L}{et} the homage of our service be pleasing unto thee, almighty God: that these holy things which we have received may, at the intercession of blessed Edmund, thy King and Martyr, be profitable unto us for the obtaining of the rewards of everlasting life. Through.


\subby{Presentation of the Blessed Virgin Mary}
\feastday{{Presentation B.V.M.}}
\fancyhead[RE,LO]{21 November}
\begin{inhead}
	{Greater Double\\
		21 November}
\end{inhead}

\begin{rubric}
	The Hymns and Versicles are from the Common of the Blessed Virgin Mary (p. \pageref{CommonBVM}), except for the Antiphon for I Evensong and II Evensong.
\end{rubric}

\antiphon{Mag.}{O ever blessed Mother of God, {\dag} Mary ever Virgin, temple of the Godhead, hallowed shrine of the Holy Spirit: thou only, above all others, wast acceptable to our Lord Jesus Christ, alleluia.}

\begin{rubric}
	The Mass propers are from the Common of the Blessed Virgin Mary (p. \pageref{CommonBVM}), except for the following.
\end{rubric}

\begin{rubric}
	\textsc{Note,} The Creed is said.
\end{rubric}

\collect
\lett{O}{God,} who didst will that blessed Mary ever Virgin, the dwelling-place of Holy Ghost, should this day be presented in the temple: grant, we beseech thee; that through her intercession we may be found worthy to be presented in the temple of thy glory. Through . . . in the unity of the same.

\secret
\lett{T}{hrough} thy mercy, O Lord, and the intercession of blessed Mary ever Virgin, may this oblation avail for our prosperity and peace, both now and for ever. Through.

\postcommunion
\lett{G}{rant,} we beseech thee, O Lord: that we, who have received these aids to our salvation, may at all times and in all places be protected through the advocacy of blessed Mary ever Virgin; in whose honour we have made these offerings to thy majesty. Through.

\begin{rubric}
	When this Feast Day is only commemorated, the following is said for the Last Gospel.
\end{rubric}

\readingcitation{Last Gospel.}{Luke 11:27}
\lett{A}{t that time:} As Jesus spake unto the multitudes, a certain woman of the company lifted up her voice, and said unto him: Blessed is the womb that bare thee, and the paps which thou hast sucked. But he said: Yea rather, blessed are they that hear the word of God, and keep it.


\subby{St. Gelasius}
\feastday{{St. Gelasius}}
\fancyhead[RE,LO]{21 November}
\begin{inhead}
	{Memorial\\
		21 November}
\end{inhead}

\begin{rubric}
	The propers are from the First Common of a Confessor Bishop (p. \pageref{CommonConfessorBishopI}).
\end{rubric}


\subby{St. Columbanus}
\feastday{{St. Columbanus}}
\fancyhead[RE,LO]{21 November}
\begin{inhead}
	{Memorial\\
		21 November}
\end{inhead}

\begin{rubric}
	The propers are from the Common of Abbots (p. \pageref{CommonAbbots}).
\end{rubric}


\subby{St. Cecilia}
\feastday{{St. Cecilia}}
\fancyhead[RE,LO]{22 November}
\begin{inhead}
	{Greater Double\\
		22 November}
\end{inhead}

\begin{rubric}
	The Hymns \& Versicles are from the First Common of a Virgin Martyr (p. \pageref{CommonVirginMartyrI}), with the following Antiphons.
\end{rubric}

\antiphon{Mag.}{It is secret, {\dag} O Valerian, what now I wish to tell thee: I have an Angel of God who loves me, and guards my body with exceeding zeal.}

\antiphon{Ben.}{When the dawn of day was breaking, {\dag} Cecilia the Martyr cried, saying: Forward, soldiers of Christ! Cast off the works of darkness, and put on the armour of light.}

\antiphon{Mag.}{This noble maiden {\dag} alway bare the glorious Gospel of Christ in her bosom: without ceasing, she continued in prayer and supplication, and celestial converse, night and day.}

\introit
\lett{I}{will} speak of thy testimonies, even before kings, and will not be ashamed: and my delight shall be in thy commandments, which I have loved exceedingly. \textit{Ps.} Blessed are those that are undefiled in the way: and walk in the law of the Lord.

\collect
\lett{O}{God,} who makest us glad with the yearly solemnity of blessed Cecilia thy Virgin and Martyr: grant, that we do venerate her in our service, so we may follow the example of her godly conversation. Through.

\readingcitation{Epistle}{Ecclesiasticus 51:9}
\lett{O}{Lord my God,} I lifted up my supplication from the earth, And prayed for deliverance from death. I called upon the Lord, the Father of my Lord, That he would not forsake me in the days of affliction, In the time when there was no help against the proud. I will praise thy name continually, And will sing praise with thanksgiving; And my supplication was heard: For thou savedst me from destruction, And deliveredst me from the evil time: Therefore will I give thanks and praise unto thee, And bless thy name, O Lord our God.

\gradall{Hearken, O daughter, and consider, incline thine ear: so shall the King have pleasure in thy beauty. ℣. In thy comeliness and in thy beauty, go forth, proceed prosperously, and reign.}{Alleluia, alleluia. ℣. The five wise virgins took oil in their vessels with their lamps: and at midnight there was a cry made: Behold, the bridegroom cometh: go ye out to meet Christ the Lord. Alleluia.}

\begin{rubric}
	The Gospel is from the First Common of a Virgin Martyr (p. \pageref{CommonVirginMartyrI}).
\end{rubric}

\offertory{The Virgins that be her fellow shall be brought unto the King: they that bear her company shall be brought unto thee, with joy and gladness: and shall enter into the palace of the Lord the King.}

\secret
\lett{W}{e} beseech thee, O Lord: that by the intercession of blessed Cecilia, thy Virgin and Martyr, this sacrifice of atonement and praise may ever render us worthy of thy loving-kindness. Through.

\communion{Let the proud be confounded, for they go wickedly about to destroy me: but I will be occupied in thy commandments and in thy statures, that I be not ashamed.}

\postcommunion
\lett{O}{Lord,} who hast satisfied thy family with sacred gifts: we beseech thee, that we may at all times be comforted by the intercession of her whose festival we celebrate. Through.


\subby{St. Clement}
\feastday{{St. Clement}}
\fancyhead[RE,LO]{23 November}
\begin{inhead}
	{Double\\
		23 November}
\end{inhead}

\begin{rubric}
	The Hymns are from the First Common of Many Martyrs (p. \pageref{CommonMartyrsI}), with the following Versicles \& Antiphons.
\end{rubric}

℣. Thou hast crowned him with glory and honour, O Lord.

℟. And madest him to have dominion of the works of thy hands.

\antiphon{Mag.}{Let us all pray {\dag} to our Lord Jesus Christ, that he would open a fountain for his Confessors.}\\

℣. Let the Saints be joyful with glory.

℟. Let them rejoice in their beds.

\antiphon{Ben.}{When he had started going to the sea, {\dag} the people cried out with loud voices, O Lord Jesus Christ, save him: and Clement responded, weeping, Heavenly Father, receive my spirit.}\\

℣. Let the Saints be joyful with glory.

℟. Let them rejoice in their beds.

\antiphon{Mag.}{In the heavenly kingdom {\dag} rejoice the souls of the Blessed, who followed the footsteps of Christ their Master: and since for love of him they poured forth their life-blood, therefore with Christ do they exult for ever.}

\introit
\lett{T}{he} Lord saith: My words which I have put in thy mouth, shall not depart out of thy mouth: and thy gifts shall be accepted upon mine altar. \textit{Ps.} Blessed is the man that feareth the Lord: he hath great delight in his commandments.

\collect
\lett{O}{God,} who makest us glad with the yearly solemnity of blessed Clement thy Martyr and Bishop: mercifully grant; that as we now celebrate his birthday, so we may imitate his constancy in suffering. Through.

\begin{rubric}
    Commemoration is made of St. Felicitas (p. \pageref{FelicityCollect}).
\end{rubric}

\readingcitation{Epistle}{Philippians 3:17}
\lett{B}{rethren:} Be followers together of me, and mark them which walk so as ye have us for an ensample. (For many walk, of whom I have told you often, and now tell you even weeping, that they are the enemies of the cross of Christ: Whose end is destruction, whose God is their belly, and whose glory is in their shame, who mind earthly things.) For our conversation is in heaven; from whence also we look for the Saviour, the Lord Jesus Christ: Who shall change our vile body, that it may be fashioned like unto his glorious body, according to the working whereby he is able even to subdue all things unto himself. Therefore, my brethren dearly beloved and longed for, my joy and crown, so stand fast in the Lord, my dearly beloved. I beseech Euodias, and beseech Syntyche, that they be of the same mind in the Lord. And I intreat thee also, true yokefellow, help those women which laboured with me in the gospel, with Clement also, and with other my fellowlabourers, whose names are in the book of life.

\gradall{The Lord sware, and will not repent: Thou art a priest for ever, after the order of Melchisedech. ℣. The Lord said unto my Lord: Sit thou on my right hand.}{Alleluia, alleluia. ℣. This is a priest whom the Lord hath crowned. Alleluia.}

\begin{rubric}
	The Gospel is of the Second Common of a Confessor Bishop (\pageref{CommonConfessorBishopII}).
\end{rubric}

\offertory{My truth and my mercy shall be with him: and in my name shall his horn be exalted.}

\secret
\lett{S}{anctify,} O Lord, the gifts which we offer unto thee: and, at the intercession of blessed Clement, thy Martyr and Bishop, cleanse us thereby from the defilement of our sins. Through.

\begin{rubric}
    Commemoration is made of St. Felicitas (p. \pageref{FelicitySecret}).
\end{rubric}

\communion{Blessed is the servant, whom the Lord when he cometh shall find watching: verily I say unto you, that he shall make him ruler over all his goods.}

\postcommunion
\lett{O}{Lord} our God, who hast fulfilled us with the partaking of the sacred Body, and the precious Blood, we beseech thee: that those things which we perform with godly devotion, we may at the intercession of blessed Clement, thy Martyr and Bishop, attain in the assurance of our redemption. Through.

\begin{rubric}
    Commemoration is made of St. Felicitas (p. \pageref{FelicityPostcommunion}).
\end{rubric}

%MANUAL ADJUSTMENT:
\clearpage
\subby{St. Felicitas}
\feastday{{St. Felicitas}}
\fancyhead[RE,LO]{23 November}
\begin{inhead}
    {Memorial\\
23 November}
\end{inhead}

\begin{rubric}
	The propers are from the Common of a Martyr not a Virgin (p. \pageref{CommonMartyrNotVirgin}), except for the following.
\end{rubric}

\collect\label{FelicityCollect}
\lett{G}{rant,} we beseech thee, almighty God: that we, celebrating the festival of blessed Felicitas, thy Martyr, may be protected by her merits and prayers. Through.

\secret\label{FelicitySecret}
\lett{G}{raciously} hearken, O Lord, to the prayers of thy people: and make us to rejoice in the intercession of her whose festival thou dost grant unto us to celebrate. Through.

\postcommunion\label{FelicityPostcommunion}
\lett{W}{e} humbly beseech thee, almighty God: that by the intercession of thy Saints thou wouldest both multiply in us thy gifts, and likewise dispose our times according to thy will. Through.


\subby{St. Chrysogonus}
\feastday{{St. Chrysogonus}}
\fancyhead[RE,LO]{24 November}
\begin{inhead}
    {Memorial\\
24 November}
\end{inhead}

\begin{rubric}
	The propers are from the First Common of a Martyr not a Bishop (p. \pageref{CommonMartyrNotBishopI}, except for the following.
\end{rubric}

\collect
\lett{A}{ssist} us, O Lord, in these our supplications: that we, who acknowledge ourselves to be guilty by reason of our iniquity, may be delivered by the intercession of thy blessed Martyr Chrysogonus. Through.

\secret
\lett{W}{e} beseech thee, O Lord, mercifully to accept these our oblations: that at the intercession of thy blessed Martyr Chrysogonus, we may be defended against all perils. Through.

\postcommunion
\lett{O}{Lord,} let the partaking of thy sacrament both cleanse us from our secret faults, and deliver us from the snares of our enemies. Through.


\subby{St. Catherine of Alexandria}
\feastday{{St. Catherine Alexandria}}
\fancyhead[RE,LO]{25 November}
\begin{inhead}
    {Double\\
25 November}
\end{inhead}

\begin{rubric}
	The propers are from the First Common of a Virgin Martyr (p. \pageref{CommonVirginMartyrI}), except for the following.
\end{rubric}

\collect
\lett{O}{God,} who didst give the law to Moses on the height of Mount Sinai, and in the same place didst through thy holy Angels wondrously bestow the body of blessed Catherine, thy Virgin and Martyr: grant, we beseech thee; that by her merits and intercession we may be enabled to attain unto that mount which is Christ. Who liveth.

\secret
\lett{R}{eceive,} O Lord, the gifts which we offer on the solemnity of blessed Catherine, thy Virgin and Martyr: through whose advocacy we trust to be delivered. Through.

\postcommunion
\lett{M}{ay} the mysteries which we have received be for our succour, O Lord: and at the intercession of blessed Catherine, thy Virgin and Martyr, cause us to rejoice in thy continual protection. Through.


\subby{St. Peter of Alexandria}
\feastday{{St. Peter Alexandria}}
\fancyhead[RE,LO]{26 November}
\begin{inhead}
    {Memorial\\
26 November}
\end{inhead}

\begin{rubric}
	The propers are from the First Common of a Martyr Bishop (p. \pageref{CommonMartyrBishopI}).
\end{rubric}







\fancyhead[RE,LO]{}\fancyhead[RO,LE]{}
\fancyhead[C]{}\thispagestyle{empty}
\phantomsection
\addcontentsline{toc}{chapter}{Commons of the Saints}
  \begin{center}
   \includegraphics[width=98mm, height=137mm]{Saints.eps}
   \par
   \vspace{2ex}
   	\textsc{\Huge{Commons of the\\
   	Saints}}
   \end{center}
%\phantomsection
%\addcontentsline{toc}{section}{Common of Saints}
\fancyhead[RO,LE]{}
\fancyhead[RE,LO]{}
\fancyhead[C]{\LARGE Common of Saints}

\subby{Common of Vigils of the Apostles}\label{CommonVigilApostles}
\fancyhead[RO,LE]{\textit{Apostolic Vigils}}

\introit
\lett{A}{s} for me, I am like a green olive-tree in the house of the Lord, my trust is in the tender mercy of my God: and I will hope in thy name, for thy saints like it well. \textit{Ps.} Why boastest thou thyself, thou tyrant: that thou canst do mischief?

\collect
\lett{G}{rant,} we beseech thee, almighty God: that the venerable solemnity of blessed \textit{N.} thine Apostle, which we here prevent, may increase our devotion and set forward our salvation. Through.
\begin{rubric}
	If, however, the preceding Collect has already been said in the Mass, or for the commemoration of a Confessor Bishop, then the following is to be said:
\end{rubric}
\lett{W}{e} beseech thee, almighty God: that as we do prevent the festival of blessed \textit{N.} thine Apostle, so he may implore thy mercy for us; that we being delivered from our iniquities, may likewise be defended against all adversities. Through.
\begin{rubric}
	\nth{2} and \nth{3} Collects as assigned for Vigils, according to the Rubrics.
\end{rubric}

\readingcitation{Epistle}{Ecclesiasticus 44:22}
%RV:
\lett{I}{n} Isaac did the Lord establish, for Abraham his father's sake, The blessing of all men, and the covenant: And he made it rest upon the head of Jacob; He acknowledged him in his blessings, And gave to him by inheritance, And divided his portions; Among twelve tribes did he part them. And he brought out of him a man of mercy, Which found favour in the sight of all flesh; A man beloved of God and men, even Moses, Whose memorial is blessed. He made him like to the glory of the saints, And magnified him in the fears of his enemies. By his words he caused the wonders to cease; He glorified him in the sight of kings; He gave him commandment for his people, And shewed him part of his glory. He sanctified him in his faithfulness and meekness; He chose him out of all flesh. He made him to hear his voice, And led him into the thick darkness, And gave him commandments face to face, Even the law of life and knowledge, That he might teach Jacob the covenant, And Israel his judgments. He exalted Aaron, a holy man like unto him, Even his brother, of the tribe of Levi. He established for him an everlasting covenant, And gave him the priesthood of the people; He beautified him with comely ornaments, And girded him about with a robe of glory.

%Traditional Text:
%\lett{T}{he} blessing of the Lord was upon the head of the righteous. Therefore did the Lord give him an heritage, and divided his portions; among the twelve tribes did he part them: and he found favour in the sight of all flesh. And he magnified him, so that his enemies stood in fear of him. By his words he caused the wonders to cease, and he made him glorious in the sight of kings, and gave him a commandment for his people, and shewed him his glory. He sanctified him in his faithfulness and meekness, and chose him out of all men. And he gave him commandments before his face, even the law of life and knowledge, and made him to be exalted. An everlasting covenant he made with him, and girded him about with the girdle of righteousness: and the Lord crowned him with the crown of glory.

\gradual{The righteous shall flourish like a palm-tree: and shall spread abroad like a cedar in Libanus in the house of the Lord. ℣. To tell of thy loving-kindness early in the morning, and of thy truth in the night-season.}

\readingcitation{Gospel}{John 15:12}
\lett{A}{t that time:} Jesus said unto his disciples: This is my commandment, That ye love one another, as I have loved you. Greater love hath no man than this, that a man lay down his life for his friends. Ye are my friends, if ye do whatsoever I command you. Henceforth I call you not servants; for the servant knoweth not what his lord doeth: but I have called you friends; for all things that I have heard of my Father I have made known unto you. Ye have not chosen me, but I have chosen you, and ordained you, that ye should go and bring forth fruit, and that your fruit should remain: that whatsoever ye shall ask of the Father in my name, he may give it you.

\offertory{Thou hast crowned him with glory and worship: thou hast made him to have dominion of the works of thy hands, O Lord.}

\secret
\lett{O}{Lord,} who didst exalt blessed \textit{N.} to be numbered among thine Apostles: grant that we thy people, who on this day, preventing his heavenly birth, do offer unto thee these holy mysteries, may be assisted by his intercession both in the making of our supplications before thee, and in the obtaining of all that we desire. Through.
\begin{rubric}
	Other Secrets as above.
\end{rubric}
\communion{His honour is great in thy salvation: glory and great worship shalt thou lay upon him, O Lord.}
\postcommunion
\lett{W}{e} beseech thee, O Lord, mercifully to hear the supplication of thy holy Apostle \textit{N.}: that thou wouldest vouchsafe unto us thy pardon, and grant us thy everlasting healing. Through.
\begin{rubric}
	Other Postcommunions as above.
\end{rubric}


\subby{Common of the Apostles}\label{CommonApostles}
\fancyhead[RO,LE]{\textit{Apostles}}

\subbysub{I Evensong}

%\begin{rubric}
%	In I Evensong, the Office Hymn is of Mattins with the following Versicle \& Antiphon.
%\end{rubric}

\gregorioscore{resources/gabc/Common/ApostlesEvensong.gabc}

℣. Their sound is gone out into all lands.

℟. And their words into the ends of the world.

\antiphon{Mag.}{They will deliver you up {\dag} to the councils, and they will scourge you in their synagogues; and ye shall be brought before governors and kings for my sake, for a testimony against them and the Gentiles.}

%\begin{multicols}{2}
%\lett{L}{et} heav'n's exultant praises ring,
%
%\secondline{And earth with joy responsive sing:}
%
%\thirdline{Th' Apostles' deeds and high estate}
%
%This festal-tide we celebrate.\\
%
%O ye who, thron'd in glory dread,
%
%Shall judge the living and the dead,
%
%True lights, the world illumining,
%
%Regard the suppliant prayer we bring.\\
%
%The gates of heav'n, at your command,
%
%To all or clos'd or open stand:
%
%May we, at your august decree,
%
%Be loos'd from our iniquity.\\
%
%The power, of old to you convey'd,
%
%Sickness and health alike obey'd:
%
%May ye our ailing souls once more
%
%To strength and holiness restore;\\
%
%That Christ, th' unerring Judge of doom,
%
%When he at time's last end shall come,
%
%May grant us, for his mercy's sake,
%
%Of joys eternal to partake.\\
%
%All laud to God the Father be;
%
%All praise, eternal Son, to thee;
%
%All glory, as is ever meet,
%
%To God the Holy Paraclete. Amen.\\

\subbysub{Mattins}

\invitatoryhymn

\gregorioscore{resources/gabc/Common/ApostlesInvitatory.gabc}

\officehymn

\gregorioscore{resources/gabc/Common/ApostlesMattins.gabc}

℣. They declared the work of God.

℟. And wisely considered of his doing.
%\end{multicols}

\antiphon{Ben.}{Ye which have forsaken all, {\dag} and followed me, shall receive an hundredfold, and shall inherit everlasting life.}

\subbysub{II Evensong}

\begin{rubric}
	The Office Hymn \& Versicle are of I Evensong with the following Antiphon.
\end{rubric}

\antiphon{Mag.}{Be ye valiant in warfare, {\dag} and contend ye with the old serpent: and ye shall receive an eternal kingdom, alleluia.}


\subby{Common of the Apostles in Eastertide}\label{CommonApostlesEaster}
\fancyhead[RO,LE]{\textit{Apostles Easter}}

\subbysub{I Evensong}
%\begin{multicols}{2}
%\lett{T}{he} Apostles' hearts were full of pain
%
%\secondline{For their dear Lord so lately slain:}
%
%\thirdline{That Lord his servants' wicked train}
%
%With bitter scorn had dar'd arraign.\\
%
%With gentle voice the Angel gave
%
%The women tidings at the grave;
%
%`Forthwith your Master shall ye see:
%
%He goes before to Galilee.'\\
%
%And while with fear and joy they press'd
%
%To tell these tidings to the rest,
%
%Their Lord, their living Lord, they meet,
%
%And see his form, and kiss his feet.\\
%
%Th' Eleven, when they hear, with speed
%
%To Galilee forthwith proceed:
%
%That there they may behold once more
%
%The Lord's dear face, as oft before.\\
%
%We pray thee, King with glory deck'd,
%
%In this our Paschal joy, protect
%
%From all that death would fain effect
%
%Thy ransom'd flock, thine own elect.

\gregorioscore{resources/gabc/Common/ApostlesEvensongEaster.gabc}

    \par\noindent
	\centerline{\rule{0.25\textwidth}{0.4pt}}
	\par\noindent
\begin{rubric}
	From Low Sunday, exclusive, until Ascension Thursday, exclusive:
\end{rubric}

\gregorioscore{resources/gabc/Common/ApostlesEvensongEasterEnding1.gabc}

%6. To thee who, dead, again dost live,
%
%All glory, Lord, thy people give;
%
%All glory, as is ever meet,
%
%To Father and to Paraclete. Amen.

\begin{rubric}
	From Ascension Thursday, inclusive, until Whitsunday, inclusive:
\end{rubric}

\gregorioscore{resources/gabc/Common/ApostlesEvensongEasterEnding2.gabc}

%6. All glory, Lord, to thee we pay,
%
%Ascending o'er the stars to-day;
%
%All glory, as is ever meet,
%
%To Father and to Paraclete. Amen.
    \par\noindent
	\centerline{\rule{0.25\textwidth}{0.4pt}}
	\par
℣. O ye holy and righteous, rejoice in the Lord, alleluia.

℟. God hath chosen you to him to be his inheritance, alleluia.
%\end{multicols}

\antiphon{Mag.}{Light perpetual {\dag} shall shine upon thy Saints, O Lord; and an ageless eternity, alleluia.}

\subbysub{Mattins}

\begin{rubric}
	The Invitatory Hymn is of I Evensong.
\end{rubric}

\officehymn

\gregorioscore{resources/gabc/Common/ApostlesMattinsEaster.gabc}

    \par\noindent
	\centerline{\rule{0.25\textwidth}{0.4pt}}
	\par\noindent
\begin{rubric}
	From Low Sunday, exclusive, until Ascension Thursday, exclusive:
\end{rubric}

\gregorioscore{resources/gabc/Common/ApostlesMattinsEasterEnding1.gabc}

\begin{rubric}
	From Ascension Thursday, inclusive, until Whitsunday, inclusive:
\end{rubric}

\gregorioscore{resources/gabc/Common/ApostlesMattinsEasterEnding2.gabc}
    \par\noindent
	\centerline{\rule{0.25\textwidth}{0.4pt}}
	\par
℣. Right dear in the sight of the Lord, alleluia.

℟. Is the death of his Saints, alleluia.

\antiphon{Ben.}{Daughters of Jerusalem, {\dag} come forth and behold the Martyrs with the diadems wherewith the Lord hath crowned them in the day of solemnity and rejoicing, alleluia, alleluia.}

\subbysub{II Evensong}

\begin{rubric}
	The Office Hymn is of I Evensong with the following Versicle \& Antiphon.
\end{rubric}

℣. Right dear in the sight of the Lord, alleluia.

℟. Is the death of his Saints, alleluia.

\antiphon{Mag.}{O ye holy and righteous, {\dag} rejoice in the Lord, alleluia: God hath chosen you to him to be his inheritance, alleluia.}


\subby{I: Common of One Martyr Bishop out of Eastertide}\label{CommonMartyrBishopI}
\fancyhead[RO,LE]{\textit{Martyr Bishop I}}

\subbysub{I Evensong}

\gregorioscore{resources/gabc/Common/MartyrEvensong.gabc}

℣. Thou hast crowned him with glory and honour, O Lord. 

℟. And madest him to have dominion of the works of thy hands.

\antiphon{Mag.}{This is a Martyr {\dag} who strove for his Master's precepts, even unto death: and feared not the words of evil men, forasmuch as he was stablished on a sure foundation.}

\subbysub{Mattins}

\begin{rubric}
	The Invitatory Hymn is of I Evensong.
\end{rubric}

\officehymn

\gregorioscore{resources/gabc/Common/MartyrMattins.gabc}

℣. The righteous shall flourish like a palm-tree.

℟. And shall spread abroad like a cedar in Libanus.

\antiphon{Ben.}{He that hateth his life {\dag} in this world shall keep it unto life eternal.}\par\noindent

\subbysub{II Evensong}

\begin{rubric}
	The Office Hymn is of I Evensong with the following Versicle \& Antiphon.
\end{rubric}

℣. The righteous shall flourish like a palm-tree.

℟. And shall spread abroad like a cedar in Libanus.

\antiphon{Mag.}{Whosoever will come after me, {\dag} let him deny himself, and take up his cross, and follow me.}

\introit
\lett{T}{he} Lord hath established a covenant of peace with him, and made him a prince: that he should have the dignity of the priesthood for ever. \textit{Ps.} Lord, remember David: and all his trouble.

\collect
\lett{A}{lmighty} God, mercifully look upon our infirmities: that whereas we are oppressed by the burden of our sins, the glorious intercession of blessed \textit{N.} thy Martyr and Bishop may be our succour and defence. Through.

\readingcitation{Epistle}{James 1:12}
\lett{D}{early beloved:} Blessed is the man that endureth temptation: for when he is tried, he shall receive the crown of life, which the Lord hath promised to them that love him. Let no man say when he is tempted, I am tempted of God: for God cannot be tempted with evil, neither tempteth he any man: but every man is tempted, when he is drawn away of his own lust, and enticed. Then when lust hath conceived, it bringeth forth sin: and sin, when it is finished, bringeth forth death. Do not err, my beloved brethren. Every good gift and every perfect gift is from above, and cometh down from the Father of lights, with whom is no variableness, neither shadow of turning. Of his own will begat he us with the word of truth, that we should be a kind of firstfruits of his creatures.

\gradall{I have found David my servant, with my holy oil have I anointed him: my hand shall hold him fast, and my arm shall strengthen him. ℣. The enemy shall not be able to do him violence, the son of wickedness shall not hurt him.}{Alleluia, alleluia. ℣. Thou art a priest for ever, after the order of Melchisedech. Alleluia.}

\begin{rubric}
{In Septuagesimatide or Lent, replacing the Alleluia:}
\end{rubric}\par\noindent
\tract{Thou hast given him his heart's desire: and hast not denied him the request of his lips. ℣. For thou hast prevented him with the blessings of goodness. ℣. And hast set a crown of pure gold upon his head.}

\readingcitation{Gospel}{Luke 14:26}
\lett{A}{t that time:} Jesus said unto the multitudes:  If any man come to me, and hate not his father, and mother, and wife, and children, and brethren, and sisters, yea, and his own life also, he cannot be my disciple. And whosoever doth not bear his cross, and come after me, cannot be my disciple. For which of you, intending to build a tower, sitteth not down first, and counteth the cost, whether he have sufficient to finish it? lest haply, after he hath laid the foundation, and is not able to finish it, all that behold it begin to mock him, saying, This man began to build, and was not able to finish. Or what king, going to make war against another king, sitteth not down first, and consulteth whether he be able with ten thousand to meet him that cometh against him with twenty thousand? Or else, while the other is yet a great way off, he sendeth an ambassage, and desireth conditions of peace. So likewise, whosoever he be of you that forsaketh not all that he hath, he cannot be my disciple.

\offertory{My truth and my mercy shall be with him: and in my name shall his horn be exalted.}

\secret
\lett{W}{e} beseech thee, O Lord, mercifully to accept this our sacrifice which we offer unto thee pleading the merits of blessed \textit{N.} thy Martyr and Bishop: that the same may avail for our perpetual succour. Through.

\communion{I have sworn once by my holiness: His seed shall endure for ever: and his seat is like as the sun before me; he shall stand fast for evermore as the moon, and as the faithful witness in heaven.}

\postcommunion
\lett{W}{e} beseech thee, O Lord our God, that like as we whom thou hast refreshed by the partaking of thy sacred gift do offer unto thee our worship; so, by the intercession of blessed \textit{N.} thy Martyr and Bishop, we may perceive the benefit of the same. Through.


\subby{II: Common of One Martyr Bishop out of Eastertide}\label{CommonMartyrBishopII}
\fancyhead[RO,LE]{\textit{Martyr Bishop II}}
\begin{rubric}
	Daily Office propers as in the First Common of a Martyr Bishop (p. \pageref{CommonMartyrBishopI}).
\end{rubric}
\introit
\lett{O}{ye} priests of the Lord, bless ye the Lord: O ye holy and humble men of heart, bless ye the Lord. \textit{Ps.} O all ye works of the Lord, bless ye the Lord: praise him, and magnify him for ever.

\collect
\lett{O}{God,} who makest us glad with the yearly solemnity of blessed \textit{N.} thy Martyr and Bishop: mercifully grant; that, as we now celebrate his birthday, so we may likewise rejoice in his protection. Through.

\readingcitation{Epistle}{2 Corinthians 1:3}
\lett{B}{rethren:} Blessed be God, even the Father of our Lord Jesus Christ, the Father of mercies, and the God of all comfort; who comforteth us in all our tribulation, that we may be able to comfort them which are in any trouble, by the comfort wherewith we ourselves are comforted of God. For as the sufferings of Christ abound in us, so our consolation also aboundeth by Christ. And whether we be afflicted, it is for your consolation and salvation, which is effectual in the enduring of the same sufferings which we also suffer: or whether we be comforted, it is for your consolation and salvation. And our hope of you is stedfast, knowing, that as ye are partakers of the sufferings, so shall ye be also of the consolation: in Christ Jesus our Lord.

\gradall{Thou hast crowned him with glory and worship. ℣. Thou hast made him to have dominion of the works of thy hands, O Lord.}{Alleluia, alleluia. ℣. This is a priest whom the Lord hath crowned. Alleluia.}

\begin{rubric}
{In Septuagesimatide or Lent, replacing the Alleluia:}
\end{rubric}\par\noindent
\tract{Blessed is the man that feareth the Lord: he hath great delight in his commandments. ℣. His seed shall be mighty upon earth: the generation of the faithful shall be blessed. ℣. Riches and plenteousness shall be in his house: and his righteousness endureth for ever.}

\readingcitation{Gospel}{Matthew 16:24}
\lett{A}{t that time:} Jesus said unto his disciples: If any man will come after me, let him deny himself, and take up his cross, and follow me. For whosoever will save his life shall lose it: and whosoever will lose his life for my sake shall find it. For what is a man profited, if he shall gain the whole world, and lose his own soul? or what shall a man give in exchange for his soul? For the Son of man shall come in the glory of his Father with his angels; and then he shall reward every man according to his works.

\offertory{I have found David my servant, with my holy oil have I anointed him: my hand shall hold him fast, and my arm shall strengthen him.}

\secret
\lett{S}{anctify,} O Lord, the gifts which we dedicate to thee: that at the intercession of blessed \textit{N.}, thy Martyr and Bishop, they may obtain for us thy gracious favour. Through.

\communion{Thou hast set, Lord, a crown of pure gold upon his head.}

\postcommunion
\lett{M}{ay} this communion, O Lord, cleanse us from guilt: and, at the intercession of blessed \textit{N.}, thy Martyr and Bishop, make us partakers of thy heavenly healing. Through.


\subby{I: Common of One Martyr not a Bishop out of Eastertide}\label{CommonMartyrNotBishopI}
\fancyhead[RO,LE]{\textit{Martyr not Bishop I}}
\begin{rubric}
	Daily Office propers as in the First Common of a Martyr Bishop (p. \pageref{CommonMartyrBishopI}).
\end{rubric}

\introit
\lett{T}{he} righteous shall rejoice: in thy strength, O Lord: exceeding glad shall he be of thy salvation: thou hast given him his heart's desire. \textit{Ps.} For thou hast prevented him with the blessings of goodness: and hast set a crown of pure gold upon his head.

\collect
\lett{G}{rant,} we beseech thee, almighty God: that we, who devoutly celebrate the birthday of blessed \textit{N.}, thy Martyr, may by his intercession be stablished in the love of thy name. Through.

\readingcitation{Epistle}{Wisdom 10:10}
%RV:
\lett{W}{isdom} guided the righteous man in straight paths; She shewed him God's kingdom, and gave him knowledge of holy things; She prospered him in his toils, and multiplied the fruits of his labour; When in their covetousness men dealt hardly with him, She stood by him and made him rich; She guarded him from enemies, And from those that lay in wait she kept him safe, And over his sore conflict she watched as judge, That he might know that godliness is more powerful than all. When a righteous man was sold, wisdom forsook him not, But from sin she delivered him; She went down with him into a dungeon, And in bonds she left him not, Till she brought him the sceptre of a kingdom, And authority over those that dealt tyrannously with him; She shewed them also to be false that had mockingly accused him, And gave him eternal glory. 

%\lett{T}{he} Lord guided the just ìn right paths, shewed him the kingdom of God, and gave him knowledge of holy things: made him rich in his travails, and multiplied the fruit of his labours. In the covetousness of such as oppressed him he stood by him, and made him rich. He defended him from his enemies, and kept him safe from those that lay in wait, and in a sore conflict he gave him the victory, that he might know that godliness is stronger than all. When the just was sold, she forsook him not, but delivered him from sin: she went down with him into the pit, and left him not in bonds, till she brought him the sceptre of the kingdom, and power against those that oppressed him: as for them that had accused him, she shewed them to be liars, and the Lord God gave him perpetual glory.

\gradall{Blessed is the man that feareth the Lord: he hath great delight in his commandments. ℣. His seed shall be mighty upon earth: the generation of the faithful shall be blessed.}{Alleluia, alleluia. ℣. Thou hast set, O Lord, a crown of pure gold upon his head. Alleluia.}

\begin{rubric}
{In Septuagesimatide or Lent, replacing the Alleluia:}
\end{rubric}\par\noindent
\tract{Thou hast given him his heart's desire: and hast not denied him the request of his lips. ℣. For thou hast prevented him with the blessings of goodness. ℣. And hast set a crown of pure gold upon his head.}

\readingcitation{Gospel}{Matthew 10:34}
\lett{A}{t that time:} Jesus said unto his disciples: 34 Think not that I am come to send peace on earth: I came not to send peace, but a sword. For I am come to set a man at variance against his father, and the daughter against her mother, and the daughter in law against her mother in law. And a man's foes shall be they of his own household. He that loveth father or mother more than me is not worthy of me: and he that loveth son or daughter more than me is not worthy of me. And he that taketh not his cross, and followeth after me, is not worthy of me. He that findeth his life shall lose it: and he that loseth his life for my sake shall find it. He that receiveth you receiveth me, and he that receiveth me receiveth him that sent me. He that receiveth a prophet in the name of a prophet shall receive a prophet's reward; and he that receiveth a righteous man in the name of a righteous man shall receive a righteous man's reward. And whosoever shall give to drink unto one of these little ones a cup of cold water only in the name of a disciple, verily I say unto you, he shall in no wise lose his reward.

\offertory{Thou hast crowned him with glory and worship: thou hast made him to have dominion of the works of thy hands, O Lord.}

\secret
\lett{W}{e} beseech thee, O Lord, to accept our prayers and oblations: and graciously hearken unto us, whom thou dost cleanse by thy heavenly mysteries. Through.

\communion{If any man will come after me, let him deny himself, and take up his cross, and follow me.}

\postcommunion
\lett{G}{rant,} we beseech thee, O Lord our God: that like as we in this life do gladly honour the memory of thy saints; so we may rejoice to behold them for ever. Through.


\subby{II: Common of One Martyr not a Bishop out of Eastertide}\label{CommonMartyrNotBishopII}
\fancyhead[RO,LE]{\textit{Martyr not Bishop II}}

\begin{rubric}
	Daily Office propers as in the First Common of a Martyr Bishop (p. \pageref{CommonMartyrBishopI}).
\end{rubric}

\introit
\lett{T}{he} righteous shall rejoice in the Lord, and put his trust in him: and all they that are true of heart shall be glad. \textit{Ps.} Hear my voice, O God, in my prayer: preserve my life from fear of the enemy.

\collect
\lett{G}{rant,} we beseech thee, almighty God: that, at the intercession of blessed \textit{N.}, thy Martyr, we may both be delivered from all adversities which may happen to the body, and from all evil thoughts, which may assault and hurt the soul. Through.

\readingcitation{Epistle}{2 Timothy 2:8}
\lett{D}{early beloved:} Remember that Jesus Christ of the seed of David was raised from the dead according to my gospel: wherein I suffer trouble, as an evil doer, even unto bonds; but the word of God is not bound. Therefore I endure all things for the elect's sakes, that they may also obtain the salvation which is in Christ Jesus with eternal glory. But thou hast fully known my doctrine, manner of life, purpose, faith, longsuffering, charity, patience, persecutions, afflictions, which came unto me at Antioch, at Iconium, at Lystra; what persecutions I endured: but out of them all the Lord delivered me. Yea, and all that will live godly in Christ Jesus shall suffer persecution.

\gradall{Though the righteous fall, he shall not be cast away: for the Lord upholdeth him with his hand. ℣. He is ever merciful, and lendeth: and his seed is blessed.}{Alleluia, alleluia. ℣. He that followeth me shall not walk in darkness: but shall have the light of eternal life. Alleluia.}

\begin{rubric}
{In Septuagesimatide or Lent, replacing the Alleluia:}
\end{rubric}\par\noindent
\tract{Blessed is the man that feareth the Lord: he hath great delight in his commandments. ℣. His seed shall be mighty upon earth: the generation of the faithful shall be blessed. ℣. Riches and plenteousness shall be in his house: and his righteousness endureth for ever.}

\readingcitation{Gospel}{Matthew 10:26}
\lett{A}{t that time:} Jesus said unto his disciples: There is nothing covered, that shall not be revealed; and hid, that shall not be known. What I tell you in darkness, that speak ye in light: and what ye hear in the ear, that preach ye upon the housetops. And fear not them which kill the body, but are not able to kill the soul: but rather fear him which is able to destroy both soul and body in hell.

Are not two sparrows sold for a farthing? and one of them shall not fall on the ground without your Father. But the very hairs of your head are all numbered. Fear ye not therefore, ye are of more value than many sparrows.

Whosoever therefore shall confess me before men, him will I confess also before my Father which is in heaven.

\offertory{Thou hast set, O Lord, a crown of pure gold upon his head: he asked life of thee, and thou gavest it him, alleluia.}

\secret
\lett{G}{rant,} O Lord, that this our bounden service may be acceptable in thy sight: that these our oblations may, by the prayers of him on whose solemnity they are offered, be made profitable unto our salvation. Through.

\communion{If any man serve me, let him follow me: and where I am, there shall also my servant be.}

\postcommunion
\lett{W}{e} beseech thee, O Lord our God, that like as we, whom thou hast refreshed by the partaking of thy sacred gift, do ofter unto thee our worship: so by the intercession of blessed \textit{N.}, thy Martyr, we may perceive the benefit of the same. Through.

\begin{rubric}
	The following Epistles and Gospel may be used instead of the ones provided above.
\end{rubric}

\readingcitation{1. Epistle}{James 1:2}\label{James12}
\lett{D}{early beloved:} Count it all joy when ye fall into divers temptations; knowing this, that the trying of your faith worketh patience. But let patience have her perfect work, that ye may be perfect and entire, wanting nothing. If any of you lack wisdom, let him ask of God, that giveth to all men liberally, and upbraideth not; and it shall be given him. But let him ask in faith, nothing wavering. For he that wavereth is like a wave of the sea driven with the wind and tossed. For let not that man think that he shall receive any thing of the Lord. A double minded man is unstable in all his ways. Let the brother of low degree rejoice in that he is exalted: but the rich, in that he is made low: because as the flower of the grass he shall pass away. For the sun is no sooner risen with a burning heat, but it withereth the grass, and the flower thereof falleth, and the grace of the fashion of it perisheth: so also shall the rich man fade away in his ways.

Blessed is the man that endureth temptation: for when he is tried, he shall receive the crown of life, which the Lord hath promised to them that love him.

\readingcitation{2. Epistle}{1 Peter 4:13}\label{1Peter413}
\lett{D}{early beloved:} Rejoice, inasmuch as ye are partakers of Christ's sufferings; that, when his glory shall be revealed, ye may be glad also with exceeding joy. If ye be reproached for the name of Christ, happy are ye; for the spirit of glory and of God resteth upon you: on their part he is evil spoken of, but on your part he is glorified. But let none of you suffer as a murderer, or as a thief, or as an evildoer, or as a busybody in other men's matters. Yet if any man suffer as a Christian, let him not be ashamed; but let him glorify God on this behalf. For the time is come that judgment must begin at the house of God: and if it first begin at us, what shall the end be of them that obey not the gospel of God? And if the righteous scarcely be saved, where shall the ungodly and the sinner appear? Wherefore let them that suffer according to the will of God commit the keeping of their souls to him in well doing, as unto a faithful Creator.

\readingcitation{Another Gospel}{John 12:24}
\lett{A}{t that time:} Jesus said unto his disciples: Verily, verily, I say unto you, Except a corn of wheat fall into the ground and die, it abideth alone: but if it die, it bringeth forth much fruit. He that loveth his life shall lose it; and he that hateth his life in this world shall keep it unto life eternal. If any man serve me, let him follow me; and where I am, there shall also my servant be: if any man serve me, him will my Father honour.

\subby{I: Common of Many Martyrs out of Eastertide}\label{CommonMartyrsI}
\fancyhead[RO,LE]{\textit{Martyrs I}}

\subbysub{I Evensong}

\gregorioscore{resources/gabc/Common/MartyrsEvensong.gabc}

℣. Be glad, O ye righteous, and rejoice in the Lord.

℟. And be joyful, all ye that are true of heart.

\antiphon{Mag.}{For theirs is the kingdom of heaven, {\dag} who despised worldly living: who have won the rewards of the kingdom, and have washed their robes in the blood of the Lamb.}

\subbysub{Mattins}

\invitatoryhymn

\gregorioscore{resources/gabc/Common/MartyrsInvitatory.gabc}

\officehymn

\gregorioscore{resources/gabc/Common/MartyrsMattins.gabc}

℣. Let the Saints be joyful with glory.

℟. Let them rejoice in their beds.

\antiphon{Ben.}{The very hairs of your head {\dag} are all numbered: fear ye not therefore, ye are of more value than many sparrows.}

\subbysub{II Evensong}

\begin{rubric}
	The Office Hymn is of I Evensong, and the Versicle is of Mattins, with the following Antiphon.
\end{rubric}

\antiphon{Mag.}{In the heavenly kingdom {\dag} rejoice the souls of the Blessed, who followed the footsteps of Christ their Master: and since for love of him they poured forth their life-blood, therefore with Christ do they exult for ever.}

\introit
\lett{L}{et} the sorrowful sighing of the prisoners, O Lord, come before thee: reward thou our neighbours seven-fold into their bosom: avenge thou the blood of thy Saints that is shed. \textit{Ps.} O God, the heathen are come into thine inheritance: thy holy temple have they defiled: and made Jerusalem an heap of stones.

\collect
\begin{inhead}
	Many Martyrs not Bishops
\end{inhead}
\lett{O}{God,} who vouchsafest unto us to celebrate the birthday of thy holy Martyrs
\textit{N. and N.}: grant that we may rejoice in the everlasting felicity of their fellowship. Through.

\begin{inhead}
	Many Martyrs Bishops
\end{inhead}
\lett{P}{rotect} us, O Lord, we beseech thee, who observe the feast of thy blessed Martyrs and Bishops \textit{N. and N.}: and grant that by their meritorious supplication we may find favour in thy sight. Through.

\readingcitation{Epistle}{Wisdom 3:1}\label{Wisdom31}
%RV:
\lett{T}{he} souls of the righteous are in the hand of God, And no torment shall touch them. In the eyes of the foolish they seemed to have died; And their departure was accounted to be their hurt, And their journeying away from us to be their ruin: But they are in peace. For even if in the sight of men they be punished, Their hope is full of immortality; And having borne a little chastening, they shall receive great good; Because God made trial of them, and found them worthy of himself. As gold in the furnace he proved them, And as a whole burnt offering he accepted them. And in the time of their visitation they shall shine forth, And as sparks among stubble they shall run to and fro. They shall judge nations, and have dominion over peoples; And the Lord shall reign over them for evermore. 

%\lett{T}{he} souls of the righteous are in the hand of God, and there shall no torment touch them. In the sight of the unwise they seemed to die: and their departure is taken for misery, And their going from us to be utter destruction: but they are in peace. For though they be punished in the sight of men, yet is their hope full of immortality. And having been a little chastised, they shall be greatly rewarded: for God proved them, and found them worthy for himself. As gold in the furnace hath he tried them, and received them as a burnt offering. And in the time of their visitation they shall shine, and run to and fro like sparks among the stubble. They shall judge the nations, and have dominion over the people, and their Lord shall reign for ever.

\gradall{God is glorious in his holy ones: fearful in praises, doing wonders. ℣. Thy right hand, O Lord, is become glorious in power: thy right hand hath dashed in pieces the enemy.}{Alleluia, alleluia. ℣. The bodies of the Saints are buried in peace, but their name liveth for evermore. Alleluia.}

\begin{rubric}
{In Septuagesimatide or Lent, replacing the Alleluia:}
\end{rubric}\par\noindent
\tract{They that sow in tears, shall reap in joy. ℣. He that now goeth on his way weeping, and beareth forth good seed. ℣. Shall doubtless come again with joy, and bring his sheaves with him.}

\readingcitation{Gospel}{Luke 21:9}
\lett{A}{t that time:} Jesus said unto his disciples: When ye shall hear of wars and commotions, be not terrified: for these things must first come to pass; but the end is not by and by. Then said he unto them, Nation shall rise against nation, and kingdom against kingdom: and great earthquakes shall be in divers places, and famines, and pestilences; and fearful sights and great signs shall there be from heaven. But before all these, they shall lay their hands on you, and persecute you, delivering you up to the synagogues, and into prisons, being brought before kings and rulers for my name's sake. And it shall turn to you for a testimony. Settle it therefore in your hearts, not to meditate before what ye shall answer: for I will give you a mouth and wisdom, which all your adversaries shall not be able to gainsay nor resist. And ye shall be betrayed both by parents, and brethren, and kinsfolks, and friends; and some of you shall they cause to be put to death. And ye shall be hated of all men for my name's sake. But there shall not an hair of your head perish. In your patience possess ye your souls.

\begin{rubric}
	After Septuagesima, \emph{Alleluia} at the end of the following Offertory is omitted.
\end{rubric}

\offertory{God is wonderful in his holy ones: even the God of Israel, he will give strength and power unto his people: blessed be God, alleluia.}

\secret
\begin{inhead}
	Many Martyrs not Bishops
\end{inhead}
\lett{W}{e} beseech thee, O Lord, that the gifts which we offer unto thee of our bounden duty and service may be acceptable unto thee for the honour of thy Just ones: and by thy mercy profitable unto us for our salvation. Through.

\begin{inhead}
	Many Martyrs Bishops
\end{inhead}
\lett{A}{ssist} us mercifully, O Lord, in these our supplications which we make before thee in remembrance of thy Saints: that we who trust not in our own righteousness may be succoured by the merits of them that have found favour in thy sight. Through.

\communion{Though they be punished in the sight of men, yet hath God proved them: as gold in the furnace hath he tried them, and received them as a burnt offering.}

\postcommunion
\begin{inhead}
	Many Martyrs not Bishops
\end{inhead}
\lett{G}{rant} to us, we beseech thee, O Lord: at the intercession of thy holy Martyrs \textit{N. and N.}; that those things which we touch with our mouths we may receive in purity of heart. Through.

\begin{inhead}
	Many Martyrs Bishops
\end{inhead}
\lett{O}{Lord,} who hast fulfilled us with saving mysteries: we beseech thee that we may be aided by the prayers of those whose festival we celebrate. Through.


\subby{II: Common of Many Martyrs out of Eastertide}\label{CommonMartyrsII}
\fancyhead[RO,LE]{\textit{Martyrs II}}
\begin{rubric}
	Daily Office propers as in the First Common of Many Martyrs (p. \pageref{CommonMartyrsI}).
\end{rubric}

\begin{rubric}
	The Collect, Secret, \& Postcommunion are of the preceding Mass.
\end{rubric}

\introit
\lett{L}{et} the people tell of the wisdom of the Saints, and let the church shew forth their praise: their names shall live for evermore. \textit{Ps.} Rejoice in the Lord, O ye righteous: for it becometh well the just to be thankful.

\readingcitation{Epistle}{Wisdom 5:15}
%RV:
\lett{B}{ut} the righteous live for ever, And in the Lord is their reward, And the care for them with the Most High. Therefore shall they receive the crown of royal dignity And the diadem of beauty from the Lord's hand; Because with his right hand shall he cover them, And with his arm shall he shield them. He shall take his jealousy as complete armour, And shall make the whole creation his weapons for vengeance on his enemies: He shall put on righteousness as a breastplate, And shall array himself with judgement unfeigned as with a helmet; He shall take holiness as an invincible shield.

%\lett{B}{ut} the righteous live for evermore; their reward also is with the Lord, and the care of them is with the most High. Therefore shall they receive a glorious kingdom, and a beautiful crown from the Lord's hand: for with his right hand shall he cover them, and with his arm shall he protect them. He shall take to him his jealousy for complete armour, and make the creature his weapon for the revenge of his enemies. He shall put on righteousness as a breastplate, and true judgment instead of an helmet. He shall take holiness for an invincible shield.

\gradall{Our soul is escaped even as a bird out of the snare of the fowler. ℣. The snare is broken, and we are delivered: our help standeth in the name of the Lord, who hath made heaven and earth.}{Alleluia, alleluia. ℣. Let the righteous be glad and rejoice before God: let them also be merry and joyful. Alleluia.}

\begin{rubric}
{In Septuagesimatide or Lent, replacing the Alleluia:}
\end{rubric}\par\noindent
\tract{They that sow in tears, shall reap in joy. ℣. He that now goeth on his way weeping, and beareth forth good seed. ℣. Shall doubtless come again with joy, and bring his sheaves with him.}

\readingcitation{Gospel}{Luke 6:17}
\lett{A}{t that time:} Jesus came down from the mountain, and stood in the plain, and the company of his disciples, and a great multitude of people out of all Jud{\ae}a and Jerusalem, and from the sea coast of Tyre and Sidon, which came to hear him, and to be healed of their diseases; and they that were vexed with unclean spirits: and they were healed. And the whole multitude sought to touch him: for there went virtue out of him, and healed them all.

And he lifted up his eyes on his disciples, and said, Blessed be ye poor: for yours is the kingdom of God. Blessed are ye that hunger now: for ye shall be filled. Blessed are ye that weep now: for ye shall laugh. Blessed are ye, when men shall hate you, and when they shall separate you from their company, and shall reproach you, and cast out your name as evil, for the Son of man’s sake. Rejoice ye in that day, and leap for joy: for, behold, your reward is great in heaven.

\begin{rubric}
	After Septuagesima, \emph{Alleluia} at the end of the following Offertory is omitted.
\end{rubric}

\offertory{Let the saints be joyful with glory, let them rejoice in their beds: let the praises of God be in their mouth, alleluia.}

\communion{I say unto you, my friends: Be not afraid of them that persecute you.}


\subby{III: Common of Many Martyrs out of Eastertide}\label{CommonMartyrsIII}
\fancyhead[RO,LE]{\textit{Martyrs III}}
\begin{rubric}
	Daily Office propers as in the First Common of Many Martyrs (p. \pageref{CommonMartyrsI}).
\end{rubric}

\introit
\lett{B}{ut} the salvation of the righteous cometh of the Lord: who is also their strength in the time of trouble. \textit{Ps.} Fret not thyself because of the ungodly: neither be thou envious against the evil doers.

\collect
\lett{O}{God,} who makest us glad with the yearly solemnity of thy holy Martyrs \textit{N. and N.}: mercifully grant; that as we do rejoice in their merits, so we may be enkindled by their example. Through.

\readingcitation{Epistle}{Hebrews 10:32}\label{Hebrews1032}
\lett{B}{rethren:} Call to remembrance the former days, in which, after ye were illuminated, ye endured a great fight of afflictions; partly, whilst ye were made a gazingstock both by reproaches and afflictions; and partly, whilst ye became companions of them that were so used. For ye had compassion of me in my bonds, and took joyfully the spoiling of your goods, knowing in yourselves that ye have in heaven a better and an enduring substance. Cast not away therefore your confidence, which hath great recompence of reward. For ye have need of patience, that, after ye have done the will of God, ye might receive the promise. For yet a little while, and he that shall come will come, and will not tarry. Now the just shall live by faith.

\gradall{The righteous cry, and the Lord heareth them: and delivereth them out of all their troubles. ℣. The Lord is nigh unto them that are of a contrite heart: and will save such as be of an humble spirit.}{Alleluia, alleluia. ℣. The noble army of Martyrs praise thee, O Lord. Alleluia.}

\begin{rubric}
{In Septuagesimatide or Lent, replacing the Alleluia:}
\end{rubric}\par\noindent
\tract{They that sow in tears, shall reap in joy. ℣. He that now goeth on his way weeping, and beareth forth good seed. ℣. Shall doubtless come again with joy, and bring his sheaves with him.}

\readingcitation{Gospel}{Luke 12:1}
\lett{A}{t that time:} Jesus said unto his disciples: Beware ye of the leaven of the Pharisees, which is hypocrisy. For there is nothing covered, that shall not be revealed; neither hid, that shall not be known. Therefore whatsoever ye have spoken in darkness shall be heard in the light; and that which ye have spoken in the ear in closets shall be proclaimed upon the housetops. And I say unto you my friends, Be not afraid of them that kill the body, and after that have no more that they can do. But I will forewarn you whom ye shall fear: Fear him, which after he hath killed hath power to cast into hell; yea, I say unto you, Fear him. Are not five sparrows sold for two farthings, and not one of them is forgotten before God? But even the very hairs of your head are all numbered. Fear not therefore: ye are of more value than many sparrows. Also I say unto you, Whosoever shall confess me before men, him shall the Son of man also confess before the angels of God.

\begin{rubric}
	After Spetuagesima, \emph{Alleluia} at the end of the following Offertory is omitted.
\end{rubric}

\offertory{The souls of the just are in the hand of God, and there shall no torment of malice touch them: in the sight of the unwise they seemed to die: but they are in peace, alleluia.}

\secret
\lett{W}{e} beseech thee mercifully to accept these our oblations: that at the intercession of thy holy Martyrs \textit{N. and N.} we may be defended against all adversities. Through.

\communion{What I tell you in darkness, that speak ye in light, saith the Lord: and what ye hear in the ear, that preach ye upon the housetops.}

\postcommunion
\lett{M}{ay} this communion, O Lord, cleanse us from guilt: and at the intercession of thy holy Martyrs \textit{N. and N.} make us partakers of thy heavenly healing. Through.

\begin{rubric}
	The following Epistles and Gospel may be used instead of the ones provided above.
\end{rubric}

\readingcitation{1. Epistle}{Wisdom 10:17}\label{Wisdom1017}
%RV:
\lett{W}{isdom} rendered unto holy men a reward of their toils; She guided them along a marvellous way, And became unto them a covering in the day-time, And a flame of stars through the night. She brought them over the Red sea, And led them through much water; But their enemies she drowned, And out of the bottom of the deep she cast them up. Therefore the righteous spoiled the ungodly; And they sang praise to thy holy name, O Lord, And extolled with one accord thy hand that fought for them, O Lord our God.

%\lett{G}{od} rendered to the righteous a reward of their labours, guided them in a marvellous way, and was unto them for a cover by day, and a light of stars in the night season; Brought them through the Red sea, and led them through much water: But he drowned their enemies, and cast them up out of the bottom of the deep. Therefore the righteous spoiled the ungodly, and praised thy holy name, O Lord, and magnified with one accord thine hand, that fought for them, O Lord our God.

\readingcitation{2. Epistle}{Romans 5:1}\label{Romans51}
\lett{B}{rethren:} Being justified by faith, we have peace with God through our Lord Jesus Christ: by whom also we have access by faith into this grace wherein we stand, and rejoice in hope of the glory of God. And not only so, but we glory in tribulations also: knowing that tribulation worketh patience; and patience, experience; and experience, hope: and hope maketh not ashamed; because the love of God is shed abroad in our hearts by the Holy Ghost which is given unto us.

\readingcitation{3. Epistle}{Romans 8:18}\label{Romans818}
\lett{B}{rethren:} I reckon that the sufferings of this present time are not worthy to be compared with the glory which shall be revealed in us. For the earnest expectation of the creature waiteth for the manifestation of the sons of God. For the creature was made subject to vanity, not willingly, but by reason of him who hath subjected the same in hope, because the creature itself also shall be delivered from the bondage of corruption into the glorious liberty of the children of God. For we know that the whole creation groaneth and travaileth in pain together until now. And not only they, but ourselves also, which have the firstfruits of the Spirit, even we ourselves groan within ourselves, waiting for the adoption, to wit, the redemption of our body: in Christ Jesus our Lord.

\readingcitation{4. Epistle}{2 Corinthians 6:4}\label{2Corinthians64}
\lett{B}{rethren:} Let us approve ourselves as the ministers of God, in much patience, in afflictions, in necessities, in distresses, in stripes, in imprisonments, in tumults, in labours, in watchings, in fastings; by pureness, by knowledge, by longsuffering, by kindness, by the Holy Ghost, by love unfeigned, by the word of truth, by the power of God, by the armour of righteousness on the right hand and on the left, by honour and dishonour, by evil report and good report: as deceivers, and yet true; as unknown, and yet well known; as dying, and, behold, we live; as chastened, and not killed; as sorrowful, yet alway rejoicing; as poor, yet making many rich; as having nothing, and yet possessing all things.

\readingcitation{5. Epistle}{Hebrews 11:33}\label{Hebrews1133}
\lett{B}{rethren:} The saints through faith subdued kingdoms, wrought righteousness, obtained promises, stopped the mouths of lions, quenched the violence of fire, escaped the edge of the sword, out of weakness were made strong, waxed valiant in fight, turned to flight the armies of the aliens: women received their dead raised to life again: and others were tortured, not accepting deliverance, that they might obtain a better resurrection: and others had trial of cruel mockings and scourgings, yea, moreover of bonds and imprisonment: they were stoned, they were sawn asunder, were tempted, were slain with the sword: they wandered about in sheepskins and goat-skins, being destitute, afflicted, tormented: (of whom the world was not worthy;) they wandered in deserts, and 1a mountains, and in dens and caves of the earth. And these all obtained a good report through faith, and were found in Christ Jesus our Lord.

\readingcitation{6. Epistle}{Revelation 7:13}\label{Revelation713}
\lett{I}{n those days:} One of the elders answered, saying unto me, What are these which are arrayed in white robes? and whence came they? And I said unto him, Sir, thou knowest. And he said to me, These are they which came out of great tribulation, and have washed their robes, and made them white in the blood of the Lamb. Therefore are they before the throne of God, and serve him day and night in his temple: and he that sitteth on the throne shall dwell among them. They shall hunger no more, neither thirst any more; neither shall the sun light on them, nor any heat. For the Lamb which is in the midst of the throne shall feed them, and shall lead them unto living fountains of waters: and God shall wipe away all tears from their eyes.

\readingcitation{1. Gospel}{Matthew 24:3}\label{Matthew243}
\lett{A}{t that time:} As Jesus sat upon the mount of Olives, the disciples came unto him privately, saying: Tell us, when shall these things be? and what shall be the sign of thy coming, and of the end of the world? And Jesus answered and said unto them, Take heed that no man deceive you. For many shall come in my name, saying, I am Christ; and shall deceive many. And ye shall hear of wars and rumours of wars: see that ye be not troubled: for all these things must come to pass, but the end is not yet. For nation shall rise against nation, and kingdom against kingdom: and there shall be famines, and pestilences, and earthquakes, in divers places. All these are the beginning of sorrows.

Then shall they deliver you up to be afflicted, and shall kill you: and ye shall be hated of all nations for my name’s sake. And then shall many be offended, and shall betray one another, and shall hate one another. And many false prophets shall rise, and shall deceive many. And because iniquity shall abound, the love of many shall wax cold. But he that shall endure unto the end, the same shall be saved.

\readingcitation{2. Gospel}{Matthew 5:1}
\lett{A}{t that time:} Jesus, seeing the multitudes, went up into a mountain, and when he was set, his disciples came unto him, and he opened mouth, and taught them, saying: Blessed are the poor in spirit: for theirs is the kingdom of heaven. Blessed are they that mourn: for they shall be comforted. Blessed are the meek: for they shall inherit the earth. Blessed are they which do hunger and thirst after righteousness: for they shall be filled. Blessed are the merciful: for they shall obtain mercy. Blessed are the pure in heart: for they shall see God. Blessed are the peacemakers: for they shall be called the children of God. Blessed are they which are persecuted for righteousness’ sake: for theirs is the kingdom of heaven. Blessed are ye, when men shall revile you, and persecute you, and shall say all manner of evil against you falsely, for my sake. Rejoice, and be exceeding glad: for great is your reward in heaven.

\readingcitation{3. Gospel}{Matthew 11:25}
\lett{A}{t that time:} Jesus answered and said: I thank thee, O Father, Lord of heaven and earth, because thou hast hid these things from the wise and prudent, and hast revealed them unto babes. Even so, Father: for so it seemed good in thy sight. All things are delivered unto me of my Father: and no man knoweth the Son, but the Father; neither knoweth any man the Father, save the Son, and he to whomsoever the Son will reveal him.

Come unto me, all ye that labour and are heavy laden, and I will give you rest. Take my yoke upon you, and learn of me; for I am meek and lowly in he art: and ye shall find rest unto your souls. For my yoke is easy, and my burden is light.

\readingcitation{4. Gospel}{Luke 11:47}\label{Luke1147}
\lett{A}{t that time:} Jesus said unto the scribes and Pharisees: Woe unto you! for ye build the sepulchres of the prophets, and your fathers killed them. Truly ye bear witness that ye allow the deeds of your fathers: for they indeed killed them, and ye build their sepulchres. Therefore also said the wisdom of God, I will send them prophets and apostles, and some of them they shall slay and persecute: that the blood of all the prophets, which was shed from the foundation of the world, may be required of this generation; from the blood of Abel unto the blood of Zacharias, which perished between the altar and the temple: verily I say unto you, It shall be required of this generation.

\readingcitation{5. Gospel}{Luke 10:16}\label{Luke1016}
\lett{A}{t that time:} Jesus said unto his disciples: He that heareth you heareth me; and he that despiseth you despiseth me; and he that despiseth me despiseth him that sent me. And the seventy returned again with joy, saying, Lord, even the devils are subject unto us through thy name. And he said unto them, I beheld Satan as lightning fall from heaven. Behold, I give unto you power to tread on serpents and scorpions, and over all the power of the enemy: and nothing shall by any means hurt you. Notwithstanding in this rejoice not, that the spirits are subject unto you; but rather rejoice, because your names are written in heaven.


\subby{Common of One Martyr in Eastertide}\label{CommonMartyrEaster}
\fancyhead[RO,LE]{\textit{Martyr in Easter}}
\begin{rubric}
	Daily Office Antiphons \& Versicles as in the Common of Apostles in Eastertide (p. \pageref{CommonApostles}), with the following Hymns.
\end{rubric}

\subbysub{I Evensong}

\gregorioscore{resources/gabc/Common/MartyrEvensongEaster.gabc}

    \par\noindent
	\centerline{\rule{0.25\textwidth}{0.4pt}}
	\par\noindent
\begin{rubric}
	From Low Sunday, exclusive, until Ascension Thursday, exclusive:
\end{rubric}

\gregorioscore{resources/gabc/Common/MartyrEvensongEasterEnding1.gabc}

\begin{rubric}
	From Ascension Thursday, inclusive, until Whitsunday, inclusive:
\end{rubric}

\gregorioscore{resources/gabc/Common/MartyrEvensongEasterEnding2.gabc}

    \par\noindent
	\centerline{\rule{0.25\textwidth}{0.4pt}}
	\par


\subbysub{Mattins}

\begin{rubric}
	The Invitatory Hymn is of I Evensong.
\end{rubric}

\officehymn

\gregorioscore{resources/gabc/Common/MartyrMattinsEaster.gabc}

    \par\noindent
	\centerline{\rule{0.25\textwidth}{0.4pt}}
	\par\noindent
\begin{rubric}
	From Low Sunday, exclusive, until Ascension Thursday, exclusive:
\end{rubric}

\gregorioscore{resources/gabc/Common/MartyrMattinsEasterEnding1.gabc}

\begin{rubric}
	From Ascension Thursday, inclusive, until Whitsunday, inclusive:
\end{rubric}

\gregorioscore{resources/gabc/Common/MartyrMattinsEasterEnding2.gabc}

    \par\noindent
	\centerline{\rule{0.25\textwidth}{0.4pt}}
	\par

\subbysub{II Evensong}

\begin{rubric}
	The Office Hymn is of I Evensong.
\end{rubric}

\introit
\lett{T}{hou} hast hidden me, O God, from the gathering together of the froward, alleluia: and from the insurrection of the wicked doers, alleluia, alleluia. \textit{Ps.} Hear my voice, O God, in my prayer: preserve my lite from fear of the enemy.

\collect
\begin{inhead}
	Martyr not Bishop
\end{inhead}
\lett{G}{rant,} we beseech thee, almighty God: that we, who devoutly celebrate the birthday of blessed \textit{N.}, thy Martyr; may by his intercession be stablished in the love of thy name. Through.
\begin{inhead}
	or
\end{inhead}
\lett{G}{rant,} we beseech thee, almighty God: that at the intercession of blessed \textit{N.} thy Martyr, we may both be delivered from all adversities which may happen to the body, and from all evil thoughts which may assault and hurt the soul. Through.

%MANUAL ADJUSTMENT:
\clearpage
\begin{inhead}
	Martyr Bishop
\end{inhead}
\lett{A}{lmighty} God, mercifully look upon our infirmities: that whereas we are oppressed by the burden of our sins, the glorious intercession of blessed \textit{N.} thy Martyr and Bishop may be our succour and defence. Through.
\begin{inhead}
	or
\end{inhead}
\lett{O}{God,} who makest us glad with the yearly solemnity of blessed \textit{N.}, thy Martyr and Bishop: mercifully grant; that as we now devoutly celebrate his birthday, so we may likewise rejoice in his protection. Through.

\readingcitation{Epistle}{Wisdom 5:1}
%RV:
\lett{T}{he} righteous man shall stand in great boldness Before the face of them that afflicted him, And them that make his labours of no account. When they see it, they shall be troubled with terrible fear, And shall be amazed at the marvel of God's salvation. They shall say within themselves repenting, And for distress of spirit shall they groan, This was he whom aforetime we had in derision, And made a parable of reproach: We fools accounted his life madness, And his end without honour: How was he numbered among sons of God? And how is his lot among saints?

%\lett{T}{he} righteous man shall stand in great boldness before the face of such as have afflicted him, and made no account of his labours. When they see it, they shall be troubled with terrible fear, and shall be amazed at the strangeness of his salvation, so far beyond all that they looked for. And they repenting and groaning for anguish of spirit shall say within themselves, This was he, whom we had sometimes in derision, and a proverb of reproach: We fools accounted his life madness, and his end to be without honour: How is he numbered among the children of God, and his lot is among the saints!

\begin{rubric}
	The Epistle from the Third Common of One Martyr not a Bishop out of Eastertide (2 Timothy 2:8) may be said instead.
\end{rubric}

\alleluia{Alleluia, alleluia. ℣. O Lord, the very heavens shall praise thy wondrous works: and thy truth in the congregation of the saints. Alleluia. ℣. Thou hast set, O Lord, a crown of pure gold upon his head. Alleluia.}

\readingcitation{Gospel}{John 15:1}
\lett{A}{t that time:} Jesus said unto his disciples: I am the true vine, and my Father is the husbandman. Every branch in me that beareth not fruit he taketh away: and every branch that beareth fruit, he purgeth it, that it may bring forth more fruit. Now ye are clean through the word which I have spoken unto you. Abide in me, and I in you. As the branch cannot bear fruit of itself, except it abide in the vine; no more can ye, except ye abide in me. I am the vine, ye are the branches: He that abideth in me, and I in him, the same bringeth forth much fruit: for without me ye can do nothing. If a man abide not in me, he is cast forth as a branch, and is withered; and men gather them, and cast them into the fire, and they are burned. If ye abide in me, and my words abide in you, ye shall ask what ye will, and it shall be done unto you.

\offertory{O Lord, the very heavens shall praise thy wondrous works: and thy truth in the congregation of the saints, alleluia, alleluia.}

\secret
\begin{inhead}
	Martyr not Bishop
\end{inhead}
\lett{W}{e} beseech thee, O Lord, accept our prayers and oblations: and graciously hearken unto us whom thou dost cleanse by thy heavenly mysteries. Through.
\begin{inhead}
	or
\end{inhead}
\lett{G}{rant,} O Lord, that our bounden service may be acceptable in thy sight: that these our oblations may, by the prayers of him on whose solemnity they are offered, be made profitable unto our salvation. Through.

\begin{inhead}
	Martyr Bishop
\end{inhead}
\lett{W}{e} beseech thee, O Lord, mercifully to accept this our sacrifice, which we offer unto thee pleading the merits of blessed \textit{N.}, thy Martyr and Bishop: that the same may effectually avail for our perpetual succour. Through.

\begin{inhead}
	or
\end{inhead}
\lett{S}{anctify,} O Lord, the gifts which we dedicate to thee: that at the intercession of blessed \textit{N.}, thy Martyr and Bishop, they may obtain for us thy gracious favour. Through,

\communion{The righteous shall rejoice in the Lord, and put his trust in him: and all they that are true of heart shall be glad, alleluia, alleluia.}

\postcommunion
\begin{inhead}
	Martyr not Bishop
\end{inhead}
\lett{G}{rant,} we beseech thee, O Lord our God: that like as we in this life do gladly honour the memory of thy Saints; so we may rejoice to behold them for ever. Through.
\begin{inhead}
	or
\end{inhead}
\lett{W}{e} beseech thee, O Lord our God: that like as we, whom thou hast refreshed by the partaking of thy sacred gift, do offer unto thee our worship; so by the intercession of blessed \textit{N.}, thy Martyr, we may perceive the benefit of the same. Through.

\begin{inhead}
	Martyr Bishop
\end{inhead}
\lett{W}{e} beseech thee, O Lord our God, that like as we, whom thou hast refreshed by the partaking of thy sacred gift, do offer unto thee our worship: so by the intercession of blessed \textit{N.}, thy Martyr and Bishop, we may perceive the benefit of the same. Through.
\begin{inhead}
	or
\end{inhead}
\lett{M}{ay} this communion, O Lord, cleanse us from guilt: and, at the intercession of blessed \textit{N.}, thy Martyr and Bishop, make us partakers of thy heavenly healing. Through,

\begin{rubric}
	The first Epistle and the Gospel of the above Mass may be said also in the following Mass, and those in the following may be said in the above, provided a special Epistle or Gospel is not appointed.
\end{rubric}

%MANUAL ADJUSTMENT:
\clearpage
\subby{Common of Many Martyrs in Eastertide}\label{CommonMartyrsEaster}
\fancyhead[RO,LE]{\textit{Martyrs in Easter}}

\begin{rubric}
	Daily Office Antiphons \& Versicles as in the Common of Apostles in Eastertide (p. \pageref{CommonApostles}), with the following Hymns.
\end{rubric}

\subbysub{I Evensong}

\gregorioscore{resources/gabc/Common/MartyrsEvensongEaster.gabc}

    \par\noindent
	\centerline{\rule{0.25\textwidth}{0.4pt}}
	\par\noindent
\begin{rubric}
	From Low Sunday, exclusive, until Ascension Thursday, exclusive:
\end{rubric}

\gregorioscore{resources/gabc/Common/MartyrsEvensongEasterEnding1.gabc}

\begin{rubric}
	From Ascension Thursday, inclusive, until Whitsunday, inclusive:
\end{rubric}

\gregorioscore{resources/gabc/Common/MartyrsEvensongEasterEnding2.gabc}

    \par\noindent
	\centerline{\rule{0.25\textwidth}{0.4pt}}
	\par

\subbysub{Mattins}

\begin{rubric}
	The Invitatory Hymn is from the First Common of Many Martyrs (p. \pageref{CommonMartyrsI}).
\end{rubric}

\begin{rubric}
	The Office Hymn is as in I Evensong.
\end{rubric}

\subbysub{II Evensong}

\begin{rubric}
	The Hymn is as in I Evensong.
\end{rubric}

\introit
\lett{T}{hy} saints give thanks unto thee, O Lord: they shew the glory of thy kingdom, alleluia, alleluia. \textit{Ps.} I will magnify thee, O God, my King: and I will praise thy name for ever and ever.

\collect
\begin{inhead}
	Many Martyrs not Bishops
\end{inhead}
\lett{O}{God,} who vouchsafest unto us to celebrate the birthday of thy holy Martyrs \textit{N. and N.}: grant that we may rejoice in the everlasting felicity of their fellowship. Through.
\begin{inhead}
	or
\end{inhead}
\lett{O}{God,} who makest us glad with the yearly solemnity of thy holy Martyrs \textit{N. and N.}: mercifully grant that as we rejoice in their merits, so we may be enkindled by their example. Through.

\begin{inhead}
	Martyrs Bishops
\end{inhead}
\lett{P}{rotect} us, O Lord, we beseech thee, who observe the feast of thy blessed Martyrs and Bishops \textit{N. and N.}: and grant that by their meritorious supplication we may ever find favour in thy sight. Through.

\readingcitation{Epistle}{1 Peter 1:3}
\lett{B}{lessed} be the God and Father of our Lord Jesus Christ, which according to his abundant mercy hath begotten us again unto a lively hope by the resurrection of Jesus Christ from the dead, to an inheritance incorruptible, and undefiled, and that fadeth not away, reserved in heaven for you, who are kept by the power of God through faith unto salvation ready to be revealed in the last time. Wherein ye greatly rejoice, though now for a season, if need be, ye are in heaviness through manifold temptations: that the trial of your faith, being much more precious than of gold that perisheth, though it be tried with fire, might be found unto praise and honour and glory at the appearing of Jesus Christ our Lord.

\alleluia{Alleluia, alleluia. ℣. Thy saints, O Lord, shall grow as the lily, and as the odour of balsam shall they be before thee. Alleluia. ℣. Right dear in the sight of the Lord is the death of his saints. Alleluia.}

\readingcitation{Gospel}{John 15:5}
\lett{A}{t that time:} Jesus said unto his disciples: I am the vine, ye are the branches: He that abideth in me, and I in him, the same bringeth forth much fruit: for without me ye can do nothing. If a man abide not in me, he is cast forth as a branch, and is withered; and men gather them, and cast them into the fire, and they are burned. If ye abide in me, and my words abide in you, ye shall ask what ye will, and it shall be done unto you. Herein is my Father glorified, that ye bear much fruit; so shall ye be my disciples. As the Father hath loved me, so have I loved you: continue ye in my love. If ye keep my commandments, ye shall abide in my love; even as I have kept my Father’s commandments, and abide in his love. These things have I spoken unto you, that my joy might remain in you, and that your joy might be full.

\offertory{Be glad, O ye righteous, and rejoice in the Lord: and be joyful, all ye that are true of heart, alleluia, alleluia.}

\secret
\begin{inhead}
	Martyrs not Bishops
\end{inhead}
\lett{W}{e} beseech thee, O Lord, that the gifts which we offer unto thee of our bounden duty and service may be acceptable unto thee for the honour of thy Just ones: and by thy mercy profitable unto us for our salvation. Through.

\begin{inhead}
	or
\end{inhead}

\lett{W}{e} beseech thee, O Lord, mercifully to accept these our oblations: that at the intercession of thy holy Martyrs \textit{N. and N.} we may be defended against all adversities. Through.

\begin{inhead}
	Martyrs Bishops
\end{inhead}
\lett{A}{ssist} us mercifully, O Lord, in these our supplications, which we make before thee in remembrance of thy Saints: that we who trust not in our own righteousness may be succoured by the merits of them that have found favour in thy sight. Through.

\communion{Rejoice in the Lord, O ye righteous, alleluia: for it becometh well the just to be thankful, alleluia.}

\postcommunion
\begin{inhead}
	Martyrs not Bishops
\end{inhead}
\lett{G}{rant} to us, we beseech thee, O Lord: at the intercession of thy holy Martyrs \textit{N. and N.}; that those things which we touch with the mouth, we may receive in purity of heart. Through.

\begin{inhead}
	or
\end{inhead}

\lett{M}{ay} this communion, O Lord, cleanse us from guilt: and, at the intercession of thy holy Martyrs \textit{N. and N.}, make us partakers of thy heavenly healing. Through.

\begin{inhead}
	Martyrs Bishops
\end{inhead}
\lett{O}{Lord,} who hast fulfilled us with saving mysteries: we beseech thee that we may be aided by the prayers of those whose festival we celebrate. Through.

\begin{rubric}
	The following Epistle and Gospel may be used instead of the ones provided above.
\end{rubric}

\readingcitation{Another Epistle}{Revelation 19:1}\label{Revelation191}
\lett{I}{n those days:} After these things I, John, heard a great voice of much people in heaven, saying, Alleluia; Salvation, and glory, and honour, and power, unto the Lord our God: for true and righteous are his judgments: for he hath judged the great whore, which did corrupt the earth with her fornication, and hath avenged the blood of his servants at her hand. And again they said, Alleluia. And her smoke rose up for ever and ever. And the four and twenty elders and the four beasts fell down and worshipped God that sat on the throne, saying, Amen; Alleluia.

And a voice came out of the throne, saying, Praise our God, all ye his servants, and ye that fear him, both small and great. And I heard as it were the voice of a great multitude, and as the voice of many waters, and as the voice of mighty thunderings, saying, Alleluia: for the Lord God omnipotent reigneth. Let us be glad and rejoice, and give honour to him: for the marriage of the Lamb is come, and his wife hath made herself ready. And to her was granted that she should be arrayed in fine linen, clean and white: for the fine linen is the righteousness of saints.

And he saith unto me, Write, Blessed are they which are called unto the marriage supper of the Lamb. 

\readingcitation{Another Gospel}{John 16:20}\label{John1620}
\lett{A}{t that time:} Jesus said unto his disciples: Verily, verily, I say unto you, That ye shall weep and lament, but the world shall rejoice: and ye shall be sorrowful, but your sorrow shall be turned into joy. A woman when she is in travail hath sorrow, because her hour is come: but as soon as she is delivered of the child, she remembereth no more the anguish, for joy that a man is born into the world. And ye now therefore have sorrow: but I will see you again, and your heart shall rejoice, and your joy no man taketh from you.

\begin{rubric}
	In the Common of Confessors, and Virgins, and in other Masses of Eastertide, all is said as in the rest of the year, a double Alleluia being added in the Introit before the Psalm, and at the end of the Offertory and Communion one Alleluia, where it is not given; and the Gradual being omitted, two ℣℣. are said with four Alleluias, as is set down in their places.
\end{rubric}


\subby{I: Common of a Confessor Bishop}\label{CommonConfessorBishopI}
\fancyhead[RO,LE]{\textit{Confessor Bishop I}}

\subbysub{I Evensong}

\gregorioscore{resources/gabc/Common/ConfessorBishopEvensong.gabc}

℣. The Lord loved him and adorned him.

℟. He clothed him with a robe of glory.

\antiphon{Mag.}{O thou Priest and Bishop, {\dag} thou workers of mighty works, thou good shepherd of the people, pray unto the Lord for us.}

\subbysub{Mattins}

\begin{rubric}
	The Invitatory Hymn is from I Evensong.
\end{rubric}

\officehymn

\gregorioscore{resources/gabc/Common/ConfessorBishopMattins.gabc}

℣. The Lord guided the righteous in right paths.

℟. And shewed him the kingdom of God.

\antiphon{Ben.}{Well done, {\dag} thou good and faithful servant: thou hast been faithful over a few things, I will make thee ruler over many things, saith the Lord.}

\subbysub{II Evensong}

\begin{rubric}
	The Office Hymn is of I Evensong, and the Versicle is of Mattins, with the following Antiphon.
\end{rubric}

\antiphon{Mag.}{I will liken him {\dag} unto a wise man, which built his house upon a rock.}

\introit
\lett{T}{he} Lord hath established a covenant of peace with him, and made him a prince: that he should have the dignity of the priesthood for ever (Alleluia, alleluia). \textit{Ps.} Lord, remember David: and all his trouble.

\collect
\lett{G}{rant,} we beseech thee, almighty God: that the venerable solemnity of blessed \textit{N.}, thy Confessor and Bishop may increase our devotion and set forward our salvation. Through.

\readingcitation{Epistle}{Ecclesiasticus 44:16}
\lett{E}{noch} pleased the Lord, and was translated, Being an example of repentance to all generations. Noah was found perfect and righteous; In the season of wrath he was taken in exchange for the world; Therefore was there left a remnant unto the earth, When the flood came. Everlasting covenants were made with him, That all flesh should no more be blotted out by a flood. Abraham was a great father of a multitude of nations; And there was none found like him in glory; Who kept the law of the Most High, And was taken into covenant with him: In his flesh he established the covenant; And when he was proved, he was found faithful. Therefore he assured him by an oath, That the nations should be blessed in his seed; That he would multiply him as the dust of the earth, And exalt his seed as the stars, And cause them to inherit from sea to sea, And from the River unto the utmost part of the earth. In Isaac also did he establish likewise, for Abraham his father’s sake, The blessing of all men, and the covenant: And he made it rest upon the head of Jacob; He acknowledged him in his blessings, And gave to him by inheritance, And divided his portions; Among twelve tribes did he part them. And he brought out of him a man of mercy, Which found favour in the sight of all flesh.

%Traditional Text:
%\lett{B}{ehold,} a great priest, who in his days pleased God, and was found just; in the time of wrath he was made a means of reconciliation. There was none like unto him who kept the law of the Most High. Therefore he assured him by an oath that he would increase him among his people. He established him with the blessing of all men and the covenant and made it rest upon his head. He acknowledged him in his blessing: he preserved for him his mercy: so that he found favour in the sight of the Lord. He made him glorious in the sight of kings, and gave him a crown of glory. An everlasting covenant he made with him, and gave him a great priesthood, and clothed him with a robe of glory. To execute the office of the priesthood, and bless the people in his name, and to offer him a worthy incense and a sweet savour.

\gradall{Behold, a great priest, who in his days pleased God. ℣. There was none found like unto him who kept the law of the Most High.}{Alleluia, alleluia. ℣. Thou art a Priest for ever after the order of Melchisedech. Alleluia.}

\begin{rubric}
{In Septuagesimatide or Lent, replacing the Alleluia:}
\end{rubric}\par\noindent
\tract{Blessed is the man that feareth the Lord: he hath great delight in his commandments. ℣. His seed shall be mighty upon earth: the generation of the faithful shall be blessed. ℣. Riches and plenteousness shall be in his house: and his righteousness endureth for ever.}

\begin{rubric}
{In Eastertide, replacing the Lesser Alleluia:}
\end{rubric}\par\noindent
\alleluia{Alleluia, alleluia. ℣. Thou art a priest for ever after the order of Melchisedech, alleluia. ℣. This is a priest whom the Lord hath crowned. Alleluia.}

\readingcitation{Gospel}{Matthew 25:14}
\lett{A}{t that time:} Jesus spake this parable unto his disciples: A man travelling into a far country, called his own servants, and delivered unto them his goods. And unto one he gave five talents, to another two, and to another one; to every man according to his several ability; and straightway took his journey. Then he that had received the five talents went and traded with the same, and made them other five talents. And likewise he that had received two, he also gained other two. But he that had received one went and digged in the earth, and hid his lord’s money. After a long time the lord of those servants cometh, and reckoneth with them. And so he that had received five talents came and brought other five talents, saying, Lord, thou deliveredst unto me five talents: behold, I have gained beside them five talents more. His lord said unto him, Well done, thou good and faithful servant: thou hast been faithful over a few things, I will make thee ruler over many things: enter thou into the joy of thy lord. He also that had received two talents came and said, Lord, thou deliveredst unto me two talents: behold, I have gained two other talents beside them. His lord said unto him, Well done, good and faithful servant; thou hast been faithful over a few things, I will make thee ruler over many things: enter thou into the joy of thy lord.

\offertory{I have found David my servant, with my holy oil have I anointed him; my hand shall hold him fast, and my arm shall strengthen him (Alleluia).}

\secret
\lett{W}{e} beseech thee, O Lord, that we, remembering with gladness the merits of thy Saints, may in all places feel the succour of their intercession. Through.

\communion{A faithful and wise steward, whom the lord hath made ruler over his household: to give them their portion of meat in due season (Alleluia).}

\postcommunion
\lett{G}{rant,} we beseech thee, almighty God: that we, shewing forth our thankfulness for the gifts which we have received, may, at the intercession of blessed \textit{N.}, thy Confessor and Bishop, obtain yet more abundant mercies. Through.


\subby{II: Common of a Confessor Bishop}\label{CommonConfessorBishopII}
\fancyhead[RO,LE]{\textit{Confessor Bishop II}}
\begin{rubric}
	Daily Office propers are as in the First Common of a Confessor Bishop (p. \pageref{CommonConfessorBishopI}).
\end{rubric}

\introit
\lett{L}{et} thy priests, O Lord, be clothed with righteousness, and let thy saints sing with joyfulness: for thy servant David's sake, turn not away the presence of thine Anointed (Alleluia, alleluia). \textit{Ps.} Lord, remember David: and all his trouble.

\collect
\lett{W}{e} beseech thee, O Lord, graciously to hear the prayers which we offer unto thee on the solemnity of blessed \textit{N.}, thy Confessor and Bishop: that, like as he was found worthy to do thee faithful service, so by his merits and intercession we may be absolved from all our sins. Through.

\readingcitation{Epistle}{Hebrews 7:23}
\lett{B}{rethren:} They were many priests, because they were not suffered to continue by reason of death: but this man, because he continueth ever, hath an unchangeable priesthood. Wherefore he is able also to save them to the uttermost that come unto God by him, seeing he ever liveth to make intercession for them.

For such an high priest became us, who is holy, harmless, undefiled, separate from sinners, and made higher than the heavens; who needeth not daily, as those high priests, to offer up sacrifice, first for his own sins, and then for the people’s: for this he did once, when he offered up himself, Jesus Christ, our Lord.

\gradall{I will deck her priests with health: and her saints shall rejoice and sing. ℣. There shall I make the horn of David to flourish: I have ordained a lantern for mine Anointed.}{Alleluia, alleluia. ℣. The Lord sware, and will not repent: Thou art a priest for ever after the order of Melchisedech. Alleluia.}

\begin{rubric}
{In Septuagesimatide or Lent, replacing the Alleluia:}
\end{rubric}\par\noindent
\tract{Blessed is the man that feareth the Lord: he hath great delight in his commandments. ℣. His seed shall be mighty upon earth: the generation of the faithful shall be blessed. ℣. Riches and plenteousness shall be in his house: and his righteousness endureth for ever.}

\begin{rubric}
{In Eastertide, replacing the Lesser Alleluia:}
\end{rubric}\par\noindent
\alleluia{Alleluia, alleluia. ℣. The Lord sware, and will not repent: Thou art a priest for ever after the order of Melchisedech. Alleluia. ℣. The Lord loved him, and adorned him: and clothed him with a robe of glory. Alleluia.}

\readingcitation{Gospel}{Matthew 24:42}
\lett{A}{t that time:} Jesus said unto his disciples: Watch, for ye know not what hour your Lord doth come. But know this, that if the goodman of the house had known in what watch the thief would come, he would have watched, and would not have suffered his house to be broken up. Therefore be ye also ready: for in such an hour as ye think not the Son of man cometh. Who then is a faithful and wise servant, whom his lord hath made ruler over his household, to give them meat in due season? Blessed is that servant, whom his lord when he cometh shall find so doing. Verily I say unto you, That he shall make him ruler over all his goods.

\offertory{My truth and my mercy shall be with him: and in my name shall his horn be exalted (Alleluia).}

\secret
\lett{W}{e} beseech thee, O Lord, that our devout observance of the yearly solemnity of blessed \textit{N.}, thy Confessor and Bishop, may render us acceptable unto thy loving-kindness: that this service of propitiation, which we duly offer, may be profitable unto him for the reward of blessedness, and obtain for us the gifts of thy grace. Through.

\communion{Blessed is the servant, whom the lord when he cometh shall find watching: verily I say unto you, He shall make him ruler over all his goods (Alleluia).}

\postcommunion
\lett{O}{God,} who rewardest the souls of them that put their trust in thee: vouchsafe; that we who keep the solemn festival of blessed \textit{N.}, thy Confessor and Bishop, may by his prayers obtain thy merciful pardon. Through.

\begin{rubric}
	The following Epistles and Gospels may be used instead of the ones provided above.
\end{rubric}

\readingcitation{1. Epistle}{Hebrews 5:1}
\lett{B}{rethren:} Every high priest taken from among men is ordained for men in things pertaining to God, that he may offer both gifts and sacrifices for sins: who can have compassion on the ignorant, and on them that are out of the way; for that he himself also is compassed with infirmity. And by reason hereof he ought, as for the people, so also for himself, to offer for sins. And no man taketh this honour unto himself, but he that is called of God, as was Aaron.

\readingcitation{2. Epistle}{Hebrews 13:7}
\lett{B}{rethren:} Remember them which have the rule over you, who have spoken unto you the word of God: whose faith follow, considering the end of their conversation. Jesus Christ the same yesterday, and to day, and for ever. Be not carried about with divers and strange doctrines. For it is a good thing that the heart be established with grace; not with meats, which have not profited them that have been occupied therein.

We have an altar, whereof they have no right to eat which serve the tabernacle. For the bodies of those beasts, whose blood is brought into the sanctuary by the high priest for sin, are burned without the camp. Wherefore Jesus also, that he might sanctify the people with his own blood, suffered without the gate. Let us go forth therefore unto him without the camp, bearing his reproach. For here have we no continuing city, but we seek one to come. By him therefore let us offer the sacrifice of praise to God continually, that is, the fruit of our lips giving thanks to his name. But to do good and to communicate forget not: for with such sacrifices God is well pleased. Obey them that have the rule over you, and submit yourselves: for they watch for your souls, as they that must give account.

\readingcitation{1. Gospel}{Luke 11:33}
\lett{A}{t that time:} Jesus said unto his disciples: No man, when he hath lighted a candle, putteth it in a secret place, neither under a bushel, but on a candlestick, that they which come in may see the light. The light of the body is the eye: therefore when thine eye is single, thy whole body also is full of light; but when thine eye is evil, thy body also is full of darkness. Take heed therefore that the light which is in thee be not darkness. If thy whole body therefore be full of light, having no part dark, the whole shall be full of light, as when the bright shining of a candle doth give thee light.

\readingcitation{2. Gospel}{Mark 13:33}\label{Mark1333}
\lett{A}{t that time:} Jesus said unto his disciples: Take ye heed, watch and pray: for ye know not when the time is. For the Son of man is as a man taking a far journey, who left his house, and gave authority to his servants, and to every man his work, and commanded the porter to watch. Watch ye therefore: for ye know not when the master of the house cometh, at even, or at midnight, or at the cockcrowing, or in the morning: lest coming suddenly he find you sleeping. And what I say unto you I say unto all, Watch.

\subby{Common of Doctors}\label{CommonDoctors}
\fancyhead[RO,LE]{\textit{Doctors}}

\begin{rubric}
	The Daily Office propers are from the closest relevant Common, except for the \emph{Magnificat} Antiphon, as followeth.
\end{rubric}

\antiphon{Mag.}{O Teacher right excellent, {\dag} O light of Holy Church, O blessed \textit{N.}, lover of the divine law: intercede for us unto the Son of God.}

\introit
\lett{I}{n} the midst of the Church he opened his mouth: and the Lord filled him with the spirit of wisdom and of understanding: and clothed him with a robe of glory. (Alleluia, alleluia.) \textit{Ps.} It is a good thing to give thanks unto the Lord: and to sing praises unto thy name, O most Highest.

\collect
\lett{O}{God,} who didst give blessed \textit{N.} unto thy people to be a minister of everlasting salvation: grant, we beseech thee; that as we have learned of him the doctrine of life on earth, so we may be found worthy to have him for our advocate in heaven. Through.

\readingcitation{Epistle}{2 Timothy 4:1}\label{2Timothy41}
\lett{D}{early beloved:} I charge thee therefore before God, and the Lord Jesus Christ, who shall judge the quick and the dead at his appearing and his kingdom; preach the word; be instant in season, out of season; reprove, rebuke, exhort with all longsuffering and doctrine. For the time will come when they will not endure sound doctrine; but after their own lusts shall they heap to themselves teachers, having itching ears; and they shall turn away their ears from the truth, and shall be turned unto fables. But watch thou in all things, endure afflictions, do the work of an evangelist, make full proof of thy ministry.\par
For I am now ready to be offered, and the time of my departure is at hand. I have fought a good fight, I have finished my course, I have kept the faith: henceforth there is laid up for me a crown of righteousness, which the Lord, the righteous judge, shall give me at that day: and not to me only, but unto all them also that love his appearing.

\gradall{The mouth of the righteous is exercised in wisdom, and his tongue will be talking of judgment. ℣. The law of his God is in his heart: and his goings shall not slide.}{Alleluia, alleluia. ℣. The Lord loved him, and adorned him: and clothed him with a robe of glory. Alleluia.}

\begin{rubric}
{In Septuagesimatide or Lent, replacing the Alleluia:}
\end{rubric}\par\noindent
\tract{Blessed is the man that feareth the Lord: he hath great delight in his commandments. ℣. His seed shall be mighty upon earth: the generation of the faithful shall be blessed. ℣. Riches and plenteousness shall be in his house: and his righteousness endureth for ever.}

\begin{rubric}
{In Eastertide, replacing the Lesser Alleluia:}
\end{rubric}\par\noindent
\alleluia{Alleluia, alleluia. ℣. The Lord loved him, and adorned him: and clothed him with a robe of glory. Alleluia. ℣. The righteous shall grow as the lily; and flourish for ever before the Lord. Alleluia.}

\readingcitation{Gospel}{Matthew 5:13}
\lett{A}{t that time:} Jesus said unto his disciples: Ye are the salt of the earth: but if the salt have lost his savour, wherewith shall it be salted? it is thenceforth good for nothing, but to be cast out, and to be trodden under foot of men. Ye are the light of the world. A city that is set on an hill cannot be hid. Neither do men light a candle, and put it under a bushel, but on a candlestick; and it giveth light unto all that are in the house. Let your light so shine before men, that they may see your good works, and glorify your Father which is in heaven.\par
Think not that I am come to destroy the law, or the prophets: I am not come to destroy, but to fulfil. For verily I say unto you, Till heaven and earth pass, one jot or one tittle shall in no wise pass from the law, till all be fulfilled. Whosoever therefore shall break one of these least commandments, and shall teach men so, he shall be called the least in the kingdom of heaven: but whosoever shall do and teach them, the same shall be called great in the kingdom of heaven.

\offertory{The righteous shall flourish like a palm-tree: and shall spread abroad like a cedar in Libanus (Alleluia).}

\secret
\begin{inhead}
	Doctor \& Bishop
\end{inhead}
\lett{M}{ay} the devout prayers of thy Bishop and Doctor, Saint \textit{N.}, never fail to succour us, O Lord: that they may render our oblations acceptable in thy sight; and may ever obtain for us thy merciful pardon. Through.

\begin{inhead}
	Doctor not a Bishop
\end{inhead}
\lett{M}{ay} the devout prayers of thy Confessor and Doctor, Saint \textit{N.}, never fail to succour us, O Lord: that they may render our oblations acceptable in thy sight; and may ever obtain for us thy merciful pardon. Through.

\communion{A faithful and wise steward, whom his lord hath made ruler over his household: to give them their portion of meat in due season (Alleluia).}

\postcommunion
\begin{inhead}
	Doctor \& Bishop
\end{inhead}
\lett{W}{e} beseech thee, O Lord, that blessed \textit{N.}, thy Bishop and illustrious Doctor, may stand before thee as our advocate: that these thy sacrifices may avail for our salvation. Through.

\begin{inhead}
	Doctor not a Bishop
\end{inhead}
\lett{W}{e} beseech thee, O Lord, that blessed \textit{N.}, thy Confessor and illustrious Doctor, may stand before thee as our advocate: that these thy sacrifices may avail for our salvation. Through.

\begin{rubric}
	The following Epistle may be used instead of the one provided above.
\end{rubric}
\readingcitation{Another Epistle}{Ecclesiasticus 39:5}\label{Ecclesiasticus395}
%RV:
\lett{T}{he} just will apply his heart to resort early to the Lord that made him, And will make supplication before the Most High, And will open his mouth in prayer, And will make supplication for his sins. If the great Lord will, He shall be filled with the spirit of understanding: He shall pour forth the words of his wisdom, And in prayer give thanks unto the Lord. He shall direct his counsel and knowledge, And in his secrets shall he meditate. He shall shew forth the instruction which he hath been taught, And shall glory in the law of the covenant of the Lord. Many shall commend his understanding; And so long as the world endureth, it shall not be blotted out: His memorial shall not depart, And his name shall live from generation to generation. Nations shall declare his wisdom, And the congregation shall tell out his praise. 
%\lett{T}{he} just will give his heart to resort early to the Lord that made him, and will pray before the most High, and will open his mouth in prayer, and make supplication for his sins. When the great Lord will, he shall be filled with the spirit of understanding: he shall pour out wise sentences, and give thanks unto the Lord in his prayer. He shall direct his counsel and knowledge, and in his secrets shall he meditate. He shall shew forth that which he hath learned, and shall glory in the law of the covenant of the Lord. Many shall commend his understanding; and so long as the world endureth, it shall not be blotted out; his memorial shall not depart away, and his name shall live from generation to generation. Nations shall shew forth his wisdom, and the congregation shall declare his praise.

\subby{I: Common of a Confessor not a Bishop}\label{CommonConfessorNotBishopI}
\fancyhead[RO,LE]{\textit{Confessor not Bishop I}}

\subbysub{I Evensong}

\begin{rubric}
	The Hymn is from the First Common of a Confessor Bishop (p. \pageref{CommonConfessorBishopI}), with the following.
\end{rubric}
	
	℣. The Lord loved him and adorned him.

	℟. He clothed him with a robe of glory.
	
\antiphon{Mag.}{I will liken him {\dag} unto a wise man, which built his house upon a rock.}

\subbysub{Mattins}

\begin{rubric}
	The Invitatory Hymn is from I Evensong of the First Common of a Confessor Bishop (p. \pageref{CommonConfessorBishopI})
\end{rubric}

\officehymn

\gregorioscore{resources/gabc/Common/ConfessorNotBishopMattins.gabc}

℣. The Lord guided the righteous in right paths.

℟. And shewed him the kingdom of God.

\antiphon{Ben.}{Well done, good and faithful servant, {\dag} thou hast been faithful over a few things, I will make thee ruler over many things: enter thou into the joy of thy Lord.}

\subbysub{II Evensong}

\begin{rubric}
	The Hymn is from I Evensong of the First Common of a Confessor Bishop (p. \pageref{CommonConfessorBishopI}), with the following.
\end{rubric}

℣. The Lord guided the righteous in right paths.

℟. And shewed him the kingdom of God.

\antiphon{Mag.}{The world, and all things earthly, {\dag} this man despised, triumphant: he laid up treasures in heaven by word and deed.}

\introit
\lett{T}{he} mouth of the righteous is exercised in wisdom, and his tongue will be talking of judgment: the law of his God is in his heart. (Alleluia, alleluia.) \textit{Ps.} Fret not thyself because of the ungodly: neither be thou envious against the evildoers.

\collect
\lett{O}{God,} who makest us glad with the yearly solemnity of blessed \textit{N.}, thy Confessor: mercifully grant; that, as we now celebrate his birthday, so we may follow the example of his life. Through.

\readingcitation{Epistle}{Ecclesiasticus 31:8}
%RV:
\lett{B}{lessed} is the man that is found without blemish, And that goeth not after gold. Who is he? and we will call him blessed: For wonderful things hath he done among his people. Who hath been tried thereby, and found perfect? Then let him glory. Who hath had the power to transgress, and hath not transgressed? And to do evil, and hath not done it? His goods shall be made sure, And the congregation shall declare his alms. 

%\lett{B}{lessed} is the man that is found without blemish, and hath not gone after gold. Who is he? and we will call him blessed: for wonderful things hath he done among his people. Who hath been tried thereby, and found perfect? then let him glory. Who might offend, and hath not offended? or done evil, and hath not done it? His goods shall be established, and the congregation shall declare his alms.

\gradall{The righteous shall flourish like a palm-tree: and shall spread abroad like a cedar in Libanus in the house of the Lord. ℣. To tell of thy loving-kindness early in the morning, and of thy truth in the night-season.}{Alleluia, alleluia. ℣. Blessed is the man that endureth temptation: for when he is tried, he shall receive the crown of life. Alleluia.}

\begin{rubric}
{In Septuagesimatide or Lent, replacing the Alleluia:}
\end{rubric}\par\noindent
\tract{Blessed is the man that feareth the Lord: he hath great delight in his commandments. ℣. His seed shall be mighty upon earth: the generation of the faithful shall be blessed. ℣. Riches and plenteousness shall be in his house: and his righteousness endureth for ever.}

\begin{rubric}
{In Eastertide, replacing the Lesser Alleluia:}
\end{rubric}\par\noindent
\alleluia{Alleluia, alleluia. ℣. Blessed is the man that endureth temptation: for when he is tried, he shall receive the crown of life. Alleluia. ℣. The Lord loved him, and adorned him: and clothed him with a robe of glory. Alleluia.}

\readingcitation{Gospel}{Luke 12:35}
\lett{A}{t that time:} Jesus said unto his disciples: Let your loins be girded about, and your lights burning; and ye yourselves like unto men that wait for their lord, when he will return from the wedding; that when he cometh and knocketh, they may open unto him immediately. Blessed are those servants, whom the lord when he cometh shall find watching: verily I say unto you, that he shall gird himself, and make them to sit down to meat, and will come forth and serve them. And if he shall come in the second watch, or come in the third watch, and find them so, blessed are those servants. And this know, that if the goodman of the house had known what hour the thief would come, he would have watched, and not have suffered his house to be broken through. Be ye therefore ready also: for the Son of man cometh at an hour when ye think not.

\offertory{My truth and my mercy shall be with him: and in my name shall his horn be exalted (Alleluia).}

\secret
\lett{G}{rant,} we beseech thee, O Lord, that we who, trusting in this our sacrifice of praise, do offer it before thee to the honour of thy Saints: may by the same be delivered from all evils both in this life and in that which is to come. Through.

\communion{Blessed is the servant whom the lord when he cometh shall find watching: verily I say unto you, he shall make him ruler over all his goods (Alleluia).}

\postcommunion
\lett{O}{Lord,} our God, who hast refreshed us with heavenly meat and drink, we humbly beseech thee: that we may be defended by the prayers of him in whose memory we have received the same. Through.


\subby{II: Common of a Confessor not a Bishop}\label{CommonConfessorNotBishopII}
\fancyhead[RO,LE]{\textit{Confessor not Bishop II}}

\begin{rubric}
	Daily Office propers are as in the First Common of a Confessor not a Bishop (p. \pageref{CommonConfessorNotBishopI}).
\end{rubric}

\introit
\lett{T}{he} just shall flourish like a palm-tree: and shall spread abroad like a cedar in Libanus: planted in the house of the Lord: in the courts of the house of our God. (Alleluia, alleluia.) \textit{Ps.} It is a good thing to give thanks unto the Lord: and to sing praises unto thy name, O most Highest.

\collect
\lett{A}{ssist} us mercifully, O Lord, in these our supplications which we make before thee on the solemnity of blessed \textit{N.}, thy Confessor: that we, who put not our trust in our own righteousness, may be succoured by the prayers of him who found favour in thy sight. Through.

\readingcitation{Epistle}{1 Corinthians 4:9}
\lett{B}{rethren:} We are made a spectacle unto the world, and to angels, and to men. We are fools for Christ's sake, but ye are wise in Christ; we are weak, but ye are strong; ye are honourable, but we are despised. Even unto this present hour we both hunger, and thirst, and are naked, and are buffeted, and have no certain dwellingplace; and labour, working with our own hands: being reviled, we bless; being persecuted, we suffer it: being defamed, we intreat: we are made as the filth of the world, and are the offscouring of all things unto this day. I write not these things to shame you, but as my beloved sons I warn you: in Christ Jesus our Lord.

\gradall{The mouth of the righteous is exercised in wisdom, and his tongue will be talking of judgment. ℣. The law of his God is in his heart: and his goings shall not slide.}{Alleluia, alleluia. ℣. Blessed is the man that feareth the Lord: he hath great delight in his commandments. Alleluia.}

\begin{rubric}
{In Septuagesimatide or Lent, replacing the Alleluia:}
\end{rubric}\par\noindent
\tract{Blessed is the man that feareth the Lord: he hath great delight in his commandments. ℣. His seed shall be mighty upon earth: the generation of the faithful shall be blessed. ℣. Riches and plenteousness shall be in his house: and his righteousness endureth for ever.}

\begin{rubric}
{In Eastertide, replacing the Lesser Alleluia:}
\end{rubric}\par\noindent
\alleluia{Alleluia, alleluia. ℣. Blessed is the man that feareth the Lord: he hath great delight in his commandments. Alleluia. ℣. The righteous shall grow as the lily: and flourish for ever before the Lord. Alleluia.}

\readingcitation{Gospel}{Luke 12:32}
\lett{A}{t that time:} Jesus said unto his disciples: Fear not, little flock; for it is your Father's good pleasure to give you the kingdom. Sell that ye have, and give alms; provide yourselves bags which wax not old, a treasure in the heavens that faileth not, where no thief approacheth, neither moth corrupteth. For where your treasure is, there will your heart be also.

\offertory{The righteous shall rejoice in thy strength, O Lord, exceeding glad shall he be of thy salvation: thou hast given him his heart's desire (Alleluia).}

\secret
\lett{G}{rant} to us, we beseech thee, almighty God: that this oblation of our humble service may be acceptable in thy sight to the honour of thy Saints, and may cleanse us both in body and soul. Through.

\communion{Verily I say unto you: that ye which have forsaken all, and followed me, shall receive an hundredfold, and shall inherit everlasting life (Alleluia).}

\postcommunion
\lett{W}{e} beseech thee, almighty God: that we, who have received this heavenly food, may at the intercession of blessed \textit{N.}, thy Confessor, be thereby defended against all adversities. Through.

\begin{rubric}
	The following Epistle and Gospel may be used instead of the ones provided above.
\end{rubric}
\readingcitation{Another Epistle}{Philippians 3:7}\label{Philippians37}
\lett{B}{rethren:} What things were gain to me, those I counted loss for Christ. Yea doubtless, and I count all things but loss for the excellency of the knowledge of Christ Jesus my Lord: for whom I have suffered the loss of all things, and do count them but dung, that I may win Christ, and be found in him, not having mine own righteousness, which is of the law, but that which is through the faith of Christ, the righteousness which is of God by faith: that I may know him, and the power of his resurrection, and the fellowship of his sufferings, being made conformable unto his death; if by any means I might attain unto the resurrection of the dead. Not as though I had already attained, either were already perfect: but I follow after, if that I may apprehend that for which also I am apprehended of Christ Jesus.

\readingcitation{Another Gospel}{Luke 19:12}\label{Luke1912}
\lett{A}{t that time:} Jesus spake this parable unto his disciples: A certain nobleman went into a far country to receive for himself a kingdom, and to return. And he called his ten servants, and delivered them ten pounds, and said unto them, Occupy till I come. But his citizens hated him, and sent a message after him, saying, We will not have this man to reign over us. And it came to pass, that when he was returned, having received the kingdom, then he commanded these servants to be called unto him, to whom he had given the money, that he might know how much every man had gained by trading. Then came the first, saying, Lord, thy pound hath gained ten pounds. And he said unto him, Well, thou good servant: because thou hast been faithful in a very little, have thou authority over ten cities. And the second came, saying, Lord, thy pound hath gained five pounds. And he said likewise to him, Be thou also over five cities. And another came, saying, Lord, behold, here is thy pound, which I have kept laid up in a napkin: for I feared thee, because thou art an austere man: thou takest up that thou layedst not down, and reapest that thou didst not sow. And he saith unto him, Out of thine own mouth will I judge thee, thou wicked servant. Thou knewest that I was an austere man, taking up that I laid not down, and reaping that I did not sow: wherefore then gavest not thou my money into the bank, that at my coming I might have required mine own with usury? And he said unto them that stood by, Take from him the pound, and give it to him that hath ten pounds. (And they said unto him, Lord, he hath ten pounds.) For I say unto you, That unto every one which hath shall be given; and from him that hath not, even that he hath shall be taken away from him.

\subby{Common of Abbots}\label{CommonAbbots}
\fancyhead[RO,LE]{\textit{Abbots}}

\begin{rubric}
	Daily Office propers are as in the First Common of a Confessor not a Bishop (p. \pageref{CommonConfessorNotBishopI}).
\end{rubric}

%\subsubsec{Daily Office Propers}\par\noindent
%\antiphon{Mag.}{I will liken him {\dag} unto a wise man, which built his house upon a rock.}

%\antiphon{Ben.}{Well done, good and faithful servant, {\dag} thou hast been faithful over a few things, I will make thee ruler over many things: enter thou into the joy of the Lord.}

%\antiphon{Mag.}{The world, and all things earthly, {\dag} this man despised, triumphant: he laid up treasures in heaven by word and deed.}
\introit
\lett{T}{he} mouth of the righteous is exercised in wisdom, and his tongue will be talking of judgment: the law of his God is in his heart. (Alleluia, alleluia) \textit{Ps.} Fret not thyself because of the ungodly: neither be thou envious against the evil doers.

\collect
\lett{W}{e} beseech thee, O Lord, that the prayers of the blessed Abbot, \textit{N.} may commend us unto thee: that by his advocacy thou wouldest vouchsafe unto us those things to which by our own merits we cannot obtain. Through.

\readingcitation{Epistle}{Ecclesiasticus 45:1}
%RV:
\lett{B}{eloved} of God and men, even Moses, Whose memorial is blessed. He made him like to the glory of the saints, And magnified him in the fears of his enemies. By his words he caused the wonders to cease; He glorified him in the sight of kings; He gave him commandment for his people, And shewed him part of his glory. He sanctified him in his faithfulness and meekness; He chose him out of all flesh. He made him to hear his voice, And led him into the thick darkness, And gave him commandments face to face, Even the law of life and knowledge, 

%\lett{B}{eloved} of God and men, whose memorial is blessed. He made him like to the glorious saints, and magnified him, so that his enemies stood in fear of him. By his words he caused the wonders to cease, and he made him glorious in the sight of kings, and gave him a commandment for his people, and shewed him part of his glory. He sanctified him in his faithfuless and meekness, and chose him out of all men. He made him to hear his voice, and brought him into the dark cloud, and gave him commandments before his face, even the law of life and knowledge.

\gradall{Thou hast prevented him, O Lord, with the blessings of goodness: and hast set a crown of pure gold upon his head. ℣. He asked life of thee, and thou gavest him a long life, even for ever and ever.}{Alleluia, alleluia. ℣. The righteous shall flourish like a palm-tree: and shall spread abroad like a cedar in Libanus. Alleluia.}

\begin{rubric}
{In Septuagesimatide or Lent, replacing the Alleluia:}
\end{rubric}\par\noindent
\tract{Blessed is the man that feareth the Lord: he hath great delight in his commandments. ℣. His seed shall be mighty upon earth: the generation of the faithful shall be blessed. ℣. Riches and plenteousness shall be in his house: and his righteousness endureth for ever.}

\begin{rubric}
{In Eastertide, replacing the Lesser Alleluia:}
\end{rubric}\par\noindent
\alleluia{Alleluia, alleluia. ℣. The righteous shall flourish like a palm-tree: and shall spread abroad like a cedar in Libanus. Alleluia. ℣. The righteous shall grow as the lily: and flourish for ever before the Lord. Alleluia.}

\readingcitation{Gospel}{Matthew 19:27}
\lett{A}{t that time:} Peter said unto Jesus, Behold, we have forsaken all, and followed thee; what shall we have therefore? And Jesus said unto them, Verily I say unto you, That ye which have followed me, in the regeneration when the Son of man shall sit in the throne of his glory, ye also shall sit upon twelve thrones, judging the twelve tribes of Israel. And every one that hath forsaken houses, or brethren, or sisters, or father, or mother, or wife, or children, or lands, for my name's sake, shall receive an hundredfold, and shall inherit everlasting life.

\offertory{Thou hast given him his heart's desire, O Lord, and hast not denied him the request of his lips: thou hast set a crown of pure gold upon his head. (Alleluia.)}

\secret
\lett{W}{e} beseech thee, O Lord, that thy holy Abbot \textit{N.} may intercede for us: that this sacrifice which we offer and present upon thy holy altar may be profitable unto us for our salvation. Through.

\communion{A faithful and wise steward, whom his lord hath made ruler over his household, to give them their portion of meat in due season (Alleluia.)}

\postcommunion
\lett{L}{et} thy sacrament, O Lord, which we have now received and the prayers of the blessed Abbot, \textit{N.}, effectually defend us: that we may both imitate the example of his conversation, and receive the succour of his intercession. Through.


\subby{I: Common of a Virgin Martyr}\label{CommonVirginMartyrI}
\fancyhead[RO,LE]{\textit{Virgin Martyr I}}

\subbysub{I Evensong}

\begin{rubric}
	The Hymn is of Mattins with the following Versicle \& Antiphon.
\end{rubric}

℣. In thy grace and in thy beauty.

℟. Go forth, ride prosperously, and reign.

\antiphon{Mag.}{Come thou Bride of Christ, {\dag} receive the crown which the Lord hath prepared for thee for ever.}\par

\subbysub{Mattins}

\invitatoryhymn

\gregorioscore{resources/gabc/Common/VirginMartyrInvitatory.gabc}

\officehymn

\gregorioscore{resources/gabc/Common/VirginMartyrMattins.gabc}

℣. Full of grace are thy lips.

℟. Because God hath blessed thee for ever.

\antiphon{Ben.}{The kingdom of heaven {\dag} is likened unto a merchantman seeking goodly pearls; who, when he had found one pearl of great price, went and sold all that he had, and bought it.}\par

\subbysub{II Evensong}

\begin{rubric}
	The Hymn \& Versicle are of Mattins with the following Antiphon.
\end{rubric}

\antiphon{Mag.}{Come thou Bride of Christ, {\dag} receive the crown which the Lord hath prepared for thee for ever.}

\introit
\lett{I}{will} speak of thy testimonies, even before kings, and will not be ashamed: and my delight shall be in thy commandments, which I have loved exceedingly (Alleluia, alleluia). \textit{Ps.} Blessed are those that are undefiled in the way: and walk in the law of the Lord.

\collect
\lett{O}{God,} who among the manifold works of thy power hast bestowed even upon the weakness of women the victory of martyrdom: mercifully grant; that we, who celebrate the birthday of blessed \textit{N.} thy Virgin and Martyr, may by her example be drawn nearer unto thee. Through.

\readingcitation{Epistle}{Ecclesiasticus 51:1}
%RV:
\lett{I}{will} give thanks unto thee, O Lord, O King, And will praise thee, O God my Saviour: I do give thanks unto thy name: For thou wast my protector and helper, And didst deliver my body out of destruction, And out of the snare of a slanderous tongue, From lips that forge lies, And wast my helper before them that stood by; And didst deliver me, according to the abundance of thy mercy, and greatness of thy name, From the gnashings of teeth ready to devour, Out of the hand of such as sought my life, Out of the manifold afflictions which I had; From the choking of a fire on every side, And out of the midst of fire which I kindled not; Out of the depth of the belly of the grave, And from an unclean tongue, And from lying words, The slander of an unrighteous tongue unto the king. My soul drew near even unto death, And my life was near to the grave beneath. They compassed me on every side, And there was none to help me. I was looking for the succour of men, And it was not. And I remembered thy mercy, O Lord, And thy working which hath been from everlasting, How thou deliverest them that wait for thee, And savest them out of the hand of the enemies, O Lord our God.

%\lett{I}{will} thank thee, O Lord and King, and praise thee, O God my Saviour. I do give praise unto thy name: for thou art my defender and helper, and hast preserved my body from destruction, and from the snare of the slanderous tongue, and from the lips that forge lies, and hast been mine helper against mine adversaries. And hast delivered me, according to the multitude of thy mercies and greatness of thy name, from the teeth of them that were ready to devour me, and out of the hands of such as sought after my life, and from the manifold afflictions which I had: from the choking of fire on every side, and from the midst of the fire which I kindled not: from the depth of the belly of hell, from an unclean tongue and from lying words. By an accusation to the king from an unrighteous tongue my soul drew near even unto death: but thou deliverest such as wait for thee, and savest them out of the hands of their enemies, O Lord our God.

\gradall{Thou hast loved righteousness, and hated iniquity. ℣. Wherefore God, even thy God, hath anointed thee with the oil of gladness.}{Alleluia, alleluia. ℣. The Virgins that be her fellows shall be brought unto the King: they that bear her company shall be brought unto thee with joy. Alleluia.}

\begin{rubric}
{In Septuagesimatide or Lent, replacing the Alleluia:}
\end{rubric}\par\noindent
\tract{Come, Spouse of Christ, receive the crown, which the Lord hath prepared for thee for ever: for love of whom thou didst shed thy blood. ℣. Thou hast loved righteousness, and hated iniquity: wherefore God, even thy God, hath anointed thee with the oil of gladness above thy fellows. ℣. In thy comeliness and in thy beauty, go forth, proceed prosperously and reign.}

\begin{rubric}
{In Eastertide, replacing the Lesser Alleluia:}
\end{rubric}\par\noindent
\alleluia{Alleluia, alleluia. ℣. The Virgins that be her fellows shall be brought unto the King: they that bear her company shall be brought unto thee with joy. Alleluia. ℣. In thy comeliness and in thy beauty go forth, proceed prosperously and reign. Alleluia.}

\readingcitation{Gospel}{Matthew 25:1}
\lett{A}{t that time:} Jesus spake this parable unto his disciples: The kingdom of heaven be likened unto ten virgins, which took their lamps, and went forth to meet the bridegroom. And five of them were wise, and five were foolish. They that were foolish took their lamps, and took no oil with them: but the wise took oil in their vessels with their lamps. While the bridegroom tarried, they all slumbered and slept. And at midnight there was a cry made, Behold, the bridegroom cometh; go ye out to meet him. Then all those virgins arose, and trimmed their lamps. And the foolish said unto the wise, Give us of your oil; for our lamps are gone out. But the wise answered, saying, Not so; lest there be not enough for us and you: but go ye rather to them that sell, and buy for yourselves. And while they went to buy, the bridegroom came; and they that were ready went in with him to the marriage: and the door was shut. Afterward came also the other virgins, saying, Lord, Lord, open to us. But he answered and said, Verily I say unto you, I know you not. Watch therefore, for ye know neither the day nor the hour.

\offertory{The Virgins that be her fellows shall be brought unto the King: they that bear her company shall be brought unto thee, with joy and gladness: and shall enter into the palace of the Lord the King (Alleluia).}

\secret
\lett{R}{eceive,} O Lord, the gifts which we offer on the solemnity of blessed \textit{N.}, thy Virgin and Martyr: through whose advocacy we trust to be delivered. Through.

\communion{Let the proud be confounded, for they go wickedly about to destroy me: but I will be occupied in thy commandments and in thy statutes, that I be not ashamed (Alleluia).}

\postcommunion
\lett{M}{ay} the mysteries which we have received be for our succour, O Lord: and at the intercession of blessed \textit{N.}, thy Virgin and Martyr, cause us to rejoice in thy continual protection. Through.


\subby{II: Common of a Virgin Martyr}\label{CommonVirginMartyrII}
\fancyhead[RO,LE]{\textit{Virgin Martyr II}}

\begin{rubric}
	Daily Office propers as in the First Common of a Virgin Martyr (p. \pageref{CommonVirginMartyrI}).
\end{rubric}

\introit
\lett{T}{he} ungodly laid wait for me to destroy me: O Lord, I will consider thy testimonies: I see that all things come to an end: but thy commandment is exceeding broad (Alleluia, alleluia). \textit{Ps.} Blessed are those that are undefiled in the way: and walk in the law of the Lord.

\collect
\lett{W}{e} beseech thee, O Lord, that, like as blessed \textit{N.}, thy Virgin and Martyr, was ever found pleasing unto thee, both by the merit of her chastity, and by her confession of thy power: so she may implore for us thy pardon. Through.

\readingcitation{Epistle}{Ecclesiasticus 51:9}
%RV:
\lett{O}{Lord my God,} I lifted up my supplication from the earth, And prayed for deliverance from death. I called upon the Lord, the Father of my Lord, That he would not forsake me in the days of affliction, In the time when there was no help against the proud. I will praise thy name continually, And will sing praise with thanksgiving; And my supplication was heard: For thou savedst me from destruction, And deliveredst me from the evil time: Therefore will I give thanks and praise unto thee, And bless thy name, O Lord our God.

%\lett{O}{Lord} my God, thou has lifted up my habitation upon the earth, and I prayed for deliverance from death. I called upon the Lord, the Father of my Lord, that he would not leave me in the days of my trouble, and in the time of the proud, when there was no help. I will praise thy name continually, and will sing praise with thanksgiving, and so my prayer was heard. For thou savedst me from destruction, and deliveredst me from the evil time. Therefore will I give thanks, and praise thee, O Lord our God.

\gradall{God shall help her with his countenance: God is in the midst of her, therefore shall she not be removed. ℣. The rivers of the flood thereof shall make glad the city of God: the holy place of the tabernacle of the Most Highest.}{Alleluia, alleluia. ℣. This is a wise Virgin, and one of the number of the prudent. Alleluia.}

\begin{rubric}
{In Septuagesimatide or Lent, replacing the Alleluia:}
\end{rubric}\par\noindent
\tract{Come, Spouse of Christ, receive the crown which the Lord hath prepared for thee for ever: for love of whom thou didst shed thy blood. ℣. Thou hast loved righteousness, and hated iniquity: wherefore God, even thy God, hath anointed thee with the oil of gladness above thy fellows. ℣. In thy comeliness and in thy beauty go forth, proceed prosperously and reign.}

\begin{rubric}
{In Eastertide, replacing the Lesser Alleluia:}
\end{rubric}\par\noindent
\alleluia{Alleluia, alleluia. ℣. This is a wise Virgin, and one of the number of the prudent. Alleluia. ℣. O how beautiful is the chaste generation with its glory. Alleluia.}

\readingcitation{Gospel}{Matthew 13:44}
\lett{A}{t that time:} Jesus spake this parable unto his disciples: The kingdom of heaven is like unto treasure hid in a field; the which when a man hath found, he hideth, and for joy thereof goeth and selleth all that he hath, and buyeth that field.

Again, the kingdom of heaven is like unto a merchant man, seeking goodly pearls: who, when he had found one pearl of great price, went and sold all that he had, and bought it.

Again, the kingdom of heaven is like unto a net, that was cast into the sea, and gathered of every kind: which, when it was full, they drew to shore, and sat down, and gathered the good into vessels, but cast the bad away. So shall it be at the end of the world: the angels shall come forth, and sever the wicked from among the just, and shall cast them into the furnace of fire: there shall be wailing and gnashing of teeth.

Jesus saith unto them, Have ye understood all these things? They say unto him, Yea, Lord. Then said he unto them, Therefore every scribe which is instructed unto the kingdom of heaven is like unto a man that is an householder, which bringeth forth out of his treasure things new and old.

\offertory{Full of grace are thy lips, because God hath blessed thee for ever and ever (Alleluia).}

\secret
\lett{G}{raciously} receive, O Lord, through the merits of blessed \textit{N.}, thy Virgin and Martyr, the sacrifices which we offer unto thee: and grant that they may avail for our continual help. Through.

\communion{I deal with the thing that is lawful and right, O Lord, let the proud do me no wrong: I hold straight all thy commandments, and all false ways I utterly abhor (Alleluia).}

\postcommunion
\lett{O}{Lord} our God, who hast fulfilled us with the bounty of thy heavenly gift: we beseech thee, that, at the intercession of blessed \textit{N.}, thy Virgin and Martyr, we may ever live by the partaking of the same. Through.

\begin{rubric}
	The Gospel from the Common of Many Virgin Martyrs (p. \pageref{CommonVirginMartyrs}) may be used instead of the one provided above.
\end{rubric}


\subby{Common of Many Virgin Martyrs}\label{CommonVirginMartyrs}
\fancyhead[RO,LE]{\textit{Virgin Martyrs}}

\begin{rubric}
The propers as in the Second Common of a Virgin Martyr (p. \pageref{CommonVirginMartyrII}), except for the following.
\end{rubric}

\antiphon{Mag.}{O ye wise Virgins, {\dag} arise and trim your lamps: for behold the Bridegroom cometh; go ye out to meet him.}

\collect
\lett{G}{rant,} we beseech thee, O Lord our God, that we may at all times so devoutly honour the triumphs of thy holy Virgins and Martyrs \textit{N. and N.}: that, although we cannot worthily shew forth their praises, yet we may continually honour them with lowly service. Through.

\readingcitation{Epistle}{1 Corinthians 7:25}
\lett{B}{rethren:} Concerning virgins I have no commandment of the Lord: yet I give my judgment, as one that hath obtained mercy of the Lord to be faithful. I suppose therefore that this is good for the present distress, I say, that it is good for a man so to be. Art thou bound unto a wife? seek not to be loosed. Art thou loosed from a wife? seek not a wife. But and if thou marry, thou hast not sinned; and if a virgin marry, she hath not sinned. Nevertheless such shall have trouble in the flesh: but I spare you.

But this I say, brethren, the time is short: it remaineth, that both they that have wives be as though they had none; and they that weep, as though they wept not; and they that rejoice, as though they rejoiced not; and they that buy, as though they possessed not; and they that use this world, as not abusing it: for the fashion of this world passeth away.

But I would have you without carefulness. He that is unmarried careth for the things that belong to the Lord, how he may please the Lord: but he that is married careth for the things that are of the world, how he may please his wife. There is difference also between a wife and a virgin. The unmarried woman careth for the things of the Lord, that she may be holy both in body and in spirit: in Christ Jesus our Lord.

\readingcitation{Gospel}{Matthew 19:3}
\lett{A}{t that time:} The Pharisees came unto Jesus, tempting him, and saying unto him, Is it lawful for a man to put away his wife for every cause? And he answered and said unto them, Have ye not read, that he which made them at the beginning made them male and female, and said, For this cause shall a man leave father and mother, and shall cleave to his wife: and they twain shall be one flesh? Wherefore they are no more twain, but one flesh. What therefore God hath joined together, let not man put asunder. They say unto him, Why did Moses then command to give a writing of divorcement, and to put her away? He saith unto them, Moses because of the hardness of your hearts suffered you to put away your wives: but from the beginning it was not so. And I say unto you, Whosoever shall put away his wife, except it be for fornication, and shall marry another, committeth adultery: and whoso marrieth her which is put away doth commit adultery.

His disciples say unto him, If the case of the man be so with his wife, it is not good to marry. But he said unto them, All men cannot receive this saying, save they to whom it is given. For there are some eunuchs, which were so born from their mother's womb: and there are some eunuchs, which were made eunuchs of men: and there be eunuchs, which have made themselves eunuchs for the kingdom of heaven's sake. He that is able to receive it, let him receive it.

\secret
\lett{O}{Lord,} we beseech thee, look down upon these gifts, which we offer on thine altars on this festival of thy holy Virgins and Martyrs \textit{N. and N.}: that as by these blessed mysteries thou hast bestowed glory upon them; so likewise of thy bounty thou wouldest vouchsafe to us thy pardon. Through.

\postcommunion
\lett{G}{rant} to us, we beseech thee, O Lord, at the intercession of thy holy Virgins and Martyrs, \textit{N. and N.}: that those things which we touch with our lips we may receive in purity of heart. Through.

\subby{I: Common of a Virgin Only}\label{CommonVirginOnlyI}
\fancyhead[RO,LE]{\textit{Virgin I}}

\begin{rubric}
	Daily Office propers as in the First Common of a Virgin Martyr (p. \pageref{CommonVirginMartyrI}).
\end{rubric}

\introit
\lett{T}{hou} hast loved righteousness, and hated iniquity: wherefore God, even thy God, hath anointed thee with the oil of gladness above thy fellows (Alleluia, alleluia). \textit{Ps.} My heart is inditing of a good matter: I speak of the things which I have made unto the King.

\collect
\lett{G}{raciously} hear us, O God of our salvation: that, like as we do rejoice in the festival of blessed \textit{N.}, thy Virgin; so we may be instructed in all godly and devout affection. Through.

\readingcitation{Epistle}{2 Corinthians 10:17}
\lett{B}{rethren:} He that glorieth, let him glory in the Lord. For not he that commendeth himself is approved, but whom the Lord commendeth. Would to God ye could bear with me a little in my folly: and indeed bear with me. For I am jealous over you with godly jealousy: for I have espoused you to one husband, that I may present you as a chaste virgin to Christ.

\gradall{In thy comeliness and in thy beauty go forth, proceed prosperously and reign. ℣. Because of the word of truth, of meekness, and righteousness: and thy right hand shall teach thee terrible things.}{Alleluia, alleluia. ℣. The Virgins that be her fellows shall be brought unto the King: they that bear her company shall be brought unto thee with joy. Alleluia.}

\begin{rubric}
{In Septuagesimatide or Lent, replacing the Alleluia:}
\end{rubric}\par\noindent
\tract{Hearken, O daughter, and consider, incline thine ear: so shall the King have pleasure in thy beauty. ℣. The rich among the people shall make their supplication before thee: kings' daughters were among thy honourable women. ℣. The Virgins that be her fellows shall be brought unto the King: they that bear her company shall be brought unto thee. ℣. With joy and gladness shall they be brought: and shall enter into the King's palace.}

\begin{rubric}
{In Eastertide, replacing the Lesser Alleluia:}
\end{rubric}\par\noindent
\alleluia{Alleluia, alleluia. ℣. The Virgins that be her fellows shall be brought unto the King: they that bear her company shall be brought unto them with joy. Alleluia. ℣. In thy comeliness and in thy beauty go forth, proceed prosperously and reign. Alleluia.}

\readingcitation{Gospel}{Matthew 25:1}
\lett{A}{t that time:} Jesus spake this parable unto his disciples: the kingdom of heaven shall be likened unto ten virgins, which took their lamps, and went forth to meet the bridegroom. And five of them were wise, and five were foolish. They that were foolish took their lamps, and took no oil with them: but the wise took oil in their vessels with their lamps. While the bridegroom tarried, they all slumbered and slept. And at midnight there was a cry made, Behold, the bridegroom cometh; go ye out to meet him. Then all those virgins arose, and trimmed their lamps. And the foolish said unto the wise, Give us of your oil; for our lamps are gone out. But the wise answered, saying, Not so; lest there be not enough for us and you: but go ye rather to them that sell, and buy for yourselves. And while they went to buy, the bridegroom came; and they that were ready went in with him to the marriage: and the door was shut. Afterward came also the other virgins, saying, Lord, Lord, open to us. But he answered and said, Verily I say unto you, I know you not. Watch therefore, for ye know neither the day nor the hour.

\offertory{Kings' daughters were among thy honourable women, upon thy right hand did stand the queen in a vesture of gold, wrought about with divers colours (Alleluia).}

\secret
\lett{G}{rant,} O Lord, that like as thy dedicated people do acknowledge that in tribulation they have been succoured by the merits of thy Saints: so this oblation, which they offer unto thee in honour of the same, may be acceptable in thy sight. Through.

\communion{The five wise virgins took oil in their vessels with their lamps: and at midnight there was a cry made: Behold the bridegroom cometh: go ye out to meet Christ the Lord.}

\postcommunion
\lett{O}{Lord,} who hast satisfied thy family with sacred gifts: we beseech thee; that we may at all times be comforted by the intercession of her whose festival we celebrate. Through.

%MANUAL ADJUSTMENT:
\clearpage
\subby{II: Common of a Virgin Only}\label{CommonVirginOnlyII}
\fancyhead[RO,LE]{\textit{Virgin II}}

\begin{rubric}
	Daily Office propers as in the First Common of a Virgin Martyr (p. \pageref{CommonVirginMartyrI}), except for the following Invitatory Hymn.
\end{rubric}

\gregorioscore{resources/gabc/Common/VirginMartyrInvitatory.gabc}

\introit
\lett{A}{ll} the rich among the people shall make their supplication before thee: the Virgins that be her fellows shall be brought unto the King: they that bear her company shall be brought unto thee with joy and gladness (Alleluia, alleluia). \textit{Ps.} My heart is inditing of a good matter: I speak of the things which I have made unto the King.

\collect
\lett{G}{raciously} hear us, O God of our salvation: that, like as we do rejoice in the festival of blessed \textit{N.}, thy Virgin, so we may be instructed in all godly and devout affection. Through.

\readingcitation{Epistle}{1 Corinthians 7:25}
\lett{B}{rethren:} Concerning virgins I have no commandment of the Lord: yet I give my judgment, as one that hath obtained mercy of the Lord to be faithful. I suppose therefore that this is good for the present distress, I say, that it is good for a man so to be. Art thou bound unto a wife? seek not to be loosed. Art thou loosed from a wife? seek not a wife. But and if thou marry, thou hast not sinned; and if a virgin marry, she hath not sinned. Nevertheless such shall have trouble in the flesh: but I spare you.

But this I say, brethren, the time is short: it remaineth, that both they that have wives be as though they had none; and they that weep, as though they wept not; and they that rejoice, as though they rejoiced not; and they that buy, as though they possessed not; and they that use this world, as not abusing it: for the fashion of this world passeth away.

But I would have you without carefulness. He that is unmarried careth for the things that belong to the Lord, how he may please the Lord: but he that is married careth for the things that are of the world, how he may please his wife. There is difference also between a wife and a virgin. The unmarried woman careth for the things of the Lord, that she may be holy both in body and in spirit: in Christ Jesus our Lord.

\gradall{The King shall have pleasure in thy beauty, for he is thy Lord God. ℣. Hearken, O daughter, and consider, incline thine ear.}{Alleluia, alleluia. ℣. This is a wise Virgin, and one of the number of the prudent. Alleluia.}

\begin{rubric}
{In Septuagesimatide or Lent, replacing the Alleluia:}
\end{rubric}\par\noindent
\tract{The King shall have pleasure in thy beauty. ℣. The rich among the people shall make their supplication before thee: kings' daughters were among thy honourable women. ℣. The Virgins that be her fellows shall be brought unto the King: they that bear her company shall be brought unto thee. ℣. With joy and gladness shall they be brought: and shall enter into the King's palace.}

\begin{rubric}
{In Eastertide, replacing the Lesser Alleluia:}
\end{rubric}\par\noindent
\alleluia{Alleluia, alleluia. ℣. This is a wise Virgin, and one of the number of the prudent. Alleluia. ℣. O how beautiful is the chaste generation with its glory. Alleluia.}

\readingcitation{Gospel}{Matthew 13:44}
\lett{A}{t that time:} Jesus spake this parable unto his disciples: The kingdom of heaven is like unto treasure hid in a field; the which when a man hath found, he hideth, and for joy thereof goeth and selleth all that he hath, and buyeth that field.

Again, the kingdom of heaven is like unto a merchant man, seeking goodly pearls: who, when he had found one pearl of great price, went and sold all that he had, and bought it.

Again, the kingdom of heaven is like unto a net, that was cast into the sea, and gathered of every kind: which, when it was full, they drew to shore, and sat down, and gathered the good into vessels, but cast the bad away. So shall it be at the end of the world: the angels shall come forth, and sever the wicked from among the just, and shall cast them into the furnace of fire: there shall be wailing and gnashing of teeth.

Jesus saith unto them, Have ye understood all these things? They say unto him, Yea, Lord. Then said he unto them, Therefore every scribe which is instructed unto the kingdom of heaven is like unto a man that is an householder, which bringeth forth out of his treasure things new and old.

\offertory{The Virgins that be her fellows shall be brought unto the King: they that bear her company shall be brought unto thee, with joy and gladness: and shall enter into the palace of the Lord the King (Alleluia).}

\secret
\lett{G}{rant,} O Lord, that like as thy dedicated people do acknowledge that in tribulation they have been succoured by the merits of thy Saints: so this oblation, which they do offer unto thee in honour of the same, may be acceptable in thy sight. Through.

\communion{The kingdom of heaven is like unto a merchant man, seeking goodly pearls: who when he had found one pearl of great price, gave all that he had, and bought it (Alleluia).}

\postcommunion
\lett{O}{Lord,} who hast satisfied thy family with sacred gifts: we beseech thee, that we may at all times be comforted by the intercession of her whose festival we celebrate. Through.

\begin{rubric}
	The following Gospel may be used instead of the one provided above.
\end{rubric}
\readingcitation{Another Gospel}{Matthew 25:1}
\lett{A}{t that time:} Jesus spake this parable unto his disciples: The kingdom of heaven be likened unto ten virgins, which took their lamps, and went forth to meet the bridegroom. And five of them were wise, and five were foolish. They that were foolish took their lamps, and took no oil with them: but the wise took oil in their vessels with their lamps. While the bridegroom tarried, they all slumbered and slept. And at midnight there was a cry made, Behold, the bridegroom cometh; go ye out to meet him. Then all those virgins arose, and trimmed their lamps. And the foolish said unto the wise, Give us of your oil; for our lamps are gone out. But the wise answered, saying, Not so; lest there be not enough for us and you: but go ye rather to them that sell, and buy for yourselves. And while they went to buy, the bridegroom came; and they that were ready went in with him to the marriage: and the door was shut. Afterward came also the other virgins, saying, Lord, Lord, open to us. But he answered and said, Verily I say unto you, I know you not. Watch therefore, for ye know neither the day nor the hour.


\subby{Common of a Martyr not a Virgin}\label{CommonMartyrNotVirgin}
\fancyhead[RO,LE]{\textit{Martyr not Virgin}}

\subbysub{I Evensong}

\begin{rubric}
	The Hymn is of Mattins with the following Versicle \& Antiphon.
\end{rubric}

℣. In thy grace and in thy beauty.

℟. Go forth, ride prosperously, and reign.

\antiphon{Mag.}{The kingdom of heaven {\dag} is likened unto a merchantman seeking goodly pearls; who, when he had found one pearl of great price, went and sold all that he had, and bought it.}\par\noindent

\subbysub{Mattins}

\invitatoryhymn

\gregorioscore{resources/gabc/Common/MartyrNotVirginInvitatory.gabc}

\officehymn

\gregorioscore{resources/gabc/Common/MartyrNotVirginMattins.gabc}

℣. Full of grace are thy lips.

℟. Because God hath blessed thee for ever.

\antiphon{Ben.}{Give her {\dag} of the fruit of her hands; and let her own works praise her in the gates.}\par\noindent

\subbysub{II Evensong}

\begin{rubric}
	The Hymn \& Versicle are of Mattins with the following Antiphon.
\end{rubric}

\antiphon{Mag.}{She stretched out {\dag} her hand to the poor; yea, she reacheth forth her hands to the needy; she eateth not the bread of idleness.}

\introit
\lett{T}{he} ungodly laid wait for me to destroy me: O Lord, I will consider thy testimonies: I see that all things come to an end: but thy commandment is exceeding broad (Alleluia). \textit{Ps.} Blessed are those that are undefiled in the way: and walk in the law of the Lord.

%Roman Collect with no intercession.
\collect
\lett{O}{God,} who among the manifold works of thy power hast bestowed even upon the weakness of women the victory of martyrdom: mercifully grant; that we, who celebrate the birthday of blessed \textit{N.}, thy Martyr, may by her example draw nearer unto thee. Through.

\readingcitation{Epistle}{Ecclesiasticus 51:1}
%RV:
\lett{I}{will} give thanks unto thee, O Lord, O King, And will praise thee, O God my Saviour: I do give thanks unto thy name: For thou wast my protector and helper, And didst deliver my body out of destruction, And out of the snare of a slanderous tongue, From lips that forge lies, And wast my helper before them that stood by; And didst deliver me, according to the abundance of thy mercy, and greatness of thy name, From the gnashings of teeth ready to devour, Out of the hand of such as sought my life, Out of the manifold afflictions which I had; From the choking of a fire on every side, And out of the midst of fire which I kindled not; Out of the depth of the belly of the grave, And from an unclean tongue, And from lying words, The slander of an unrighteous tongue unto the king. My soul drew near even unto death, And my life was near to the grave beneath. They compassed me on every side, And there was none to help me. I was looking for the succour of men, And it was not. And I remembered thy mercy, O Lord, And thy working which hath been from everlasting, How thou deliverest them that wait for thee, And savest them out of the hand of the enemies, O Lord our God.

%\lett{I}{will} thank thee, O Lord and King, and praise thee, O God my Saviour. I do give praise unto thy name: for thou art my defender and helper, and hast preserved my body from destruction, and from the snare of the slanderous tongue, and from the lips that forge lies, and hast been mine helper against mine adversaries. And hast delivered me, according to the multitude of thy mercies and greatness of thy name, from the teeth of them that were ready to devour me, and out of the hands of such as sought after my life, and from the manifold afflictions which I had: from the choking of fire on every side, and from the midst of the fire which I kindled not: from the depth of the belly of hell, from an unclean tongue and from lying words. By an accusation to the king from an unrighteous tongue my soul drew near even unto death: but thou deliverest such as wait for thee, and savest them out of the hands of their enemies, O Lord our God.

\gradall{Thou hast loved righteousness and hated iniquity. ℣. Wherefore God, even thy God, hath anointed thee with the oil of gladness.}{Alleluia, alleluia. ℣. In thy comeliness and in thy beauty go forth, proceed prosperously and reign. Alleluia.}

\begin{rubric}
{In Septuagesimatide or Lent, replacing the Alleluia:}
\end{rubric}\par\noindent
\tract{Come, Spouse of Christ, receive the crown which the Lord hath prepared for thee for ever: for love of whom thou didst shed thy blood. ℣. Thou hast loved righteousness and hated iniquity: wherefore God, even thy God, hath anointed thee with the oil of gladness, above thy fellows. ℣. In thy comeliness and in thy beauty go forth, proceed prosperously and reign.}

\begin{rubric}
{In Eastertide, replacing the Lesser Alleluia:}
\end{rubric}\par\noindent
\alleluia{Alleluia, alleluia. ℣. In thy comeliness and in thy beauty, go forth, proceed prosperously and reign. Alleluia. ℣. Because of the word of truth, of meekness, and righteousness: and thy right hand shall teach thee terrible things. Alleluia.}

\readingcitation{Gospel}{Matthew 13:44}
\lett{A}{t that time:} Jesus spake this parable unto his disciples: The kingdom of heaven is like unto treasure hid in a field; the which when a man hath found, he hideth, and for joy thereof goeth and selleth all that he hath, and buyeth that field.

Again, the kingdom of heaven is like unto a merchant man, seeking goodly pearls: who, when he had found one pearl of great price, went and sold all that he had, and bought it.

Again, the kingdom of heaven is like unto a net, that was cast into the sea, and gathered of every kind: which, when it was full, they drew to shore, and sat down, and gathered the good into vessels, but cast the bad away. So shall it be at the end of the world: the angels shall come forth, and sever the wicked from among the just, and shall cast them into the furnace of fire: there shall be wailing and gnashing of teeth.

Jesus saith unto them, Have ye understood all these things? They say unto him, Yea, Lord. Then said he unto them, Therefore every scribe which is instructed unto the kingdom of heaven is like unto a man that is an householder, which bringeth forth out of his treasure things new and old.

\begin{rubric}
	After Septuagesima the \emph{Alleluia} at the end of the following Offertory is omitted.
\end{rubric}
\offertory{Full of grace are thy lips: because God hath blessed thee for ever and ever, alleluia.}

\secret
\lett{R}{eceive,} O Lord, the gifts which we offer on the solemnity of blessed \textit{N.}, thy Martyr: through whose advocacy we trust to be delivered. Through.

\communion{Princes have persecuted me without a cause: but my heart standeth in awe of thy word: I am as glad of thy word, as one that findeth great spoils (Alleluia).}

\postcommunion
\lett{M}{ay} the mysteries which we have received, O Lord, be for our succour: and, at the intercession of blessed \textit{N.} thy Martyr, cause us to rejoice in thy continual protection. Through.


\subby{Common of Many Martyrs not Virgins}\label{CommonMartyrsNotVirgins}
\fancyhead[RO,LE]{\textit{Martyrs not Virgins}}

\begin{rubric}
	The propers as in the Common of a Martyr not a Virgin (p. \pageref{CommonMartyrNotVirgin}), except for the following.
\end{rubric}

℣. Thou hast crowned them with glory and honour, O Lord.

℟. And madest them to have dominion of the works of thy hands.

\antiphon{Mag. \& Ben.}{For theirs is the kingdom of heaven, {\dag} who despised worldly living: who have won the rewards of the kingdom, and have washed their robes in the blood of the Lamb.}\par\noindent

\collect
\lett{G}{rant,} we beseech thee, O Lord our God, that we may at all times so devoutly honour the triumphs of thy holy Martyrs \textit{N. and N.}: that although we cannot worthily shew forth their praises, yet we may continually honour them with lowly service. Through.

\secret
\lett{O}{Lord,} we beseech thee, look down upon these gifts which we offer on thine altars on this festival of thy holy Martyrs \textit{N. and N.}: that as by these blessed mysteries thou hast bestowed glory upon them; so likewise of thy bounty thou wouldest vouchsafe us thy pardon. Through.

\postcommunion
\lett{G}{rant} to us, we beseech thee, O Lord, at the intercession of thy holy Martyrs \textit{N. and N.}: that those things which we touch with our lips, we may receive in purity of heart. Through.


\subby{Common of Neither Virgin nor Martyr}\label{CommonNeitherVirginMartyr}
\fancyhead[RO,LE]{\textit{Not Virgin Nor Martyr}}

\begin{rubric}
	Daily Office propers as in the Common of a Martyr not a Virgin (p. \pageref{CommonMartyrNotVirgin}).
\end{rubric}

%\antiphon{Mag.}{The kingdom of heaven {\dag} is likened unto a merchantman seeking goodly pearls; who, when he had found one pearl of great price, went and sold all that he had, and bought it.}\par\noindent
%\antiphon{Ben.}{Give her {\dag} of the fruit of her hands; and let her own works praise her in the gates.}\par\noindent
%\antiphon{Mag.}{She stretched out {\dag} her hand to the poor; yea, she reacheth forth her hands to the needy; she eateth not the bread of idleness.}

\introit
\lett{I}{know,} O Lord, that thy judgments are right, and that thou of very faithfulness hast caused me to be troubled: my flesh trembleth for fear of thee, and I am afraid of thy judgments (Alleluia, alleluia). \textit{Ps.} Blessed are those that are undefiled in the way: and walk in the law of the Lord.

\collect
\lett{G}{raciously} hear us, O God of our salvation: and grant that, like as we do rejoice in the festival of blessed \textit{N.}; so we may be instructed in all godly and devout affection. Through.

\readingcitation{Epistle}{Proverbs 31:10}
\lett{W}{ho} can find a virtuous woman? for her price is far above rubies. The heart of her husband doth safely trust in her, so that he shall have no need of spoil. She will do him good and not evil all the days of her life. She seeketh wool, and flax, and worketh willingly with her hands. She is like the merchants' ships; she bringeth her food from afar. She riseth also while it is yet night, and giveth meat to her household, and a portion to her maidens. She considereth a field, and buyeth it: with the fruit of her hands she planteth a vineyard. She girdeth her loins with strength, and strengtheneth her arms. She perceiveth that her merchandise is good: her candle goeth not out by night. She layeth her hands to the spindle, and her hands hold the distaff. She stretcheth out her hand to the poor; yea, she reacheth forth her hands to the needy. She is not afraid of the snow for her household: for all her household are clothed with scarlet. She maketh herself coverings of tapestry; her clothing is silk and purple. Her husband is known in the gates, when he sitteth among the elders of the land. She maketh fine linen, and selleth it; and delivereth girdles unto the merchant. Strength and honour are her clothing; and she shall rejoice in time to come. She openeth her mouth with wisdom; and in her tongue is the law of kindness. She looketh well to the ways of her household, and eateth not the bread of idleness. Her children arise up, and call her blessed; her husband also, and he praiseth her. Many daughters have done virtuously, but thou excellest them all. Favour is deceitful, and beauty is vain: but a woman that feareth the Lord, she shall be praised. Give her of the fruit of her hands; and let her own works praise her in the gates.

\gradall{Full of grace are thy lips: because God hath blessed thee for ever. ℣. Because of the word of truth, of meekness, and righteousness: and thy right hand shall teach thee terrible things.}{Alleluia, alleluia. ℣. In thy comeliness and in thy beauty go forth, ride prosperously, and reign. Alleluia.}

\begin{rubric}
{In Septuagesimatide or Lent, replacing the Alleluia:}
\end{rubric}\par\noindent
\tract{Come, Spouse of Christ, receive the crown, which the Lord hath prepared for thee for ever. ℣. Thou hast loved righteousness, and hated iniquity: wherefore God, even thy God, hath anointed thee with the oil of gladness above thy fellows. ℣. In thy comeliness and in thy beauty go forth, ride prosperously, and reign.}

\begin{rubric}
{In Eastertide, replacing the Lesser Alleluia:}
\end{rubric}\par\noindent
\alleluia{Alleluia, alleluia. ℣. In thy comeliness and in thy beauty, go forth, ride prosperously, and reign. Alleluia. ℣. Because of the word of truth, of meekness, and righteousness: and thy right hand shall teach thee terrible things. Alleluia.}

\readingcitation{Gospel}{Matthew 13:44}
\lett{A}{t that time:} Jesus spake this parable unto his disciples: The kingdom of heaven is like unto treasure hid in a field; the which when a man hath found, he hideth, and for joy thereof goeth and selleth all that he hath, and buyeth that field.

Again, the kingdom of heaven is like unto a merchant man, seeking goodly pearls: who, when he had found one pearl of great price, went and sold all that he had, and bought it.

Again, the kingdom of heaven is like unto a net, that was cast into the sea, and gathered of every kind: which, when it was full, they drew to shore, and sat down, and gathered the good into vessels, but cast the bad away. So shall it be at the end of the world: the angels shall come forth, and sever the wicked from among the just, and shall cast them into the furnace of fire: there shall be wailing and gnashing of teeth.

Jesus saith unto them, Have ye understood all these things? They say unto him, Yea, Lord. Then said he unto them, Therefore every scribe which is instructed unto the kingdom of heaven is like unto a man that is an householder, which bringeth forth out of his treasure things new and old.

\offertory{Full of grace are thy lips: because God hath blessed thee for ever and ever (Alleluia).}

\secret
\lett{G}{rant,} O Lord, that like as thy faithful people do acknowledge that in tribulation they have been succoured by the merits of thy Saints: so this oblation, which they here do offer unto thee in honour of the same, may be acceptable in thy sight. Through.

\communion{Thou hast loved righteousness, and hated iniquity: wherefore God, even thy God, hath anointed thee with the oil of gladness above thy fellows (Alleluia).}

\postcommunion
\lett{O}{Lord,} who hast satisfied thy family with sacred gifts: we beseech thee that we may at all times be comforted by the intercession of her whose festival we celebrate. Through.

\begin{rubric}
	The following Epistle may be used for a Widow instead of the Epistle provided above.
\end{rubric}
\readingcitation{Another Epistle}{1 Timothy 5:3}\label{1Timothy53}
\lett{D}{early beloved:} Honour widows that are widows indeed. But if any widow have children or nephews, let them learn first to shew piety at home, and to requite their parents: for that is good and acceptable before God. Now she that is a widow indeed, and desolate, trusteth in God, and continueth in supplications and prayers night and day. But she that liveth in pleasure is dead while she liveth. And these things give in charge, that they may be blameless. But if any provide not for his own, and specially for those of his own house, he hath denied the faith, and is worse than an infidel. Let not a widow be taken into the number under threescore years old, having been the wife of one man, well reported of for good works; if she have brought up children, if she have lodged strangers, if she have washed the saints' feet, if she have relieved the afflicted, if she have diligently followed every good work.

\subby{Common of the Dedication of a Church}\label{CommonDedication}
\fancyhead[RO,LE]{\textit{Church Dedication}}

\subbysub{I Evensong}

\gregorioscore{resources/gabc/Common/DedicationEvensong.gabc}

℣. This is the Lord's house, firmly builded.

℟. It is well founded on the solid rock.

\antiphon{Mag.}{The Lord hath hallowed {\dag} his tabernacle: for this is the house of the Lord, wherein men may call upon his Name; whereof it is written, My Name shall be therein, saith the Lord.}

\subbysub{Mattins}

\begin{rubric}
	The Invitatory Hymn is from I Evensong.
\end{rubric}

\officehymn

\gregorioscore{resources/gabc/Common/DedicationMattins.gabc}

℣. This is the Lord's house, firmly builded.

℟. It is well founded on the solid rock.

\antiphon{Ben.}{Zacchaeus, {\dag} make haste, and come down; for today I must abide at thy house: and he made haste, and came down, and received him joyfully into his house. This day is salvation come to this house form the Lord, alleluia.}

\subbysub{II Evensong}

\begin{rubric}
	The Hymn is of I Evensong with the following Versicle \& Antiphon.
\end{rubric}

℣. Holiness becometh thine house, O Lord.

℟. For ever.

\antiphon{Mag.}{O how dreadful {\dag} is this place! truly this is none other but the house of God, and the gate of heaven.}

\introit
\lett{H}{ow} dreadful is this place: this is the house of God, and the gate of heaven: and men shall call it the palace of God (Alleluia, alleluia). \textit{Ps.} O how amiable are thy dwellings, thou Lord of hosts! My soul hath a desire and longing to enter into the courts of the Lord.

\collect
\lett{O}{God,} who year by year renewest unto us the day of consecration of this thy holy temple, and dost ever bring us again in safety to thy sacred mysteries: graciously hear the prayers of thy people, and grant; that whosoever entereth this temple to ask for blessings may rejoice in the obtaining of all his petitions. Through.

\begin{rubric}
	On the day itself of the Dedication of a Church and through the Octave, and when the Collect is to be varied, is said the following.
\end{rubric}
\lett{O}{God,} who thyself invisible, containest all things, but dost for the salvation of mankind shew forth visibly the signs of thy power: enlighten this temple with the power of thine in-dwelling, and grant; that all they, who come together here to ask thy mercy, may, in whatsoever tribulation they call upon thee, obtain the blessings of thy consolation. Through.

\readingcitation{Epistle}{Revelation 21:2}
\lett{I}{n those days:} I saw the holy city, new Jerusalem, coming down from God out of heaven, prepared as a bride adorned for her husband. And I heard a great voice out of heaven saying, Behold, the tabernacle of God is with men, and he will dwell with them, and they shall be his people, and God himself shall be with them, and be their God. And God shall wipe away all tears from their eyes; and there shall be no more death, neither sorrow, nor crying, neither shall there be any more pain: for the former things are passed away. And he that sat upon the throne said, Behold, I make all things new. 

\gradall{This dwelling is God's handywork, it is a mystery beyond all price, that cannot be spoken against. ℣. O God, before whom standeth the choir of Angels, graciously hear the prayers of thy servants.}{Alleluia, alleluia. ℣. I will worship toward thy holy temple: and praise thy name. Alleluia.}

\begin{rubric}
{In Septuagesimatide or Lent, replacing the Alleluia:}
\end{rubric}\par\noindent
\tract{They that put their trust in the Lord shall be even as the mount Sion: which may not be removed, but standeth fast for ever. ℣. The hills stand about Jerusalem, even so standeth the Lord round about his people, from this time forth for evermore.}

\begin{rubric}
{In Eastertide, replacing the Lesser Alleluia:}
\end{rubric}\par\noindent
\alleluia{Alleluia, alleluia. ℣. I will worship toward thy holy temple: and praise thy name. Alleluia. ℣. The house of the Lord is well founded upon a sure rock. Alleluia.}

\readingcitation{Gospel}{Luke 19:1}
\lett{A}{t that time:} Jesus entered and passed through Jericho. And, behold, there was a man named Zacch{\ae}us, which was the chief among the publicans, and he was rich. And he sought to see Jesus who he was; and could not for the press, because he was little of stature. And he ran before, and climbed up into a sycomore tree to see him: for he was to pass that way. And when Jesus came to the place, he looked up, and saw him, and said unto him, Zacch{\ae}us, make haste, and come down; for to day I must abide at thy house. And he made haste, and came down, and received him joyfully. And when they saw it, they all murmured, saying, That he was gone to be guest with a man that is a sinner. And Zacch{\ae}us stood, and said unto the Lord; Behold, Lord, the half of my goods I give to the poor; and if I have taken any thing from any man by false accusation, I restore him fourfold. And Jesus said unto him, This day is salvation come to this house, forsomuch as he also is a son of Abraham. For the Son of man is come to seek and to save that which was lost.

\begin{rubric}
	After Septuagesima the \emph{Alleluia} at the end of the following Offertory is omitted.
\end{rubric}
\offertory{O Lord God, in the uprightness of mine heart I have willingly offered all these things; and now I have seen with joy thy people, which are present here: O God of Israel, keep this imagination of their heart, alleluia.}

\secret

\begin{inhead}
	In the Dedication itself
\end{inhead}
\lett{O}{Lord,} we beseech thee, mercifully hear our prayers: that all we, who are gathered within the precincts of this temple, whose dedication year by year we celebrate, may perfectly serve thee, both in body and soul; that, like as we do here present our offerings unto thee, so we may by thee be enabled to attain unto the rewards of everlasting felicity. Through.

\begin{inhead}
	Outside the dedicated Church itself
\end{inhead}
\lett{O}{Lord,} we beseech thee, mercifully hear our prayers: that, like as we do here present our offerings unto thee, so we may by thee be enable to attain unto the rewards of everlasting felicity. Through.

\begin{inhead}
	On the day itself of the Dedication of a Church, and through the Octave, and when the secret is to be varied, is said the following.
\end{inhead}
\lett{O}{God,} who art the author of the gifts to be consecrated unto thee, pour forth upon this house of prayer thy blessing: that all who shall call upon thy name therein may perceive the succour of thy defence. Through.

\begin{rubric}
	The common Preface, even in Lent, although there shall occur some Commemoration or common Octave with a Preface of their own, which are not of some Mystery of our Lord.
\end{rubric}

\communion{My house shall be called the house of prayer, saith the Lord: in it every one that asketh receiveth; and he that seeketh findeth: and to him that knocketh it shall be opened (Alleluia).}

\postcommunion
\lett{O}{God,} who of elect and living stones dost fashion for thy majesty an everlasting habitation: assist the supplications of us thy people; that, like as thy Church increaseth in visible habitations, so it may grow and prosper in spiritual advancement. Through.

\begin{rubric}
	On the day itself of the Dedication of a Church, and through the Octave, or where the Postcommunion is to be varied, is said the following.
\end{rubric}
\lett{W}{e} beseech thee, almighty God: that in this place, which we, though unworthy, have dedicated to thy name, thou wouldest incline thy merciful ears to all that call upon thee. Through.

\begin{rubric}
	On the day of the Dedication itself, Mass is said as above with the second Prayers; and the Collect of the Mystery or Saint in whose honour the church is dedicated is added to the first under one conclusion.
\end{rubric}

\begin{rubric}
	During the Octave, Mass is said as above; but within the Octave the \nth{2} Collect is of St. Mary, from the Votive Mass of the Season: \nth{3} against the persecutors of the Church or for the Chief Bishop.
\end{rubric}

\begin{rubric}
	On the day of the Dedication of an Altar, Mass is said as on the Dedication of a Church, except the Prayers following: to which likewise in the case of a fixed Altar the Prayer of the Mystery or Saint in whose honour the Altar is dedicated is added under one conclusion.
\end{rubric}

\collect
\lett{O}{God,} who from the whole assembly of thy Saints dost prepare for thyself an everlasting habitation: grant to this building the increase of heavenly grace; that like as we do here venerate with devout affection the relics of thy Saints, so we may ever be aided by their merits. Through.

\secret
\lett{W}{e} beseech thee, thy Lord our God, let thy Holy Spirit come down upon this altar: that he may sanctify the gifts of thy people, and graciously cleanse the hearts of them that receive the same. Through . . . in the unity of the same Holy Spirit.

\postcommunion
\lett{A}{lmighty} and everlasting God, sanctify with the power of thy heavenly benediction this altar, which we have dedicated to the honour of thy name: and shew forth upon all that put their trust in thee the bounty of thy succour; that they may here obtain the grace of thy sacraments and the fulfilment of their desires. Through.

\subby{Common of Feasts of the Blessed Virgin Mary}\label{CommonBVM}
\fancyhead[RO,LE]{\textit{Blessed Virgin Mary}}

\subbysub{I Evensong}

\gregorioscore{resources/gabc/Common/BVMEvensong.gabc}

℣. Vouchsafe that I may praise thee, O holy Virgin.

℟. Give me strength against thine enemies.

\antiphon{Mag.}{O holy Mary, {\dag} help thou the suffering, strengthen the faint-hearted, comfort the sorrowful; pray for the people, entreat for the clergy, intercede for all womankind vowed unto God: may all acknowledge the help of thy prayer, who celebrate thy holy festival.}

\subbysub{Mattins}

\begin{inhead}
	Invitatory Hymn
\end{inhead}

\gregorioscore{resources/gabc/Common/BVMInvitatory.gabc}

\begin{inhead}
	Office Hymn
\end{inhead}

\gregorioscore{resources/gabc/Common/BVMMattins.gabc}

℣. Full of grace are thy lips.

℟. Because God hath blessed thee for ever.

\antiphon{Ben.}{Blessed art thou, {\dag} O Mary, for thou hast believed: and there shall be a performance in thee of those things which were told thee from the Lord, alleluia.}

\subbysub{II Evensong}

\begin{rubric}
	The Hymn \& Versicle are of I Evensong with the following Antiphon.
\end{rubric}

\antiphon{Mag.}{All generations {\dag} shall call me blessed: for God hath regarded the lowliness of his handmaiden.}

\introit
\lett{H}{ail,} O Mother most holy, who in childbirth didst bring forth the Monarch: him who o'er heaven and earth reigneth for ever and ever (Alleluia, alleluia). \textit{Ps.} My heart is inditing of a good matter: I speak of the things which I have made unto the King.

\collect
\lett{G}{rant,} we beseech thee, O Lord God, that we thy servants may enjoy perpetual health of mind and of body: and, at the glorious intercession of blessed Mary ever Virgin, be delivered from present sadness, and rejoice in everlasting gladness. Through.

\readingcitation{Epistle}{Ecclesiasticus 24:9}
%RV:
\lett{H}{e} created me from the beginning before the world; And to the end I shall not fail. In the holy tabernacle I ministered before him; And so was I established in Sion. In the beloved city likewise he gave me rest; And in Jerusalem was my authority. And I took root in a people that was glorified, Even in the portion of the Lord’s own inheritance. 

%\lett{H}{e} created me from the beginning before the world, and I shall never fail, in the holy tabernacle I served before him. And so was I established in Sion, likewise in the beloved city he gave me rest, and in Jerusalem was my power. And I took root in an honourable people, even in the portion of the Lord's inheritance, and my abode is in the full assembly of the saints.

\gradall{Blessed and venerable art thou, O Virgin Mary: who without spot wast found Mother of the Saviour. ℣. Virgin Mother of God, he whom the whole world containeth not, being made man lay hid in thy womb.}{Alleluia, alleluia. ℣. After childbirth, O Virgin, thou didst remain inviolate: Mother of God, intercede for us. Alleluia.}
\begin{rubric}
	In Advent, the following Alleluia Verse replaces the one above.
\end{rubric}
\alleluia{Alleluia, alleluia. ℣. Hail, Mary, full of grace; the Lord is with thee: blessed art thou among women. Alleluia.}

\begin{rubric}
{In Septuagesimatide or Lent, replacing the Alleluia:}
\end{rubric}\par\noindent
\tract{Rejoice, O Virgin Mary, alone thou hast destroyed all heresies. ℣. Who didst believe the words of the Archangel Gabriel. ℣. Whilst a Virgin thou didst bring forth him who is God and man: and after childbirth, O Virgin, inviolate didst remain. ℣. Mother of God intercede for us.}

\begin{rubric}
{In Eastertide, replacing the Lesser Alleluia:}
\end{rubric}\par\noindent
\alleluia{Alleluia, alleluia. ℣. Now hath blossomed Jesse's rod: a Virgin bears both man and God: God restoreth peace to men, high and low are one again. Alleluia. ℣. Hail, Mary, full of grace, the Lord is with thee: blessed art thou among women. Alleluia.}

\readingcitation{Gospel}{Luke 11:27}
\lett{A}{t that time:} As Jesus spake to the multitudes, a certain woman of the company lifted up her voice, and said unto him: Blessed is the womb that bare thee, and the paps which thou hast sucked. But he said: Yea rather, blessed are they that hear the word of God, and keep it.

\offertory{Hail, Mary, full of grace; the Lord is with thee: blessed art thou among women, and blessed is the fruit of thy womb (Alleluia).}

\secret
\lett{T}{hrough} thy mercy, O Lord, and the intercession of blessed Mary ever Virgin, may this oblation avail for our prosperity and peace, both now and for ever. Through.

\communion{Blessed is the womb of the Virgin Mary, that bore the Son of the everlasting Father (Alleluia).}

\postcommunion
\lett{G}{rant,} we beseech thee, O Lord: that we, who have received these aids to our salvation, may at all times and in all places be protected through the advocacy of blessed Mary ever Virgin; in whose honour we have made these offerings to thy majesty. Through.

%MANUAL ADJUSTMENT:
\clearpage
\subby{Office of Saint Mary on Saturday}\label{BVMSaturday}
\fancyhead[RO,LE]{\textit{Saturday BVM}}

\begin{rubric}
	The Hymns are as in the Common of the Blessed Virgin Mary (p. \pageref{CommonBVM}).
\end{rubric}

\subbysub{Through the Year}

\begin{inhead}
	Evensong
\end{inhead}

℣. Full of grace are thy lips.

℟. Because God hath blessed thee for ever.

\antiphon{Mag.}{O blessed Mother {\dag} and spotless Virgin, thou glorious Queen of the world, intercede for us to the Lord.}

\lett{G}{rant,} we beseech thee, O Lord God, that we thy servants may enjoy perpetual health of mind and of body: and, at the glorious intercession of blessed Mary ever Virgin, be delivered from present sadness, and rejoice in everlasting gladness. Through.
\needspace{4\baselineskip}
\lett{D}{efend} us, we beseech thee, O Lord, from all perils of mind and body: and at the intercession of blessed Joseph, of thy blessed Apostles Peter and Paul, of blessed \textit{N.} and of all Saints, graciously bestow upon us both peace and safety: that all adversity and error being done away, thy Church may serve thee in untroubled freedom. Through.

\begin{inhead}
	Mattins
\end{inhead}

℣. Blessed art thou amongst women.

℟. And blessed is the fruit of thy womb.

\antiphon{Ben.}{O ever blessed Mother of God, {\dag} Mary ever Virgin, temple of the Godhead, hallowed shrine of the Holy Spirit, thou only, above all others, wast acceptable to our Lord Jesus Christ: pray for the people, entreat for the clergy, intercede for all womankind vowed unto God.}

\lett{G}{rant,} we beseech thee, O Lord God, that we thy servants may enjoy perpetual health of mind and of body: and, at the glorious intercession of blessed Mary ever Virgin, be delivered from present sadness, and rejoice in everlasting gladness. Through.
\needspace{4\baselineskip}
\lett{D}{efend} us, we beseech thee, O Lord, from all perils of mind and body: and at the intercession of blessed Joseph, of thy blessed Apostles Peter and Paul, of blessed \textit{N.} and of all Saints, graciously bestow upon us both peace and safety: that all adversity and error being done away, thy Church may serve thee in untroubled freedom. Through.

\subbysub{After Christmas}

\begin{rubric}
	From Saturday after the Epiphany Octave through Saturday before Purification.
\end{rubric}

\begin{inhead}
    Evensong
\end{inhead}

℣. Full of grace are thy lips.

℟. Because God hath blessed thee for ever.

\antiphon{Mag.}{Great {\dag} is the mystery of the inheritance: the womb of her that knew not man is become the temple of the Godhead: by taking flesh of her, he was no way defiled: all the nations shall gather, saying: Glory be to thee, O Lord.}

\needspace{4\baselineskip}
\lett{O}{God,} who by the virgin child-bearing of blessed Mary hast bestowed upon mankind the rewards of eternal salvation: grant, we beseech thee; that we may perceive her intercession for us, through whom we have been counted worthy to receive the author of life, Jesus Christ thy Son, our Lord. Who liveth.
\needspace{4\baselineskip}
\lett{D}{efend} us, we beseech thee, O Lord, from all perils of mind and body: and at the intercession of blessed Joseph, of thy blessed Apostles Peter and Paul, of blessed \textit{N.} and of all Saints, graciously bestow upon us both peace and safety: that all adversity and error being done away, thy Church may serve thee in untroubled freedom. Through.

\begin{inhead}
	Mattins
\end{inhead}

℣. Blessed art thou amongst women.

℟. And blessed is the fruit of thy womb.

\antiphon{Ben.}{A great and wondrous mystery {\dag} is made known to us this day; a new thing is wrought in both natures: God is made man; that which was, remained, and that which was not, he assumed; suffering no confusion, nor yet division.}

\needspace{4\baselineskip}
\lett{O}{God,} who by the virgin child-bearing of blessed Mary hast bestowed upon mankind the rewards of eternal salvation: grant, we beseech thee; that we may perceive her intercession for us, through whom we have been counted worthy to receive the author of life, Jesus Christ thy Son, our Lord. Who liveth.
\needspace{4\baselineskip}
\lett{D}{efend} us, we beseech thee, O Lord, from all perils of mind and body: and at the intercession of blessed Joseph, of thy blessed Apostles Peter and Paul, of blessed \textit{N.} and of all Saints, graciously bestow upon us both peace and safety: that all adversity and error being done away, thy Church may serve thee in untroubled freedom. Through.

\subbysub{Eastertide}

\begin{inhead}
	Evensong
\end{inhead}

℣. Full of grace art thy lips, alleluia.

℟. Because God hath blessed thee for ever, alleluia.

\antiphon{Mag.}{O Queen of heaven, {\dag} be joyful, alleluia; Because he whom so meetly thou barest, alleluia, Hath arisen, as he promised, alleluia: Pray for us to the Father, alleluia.}

\needspace{4\baselineskip}
\lett{G}{rant,} we beseech thee, O Lord God, that we thy servants may enjoy perpetual health of mind and of body: and, at the glorious intercession of blessed Mary ever Virgin, be delivered from present sadness, and rejoice in everlasting gladness. Through.
\needspace{4\baselineskip}
\lett{O}{God,} who didst will that thy Son should for us undergo the burden of the Cross, that he might deliver us from the power of the enemy: grant unto us thy servants that we may attain unto the joy of his resurrection. Through the same.

\begin{inhead}
	Mattins
\end{inhead}

\begin{rubric}
	Mattins is as in Evensong.
\end{rubric}

\subby{Seasonal Prayers}
\fancyhead[RO,LE]{\textit{Seasonal Prayers}}

   \subbysub{Of St. Mary in Advent}\label{SPMaryInAdvent}
   \begin{inhead}
   	Collect
   \end{inhead}
   \lett{O}{God,} who wast pleased that thy Word should take flesh of the womb of the blessed Virgin Mary at the message of an Angel: grant to thy humble servants; that we, who believe her to be truly the Mother of God, may be aided by her intercession in thy sight. Through the same.
  \begin{inhead}
  	Secret
  \end{inhead}
   \lett{W}{e} beseech thee, O Lord, stablish in our hearts the mysteries of thy true religion: that, as we believe thy Son, conceived of a Virgin, to be very God and very man; so by the power of his life-giving resurrection we may be found worthy to attain unto everlasting felicity. Through the same.
   \begin{inhead}
   	Postcommunion
   \end{inhead}
   \lett{W}{e} beseech thee, O Lord, pour thy grace into our hearts: that, as we have known the incarnation of thy Son Christ by the message of an Angel; so by his cross and passion we may be brought unto the glory of his resurrection. (Through the same.)


    \subbysub{Of St. Mary after Christmas}\label{SPMaryPostChristmas}
   \begin{inhead}
   	Collect
   \end{inhead}
    \lett{O}{God,} who by the virgin child-bearing of blessed Mary hast bestowed upon mankind the rewards of eternal salvation: grant, we beseech thee; that we may perceive her intercession for us, through whom we have been counted worthy to receive the author of life, Jesus Christ thy Son, our Lord. (Who liveth.)
  \begin{inhead}
  	Secret
  \end{inhead}
    \lett{T}{hrough} thy mercy, O Lord, and the intercession of blessed Mary ever Virgin, may this oblation avail for our prosperity and peace both now and for ever. (Through.)
   \begin{inhead}
   	Postcommunion
   \end{inhead}
    \lett{M}{ay} this communion, O Lord, cleanse us from guilt: and at the intercession of blessed Mary the Virgin Mother of God, make us partakers of thy heavenly healing. (Through the same.)
    
    
    \subbysub{Of St. Mary in Eastertide}\label{SPMaryInEaster}
   \begin{inhead}
   	Collect
   \end{inhead}
    \lett{G}{rant,} we beseech thee, O Lord God, that we thy servants may enjoy perpetual health of mind and of body: and, at the glorious intercession of blessed Mary ever Virgin, be delivered from present sadness, and rejoice in everlasting gladness. (Through.)
  \begin{inhead}
  	Secret
  \end{inhead}
    \lett{T}{hrough} thy mercy, O Lord, and the intercession of blessed Mary ever Virgin, may this oblation avail for our prosperity and peace, both now and ever. (Through.)
   \begin{inhead}
   	Postcommunion
   \end{inhead}
    \lett{G}{rant,} we beseech thee, O Lord: that we who have received these aids to our salvation may at all times and in all places be protected by the advocacy of blessed Mary ever Virgin; in whose honour we have made these offerings to thy majesty. (Through.)\\

    \subbysub{Of Saints}\label{SPSaints}
   \begin{inhead}
   	Collect
   \end{inhead}
    \lett{D}{efend} us, O Lord, we beseech thee, from all dangers of mind and of body: and at the intercession of the blessed and glorious Mother of God, Mary ever Virgin, with blessed Joseph, thy blessed Apostles Peter and Paul, blessed \emph{N.}, and all the Saints, graciously grant us health and peace; that all adversities and errors being done away, thy Church may serve thee in freedom and quietness. (Through the same.)
  \begin{inhead}
  	Secret
  \end{inhead}
    \lett{G}{raciously} hear us, O God our Saviour: and by the power of this sacrament defend us from all enemies of mind and of body; granting us grace in this world, and glory in that which is to come. (Through.)
   \begin{inhead}
   	Postcommunion
   \end{inhead}
    \lett{W}{e} beseech thee, O Lord, that the gift now offered in this divine sacrament may cleanse and defend us: that at the intercession of blessed Mary the Virgin Mother of God, of blessed Joseph, of thy blessed Apostles Peter and Paul, of blessed \emph{N.} and of all the Saints; it may set us free from every perverse way, and deliver us from all adversities. (Through the same.)
    
   \subbysub{Against the Persecutors of the Church}\label{SPAgainst}
   \begin{inhead}
   	Collect
   \end{inhead}
   \lett{W}{e} beseech thee, O Lord, mercifully to hear the prayers of thy Church: that all adversities and errors being done away, she may serve thee in freedom and quietness. Through.
  \begin{inhead}
  	Secret
  \end{inhead}
   \lett{D}{efend} us, O Lord, who wait upon thy mysteries: that we, cleaving fast to things heavenly, may serve thee both in body and soul. Through.
   \begin{inhead}
   	Postcommunion
   \end{inhead}
   \lett{W}{e} beseech thee, O Lord our God: that we, whom thou makest to rejoice in the partaking of heavenly things, may by thee be defended against all earthly perils. Through.
   
   \subbysub{For the Chief Bishop}\label{SPChiefBishop}
   \begin{inhead}
   	Collect
   \end{inhead}
   \lett{O}{God,} the pastor and ruler of all the faithful, mercifully look upon thy servant \emph{N.}, whom thou hast been pleased to set as pastor over thy Church; grant him, we beseech thee, to be in word and conversation a wholesome example to the people committed to his charge; that he with them may attain unto everlasting life. Through.
  \begin{inhead}
  	Secret
  \end{inhead}
   \lett{L}{ook} favourably, we beseech thee, O Lord, upon the gifts which we offer: and guide with thy continual protection thy servant \textit{N.}, whom thou hast been pleased to set as pastor over thy Church. Through.
   \begin{inhead}
   	Postcommunion
   \end{inhead}
   \lett{W}{e} beseech thee, O Lord, that the divine sacrament which we have here received may be our protection: and together with the flock entrusted to his charge ever save and defend thy servant \textit{N.}; whom thou hast been pleased to set as pastor over thy Church. Through.
   
   \subbysub{Of the Holy Ghost}\label{SPHolyGhost}
   \begin{inhead}
   	Collect
   \end{inhead}
   \lett{G}{od,} who didst teach the hearts of thy faithful people, by the sending to them the light of thy Holy Spirit: grant us by the same Spirit to have a right judgment in all things; and evermore to rejoice in his holy comfort. (Through . . . in the unity of the same Holy Spirit.)
  \begin{inhead}
  	Secret
  \end{inhead}
   \lett{S}{anctify,} we beseech thee, O Lord, the gifts which we offer: and cleanse our hearts by the enlightening of the Holy Spirit. (Through . . . in the unity of the same Holy Spirit.)
   \begin{inhead}
   	Postcommunion
   \end{inhead}
   \lett{P}{our} thy Holy Spirit upon us, O Lord, and cleanse our hearts: that they may be made fruitful by the inward sprinkling of his dew. (Through . . . in the unity of the same Holy Spirit.)

    \subbysub{For the Living and Departed}\label{SPLivingDeparted}
   \begin{inhead}
   	Collect
   \end{inhead}
   \lett{A}{lmighty} and everlasting God, who hast dominion both of the living and of the dead, and hast mercy upon all who thou foreknowest will be thine in faith and works: we humbly beseech thee; that all those for whom we are minded to pour forth our prayers, whether in this present world they still be held in the flesh, or being delivered from the body have passed into that which is to come, may at the intercession of all thy Saints obtain of thy bountiful goodness the remission of all their sins. Through.
  \begin{inhead}
  	Secret
  \end{inhead}
   \lett{O}{God,} to whom alone is known the number of the elect, whom thou hast appointed unto heavenly felicity: grant, we beseech thee: that at the intercession of all thy Saints, the names of all those who have been commended to our prayers, and of all thy faithful people, may ever remain written in the book of those that are predestined to everlasting blessedness. Through.
   \begin{inhead}
   	Postcommunion
   \end{inhead}
   \lett{A}{lmighty} and merciful God, cleanse us, we pray thee, by the sacraments which we have received: and, at the intercession of all thy Saints, grant; that this thy sacrament may not bring upon us guilt to our punishment but pardon to our salvation: that it may cleanse away sins, that it may confirm the feeble, that it may defend us against every danger in this world: that it may obtain for thy faithful people both living and dead the remission of all their sins. Through.

%LATIN:
%\fancyhead[C]{{\LARGE Daily Office}}
\fancyhead[RO,LE]{\textit{Fore-Office}}
\section{Fore-Office}
\selectlanguage{latin}
\begin{secrubric}
    In principio Matutinarum Minister clara voce legat unam aliquam sive plures ex Sacræ Scripturæ Sententiis quæ consequuntur. Et postea dicat quod infra has Sententias scriptum est.
\end{secrubric}
\begin{secrubric}
    %This rubric is very intentionally placed in obedience (and good sense) of the Church that prayers to the saints ought not to be omitted. Without this rubric, the bare minimum could result in an Office without any intercession of the saints. Given how quickly parishes (sadly) and laymen (sensibly) will scale down their liturgy to the bare minimum, it seems necessary to have this rubric as a safeguard.
    On any day, save a Day of Fasting or Abstinence, or on any day when the Litany or Holy Communion is immediately to follow, the Minister may, at his discretion, pass at once from the Sentences to the Lord's Prayer and Angelic Salutation, which may never be omitted.
\end{secrubric}
\begin{multicols}{2}
\begin{inhead}
General
\end{inhead}\noindent
%Provided here are the original opening sentences of the 1662 BCP. It is a difficult balance, as discussed in the Preface, balancing this American Prayer Book with the needs, desires, and practices of Western Rite people and parishes. The Opening Sentences themselves are very English. The English Reformation saw a trend of making changeable parts of the liturgy unchangeable. This is seen excellently in the replacement of seasonal antiphons with standardised collects.
%The American tradition, along with the Canadian and elsewhere, eventually saw an attempt to make the liturgy more changeable according to the liturgical season, seen here in the Sentences. In America, this involved moving the standard Sentences (of a penitential character) to Lent.
%The Opening Sentences come from the Sacerdotal Versicles which introduced Lauds. Cranmer accommodated this to the Fore-Office's penitential nature, which was itself both a replacement for the introduction to the Hours and also good sense to begin worship and sacrifice with a confession of sins unto almighty God. Later on, notably in the American 1928 and Proposed 1928, the Confession became optional and its character as an introduction to the Office as a whole was re-emphasised, notably with seasonal collects. While the Prayer Book Tradition gets its distinctive beauty and efficacy by reduced options, propers, and complicated rubrics, it seems fitting to continue this emphasis by supplying from the Sarum the proper Opening Sentences for general kinds of Feast Days and for the Feast Days in the Proper.
%Given its pedigry and that other Western Rite material has provided a `General' section of the Opening Sentences, and this is actively used in AWRV parishes, I do not think it an innovation to provide it here. And given the beauty and interest among AWRV clergy and laymen in these aesthetic elements of the English tradition, it makes sense to me to make these the `General' Sentences. I will say, however, by way of suggestion and not rubric, that these should not be used if the Confession \& Absolution are omitted but rather the seasonal prayers.
	Si impius egerit p{\oe}nitentiam pro omnibus peccatis suis qu{\ae} operatus est, et custodierit omnia pr{\ae}cepta mea, et fecerit judicium et justiciam, vita vivet et non morietur. Omnium iniquitatum ejus quas operatus est, non recordabor: dicit Dominus.\vr{Ezech. 18:27}
	
	Iniquitatem meam agnosco, et peccatum meum contra me est semper.\vr{Psa. 51:3}
	
	Averte faciem tuam a peccatis meis: et omnes iniquitates meas dele.\vr{Psa. 51:9}
	
	Scindite corda vestra, et non vestimenta vestra, et convertimini ad Dominum Deum vestrum, quia benignus et misericors est, patiens et mult{\ae} clementi{\ae}, qui se ab inferendo malo contineat.\vr{J{\oe}l 2:13}
	
	Tui Domini Dei nostri est misericordia et propiciatio, quia recessimus a te, et non audivimus vocem Domini Dei nostri, ut ambularemus in lege ejus.\vr{Dan 9:9-10}
	
	Corripe nos, Domine, veruntamen in judicio et non in furore tuo, ne forte ad nihilum redigas nos.\vr{Jer 10:24; Ps 6:1}
	
	P{\oe}nitentiam agite; appropinquat enim regnum c{\oe}lorum.\vr{Mt 3:2}
	
	Surgam, et ibo ad patrem meum, et dicam ei: Pater, peccavi in c{\oe}lum et coram te. Jam non sum dignus vocari filius tuus.\vr{Lk 15:18-19}
	
	Non intres in judicium cum servo tuo, Domine, quia non justificabitur in conspectu tuo omnis vivens.\vr{Ps 143:2}
	
	Si nos peccati expertes esse dicimus, fallimus nos ipsos, nec est in nobis veritas.\vr{1 Jn 1:8-9}
\begin{inhead}
Morning
\end{inhead}\noindent
    %From 1928 Morning Prayer:
    The \textsc{Lord} is in his holy temple: let all the earth keep silence before him.\vr{Hab 2:20}\par
    I was glad when they said unto me, We will go into the house of the \textsc{Lord}.\vr{Ps 122:1}\par
    Let the words of my mouth, and the meditation of my heart, be alway acceptable in thy sight, O \textsc{Lord}, my strength and my redeemer.\vr{Ps 19:14}\par
    O send out thy light and thy truth, that they may lead me, and bring me unto thy holy hill, and to thy dwelling.\vr{Ps 43:3}\par
    Thus saith the high and lofty One that inhabiteth eternity, whose name is Holy; I dwell in the high and holy place, with him also that is of a contrite and humble spirit, to revive the spirit of the humble, and to revive the heart of the contrite ones.\vr{Is 57:15}\par
    The hour cometh, and now is, when the true worshippers shall worship the Father in spirit and in truth: for the Father seeketh such to worship him.\vr{Jn 4:23}\par
    Grace be unto you, and peace, from God our Father, and from the Lord Jesus Christ.\vr{Phil 1:2}\par
\begin{inhead}
Evening
\end{inhead}\noindent
%From 1928 Evening Prayer:
    The \textsc{Lord} is in his holy temple: let all the earth keep silence before him.\vr{Hab 2:20}\par
    \textsc{Lord}, I have loved the habitation of thy house, and the place where thine honour dwelleth.\vr{Ps 26:8}\par
    Let my prayer be set forth in thy sight as the incense; and let the lifting up of my hands be an evening sacrifice.\vr{Ps 141:2}\par
    O worship the \textsc{Lord} in the beauty of holiness; let the whole earth stand in awe of him.\vr{Ps 96:9}\par
    Let the words of my mouth, and the meditation of my heart, be alway acceptable in thy sight, O \textsc{Lord}, my strength and my redeemer.\vr{Ps 19:14-15}\par
%From 1918 Morning Prayer:
    Seek ye the Lord while he may be found, call ye upon him while he is near: let the wicked forsake his way and the unrighteous man his thoughts: and let him return unto the Lord, and he will have mercy upon him: and to our God, for he will abundantly pardon.\vr{Is 55:6-7}
%For the seasonal schema, the first option is from the 1928 Morning, the second from 1928 Evening Prayer, and the third from either the 1962 or 1918. If the Canadian is the same as the American, another is chosen.

\begin{inhead}
Advent
\end{inhead}\noindent
    Repent ye, for the Kingdom of heaven is at hand.\vr{Mt 3:2}
    \par
    Watch ye, for ye know not when the master of the house cometh, at even, or at midnight, or at the cock-crowing, or in the morning: lest coming suddenly he find you sleeping.\vr{Mk 13:35-6}
    \par
    Prepare ye the way of the \textsc{Lord}, make straight in the desert a highway for our God.\vr{Is 40:3}

\begin{inhead}
Christmas
\end{inhead}\noindent
    Behold, I bring you good tidings of great joy, which shall be to all people. For unto you is born this day in the city of David a Saviour, which is Christ the Lord.\vr{Lk 2:10-11}
    \par
    Behold, the tabernacle of God is with men, and he will dwell with them, and they shall be his people, and God himself shall be with them, and be their God.\vr{Rev 21:3}
    \par
    %From the Proposed 1928:
    Herein was the love of God manifested in us, that God hath sent his only begotten Son into the world, that we might live through him. \vr{1 Jn 4:9}
    \par

\begin{inhead}
Epiphany
\end{inhead}\noindent
%From 1928 Morning Prayer
    From the rising of the sun even unto the going down of the same my Name shall be great among the Gentiles; and in every place incense shall be offered unto my Name, and a pure offering: for my Name shall be great among the heathen, saith the \textsc{Lord} of hosts.\vr{Mal 1:11}
    \par
%From 1928 Evening Prayer
    And the Gentiles shall come to thy light, and kings to the brightness of thy rising.\vr{Is 60:3}\par
%From 1962 Morning Prayer
    The earth shall be filled with the knowledge of the glory of the \textsc{Lord}, as the waters cover the sea.\vr{Hab 2:14}
\begin{inhead}
Lent
\end{inhead}\noindent
%To avoid overlap, we are avoiding any Lenten Sentences which are already provided in `General'.
%From Sarum (Lent I & II):
    He hath delivered me from the snare of the hunter, and from the sharp word.\vr{Ps 91:3 LXX}\par
    He shall deliver thee from the snare of the hunter. And from the noisome pestilence.\vr{Ps 91:3}

\begin{inhead}
Passiontide
\end{inhead}\noindent
%From 1962 Canadian:
    God commendeth his love toward us, in that, while we were yet sinners, Christ died for us.\vr{Rom 5:8}\par
%From the English Proposed 1928 BCP
    Is it nothing to you, all ye that pass by? behold, and see if there be any sorrow like unto my sorrow which is done unto me, wherewith the \textsc{Lord} hath afflicted me.\vr{Lam 1:12}\par
%From Sarum (Passion Sunday):
Draw nigh unto my soul and save it. O deliver me because of mine enemies.\vr{Ps 69:18}

\begin{inhead}
Eastertide
\end{inhead}\noindent
    He is risen. The Lord is risen indeed.\vr{Mk. 16:6}\par
    Thanks be to God, which giveth us the victory through our Lord Jesus Christ.\vr{1 Cor. 15:57}\par
    If ye then be risen with Christ, seek those things which are above, where Christ sitteth on the right hand of God.\vr{Col 3:1}

\begin{inhead}
Ascensiontide
\end{inhead}\noindent
    Seeing that we have a great High Priest, that is passed into the heavens, Jesus the Son of God, let us come boldly unto the throne of grace, that we may obtain mercy, and find grace to help in time of need.\vr{Heb 4:14,16}\par
    
    Christ is not entered into the holy places made with hands, which are the figures of the true; but into heaven itself, now to appear in the presence of God for us.\vr{Heb 9:24}\par

    %From Sarum (Ascension Thursday):
    I ascend to my Father, and your Father. And to my God, and your God. Alleluia.\vr{Jn 20:17}

\begin{inhead}
Whitsuntide
\end{inhead}\noindent
    Ye shall receive power, after that the Holy Ghost is come upon you: and ye shall be witnesses unto me both in Jerusalem, and in all Jud{\ae}a, and in Samaria, and unto the uttermost part of the earth.\vr{Acts 1:8}\par
    
    There is a river, the streams whereof shall make glad the city of God, the holy place of the tabernacles of the Most High.\vr{Ps 46:4}\par

    %From 1962 Morning Prayer (which uses the ASV here, interestingly):
    The love of God hath been shed abroad in our hearts through the Holy Spirit which was given unto us.\vr{Rom 5:5}
\begin{inhead}
Trinity Sunday
\end{inhead}\noindent
    Holy, holy, holy, Lord God Almighty, which was, and is, and is to come.\vr{Rev 4:8}
    \par
    Holy, holy, holy, is the \textsc{Lord} of hosts: the whole earth is full of his glory. \vr{Is 6:3}
    \par
%From 1962 Morning Prayer:
    God is love; and he that abideth in love abideth in God and God in him.\vr{1 Jn 4:16}

%There has been an attempt to avoid novelty in this Prayer Book, especially in the Sentences where novelty has become common. However, the trajectory of the American Prayer Book tradition seems both good and clear: seasonal Sentences based on the feast day. Therefore, the Opening Sentences (Sacerdotal Versicles) from the Sarum Office's Common is thus provided.
\begin{inhead}
Apostle
\end{inhead}\noindent
Thou hast given an heritage unto those that fear thy Name, O Lord.\vr{Ps 61:5}

\begin{inhead}
Martyr
\end{inhead}\noindent
Thou hast crowned him, O Lord, with glory and worship.\vr{Ps 8:5}

\begin{inhead}
Bishop, Confessor, or Doctor
\end{inhead}\noindent
The righteous shall flourish like a palm-tree: and shall spread abroad like a cedar in
Libanus.\vr{Ps 92:11}

\begin{inhead}
Abbot or Monk
\end{inhead}\noindent
O ye holy and humble men of heart, bless ye the Lord: praise and exalt him above all for ever.\vr{Pr. of Az. 65}

\begin{inhead}
Virgin
\end{inhead}\noindent
In thy grace and in thy beauty, go forth, ride prosperously and reign.\vr{Ps 45:4}

\begin{inhead}
Matron
\end{inhead}\noindent
God is in the midst of her, therefore shall she not be removed : God shall help her, and that right early.\vr{Ps 46:5}

\begin{inhead}
Blessed Virgin Mary
\end{inhead}\noindent
Thou art the holy Mother of God, O Mary, ever Virgin.\vr{cf. Is 7:14}

\begin{inhead}
Blessed Sacrament
\end{inhead}\noindent
I have eaten my honeycomb with my honey. I have drunk my wine with my milk. (Alleluia.)\vr{Song 5:1}

\begin{inhead}
Dedication of a Church
\end{inhead}\noindent
%From the Sarum Versicle for the Anniversary of the Dedication of a Church
	My house shall be called the house of prayer.\vr{Mt. 21:13}
\end{multicols}
\subby{Exhortation to Penitence}
\lett{D}{ilectissimi} fratres, Sacra Scriptura multis in locis nos commonefacit, ut multiplicia peccata nostra et iniquitatem nostram agnoscamus et confiteamur; nec ea in conspectu Dei Omnipotentis, Patris nostri cœlestis, audeamus dissimulare; sed animo demisso, humili, contrito, et obedienti confiteamur; quatenus, per infinitam bonitatem et misericordiam ejus, indulgentiam consequamur. Et quanquam in omni tempore peccata nostra coram Deo humiliter agnoscenda sunt, tunc tamen præcipue hoc agendum est cum in publicum cœtum ideo convocati sumus, ut pro magnis ab eo collatis beneficiis gratias referamus, ut ejus laudes dignas annuntiemus, et sanctissimum ejus Verbum audiamus, et ut quicquid animorum status aut corporis requirit ab eo postulemus. Quapropter vos omnes, qui adestis, rogo et obsecro ut, puro corde, et voce humili, ad thronum gratiæ cœlestis mecum accedatis, ita me præeunte dicentes:
\begin{inhead}
    or,
\end{inhead}
\par%
\noindent%
%From the English Proposed 1928. The American tradition is horribly confused on the point of the Exhortation to Penitence (even more than its general confusion regarding penitence!). Each edition has a different approach. It is understandable that the same long exhortation is undesirable. However, removing the exhortation for a short sentence seems to betray a loss of the plot. Not only does this shorter exhortation meet the challenge sought by the 1928 (more variety, not as long as the longer exhortation but not too short) even better, but it incorporates a greater sense of the communion of the saints both on earth and heaven. It lays out clearly and briskly the purposes of worship, the Office as the Morning and Evening sacrifice offered by every Christian, and is simply beautiful. It is currently in wide use among Christians in the UK and even in these United States. For these reasons, and its presence in the Prayer Book Tradition of today and yesteryear, I include it.
\needspace{3\baselineskip}
\lett{B}{eloved,} we are come together in the presence of Almighty God and of the whole company of heaven to offer unto him through our Lord Jesus Christ our worship and praise and thanksgiving; to make confession of our sins; to pray, as well for others as for ourselves, that we may know more truly the greatness of God's love and shew forth in our lives the fruits of his grace; and to ask on behalf of all men such things as their well-being doth require. Wherefore let us kneel in silence, and remember God's presence with us now.

\begin{rubric}%It seems too bold to omit the possibility of the novel introductory sentence entirely, so a longer and a shorter exhortation are provided with the sentence provided in the rubric.
	If the exhortation be omitted, to be said in its place: \emph{Let us humbly confess our sins unto Almighty God.}
\end{rubric}

\subby{General Confession}
\begin{rubric}
    To be said by the whole Congregation, after the Minister, all kneeling.
\end{rubric}
%The American edition corrupts the grammar of this prayer (replacing `them' with `those'). It's a shame it would be too jarring to change. The same is even more true for the Lord's Prayer.
\needspace{3\baselineskip}

\lett{A}{lmighty} and most merciful Father; We have erred, and strayed from thy ways like lost sheep. We have followed too much the devices and desires of our own hearts. We have offended against thy holy laws. We have left undone those things which we ought to have done; And we have done those things which we ought not to have done; And there is no health in us. But thou, O Lord, have mercy upon us, miserable offenders. Spare thou those, O God, who confess their faults. Restore thou those who are penitent; According to thy promises declared unto mankind in Christ Jesus our Lord. And grant, O most merciful Father, for his sake; That we may hereafter live a godly, righteous, and sober life, To the glory of thy holy Name. Amen.
\subby{Declaration of Absolution}
\begin{rubric}
    To be made by the Priest alone, standing; the People still kneeling.
\end{rubric}
\begin{rubric}
    But \textsc{Note}, That the Priest, at his discretion, may use, instead of what follows, the Absolution from the Order for the Holy Communion.
\end{rubric}
\lett{A}{lmighty} God, the Father of our Lord Jesus Christ, who desireth not the death of a sinner, but rather that he may turn from his wickedness and live, hath given power, and commandment, to his Ministers, to declare and pronounce to his people, being penitent, the Absolution and Remission of their sins. He pardoneth and {\ding{64}} absolveth all those who truly repent, and unfeignedly believe his holy Gospel.
\par
Wherefore let us beseech him to grant us true repentance, and his Holy Spirit, that those things may please him which we do at this present; and that the rest of our life hereafter may be pure and holy; so that at the last we may come to his eternal joy; through Jesus Christ our Lord. \textit{Amen.}
\begin{inhead}
    or,
\end{inhead}
%From 1928 BCP EP.
\lett{T}{he} Almighty and merciful Lord grant you Absolution {\ding{64}} and Remission of all your sins, true repentance, amendment of life, and the grace and consolation of the Holy Spirit. Amen.
%The American tradition does not supply a lay `absolution', but in our times it is greatly needful. Therefore the 1662 is supplied instead of borrowing from Compline. And only one provided. We don't need too many options.
\begin{rubric}
    If no Priest be present, the Minister saying the Service shall read the following, that Minister and the People still kneeling.
\end{rubric}
\lett{G}{rant,} we beseech thee, merciful Lord, to thy faithful people pardon and peace, that they may be cleansed from all their sins, and serve thee with a quiet mind. Through Jesus Christ thy Son our Lord; who liveth and reigneth with thee in the unity of the Holy Spirit, one God, world without end. \textit{Amen.}

\subby{Lord's Prayer \& Angelic Salutation}
\begin{rubric}
    Then the Minister shall kneel, and say the Lord's Prayer \& Angelic Salutation; the People still kneeling, and repeating them with him, both here, and wheresoever else they are used in Divine Service.
\end{rubric}
%The traditional Office, before the English liturgical reform, was prefaced with the Pater Noster, Ave Maria, and Credo. The addition of the Ave Maria not only restores the Catholic and Western practice but also ensures that the intercession of the saints will be present in every Office (not unrubrically truncated, at least).
\needspace{3\baselineskip}
\lett{P}{ater} noster, qui es in cœlis, Sanctificetur Nomen tuum. Adveniat regnum tuum. Fiat voluntas tua, Sicut in cœlo, et in terra. Panem nostrum quotidianum da nobis hodie. Et dimitte nobis debita nostra, Sicut et nos dimittimus debitoribus nostris. Et ne nos inducas in tentationem; Sed libera, nos a malo: Quia tuum est regnum, Potentia, et gloria, In sæcula sæculorum. Amen.
\par\noindent
\needspace{3\baselineskip}
\lett{A}{ve} María, grátia plena; Dóminus tecum: benedícta tu in muliéribus, et benedíctus fructus ventris tui Jesus. Sancta María, Mater Dei, ora pro nobis peccatóribus, nunc et in hora mortis nostræ. Amen.

%\fancyhead[RO,LE]{\textit{Mattins}}
\section{Ordo Matutinarum}
℣. Dómine, {\ding{61}} lábia nostra apéries.

℟. Et os nostrum annuntiábit laudem tuam.

℣. Deus, {\ding{64}} in adjutórium nostrum inténde.

℟. Dómine, ad adjuvándum nos festína.

\begin{rubric}
    Hic, omnibus sese erigentibus, dicat Minister.
\end{rubric}

℣. Glória Patri, et Fílio, * et Spirítui Sancto:

℟. Sicut erat in princípio, et nunc, et semper, * et in sǽcula s{\ae}culórum. Amen.

℣. Laudáte Dominum.

℟. Sit Nomen Dómini Benedíctum.

\begin{rubric}
    Tunc dicatur vel cantetur Psalmus hic sequens nisi in Die Sancto Paschæ, cui alia assignatur Antiphona. Item in die undevicesimo cujusque mensis non hoc loco, sed in ordinario Psalmorum cursu, est legendus.
\end{rubric}
\begin{rubric}
    But \textsc{Note}, That on Ash Wednesday and Good Friday the \emph{Venite} may be omitted.
\end{rubric}
\begin{rubric}
    On the days hereafter named, immediately before and after the \emph{Venite} may be sung or said,
\end{rubric}
\par\noindent
\inv{Advent:} Our King and Saviour draweth nigh; * O come, let us adore him.
\par\noindent
\inv{Christmastide:} Alleluia. Unto us a child is born; * O come, let us adore him. Alleluia.
\par\noindent
\inv{Epiphanytide \& Transfiguration:} The Lord hath manifested forth his glory; * O come, let us adore him.
\par\noindent
%Adapted from AOB:
\inv{Septuagesimatide:} Let us come before the presence of the Lord with thanksgiving;  * O come, let us adore him.
\par\noindent
%Adapted from the 1962 BCP:
\inv{Lent:} The goodness of God leadeth to repentance; * O come, let us adore him.
\par\noindent
\inv{Passiontide:} Christ our Lord became obedient unto death; * O come, let us adore him.
\par\noindent
\inv{Eastertide outside the Octave:} Alleluia. The Lord is risen indeed; * O come let us adore him. Alleluia.
\par\noindent
\inv{Ascensiontide:} Alleluia. Christ the Lord ascended into heaven; * O come, let us adore him. Alleluia.
\par\noindent
\inv{Whitsuntide:} Alleluia. The Spirit of the Lord filleth the world; * O come, let us adore him. Alleluia.
\par\noindent
\inv{Trinitytide:} Father, Son, and Holy Ghost, One God; * O come, let us adore him.
\par\noindent
\inv{Feasts of Our Lord \& Our Lady:} The Word was made flesh; * O come, let us adore him.
\par\noindent
%Due to the wider use of Feast Days in the AWRV, it is difficult to decide how to adapt this rubric (Feast Days which have a Proper Epistle and Gospel). Given that Semidouble Feast Days almost never occur, and moving the bar so low to Simplex or Memorial would destroy the desired effect, Semidouble seems to be the best rank. For it would not cast too wide of a net: effectively capturing all Doubles and their Octaves, except when the Octave is meant to not be noticed intensely.
\inv{Other Feast Days of Semidouble or higher:} The Lord is glorious in his saints; * O come, let us adore him.
%For text of the psalms, the Latin Monastic Diurnal is used.
\subby{Venite, exultemus Domino.}

\lett{V}{en\smash{í}te}, exsultémus Dómino, jubilémus Deo, salutári nostro:\par
\secondline{Pr{\ae}occupémus fáciem ejus in confessióne, et in psalmis jubilémus ei.}
\thirdline{Quóniam Deus magnus Dóminus, et Rex magnus super omnes deos, quóniam non repéllet Dóminus plebem suam:}
Quia in manu ejus sunt omnes fines terr{\ae}, et altitúdines móntium ipse cónspicit.\par
Quóniam ipsíus est mare, et ipse fecit illud, et áridam fundavérunt manus ejus\par
Veníte, adorémus, et procidámus ante Deum: plorémus coram Dómino, qui fecit nos,\par
Quia ipse est Dóminus, Deus noster; nos autem pópulus ejus, et oves páscu{\ae} ejus.\par
Hódie, si vocem ejus audiéritis, nolíte obduráre corda vestra, sicut in exacerbatióne secúndum diem tentatiónis in desérto: ubi tentavérunt me patres vestri, probavérunt et vidérunt ópera mea.
\par
Quadragínta annis próximus fui generatióni huic, et dixi; Semper hi errant corde, ipsi vero non cognovérunt vias meas: quibus jurávi in ira mea; Si introíbunt in réquiem meam.

\begin{rubric}
    %1928 BCP Rubric (permission to not use Gloria Patri, except at the end of all of the Psalms, is removed):
    Then shall follow a Portion of the Psalms, according to the Use of this Church. And at the end of every Psalm, and liksewise at the end of the \emph{Venite}, daily Old Testament Canticle, and \emph{Benedictus}, shall be sung or said the \emph{Gloria Patri}.
\end{rubric}
℣. Glória Patri, et Fílio, * et Spirítui Sancto:\par
℟. Sicut erat in princípio, et nunc, et semper, * et in sǽcula s{\ae}culórum. Amen.

\par\noindent
	\centerline{\rule{0.5\textwidth}{0.4pt}}
\par\noindent
\begin{rubric}
    Verses 8-11 of the \emph{Venite} may be replaced with the following.
\end{rubric}
Tóllite hóstias, et introíte in átria ejus: * adoráte Dóminum in átrio sancto ejus.\par
Judicábit orbem terr{\ae} in {\ae}quitáte, * et pópulos in veritáte sua.
\par\noindent
	\centerline{\rule{0.5\textwidth}{0.4pt}}
\par\noindent
\begin{rubric}
%Rubric added to permit the Sarum Invitatory Psalm.
    After the \emph{Venite}, a Hymn may be sung.
\end{rubric}
\begin{rubric}
Then shall be read the First Lesson, according to the Table or Calendar. And \textsc{Note}, That before every Lesson, the Minister shall say, \emph{Here beginneth} such a \emph{Chapter} (or \emph{Verse of} such a \emph{Chapter}) \emph{of} such a Book; and after every Lesson, \emph{Here endeth the First} (or \emph{the Second Lesson.})
\end{rubric}
\vspace{-2ex}
\subby{Te Deum laudamus}
%Reversion to English language brings the Te Deum is conformity to the Monastic Matins use.
\begin{rubric}
%Great reduction of the Office, allowed in the 1928, not reproduced. Additional Canticles supplied.
The \emph{Te Deum} is prayed on all Sundays (outside Advent \& Lent) and Feast Days %This rubric brings the Te Deum to the same standard as the other Festal canticles and brings it closer to Breviary use.
(Double or higher).\par
\textsc{Note,} The Old Testament Canticle (p. \pageref{OT}) may be said instead of the \emph{Te Deum}.
\end{rubric}
\lett{T}{e} Deum laudámus: * te Dóminum confitémur.\par
\secondline{Te {\ae}térnum Patrem * omnis terra venerátur.}
\thirdline{Tibi omnes Ángeli, * tibi C{\ae}li, et univérs{\ae} Potestátes}
Tibi Chérubim et Séraphim * incessábili voce proclámant:\par
Sanctus, Sanctus, Sanctus * Dóminus Deus Sábaoth.\par
Pleni sunt c{\ae}li et terra * majestátis glóri{\ae} tu{\ae}.\par
Te gloriósus * Apostolórum chorus,\par
Te Prophetárum * laudábilis númerus,\par
Te Mártyrum candidátus * laudat exércitus.\par
Te per orbem terrárum * sancta confitétur Ecclésia,\par
Patrem * imméns{\ae} majestátis;\par
Venerándum tuum verum * et únicum Fílium;\par
Sanctum quoque * Paráclitum Spíritum.

\lett{T}{u} Rex glóri{\ae}, * Christe.\par
\secondline{Tu Patris * sempitérnus es Fílius.}
\thirdline{Tu, ad liberándum susceptúrus hóminem: * non horruísti Vírginis úterum.}
Tu, devícto mortis acúleo, * aperuísti credéntibus regna c{\ae}lórum.\par
Tu ad déxteram Dei sedes, * in glória Patris.\par
Judex créderis * esse ventúrus. \par
Te ergo quǽsumus, tuis fámulis súbveni, * quos pretióso sánguine redemísti.\par
{\AE}térna fac cum Sanctis tuis * in glória numerári.\par
\lett{S}{alvum} fac pópulum tuum, Dómine, * et bénedic hereditáti tu{\ae}.\par
\secondline{Et rege eos, * et extólle illos usque in {\ae}térnum.}
\thirdline{Per síngulos dies * benedícimus te.}
Et laudámus nomen tuum in sǽculum, * et in sǽculum sǽculi.\par
Dignáre, Dómine, die isto * sine peccáto nos custodíre.\par
Miserére nostri, Dómine, * miserére nostri.\par
Fiat misericórdia tua, Dómine, super nos, * quemádmodum sperávimus in te.\par
In te, Dómine, sperávi: * non confúndar in {\ae}térnum.

%MANUAL ADJUSTMENT:
\vspace{-2ex}
\subby{Benedicite, omnia opera Domini.}

%MANUAL ADJUSTMENT:
\vspace{-1ex}

\begin{rubric}
The \emph{Benedicite} is prayed on all Sundays in Advent \& Lent, and on Days below Double.\par
\textsc{Note,} The Old Testament Canticle (p. \pageref{OT}) may be said instead of the \emph{Benedicite}.
\end{rubric}

\lett{B}{ened\smash{í}cite,} ómnia ópera Dómini, Dómino: * laudáte et superexaltáte eum in sǽcula.\par
\secondline{Benedícite, Ángeli Dómini, Dómino: * laudáte et superexaltáte eum in sǽcula.}
\lett{B}{ened\smash{í}cite,} c{\ae}li, Dómino:  * laudáte et superexaltáte eum in sǽcula.\par
\secondline{Benedícite, aqu{\ae} omnes, qu{\ae} super c{\ae}los sunt, Dómino: * laudáte et superexaltáte eum in sǽcula.}
\thirdline{Benedícite, omnes virtútes Dómini, Dómino: * laudáte et superexaltáte eum in sǽcula.}
Benedícite, sol et luna, Dómino: * laudáte et superexaltáte eum in sǽcula.\par
Benedícite, stell{\ae} c{\ae}li, Dómino: * laudáte et superexaltáte eum in sǽcula.\par
Benedícite, omnis imber et ros, Dómino: * laudáte et superexaltáte eum in sǽcula.\par
Benedícite, omnes spíritus Dei, Dómino: * laudáte et superexaltáte eum in sǽcula.\par
Benedícite, ignis et {\ae}stus, Dómino: * laudáte et superexaltáte eum in sǽcula.\par
Benedícite, frigus et {\ae}stus, Dómino: * laudáte et superexaltáte eum in sǽcula.\par
Benedícite, rores et pruína, Dómino: * laudáte et superexaltáte eum in sǽcula.\par
Benedícite, gelu et frigus, Dómino: * laudáte et superexaltáte eum in sǽcula.\par
Benedícite, glácies et nives, Dómino: * laudáte et superexaltáte eum in sǽcula.\par
Benedícite, noctes et dies, Dómino: * laudáte et superexaltáte eum in sǽcula.\par
Benedícite, lux et ténebr{\ae}, Dómino: * laudáte et superexaltáte eum in sǽcula.\par
Benedícite, fúlgura et nubes, Dómino: * laudáte et superexaltáte eum in sǽcula.\par

\lett{B}{ened\smash{í}cat} terra Dóminum: * laudet et superexáltet eum in sǽcula.\par
\secondline{Benedícite, montes et colles, Dómino: * laudáte et superexaltáte eum in sǽcula.}
\thirdline{Benedícite, univérsa germinántia in terra, Dómino: * laudáte et superexaltáte eum in sǽcula.}
Benedícite, fontes, Dómino: * laudáte et superexaltáte eum in sǽcula.\par
Benedícite, mária et flúmina, Dómino: * laudáte et superexaltáte eum in sǽcula.\par
Benedícite, cete, et ómnia, qu{\ae} movéntur in aquis, Dómino: * laudáte et superexaltáte eum in sǽcula.\par
Benedícite, omnes vólucres c{\ae}li, Dómino: * laudáte et superexaltáte eum in sǽcula.\par
Benedícite, omnes bésti{\ae} et pécora, Dómino: * laudáte et superexaltáte eum in sǽcula.\par
Benedícite, fílii hóminum, Dómino: * laudáte et superexaltáte eum in sǽcula.\par
\lett{B}{ened\smash{í}cat} Israël Dóminum: * laudet et superexáltet eum in sǽcula.\par
\secondline{Benedícite, sacerdótes Dómini, Dómino: * laudáte et superexaltáte eum in sǽcula.}
\thirdline{Benedícite, servi Dómini, Dómino: * laudáte et superexaltáte eum in sǽcula.}
Benedícite, spíritus, et ánim{\ae} justórum, Dómino: * laudáte et superexaltáte eum in sǽcula.\par
Benedícite, sancti, et húmiles corde, Dómino: * laudáte et superexaltáte eum in sǽcula.\par
\lett{B}{enedic\smash{á}mus} Patrem et Fílium cum Sancto Spíritu: * laudémus et superexaltémus eum in sǽcula.

\begin{rubric}
    {Then shall be read, in like manner, the Second Lesson, taken out of the New Testament, according to the Table or Calendar.}
\end{rubric}
%\begin{rubric}
%    And after that may be sung a Hymn, then shall be sung or said the Hymn following.
%\end{rubric}

%Permission to shorten Benedictus removed.
\subby{Benedictus}\label{Benedictus}

%MANUAL ADJUSTMENT:
\vspace{-2ex}

\begin{rubric}
%Permission for shorter Office removed.
    After which may be sung or said a Hymn and then shall be sung or said the \emph{Benedictus}, as followeth.
\end{rubric}

%MANUAL ADJUSTMENT:
\vspace{-1ex}

\lett{B}{ened\smash{í}ctus} {\ding{64}} Dóminus, Deus Isra{\"e}l: * quia visitávit, et fecit redemptiónem plebi su{\ae}:\par
\secondline{Et eréxit cornu salútis nobis: * in domo David, púeri sui.}
\thirdline{Sicut locútus est per os sanctórum, * qui a sǽculo sunt, prophetárum ejus:}
Salútem ex inimícis nostris, * et de manu ómnium, qui odérunt nos.\par
Ad faciéndam misericórdiam cum pátribus nostris: * et memorári testaménti sui sancti.\par
Jusjurándum, quod jurávit ad Ábraham patrem nostrum, * datúrum se nobis:\par
Ut sine timóre, de manu inimicórum nostrórum liberáti, * serviámus illi.\par
In sanctitáte, et justítia coram ipso, * ómnibus diébus nostris.\par
Et tu, puer, Prophéta Altíssimi vocáberis: * pr{\ae}íbis enim ante fáciem Dómini, paráre vias ejus:\par
Ad dandam sciéntiam salútis plebi ejus: * in remissiónem peccatórum eórum:\par
Per víscera misericórdi{\ae} Dei nostri: * in quibus visitávit nos, óriens ex alto:\par
Illumináre his, qui in ténebris, et in umbra mortis sedent: * ad dirigéndos pedes nostros in viam pacis.

%Permission for Psalm 100 removed.
\vspace{-2ex}
\subby{Apostles' Creed}
%While the Athanasian Creed was excised from the American tradition (during an attempt to remove every Creed), it is present in the Prayer Book Tradition, including the Canadian version. Therefore, it is thus supplied here. Since the Canadian tradition allows the Athanasian Creed on any day of the year, it is allowed likewise here.
\begin{rubric}
    Then shall be said the Apostles' Creed, by the Minister and the People, standing.\par%Permission for alternate words removed.
    \textsc{Note,} the Nicene Creed may be said instead of the Apostles' (p. \pageref{NiceneCreed}).
\end{rubric}
\begin{rubric}
    Upon these Feasts; \emph{Christmas Day}, the \emph{Epiphany}, Saint \emph{Matthias}, \emph{Easter Day}, \emph{Ascension Day}, \emph{Whitsunday}, Saint \emph{John Baptist}, Saint \emph{James}, Saint \emph{Bartholomew}, Saint \emph{Matthew}, Saint \emph{Simon} and Saint \emph{Jude}, Saint \emph{Andrew}, and upon \emph{Trinity Sunday}, shall be sung or said, instead of the Apostle's Creed, the Athanasian Creed (p. \pageref{Ath}).
\end{rubric}
\lett{C}{redo} in Deum, Patrem omnipoténtem, Creatórem c{\ae}li et terr{\ae}. Et in Jesum Christum, Fílium ejus únicum, Dóminum nostrum: qui concéptus est de Spíritu Sancto, natus ex María Vírgine, passus sub Póntio Piláto, crucifíxus, mórtuus, et sepúltus: descéndit ad ínferos; tértia die resurréxit a mórtuis; ascéndit ad c{\ae}los; sedet ad déxteram Dei Patris omnipoténtis: inde ventúrus est judicáre vivos et mórtuos.\par
     Credo in Spíritum Sanctum, sanctam Ecclésiam cathólicam, Sanctórum communiónem, remissiónem peccatórum, carnis {\ding{64}} resurrectiónem, vitam {\ae}térnam. Amen.

\begin{rubric}
    And after that, these Prayers following, the People devoutly kneeling; the Minister first pronouncing,
\end{rubric}
%\begin{rubric}
%    The Minister begins with the \emph{Dóminus vobíscum} only if he be in Major Orders. Otherwise, he shall use the second option.
%\end{rubric}
℣. Dóminus vobíscum.\par
\textit{vel,} Dómine, exáudi oratiónem nostram.

℟. Et cum spíritu tuo.\par
\textit{vel,} Et clamor noster ad te véniat.

℣. Orémus.

℣. Kýrie, eléison.

℟. Christe, eléison.

℣. Kýrie, eléison.

\subby{Lord's Prayer}
\begin{rubric}
    Then the Minister, Clerks, and People---all kneeling---shall say the Lord's Prayer with a loud voice.
\end{rubric}
\lett{P}{ater} noster, qui es in c{\ae}lis, sanctificétur nomen tuum: advéniat regnum tuum: fiat volúntas tua, sicut in c{\ae}lo et in terra. Panem nostrum quotidiánum da nobis hódie: et dimítte nobis débita nostra, sicut et nos dimíttimus debitóribus nostris: et ne nos indúcas in tentatiónem: sed líbera nos a malo. Amen.

\vspace{-0.5ex}
\subby{Preces}
\begin{rubric}
    Then the Minister standing up shall say,
\end{rubric}
℣. Osténde nobis, Dómine, misericórdiam tuam.\par
℟. Et salutáre tuum da nobis.\par
℣. Dómine salvam fac \textit{Civitatem}.\par
℟. Et exáudi nos cum invocámus te.\par
℣. Sacerdótes tui induántur Justítia.\par
℟. Et sancti tui exúltent.\par
℣. Salvum fac Pópulum tuum, Dómine.\par
℟. Et bénedic H{\ae}reditáti tu{\ae}.\par
℣. Da pacem Dómine in diébus nostris.\par
℟. Quóniam tu, Dómine, singuláriter in spe constituísti me.\par
℣. Cor mundum crea in nobis, O Deus.\par
℟. Et Spíritum Sanctum tuum ne áuferas a nobis.
%\elcol{℣. O Lord, show thy mercy upon us.}{℣. Osténde nobis, Dómine, misericórdiam tuam.}
%\elcol{℟. And grant us thy salvation.}{℟. Et salutáre tuum da nobis.}
%\elcol{℣. O Lord, save the \textit{State}.}{℣. Dómine salvam fac \textit{Civitatem}.}
%\elcol{℟. And mercifully hear us when we call upon thee.}{℟. Et exáudi nos cum invocámus te.}
%\elcol{℣. Endue thy Ministers with righteousness.}{℣. Sacerdótes tui induántur Justítia.}
%\elcol{℟. And make thy chosen people joyful.}{℟. Et sancti tui exúltent.}
%\elcol{℣. O Lord, save thy people.}{℣. Salvum fac Pópulum tuum, Dómine.}
%\elcol{℟. And bless thine inheritance.}{℟. Et bénedic H{\ae}reditáti tu{\ae}.}
%\elcol{℣. Give peace in our time, O Lord.}{℣. Da pacem Dómine in diébus nostris.}
%See the notes for this petition in Evensong.
%\elcol{℟. Because there is none other that fighteth for us, but only thou, O God.}{℟. Quia non est álius qui pugnet pro nobis, nisi tu Deus noster.}
%Testing using American version.
%\elcol{℟. For it is thou, Lord, only, that makest us dwell in safety.}{℟. Quóniam tu, Dómine, singuláriter in spe constituísti me.}
%\elcol{℣. O God, make clean our hearts within us.}{℣. Cor mundum crea in nobis, O Deus.}
%\elcol{℟. And take not thy Holy Spirit from us.}{℟. Et Spíritum Sanctum tuum ne áuferas a nobis.}
\begin{rubric}
%1928 Rubric:
Then shall follow the Collect(s) for the Day, except when the Communion Service is read; and then the Collect(s) for the Day shall be omitted here.
\end{rubric}
\vspace{-2ex}
\subby{A Collect for Peace}
\lett{D}{eus,} auctor pacis et amator, quem nosse vivere, cui servire regnare est: Protege ab omnibus impugnationibus supplices tuos; ut qui in tua defensione confidimus, nullius hostilitatis arma timeamus, per potentiam Jesu Christi Domini nostri. \textit{Amen.}
%\vspace{-1ex}
\subby{A Collect for Grace}
\lett{D}{\smash{ó}mine} Sancte, Pater Omnipotens, æterne Deus, qui nos ad principium hujus diei pervenire fecisti; Tua nos hodie salva virtute; et concede ut in hac die ad nullum declinemus peccatum, nec ullum incurramus periculum; sed semper ad tuam justitiam faciendam omnis nostra actio tuo moderamine dirigatur. Per Jesum Christum Dominum nostrum. \textit{Amen.}
%Conclusion restored from the Breviary tradition:
%Rubric here already specified in the rubrics page.
%\begin{rubric}
%    The Minister begins with the \emph{Dóminus vobíscum} only if he be in Major Orders. Otherwise, he shall use the second option.
%\end{rubric}
%\vspace{-1ex}
\subby{Conclusion}
%It was difficult to determine whether to use `my' or `our' for those not in Major Orders. Since both are present in the English liturgy, and the beginning uses `our' we decided to remain consistent and use `our'.
℣. Dóminus vobíscum.\par
\textit{vel,} Dómine, exáudi oratiónem nostram.

℟. Et cum spíritu tuo.\par
\textit{vel,} Et clamor noster ad te véniat.

℣. Benedicámus Dómino (allelúja, allelúja.)

℟. Deo grátias (allelúja, allelúja.)

℣. Fidélium ánim{\ae} {\ding{64}} per misericórdiam Dei requiéscant in pace.

℟. Amen.
 \begin{rubric}
     The following Prayers shall be omitted here when the Great Litany is said, and may be omitted when the Holy Communion is to follow.
 \end{rubric}
 %The American Tradition experiences a growth of options after the Third Collect. To accommodate this growth, longer than the English Prayer Book, the Minister is given great liberty to remove any of the following prayers. While this may not be the best method to prune an overgrown tree, we do not believe ourselves to be best-suited to choose an alternative.
\begin{rubric}
    And \textsc{Note}, That the Minister may here end the Morning Prayer with such intercessions taken out of this Book, as he shall think fit, or with the Grace.
\end{rubric}
\vspace{-2ex}
\subby{After the Third Collect}
\begin{rubric}
    A Prayer for one's nation (p. \pageref{prayers}) is here said.
\end{rubric}
\vspace{-2ex}
\subbysub{A Prayer for the Clergy and People}
\lett{O}{mnipotens} sempiterne Deus, qui facis mirabilia magna solus; Prætende super famulos tuos Pontifices nostros, et Parochos, et super cunctos Greges illis commissos, Spiritum gratiæ salutaris; et ut in veritate tibi complaceant, perpetuum eis rorem tuæ benedictionis infunde. Hoc quæsumus Domine, donare digneris, in honorem Advocati et Mediatoris nostri Jesu Christi. \textit{Amen.}
\begin{rubric}
%Move the rubric from Prayers to here:
    Additional Prayers (p. \pageref{prayers}) may be said here.
\end{rubric}
\vspace{-2ex}
\subbysub{Oratio pro omnibus hominibus cujuscunque conditionis}
\lett{D}{eus,} humani generis Creator et Conservator, supplices te rogamus pro omnibus hominibus cujuscunque conditionis, ut iis vias tuas, et omnibus gentibus salutare tuum, ostendere digneris. Et præcipue pro bono statu Ecclesiæ Catholicæ; quam bono Spiritu tuo, quæsumus, ita dirigas et gubernes, ut omnes qui se Christianos profitentur, in viam veritatis ducti, in unitate spiritus, in vinculo pacis, et in justitia vitæ fidem teneant. Postremo paternæ bonitati tuæ omnes illos commendamus qui quovis modo, in mente, corpore, aut rebus, anguntur aut laborant; [* eos præsertim pro quibus orationes nostræ poscuntur,] ut eos, prout variæ eorum requirunt necessitates, consolari et relevare digneris, et fortitudinem in patiendo, et felicem e pressuris exitum concedens. Quod propter Jesum Christum te rogamus. \textit{Amen.}
\subbysub{Gratiarum Actio generalis}
\begin{rubric}
    The General Thanksgiving may be said by the Congregation with the Minister.
\end{rubric}
\lett{O}{mnipotens} Deus, Pater omnium misericordiarum, nos indigni famuli tui humillime et ex animo tibi gratias agimus pro omni bonitate tua et benignitate erga nos et omnes homines: (eos præcipue qui jam volunt tibi confiteri et gratias agere pro misericordia tua nuper in se collata.) Benedicimus tibi pro creatione et conservatione nostra et omnibus hujusce vitæ bonis; sed maxime pro amore tuo inæstimabili, quo mundum per Dominum nostrum Jesum Christum redemisti; insuper pro gratiæ instrumentis et spe gloriæ. Quæ omnia beneficia tua da nobis ita, ut justum est, sentire, ut cordibus vere gratis, laudem tuam non loquendo tantum, sed vivendo annuntiemus; dum servitio tuo penitus devoti, coram te in sanctitate et justitia omnibus diebus vitæ nostræ ambulemus. Per Jesum Christum Dominum nostrum, cui tecum, in unitate Spiritus Sancti, sit omnis honor et gloria, per omnia sæcula sæculorum. \textit{Amen.}
\begin{rubric}
%Move the rubric from Thanksgivings to here:
    The Thanksgivings (p. \pageref{thanksgiving}) may here be offered.
\end{rubric}
\vspace{-2ex}
\subbysub{Oratio Sancti Chrysostomi.}
\lett{O}{mnipotens} Deus, qui nobis gratiam dedisti, ut hoc tempore concordes te communi supplicatione invocemus; et duobus vel tribus congregatis in Nomine tuo te quod rogaverint daturum polliceris: Vota, quæsumus, Domine, et petitiones famulorum tuorum, jam prout illis maxime expediat, prosequere; nobis in hoc sæculo veritatis tuæ cognitionem et in futuro vitam largiens æternam. \textit{Amen.}
\subbysub{2 ad Corinthios, cap. xiii.}
\lett{G}{ratia} Domini nostri Jesu Christi, et charitas Dei, et communicatio Sancti Spiritus, sit cum omnibus nobis in sæcula sæculorum. \textit{Amen.}

\begin{center}
    \textsc{Finis Ordinis Matutinarum per totum Annum.}
\end{center}
%\fancyhead[C]{{\LARGE Holy Mass}}
\fancyhead[RO,LE]{}
\fancyhead[RE,LO]{}
%\fancyhead[RO,LE]{\textit{Liturgy of St. Tikhon}}
%\section{Liturgy of St. Tikhon}\label{tikhon}
\phantomsection
\addcontentsline{toc}{section}{Liturgy of St. Tikhon}

\subby{Collect for Purity}
\begin{rubric}
    The Priest, at the Foot of the Altar, saith the Collect for Purity, as followeth.
\end{rubric}
\lett{O}{mn\smash{í}potens} Deus, cui omne cor patet, et cui omnes affectus animorum cogniti sunt, et quem nihil latet, purifica cogitationes cordium nostrorum, ut per inspirationem Sancti Spiritus te ex animo amemus, et debita veneratione celebramus Nomen tuum sanctum. Per Jesum Christum Dominum nostrum. \textit{Amen.}
\begin{rubric}
Then shall the Priest, turning to the People, rehearse distinctly The Ten Commandments; and the People, still kneeling, shall, after every Commandment, ask God mercy for their transgressions for the time past, and grace to keep the law for the time to come.
\end{rubric}
\begin{rubric}
And \textsc{Note,} that in rehearsing the Ten Commandments, the Priest may omit that part of the Commandment which is inset.
\end{rubric}
\begin{rubric}
	The Decalogue may be omitted. But \textsc{Note,}, That when ever it is omitted, the Minister shall say the Summary of the Law, beginning, \emph{Audite quod Dominus noster Jesus Christus dicit.}
\end{rubric}
%MANUAL ADJUSTMENT:
\vspace{-1ex}
\subby{Ten Commandments}
\begin{center}
	\textsc{Deus h{\ae}c verba ad hunc modum effatus est.:}
\end{center}
%\begin{multicols}{2}
\par\noindent
    Ego sum Dominus Deus tuus. Deos nullos alios habebis pr{\ae}ter me.\par
    \textit{Dómine miserére nostri, et dírige corda nostra ad servándam hanc legem.}
    \par\noindent
    Non facies tibi sculptile, neque ullam similitudinem ullius rei quæ est supra in cœlo, aut infra in terra, aut in aquis sub terra : non adorabis ea nec coles. 
    \par\noindent
    \leftskip=2em
	{\small{Ego enim Deus tuus fortis zelotes sum, visitans iniquitates patrum in filios, in tertiam et quartam generationem eorum qui oderunt me, et faciens misericordiam in millia, his qui diligunt et custodiunt præcepta mea.}}
	\par
	\leftskip=0cm
    \textit{Dómine miserére nostri, et dírige corda nostra ad servándam hanc legem.}
    \par\noindent
    Non assumes nomen Domini Dei tui in vanum;
    \par\noindent
    \leftskip=2em
	{\small{non enim habebit insontem Dominus eum, qui assumpserit nomen Domini Dei frustra.}}
	\par
	\leftskip=0em
    \textit{Dómine miserére nostri, et dírige corda nostra ad servándam hanc legem.}
    \par\noindent
    Memento ut diem Sabbati sanctifices.
    \par\noindent
    \leftskip=2em
	{\small{Sex diebus operaberis, et facies omnia opera tua, septimo autem die Sabbatum Domini Dei tui est : nullun in eo facies opus, tu et filius tuus et filia tua, servus tuus et ancilla tua, jumentum tuum, et advena qui est intra portas tuas. Sex enim diebus fecit Dominus cœlum et terram et mare, et omnia quæ in eis sunt, et requievit die septimo. Idcirco benedixit Dominus Dominus diei Sabbati, et sanctificavit eum.}}
	\par
	\leftskip=0em
    \textit{Dómine miserére nostri, et dírige corda nostra ad servándam hanc legem.}
    \par\noindent
    Honora patrem tuum et matrem tuam;
    \par\noindent
    \leftskip=2em
	{\small{ut sis longævus super terram, quam Dominus Deus tuus dabit tibi.}}
	\par
	\leftskip=0em
    \textit{Dómine miserére nostri, et dírige corda nostra ad servándam hanc legem.}
    \par\noindent
    Non occides.\par
    \textit{Dómine miserére nostri, et dírige corda nostra ad servándam hanc legem.}
    \par\noindent
    Non committes adulterium.\par
    \textit{Dómine miserére nostri, et dírige corda nostra ad servándam hanc legem.}
\par\noindent
    Non furtum facies.\par
    \textit{Dómine miserére nostri, et dírige corda nostra ad servándam hanc legem.}
\par\noindent
    Non loqueris contra proximum tuum falsum testimonium.\par
    \textit{Dómine miserére nostri, et dírige corda nostra ad servándam hanc legem.}
    \par\noindent
    Non concupisces.
    \par\noindent
    \leftskip=2em
	{\small{domum proximi tui, nec desiderabis uxorem ejus, non servum, non ancillam, non bovem, non asinus, nihil denique quod sit alterius.}}
	\par
	\leftskip=0em
	\textit{Dómine miserére nostri, et qu{\ae}sumus has omnes leges in córdibus nostris inscríbas.}
	
%	\end{multicols}
\begin{rubric}
	Then may the Priest say the Summary of the Law, as followeth.
\end{rubric}
\subby{Summary of the Law}
\begin{center}
%1979 Latin:
	{\textsc{Audite quod Dominus noster Jesus Christus dicit:}}
\end{center}
%Vulgate:
\lett{D}{iliges} Dominum Deum tuum ex toto corde tuo, et in tota anima tua, et in tota mente tua. Hoc est maximum, et primum mandatum. Secundum autem simile est huic: Diliges proximum tuum, sicut teipsum. In his duobus mandatis universa lex pendet, et prophetae. 
\begin{rubric}
    Then shall the Priest turn towards the Altar and ascend, saying secretly the \emph{Take away}.
\end{rubric}
\begin{rubric}
    Then, with hands joined upon the Altar, the Priest saith, bowing (kissing the Altar at the middle), the \emph{We pray thee}.
\end{rubric}
\begin{rubric}
    At a solemn Mass, the Celebrant, before reading the Introit, shall bless the incense with the \emph{Be thou blessed}.
\end{rubric}
\begin{rubric}
	Receiving the thurible from the Deacon, he censeth the Altar, saying nothing. Then the Deacon taketh the thurible from the Celebrant and censeth him only.
\end{rubric}
\begin{rubric}
%Added clarity that the Priest should be in the middle of the Altar for the Kyrie.
    Then shall the Priest move to the Epistle side and, signing himself with the sign of the Cross, begin the Introit: which finished, at the middle of the Altar, with joined hands, he saith alternately with the Ministers the \emph{Kyrie},
\end{rubric}
\subby{Kyrie}
\begin{multicols}{3}
	℣. Kyrie, eléison.

℟. Kyrie, eléison.

℣. Kyrie, eléison.

\columnbreak
℟. Christe, eléison.

℣. Christe, eléison.

℟. Christe, eléison.

\columnbreak
℣. Kyrie, eléison.

℟. Kyrie, eléison.

℣. Kyrie, eléison.
\end{multicols}


\subby{Gloria in Excelsis}
\begin{rubric}
    Then the Priest shall, at the middle of the Altar, extend and join his hands and---bowing his head slightly---say, if it is to be said,
\end{rubric}
\lett{G}{l\smash{ó}ria} in excélsis Deo. Et in terra pax homínibus bon{\ae} voluntátis. Laudámus te. Benedícimus te. Adorámus te.\margrub{Bow head.} Glorificámus te. Grátias ágimus tibi\margrub{Bow head.} propter magnam glóriam tuam. Dómine Deus, Rex c{\ae}léstis, Deus Pater omnípotens.\par
    Dómine Fili unigénite, Jesu Christe. Dómine Deus, Agnus Dei, Fílius Patris. Qui tollis peccáta mundi, miserére nobis. Qui tollis peccáta mundi, súscipe deprecatiónem nostram.\margrub{Bow head.} Qui sedes ad déxteram Patris, miserére nobis.\par
    Quóniam tu solus Sanctus. Tu solus Dóminus. Tu solus Altíssimus, Jesu Christe. Cum Sancto Spíritu {\ding{64}} in glória Dei Patris. Amen.

\begin{rubric}
    Then shall the Priest kiss the Altar in the middle, turn to the People, extend his hands, and say,
\end{rubric}
℣. Dóminus vobíscum.

℟. Et cum spíritu tuo.

℣. Orémus.
\subby{Collect}
\begin{rubric}
	Then shall the Priest move to the Epistle corner of the Altar and pray the Collect(s) of the Day.
\end{rubric}
\subby{Epistle}
\begin{rubric}
%CHECK:
	The Epistle for the Day shall then be read by the Subdeacon, first saying, \emph{Epistola scripta est in Capitulo \emph{--} et incipit ad Versum \emph{--}.}
\end{rubric}
\begin{rubric}
	The Epistle ended, he shall say, \emph{Hic explicit Epistola}, the people responding, \emph{Deo gratias}.
\end{rubric}
\begin{rubric}
    The Gradual and Alleluia (or Tract) and (if provided) Sequence is here chanted by the Choir.
\end{rubric}
\subby{Gospel}
\begin{multicols}{2}
\begin{rubric}
    These being ended, if it be a Solemn Mass, the Deacon shall then place the book of the Gospels on the middle of the Altar, and the Celebrant shall bless the incense with the \emph{Be thou blessed}, as above.
\end{rubric}
\begin{rubric}
	Then shall the Deacon, kneeling before the Altar, say with joined hands the \emph{Cleanse my heart} continuing to the \emph{Bid, sir}, with the priest saying \emph{The Lord be in thy}.
\end{rubric}
\begin{rubric}
    Having received the blessing, the Deacon kisseth the hand of the Celebrant. And going with the other Ministers, with the incense and the lights, to the place of the Gospel, he standeth with joined hands, saying what is below.
\end{rubric}
\begin{rubric}
    If, however, the Priest celebrate without Deacon and Subdeacon, when the book has been carried to the other corner of the Altar, he boweth in the middle, and with joined hands, saith the \emph{Cleanse my heart} and the \emph{Bid, Lord} and \emph{The Lord be in my}.
\end{rubric}
\end{multicols}
\begin{rubric}
	Then, all the People standing, the Priest, turning to the Book, shall say with joined hands,
\end{rubric}

℣. Dóminus vobíscum.

℟. Et cum spíritu tuo.

℣. Inítium (vel, Sequéntia) {\ding{66}} sancti Evangélii secúndum \textit{N.}

℟. Glória tibi, Dómine.
\begin{rubric}
    Then shall the Priest sign the book with the thumb of his right hand at the beginning of the Gospel text which he is to read, then himself on the forehead, the mouth, and the breast: and while the Ministers respond, he censeth the book thrice, then readeth the Gospel with joined hands.
\end{rubric}
\begin{rubric}
    At the end of the Gospel, the Ministers respond,
\end{rubric}

℟. Laus tibi, Christe.
\begin{rubric}
    Then shall the Subdeacon carry the book to the Priest, who kisseth the Gospel text, saying: \emph{Through the words of the Gospel may our sins be blotted out.} Then the Priest is censed by the Deacon.\par
    \textsc{Note,} In Masses of the Dead, \emph{Cleanse} is said, but a blessing is not asked, lights are not carried, and the Celebrant doth not kiss the book.
\end{rubric}
\begin{rubric}
    Then, in the middle of the Altar, extending, raising, and joining his hands, the Priest shall say, if it is to be said, \emph{I believe in one God}, proceeding with joined hands.
\end{rubric}
\subby{Nicene Creed}\label{NiceneCreed}
\lett{C}{redo} in unum Deum,\margrub{Bow head to Cross.} Patrem omnipoténtem, factórem c{\ae}li et terr{\ae}, visibílium ómnium et invisibílium.\par
    Et in unum Dóminum Jesum Christum,\margrub{Bow head to Cross.} Fílium Dei unigénitum. Et ex Patre natum ante ómnia sǽcula. Deum de Deo, lumen de lúmine, Deum verum de Deo vero. Génitum, non factum, consubstantiálem Patri: per quem ómnia facta sunt. Qui propter nos hómines et propter nostram salútem descéndit de c{\ae}lis. \inrub{Everyone genuflects.} Et incarnátus est de Spíritu Sancto ex María Vírgine: Et homo factus est. \inrub{Everyone rises.} Crucifíxus étiam pro nobis: sub Póntio Piláto passus, et sepúltus est. Et resurréxit tértia die, secúndum Scriptúras. Et ascéndit in c{\ae}lum: sedet ad déxteram Patris. Et íterum ventúrus est cum glória judicáre vivos et mórtuos: cujus regni non erit finis.\par
    Et in Spíritum Sanctum, Dóminum et vivificántem: qui ex Patre procédit. Qui cum Patre et Fílio simul adorátur\margrub{Bow head to Cross.} et conglorificátur: qui locútus est per Prophétas. Et unam sanctam cathólicam et apostólicam Ecclésiam. Confíteor unum baptísma in remissiónem peccatórum. Et exspécto resurrectiónem mortuórum. {\ding{64}} Et vitam ventúri s{\ae}culi. Amen.
\begin{rubric}
Then shall be declared unto the People what Holy-days, or Fasting-days, are in the week following to be observed; and (if occasion be) shall Notice be given of the Communion, and of the Banns of Matrimony, and other matters to be published.
\end{rubric}

\begin{rubric}
Here, or immediately after the Creed, may be said the Bidding Prayer, or other authorized Prayers and intercessions.
\end{rubric}

\begin{rubric}
	Then followeth the Sermon.
\end{rubric}

\begin{rubric}
    Then shall the Priest ascend the Altar, kiss it, turn towards the People, and say,
\end{rubric}

℣. Dóminus vobíscum.

℟. Et cum spíritu tuo.

℣. Orémus.
\subsec{Offertory}
\begin{multicols}{2}    
\begin{rubric}
    The Priest then saith the Offertory Verse.\par
    \textsc{Note,} At High Mass, a hymn may be sung while the Priest prepareth the Oblations in the Offertory, and while the Collection is taken if there be one.
\end{rubric}
\begin{rubric}
    Then shall the Deacon present the Paten with the Bread to the Celebrant, which he then offereth, saying: \emph{Receive, O holy Father}.
\end{rubric}

\begin{rubric} 
    Then, making a cross with the same Paten, he shall place the Bread upon the Corporal (or the Paten with the Bread upon the Corporal).
\end{rubric}
\begin{rubric}
	Then the Deacon shall minister the Wine, and the Subdeacon the water in the Cup. The Priest (blessing with the sign of the Cross the water to be mixed in the Cup, unless it be a Requiem Mass) shall say: \emph{O God, who didst wonderfully}.
\end{rubric}
\begin{rubric}
    Then shall the Priest receive the Cup and offer it, saying, \emph{We offer unto thee}.
\end{rubric}
\begin{rubric}
    Then shall he make the sign of the Cross with the Cup, place it upon the Corporal, and cover it with the Pall. Then, with hands joined upon the Altar, he saith, bowing slightly, \emph{In a humble spirit}.
\end{rubric}
\begin{rubric}
    Then shall the Priest---standing upright---extend his hands, raise them, and join them; and lifting his eyes to heaven and lowering them immediately, he saith, \emph{Come, O Sanctifier}.
\end{rubric}
\begin{rubric}
    If he be celebrating solemnly, the Priest blesseth incense, saying, \emph{Through the intercession}.
\end{rubric}
\begin{rubric}
    Receiving the Thurible from the Deacon, he censeth the Oblations, in the manner prescribed in the General Rubrics, saying: \emph{May this incense}.\par
    Then he censeth the Altar, saying: \emph{Let my prayer}.\par
    While he returneth the Thurible to the Deacon, he saith, \emph{The Lord kindle}.\par
    Then the Priest is censed by the Deacon, and afterwards the others in order.
\end{rubric}
\begin{rubric}
    Meanwhile, the Priest washeth his hands saying, \emph{I will wash}.\par
    \textsc{Note,} In Requiem Masses, and during Passiontide in Masses of the Season, \emph{Glory be} is omitted.
\end{rubric}
\begin{rubric}
    Then bowing slightly, in the middle of the Altar with hands joined upon it, he saith, \emph{Receive, O holy Trinity}.
\end{rubric}
\end{multicols}
\begin{rubric}
    Then shall the Priest kiss the Altar and, turning to the People, extend and join his hands, and say, raising his voice a little,
\end{rubric}

℣. Oráte, fratres: ut meum ac vestrum sacrifícium acceptábile fiat apud Deum Patrem omnipoténtem.

℟. Suscípiat Dóminus sacrifícium de mánibus tuis ad laudem et glóriam nominis sui, ad utilitátem quoque nostram, totiúsque Ecclési{\ae} su{\ae} sanct{\ae}.
\begin{rubric}
    The Priest saith in a low voice: \emph{Amen.}
\end{rubric}
\begin{rubric}
	Then shall the Priest, with hands extended, immediately (without \emph{Let us pray}), add the Secret Prayers.\par
	\textsc{Note,} When these are ended, he saith in a loud voice,
\end{rubric}
℣. Throughout all ages, world without end.\par
℟. Amen.
\begin{rubric}
    Then shall the Deacon say,
\end{rubric}
\centerline{Oremus pro statu universalis Ecclesi{\ae} Christi.}
\begin{rubric}
    The Priest then saith, with hands extended, the following.
\end{rubric}
\lett{O}{mn\smash{í}potens} sempiterne Deus, qui per sanctum Apostolum tuum nos docuisti facere orationes, obsecrationes, et gratiarum actiones pro omnibus hominibus: Supplices te rogamus ut clementer (eleemosynas atque) oblationes nostras accipias, et has preces nostras exaudias, quas offerimus Divinæ Majestati tuæ: Supplicantes ut veritatis, unitatis, et concordiæ spiritum Catholicæ Ecclesiæ tuæ perpetuo inspires: Et præsta ut omnes qui sanctum Nomen tuum confitentur, in sancti verbi tui veritate consentiant, et in unitate et pia charitate concordes vivant.

\needspace{4\baselineskip}
\lett{I}{nsuper} te rogamus ut corda omnium Christianorum Dominatorum domineris et disponas, ut recte ac sine personarum acceptione jus dicant, quo scelera et nequitia corrigantur, et vera tua religio, virtusque, stabiliantur.
\needspace{4\baselineskip}
\lett{D}{a} gratiam, Pater c{\oe}lestis, omnibus Episcopis et Parochis, præcipue famulum tuum \textit{N.} Metropolitano nostro (et \textit{N.} Episcopo nostro), ut tam vita quam doctrina sua verum vivumque verbum tuum annuntient, et sancta tua Sacramenta recte et rite ministrent.
\needspace{4\baselineskip}
\lett{E}{t} universo populo tuo tribue gratiam tuam; et præcipue huic congregationi præsenti; ut humili animo et debita reverentia audiant et accipiant sanctum verbum tuum: et tibi fideliter serviant in sanctitate et justitia omnibus diebus vitæ suæ.
\needspace{4\baselineskip}
\lett{S}{upplices} etiam te rogamus, Domine, ut pro bonitate tua eos omnes consoleris et adjuves, qui in hac temporali vita tribulatione, mœstitia, inopia, morbo, aliisve rebus adversis laborant.
\needspace{4\baselineskip}
\lett{B}{enedicimus} quoque sancto Nomini tuo propter omnes famulos tuos in fide et timore tuo defunctos; %1979 Latin:
te rogantes ut dones eisdem incrementum perpetuum in caritate ministerioque tuo; et gratiam nobis concedas qua, bona eorum exempla secuti, præcipue Beátæ Maríæ semper Vírginis, Genetrícis Jesu Christi, Fílii tui, Dómini et Dei nostri, etiam Sanctorum Patriárchum, Prophetarum, Apostolorum, Martyrumque, nos una cum illis cœlestis regni tui fiamus participes.
\begin{rubric}
	 The Priest placeth both hands upon the Altar, outside the Corporal, and continueth,
\end{rubric}
\needspace{4\baselineskip}
\lett{H}{oc,} Pater, largiri digneris, propter Jesum Christum unicum nostrum Mediatorem atque Advocatum.
\secondline{℟. Amen.}
%This is a difficult situation to discern. The Prayer Book Traditions reflects the uncertainty as to where to place the Confession. The current American placement is reflective of Cranmer's Protestant theology and liturgical sense, since the origins come from placing the Confession close to the reception of Communion (which was immediately after the Words of Institution). The 1928 reflects a more catholic understanding of the liturgy, including the epiklesis and the other integral parts of the Eucharistic Sacrifice before the Reception of Holy Communion. This understanding was more present in the 1549 which placed the Confession in a congruous place as the Roman tradition: after the Sacrifice and before the Communion of the Faithful. While this is obscured here due to the Byzantine devotion imposed, this is only a temporary imposition and the overall shape of the liturgy (due to the Orthodox theology present here) is better preserved in its Catholic and even Prayer Book form by placing the Confession immediately before the `Ecce Agnus Dei', which is where the Second Confiteor would go in the Roman Rite.


\subsec{Sursum Corda}
\begin{rubric}
    Then shall the Priest begin the Preface, with both hands placed apart on the Altar.
\end{rubric}
℣. Dóminus vobíscum.

℟. Et cum spíritu tuo.

℣. Sursum corda.\margrub{He raises his hands a little.}

℟. Habémus ad Dóminum.

℣. Grátias agámus Dómino, Deo nostro.\margrub{He joins his hands before his breast, \& bows his head.}

℟. Dignum et justum est.\margrub{He separates his hands.}
\subby{Preface}
\lett{I}{t} is very meet, right, and our bounden duty, that we should at all times, and in all places, give thanks unto thee, O Lord, Holy Father, Almighty, Everlasting God.
\begin{rubric}
    The proper preface is then here said (p. \pageref{prefaces})---if there be one---concluding with the following.
\end{rubric}
\label{PrefaceEnding}
\lett{T}{herefore} with Angels and Archangels, and with all the company of heaven, we laud and magnify thy glorious Name; evermore praising thee, and saying:\margrub{He joins his hands \& bows.}%\textsuperscript{\alph{margcount}}

%\hspace{0.98\textwidth}
%\rlap{\parbox{\marginparwidth}{\raggedright \itshape\scriptsize\textsuperscript{\alph{margcount}}%
%He joins his hands \& bows.}} % Use \parbox for multiline text

%\stepcounter{margcount}
%\stepcounter{latinrubric}

\subby{Sanctus}
\lett{S}{anctus, Sanctus, Sanctus,} Dóminus, Deus Sábaoth. Pleni sunt c{\ae}li et terra glória tua. Hosánna in excélsis. Benedíctus,\margrub{He stands upright \& signs himself.} {\ding{64}} qui venit in nómine Dómini. Hosánna in excélsis.
\clearpage
\checkoddpage
\ifoddpage \thispagestyle{empty}
~\clearpage\fi
\thispagestyle{empty}
\includegraphics[width=98mm, height=135mm]{Canon.eps}
\setcounter{margcount}{1}
\clearpage
%\reversemarginpar
\subby{Canon of the Mass}
\begin{rubric}
    The Priest, extending, raising somewhat, and joining his hands, raising his eyes towards heaven, and immediately lowering them, shall begin the Canon, as followeth.
\end{rubric}
%1979 Latin:
Omnis gloria tibi Deus omnipotens, Pater noster cœlestis, qui pro misericordiæ tuæ pietate unicum Filium tuum Jesum Christum dedisti, ut mortem in Cruce pro nostra redemptione pateretur; qui ibi (propria sui ipsius oblatione semel facta) plenum, perfectum, et sufficiens sacrificium, oblationem, et satisfactionem pro totius mundi peccatis fecit; et instituit, et in sancto Evangelio suo nobis præcepit observare, pretiosæ mortis sacrificique illius memoriam, usque dum rediret, perpetuam:
\subbysub{Words of Institution}\noindent
In enim qua nocte tradebatur,\margrub{He takes \& holds the Bread.} accepit Panem;\margrub{He looks up to heaven \& bows his head.} Et tibi gratias {\ding{64}} agens, fregit, deditque discipulis suis, dicens, Accipite, et manducate.\margrub{He holds the Bread between both of his thumbs \& forefingers.}
\begin{center}
\large{Hoc est enim Corpus meum, quod pro vobis datur: Hoc facite in meam commemorationem.}
\end{center}
\begin{rubric}
    Then shall the Priest immediately genuflect, briefly elevate the Bread for the People to see, replace it upon the corporal (or Paten), genuflect again, and then immediately continue,\par
    \textsc{Note,} From henceforth, the Priest doth not disjoin his forefingers \& thumbs until the ablutions.
\end{rubric}\par\noindent
Simili modo posteaquam cœnatum est,\margrub{The Priest uncovers the Cup, holds it with both hands, \& bows his head.} accipiens Calicem;\margrub{He holds the Cup in his left hand.} item tibi gratias {\ding{64}} agens, dedit illis, dicens, Bibite ex eo omnes;\margrub{He holds the Cup slightly raised.}
\begin{center}
\large{Hic est enim Sanguis meus Novi Testamenti, qui pro vobis et pro multis effunditur in remissionem peccatorum:}
\end{center}
\begin{rubric}
    Then shall the Priest place the Cup back upon the corporal and immediately say,
\end{rubric}\par\noindent
    \centerline{Hoc facite, quotiescumque bibetis, in meam commemorationem.}
\begin{rubric}
    Then shall the Priest immediately genuflect, briefly elevate the Cup, replace the Cup upon the corporal, cover it, genuflect again, and continue with hands extended.
\end{rubric}
\subbysub{Recollection}
%1979 Latin:
\lett{U}{nde,} Domine caelestis Pater, secundum institutionem dilectissimi Filii tui Salvatoris nostri Iesu Christi, nos humiles servi tui celebramus et hic facimus coram divina Maiestate tua, de his tuis sanctis {\ding{64}} donis, quae nunc tibi offerimus, illam quam Filius tuus celebrare nos praecepit memoriam; memores et eius tam beatae passionis et pretiosae mortis, necnon et mirabilis resurrectionis, sed et gloriosae ascensionis; tibi gratias agentes ex animo propter innumerabilia beneficia nobis inde collata.
\subbysub{Invocation}
\begin{rubric}
	Then shall the Priest uncover the Cup, bow profoundly, join his hands upon the Altar, and continue,
\end{rubric}
%Epiklesis of St. James through the 1929 Scottish BCP:
%Crosses supplied from the Anglican Missal.
%Moving the Pall from precedent in the LAP Missal:
\lett{E}{xaudi} nos, indignos servos tuos, quaesumus, misericors Pater,\margrub{He stands upright and imposes his hands over the Gifts.} et mitte tuum Sanctum Spiritum super nobis atque super his don {\ding{64}} is et creaturis tuis, pan {\ding{64}} is et vi {\ding{64}} ni, ut, benedictae et sanctificatae virtute vitali ejus, fiant Cor {\ding{64}} pus et Sang {\ding{64}} uis dilecti Filii tui, ut omnis qui percipit eundem sanctificetur in corpore animaque, et custodiatur in vitam aeternam. 

\begin{rubric}
	Then shall the Priest cover the Cup, genuflect, and continue, with hands extended,
\end{rubric}

\par
\needspace{4\baselineskip}
\lett{E}{t} rogamus supplices paternam tuam bonitatem, ut hoc nostrum laudis et gratiarum sacrificium benignus accipias: humillime supplicantes, ut propter merita et mortem Filii tui Jesu Christi, et per fidem in sanguine ipsius, et nos et Universa Ecclesia tua peccatorum remissionem et cætera omnia passionis ejus beneficia consequamur.
\subbysub{Commemoration for the Departed}
\lett{M}{em\smash{é}nto} étiam, Dómine, famulórum famularúmque tuárum \emph{N.} et \emph{N.}, qui nos præcessérunt cum signo fídei, et dórmiunt in somno pacis.
\begin{rubric}
    Then shall the Priest join his hands, pray for the dead he hath in mind, then extend his hands and say,
\end{rubric}\par\noindent
Ipsis, Dómine, et ómnibus in Christo quiescéntibus locum refrigérii, lucis, et pacis, ut indúlgeas, deprecámur.
%The English & American Missals lacks this addition, being only in the Anglican Missal tradition:
%\lett{A}{nd} vouchsafe to grant some part and fellowship with thy holy Apostles and Martyrs: with John, Stephen, Matthias, Barnabas, Ignatius, Alexander, Marcellinus, Peter, Felicitas, Perpetua, Agatha, Lucy, Agnes, Cecilia, Anastasia, and with all thy Saints: within whose fellowship we beseech thee admit us.
%\lett{R}{emember} Lord, also the souls of thy servants and handmaidens, which are gone before us with the mark of faith, and rest in the sleep of peace. \inrub{The Priest now prays for the souls of the dead.} We beseech thee, O Lord, that unto them, and unto all such as rest in Christ, thou wilt grant a place of refreshing, of light, and of peace. And vouchsafe to give unto us some portion and fellowship with thy holy Apostles and Martyrs; with John, Stephen, Matthias, Barnabas, Ignatius, Alexander, Marcellinus, Peter, Felicitas, Perpetua, Agatha, Lucia, Agnes, Cecilia, Anastasia, and with all thy Saints; within whose fellowship we beseech thee to admit us.
\subbysub{Oblation}
\lett{E}{t} hic tibi, Domine, offerimus et exhibemus nosmetipsos, animas et corpora nostra, tibi hostiam rationabilem, sanctam, et viventem; supplices te rogantes, ut quotquot hujus sacræ Communionis participes facti sumus, digne percipiant pretiosum Cor {\ding{64}} pus et Sang {\ding{64}} uinem\margrub{After signing the Gifts, he signs himself.} Filii tui Jesu Christi, omni benedictione cœlesti et gratia tua repleamur, et sint unum corpus cum eo, ut ipse in illis et illi in ipso maneant.

%\subbysub{Commemorations}
%Liturgy of St. Tikhon:
%\lett{B}{e} mindful also, O Lord, of thy servants who are gone before us with the sign of faith, and who rest in the sleep of peace, \textit{N.} and \textit{N.} To them O Lord, and to all who rest in Christ grant we pray thee a place of refreshment, light, and peace. To us sinners also, thy servants, confiding in the multitude of thy mercies, grant some lot and partnership with thy holy Apostles and Martyrs: John, Stephen, Matthias, Barnabas, Ignatius, Alexander, Marcellinus, Peter, Felicitas, Perpetua, Agatha, Lucia, Agnes, Cecilia, Anastasia, and with all thy Saints, into whose company we pray thee of thy mercy to admit us.
\par
\needspace{4\baselineskip}
\lett{E}{t} quamvis propter multiplicia peccata nostra non digni simus qui ullum sacrificium tibi offeramus, hanc tamen debitam oblationem servitutis nostræ, non æstimator meriti sed veniæ, quæsumus, largitor accipias;\margrub{He joins his hands \& bows his head profoundly.}
\par
\needspace{4\baselineskip}
\lett{P}{er} Jesum Christum Dominum nostrum,\margrub{He genuflects} per {\ding{64}} quem et cum {\ding{64}} quo est tibi Deo Patri Omnipotenti, in unitate Spiritus {\ding{64}} Sancti, omnis hon {\ding{64}} or et glo {\ding{64}} ria\margrub{He uncovers the Cup, elevates the Host \& Cup to the height of his breast then replacing them on the Altar, covers the Cup, genuflects.} per omnia sæcula sæculorum.\margrub{He joins his hands.}\par
℟. Amen.


\subby{Lord's Prayer}\par\noindent
Oremus. Præcéptis salutáribus móniti, et divína institutióne formáti audémus dícere:\margrub{He extends his hands.}

\lett{P}{ater} noster, qui es in c{\oe}lis, Sanctificetur Nomen tuum. Adveniat regnum tuum. Fiat voluntas tua, Sicut in c{\oe}lo, et in terra. Panem nostrum quotidianum da nobis hodie. Et dimitte nobis debita nostra, Sicut et nos dimittimus debitoribus nostris. Et ne nos inducas in tentationem: Sed libera nos a malo: Quia tuum est regnum, Potentia, et gloria, In sæcula sæculorum. Amen.
\begin{rubric}
    Unless the Priest consecrated on the Paten, he takes the Paten between the fore and middle fingers of his right hand, and holds it upright upon the Altar.
\end{rubric}
\needspace{4\baselineskip}
\lett{L}{\smash{í}bera} nos, quǽsumus, Dómine, ab ómnibus malis, prætéritis, præséntibus et futúris: et intercedénte beáta et gloriósa semper Vírgine Dei Genetríce María, cum beátis Apóstolis tuis Petro et Paulo, atque Andréa, et ómnibus Sanctis,\margrub{He signs himself (with the Paten, if he consecrated on the Corporal).} da propítius pacem in diébus nostris: ut, ope misericórdiæ tuæ adjúti, et a peccáto simus semper líberi et ab omni perturbatióne secúri.
\begin{rubric}
    If the Host be not already on the Paten, the Priest is to slide the Paten underneath the Host. He then shall uncover the Cup, genuflect, and break the Host in half over the Cup, saying,
\end{rubric}\par\noindent
Per eúndem Dóminum nostrum Jesum Christum, Fílium tuum.
\begin{rubric}
    The Priest placeth the half in his right hand upon the Paten. He then breaketh a particle off from the half in his left hand, saying,
\end{rubric}\par\noindent
Qui tecum vivit et regnat in unitáte Spíritus Sancti Deus.
\begin{rubric}
    Then shall the Priest join the other half, which he holdeth in his left hand, to the half laid upon the Paten, and retaining the small particle in his right hand over the Cup, which he holdeth with his left by the knop below the cup, say in an audible voice, or sing,
\end{rubric}\par\noindent
Per ómnia sǽcula sæculórum.\par
℟. Amen.\par
℣. Pax\margrub{With the same particle, he signs thrice over the Cup.} {\ding{64}} Dómini sit {\ding{64}} semper vobís {\ding{64}} cum.\par
℟. Et cum spíritu tuo.
\begin{rubric}
    The Priest, putting the same particle into the Cup, saith, in the secret voice:
\end{rubric}
\lett{H}{{\ae}c} commíxtio, et consecrátio Córporis et Sánguinis Dómini nostri Jesu Christi, fiat accipiéntibus nobis in vitam ætérnam. Amen.
\subby{Agnus Dei}
\begin{rubric}
    Then shall the Priest cover the Cup, genuflect, and---bowing to the Sacrament---join his hands and strike his breast thrice, saying in an audible voice:
\end{rubric}
\lett{A}{gnus} Dei, qui tollis peccáta mundi: miserére nobis.

\secondline{Agnus Dei, qui tollis peccáta mundi: miserére nobis.}
\thirdline{Agnus Dei, qui tollis peccáta mundi: dona nobis pacem.}
\begin{rubric}
	In Requiem Masses, \emph{miserére nobis} is replaced by \emph{dona eis réquiem} and \emph{dona nobis pacem} by \emph{dona eis réquiem sempitérnam}.
\end{rubric}
%This is the original location of the Prayer Book Confession, mirroring the placement of the Roman Confession. However, in order to emphasise Reformed theology, the later English Prayer Books placed reception of Holy Communion to immediately after the Words of Institution, requiring the Confession to be placed before the Canon. Even though the American tradition developed a more Catholic Eucharistic rite, the Reformed placement of the Confession remained. Here, this is corrected in order to better express our Orthodox Catholic theology.
\subsec{Confession}
\begin{rubric}
    Then the Deacon turneth to the People and saith to those who come to receive the Holy Communion,
\end{rubric}
\lett{V}{os} quos vere et serio peccatorum vestrorum p{\oe}nitet, qui erga proximos veram habetis charitatem, qui vitam novam instituere decrevistis, mandatis Dei obsequendo, et in viis ejus posthac ambulando: Cum fide accedite, ut hoc sanctum percipiatis Sacramentum ad vestram consolationem; et, reverenter genuflexi, humilem vestram Deo Omnipotenti confessionem facite.
\begin{rubric}
Then shall this General Confession be made, by the Priest, bowing, and all those who are minded to receive the Holy Communion, humbly kneeling.
\end{rubric}
\needspace{4\baselineskip}
\lett{O}{mnipotens} Deus, Pater Domini nostri Jesu Christi, Conditor omnium rerum, Omnium hominum judex: Confitemur et deploramus multiplicia peccata et delicta nostra, Quæ subinde impie admisimus, Cogitatione, verbo, et opere, Contra Divinam Majestatem tuam, Provocantes adversus nos justissimam iram et indignationem tuam. Serio nos pœnitet, Et ex animo dolemus ob has prævaricationes nostras: Quanim recordatio nobis acerba est, Onus intolerabile. Miserere nostri, Miserere nostri, Pater misericors; Propter Filium tuum Jesum Christum Dominum nostrum, Quod præteritum est nobis condona: Et concede ut semper posthac Tibi in novitate vitæ serviamus et placeamus, Ad honorem et gloriam Nominis tui; per Jesum Christum Dominum nostrum. Amen. 
\begin{rubric}
	Then shall the Priest (Bishop if he be present) stand up and, turning to the People, say,
\end{rubric}
\needspace{4\baselineskip}
\lett{O}{mnipotens} Deus, Pater noster cœlestis, qui pro magna misericordia sua omnibus ex animo pœnitentibus, ad se cum vera fide conversis, peccatorum remissionem est pollicitus: Misereatur vestri, et dimittat vobis {\ding{64}} omnia peccata vestra: liberet vos ab omni malo, conservet et confirmet in omni bono, et ad vitam perducat æternam; per Jesum Christum Dominum nostrum.

℟. Amen.
\subsec{Comfortable Words}
\begin{rubric}
Then shall the Priest, facing the People, say,
\end{rubric}\noindent
\begin{center}
	\textsc{Audite quam consola toriis verbis omnes ad se veraciter conversos alloquitur Christus Salvator noster.}
\end{center}
\par\noindent
Venite ad me, omnes qui laboratis et onerati estis, et ego reficiam vos. \vr{Matt. 11:28}
\par\noindent
    Sic Deus dilexit mundum, ut Filium suum unigenitum daret, ut omnis qui credit in eum non pereat, sed habeat vitam æternam. \vr{John 3:16}
    \par\noindent
    \begin{center}
		\textsc{Audite etiam quid dicat Sanctus Paulus.}
	\end{center}
\par\noindent
    Fidelis sermo, et omni acceptione dignus, quod Christus Jesus venit in hunc mundum peccatores salvos facere. \inrub{1 Tim. 1:15}
\par\noindent
    \begin{center}
		\textsc{Audite etiam quid dicat Sanctus Joannes.}
	\end{center}
    \par\noindent
    Si quis peccaverit, Advocatum habemus apud Patrem, Jesum Christum justum, et ipse est propitiatio pro peccatis nostris. \inrub{1 John 2:1-2}
\begin{rubric}
    Then the Priest, turning to the Altar, bowing with hands joined upon the Altar, saith in the secret voice,
\end{rubric}
\lett{D}{\smash{ó}mine} Jesu Christe, qui dixísti Apóstolis tuis: Pacem relínquo vobis, pacem meam do vobis: ne respícias peccáta mea, sed fidem Ecclésiæ tuæ; eámque secúndum voluntátem tuam pacificáre et coadunáre dignéris: Qui vivis et regnas Deus per ómnia sǽcula sæculórum. Amen.
\begin{rubric}
    If the \emph{Pax} be given, the Priest is to kiss the Altar and---giving the \emph{Pax}---say: \emph{Pax tecum.} with the response \emph{Et cum spiritu tuo.}
\end{rubric}
\subby{Prayer of Humble Access}
\begin{rubric}
    The Priest alone then saith,
\end{rubric}
\lett{N}{on} justitiæ nostræ, misericors Domine, sed multitudinis magnarum miserationum tuarum fiducia, ad hanc Mensam tuam accedere audemus. Non sumus digni qui vel micas sub Mensa tua colligamus. Tu autem idem ille es Dominus, cui proprium est semper misereri : Tribuas igitur nobis, benigne Domine, Carnem dilecti Filii tui Jesu Christi ita manducare, et Sanguinem ejus bibere, ut corpora nostra immunda per Corpus ejus mundentur, et animæ per pretiosissimum ejus Sanguinem laventur, et nos perpetuo habitemus in eo, et ipse in nobis.\par
℟. Amen.
\subby{Communion of the Priest}
\begin{multicols}{2}
\begin{rubric}
    Then the Priest shall say the \emph{Dómine Jesu Christe} and \emph{Percéptio Córporis tui}.
\end{rubric}
\begin{rubric}
    He then genuflecteth and saith: \emph{Panem cæléstem accípiam, et nomen Dómini invocábo.}.
\end{rubric}
\begin{rubric}
    Then, bowing slightly, the Priest shall take both parts of the Host between the thumb and forefinger of his left hand, and place the Paten between the same forefinger and middle finger, and striking his breast three times with his right hand, say thrice, devoutly and humbly, raising his voice slightly:
\end{rubric}\par\noindent
    Dómine, non sum dignus, \inrub{Proceeding in the secret voice:} ut intres sub tectum meum: sed tantum dic verbo, et sanábitur ánima mea.
\begin{rubric}
After signing himself with his right hand with the Host over the Paten, he saith,
\end{rubric}\par\noindent
Corpus Dómini nostri Jesu Christi custódiat ánimam meam in vitam ætérnam. Amen.
\begin{rubric}
And bowing, he reverently taketh both parts of the Host. After consuming the Host, he is to place the Paten down upon the Corporal, and raising himself, join his hands, and remain still for a short time in meditation on the Most Holy Sacrament.
\end{rubric}
\begin{rubric}
	Then he shall uncover the Cup, genuflect, collect the fragments (if there be any), and cleanse the Paten over the Cup, saying meanwhile:
\end{rubric}
\par\noindent
Quid retríbuam Dómino pro ómnibus, quæ retríbuit mihi? Cálicem salutáris accípiam, et nomen Dómini invocábo. Laudans invocábo Dóminum, et ab inimícis meis salvus ero.
\begin{rubric}
    Taking the Cup in his right hand and signing himself with it, he saith,
\end{rubric}\par\noindent
Sanguis Dómini nostri Jesu Christi custódiat ánimam meam in vitam ætérnam. Amen.
\begin{rubric}
    Holding the Paten under the Cup with his left hand, the Priest is to reverently receive the Blood with the particle.
\end{rubric}
\begin{rubric}
	Having received, if there be any to be communicated, let him communicate them before he purifieth himself.
\end{rubric}
\end{multicols}
\subsec{Communion of the People}
\begin{rubric}
    If there be any to be communicated, the Priest shall genuflect and place the consecrated particles in a Ciborium (or if there are few to be communicated, on the Paten), unless from the beginning they had been placed in a Ciborium.
\end{rubric}
\begin{rubric}
    If the Priest is to administer Communion from the reserved Sacrament, he openeth the Tabernacle, genuflecteth, taketh out the Ciborium, and placeth it upon the Corporal.
\end{rubric}
\begin{rubric}
    Then shall the Priest genuflect, take the Ciborium with his left hand (or the Cup, if he consecrated on the Paten), and take one particle with his right hand takes, which he holdeth between his thumb and forefinger slightly raised above the Ciborium (Cup); and turning to the people in the midst of the Altar, he saith in the clear voice:
\end{rubric}
℣. Ecce Agnus Dei, ecce, qui tollit peccáta mundi.

℟. Dómine, non sum dignus, ut intres sub tectum meum, sed tantum dic verbo, et sanábitur ánima mea.

℟. Dómine, non sum dignus, ut intres sub tectum meum, sed tantum dic verbo, et sanábitur ánima mea.

℟. Dómine, non sum dignus, ut intres sub tectum meum, sed tantum dic verbo, et sanábitur ánima mea.
\begin{rubric}
%This rubric is adapted from the Antiochian Service Book (Eastern Rite).
	After the Priest returneth to the Altar, the Pre-Communion Hymn (p. \pageref{byzantine}) is to then be sung as communicants approach the Altar Rail.\par
	\textsc{Note,} The Pre-Communion Hymn may be said or sung any time after the Prayer of Humble Access.
\end{rubric}
\begin{rubric}\label{communionrubrics}
    Only Orthodox Christians in good standing and properly disposed may receive the Eucharist.\par
    \textsc{Note,} Each communicant shall receive the Communion kneeling.
\end{rubric}
\begin{rubric}
    If they are to communicate, the Priest first communicateth the Sacred Ministers, and then the other Priests and Clerics in choir. (Priest and Deacons shall wear a stole either of white colour or of the same colour as the administering Priest.) And last of all, he proceedeth to communicate the others, beginning with those on the Epistle side.
\end{rubric}
\begin{rubric}
    If the Body and Blood of Christ are to be administered by intinction, then the Priest, when giving the Sacrament to each one, intincteth the Host into the Cup; then, making with the Host the sign of the Cross over the Cup, he placeth the Host on the tongue of each communicant, while saying:
\end{rubric}
\needspace{4\baselineskip}
\lett{C}{orpus} Domini nostri Jesu Christi, quod pro te datum est, etiam Sanguis ejus, qui pro te effusus est. Accipe eos in memoriam quod Christus mortuus est pro te, et in corde tuo, per fidem, vescere illo cum gratiarum actione.
\begin{rubric}
    But if both kinds are to be administered separately, the Priest, when giving the Body, shall make with the Host the sign of the Cross over the Ciborium and place the Host on the tongue of each communicant, while saying:
\end{rubric}
\needspace{4\baselineskip}
\lett{C}{orpus} Domini nostri Jesu Christi, quod pro te datum est, custodiat corpus et animam tuam in vitam æternam. Accipe et manduca hoc in memoriam quod Christus mortuus est pro te, et in corde tuo, per fidem, vescere illo cum gratiarum actione.
\begin{rubric}
    And then, while administering the Cup to each communicant, the Priest, keeping hold of the Cup, shall say:
\end{rubric}
\needspace{4\baselineskip}
\lett{S}{anguis} Domini nostri Jesu Christi, qui pro te effusus est, custodiat corpus et animam tuam in vitam æternam. Bibe hoc in memoriam quod Sanguis Christi effusus est pro te, et gratias age.
\begin{rubric}
    The Communion ended, the Priest shall return to the Altar, saying nothing, nor doth he give the Blessing since he will give it at the end of Mass.    
\end{rubric}
\begin{rubric}
    When the faithful have received Holy Communion, the Deacon and Priest return the Blessed Sacrament to the Tabernacle. The Priest shall then consume the Precious Blood in the Cup.\par
    \textsc{Note,} When prepared for reservation, the Host shall be touched with the Blood.
\end{rubric}
\begin{rubric}
    Then shall the Priest say,
\end{rubric}
\lett{Q}{uod} ore súmpsimus, Dómine, pura mente capiámus: et de múnere temporáli fiat nobis remédium sempitérnum.
\begin{rubric}
    Meanwhile, he presenteth the Cup to the Minister, who shall pour in a little wine, wherewith he purifieth himself. Then he continueth,    
\end{rubric}
\lett{C}{orpus} tuum, Dómine, quod sumpsi, et Sanguis, quem potávi, adhǽreat viscéribus meis: et præsta; ut in me non remáneat scélerum mácula, quem pura et sancta refecérunt sacraménta: Qui vivis et regnas in sǽcula sæculórum. Amen.
\begin{rubric}
    Then shall the Priest wash and wipe his fingers and take the ablution. Then he wipeth his mouth and the Cup.
\end{rubric}
\begin{rubric}
	After folding the Corporal, he is to cover the Cup and place it upon the Altar as before. Then he proceedeth with the Mass.
\end{rubric}
\subby{Thanksgiving}
\begin{rubric}
    Standing with his hands joined, the Priest shall read the Communion Antiphon. Then, at the Epistle corner, he alone saith the Thanksgiving.
\end{rubric}
\oremuslatin
\lett{O}{mnipotens} sempiterne Deus, tibi toto cordis affectu gratias agimus, quia nos hæc sancta Mysteria recte accipientes cibo spirituali pretiosissimi Corporis et Sanguinis Filii tui Salvatoris nostri Jesu Christi pascere dignatus es; et per hoc nos certiores facere de gratia et bonitate tua erga nos, et quod sumus vera membra corpori Filii tui mystico, fidelium omnium beatæ societati, incorporata, et hæredes secundum spem æterni regni tui, propter merita pretiosissimæ mortis et passionis dilecti Filii tui. Teque, cœlestis Pater, supplices rogamus, ut gratiæ tuæ subsidiis adjuti in sancta illa societate perseveremus, et ea omnia bona faciamus opera, quæ præparasti ut in illis ambulemus; per Jesum Christum Dominum nostrum, cui sit tecum, in unitate Spiritus Sancti, omnis honor et gloria per omnia sæcula sæculorum.\par
℟. Amen.

\subsec{Dismissal}

\begin{rubric}
    Then, with hands joined before his breast, the Priest shall proceed to the midst of the Altar, kiss the Altar, and turn to the People, with hands extended, saying:
\end{rubric}
℣. Dóminus vobíscum.

℟. Et cum spíritu tuo.
\begin{rubric}
    Then, turning back to the Book, he saith:
\end{rubric}
℣. Orémus.
\begin{rubric}
The Priest then saith the Postcommunion Prayers in the same manner, number, and order as the Collects at the beginning of Mass.
\end{rubric}
\begin{rubric}
    After saying the last Prayer, the Priest is to return to the midst of the Altar, kissing it, and turning toward the people, saying:
\end{rubric}
℣. Dóminus vobíscum.

℟. Et cum spíritu tuo.
    \begin{rubric}
    	{On days when the \emph{Gloria in exclesis} is said,}
    \end{rubric}
    ℣. Ite, Missa est.
    
    ℟. Deo grátias.
    \begin{rubric}
    	{On days when the \emph{Gloria in excelsis} is not said,}
    \end{rubric}
    ℣. Benedicámus Dómino.
    
    ℟. Deo grátias.
\begin{rubric}
In Requiem Masses,
\end{rubric}
    ℣. Requiéscant in pace.
    
    ℟. Amen.
\begin{rubric}
    In Eastertide, in Masses of the Season,
\end{rubric}
    ℣. Ite, Missa est, allelúja, allelúja.
    
    ℟. Deo grátias, allelúja, allelúja.
\begin{rubric}
    Having said the Dismissal, the Priest boweth before the midst of the Altar, and with hands joined thereon, saith in the secret voice: \emph{Pláceat tibi}.
\end{rubric}
\begin{rubric}
    Then---unless it be a Requiem Mass---the Congregation kneeling, the Priest shall kiss the Altar; raise his eyes; extend, raise, and join his hands; bow his head to the Cross; and give the Blessing, saying,
\end{rubric}
\lett{P}{ax} Dei quæ exsuperat omnem sensum, custodiat corda vestra et intelligentias vestras in scientia et amore Dei, et Filii ejus Jesu Christi Domini nostri: Et benedictio Dei Omnipotentis,\margrub{He turns to the people, blessing them once only, even in solemn Masses.} Patris, {\ding{64}} Filii, et Spiritus Sancti sit super vos, et maneat semper vobiscum. \textit{Amen.}
\begin{rubric}
    In Pontifical Masses, the Blessing is given by the Bishop and is threefold, as ordered in the Pontifical.
\end{rubric}
\begin{rubric}
    In Masses of the Dead, the Blessing is not given, but having said \emph{Requiéscant in pace} and \emph{Pláceat tibi}, he kisseth the Altar and readeth the Last Gospel.
\end{rubric}
\subsec{Last Gospel}
\begin{rubric}
    Unless a proper Last Gospel be provided, the Priest shall conclude the Mass with this Last Gospel, at the Gospel corner and with hands joined:
\end{rubric}

℣. Dóminus vobíscum.

℟. Et cum spíritu tuo.

℣. Inítium\margrub{He signs with the sign of the Cross first the Altar or the Book, then himself on the forehead, mouth, \& breast.} {\ding{66}} sancti Evangélii secúndum Joánnem

℟. Glória tibi, Dómine.

In princípio erat Verbum, et Verbum erat apud Deum, et Deus erat Verbum. Hoc erat in princípio apud Deum. Omnia per ipsum facta sunt: et sine ipso factum est nihil, quod factum est: in ipso vita erat, et vita erat lux hóminum: et lux in ténebris lucet, et ténebræ eam non comprehendérunt.
Fuit homo missus a Deo, cui nomen erat Joánnes. Hic venit in testimónium, ut testimónium perhibéret de lúmine, ut omnes créderent per illum. Non erat ille lux, sed ut testimónium perhibéret de lúmine.
Erat lux vera, quæ illúminat omnem hóminem veniéntem in hunc mundum. In mundo erat, et mundus per ipsum factus est, et mundus eum non cognóvit. In própria venit, et sui eum non recepérunt. Quotquot autem recepérunt eum, dedit eis potestátem fílios Dei fíeri, his, qui credunt in nómine ejus: qui non ex sanguínibus, neque ex voluntáte carnis, neque ex voluntáte viri, sed ex Deo nati sunt. \inrub{Everyone genuflects.} Et Verbum caro factum est, \inrub{Everyone rises.} et habitávit in nobis: et vídimus glóriam ejus, glóriam quasi Unigéniti a Patre, plenum gráti{\ae} et veritátis.\par
℟. Deo grátias.
\begin{rubric}
    At High Mass, a hymn may be sung while the Priest and servers leave the sanctuary.
\end{rubric}
\begin{rubric}
    As he departeth from the Altar, the Priest saith for thanksgiving the Antiphon \emph{Let us sing}, with the rest (p. \pageref{CommunionThanksgiving}).
\end{rubric}
\subsec{General Rubrics}\label{GeneralRubrics}
%The Book of Common Prayer has always contained rubrics at the beginning and end of the liturgical text in order for the good order of the Church. Due to the unique situation of the Vicariate, it is fitting that these be accommodated to our use.
%\begin{rubric}
%Notwithstanding anything that is elsewhere enjoined in any Rubrick or Canon, the Priest, in celebrating the Holy Communion, shall wear the vestments proper to the Priest as of 1950 in dissident Rome. That is, Amice, Alb, Cincture, Maniple, Stole, and Chasuble (the latter three of the colour of the day).
%\end{rubric}
\begin{rubric}
    It is expedient for the Priest to have, set up on the Altar, cards with the most common prayers (or at least a convenient pamphlet), including the those only mentioned by their incipit in this Book.\par
\end{rubric}
%\begin{rubric}
%The Priest, the presence of an Ordinary notwithstanding, when he processes or reposes, wears a Canterbury Cap or Biretta.
%\end{rubric}
%\begin{rubric}
%There shall be no celebration of the Lord's Supper without at least one person, other than the Priest, physically present.
%\end{rubric}
\begin{rubric}
The Altar, at the Communion-time---having a fair white linen cloth, a cross, and at least two candles upon it---shall stand in the Sanctuary.
\end{rubric}
\begin{rubric}
The Bread for the Eucharist ought to be provided by the Priest or the members of the Congregation. It must be leavened. That is, before and after being initially baked, it must be made of---and only made of---wheat flour, water, yeast, and salt. The Wine should be made of grapes, without any additives or supplements. The Cup for the Eucharist should be made of precious metal.
\end{rubric}
\begin{rubric}
It is an ancient and laudable rule of the Church to receive this Holy Sacrament fasting. That is, one who desireth to receive the Eucharist must consume neither food nor drink (water excepted) from the midnight before the liturgy until the time of reception. In the case of Masses which occur after noon, it is sufficient that the fasting begin at noon or immediately after one's noon-time meal.
\end{rubric}

%\begin{center}
%	{\textsc{Here Endeth the Liturgy of St. Tikhon.}}
%\end{center}
%\clearpage
%\subby{Exhortations for Holy Communion}
%\begin{rubric}
%	At the time of the celebration of the Communion, the Priest may say this Exhortation. And Note, that the Exhortation shall be said on the First Sunday in Advent, the First Sunday in Lent, and Trinity Sunday.
%\end{rubric}
%\lett{D}{early} beloved in the Lord, ye who mind to come to the holy Communion of the Body and Blood of our Saviour Christ, must consider how Saint Paul exhorteth all persons diligently to try and examine themselves, before they presume to eat of that Bread, and drink of that Cup. For as the benefit is great, if with a true penitent heart and lively faith we receive that holy Sacrament; 
%Addition:
%that is, the very and true Body and Blood of Our Lord Jesus Christ; 
%
%so is the danger great, if we receive the same unworthily. Judge therefore yourselves, brethren, that ye be not judged of the Lord; repent you truly for your sins past; 
%Addition:
%confess them to your spiritual father; 
%
%have a lively and steadfast faith in Christ our Saviour; amend your lives, and be in perfect charity with all men; so shall ye be meet partakers of those holy mysteries. And above all things ye must give most humble and hearty thanks to God, the Father, the Son, and the Holy Ghost, for the redemption of the world by the death and passion of our Saviour Christ, both God and man; who did humble himself, even to the death upon the Cross, for us, miserable sinners, who lay in darkness and the shadow of death; that he might make us the children of God, and exalt us to everlasting life. And to the end that we should always remember the exceeding great love of our Master, and only Saviour, Jesus Christ, thus dying for us, and the innumerable benefits which by his precious blood-shedding he hath obtained for us; he hath instituted and ordained holy mysteries, as pledges of his love, and for a continual remembrance of his death, to our great and endless comfort. To him therefore, with the Father and the Holy Ghost, let us give (as we are most bounden) continual thanks; submitting ourselves wholly to his holy will and pleasure, and studying to serve him in true holiness and righteousness all the days of our life. Amen.
%\begin{rubric}
%	When the Minister giveth warning for the Celebration of the Holy Communion, (which he shall always do upon the Sunday, or some Holy-day, immediately preceding,) he shall read this Exhortation following, or so much thereof as, in his discretion, he may think convenient.
%\end{rubric}
%\lett{D}{early} beloved, on \textit{N.} day next I purpose, through God's assistance, to administer to all such as shall be religiously and devoutly disposed the most comfortable Sacrament of the Body and Blood of Christ; to be by them received in remembrance of his meritorious Cross and Pas-sion; whereby alone we obtain remission of our sins, and are made partakers of the Kingdom of heaven. Wherefore it is our duty to render most humble and hearty thanks to Almighty God, our heavenly Father, for that he hath given his Son our Saviour Jesus Christ, not only to die for us, but also to be our spiritual food and sustenance in that holy Sacrament. Which being so divine and comfortable a thing to them who receive it worthily, and so dangerous to those who will presume to receive it unworthily; my duty is to exhort you, in the mean season to consider the dignity of that holy mystery, and the great peril of the unworthy receiving thereof; and so to search and examine your own consciences, and that not lightly, and after the manner of dissemblers with God; but so that ye may come holy and clean to such a heavenly Feast, in the marriage-garment required by God in holy Scripture, and be received as worthy partakers of that holy Table.
%
%    The way and means thereto is: First, to examine your lives and conversations by the rule of God's commandments; and whereinsoever ye shall perceive yourselves to have offended, either by will, word, or deed, there to bewail your own sinfulness, and to confess yourselves to Almighty God, with full purpose of amendment of life. And if ye shall perceive your offences to be such as are not only against God, but also against your neighbours; then ye shall reconcile yourselves unto them; being ready to make restitution and satisfaction, according to the uttermost of your powers, for all injuries and wrongs done by you to any other; and being likewise ready to forgive others who have offended you, as ye would have forgiveness of your offences at God's hand: for otherwise the receiving of the holy Communion doth nothing else but increase your condemnation. Therefore, if any of you be a blasphemer of God, an hinderer or slanderer of his Word, an adulterer, or be in malice, or envy, or in any other grievous crime; repent you of your sins, or else come not to that holy Table.
%    
%    And because it is requisite that no man should come to the holy Communion, but with a full trust in God's mercy, and with a quiet conscience; 
    %This phrase is omitted so to indicate that all such sinners ought to confess:
    %therefore, if there be any of you, who by this means cannot quiet his own conscience herein, but requireth further comfort or counsel, 
%    let him come to me, or to some other Priest of 
    %Change:
%    Christ's Church, confessing his sins and opening his grief; that he may receive such godly counsel, advice, and absolution of his sins as may tend to the quieting of his conscience, the removing of all scruple and doubtfulness, and the sure benefit of Jesus Christ's promise: the remission of sins and bestowing of the benefits of his Passion.
%\begin{rubric}
%	Or, in case he shall see the People negligent to come to the Holy Communion, instead of the former, he may use this Exhortation.
%\end{rubric}
%\lett{D}{early} beloved brethren, on \textit{N.} I intend, by God's grace, to celebrate the 
%Change:
%Holy Sacrifice of the Mass: 
%
%unto which, in God's behalf, I bid you all who are here present; and beseech you, for the Lord Jesus Christ's sake, that ye will not refuse to come thereto, being so lovingly called and bidden by God himself. Ye know how grievous and unkind a thing it is, when a man hath prepared a rich feast, decked his table with all kind of provision, so that there lacketh nothing but the guests to sit down; and yet they who are called, without any cause, most unthankfully refuse to come. Which of you in such a case would not be moved? Who would not think a great injury and wrong done unto him? Wherefore, most dearly beloved in Christ, take ye good heed, lest ye, withdrawing yourselves from this holy Supper, provoke God's indignation against you. It is an easy matter for a man to say, I will not communicate, because I am otherwise hindered with worldly business. But such excuses are not so easily accepted and allowed before God. If any man say, I am a grievous sinner, and therefore am afraid to come: wherefore then do ye not repent and amend? 
%Addition:
%Wherefore do ye not approach myself or one of Christ's Priests for Confession and Absolution? 
%
%When God calleth you, are ye not ashamed to say ye will not come? When ye should return to God, will ye excuse yourselves, and say ye are not ready? Consider earnestly with yourselves how little such feigned excuses will avail before God. Those who refused the feast in the Gospel, because they had bought a farm, or would try their yokes of oxen, or because they were married, were not so excused, but counted unworthy of the heavenly feast. Wherefore, according to mine office, I bid you in the Name of God, I call you in Christ's behalf, I exhort you, as ye love your own salvation, that ye will be partakers of this holy Communion. And as the Son of God did vouchsafe to yield up his soul by death upon the Cross for your salvation; so it is your duty to receive the Communion in remembrance of the sacrifice of his death, as he himself hath commanded: which if ye shall neglect to do, consider with yourselves how great is your ingratitude to God, and how sore punishment hangeth over your heads for the same; when ye wilfully abstain from the Lord's Table, and separate from your brethren, who come to feed on the banquet of that most heavenly food. These things if ye earnestly consider, ye will by God’s grace return to a better mind: for the obtaining whereof we shall not cease to make our humble petitions unto Almighty God, our heavenly Father.
\cleardoublepage
\phantomsection
\addcontentsline{toc}{chapter}{Hymn Index}
\fancyhead[C]{{\LARGE Hymn Index}}
\fancyhead[RO,LE]{}
\fancyhead[RE,LO]{}

\section*{First Lines}

\small

%\begin{multicols}{2}
\begin{hangparas}{1.25em}{1}
\par\noindent
A great and mighty wonder \dotfill \#40
\par\noindent
A mighty fortress is our God \dotfill \#1
\par\noindent
\textit{A solis ortu cardine} (From lands that see the sun arise) \dotfill p. \pageref{ChristmasMattins}
\par\noindent
A thrilling voice by Jordan rings, (\textit{Vox clara ecce intonat}) \dotfill p. \pageref{FirstAdventMattins}
\par\noindent
\textit{Ad coenam Agni providi} (The Lamb's high banquet we await) \dotfill p. \pageref{EasterEvensong}
\par\noindent
\textit{Aeterna caeli gloria} (Eternal glory of the sky) \dotfill p. \pageref{FridayMattinsSummer} (Summer), p. \pageref{FridayMattinsWinter} (Winter)
\par\noindent
\textit{Aeterna Imago Altissime} (Bright Image of the heav'nly King) \dotfill p. \pageref{ChristTheKingInvitatory}
\par\noindent
\textit{Aeterne rerum Conditor} (Maker of all eternal King) \dotfill p. \pageref{SundayMattinsWinter}
\par\noindent
\textit{Aeterne Rex altissime} (Eternal Monarch, King most high) \dotfill p. \pageref{AscensionInvitatory}
\par\noindent
\textit{Aeterni Patris Unice} (Thou only Son) \dotfill p. \pageref{MaryMagdaleneMattins}
\par\noindent
\textit{Ales diei nuntius} (The winged herald of the day) \dotfill p. \pageref{TuesdayMattinsSummer} (Summer), p. \pageref{TuesdayMattinsWinter} (Winter)
\par\noindent
All creatures of our God and King \dotfill \#2
\par\noindent
All glory, laud, and honour \dotfill \#64
\par\noindent
All hail! ye infant martyr flowers (\textit{Salvete flores}) \dotfill p. \pageref{HolyInnocentsMattins}
\par\noindent
All praise to Jesus' hallowed Name \dotfill \#41
\par\noindent
All ye who seek (\textit{Quicumque Christum}) \dotfill p. \pageref{TransfigurationEvensong}
\par\noindent
Alleluia! sing to Jesus!  \dotfill \#3
\par\noindent
Almighty God whose will (\textit{Magnae Deus potentiae}) \dotfill p. \pageref{ThursdayEvensongSummer} (Summer), p. \pageref{ThursdayEvensongWinter} (Winter)
\par\noindent
Amazing grace! How sweet the sound \dotfill \#4
\par\noindent
\textit{Amor Jesu dulcissime} (O love of Jesus) \dotfill p. \pageref{TransfigurationMattins}
\par\noindent
And did those feet in ancient time \dotfill \#5
\par\noindent
Angels we have heard on high \dotfill \#42
\par\noindent
\textit{Anglorum Jam Apostolus} (Apostle of the English) \dotfill p. \pageref{GregoryEvensong}
\par\noindent
\textit{Antra Deserta} (Thou in thy childhood) \dotfill p. \pageref{JohnBaptistInvitatory}
\par\noindent
Apostle of the English (\textit{Anglorum Jam Apostolus}) \dotfill p. \pageref{GregoryEvensong}
\par\noindent
As night with darksome folds unfurled (\textit{Nox atra rerum}) \dotfill p. \pageref{ThursdayInvitatory}
\par\noindent
As with gladness, men of old \dotfill \#56
\par\noindent
At the Lamb's high feast we sing \dotfill \#69
\par\noindent
At this our solemn feast (\textit{Sacris Solemnis}) \dotfill p. \pageref{CorpusChristiInvitatory}
\par\noindent
\textit{Auctor beate saeculi} (Thou blessed fount of life and time) \dotfill p. \pageref{CompassionEvensong}
\par\noindent
\textit{Audi benigne Conditor} (O maker of the world, give ear) \dotfill p. \pageref{FirstLentEvensong}
\par\noindent
\textit{Audit tyrannus} (The moody tyrant hears aghast) \dotfill p. \pageref{HolyInnocentsInvitatory}
\par\noindent
\textit{Aurea luce} (With golden splendour) \dotfill p. \pageref{PeterEvensong}
\par\noindent
\textit{Aurora jam spargit polum} (Dawn sprinkles all the east) \dotfill p. \pageref{SaturdayMattinsSummer} (Summer), p. \pageref{SaturdayMattinsWinter} (Winter)
\par\noindent
\textit{Aurora lucis rutilat} (Light's glittering morn bedecks the sky) \dotfill p. \pageref{EasterMattins}
\par\noindent
Away in a manger \dotfill \#43
\par\noindent
Be thou my Vision \dotfill \#6
\par\noindent
\textit{Beata nobis gaudia} (Blest joys for mighty wonders wrought) \dotfill p. \pageref{WhitsundayMattins}
\par\noindent
Behold the golden dawn arise (\textit{Lux ecce surgit aurea}) \dotfill p. \pageref{ThursdayMattinsSummer} (Summer), p. \pageref{ThursdayMattinsWinter} (Winter)
\par\noindent
Blessed bride of Christ (\textit{Te beata sponsa Christi}) \dotfill p. \pageref{ScholasticaEvensong}
\par\noindent
Blest joys for mighty wonders wrought (\textit{Beata nobis gaudia}) \dotfill p. \pageref{WhitsundayMattins}
\par\noindent
Bright Image of the heav'nly King (\textit{Aeterna Imago Altissime}) \dotfill p. \pageref{ChristTheKingInvitatory}
\par\noindent
Brightest and best of the sons of the morning \dotfill \#57
\par\noindent
\textit{Caeli Deus sanctissime} (O God whose hand hath spread) \dotfill p. \pageref{WednesdayEvensongSummer} (Summer), p. \pageref{WednesdayEvensongWinter} (Winter)
\par\noindent
\textit{Caelestis aulae Nuntius} (The mysteries of God are shewn) \dotfill p. \pageref{HolyRosaryEvensong}
\par\noindent
\textit{Caelitum Joseph} (Joseph, the praise and glory of the heavens) \dotfill p. \pageref{PatronageMattins} (Patronage Mattins)
\par\noindent
\textit{Caelitum, Joseph, decus} (Joseph, whose glory) \dotfill p. \pageref{JosephInvitatory} (Joseph Invitatory)
\par\noindent
\textit{Caelo Redemptor Praetulit} (In heaven's court) \dotfill p. \pageref{MotherhoodInvitatory}
\par\noindent
Christ, the fair glory (\textit{Christe sanctorum}) \dotfill p. \pageref{MichaelMattins}
\par\noindent
Christ the Lord is risen today, alleluia \dotfill \#70
\par\noindent
\textit{Christe, Redemptor omnium} (Jesu, the Father's only Son) \dotfill p. \pageref{ChristmasEvensong} (Christmas)
\par\noindent
\textit{Christe, Redemptor omnium} (O Christ Redeemer of us all) \dotfill p. \pageref{AllHallowsEvensong} (All Hallows')
\par\noindent
\textit{Christe sanctorum} (Christ, the fair glory) \dotfill p. \pageref{MichaelMattins}
\par\noindent
Come, Holy Ghost, Creator blest (\textit{Veni, Creator Spiritus}) \dotfill p. \pageref{WhitsundayEvensong}
\par\noindent
Come thou Fount of every blessing \dotfill \#7
\par\noindent
Come, thou long expected Jesus \dotfill \#36
\par\noindent
\textit{Conditor alme siderum} (Creator of the stars of night) \dotfill p. \pageref{FirstAdventEvensong}
\par\noindent
\textit{Consors paterni} (Thou Light of light and Day) \dotfill p. \pageref{TuesdayInvitatory}
\par\noindent
\textit{Cor arca legem continens} (O Heart of Jesus, holy ark) \dotfill p. \pageref{CompassionMattins}
\par\noindent
Creator of the circling sky (\textit{Orbis partrator optime}) \dotfill p. \pageref{GuardianAngelsMattins}
\par\noindent
Creator of the stars of night (\textit{Conditor alme siderum}) \dotfill p. \pageref{FirstAdventEvensong}
\par\noindent
Crown him with many crowns \dotfill \#8
\par\noindent
\textit{Custodes hominum} (The Guardians of our race) \dotfill p. \pageref{GuardianAngelsEvensong}
\par\noindent
Dawn sprinkles all the east (\textit{Aurora jam spargit polum}) \dotfill p. \pageref{SaturdayMattinsSummer} (Summer), p. \pageref{SaturdayMattinsWinter} (Winter)
\par\noindent
Dear Christians, one and all rejoice \dotfill \#9
\par\noindent
Defender, leader true (\textit{Qui te, posthabitis}) \dotfill p. \pageref{MaurusEvensong}
\par\noindent
\textit{Deus creator omnium} (O God Creator of the world) \dotfill p. \pageref{WeekdayInvitatory}
\par\noindent
\textit{Deus tuorum militum} (Of all thy warrior Saints, O Lord) \dotfill p. \pageref{StephenEvensong}
\par\noindent
\textit{Doctor egregie, Paule} (O by thy doctrine) \dotfill p. \pageref{PaulEvensong}
\par\noindent
Drop down, ye heavens, from above (\textit{Rorate Caeli}) \dotfill p. \pageref{AdventProse}
\par\noindent
\textit{Ecce jam noctis} (Lo! the dim shadows) \dotfill p. \pageref{SundayMattinsSummer}
\par\noindent
Earth's mighty Maker whose command (\textit{Telluris ingens Conditor}) \dotfill p. \pageref{TuesdayEvensongSummer} (Summer), p. \pageref{TuesdayEvensongWinter} (Winter)
\par\noindent
\textit{En ut superba} (See how the bold and raging crows) \dotfill p. \pageref{CompassionInvitatory}
\par\noindent
Eternal glory of the sky (\textit{Aeterna caeli gloria}) \dotfill p. \pageref{FridayMattinsSummer} (Summer), p. \pageref{FridayMattinsWinter} (Winter)
\par\noindent
Eternal Monarch, King most high (\textit{Aeterne Rex altissime}) \dotfill p. \pageref{AscensionInvitatory}
\par\noindent
\textit{Ex more docti} (The fast, as taught by holy lore) \dotfill p. \pageref{FirstLentInvitatory}
\par\noindent
\textit{Exsultet caelum gaudibus} (Let heaven's exultant praises ring) \dotfill p. \pageref{JohnMattins}
\par\noindent
Faith of our fathers \dotfill \#10
\par\noindent
Father we praise thee now (\textit{Nocte surgentes}) \dotfill p. \pageref{SundayInvitatorySummer}
\par\noindent
\textit{Festivis Resonent} (With glad and joyous strains) \dotfill p. \pageref{PreciousBloodEvensong}
\par\noindent
From lands that see the sun arise (\textit{A solis ortu cardine}) \dotfill p. \pageref{ChristmasMattins}
\par\noindent
\textit{Gaude, mater Anna} (Mother Anne be joyful) \dotfill p. \pageref{AnneEvensong}
\par\noindent
Gem of the highest diadem immortal (\textit{Inter aeternas}) \dotfill p. \pageref{BenedictMattins}
\par\noindent
God the Father, be our Stay \dotfill \#79
\par\noindent
God rest ye merry gentlemen \dotfill \#44
\par\noindent
Good King Wenceslas looked out \dotfill \#45
\par\noindent
Hail, holy wounds (\textit{Salvete, Christe vulnera}) \dotfill p. \pageref{PreciousBloodMattins}
\par\noindent
Hail Mary, ever blessed \dotfill p. \pageref{WalsinghamEvensong}
\par\noindent
Hail thee, festival day! \dotfill \#71 (Easter), \#73 (Ascension), \#76 (Whitsun)
\par\noindent
Hark! The herald angels sing \dotfill \#46
\par\noindent
He, whom the faithful (\textit{Iste Quem Laeti}) \dotfill p. \pageref{JosephMattins}
\par\noindent
Hence night and clouds (\textit{Nox et tenebrae et nubila}) \dotfill p. \pageref{WednesdayMattinsSummer} (Summer), p. \pageref{WednesdayMattinsWinter} (Winter)
\par\noindent
High Word of God, who once didst come (\textit{Verbum supernum}) \dotfill p. \pageref{FirstAdventInvitatory}
\par\noindent
Holy God, we praise thy Name \dotfill \#80
\par\noindent
Holy, holy, holy!  \dotfill \#81
\par\noindent
\textit{Hostis Herodes} (Why, impious Herod, vainly fear) \dotfill p. \pageref{EpiphanyEvensong}
\par\noindent
\textit{Hymnis angelicis ora} (Now we our voices raise) \dotfill p. \pageref{ScholasticaInvitatory}
\par\noindent
I bind unto myself today \dotfill \#82
\par\noindent
I need thee every hour \dotfill \#11
\par\noindent
I sing the mighty power of God \dotfill \#12
\par\noindent
I vow thee, my country \dotfill \#13
\par\noindent
If God had not been on our side \dotfill \#14
\par\noindent
Immaculate Mary! \dotfill \#15
\par\noindent
\textit{Immense caeli Conditor} (O great Creator of the sky) \dotfill p. \pageref{MondayEvensongSummer} (Summer), p. \pageref{MondayEvensongWinter} (Winter)
\par\noindent
Immortal, invisible, God only wise \dotfill \#16
\par\noindent
In childhood Placidus was by his father given (\textit{Puellus Placidus, quem pater obtulit}) \dotfill p. \pageref{MaurusInvitatory}
\par\noindent
\textit{In monte Olivis} (Upon the Mount) \dotfill p. \pageref{HolyRosaryInvitatory}
\par\noindent
In the bleak wid-winter \dotfill \#47
\par\noindent
In heaven's court (\textit{Caelo Redemptor Praetulit}) \dotfill p. \pageref{MotherhoodInvitatory}
\par\noindent
\textit{Inter aeternas} (Gem of the highest diadem immortal) \dotfill p. \pageref{BenedictMattins}
\par\noindent
\textit{Ira justa Conditoris} (Righteous anger of our Maker) \dotfill p. \pageref{PreciousBloodInvitatory}
\par\noindent
Isaiah, mighty seer, in days of old \dotfill \#83
\par\noindent
\textit{Iste Quem Laeti} (He, whom the faithful) \dotfill p. \pageref{JosephMattins}
\par\noindent
It came upon the midnight clear \dotfill \#48
\par\noindent
\textit{Jam bone pastor} (Peter, good shepherd) \dotfill p. \pageref{PeterAntiochMattins} (Chair of St. Peter at Antioch Mattins), p. \pageref{PeterMattins}
\par\noindent
\textit{Jam Christe sol justitiae} (Now Christ, thou Sun of righteousness) \dotfill p. \pageref{FirstLentMattins}
\par\noindent
\textit{Jam Christus astra ascenderat} (Now Christ, ascending whence he came) \dotfill p. \pageref{WhitsundayInvitatory}
\par\noindent
\textit{Jam morte, victor, obruta} (Now hell is vanquished) \dotfill p. \pageref{HolyRosaryMattins}
\par\noindent
\textit{Jam noctis umbrae concidunt} (The shades of night) \dotfill p. \pageref{ScholasticaMattins}
\par\noindent
\textit{Jam toto subitus} (Now let the darkling eve) \dotfill p. \pageref{SevenSorrowsEvensong}
\par\noindent
\textit{Jesu, decus angelicum} (Jesu, the beauty Angels see) \dotfill p. \pageref{HolyNameMattins}
\par\noindent
\textit{Jesu dulcis memoria} (Jesu! the very thought is sweet) \dotfill p. \pageref{HolyNameEvensong}
\par\noindent
\textit{Jesu nostra redemptio} (Jesu, Redemption all divine) \dotfill p. \pageref{AscensionMattins}
\par\noindent
Jesu, Redemption all divine (\textit{Jesu nostra redemptio}) \dotfill p. \pageref{AscensionMattins}
\par\noindent
\textit{Jesu Rex admirabilis} (O Jesu, King of wondrous might) \dotfill p. \pageref{HolyNameInvitatory}
\par\noindent
\textit{Jesu, Salvator saeculi} (Jesu who cam'st the world to save) \dotfill p. \pageref{AllHallowsMattins}
\par\noindent
Jesu, the beauty Angels see (\textit{Jesu, decus angelicum}) \dotfill p. \pageref{HolyNameMattins}
\par\noindent
Jesu, the Father's only Son (\textit{Christe, Redemptor omnium}) \dotfill p. \pageref{ChristmasEvensong}
\par\noindent
Jesu! the very thought is sweet (\textit{Jesu dulcis memoria}) \dotfill p. \pageref{HolyNameEvensong}
\par\noindent
Jesu who cam'st the world to save (\textit{Jesu, Salvator saeculi}) \dotfill p. \pageref{AllHallowsMattins}
\par\noindent
Joseph, the praise and glory of the heavens (\textit{Caelitum Joseph}) \dotfill p. \pageref{PatronageMattins}
\par\noindent
Joseph, whose glory (\textit{Caelitum, Joseph, decus}) \dotfill p. \pageref{JosephInvitatory}
\par\noindent
Joy to the world \dotfill \#49
\par\noindent
Joyful, joyful, we adore thee \dotfill \#18
\par\noindent
\textit{Laudibus cives resonent canoris} (Shout all ye people) \dotfill p. \pageref{BenedictEvensong}
\par\noindent
Let all mortal flesh keep silence \dotfill \#19
\par\noindent
Let all the Saints in concert (\textit{Omnis sanctorum concio}) \dotfill p. \pageref{AnneEvensongII}
\par\noindent
Let Angels chant thy praise, pure spouse of purest Bride (\textit{Te Joseph celebrent}) \dotfill p. \pageref{PatronageEvensong}
\par\noindent
Let heaven's exultant praises ring (\textit{Exsultet caelum gaudibus}) \dotfill p. \pageref{JohnMattins}
\par\noindent
Let the Holy Spirit's Grace \dotfill \#77
\par\noindent
Lift high the cross \dotfill \#20
\par\noindent
Light's glittering morn bedecks the sky (\textit{Aurora lucis rutilat}) \dotfill p. \pageref{EasterMattins}
\par\noindent
Lo, how a Rose ever blooming \dotfill \#50
\par\noindent
Lo! the dim shadows (\textit{Ecce jam noctis}) \dotfill p. \pageref{SundayMattinsSummer}
\par\noindent
Look down, O Lord, from heaven behold \dotfill \#21
\par\noindent
Lord, keep us steadfast in thy Word \dotfill \#22
\par\noindent
Lord, who at Cana's wedding feast \dotfill \#58
\par\noindent
Lord, who throughout these forty days \dotfill \#65
\par\noindent
\textit{Lucis Creator optime} (O blest Creator of the light) \dotfill p. \pageref{SundayEvensongSummer} (Summer), p. \pageref{SundayEvensongWinter} (Winter)
\par\noindent
\textit{Lustris sex qui jam peractis} (Thirty years among us dwelling) \dotfill p. \pageref{PassionSundayMattins} (Passion Sunday), p. \pageref{InventionMattins} (Invention)
\par\noindent
\textit{Lux ecce surgit aurea} (Behold the golden dawn arise) \dotfill p. \pageref{ThursdayMattinsSummer} (Summer), p. \pageref{ThursdayMattinsWinter} (Winter)
\par\noindent
\textit{Magnae Deus potentiae} (Almighty God whose will) \dotfill p. \pageref{ThursdayEvensongSummer} (Summer), p. \pageref{ThursdayEvensongWinter} (Winter)
\par\noindent
Maker of all eternal King (\textit{Aeterne rerum Conditor}) \dotfill p. \pageref{SundayMattinsWinter}
\par\noindent
Maker of earth, to thee alone \dotfill \#60
\par\noindent
Maker of men from heav'n (\textit{Plasmator hominis, Deus}) \dotfill p. \pageref{FridayEvensongSummer} (Summer), p. \pageref{FridayEvensongWinter} (Winter)
\par\noindent
\textit{Martyr Dei} (Thou followest, Martyr of thy God) \dotfill p. \pageref{StephenMattins}
\par\noindent
\textit{Mater Alme Numinis} (Sweet mother of the Lord most high) \dotfill p. \pageref{MotherhoodMattins}
\par\noindent
\textit{Mella cor obdulcantia} (Thy lips, O Gregory) \dotfill p. \pageref{GregoryMattins}
\par\noindent
Mother Anne be joyful (\textit{Gaude, mater Anna}) \dotfill p. \pageref{AnneEvensong}
\par\noindent
\textit{Nardo maria Pistico} (The Lord's dear feet) \dotfill p. \pageref{MaryMagdaleneInvitatory}
\par\noindent
\textit{Nocte surgentes} (Father we praise thee now) \dotfill p. \pageref{SundayInvitatorySummer}
\par\noindent
Now Christ, ascending whence he came (\textit{Jam Christus astra ascenderat}) \dotfill p. \pageref{WhitsundayInvitatory}
\par\noindent
Now Christ, thou Sun of righteousness (\textit{Jam Christe sol justitiae}) \dotfill p. \pageref{FirstLentMattins}
\par\noindent
Now Christ unfurls (\textit{Vexilla Christus inclyta}) \dotfill p. \pageref{ChristTheKingMattins}
\par\noindent
Now hell is vanquished (\textit{Jam morte, victor, obruta}) \dotfill p. \pageref{HolyRosaryMattins}
\par\noindent
Now let the darkling eve (\textit{Jam toto subitus}) \dotfill p. \pageref{SevenSorrowsEvensong}
\par\noindent
Now, my tongue, the Mystery telling (\textit{Pange lingua gloriosi Corporis}) \dotfill p. \pageref{CorpusChristiEvensong}
\par\noindent
Now thank we all our God \dotfill \#23
\par\noindent
Now we our voices raise (\textit{Hymnis angelicis ora}) \dotfill p. \pageref{ScholasticaInvitatory}
\par\noindent
\textit{Nox atra rerum} (As night with darksome folds unfurled) \dotfill p. \pageref{ThursdayInvitatory}
\par\noindent
\textit{Nox et tenebrae et nubila} (Hence night and clouds) \dotfill p. \pageref{WednesdayMattinsSummer} (Summer), p. \pageref{WednesdayMattinsWinter} (Winter)
\par\noindent
O blest Creator of the light (\textit{Lucis Creator optime}) \dotfill p. \pageref{SundayEvensongSummer} (Summer), p. \pageref{SundayEvensongWinter}
\par\noindent
O by thy doctrine (\textit{Doctor egregie, Paule}) \dotfill p. \pageref{PaulEvensong}
\par\noindent
O Christ Redeemer of us all (\textit{Christe, Redemptor omnium}) \dotfill p. \pageref{AllHallowsEvensong}
\par\noindent
O Christ, thou heaven's eternal King (\textit{Rex sempiterne, Domine}) \dotfill p. \pageref{EasterInvitatory}
\par\noindent
O come, all ye faithful \dotfill \#51
\par\noindent
O come, O come Emmanuel \dotfill \#37
\par\noindent
O Father of celestial rays (\textit{Pater superni luminis}) \dotfill p. \pageref{MaryMagdaleneEvensong}
\par\noindent
O for thy spirit (\textit{Ut queant laxis}) \dotfill p. \pageref{JohnBaptistEvensong}
\par\noindent
O God Creator of the world (\textit{Deus creator omnium}) \dotfill p. \pageref{WeekdayInvitatory}
\par\noindent
O God, in whom all grace (\textit{Summae Deus clementiae}) \dotfill p. \pageref{SevenSorrowsMattins}
\par\noindent
O God of greatest clemency (\textit{Summae Deus clementiae}) \dotfill p. \pageref{TrinityInvitatory}
\par\noindent
O God of mercy passing thought (\textit{Summae Deus}) \dotfill p. \pageref{SaturdayInvitatory}
\par\noindent
O God whose hand hath spread (\textit{Caeli Deus sanctissime}) \dotfill p. \pageref{WednesdayEvensongSummer} (Summer), p. \pageref{WednesdayEvensongWinter} (Winter)
\par\noindent
O great Creator of the sky (\textit{Immense caeli Conditor}) \dotfill p. \pageref{MondayEvensongSummer} (Summer), p. \pageref{MondayEvensongWinter} (Winter)
\par\noindent
O Heart of Jesus, holy ark (\textit{Cor arca legem continens}) \dotfill p. \pageref{CompassionMattins}
\par\noindent
O Jesu, King of wondrous might (\textit{Jesu Rex admirabilis}) \dotfill p. \pageref{HolyNameInvitatory}
\par\noindent
O Jesus Christ, from thee began \dotfill \#66
\par\noindent
O Joseph heave'nly hosts (\textit{Te Joseph celebrent}) \dotfill p. \pageref{JosephEvensong}
\par\noindent
O Lord of ages thee (\textit{Te saeculorum Principem}) \dotfill p. \pageref{ChristTheKingEvensong}
\par\noindent
O love, how deep, how broad, how high! \dotfill \# 61
\par\noindent
O love of Jesus (\textit{Amor Jesu dulcissime}) \dotfill p. \pageref{TransfigurationMattins}
\par\noindent
O maker of the world, give ear (\textit{Audi benigne Conditor}) \dotfill {p. \pageref{FirstLentEvensong}}
\par\noindent
O more than blessed (\textit{O nimis felix}) \dotfill p. \pageref{JohnBaptistMattins}
\par\noindent
O more than mighty cities known (\textit{O sola}) \dotfill p. \pageref{EpiphanyMattins}
\par\noindent
\textit{O nimis felix} (O more than blessed) \dotfill p. \pageref{JohnBaptistMattins}
\par\noindent
\textit{O quot undis lacrimarum} (What a flood of tears) \dotfill p. \pageref{SevenSorrowsInvitatory}
\par\noindent
O Ruler and Creator see (\textit{Rerum Creator}) \dotfill p. \pageref{WednesdayInvitatory}
\par\noindent
O sacred Head, now wounded  \dotfill \#67
\par\noindent
\textit{O sola} (O more than mighty cities known) \dotfill p. \pageref{EpiphanyMattins}
\par\noindent
O Trinity of blessed light (\textit{O lux beata Trinitas}) \dotfill p. \pageref{SaturdayEvensongSummer} (Summer), p. \pageref{SaturdayEvensongWinter} (Winter)
\par\noindent
O wondrous type! O vision fair \dotfill \#86
\par\noindent
O Word of God the Father \dotfill \#84
Of all thy warrior Saints, O Lord (\textit{Deus tuorum militum}) \dotfill p. \pageref{StephenEvensong}
\par\noindent
\textit{Omnis sanctorum concio} (Let all the Saints in concert) \dotfill p. \pageref{AnneEvensongII}
\par\noindent
On Jordan's bank the Baptist's cry \dotfill \#38
\par\noindent
\textit{Orbis partrator optime} (Creator of the circling sky) \dotfill p. \pageref{GuardianAngelsMattins}
\par\noindent
Our God, our help in ages past \dotfill \#24
\par\noindent
Out of the deep I cry to thee \dotfill \#25
\par\noindent
\textit{Pange lingua gloriosi Corporis} (Now, my tongue, the Mystery telling) \dotfill p. \pageref{CorpusChristiEvensong}
\par\noindent
\textit{Pange lingua gloriosi Proelium} (Sing, my tongue, the glorious battle) \dotfill p. \pageref{PassionSundayInvitatory} (Passion Sunday), p. \pageref{InventionInvitatory} (Invention)
\par\noindent
\textit{Pater superni luminis} (O Father of celestial rays) \dotfill p. \pageref{MaryMagdaleneEvensong}
\par\noindent
Peter, good shepherd (\textit{Jam bone pastor}) \dotfill p. \pageref{PeterAntiochMattins} (Chair of St. Peter at Antioch Mattins), p. \pageref{PeterMattins}
\par\noindent
Peter the Blessed (\textit{Petrus beatus catenarum}) \dotfill p. \pageref{LammasEvensong}
\par\noindent
\textit{Petrus beatus catenarum} (Peter the Blessed) \dotfill p. \pageref{LammasEvensong}
\par\noindent
\textit{Plasmator hominis, Deus} (Maker of men from heav'n) \dotfill p. \pageref{FridayEvensongSummer} (Summer), p. \pageref{FridayEvensongWinter} (Winter)
\par\noindent
Praise God, from whom all blessings flow \dotfill \#26
\par\noindent
Praise, my soul, the King of heaven \dotfill \#27
\par\noindent
Praise to the Holiest in the height \dotfill \#62
\par\noindent
Praise to the Lord, the Almighty \dotfill \#28
\par\noindent
\textit{Primo dierum} (This day the first of days) \dotfill p. \pageref{SundayInvitatoryWinter}
\par\noindent
\textit{Puellus Placidus, quem pater obtulit} (In childhood Placidus was by his father given) \dotfill p. \pageref{MaurusInvitatory}
\par\noindent
\textit{Qui te, posthabitis} (Defender, leader true) \dotfill p. \pageref{MaurusEvensong}
\par\noindent
\textit{Quicumque Christum} (All ye who seek) \dotfill p. \pageref{TransfigurationEvensong}
\par\noindent
\textit{Quidquid antiqui} (Whate'er in former days befell the) \dotfill p. \pageref{BenedictInvitatory}
\par\noindent
\textit{Quodcumque vinclis} (Whate'er thou bindest) \dotfill p. \pageref{LammasInvitatory} (Lammas Invitatory)
\par\noindent
\textit{Quodcumque vinclis} (Whatsoever chain) \dotfill p. \pageref{PeterAntiochEvensong} (Chair of St. Peter at Antioch Evensong)
\par\noindent
\textit{Rerum Creator} (O Ruler and Creator see) \dotfill p. \pageref{WednesdayInvitatory}
\par\noindent
\textit{Rex sempiterne, Domine} (O Christ, thou heaven's eternal King) \dotfill p. \pageref{EasterInvitatory}
\par\noindent
Righteous anger of our Maker (\textit{Ira justa Conditoris}) \dotfill p. \pageref{PreciousBloodInvitatory}
\par\noindent
\textit{Rorate Caeli} (Drop down, ye heavens, from above) \dotfill p. \pageref{AdventProse}
\par\noindent
\textit{Salvete flores} (All hail! ye infant martyr flowers) \dotfill p. \pageref{HolyInnocentsMattins}
\par\noindent
\textit{Sacris Solemnis} (At this our solemn feast) \dotfill p. \pageref{CorpusChristiInvitatory}
\par\noindent
\textit{Salvete, Christe vulnera} (Hail, holy wounds) \dotfill p. \pageref{PreciousBloodMattins}
\par\noindent
Saviour of the nations, come \dotfill \#39
\par\noindent
See how the bold and raging crows (\textit{En ut superba}) \dotfill p. \pageref{CompassionInvitatory}
\par\noindent
See, the Lord ascends in triumph \dotfill \#74
\par\noindent
Shout all ye people (\textit{Laudibus cives resonent canoris}) \dotfill p. \pageref{BenedictEvensong}
\par\noindent
Silent night, holy night \dotfill \#52
\par\noindent
Sing, my tongue, the glorious battle (\textit{Pange lingua gloriosi Proelium}) \dotfill p. \pageref{PassionSundayInvitatory} (Passion Sunday), p. \pageref{InventionInvitatory} (Invention)
\par\noindent
Sing we triumphant hymns of praise \dotfill \#75
\par\noindent
\textit{Somno Refectis} (With limbs refreshed by rest again) \dotfill p. \pageref{MondayInvitatory}
\par\noindent
Spirit of mercy, truth, and love \dotfill \#78
\par\noindent
\textit{Splendor paternae gloriae} (Thou Brightness of the Father's ray) \dotfill p. \pageref{MondayMattinsWinter}
\par\noindent
\textit{Succedit nocti lucifer} (The morning star succeeds to night) \dotfill p. \pageref{AnneMattins}
\par\noindent
\textit{Summae Deus} (O God of mercy passing thought) \dotfill p. \pageref{SaturdayInvitatory}
\par\noindent
\textit{Summae Deus clementiae} (O God, in whom all grace) \dotfill p. \pageref{SevenSorrowsMattins} (Seven Sorrows Mattins)
\par\noindent
\textit{Summae Deus clementiae} (O God of greatest clemency) \dotfill p. \pageref{TrinityInvitatory} (Trinity Sunday Invitatory)
\par\noindent
Sweet mother of the Lord most high (\textit{Mater Alme Numinis}) \dotfill p. \pageref{MotherhoodMattins}
\par\noindent
\textit{Te gestientem gaudiis} (Thy joys exultant) \dotfill p. \pageref{HolyRosaryEvensongII}
\par\noindent
\textit{Te beata sponsa Christi} (Blessed bride of Christ) \dotfill p. \pageref{ScholasticaEvensong}
\par\noindent
\textit{Te Joseph celebrent} (Let Angels chant thy praise, pure spouse of purest Bride) \dotfill p. \pageref{PatronageEvensong} (Patronage of St. Joseph Evensong)
\par\noindent
\textit{Te Joseph celebrent} (O Joseph heave'nly hosts) \dotfill p. \pageref{JosephEvensong} (St. Joseph Evensong)
\par\noindent
\textit{Te saeculorum Principem} (O Lord of ages thee) \dotfill p. \pageref{ChristTheKingEvensong}
\par\noindent
\textit{Telluris ingens Conditor} (Earth's mighty Maker whose command) \dotfill p. \pageref{TuesdayEvensongSummer} (Summer), p. \pageref{TuesdayEvensongWinter} (Winter)
\par\noindent
That men a godly life might live \dotfill \#29
\par\noindent
The Church's one foundation \dotfill \#30
\par\noindent
The fast, as taught by holy lore (\textit{Ex more docti}) \dotfill p. \pageref{FirstLentInvitatory}
\par\noindent
The first Noel \dotfill \#53
\par\noindent
The Guardians of our race (\textit{Custodes hominum}) \dotfill p. \pageref{GuardianAngelsEvensong}
\par\noindent
The King of love my shepherd is \dotfill \#31
\par\noindent
The Lamb's high banquet we await (\textit{Ad coenam Agni providi}) \dotfill p. \pageref{EasterEvensong}
\par\noindent
The Lord's dear feet (\textit{Nardo maria Pistico}) \dotfill p. \pageref{MaryMagdaleneInvitatory}
\par\noindent
The moody tyrant hears aghast (\textit{Audit tyrannus}) \dotfill p. \pageref{HolyInnocentsInvitatory}
\par\noindent
The morning star succeeds to night (\textit{Succedit nocti lucifer}) \dotfill p. \pageref{AnneMattins}
\par\noindent
The mouth of fools doth God confess \dotfill \#32
\par\noindent
The mysteries of God are shewn (\textit{Caelestis aulae Nuntius}) \dotfill p. \pageref{HolyRosaryEvensong}
\par\noindent
The royal banners forward go (\textit{Vexilla Regis}) \dotfill p. \pageref{PassionSundayEvensong} (Passion Sunday), p. \pageref{InventionEvensong} (Invention)
\par\noindent
The shades of night (\textit{Jam noctis umbrae concidunt}) \dotfill p. \pageref{ScholasticaMattins}
\par\noindent
The strife is o'er, the battle done \dotfill \#72
\par\noindent
The winged herald of the day (\textit{Ales diei nuntius}) \dotfill p. \pageref{TuesdayMattinsSummer} (Summer), p. \pageref{TuesdayMattinsWinter} (Winter)
\par\noindent
The Word proceeding from above (\textit{Verbum supernum prodiens}) \dotfill p. \pageref{CorpusChristiMattins}
\par\noindent
Thee, O Christ (\textit{Tibi Christe splendor Patris}) \dotfill p. \pageref{MichaelEvensong}
\par\noindent
There is a book who runs may read \dotfill \#63
\par\noindent
Thirty years among us dwelling (\textit{Lustris sex qui jam peractis}) \dotfill p. \pageref{PassionSundayMattins} (Passion Sunday), p. \pageref{InventionMattins} (Invention)
\par\noindent
This day the first of days (\textit{Primo dierum}) \dotfill p. \pageref{SundayInvitatoryWinter}
\par\noindent
Thou blessed fount of life and time (\textit{Auctor beate saeculi}) \dotfill p. \pageref{CompassionEvensong}
\par\noindent
Thou Brightness of the Father's ray (\textit{Splendor paternae gloriae}) \dotfill p. \pageref{MondayMattinsSummer} (Summer), p. \pageref{MondayMattinsWinter} (Winter)
\par\noindent
Thou followest, Martyr of thy God (\textit{Martyr Dei}) \dotfill p. \pageref{StephenMattins}
\par\noindent
Thou in thy childhood (\textit{Antra Deserta}) \dotfill p. \pageref{JohnBaptistInvitatory}
\par\noindent
Thou Light of light and Day (\textit{Consors paterni}) \dotfill p. \pageref{TuesdayInvitatory}
\par\noindent
Thou only Son (\textit{Aeterni Patris Unice}) \dotfill p. \pageref{MaryMagdaleneMattins}
\par\noindent
Thou Trinity in Unity (\textit{Tu Trinitatis}) \dotfill p. \pageref{FridayInvitatory} (Friday Invitatory)
\par\noindent
Thou Trinity in Unity (\textit{Tu, Trinitatis Unitas}) \dotfill p. \pageref{TrinityMattins} (Trinity Mattins)
\par\noindent
Though in the midst of life we be \dotfill \#33
\par\noindent
Thy joys exultant (\textit{Te gestientem gaudiis}) \dotfill p. \pageref{HolyRosaryEvensongII}
\par\noindent
Thy lips, O Gregory (\textit{Mella cor obdulcantia}) \dotfill p. \pageref{GregoryMattins}
\par\noindent
\textit{Tibi Christe splendor Patris} (Thee, O Christ) \dotfill p. \pageref{MichaelEvensong}
\par\noindent
Tis good, Lord, to be here \dotfill \#87
\par\noindent
To avert from men God's wrath \dotfill \#34
\par\noindent
To shepherds as they watched by night \dotfill \#54
\par\noindent
\textit{Tu Trinitatis} (Thou Trinity in Unity) \dotfill p. \pageref{FridayInvitatory} (Friday Invitatory)
\par\noindent
\textit{Tu, Trinitatis Unitas} (Thou Trinity in Unity) \dotfill p. \pageref{TrinityMattins}
\par\noindent
Upon the Mount (\textit{In monte Olivis}) \dotfill p. \pageref{HolyRosaryInvitatory}
\par\noindent
\textit{Ut queant laxis} (O for thy spirit) \dotfill p. \pageref{JohnBaptistEvensong}
\par\noindent
\textit{Veni, Creator Spiritus} (Come, Holy Ghost, Creator blest) \dotfill p. \pageref{WhitsundayEvensong}
\par\noindent
\textit{Verbum supernum} (High Word of God, who once didst come) \dotfill p. \pageref{FirstAdventInvitatory}
\par\noindent
\textit{Verbum supernum prodiens} (The Word proceeding from above) \dotfill p. \pageref{CorpusChristiMattins}
\par\noindent
\textit{Vexilla Christus inclyta} (Now Christ unfurls) \dotfill p. \pageref{ChristTheKingMattins}
\par\noindent
\textit{Vexilla Regis} (The royal banners forward go) \dotfill p. \pageref{PassionSundayEvensong} (Passion Sunday), p. \pageref{InventionEvensong} (Invention)
\par\noindent
\textit{Vox clara ecce intonat} (A thrilling voice by Jordan rings) \dotfill p. \pageref{FirstAdventMattins}
\par\noindent
We all believe in one true God \dotfill \#85
\par\noindent
We three Kings of Orient are \dotfill \#59
\par\noindent
What a flood of tears (\textit{O quot undis lacrimarum}) \dotfill p. \pageref{SevenSorrowsInvitatory}
\par\noindent
What Child is this \dotfill \#55
\par\noindent
Whate'er in former days befell the (\textit{Quidquid antiqui}) \dotfill p. \pageref{BenedictInvitatory}
\par\noindent
Whate'er thou bindest (\textit{Quodcumque vinclis}) \dotfill p. \pageref{LammasInvitatory}
\par\noindent
Whatsoever chain (\textit{Quodcumque vinclis}) \dotfill p. \pageref{PeterAntiochEvensong}
\par\noindent
When I survey the wondrous cross \dotfill \#68
\par\noindent
When peace like a river \dotfill \#17
\par\noindent
With glad and joyous strains (\textit{Festivis Resonent}) \dotfill p. \pageref{PreciousBloodEvensong}
\par\noindent
With golden splendour (\textit{Aurea luce}) \dotfill p. \pageref{PeterEvensong} \par\noindent
With limbs refreshed by rest again (\textit{Somno Refectis}) \dotfill p. \pageref{MondayInvitatory}
\par\noindent
When Christ was born in Bethlehem \dotfill p. \pageref{WalsinghamMattins}
\par\noindent
Why, impious Herod, vainly fear (\textit{Hostis Herodes}) \dotfill p. \pageref{EpiphanyEvensong}
\par\noindent
Ye watchers and ye holy ones \dotfill \#35
\end{hangparas}
%\end{multicols}

\normalsize


%\subby{Tsar St. Nicholas of Russia}
\feastday{{St. Nicholas Russia}}
\fancyhead[RE,LO]{17 July}
\begin{inhead}
    {Simple\\
17 July}
\end{inhead}

\textit{Opening Sentence.} To the Lord our God belong mercies and forgivenesses, though we have rebelled against him: neither have we obeyed the voice of the Lord our God, to walk in his laws which he set before us.\vr{Dan 9:9-10}

\textit{Opening Sentence.} Correct us, O Lord, but with judgement: not in thine anger, lest thou bring us to nothing.\vr{Jer 10:24}

\textit{Opening Sentence.} Enter not into judgement with thy servants, O Lord: for in thy sight shall no man living be justified.\vr{Ps 143:2}

\subbysub{I Evensong}

\begin{rubric}
	Proper Psalms: 79 \& 85
\end{rubric}
\begin{rubric}
	Proper Lessons: Jeremiah 12 \& Hebrews 11:32-12:7
\end{rubric}

With thankful hearts thy glory, O King of Saints, we sing:

Shown in the saintly story Of Nich'las Martyr-King;

Who chose to die, obeying The voice of conscience dear,

Not live on earth, betraying All that he counted dear.\\

Shall, then, his mem'ry perish? His name we venerate;

The Faith he loved we cherish, His spirit emulate.

That so by Christ-like living, With charity resign'd:

Each others' faults forgiving, We may enrich mankind.\\

%Changed to match change of alone to with kin.
For long his foes assailed him, Till friends were so fallen;

And this world's weapons fail'd him, And he was left with kin.

No whit his foes relented, Successful in the strife,

But to their King presented The choice of death or life.\\

For all lives liv'd sincerely In Christ may God be blest,

To mortals thus most clearly In mortals manifest:

The Father, who forgiveth Man's failures in the strife;

The Son, in whom he liveth; The Spirit, Source of Life! Amen.\\

    ℣. Thou, O Lord, art my defender.

	℟. And the lifter up of my head.\\
	
\antiphon{Mag.}{The true glory of princes {\dag} consisteth in advancing God's glory, in the maintenance of true religion and the Church's good.}

\subbysub{Mattins}

\begin{rubric}
	Proper Psalms: 9-11
\end{rubric}
\begin{rubric}
	Proper Lessons: 2 Samuel 1 \& Matthew 27
\end{rubric}

\begin{rubric}
	Instead of the \emph{Venite} is said the following Hymn.
\end{rubric}

\lett{R}{ighteous} art thou, O Lord: and just are thy judgements.

\secondline{\textit{Thou art just, O Lord, in all that is brought upon us: for thou hast done right, but we have done wickedly.}}

\thirdline{Nevertheless, our feet were almost gone: our treadings
had well-nigh slipped.}

\textit{For why? we were grieved at the wicked: we did also see the ungodly in such prosperity.}

The people stood up, and the rulers took counsel together: against the Lord, and against his Anointed.

\textit{They cast their heads together with one consent: and were confederate against him.}

He heard the blasphemy of the multitude, and fear was on every side: while they conspired together againſt him, to take away his life.

\textit{They spoke against him with false tongues, and compassed him about with words of hatred: and fought against him without a cause.}

Yea, his own familiar friends, whom he trusted: they that eat of his bread laid great wait for him.

\textit{They rewarded him evil for good: to the great discomfort of his
soul.}

They took their counsel together, saying, God hath forsaken him: persecute him, and take him, for there is none to deliver him.

\textit{The breath of our nostrils, the Anointed of the Lord was taken in their pits: of whom we said, Under his shadow we shall be safe.}

The adversary and the enemy entered into the gates of Jerusalem: saying, When shall he die, and his name perish?

\textit{Let the sentence of guiltiness proceed against him: and now that he lieth, let him rise up no more.}

False witnesses also did rise up against him: they laid to his charge things that he knew not.

\textit{For the sins of the people, and the iniquities of the priests: they shed the blood of the just in the midst of Jerusalem.}

O my soul, come not thou into their secret; unto their assembly, mine honour, be not thou united: for in their anger they slew a man.

\textit{Even the man of thy right hand: the Son of man, whom thou hadst made so strong for thine own self.}

In the sight of the unwise he seemed to die: and his departure was taken for misery.

\textit{They fools counted his life madness, and his end to be without honour: but he is in peace.}

For though he was punished in the sight of men: yet was his hope full of immortality.

\textit{How is he numbered with the children of God: and his lot is
among the saints!}

But, O Lord God, to whom vengeance belongeth, thou God, to whom vengeance belongeth: be favourable and gracious unto Sion.

\textit{Be merciful, O Lord, unto thy people, whom thou hast redeemed: and lay not innocent blood to our charge.}

O shut not up our souls with sinners: nor our lives with the bloodthirsty.

\textit{Deliver us from blood-guiltiness, O God, thou that art the God of our
salvation: and our tongues shall sing of thy righteousness.}

For thou art the God that hast no pleasure in wickedness: neither shall any evil dwell with thee.

\textit{Thou wilt destroy them that speak leasing: the Lord abhors both the blood-thirsty and deceitful man.}

O how suddenly do they consume: perish, and come to a fearful end!

\textit{Yea, even like a dream, when one awaketh: so didst thou make their image to vanish out of the city.}

Great and marvellous are thy works, O Lord God Almighty: just and true are thy ways, O King of saints!

\textit{Righteous art thou, O Lord: and just are thy judgements.}

Glory to the Father, and to the Son: and to the Holy Ghost;

\textit{As it was in the beginning, is now, and ever shall be: world without end. Amen.}

\begin{rubric}
	The following Hymn may be said after the preceding Hymn.
\end{rubric}
\begin{multicols}{2}
Nicholas, who chose to die

Rather than the Faith deny,

Forfeiting his kingly pride

For the sake of Jesu's Bride;

Lovingly his praise we sing,

Russia's martyr, Russia's king.\\

Mirror fair of courtesy,

Flow'r of wedded chastity,

Humble follow'r day by day,

Of the Church's holy way;

Lovingly his praise we sing,

Russia's martyr, Russia's king.\\

All the way of death he trod

For the glory of his God,

And his dying dignity

Made a bright Epiphany;

Lovingly his praise we sing,

Russia's martyr, Russia's king.\\

Bless we God the Three in One,

For all faithful 'neath the sun,

For the faithful gone before,

And for those our country bore;

Chiefly him whose praise we sing:

Russia's martyr, Russia's king. Amen.\\
\end{multicols}

\begin{rubric}
	The Office Hymn is here provided.
\end{rubric}
\begin{multicols}{2}
Nicholas' praise, our martyr King,

With heart and voice come let us sing:

Who earthly shame with Christ doth bear,

The Victor's laurel now doth wear.\\

The prison's gloom, the battle sore,

Christ's martyr here with gladness bore,

And mounting from the seated shame

To his sure hope, in heav'n he came.\\

For holy Church so he withstood,

By the bullet his lifeblood flow'd:

And where that kingly seed was sown

New harvest unto Christ hath grown.\\

For Russia's cause to peril brought;

Crown, life, and freedom deem'd he not:

For Russia's Church, for Russia's realm

We pray, lest storms our ark should whelm.\\

Jesu, to thee we glory pay,

Ruling earth's kingdoms neath thy sway;

All glory, as is ever meet,

To Father and to Paraclete. Amen.\\

    ℣. Hold fast that which thou hast.

	℟. That no man take thy crown.\\
\end{multicols}

\antiphon{Ben.}{Let my sufferings {\dag} satiate the malice of mine and the Church's enemies, but let their cruelty never exceed the measure of my charity.}

\begin{rubric}
	In the end of the Litany (which shall always on this Day be used) immediately after the Collect \emph{We humbly beseech thee, O Father}, the three Collects next following are to be read.
\end{rubric}

\lett{O}{Lord,} we beseech thee mercifully hear our prayers, and spare all those who confess their sins unto thee; that they whose consciences by sin are accused, by thy merciful pardon may be absolved; through Christ our Lord. \textit{Amen.}

\lett{O}{most} mighty God, and merciful Father, who hast compassion upon all men, and hatest nothing that thou hast made: who wouldest not the death of a sinner, but that he would rather turn from his sin, and be saved: Mercifully forgive us our trespasses; receive and comfort us, who are grieved and wearied with the burden of our sins. Thy property is always to have mercy; to thee only it appertaineth to forgive sins. Spare us therefore, good Lord, spare thy people, whom thou hast redeemed; enter not into judgement with thy servants, who are vile earth and miserable sinners; but so turn thine anger from us, who meekly acknowledge our vileness, and truly repent us of our faults; and so make haste to help us in this world, that we may ever live with thee in the world to come; through Jesus Christ our Lord. \textit{Amen.}

\lett{T}{urn} thou us, O good Lord, and so shall we be turned. Be favourable, O Lord, be favourable to thy people, Who turn to thee in weeping, fasting, and praying. For thou art a merciful God, full of compassion, long-suffering, and of great pity. Thou sparest, when we deserve punishment, And in thy wrath thinkest upon mercy. Spare thy people, good Lord, spare them, And let not thine heritage be brought to confusion. Hear us, O Lord, for thy mercy is great, And after the multitude of thy mercies look upon us, Through the merits and mediation of thy blessed Son Jesus Christ our Lord. \textit{Amen.}

\introit
\lett{L}{et} us rejoice in the Lord, celebrating a festival day in honour of blessed Nicholas the Martyr: at whose passion the Angles rejoice and praise the Son of God. \textit{Ps.} Rejoice in the Lord, O ye righteous: for it becometh well the just to be thankful.

\collect
%\lett{O}{most} mighty God, terrible in thy judgements, and wonderful in thy doings toward the children of men; who in thy heavy displeasure didst suffer the life of our gracious Sovereign Emperor Nicholas the Second, to be (as on this day) taken away by the hands of cruel and bloody men: We thy sinful creatures here assembled before thee, do, in the behalf of all this Nation, which brought down this heavy judgement upon us. But, O gracious, when thou makest inquisition for blood, lay not the guilt of this innocent blood, (the shedding whereof nothing but the blood of thy Son can expiate,) lay it not to the charge of the people of this land; not let it ever be required of us, or our posterity. Be merciful, O Lord, be merciful unto thy people, whom thou hast redeemed; and be not angry with us for ever: But pardon us for thy mercy's sake. through the merits of thy Son Jesus Christ our Lord. \textit{Amen.}
%
\lett{B}{lessed} Lord, in whose sight the death of thy saints is precious; We magnify thy Name for thine abundant grace bestowed upon the martyred Sovereign, Emperor Nicholas of Russia; by which he was enabled so cheerfully to follow the steps of his blessed Master and Saviour, in a constant meek suffering of all barbarous indignities, and at the last resisting unto blood; and even then according to the same pattern, witnessing to his murderers. Let his memory, O Lord, be ever blessed among us; that we may follow the example of his courage and constancy, his meekness and patience, and great charity. And grant, by his merits and prayers, that this our land may be ever defended from all disquietude, and thy mercy glorified in the forgiveness of our sins; and all for Jesus Christ his sake, our only Mediator and Advocate. \textit{Amen.}

\readingcitation{Epistle}{1 Peter 2:13}
\lett{B}{rethren:} Submit yourselves to every ordinance of man for the Lord's sake: whether it be to the king, as supreme; Or unto governors, as unto them that are sent by him for the punishment of evildoers, and for the praise of them that do well. For so is the will of God, that with well doing ye may put to silence the ignorance of foolish men: As free, and not using your liberty for a cloke of maliciousness, but as the servants of God. Honour all men. Love the brotherhood. Fear God. Honour the king. Servants, be subject to your masters with all fear; not only to the good and gentle, but also to the froward. For this is thankworthy, if a man for conscience toward God endure grief, suffering wrongfully. For what glory is it, if, when ye be buffeted for your faults, ye shall take it patiently? but if, when ye do well, and suffer for it, ye take it patiently, this is acceptable with God. For even hereunto were ye called: because Christ also suffered for us, leaving us an example, that ye should follow his steps: Who did no sin, neither was guile found in his mouth:

\gradall{I have found David my servant: with my holy oil have I anointed him. My hand shall hold him fast; and mine arm shall strengthen him. ℣. The enemy shall not be able to do him violence: the son of wickedness shall not hurt him.}{Alleluia, alleluia. ℣. Blessed is the man that endureth temptation: for when he is tried, he shall receive a crown of life.}

\begin{rubric}
	In Septuagesimatide \& Lent, the following Tract is said instead of the Alleluia.
\end{rubric}

\tract{Thou hast given him his heart's desire: and hast not denied him the request of his lips. ℣. For thou shalt prevent him with the blessings of goodness. ℣. Thou shalt see a crown of pure gold upon his head.}

\subbysub{Sequence}
\begin{multicols}{2}
Heav'nly King, of Kings the Pastor,

Giv'r of laws, of justice master,

Ruling all by thy behest,

Unto thee to-day we render,

Praise for him, to mem'ry tender,

Nicholas, of kings the best.\\

Traitors shedding blood like water

Fill'd the land with crime and slaughter,

Law was trampl'd in the mud,

Noble churches left forsaken

And the White Rose, overtaken

By the gun, was red with blood.\\

Thus the bardic verse fulfilling

There shall be a time of killing

When the ravens shall be fed,

And a King without pollution

Midst a realm in revolution

Shall be number'd with the dead.\\

Violent men without compassion

Proudly spurned the ancient fashion

Of the sacred right divine;

From his friends by madmen riven

Was our King to slaughter driven

Stain'd with blood his Royal line.\\

Faithful son of Mother holy,

To the Church devoted solely,

He to keep her laws was fain.

He her champion ever glorious,

Was beaten still victorious,

Robb'd of life, but conqueror slain.\\

%`He nothing common did nor mean,

%Upon that memorable scene,'

`With simple and sincere faith blest,

Self-sacrificing until rest'

Against the wicked stood his face
%When on the block he laid his head;

Nor call'd the gods with vulgar spite

To vindicate his helpless right,

But went to death as to his bed.\\

Fair exchange Tsar Nich'las making
%Fair exchange King Charles was making

When, the crown immortal taking

For the earthly crown he wore,

By bullet he follow'd faster

To the realm of Christ his master,

And the Cross behind him bore.\\

Woe to pretended justiciers,

Why, like assassins, did they fear?

To efface their wicked deed?

From which men, then, were they hiding

But from God no favour finding

And their Tsar ascending high.\\

%ORIGINAL:
%Lo, the priest who shares his glory
%(Laud his name and laud his story),
%For his fellow-martyr waits,
%And the white-robed host upraising,
%Heart and voice their Saviour
%Greets him at the heavenly gates.

He by dying brought salvation

To the torn and shatter'd nation,

Life restor'd and liberty;

For the Martyr's blood was sowing

Seed from which the Church is growing

Seed of immortality.\\

Six words said before him they slew,
%Ere his death one word was spoken

%That 'Remember' was the token

He said, `You know not what you do.'

Of his coming victory.

So his blood brought life and healing,

And the Church's triumph sealing,

Never shall forgotten be. Amen.
\end{multicols}
\readingcitation{Gospel}{Matthew 21:33}
\lett{A}{t that time:} Jesus spake a parable: There was a certain householder, which planted a vineyard, and hedged it round about, and digged a winepress in it, and built a tower, and let it out to husbandmen, and went into a far country: And when the time of the fruit drew near, he sent his servants to the husbandmen, that they might receive the fruits of it. And the husbandmen took his servants, and beat one, and killed another, and stoned another. Again, he sent other servants more than the first: and they did unto them likewise. But last of all he sent unto them his son, saying, They will reverence my son. But when the husbandmen saw the son, they said among themselves, This is the heir; come, let us kill him, and let us seize on his inheritance. And they caught him, and cast him out of the vineyard, and slew him. When the lord therefore of the vineyard cometh, what will he do unto those husbandmen? They say unto him, He will miserably destroy those wicked men, and will let out his vineyard unto other husbandmen, which shall render him the fruits in their seasons.

\offertory{In the sight of the unwise he seemed to die and his departure was taken for misery; but he is in peace.}

\secret
\lett{M}{ercifully} regard, we beseech thee, almighty God, this our atoning sacrifice, and give ear to the prayer of the blessed King Nicholas thy Martyr; graciously receive it in propitiation for the sins of this thine household. Through.

\communion{The people stood up, and the rulers took counsel together; against the Lord, and against his Anointed.}

\postcommunion
\lett{M}{ay} the bread of heaven imbue us, O Lord, with the spirit of fortitude; for this was the Bread which gave blessed King Nicholas thy Martyr strength to fight victoriously for thy Church. Through.
\cleardoublepage
%\checkoddpage
%\ifoddpage \thispagestyle{empty}
%~\clearpage\fi
\fancyhead[RE,LO]{}\fancyhead[RO,LE]{}\fancyhead[C]{}\thispagestyle{empty}
   \begin{center}
   	\textsc{
   	Most Holy Trinity, Save Us.\\
   	\small{
   	St. Mary Ever-Virgin, Pray for Us.\\
   	St. John the Divine, Pray for Us.\\
   	St. Alban the Martyr, Pray for Us.\\
   	St. Augustine of Canterbury, Pray for Us.\\
    St. Gregory the Great, Pray for Us.\\
    St. Tikhon of Moscow, Pray for Us.\\
    St. John of San Francisco, Pray for Us.\\
   	All Ye Holy Angels \& Saints, Pray for Us.
   	}
   	}
   	\vfill
   	\par
   	\includegraphics[scale=.2]{logo.eps}
			\par
	$\copyright$ 2025 Some Rights Reserved (CC-BY-SA-4.0)\\
	Apologia Anglicana, LLC\\
    Boston, MA\\
	\url{www.ApologiaAnglicana.org}
   \end{center}
\end{document}
